% !TeX program = pdflatex
% Wykład (rozdziały 1--32) + ramkowany kod numeryczny.
% Zadania i audyty znajdują się w oddzielnych plikach TEX/PDF.
% Kompilacja: latexmk -pdf PDE_tom_I_wyklad.tex
\pdfobjcompresslevel=0
\documentclass[11pt,a4paper,oneside]{book}

% Wspólna preambuła skryptu „Równania różniczkowe cząstkowe”.
% Preambula włączona bezpośrednio do jednego pliku.
\usepackage[T1]{fontenc}
\usepackage[utf8]{inputenc}
\usepackage[polish]{babel}
\usepackage{lmodern}
\usepackage{amsmath,amssymb,amsthm,mathtools}
\usepackage{esint}
\usepackage{microtype}
\usepackage[a4paper,margin=2.7cm]{geometry}
\usepackage{enumitem}
\usepackage{xcolor}
\usepackage{listings}
\lstdefinestyle{pythonpde}{
 language=Python,
 basicstyle=\ttfamily\scriptsize,
 keywordstyle=\color{blue!55!black},
 commentstyle=\color{green!30!black},
 stringstyle=\color{red!55!black},
 numbers=left, numberstyle=\tiny\color{black!55},
 numbersep=9pt, frame=single, framerule=.55pt,
 rulecolor=\color{black!45}, backgroundcolor=\color{black!2},
 breaklines=true, breakatwhitespace=true, columns=fullflexible,
 keepspaces=true, showstringspaces=false, tabsize=4,
 captionpos=b, xleftmargin=8pt,xrightmargin=3pt,framexleftmargin=3pt,
 aboveskip=10pt,belowskip=8pt}
\usepackage{graphicx}
\usepackage{tikz}
\usetikzlibrary{arrows.meta,calc,positioning}
\usepackage[unicode,colorlinks=true,linkcolor=blue!45!black,urlcolor=blue!45!black]{hyperref}

\setlength{\parindent}{0pt}
\setlength{\parskip}{4pt}
\setlength{\emergencystretch}{2em}
\setlist[enumerate]{topsep=4pt,itemsep=2pt,leftmargin=*}
\setlist[itemize]{topsep=4pt,itemsep=2pt,leftmargin=*}
\numberwithin{equation}{chapter}
\renewcommand{\lstlistingname}{Kod}

\theoremstyle{plain}
\newtheorem{twierdzenie}{Twierdzenie}[chapter]
\newtheorem{stwierdzenie}[twierdzenie]{Stwierdzenie}
\newtheorem{lemat}[twierdzenie]{Lemat}
\newtheorem{wniosek}[twierdzenie]{Wniosek}
\theoremstyle{definition}
\newtheorem{definicja}[twierdzenie]{Definicja}
\newtheorem{przyklad}[twierdzenie]{Przykład}
\newtheorem{kontrprzyklad}[twierdzenie]{Kontrprzykład}
\theoremstyle{remark}
\newtheorem{uwaga}[twierdzenie]{Uwaga}

\newcommand{\R}{\mathbb{R}}
\newcommand{\C}{\mathbb{C}}
\newcommand{\N}{\mathbb{N}}
\newcommand{\T}{\mathbb{T}}
\newcommand{\eps}{\varepsilon}
\newcommand{\pd}{\partial}
\DeclareMathOperator{\supp}{supp}
\DeclareMathOperator{\tr}{tr}
\DeclareMathOperator{\diver}{div}

\title{Równania różniczkowe cząstkowe\\\large Tom I. Podstawy teorii PDE i rozwiązań słabych}
\author{}
\date{19 listopada 2026}
\begin{document}
\frontmatter
\hypersetup{pageanchor=false}
\maketitle
\hypersetup{pageanchor=true}
\tableofcontents
\chapter*{Od redakcji}
\addcontentsline{toc}{chapter}{Od redakcji}
Tom I zawiera ciągły wykład z teorii PDE: od podstaw i metody
charakterystyk, przez równania Laplace'a i Poissona oraz
równanie ciepła, po równanie falowe, metody Fouriera, teorię dystrybucji, przestrzenie Sobolewa, metody wariacyjne, teorię spektralną i eliptyczne zagadnienia brzegowe. Zachowano wspólną
numerację i odsyłacze. Zbiór zadań do rozdziałów 1--32
oraz szczegółowy audyt matematyczno-redakcyjny są
przygotowane jako osobne dokumenty, niezależne od wykładu.
W tekście przez \emph{obszar} rozumiemy niepusty otwarty
spójny podzbiór przestrzeni euklidesowej. Określenie
\emph{dziedzina} rezerwujemy dla dziedziny funkcji lub operatora;
angielskie \emph{field} w algebrze oznacza \emph{ciało},
a \emph{vector field} w geometrii — \emph{pole wektorowe}.
W nazwach matematycznych stosujemy poprawną odmianę nazwisk:
równanie Hamiltona--Jacobiego, nierówność Harnacka,
twierdzenie Laxa--Milgrama, oszacowanie Schaudera,
metoda Galerkina. Używamy \emph{iloczynu tensorowego},
a nie sztucznej nazwy odczasownikowej.
Każda całka oznaczona zawiera obszar całkowania i różniczkę
miary, na przykład $\int_\Omega f(x)\,\mathrm d\lambda_d(x)$;
na torusie $\T^d$ używamy miary Lebesgue'a o ustalonej
normalizacji, określanej przy pierwszym użyciu.
Kod numeryczny zebrano na końcu wykładu w obramowanych
listingach; ich kompletne, uruchamialne odpowiedniki (29 plików)
znajdują się w katalogu \texttt{scripts/listings/} archiwum wydania.
W kodzie do wykresów korzystamy z Plotly, a interaktywne grafiki
zapisujemy jako pliki HTML. Instrukcja uruchomienia i lista
zależności znajdują się w pliku \texttt{INSTRUKCJA\_URUCHOMIENIA.md}.
Mapa rzeczywistych odsyłaczy z bieżącego wydania jest w pliku
\texttt{MAPA\_ZALEZNOSCI\_TOM\_I.md}.

\subsection*{Jak czytać rozdziały w porządku logicznym}
Pierwsze przykłady PDE i charakterystyki można czytać niezależnie,
ale zaawansowana część teorii ma świadomy porządek zależności.
Szeregi Fouriera na torusie przedstawiamy przed teorią dystrybucji;
następnie wprowadzamy dystrybucje, transformatę Fouriera na
$\R^d$, przestrzenie Sobolewa i ich zanurzenia. Dopiero po nich
budujemy metody wariacyjne, operatory nieograniczone oraz
słabą teorię eliptyczną i ewolucyjną.

Nie każde odwołanie do dalszego twierdzenia oznacza cykl w dowodzie:
czasem jest to jawna zapowiedź późniejszego, niezależnie
udowodnionego wyniku. W szczególności wykład o funkcjach
harmonicznych odwołuje się do ich analityczności dowiedzionej
później w tym samym rozdziale. Rozdziały o operatorach
pseudoróżniczkowych korzystają z nieunitarnej konwencji
transformaty Fouriera już zdefiniowanej w części podstawowej.
Szczegółową mapę odsyłaczy i zależności zawiera dokument
\emph{MAPA\_ZALEZNOSCI\_TOM\_I.md} dołączony do archiwum.
\mainmatter

% Rozdział 1.
\chapter[Wprowadzenie do PDE]{Równania różniczkowe cząstkowe: pojęcia i pierwsze zjawiska}
\label{chap:pde-intro}

Równania różniczkowe cząstkowe (ang. \emph{partial differential equations}, PDE)
opisują związki między pochodnymi funkcji wielu zmiennych. Przy ich pomocy
formułuje się prawa przepływu ciepła, rozchodzenia się fal, ruchu płynów,
dynamiki kwantowej i geometrii rozmaitości. Ten sam operator różniczkowy
może prowadzić do zupełnie różnych zagadnień w zależności od tego, na
jakim obszarze pracujemy, jakie narzucamy warunki i w jakiej klasie
szukamy rozwiązań.

Celem rozdziału jest zbudowanie języka potrzebnego w dalszej części
skryptu. Wprowadzimy rząd i typ nieliniowości równania, symbol główny,
powierzchnie charakterystyczne i podstawowe klasy równań. Następnie,
na przykładach, oddzielimy trzy problemy: istnienie rozwiązania,
jego jednoznaczność i stabilną zależność od danych.

\section{Co jest niewiadomą w równaniu PDE?}
\label{sec:pde-unknown}

Niech $n\geq1$, $\Omega\subset\R^n$ będzie niepustym zbiorem otwartym,
a $u\colon\Omega\to\R$ funkcją niewiadomą. Równanie różniczkowe
cząstkowe wiąże wartości $u$ i jej pochodnych cząstkowych. Przykładowo:
\begin{align}
 -\Delta u&=f &&\text{(równanie Poissona)},\label{eq:poisson-intro}\\
 u_t-\kappa\Delta_xu&=0,\quad \kappa>0
 &&\text{(równanie ciepła)},\label{eq:heat-intro}\\
 u_{tt}-c^2\Delta_xu&=0,\quad c>0
 &&\text{(równanie falowe)},\label{eq:wave-intro}\\
 u_t+u\,u_x&=0 &&\text{(równanie Burgersa bez lepkości)}.
 \label{eq:burgers-intro}
\end{align}
W trzech ostatnich równaniach niewiadoma zależy od czasu $t$ i
zmiennych przestrzennych $x=(x_1,\ldots,x_d)$. Zapis $\Delta_x$ oznacza
laplasjan liczony wyłącznie względem zmiennych przestrzennych, tj.
$\Delta_xu=\sum_{j=1}^d\partial_{x_j}^2u$.

\begin{uwaga}[Rozwiązanie to równanie wraz z warunkami]
Samo równanie $-\Delta u=f$ ma zwykle wiele rozwiązań. Dopiero po
określeniu obszaru, warunków brzegowych i klasy funkcji otrzymujemy
konkretne \emph{zagadnienie}. Przykładowo każda funkcja afiniczna
$u(x)=a\cdot x+b$ jest harmoniczna w $\R^n$; nie można zatem mówić
o jednoznaczności bez dalszych założeń.
\end{uwaga}

W pierwszych rozdziałach przez rozwiązanie klasyczne równania rzędu
$m$ rozumiemy funkcję klasy $C^m$ spełniającą równanie punktowo. Dla równań
ewolucyjnych wygodniej jest osobno liczyć pochodne czasowe i
przestrzenne; na przykład $C^{1,2}$ dla kolejności argumentów $(t,x)$ oznacza jedną ciągłą
pochodną czasową i dwie przestrzenne. Rozwiązania słabe i dystrybucyjne
wprowadzimy później.

\section{Rząd równania i rodzaje nieliniowości}
\label{sec:pde-types}

Dla wielowskaźnika $\alpha=(\alpha_1,\ldots,\alpha_n)\in\N_0^n$
przyjmujemy $|\alpha|=\alpha_1+\cdots+\alpha_n$ i
$\partial^\alpha=\partial_{x_1}^{\alpha_1}\cdots
\partial_{x_n}^{\alpha_n}$. Niech $m\geq1$, a $J^m u(x)$ oznacza
uporządkowaną rodzinę wszystkich wartości $\partial^\alpha u(x)$,
$|\alpha|\leq m$ (tzw. $m$-dżet funkcji w punkcie).

\begin{definicja}[Równanie rzędu $m$]
Niech $\Omega\subset\R^n$ będzie otwarta, a $F$ funkcją określoną
na odpowiednim otwartym podzbiorze przestrzeni zmiennych
$(x,(z_\alpha)_{|\alpha|\leq m})$. Równanie
$F(x,J^m u(x))=0$ nazywamy równaniem rzędu $m$, jeżeli
$F$ zależy nietrywialnie od co najmniej jednej zmiennej $z_\alpha$
o $|\alpha|=m$. Pojęcie rzędu odnosi się do przedstawionego równania.
W szczególnych punktach część najwyższego rzędu może znikać;
na przykład operator $x\,\partial_x^2+\partial_x$ ma rząd $2$
na $\R$, choć w punkcie $x=0$ jego symbol drugiego rzędu
zeruje się dla każdego kowektora.
\end{definicja}

\begin{definicja}[Liniowość, półliniowość i quasiliniowość]
Niech $\Omega\subset\R^n$ będzie otwarta i $m\geq1$.
\begin{enumerate}
\item Równanie \emph{liniowe} rzędu $m$ ma postać
$Lu=f$, gdzie $Lu=\sum_{|\alpha|\leq m}a_\alpha(x)\partial^\alpha u$,
a współczynniki $a_\alpha$ oraz $f$ nie zależą od $u$.
Jest \emph{jednorodne}, jeśli $f=0$.
\item Równanie jest \emph{półliniowe} (semiliniowe), jeśli ma postać
$\sum_{|\alpha|=m}a_\alpha(x)\partial^\alpha u
+G(x,J^{m-1}u)=0$. Pochodne najwyższego rzędu wchodzą zatem
liniowo, a ich współczynniki zależą tylko od $x$.
\item Równanie jest \emph{quasiliniowe}, jeśli można je zapisać jako
$\sum_{|\alpha|=m}a_\alpha(x,J^{m-1}u)\partial^\alpha u
+G(x,J^{m-1}u)=0$.
\item Równanie nazywamy \emph{w pełni nieliniowym}, gdy jego zależność
od pochodnych najwyższego rzędu nie jest afiniczna.
\end{enumerate}
Terminy te dotyczą przyjętego rzędu i postaci równania; równania
liniowe są szczególnym przypadkiem półliniowych, a półliniowe
szczególnym przypadkiem quasiliniowych. Klasyfikacja odnosi się do
\emph{wyrażenia różniczkowego}, nie wyłącznie do zbioru jego
rozwiązań: równania $\Delta u=0$ i $(\Delta u)^2=0$ mają ten
sam zbiór rozwiązań klasycznych, ale inne postaci algebraiczne.
\end{definicja}

\begin{przyklad}[Rozpoznawanie typu]
Dla niewiadomej $u=u(x,y)$:
\begin{enumerate}
\item $u_{xx}+u_{yy}=\sin x$ jest liniowe i niejednorodne.
\item $u_{xx}+u_{yy}+u^3=0$ jest półliniowe.
\item $(1+u_x^2)u_{xx}+u_{yy}=0$ jest quasiliniowe: współczynnik
przy $u_{xx}$ zależy od pochodnej pierwszego rzędu.
\item $u_{xx}u_{yy}-u_{xy}^2=1$ jest w pełni nieliniowe
(równanie typu Monge'a--Ampère'a).
\item $u_t+u\,u_x=0$ jest quasiliniowe pierwszego rzędu.
\end{enumerate}
\end{przyklad}

\begin{stwierdzenie}[Zasada superpozycji]
Niech $V,W$ będą przestrzeniami wektorowymi funkcji, a
$L\colon V\to W$ operatorem liniowym. Jeśli $Lu_1=f_1$ i
$Lu_2=f_2$, to dla dowolnych $a,b\in\R$ zachodzi
$L(au_1+bu_2)=af_1+bf_2$. W szczególności zbiór rozwiązań
równania jednorodnego $Lu=0$ jest podprzestrzenią wektorową
$V$, a zbiór rozwiązań $Lu=f$, jeżeli jest niepusty,
ma postać $u_0+\ker L$.
\end{stwierdzenie}
\begin{proof}
Pierwsza równość wynika bezpośrednio z liniowości $L$.
Jeżeli $Lu_0=f$, to $Lu=f$ jest równoważne
$L(u-u_0)=0$, skąd wynika ostatnie stwierdzenie.
\end{proof}

\begin{uwaga}
Dla równania nieliniowego nie wolno zazwyczaj dodawać rozwiązań.
Rozważmy na przykład równanie $w_x=w^2$ na
$\Omega=(-1,1)^2$, gdzie $w=w(x,y)$.
Funkcje $w_1(x,y)=1/(2-x)$ i $w_2(x,y)=1/(3-x)$
są rozwiązaniami, ponieważ $w_j'=w_j^2$.
Ich suma nie jest jednak rozwiązaniem, gdyż
$(w_1+w_2)'=w_1^2+w_2^2\ne(w_1+w_2)^2$.
Podobnie $2w_1$ nie jest rozwiązaniem: pokazuje to,
że nie działa również zwykłe mnożenie rozwiązań przez skalar.
\end{uwaga}

\section{Dane początkowe i brzegowe}
\label{sec:pde-data}

Niech $\Omega\subset\R^d$ będzie obszarem ograniczonym, tj.
niepustym, spójnym zbiorem otwartym, i załóżmy, że jej brzeg
$\partial\Omega$ jest co najmniej klasy $C^1$ tam, gdzie używamy
wektora normalnego $\nu$ skierowanego na zewnątrz.

\begin{definicja}[Podstawowe warunki brzegowe]
Dla dostatecznie regularnej funkcji $u$ określonej na $\overline\Omega$ rozważamy:
\begin{enumerate}
\item \emph{warunek Dirichleta}: $u|_{\partial\Omega}=g$;
\item \emph{warunek Neumanna}: $\partial_\nu u|_{\partial\Omega}=g$,
 gdzie $\partial_\nu u=\nabla u\cdot\nu$;
\item \emph{warunek Robina}: $\partial_\nu u+\beta u=g$
 na $\partial\Omega$, gdzie $\beta$ jest zadaną funkcją.
\end{enumerate}
Warunki te rozumiemy na razie punktowo, dla funkcji o odpowiedniej
regularności aż do brzegu.
\end{definicja}

Dla zagadnień ewolucyjnych na $(0,T)\times\Omega$, gdzie $T>0$,
warunek początkowy zapisujemy na przykład jako $u(0,x)=u_0(x)$.
Równanie falowe drugiego rzędu względem czasu wymaga zwykle także
warunku $u_t(0,x)=u_1(x)$. Jeśli $\Omega$ jest ograniczona,
potrzebne są ponadto odpowiednie warunki na bocznym brzegu
$(0,T)\times\partial\Omega$.

\begin{przyklad}[Zgodność danych na krawędzi czasoprzestrzeni]
Jeżeli $u\in C([0,T]\times\overline\Omega)$ spełnia
$u(0,x)=u_0(x)$ dla $x\in\overline\Omega$ oraz
$u(t,x)=g(t,x)$ dla $(t,x)\in[0,T]\times\partial\Omega$,
to koniecznie $u_0(x)=g(0,x)$ dla każdego $x\in\partial\Omega$.
Jest to najprostszy \emph{warunek zgodności} danych; wyższa
regularność często narzuca dalsze warunki.
\end{przyklad}

\begin{uwaga}[Zagadnienie Cauchy'ego]
Przez zagadnienie Cauchy'ego rozumiemy poszukiwanie rozwiązania
na podstawie danych zadanych na odpowiedniej hiperpłaszczyźnie
lub hiperpowierzchni. Nie należy utożsamiać go z zagadnieniem
Dirichleta. Na przykład zadanie wartości $u$ i pochodnej normalnej
na części brzegu dla równania Laplace'a jest zagadnieniem Cauchy'ego,
które, jak zobaczymy w~\S\ref{sec:pde-hadamard}, jest niestabilne.
\end{uwaga}

\section{Symbol główny i powierzchnie charakterystyczne}
\label{sec:pde-symbol}

Pojęcie symbolu ujawnia, co z równania pozostaje przy badaniu
oscylacji o bardzo dużej częstości. Niech $m\geq1$,
$\Omega\subset\R^n$ będzie otwarta, a
$P=\sum_{|\alpha|\leq m}a_\alpha(x)\partial^\alpha$
liniowym operatorem różniczkowym o współczynnikach ciągłych
na $\Omega$. Zakładamy, że przynajmniej jeden współczynnik
najwyższego rzędu nie znika identycznie.

\begin{definicja}[Symbol główny]
\emph{Symbolem głównym} operatora $P$ nazywamy wielomian
$p_m\colon\Omega\times\R^n\to\C$ dany wzorem
\begin{equation}\label{eq:principal-symbol}
 p_m(x,\xi)=\sum_{|\alpha|=m}a_\alpha(x)\xi^\alpha.
\end{equation}
W naszej konwencji $\partial^\alpha$ nie zawiera czynników $-i$.
Przy konwencji $D_j=-i\partial_{x_j}$ symbol różni się o
odpowiedni czynnik potęgi $i$.
\end{definicja}

\begin{stwierdzenie}[Test wysokich częstości]
Niech $P$ będzie operatorem opisanym powyżej, a
$x\in\Omega$, $\xi\in\R^n$. Dla $\lambda>0$ oraz
$v_\lambda(z)=e^{i\lambda z\cdot\xi}$ zachodzi
\[
 e^{-i\lambda x\cdot\xi}Pv_\lambda(x)
 =i^m\lambda^m p_m(x,\xi)+O(\lambda^{m-1})
 \quad\text{gdy }\lambda\to\infty.
\]
Stała w symbolu $O(\lambda^{m-1})$ może zależeć od $x$ i $\xi$.
\end{stwierdzenie}
\begin{proof}
Mamy $\partial^\alpha v_\lambda(x)=
(i\lambda)^{|\alpha|}\xi^\alpha v_\lambda(x)$.
Po podstawieniu do definicji $P$ wyrazy z $|\alpha|=m$
dają dokładnie $i^m\lambda^mp_m(x,\xi)$, a skończona
suma pozostałych wyrazów jest rzędu co najwyżej $\lambda^{m-1}$.
\end{proof}

\begin{definicja}[Eliptyczność]
Liniowy operator skalarny $P$ rzędu $m$ na zbiorze otwartym
$\Omega\subset\R^n$ nazywamy \emph{eliptycznym w punkcie}
$x\in\Omega$, jeżeli $p_m(x,\xi)\ne0$ dla każdego
$\xi\in\R^n\setminus\{0\}$. Jest \emph{eliptyczny na $\Omega$},
gdy warunek zachodzi w każdym punkcie.
\end{definicja}

\begin{przyklad}
Dla $P=-\Delta=-\sum_{j=1}^n\partial_{x_j}^2$
przyjęta konwencja daje $p_2(x,\xi)=-|\xi|^2$.
Symbol nie zeruje się poza $\xi=0$, więc laplasjan jest
operatorem eliptycznym. Znak symbolu nie ma znaczenia
dla samej definicji eliptyczności.
\end{przyklad}

\begin{uwaga}[Eliptyczność a jednostajna eliptyczność]
Dla rzeczywistej symetrycznej macierzy współczynników
$A(x)=(a_{ij}(x))_{i,j=1}^n$ warunek jednostajnej
eliptyczności ma postać
$\xi^{\mathsf T}A(x)\xi\ge\lambda|\xi|^2$
dla wszystkich $x\in\Omega$, $\xi\in\R^n$ i pewnej
\emph{jednej} stałej $\lambda>0$. Samo niezerowanie się
symbolu głównego w poszczególnych punktach nie daje
takiej stałej. Przykładowo operator
$P=-x^2\partial_x^2$ jest eliptyczny na $(0,1)$,
ale nie jest jednostajnie eliptyczny: dla $x\downarrow0$
wartość współczynnika $x^2$ dąży do zera.
Wymóg dolnego oszacowania przez $\lambda|\xi|^2$
dotyczy operatorów drugiego rzędu w postaci dodatnio
określonej; nie jest definicją wszystkich eliptycznych
operatorów dowolnego rzędu.
\end{uwaga}

\begin{definicja}[Powierzchnia charakterystyczna]
Niech $P$ będzie liniowym operatorem skalarnym rzędu $m$ na
$\Omega\subset\R^n$, a $\Sigma\subset\Omega$ gładką
hiperpowierzchnią lokalnie opisaną jako
$\Sigma=\{x\in\Omega:\phi(x)=0\}$, gdzie
$\phi\in C^1(\Omega)$ i $\mathrm d\phi\ne0$ na $\Sigma$.
Mówimy, że $\Sigma$ jest \emph{charakterystyczna dla $P$
w punkcie} $x\in\Sigma$, jeżeli
$p_m(x,\mathrm d\phi(x))=0$. W przeciwnym razie jest
\emph{niecharakterystyczna w $x$}. Warunek nie zależy od
wyboru lokalnej funkcji definiującej $\phi$.
\end{definicja}

\begin{proof}[Uzasadnienie niezależności od funkcji definiującej]
Dwie funkcje definiujące tę samą gładką hiperpowierzchnię mają
w rozważanym punkcie niezerowe i proporcjonalne różniczki:
$\mathrm d\psi(x)=a\,\mathrm d\phi(x)$ dla pewnego $a\ne0$.
Ponieważ $p_m(x,\cdot)$ jest jednorodny stopnia $m$,
$p_m(x,\mathrm d\psi(x))=a^mp_m(x,\mathrm d\phi(x))$.
Oba wyrażenia zerują się zatem jednocześnie.
\end{proof}

\begin{przyklad}[Stożek charakterystyczny równania falowego]
Dla operatora $P=\partial_t^2-c^2\Delta_x$ w $\R^{1+d}$
symbol główny jest równy
$p_2(t,x;\tau,\xi)=\tau^2-c^2|\xi|^2$.
Wektory kowariantne $(\tau,\xi)\ne0$ spełniające
$\tau^2=c^2|\xi|^2$ są charakterystyczne. W przypadku
jednego wymiaru przestrzennego proste $t-x/c=\mathrm{const}$
i $t+x/c=\mathrm{const}$ są charakterystyczne.
Natomiast powierzchnia początkowa $t=0$ jest
niecharakterystyczna, bo $p_2(1,0)=1$.
\end{przyklad}

\begin{uwaga}[Pułapka: symbol równania ciepła]
Dla $P=\partial_t-\kappa\Delta_x$, traktowanego jako zwykły
operator drugiego rzędu we wszystkich zmiennych $(t,x)$,
symbol główny wynosi $-\kappa|\xi|^2$: nie zależy od $\tau$.
W takim sensie powierzchnia $t=0$ jest charakterystyczna
($p_2((t,x),(1,0))=0$),
chociaż zadanie początkowe dla równania ciepła jest naturalne
i stabilne w odpowiednich przestrzeniach. Dlatego równania
paraboliczne klasyfikuje się przy użyciu \emph{ważonego}
porządku pochodnych: jednej pochodnej czasowej przypisuje się
wagę $2$, a przestrzennej wagę $1$. Przy konwencji symbolu
$D_j=-i\partial_{x_j}$ prowadzi to dla operatora ciepła do
symbolu parabolicznego $i\tau+\kappa|\xi|^2$, który dla
rzeczywistych $(\tau,\xi)\ne0$ nie znika. W szczególności
nie wolno bez dodatkowych zastrzeżeń stosować do równania
ciepła intuicji charakterystyk równania falowego.
\end{uwaga}

\section{Klasyfikacja operatorów drugiego rzędu}
\label{sec:pde-classification}

Dla równania drugiego rzędu w dwóch zmiennych rzeczywistych
rozważmy część główną
\begin{equation}\label{eq:2d-principal}
 A(x,y)u_{xx}+2B(x,y)u_{xy}+C(x,y)u_{yy}.
\end{equation}
Współczynniki $A,B,C$ są rzeczywiste i ciągłe, a odpowiadająca
im macierz symetryczna to
$\mathsf A(x,y)=\begin{pmatrix}A&B\\B&C\end{pmatrix}$.

\begin{definicja}[Klasyfikacja w dwóch zmiennych]
Niech $\Omega\subset\R^2$ będzie otwarta i
$\mathsf A(x,y)\ne0$ w rozważanym punkcie $(x,y)\in\Omega$.
Oznaczmy $\delta(x,y)=B(x,y)^2-A(x,y)C(x,y)$.
Równanie z częścią główną~\eqref{eq:2d-principal} ma w tym punkcie
typ \emph{eliptyczny}, gdy $\delta<0$;
\emph{hiperboliczny}, gdy $\delta>0$;
\emph{paraboliczny w tradycyjnej klasyfikacji dwuzmiennowej},
gdy $\delta=0$ i $\mathsf A\ne0$.
Jeśli $\mathsf A=0$, część drugiego rzędu zanika i tej
klasyfikacji tam nie stosujemy.
\end{definicja}

\begin{stwierdzenie}[Interpretacja geometryczna i niezmienniczość typu]
Niech $\mathsf A$ będzie rzeczywistą, symetryczną macierzą
$2\times2$, różną od zera. Wtedy:
\begin{enumerate}
\item $B^2-AC<0$ wtedy i tylko wtedy, gdy $\mathsf A$ jest
 dodatnio albo ujemnie określona;
\item $B^2-AC>0$ wtedy i tylko wtedy, gdy $\mathsf A$ ma
 wartości własne przeciwnych znaków;
\item $B^2-AC=0$ wtedy i tylko wtedy, gdy $\mathsf A$ ma rząd $1$.
\end{enumerate}
Pod lokalną zmianą współrzędnych klasy $C^2$ o nieosobliwym
jakobianie liczba dodatnich, ujemnych i zerowych wartości własnych
macierzy części głównej nie zmienia się. W konsekwencji typ
równania jest niezmienniczy względem takich zmian współrzędnych.
\end{stwierdzenie}
\begin{proof}
Macierz $\mathsf A$ ma rzeczywiste wartości własne
$\lambda_1,\lambda_2$, przy czym
$\lambda_1\lambda_2=\det\mathsf A=AC-B^2=-\delta$.
Jeśli $\delta<0$, iloczyn jest dodatni i obie niezerowe
wartości własne mają ten sam znak. Jeśli $\delta>0$,
mają znaki przeciwne. Gdy $\delta=0$, jedna wartość własna
jest zerowa, a druga niezerowa, ponieważ $\mathsf A\ne0$.

Niech $z=\Phi(x)$ będzie lokalną zmianą współrzędnych,
$J=D\Phi(x)$ jej nieosobliwym jakobianem i
$u(x)=v(\Phi(x))$. Z reguły łańcuchowej wynika, że
macierz współczynników drugich pochodnych $v$ jest równa
$J\mathsf A J^{\mathsf T}$; pochodne drugie $\Phi$
wnoszą jedynie wyrazy pierwszego rzędu. Prawo bezwładności
Sylvestera zapewnia równość liczby dodatnich, ujemnych
i zerowych wartości własnych obu macierzy.
\end{proof}

\begin{przyklad}[Trzy podstawowe typy]
\begin{enumerate}
\item $u_{xx}+u_{yy}=0$: $A=C=1$, $B=0$, więc $\delta=-1$;
 typ eliptyczny.
\item $u_{xx}-u_{yy}=0$: $A=1$, $C=-1$, $B=0$, więc
 $\delta=1$; typ hiperboliczny.
\item $u_{xx}+2u_{xy}+u_{yy}=0$: $A=B=C=1$, więc
 $\delta=0$; typ paraboliczny w klasyfikacji dwuzmiennowej.
 Część główna ma postać $(\partial_x+\partial_y)^2$.
\end{enumerate}
\end{przyklad}

\begin{przyklad}[Równanie zmieniające typ]
Rozważmy na $\R^2$ równanie Tricomiego
\begin{equation}\label{eq:tricomi}
 u_{xx}+x\,u_{yy}=0.
\end{equation}
Mamy $A=1$, $B=0$, $C=x$, zatem $\delta=-x$. Równanie jest
eliptyczne w obszarze $x>0$, hiperboliczne w obszarze $x<0$
i typu parabolicznego na prostej $x=0$.
\end{przyklad}

\begin{center}
\begin{tikzpicture}[scale=.9,>=Stealth]
 \fill[gray!10] (-4,-1.25) rectangle (0,1.25);
 \fill[gray!23] (0,-1.25) rectangle (4,1.25);
 \draw[->] (-4.3,0)--(4.45,0) node[right] {$x$};
 \draw[->] (0,-1.5)--(0,1.5);
 \draw[line width=1.1pt] (0,-1.25)--(0,1.25);
 \node[align=center] at (-2,0.45) {hiperboliczny\\$x<0$};
 \node[align=center] at (2,0.45) {eliptyczny\\$x>0$};
 \node[above,align=center] at (0,1.6) {$x=0$: typ paraboliczny};
\end{tikzpicture}
\par\smallskip\small
Typ równania Tricomiego zmienia się na prostej $x=0$.
\end{center}

\begin{uwaga}[Dwa znaczenia słowa „paraboliczny”]
Typ paraboliczny w klasyfikacji wyrażenia
$Au_{xx}+2Bu_{xy}+Cu_{yy}$ dotyczy macierzy części
drugiego rzędu w \emph{dwóch równorzędnych zmiennych}.
Równanie ciepła $u_t-\kappa\Delta_xu=0$ jest natomiast
parabolicznym \emph{równaniem ewolucyjnym} z wyróżnionym
czasem i z innym skalowaniem pochodnych. Terminy są
powiązane historycznie, ale nie należy utożsamiać tych
definicji bez wskazania konwencji.
\end{uwaga}

W większej liczbie zmiennych sama nieokreśloność macierzy
części głównej nie wystarcza do pełnej definicji
hiperboliczności. Dla równania falowego wyróżnia się
kowektor czasowy oraz wymaga odpowiednich własności
rzeczywistych pierwiastków wielomianu symbolu w kierunku
czasowym. Systematyczną teorię rozwinie rozdział o
równaniach hiperbolicznych.

\section{Pierwsze modele: co można udowodnić od razu?}
\label{sec:pde-models}

\subsection{Laplasjan: zasada maksimum i jednoznaczność}

\begin{twierdzenie}[Słaba zasada maksimum]
\label{thm:laplace-max}
Niech $\Omega\subset\R^n$ będzie niepustym, ograniczonym
zbiorem otwartym i niech $u\in C^2(\Omega)\cap
C(\overline\Omega)$ będzie funkcją rzeczywistą. Jeśli
$\Delta u\geq0$ w $\Omega$, to
\begin{equation}\label{eq:maximum-principle}
 \max_{\overline\Omega}u\leq\max_{\partial\Omega}u.
\end{equation}
Jeśli $\Delta u\leq0$, to
$\min_{\overline\Omega}u\geq\min_{\partial\Omega}u$.
W szczególności, gdy $\Delta u=0$, oba wnioski zachodzą.
\end{twierdzenie}
\begin{proof}
Dla $\eps>0$ określmy $u_\eps(x)=u(x)+\eps|x|^2$.
W $\Omega$ zachodzi
$\Delta u_\eps=\Delta u+2n\eps>0$.
Funkcja $u_\eps$ jest ciągła na zwartym zbiorze
$\overline\Omega$, zatem osiąga maksimum. Gdyby osiągała
je w punkcie wewnętrznym, jej hesjan byłby w tym punkcie
ujemnie półokreślony, stąd $\Delta u_\eps\leq0$,
co daje sprzeczność. Maksimum $u_\eps$ leży zatem na brzegu.
Wobec tego dla każdego $x\in\overline\Omega$ mamy
$u(x)+\eps|x|^2\leq
\max_{z\in\partial\Omega}(u(z)+\eps|z|^2)$.
Ponieważ $\overline\Omega$ jest ograniczony,
możemy przejść z $\eps$ do zera i otrzymać
$u(x)\leq\max_{\partial\Omega}u$.
Dla $\Delta u\leq0$ stosujemy wykazany wynik do $-u$.
\end{proof}

\begin{wniosek}[Jednoznaczność i stabilność zagadnienia Dirichleta]
\label{cor:dirichlet-stability}
Niech $\Omega\subset\R^n$ będzie niepustym ograniczonym
zbiorem otwartym, a $u,v\in C^2(\Omega)\cap
C(\overline\Omega)$ niech spełniają $\Delta u=\Delta v$
w $\Omega$. Wówczas
\[
 \max_{x\in\overline\Omega}|u(x)-v(x)|
 \leq\max_{x\in\partial\Omega}|u(x)-v(x)|.
\]
W szczególności dwa rozwiązania klasyczne zagadnienia
$-\Delta u=f$ z takimi samymi wartościami brzegowymi
muszą być równe.
\end{wniosek}
\begin{proof}
Różnica $w=u-v$ jest harmoniczna. Twierdzenie
\ref{thm:laplace-max}, zastosowane do $w$ i $-w$,
daje odpowiednio oszacowania $w$ z góry i z dołu
przez wartości na $\partial\Omega$. Złożenie tych
oszacowań daje nierówność dla $|w|$.
\end{proof}

Zauważmy, że udowodniliśmy \emph{jednoznaczność i stabilność},
ale nie istnienie rozwiązania Dirichleta dla dowolnego obszaru
i dowolnych danych. Istnienie wymaga osobnej argumentacji.

\subsection{Warunek Neumanna: brak jednoznaczności i warunek zgodności}

\begin{stwierdzenie}[Konieczny warunek zgodności]
Niech $\Omega\subset\R^n$ będzie obszarem ograniczonym
o brzegu klasy $C^1$, a $f\in C(\overline\Omega)$ i
$g\in C(\partial\Omega)$. Jeśli
$u\in C^2(\overline\Omega)$ spełnia
$-\Delta u=f$ w $\Omega$ oraz $\partial_\nu u=g$
na $\partial\Omega$, to
\begin{equation}\label{eq:neumann-compatibility}
 \int_\Omega f(x)\,\mathrm d\lambda_n(x)
 +\int_{\partial\Omega}g(x)\,\mathrm d\sigma(x)=0.
\end{equation}
Tutaj $\lambda_n$ oznacza miarę Lebesgue'a na $\R^n$,
a $\sigma$ miarę powierzchniową na $\partial\Omega$.
Jeśli $u,v\in C^2(\overline\Omega)$ rozwiązują to samo
zagadnienie Neumanna, to $u-v$ jest stała na $\Omega$.
\end{stwierdzenie}
\begin{proof}
Z twierdzenia Gaussa wynika
$\int_\Omega\Delta u(x)\,\mathrm d\lambda_n(x)
=\int_{\partial\Omega}\partial_\nu u(x)\,\mathrm d\sigma(x)$.
Po wstawieniu równania otrzymujemy
\eqref{eq:neumann-compatibility}.

Dla $w=u-v$ zachodzi $\Delta w=0$ i $\partial_\nu w=0$.
Pierwsza tożsamość Greena daje
\[
 \int_\Omega|\nabla w(x)|^2\,\mathrm d\lambda_n(x)
 =\int_{\partial\Omega}w(x)\partial_\nu w(x)\,
 \mathrm d\sigma(x)
 -\int_\Omega w(x)\Delta w(x)\,\mathrm d\lambda_n(x)=0.
\]
Z ciągłości $\nabla w$ wynika $\nabla w=0$ w $\Omega$,
a spójność $\Omega$ implikuje stałość $w$.
\end{proof}

\subsection{Transport: rozwiązanie przesuwa się po charakterystykach}

\begin{stwierdzenie}[Transport ze stałą prędkością]
\label{prop:transport}
Niech $c\in\R$ i $u_0\in C^1(\R)$. Funkcja
$u\colon[0,\infty)\times\R\to\R$ określona wzorem
$u(t,x)=u_0(x-ct)$ jest rozwiązaniem klasycznym
$u_t+c\,u_x=0$ z warunkiem $u(0,x)=u_0(x)$.
Jest to jedyne rozwiązanie klasy $C^1$ na
$[0,\infty)\times\R$ z tym warunkiem.
\end{stwierdzenie}
\begin{proof}
Różniczkując, dostajemy $u_t=-c\,u_0'(x-ct)$
i $u_x=u_0'(x-ct)$; równanie jest spełnione.
Jeśli $v$ jest dowolnym innym rozwiązaniem, to dla
ustalonego $\xi\in\R$ funkcja
$h(t)=v(t,\xi+ct)$ spełnia
$h'(t)=v_t(t,\xi+ct)+c v_x(t,\xi+ct)=0$.
Zatem $v(t,\xi+ct)=v(0,\xi)=u_0(\xi)$, czyli
$v(t,x)=u_0(x-ct)$.
\end{proof}

\begin{center}
\begin{tikzpicture}[>=Stealth,scale=.95]
 \draw[->] (-3.5,0)--(4.3,0) node[right] {$x$};
 \draw[->] (-3.15,-.15)--(-3.15,3.1) node[above] {$t$};
 \foreach \x in {-2.4,-1.3,-.2,.9} {
   \draw[blue!55!black,thick] (\x,0)--({\x+1.3},2.6);
 }
 \draw[densely dashed] (-3.15,2.6)--(3.7,2.6);
 \node[below] at (-1.3,0) {$\xi$};
 \node[above right] at (0,2.6) {$(t,\xi+ct)$};
 \node[anchor=west] at (2.6,1) {$c>0$};
\end{tikzpicture}
\par\smallskip\small
Charakterystyki równania $u_t+c\,u_x=0$ są prostymi $x=\xi+ct$;
wartość $u$ pozostaje na każdej z nich stała (rysunek dla $c>0$).
\end{center}

\begin{przyklad}[Utrata gładkości mimo gładkich danych]
\label{ex:burgers-blowup}
Rozważmy na $[0,T)\times\R$ równanie
$u_t+u\,u_x=0$ z warunkiem $u(0,x)=-\sin x$.
Wzdłuż charakterystyki startującej z punktu $\xi$
wartość $u$ jest stała i równa $-\sin\xi$;
zatem charakterystyka ma równanie
\begin{equation}\label{eq:burgers-characteristics}
 x=\xi-t\sin\xi,\qquad u(t,x)=-\sin\xi.
\end{equation}
Dla $0\leq t<1$ pochodna odwzorowania
$\xi\mapsto\xi-t\sin\xi$ wynosi $1-t\cos\xi\geq1-t>0$.
Odwzorowanie jest bijekcją $\R\to\R$ i ma gładką odwrotność,
więc~\eqref{eq:burgers-characteristics} definiuje gładkie
rozwiązanie. Różniczkując względem $x$ i korzystając
z reguły łańcuchowej, dostajemy
\begin{equation}\label{eq:burgers-gradient}
 u_x(t,x)=\frac{-\cos\xi}{1-t\cos\xi}.
\end{equation}
Dla $x=0$ odpowiada $\xi=0$, a więc
$u_x(t,0)=-1/(1-t)\to-\infty$ gdy $t\uparrow1$.
Rozwiązanie pozostaje ograniczone, ale jego pochodna
przestrzenna staje się nieograniczona w skończonym czasie.
To najprostszy mechanizm powstawania osobliwości i
fal uderzeniowych.
\end{przyklad}

\subsection{Równanie ciepła: natychmiastowe rozprzestrzenianie}

\begin{stwierdzenie}[Jądro ciepła na całej przestrzeni]
\label{prop:heat-kernel}
Niech $d\geq1$, $\kappa>0$, $u_0\in C_c^\infty(\R^d)$.
Dla $t>0$, $x\in\R^d$ zdefiniujmy
\begin{equation}\label{eq:heat-kernel}
 \Gamma_\kappa(t,x)=\frac{1}{(4\pi\kappa t)^{d/2}}
 \exp\!\left(-\frac{|x|^2}{4\kappa t}\right),\qquad
 u(t,x)=\int_{\R^d}\Gamma_\kappa(t,x-y)u_0(y)\,
 \mathrm d\lambda_d(y).
\end{equation}
Wówczas $u\in C^\infty((0,\infty)\times\R^d)$,
$u_t-\kappa\Delta_xu=0$ i
$u(t,\cdot)\to u_0$ jednostajnie, gdy $t\downarrow0$.
Jeśli dodatkowo $u_0\geq0$ i $u_0\not\equiv0$, to
$u(t,x)>0$ dla każdego $t>0$ oraz $x\in\R^d$.
\end{stwierdzenie}
\begin{proof}
Bezpośrednie różniczkowanie funkcji $\Gamma_\kappa$
(dla $t>0$) daje
$\partial_t\Gamma_\kappa=\kappa\Delta_x\Gamma_\kappa$.
Ponieważ $u_0$ ma zwarty nośnik, na każdym zwartym podzbiorze
$(0,\infty)\times\R^d$ możemy dowolnie wiele razy
różniczkować całkę~\eqref{eq:heat-kernel} pod znakiem całki.
Stąd gładkość i równanie ciepła.

Ze wzoru na całkę Gaussa wynika
$\int_{\R^d}\Gamma_\kappa(t,z)\,\mathrm d\lambda_d(z)=1$.
Dlatego
$u(t,x)-u_0(x)=\int_{\R^d}\Gamma_\kappa(t,z)
\bigl(u_0(x-z)-u_0(x)\bigr)\,\mathrm d\lambda_d(z)$.
Funkcja $u_0$ jest jednostajnie ciągła. Ustalmy $\eps>0$
i wybierzmy $r>0$ tak, aby
$|u_0(x-z)-u_0(x)|<\eps$ dla $|z|<r$, niezależnie od $x$.
Część całki po $|z|<r$ ma moduł nie większy od $\eps$;
część po $|z|\geq r$ jest ograniczona przez
$2\|u_0\|_\infty\int_{|z|\geq r}\Gamma_\kappa(t,z)
\,\mathrm d\lambda_d(z)$, a ostatnia całka dąży do zera,
gdy $t\downarrow0$ (po podstawieniu $z=\sqrt t\,w$
i zastosowaniu zbieżności zdominowanej). Dowodzi to
zbieżności jednostajnej.

Jeżeli $u_0\geq0$ i $u_0\not\equiv0$, to z ciągłości
wynika, że $u_0$ jest dodatnia na pewnym niepustym zbiorze
otwartym. Ponieważ $\Gamma_\kappa(t,x-y)>0$ dla wszystkich
$t>0$, $x,y\in\R^d$, całka we wzorze na $u(t,x)$
jest ściśle dodatnia.
\end{proof}

Jest to pierwsza istotna różnica między równaniem ciepła
a równaniem falowym: nawet jeśli $u_0$ jest dodatnia tylko na
niewielkim obszarze, dla dowolnie małego $t>0$ rozwiązanie
równania ciepła staje się dodatnie w całej $\R^d$.
Nie oznacza to, że zaburzenie ma wszędzie porównywalną
wielkość; daleko od nośnika danych wartości są bardzo małe.

\begin{stwierdzenie}[Energetyczna jednoznaczność równania ciepła]
\label{prop:heat-uniqueness}
Niech $\Omega\subset\R^d$ będzie obszarem ograniczonym
o brzegu klasy $C^1$, $T>0$ i $\kappa>0$.
Jeśli $u,v\in C^{1,2}([0,T]\times\overline\Omega)$
spełniają $w_t-\kappa\Delta_xw=0$ dla $w=u-v$
na $(0,T)\times\Omega$, a ponadto
$u(0,\cdot)=v(0,\cdot)$ w $\Omega$ oraz
$u=v$ na $[0,T]\times\partial\Omega$, to $u=v$
na $[0,T]\times\overline\Omega$.
\end{stwierdzenie}
\begin{proof}
Niech $E(t)=\frac12\int_\Omega w(t,x)^2\,
\mathrm d\lambda_d(x)$. Różniczkowanie pod znakiem całki,
a następnie całkowanie przez części (przy $w=0$ na brzegu)
daje
\[
 E'(t)=\int_\Omega w(t,x)w_t(t,x)\,\mathrm d\lambda_d(x)
 =-\kappa\int_\Omega|\nabla_xw(t,x)|^2\,
 \mathrm d\lambda_d(x)\leq0.
\]
Ponieważ $E(0)=0$ i $E(t)\geq0$, mamy $E(t)=0$
dla każdego $t\in[0,T]$. Z ciągłości $w$ wynika
$w\equiv0$.
\end{proof}

\subsection{Równanie falowe: skończona prędkość propagacji}

\begin{twierdzenie}[Wzór d'Alemberta]
\label{thm:dalembert}
Niech $c>0$, $u_0\in C^2(\R)$ i $u_1\in C^1(\R)$.
Jedynym rozwiązaniem klasy $C^2$ na $[0,\infty)\times\R$
zagadnienia
$u_{tt}-c^2u_{xx}=0$, $u(0,x)=u_0(x)$,
$u_t(0,x)=u_1(x)$ jest
\begin{equation}\label{eq:dalembert}
 u(t,x)=\frac{u_0(x-ct)+u_0(x+ct)}{2}
 +\frac{1}{2c}\int_{x-ct}^{x+ct}u_1(s)\,\mathrm ds.
\end{equation}
Jeśli $\supp u_0\cup\supp u_1\subset[-R,R]$ dla $R>0$,
to $u(t,x)=0$ dla $|x|>R+ct$.
\end{twierdzenie}
\begin{proof}
Wprowadzamy zmienne $\xi=x+ct$ i $\eta=x-ct$.
Dla $v(\xi,\eta)=u(t,x)$ równanie falowe staje się
$\partial_\xi\partial_\eta v=0$ (z niezerowym
czynnikiem $-4c^2$ przed pochodną). Na płaszczyźnie
$(\xi,\eta)$ oznacza to $v(\xi,\eta)=F(\xi)+G(\eta)$.
Zatem $u(t,x)=F(x+ct)+G(x-ct)$.
Warunki początkowe dają
$F(x)+G(x)=u_0(x)$ i
$cF'(x)-cG'(x)=u_1(x)$.
Różniczkując pierwszą równość i rozwiązując układ względem
$F',G'$, otrzymujemy
$F'=(u_0'+u_1/c)/2$ i
$G'=(u_0'-u_1/c)/2$.
Całkowanie tych pochodnych, a następnie dodanie $F$ i $G$,
daje~\eqref{eq:dalembert}; stałe całkowania znoszą się
dzięki warunkowi $F+G=u_0$.
Otrzymany wzór definiuje rozwiązanie $C^2$, co można także
sprawdzić przez różniczkowanie, i pokazuje jego jednoznaczność.

Jeżeli $|x|>R+ct$, to oba punkty $x-ct$, $x+ct$ leżą
po tej samej stronie przedziału $[-R,R]$, a przedział
$[x-ct,x+ct]$ nie przecina nośnika $u_1$.
Wszystkie składniki~\eqref{eq:dalembert} są więc równe zeru.
\end{proof}

\begin{center}
\begin{tikzpicture}[>=Stealth,scale=1]
 \draw[->] (-4.0,0)--(4.05,0) node[right] {$x$};
 \draw[->] (0,-.2)--(0,3.9) node[above] {$t$};
 \coordinate (P) at (.7,3.0);
 \coordinate (L) at (-1.6,0);
 \coordinate (R) at (3.0,0);
 \fill[gray!14] (L)--(P)--(R)--cycle;
 \draw[blue!65!black,thick] (L)--(P)--(R);
 \draw[very thick] (L)--(R);
 \fill (P) circle (2pt);
 \node[above right] at (P) {$(x_0,t_0)$};
 \node[below] at (L) {$x_0-ct_0$};
 \node[below] at (R) {$x_0+ct_0$};
 \node[align=center] at (.7,1.0) {obszar\\zależności};
\end{tikzpicture}
\par\smallskip\small
Dla równania falowego w jednym wymiarze wartość w punkcie
$(x_0,t_0)$ zależy wyłącznie od danych początkowych na
$[x_0-ct_0,x_0+ct_0]$.
\end{center}

\begin{stwierdzenie}[Lokalny bilans energii w zwężającym się przedziale]
\label{prop:pde-v31-wave-cone}
Niech $c>0$, $a<b$, $0<T<(b-a)/(2c)$ oraz
$u\in C^2([0,T]\times\R)$ spełnia $u_{tt}-c^2u_{xx}=0$.
Dla $0\le t\le T$ połóżmy $L(t)=a+ct$, $R(t)=b-ct$ i
\[
 E(t)=\frac12\int_{L(t)}^{R(t)}
 \bigl(u_t(t,x)^2+c^2u_x(t,x)^2\bigr)\,\mathrm d\lambda_1(x).
\]
Wtedy
\begin{equation}\label{eq:pde-v31-wave-cone}
 E'(t)=-\frac c2\bigl(u_t-cu_x\bigr)^2(t,R(t))
       -\frac c2\bigl(u_t+cu_x\bigr)^2(t,L(t))\le0.
\end{equation}
Jeśli ponadto $u(0,x)=u_t(0,x)=0$ dla $x\in[a,b]$,
to $u(t,x)=0$ dla $0\le t\le T$ i $L(t)\le x\le R(t)$.
\end{stwierdzenie}
\begin{proof}
Oznaczmy $e=(u_t^2+c^2u_x^2)/2$. Z równania falowego,
reguły iloczynu i równości pochodnych mieszanych otrzymujemy
$\partial_t e=c^2\partial_x(u_tu_x)$.
Reguła Leibniza dla całki o ruchomych końcach, z
$L'=c$ oraz $R'=-c$, daje
\[
 E'(t)=c^2[u_tu_x]_{L(t)}^{R(t)}
       -ce(t,R(t))-ce(t,L(t)).
\]
Po osobnym zebraniu wyrazów przy prawym i lewym końcu
przedziału otrzymujemy dokładnie~\eqref{eq:pde-v31-wave-cone}.
Jeśli dane początkowe są zerowe na $[a,b]$, to $E(0)=0$.
Ponieważ $E\ge0$ i $E$ jest nierosnąca, mamy $E(t)=0$.
Ciągłość pochodnych daje $u_t=u_x=0$ w otwartym stożku
$0<t<T$, $L(t)<x<R(t)$. Dla każdego punktu tego stożka
odcinek pionowy od chwili $0$ pozostaje w stożku, zatem
$u(t,x)=u(0,x)=0$. Ciągłość daje też równość na jego brzegu.
\end{proof}

\begin{uwaga}[Dlaczego bilans lokalny jest ważniejszy od samego wzoru]
Oszacowanie~\eqref{eq:pde-v31-wave-cone} dowodzi skończonej
prędkości propagacji bez użycia jawnego wzoru d'Alemberta.
Ta sama idea, po modyfikacji geometrii i wyrazów brzegowych,
pojawia się przy wielowymiarowym równaniu falowym i układach
hiperbolicznych; tam jawny wzór na ogół nie jest dostępny.
\end{uwaga}

\section{Poprawność w sensie Hadamarda}
\label{sec:pde-hadamard}

Wyrażenie „rozwiązać równanie” nie jest wystarczające,
jeśli nie kontrolujemy wpływu niewielkich perturbacji danych.
Ustalmy przestrzenie danych $X$ i rozwiązań $Y$, wyposażone
w określone topologie, oraz dopuszczalny zbiór danych
$D\subset X$.

\begin{definicja}[Poprawność Hadamarda]
Zagadnienie jest \emph{poprawnie postawione w sensie Hadamarda}
w parze przestrzeni $(X,Y)$ na zbiorze danych $D\subset X$,
jeśli:
\begin{enumerate}
\item dla każdego $f\in D$ istnieje rozwiązanie $u\in Y$;
\item takie rozwiązanie jest jednoznaczne w $Y$;
\item odwzorowanie rozwiązujące $S\colon D\to Y$, $S(f)=u$,
jest ciągłe względem topologii podprzestrzeni na $D$.
\end{enumerate}
W zależności od zagadnienia „dane” mogą obejmować prawą stronę,
warunki początkowe, warunki brzegowe lub ich kombinację.
\end{definicja}

Trzeci warunek jest równie istotny jak pozostałe: dane
pochodzące z pomiarów lub obliczeń numerycznych zawsze
zawierają niewielki błąd.

\begin{kontrprzyklad}[Niestabilne zagadnienie Cauchy'ego dla laplasjanu]
\label{ex:laplace-illposed}
Niech $\Omega=(0,\pi)\times(0,1)$ i rozważmy równanie
$u_{xx}+u_{yy}=0$ w $\Omega$ z danymi na dolnym brzegu
$u(x,0)=0$ i $u_y(x,0)=g(x)$.
Dla każdej liczby naturalnej $k\geq1$ określmy
\begin{equation}\label{eq:hadamard-counterexample}
 \eps_k=e^{-k/2},\qquad
 u_k(x,y)=\frac{\eps_k}{k}\sinh(ky)\sin(kx).
\end{equation}
Wtedy $u_k$ jest funkcją harmoniczną,
$u_k(x,0)=0$ i $\partial_yu_k(x,0)=\eps_k\sin(kx)$.
Dla każdego ustalonego $r\in\N_0$ zachodzi
\[
 \|g_k\|_{C^r([0,\pi])}\leq C_r k^r e^{-k/2}\longrightarrow0,
 \qquad g_k(x)=\eps_k\sin(kx),
\]
przy stałej $C_r$ niezależnej od $k$.
A więc dane dążą do zera nawet w topologii $C^\infty$.
Tymczasem na poziomie $y=3/4$ otrzymujemy
\[
 \|u_k(\cdot,3/4)\|_{L^2(0,\pi)}
 =\sqrt{\frac\pi2}\,\frac{e^{-k/2}}{k}\sinh(3k/4)
 \underset{k\to\infty}{\sim}
 \sqrt{\frac\pi8}\,\frac{e^{k/4}}{k}
 \longrightarrow\infty.
\]
W konsekwencji nie istnieje ciągła w zerze operacja
rekonstrukcji rozwiązania z tych danych Cauchy'ego,
która wszystkie $u_k$ przypisywałaby danym $g_k$ i
której przestrzeń wynikowa kontrolowałaby powyższą normę.
Zagadnienie jest niestabilne pomimo gładkości równania
i danych, a także mimo tego, że prosta $y=0$ jest
niecharakterystyczna dla laplasjanu.
\end{kontrprzyklad}

\begin{uwaga}[Dlaczego to nie przeczy zasadzie maksimum?]
We wniosku~\ref{cor:dirichlet-stability} podawaliśmy
\emph{wartości funkcji na całym brzegu}.
W kontrprzykładzie~\ref{ex:laplace-illposed} zadaliśmy
\emph{wartość i pochodną normalną tylko na części brzegu}.
To dwa zupełnie różne zagadnienia. W szczególności
zasada maksimum nie oszacowuje rozwiązania w prostokącie
przez dane Cauchy'ego wyłącznie na dolnej krawędzi.
\end{uwaga}

\begin{kontrprzyklad}[Równanie ciepła odwrócone w czasie]
\label{ex:backward-heat}
Niech $\kappa>0$, $T>0$, $\T=\R/(2\pi\mathbb Z)$.
Rozważmy na $[0,T]\times\T$ równanie
$u_t=-\kappa u_{xx}$. Dla każdej liczby całkowitej $k\geq1$
funkcja
\[
 u_k(t,x)=\exp\!\bigl(\kappa k^2(t-T/2)\bigr)\sin(kx)
\]
jest gładkim rozwiązaniem. Jej dane początkowe
$u_k(0,x)=e^{-\kappa k^2T/2}\sin(kx)$ dążą do zera
w $C^\infty(\T)$, podczas gdy
$\|u_k(T,\cdot)\|_{L^2(\T)}=
\sqrt\pi\,e^{\kappa k^2T/2}\to\infty$.
Zatem problem ewolucji wstecznej jest niestabilny;
formalna zmiana znaku przy laplasjanie odwraca mechanizm
wygładzania w gwałtowne wzmacnianie wysokich częstości.
\end{kontrprzyklad}

\section{Skalowanie i związek typu równania z geometrią}
\label{sec:pde-scaling}

Dla równań ewolucyjnych czas i przestrzeń nie zawsze
skalują się w ten sam sposób. Współczynniki $c$ i $\kappa$
mają odpowiednio wymiary długość/czas oraz
długość$^2$/czas. To sugeruje dwa odmienne skalowania.

\begin{stwierdzenie}[Skalowania równań ciepła i falowego]
Niech $\lambda>0$ i $d\geq1$.
\begin{enumerate}
\item Jeśli $u$ jest rozwiązaniem klasycznym
$u_t-\kappa\Delta_xu=0$ na zbiorze otwartym
$Q\subset\R\times\R^d$, to
$u_\lambda(t,x)=u(\lambda^2t,\lambda x)$
rozwiązuje to samo równanie na
$Q_\lambda=\{(t,x):(\lambda^2t,\lambda x)\in Q\}$.
\item Jeśli $u$ jest rozwiązaniem klasycznym
$u_{tt}-c^2\Delta_xu=0$ na $Q$, to
$v_\lambda(t,x)=u(\lambda t,\lambda x)$
rozwiązuje to samo równanie na
$\{(t,x):(\lambda t,\lambda x)\in Q\}$.
\end{enumerate}
\end{stwierdzenie}
\begin{proof}
W pierwszym przypadku reguła łańcuchowa daje
$\partial_tu_\lambda=\lambda^2u_t(\lambda^2t,\lambda x)$
i $\Delta_xu_\lambda=\lambda^2\Delta_xu(\lambda^2t,\lambda x)$.
W drugim przypadku zarówno druga pochodna czasowa, jak
i każda druga pochodna przestrzenna wprowadza czynnik
$\lambda^2$. Po podstawieniu odpowiednie równania
są spełnione.
\end{proof}

Przy skali przestrzennej $\ell$ typowy czas dyfuzji ma więc
rzędu $\ell^2/\kappa$, natomiast czas propagacji fali
na odległość $\ell$ ma rzędu $\ell/c$.
To różnica strukturalna, nie tylko rachunkowa.

\section{Co będzie potrzebne dalej?}
\label{sec:pde-outlook}

Podsumujmy rozróżnienia, które pozostaną ważne przez cały skrypt:
\begin{enumerate}
\item \emph{Równanie a zagadnienie.} Dwa takie same operatory
mogą prowadzić do różnych problemów przy innych danych
początkowych lub brzegowych.
\item \emph{Istnienie, jednoznaczność, stabilność.} Każde z tych
zagadnień wymaga osobnego rozumowania. Udowodnienie
jednoznaczności nie implikuje istnienia.
\item \emph{Typ operatora.} Eliptyczność, hiperboliczność
i paraboliczność opisują zasadniczo różne mechanizmy.
\item \emph{Klasa rozwiązań.} Rozwiązania klasyczne są tylko
początkiem. Naturalnymi obiektami bywają funkcje Sobolewa,
dystrybucje oraz rozwiązania słabe i lepkościowe.
\item \emph{Rola geometrii.} Symbol główny, charakterystyki,
struktura brzegu i skalowanie decydują o propagacji
i regularności rozwiązań.
\end{enumerate}

W następnym rozdziale rozwiniemy metodę charakterystyk dla
równań pierwszego rzędu: od liniowego transportu przez
quasiliniowe zagadnienia Cauchy'ego po tworzenie się
fal uderzeniowych.

% Rozdział 2.
\chapter[Charakterystyki i PDE pierwszego rzędu]{Równania pierwszego rzędu i metoda charakterystyk}
\label{chap:pde-characteristics}

W tym rozdziale zamienimy równania pierwszego rzędu na układy
równań zwyczajnych. Wersja liniowa metody pozwala rozwiązywać
równania transportu i ciągłości, natomiast wersja quasiliniowa
wyjaśnia powstawanie osobliwości z gładkich danych. Dla równań
w pełni nieliniowych skonstruujemy układ charakterystyczny
Charpita i zastosujemy go do równań Hamiltona--Jacobiego.

W całym rozdziale $d\geq1$, $T>0$ oraz $\lambda_d$ oznacza
miarę Lebesgue'a na $\R^d$. Wszystkie pochodne występujące
w sformułowaniach klasycznych rozumiemy w sensie punktowym.
W szczególności przez $C^1([0,T]\times\R^d)$ rozumiemy
funkcje o pochodnych pierwszego rzędu ciągłych aż do
$t=0,T$ (jednostronnie w tych punktach).

\section{Pochodna wzdłuż krzywej i najprostsze równanie transportu}
\label{sec:pde2-basic}

Niech $u\in C^1((0,T)\times\R^d)$ i niech
$X\colon I\to\R^d$ będzie krzywą klasy $C^1$, gdzie
$I\subset(0,T)$ jest przedziałem. Reguła łańcuchowa daje
\begin{equation}\label{eq:pde2-chain}
 \frac{\mathrm d}{\mathrm dt}u(t,X(t))
 =u_t(t,X(t))+\nabla_xu(t,X(t))\cdot\dot X(t).
\end{equation}
Dobierając $\dot X(t)$ do współczynników równania,
możemy przekształcić PDE w równanie zwyczajne
wzdłuż krzywej $t\mapsto(t,X(t))$. Taką krzywą nazywamy
\emph{charakterystyką} równania transportu.

\begin{stwierdzenie}[Transport ze stałą prędkością i źródłem]
\label{prop:pde2-constant}
Niech $c\in\R^d$, $u_0\in C^1(\R^d)$,
$f\in C^1([0,T]\times\R^d)$. Istnieje dokładnie jedno
rozwiązanie $u\in C^1([0,T]\times\R^d)$ zagadnienia
\[
 u_t+c\cdot\nabla_xu=f(t,x),\qquad u(0,x)=u_0(x).
\]
Jest ono dane wzorem
\begin{equation}\label{eq:pde2-constant-formula}
 u(t,x)=u_0(x-tc)+\int_0^t f(s,x-(t-s)c)\,\mathrm ds.
\end{equation}
\end{stwierdzenie}
\begin{proof}
Charakterystyka wychodząca z $\xi$ ma postać
$X(t;\xi)=\xi+tc$. Dla każdego rozwiązania klasy $C^1$
reguła \eqref{eq:pde2-chain} pokazuje, że
$\frac{\mathrm d}{\mathrm dt}u(t,\xi+tc)=f(t,\xi+tc)$.
Całkując od $0$ do $t$ i podstawiając $\xi=x-tc$,
otrzymujemy \eqref{eq:pde2-constant-formula}; dowodzi to
jednoznaczności. Odwrotnie, wzór definiuje funkcję klasy
$C^1$. Po podstawieniu $x=\xi+tc$ jego prawa strona
jest równa $u_0(\xi)+\int_0^t f(s,\xi+sc)\,\mathrm ds$,
zatem reguła łańcuchowa daje żądane PDE i warunek początkowy.
\end{proof}

\begin{przyklad}[Tłumienie wzdłuż charakterystyk]
Niech $\gamma>0$. Dla równania
$u_t+c\cdot\nabla u+\gamma u=0$ z danymi $u_0$
otrzymujemy $u(t,x)=e^{-\gamma t}u_0(x-tc)$.
Transport przesuwa profil, podczas gdy czynnik wykładniczy
tłumi jego amplitudę.
\end{przyklad}

\section{Zmienne pole prędkości: przepływ charakterystyk}
\label{sec:pde2-linear}

Rozważmy równanie
\begin{equation}\label{eq:pde2-linear}
 u_t+b(t,x)\cdot\nabla_xu+c(t,x)u=f(t,x)
 \quad\text{na }(0,T)\times\R^d.
\end{equation}
Dla przejrzystości przyjmiemy dość mocne założenia
zapewniające globalny przepływ; później będziemy je osłabiać.

\begin{twierdzenie}[Metoda charakterystyk dla transportu liniowego]
\label{thm:pde2-linear}
Niech $b\colon[0,T]\times\R^d\to\R^d$ będzie klasy $C^1$
i niech istnieje $M>0$ takie, że
$\|D_xb(t,x)\|\leq M$ i $|b(t,x)|\leq M(1+|x|)$
dla wszystkich $(t,x)$. Niech $c,f\in C^1([0,T]\times\R^d)$
oraz $u_0\in C^1(\R^d)$. Wówczas istnieje dokładnie jedno
rozwiązanie $u\in C^1([0,T]\times\R^d)$ zagadnienia
\eqref{eq:pde2-linear} z $u(0,x)=u_0(x)$.

Niech $X(t;\xi)$ będzie rozwiązaniem
\begin{equation}\label{eq:pde2-flow-ode}
 \dot X(t;\xi)=b(t,X(t;\xi)),\qquad X(0;\xi)=\xi.
\end{equation}
Dla każdego $t\in[0,T]$ odwzorowanie
$X_t\colon\xi\mapsto X(t;\xi)$ jest dyfeomorfizmem klasy $C^1$
przestrzeni $\R^d$ na siebie. Rozwiązanie jest określone przez
\begin{equation}\label{eq:pde2-solution}
 \begin{aligned}
 u(t,X(t;\xi))={}&e^{-C(t;\xi)}u_0(\xi)\\
 &+\int_0^t e^{-(C(t;\xi)-C(s;\xi))}
 f(s,X(s;\xi))\,\mathrm ds,
 \end{aligned}
\end{equation}
gdzie $C(t;\xi)=\int_0^t c(r,X(r;\xi))\,\mathrm dr$.
\end{twierdzenie}
\begin{proof}
Pole $b$ jest lokalnie lipschitzowskie względem $x$,
a dzięki warunkowi na $D_xb$ jest jednostajnie
lipschitzowskie na całej przestrzeni. Wzrost co najwyżej
liniowy i nierówność Gronwalla wykluczają ucieczkę
rozwiązania \eqref{eq:pde2-flow-ode} do nieskończoności
w skończonym czasie. Twierdzenie o zależności rozwiązań
ODE od warunków początkowych daje $C^1$-gładkość $X$
względem $(t,\xi)$.

Dla ustalonego $t$ odwrotne odwzorowanie wyznacza
rozwiązanie tego samego ODE rozwiązywane od chwili $t$
do chwili $0$. Jednoznaczność ODE zapewnia, że jest to
$X_t^{-1}$, a regularna zależność od warunku końcowego
pokazuje, że odwrotność jest klasy $C^1$. Tak więc
$X_t$ jest dyfeomorfizmem.

Dla dowolnego rozwiązania $u$ postawmy
$U(t;\xi)=u(t,X(t;\xi))$. Z \eqref{eq:pde2-chain}
wynika $\dot U+c(t,X(t;\xi))U=f(t,X(t;\xi))$.
Mnożąc przez czynnik całkujący $e^{C(t;\xi)}$,
otrzymujemy \eqref{eq:pde2-solution}.

Odwrotnie, prawa strona \eqref{eq:pde2-solution} wyznacza
funkcję $U\in C^1$; dzięki dyfeomorfizmowi $X_t$
definiuje ona $u(t,x)=U(t;X_t^{-1}(x))$ klasy $C^1$.
Różniczkowanie po charakterystykach daje
\eqref{eq:pde2-linear}. Każde inne rozwiązanie prowadzi
do tej samej funkcji $U$, co kończy dowód jednoznaczności.
\end{proof}

\begin{przyklad}[Transport w polu dylatacji]
Niech $b(x)=x$, $c=f=0$, $u_0\in C^1(\R^d)$.
Wówczas $X(t;\xi)=e^t\xi$ i
$u(t,x)=u_0(e^{-t}x)$. Profil rozszerza się
wraz z przepływem, ale wartości wzdłuż charakterystyk
nie zmieniają się.
\end{przyklad}

\begin{uwaga}[Odwracalność przepływu a odwracalność ewolucji]
W twierdzeniu \ref{thm:pde2-linear} w skończonym czasie
przepływ przestrzenny jest odwracalny. Nie oznacza to,
że każdy operator ewolucyjny PDE jest odwracalny
w dobrze postawionej klasie funkcji: dla równania ciepła
występuje wygładzanie, a dla transportu quasiliniowego
charakterystyki mogą przestać tworzyć układ współrzędnych.
\end{uwaga}

\section{Jak zmieniają się objętość i masa?}
\label{sec:pde2-jacobian}

Transport wartości skalarnej i transport gęstości to
dwa różne zagadnienia. Dla tego samego pola $b$ porównajmy
$u_t+b\cdot\nabla u=0$ oraz
$\rho_t+\diver_x(\rho b)=0$. Drugie równanie nazywamy
\emph{równaniem ciągłości}; w nim dodatkowo występuje
czynnik $(\diver_x b)\rho$.

\begin{stwierdzenie}[Wzór Jacobiego--Liouville'a]
\label{prop:pde2-liouville}
Niech $b$ spełnia założenia twierdzenia
\ref{thm:pde2-linear}, a $X$ będzie przepływem
\eqref{eq:pde2-flow-ode}. Dla $J(t;\xi)=\det D_\xi X(t;\xi)$
zachodzi
\begin{equation}\label{eq:pde2-jacobian}
 \frac{\mathrm d}{\mathrm dt}J(t;\xi)
 =\bigl(\diver_x b\bigr)(t,X(t;\xi))J(t;\xi),
 \qquad J(0;\xi)=1.
\end{equation}
W szczególności
$J(t;\xi)=\exp\!\bigl(\int_0^t
(\diver_x b)(s,X(s;\xi))\,\mathrm ds\bigr)>0$.
\end{stwierdzenie}
\begin{proof}
Różniczkując \eqref{eq:pde2-flow-ode} względem $\xi$,
otrzymujemy dla $A(t)=D_\xi X(t;\xi)$
$\dot A(t)=D_xb(t,X(t;\xi))A(t)$ i $A(0)=I$.
Macierz $A(t)$ jest odwracalna na mocy twierdzenia
\ref{thm:pde2-linear}. Zatem wzór na pochodną wyznacznika
daje $\frac{\mathrm d}{\mathrm dt}\det A
=\det A\,\tr(A^{-1}\dot A)$.
Cykliczność śladu i $\dot A=D_xb\,A$ prowadzą do
$\tr(A^{-1}\dot A)=\tr(D_xb)=\diver_xb$.
Równanie skalarne \eqref{eq:pde2-jacobian} rozwiązujemy
przez całkowanie.
\end{proof}

\begin{twierdzenie}[Transport gęstości i zachowanie masy]
\label{thm:pde2-mass}
Niech $b\in C^2([0,T]\times\R^d;\R^d)$ spełnia
warunki wzrostu i lipschitzowskości z twierdzenia
\ref{thm:pde2-linear}. Zakładamy ponadto, że wszystkie
pochodne przestrzenne $b$ rzędu $1$ i $2$ są ograniczone.
Dla dowolnej $\rho_0\in C_c^1(\R^d)$ istnieje dokładnie jedno
rozwiązanie $\rho\in C^1([0,T]\times\R^d)$ równania
\[
 \rho_t+\diver_x(\rho b)=0,\qquad \rho(0,x)=\rho_0(x).
\]
Jest ono zadane zależnością
\begin{equation}\label{eq:pde2-mass-flow}
 \rho(t,X(t;\xi))J(t;\xi)=\rho_0(\xi).
\end{equation}
Ponadto $\int_{\R^d}\rho(t,x)\,\mathrm d\lambda_d(x)
=\int_{\R^d}\rho_0(x)\,\mathrm d\lambda_d(x)$ dla $t\in[0,T]$.
\end{twierdzenie}
\begin{proof}
Ponieważ $\diver_x(\rho b)=b\cdot\nabla\rho+
(\diver_xb)\rho$, twierdzenie \ref{thm:pde2-linear}
można zastosować z $c=\diver_xb$, $f=0$.
Założenie $b\in C^2$ zapewnia, że $\diver_x b\in C^1$,
więc wzór \eqref{eq:pde2-mass-flow} daje rozwiązanie
klasy $C^1$. Można też bezpośrednio różniczkować
iloczyn po trajektorii: z \eqref{eq:pde2-jacobian}
wynika
\[
 \frac{\mathrm d}{\mathrm dt}
 \bigl(\rho(t,X(t;\xi))J(t;\xi)\bigr)=0.
\]
Następnie stosujemy zmianę zmiennych $x=X(t;\xi)$.
Dodatniość $J$ i \eqref{eq:pde2-mass-flow} dają
\[
 \int_{\R^d}\rho(t,x)\,\mathrm d\lambda_d(x)
 =\int_{\R^d}\rho(t,X(t;\xi))J(t;\xi)
 \,\mathrm d\lambda_d(\xi)
 =\int_{\R^d}\rho_0(\xi)\,\mathrm d\lambda_d(\xi).
\]
Zwarty nośnik $\rho_0$ zapewnia, że całki są skończone.
\end{proof}

\begin{uwaga}[Warunki regularności]
Wzór \eqref{eq:pde2-mass-flow} zachowuje sens również
przy słabszych założeniach. W poprzednim dowodzie
zastosowaliśmy dodatkową regularność $b$, aby pozostać
w klasie rozwiązań klasycznych; teoria dystrybucyjna
pozwoli te wymagania osłabić.
\end{uwaga}

\begin{stwierdzenie}[Kontrola jakobianu i norm przy transporcie]
\label{prop:pde-v31-transport-lp}
Niech $T>0$, $d\ge1$, a $b$ spełnia założenia
twierdzenia~\ref{thm:pde2-linear} z $\|D_xb\|\le M$.
Niech $X_t$ i $J(t;\xi)=\det D_\xi X_t(\xi)$ będą
odpowiednio przepływem i jego jakobianem. Wówczas
\[
 e^{-dMt}\le J(t;\xi)\le e^{dMt}
 \qquad(0\le t\le T,\ \xi\in\R^d).
\]
Jeśli $u_0\in C^1(\R^d)\cap L^p(\R^d)$ dla
$1\le p<\infty$, a $u$ rozwiązuje
$u_t+b\cdot\nabla u=0$, $u(0)=u_0$, to
\begin{equation}\label{eq:pde-v31-transport-lp}
 e^{-dMt/p}\|u_0\|_p\le\|u(t,\cdot)\|_p
 \le e^{dMt/p}\|u_0\|_p.
\end{equation}
Dla ograniczonego $u_0\in C^1(\R^d)$ zachodzi ponadto
$\|u(t,\cdot)\|_\infty=\|u_0\|_\infty$.
Jeśli $\diver_xb=0$, to $J\equiv1$ i wszystkie wymienione
normy $L^p$ są zachowane.
\end{stwierdzenie}
\begin{proof}
Ponieważ każda wartość własna rzeczywistej macierzy
$D_xb$ ma moduł co najwyżej $M$, mamy
$|\tr D_xb|\le dM$ (można też użyć nierówności
$|\tr A|\le d\|A\|$ bez odwoływania się do widma).
Wzór~\eqref{eq:pde2-jacobian} i jego rozwiązanie dają
pierwsze oszacowanie. Wzdłuż charakterystyk
$u(t,X_t(\xi))=u_0(\xi)$. Ponieważ $X_t$ jest
dyfeomorfizmem, wzór na zmianę zmiennych daje
\[
 \|u(t,\cdot)\|_p^p
 =\int_{\R^d}|u_0(\xi)|^pJ(t;\xi)\,
   \mathrm d\lambda_d(\xi).
\]
Podstawiamy ograniczenia na $J$ i wyciągamy pierwiastek
stopnia $p$. Dla normy supremum korzystamy z tego,
że $X_t$ jest bijekcją. W przypadku zerowej dywergencji
równanie na $J$ ma rozwiązanie stałe równe $1$.
\end{proof}

\section{Niecharakterystyczne dane Cauchy'ego na hiperpowierzchni}
\label{sec:pde2-noncharacteristic}

Rozważmy równanie bez wyróżnionego czasu
\begin{equation}\label{eq:pde2-static}
 a(x)\cdot\nabla u(x)+c(x)u(x)=f(x),\qquad x\in\Omega,
\end{equation}
gdzie $\Omega\subset\R^n$ jest otwarta, $n\geq2$.
Gdy dane zadano na hiperpowierzchni $\Sigma$, istotne
jest, czy charakterystyki przecinają ją poprzecznie.

\begin{definicja}[Niecharakterystyczność dla transportu]
Niech $\Omega\subset\R^n$ będzie otwarta, $a\in C^1(\Omega;\R^n)$,
a $\Sigma\subset\Omega$ hiperpowierzchnią klasy $C^2$.
Mówimy, że $\Sigma$ jest niecharakterystyczna w punkcie
$x_0\in\Sigma$ dla operatora $a\cdot\nabla$, jeżeli
$a(x_0)\notin T_{x_0}\Sigma$. Jeśli
$\Sigma=\{x:\phi(x)=0\}$ lokalnie, jest to równoważne
warunkowi $a(x_0)\cdot\nabla\phi(x_0)\ne0$.
\end{definicja}

\begin{twierdzenie}[Lokalne zagadnienie Cauchy'ego]
\label{thm:pde2-cauchy}
Niech $\Omega\subset\R^n$ będzie otwarta,
$a\in C^1(\Omega;\R^n)$, $c,f\in C^1(\Omega)$,
$\Sigma\subset\Omega$
hiperpowierzchnią klasy $C^2$, a $x_0\in\Sigma$.
Niech $g\in C^1(\Sigma)$ w otoczeniu $x_0$.
Jeśli $\Sigma$ jest niecharakterystyczna w $x_0$,
to istnieje otoczenie $U\subset\Omega$ punktu $x_0$
oraz dokładnie jedno rozwiązanie $u\in C^1(U)$ zagadnienia
\eqref{eq:pde2-static} spełniające $u|_{U\cap\Sigma}=g$,
rozumiane jako jednoznaczność zarodka rozwiązania
w otoczeniu $x_0$.
\end{twierdzenie}
\begin{proof}
Wybierzmy lokalną parametryzację
$\gamma\colon V\to\Sigma$, gdzie $V\subset\R^{n-1}$
jest otwarty, $\gamma(0)=x_0$, a wektory
$\partial_{y_1}\gamma(0),\ldots,\partial_{y_{n-1}}\gamma(0)$
są liniowo niezależne. Niech $\Phi_s(x)$ oznacza lokalny
przepływ pola $a$, tj. $\partial_s\Phi_s(x)=a(\Phi_s(x))$
i $\Phi_0(x)=x$. Zdefiniujmy
$\Psi(s,y)=\Phi_s(\gamma(y))$.

Macierz pochodnej $D\Psi(0,0)$ ma kolumny
$a(x_0)$ oraz $\partial_{y_j}\gamma(0)$. Założenie
niecharakterystyczności mówi dokładnie, że tworzą bazę
$\R^n$. Z twierdzenia o funkcji odwrotnej istnieją
otoczenia $W$ punktu $(0,0)$ i $U$ punktu $x_0$,
dla których $\Psi\colon W\to U$ jest dyfeomorfizmem.
Po ewentualnym zmniejszeniu otoczeń można zapewnić,
że $\Psi(0,y)$ parametryzuje całe $U\cap\Sigma$.

Niech $v(s,y)=u(\Psi(s,y))$. Równanie
\eqref{eq:pde2-static} sprowadza się do ODE
$\partial_sv+c(\Psi(s,y))v=f(\Psi(s,y))$
z danymi $v(0,y)=g(\gamma(y))$. Dla
$C(s,y)=\int_0^s c(\Psi(r,y))\,\mathrm dr$ rozwiązanie ma
postać
\[
 v(s,y)=e^{-C(s,y)}g(\gamma(y))+
 \int_0^s e^{-(C(s,y)-C(r,y))}
 f(\Psi(r,y))\,\mathrm dr.
\]
Różniczkowalność przepływu i danych pokazuje, że $v$
jest klasy $C^1$. Złożenie z $\Psi^{-1}$ definiuje
$u\in C^1(U)$, które spełnia równanie i dane. Każde
inne rozwiązanie ma w tych samych współrzędnych
charakterystycznych identyczny problem ODE, a więc
jest równe $u$ na dostatecznie małym otoczeniu $x_0$.
\end{proof}

\begin{center}
\begin{tikzpicture}[>=Stealth,scale=.92]
 \draw[->] (-3.3,-1.35)--(3.8,-1.35) node[right] {$x_1$};
 \draw[->] (-3.1,-1.55)--(-3.1,2.5) node[above] {$x_2$};
 \draw[very thick] (-2.55,-.65)..controls (-.8,-.8) and (1.3,-.18)..(2.9,-.45);
 \node[below] at (2.6,-.6) {$\Sigma$};
 \foreach \x/\yy in {-1.75/-.64,-.65/-.52,.55/-.37,1.65/-.34} {
   \draw[blue!65!black,thick,->] (\x,\yy)--({\x+.48},{\yy+1.85});
   \fill (\x,\yy) circle (1.5pt);
 }
 \fill[red!70!black] (.55,-.37) circle (2.4pt);
 \node[below] at (.55,-.76) {$x_0$};
 \node[above] at (1.7,1.4) {charakterystyki};
\end{tikzpicture}
\par\smallskip\small
Dane są zadawane na $\Sigma$; przechodzące przez nią
poprzecznie charakterystyki tworzą lokalny układ współrzędnych.
\end{center}

\begin{kontrprzyklad}[Dane na powierzchni charakterystycznej]
\label{ex:pde2-char-data}
Rozważmy $u_x=0$ na $\R^2$, gdzie $u=u(x,y)$,
oraz dane $u(x,0)=g(x)$ na prostej $\Sigma=\{y=0\}$.
Każde rozwiązanie ma postać $u(x,y)=h(y)$, stąd
istnienie wymaga, by $g$ była stała. Gdy $g\equiv c$,
rozwiązań jest nieskończenie wiele: można wybrać
dowolne $h\in C^1(\R)$ z $h(0)=c$. Prosta $y=0$
jest charakterystyczna, bo pole $(1,0)$ leży
w jej przestrzeni stycznej.
\end{kontrprzyklad}

\section{Quasiliniowe równania pierwszego rzędu}
\label{sec:pde2-quasilinear}

Teraz prędkość charakterystyki może zależeć od niewiadomej:
\begin{equation}\label{eq:pde2-quasilinear}
 u_t+a(t,x,u)\cdot\nabla_xu=b(t,x,u),
 \qquad u(0,x)=u_0(x).
\end{equation}
W odróżnieniu od przypadku liniowego nie znamy przepływu
z góry: równocześnie musimy wyznaczać jego położenie
$X$ oraz wartość $U$ rozwiązania na charakterystyce.

\begin{twierdzenie}[Lokalna konstrukcja i jednoznaczność]
\label{thm:pde2-quasilinear}
Niech $a\colon(-T,T)\times\R^d\times\R\to\R^d$
i $b\colon(-T,T)\times\R^d\times\R\to\R$
będą klasy $C^2$, a $u_0\in C^2(\R^d)$.
Dla dowolnego $\xi_0\in\R^d$ istnieją otoczenie
$V\subset\R^d$ punktu $\xi_0$, liczba $\delta>0$
i otwarty zbiór $Q\subset(-\delta,\delta)\times\R^d$
zawierający $(0,\xi_0)$, na którym zagadnienie
\eqref{eq:pde2-quasilinear} ma rozwiązanie $u\in C^1(Q)$
z danymi $u(0,x)=u_0(x)$ dla $(0,x)\in Q$.
Rozwiązanie jest jednoznaczne jako zarodek klasy $C^1$
w otoczeniu $(0,\xi_0)$.

Dokładniej, niech $(X,U)$ rozwiązują układ ODE
\begin{equation}\label{eq:pde2-quasi-chars}
 \begin{aligned}
 \dot X(t;\xi)&=a(t,X(t;\xi),U(t;\xi)),& X(0;\xi)&=\xi,\\
 \dot U(t;\xi)&=b(t,X(t;\xi),U(t;\xi)),& U(0;\xi)&=u_0(\xi).
 \end{aligned}
\end{equation}
W dostatecznie małym otoczeniu $(0,\xi_0)$
odwzorowanie $(t,\xi)\mapsto(t,X(t;\xi))$
jest lokalnym dyfeomorfizmem; rozwiązanie spełnia
$u(t,X(t;\xi))=U(t;\xi)$.
\end{twierdzenie}
\begin{proof}
Ponieważ prawa strona układu \eqref{eq:pde2-quasi-chars}
jest klasy $C^2$, lokalne twierdzenie o istnieniu
ODE i gładkiej zależności od danych daje rozwiązania
$(X,U)$ klasy co najmniej $C^1$ względem $(t,\xi)$.
W chwili $t=0$ mamy $D_\xi X(0;\xi)=I$. Zatem
jakobian mapy $(t,\xi)\mapsto(t,X(t;\xi))$ w
$(0,\xi_0)$ jest niezerowy. Twierdzenie o funkcji odwrotnej
pozwala w odpowiednich otoczeniach odwrócić tę mapę.

Niech $\xi=\Xi(t,x)$ oznacza jej lokalną odwrotność
współrzędnych przestrzennych. Kładziemy
$u(t,x)=U(t;\Xi(t,x))$. Różniczkując tożsamość
$u(t,X(t;\xi))=U(t;\xi)$ po $t$ przy ustalonym $\xi$,
otrzymujemy
\[
 u_t(t,X)+\nabla_xu(t,X)\cdot a(t,X,U)
 =b(t,X,U).
\]
Ponieważ $U=u(t,X)$, jest to \eqref{eq:pde2-quasilinear}.
W chwili $t=0$ mamy $X(0;\xi)=\xi$ i
$U(0;\xi)=u_0(\xi)$, więc warunek początkowy zachodzi.

Jeśli $v\in C^1$ jest innym lokalnym rozwiązaniem z tymi
danymi, rozwiązujemy ODE dla $\dot Y=a(t,Y,v(t,Y))$
z $Y(0)=\xi$. Funkcja $V(t)=v(t,Y(t))$ spełnia
$\dot V=b(t,Y,V)$, więc para $(Y,V)$ rozwiązuje układ
\eqref{eq:pde2-quasi-chars}. Jednoznaczność ODE wymusza
$(Y,V)=(X,U)$ na wspólnym małym przedziale; stąd $v=u$
w dostatecznie małym otoczeniu $(0,\xi_0)$.
\end{proof}

\begin{uwaga}[Co może uniemożliwić przedłużenie?]
Lokalna konstrukcja wymaga zarówno istnienia rozwiązań
układu ODE, jak i odwracalności mapy $\xi\mapsto X(t;\xi)$.
W zagadnieniach quasiliniowych nawet przy globalnym
istnieniu charakterystyk może wystąpić
$\det D_\xi X=0$: wtedy klasyczne współrzędne
charakterystyczne zawodzą. W modelach skalarnych jest
to typowy mechanizm utraty regularności rozwiązania.
\end{uwaga}

\section{Fale proste i kryterium wybuchu pochodnej}
\label{sec:pde2-blowup}

Rozważmy w jednym wymiarze równanie autonomiczne
\begin{equation}\label{eq:pde2-simplewave}
 u_t+a(u)u_x=0,\qquad u(0,x)=u_0(x),
\end{equation}
gdzie $a\in C^2(\R)$, $u_0\in C^2(\R)$.
Wzdłuż charakterystyki wartość $u$ jest stała, więc
\begin{equation}\label{eq:pde2-simplewave-map}
 X(t;\xi)=\xi+t\,a(u_0(\xi)),\qquad
 U(t;\xi)=u_0(\xi).
\end{equation}
Różniczkując względem $\xi$, dostajemy
$X_\xi(t;\xi)=1+t\,a'(u_0(\xi))u_0'(\xi)$.

\begin{twierdzenie}[Pierwsza osobliwość gradientu]
\label{thm:pde2-breaking}
Niech $a,u_0\in C^2(\R)$ i załóżmy, że
$a\circ u_0$ jest ograniczona, a funkcja
$q(\xi)=a'(u_0(\xi))u_0'(\xi)$ jest ograniczona z dołu.
Jeżeli $m=\inf_{\xi\in\R}q(\xi)<0$, określmy
$T_*=-1/m$. Wtedy dla każdego $0\leq t<T_*$
odwzorowanie $X_t(\xi)=\xi+t a(u_0(\xi))$
jest ściśle rosnącym dyfeomorfizmem $\R\to\R$ i
\eqref{eq:pde2-simplewave-map} definiuje rozwiązanie
klasy $C^1$ na $[0,T_*)\times\R$. Jego pochodna wynosi
\begin{equation}\label{eq:pde2-gradient-formula}
 u_x(t,X(t;\xi))=
 \frac{u_0'(\xi)}{1+t a'(u_0(\xi))u_0'(\xi)}.
\end{equation}
Jeśli ponadto minimum $m$ jest osiągane w punkcie
$\xi_*$, to $|u_x(t,X(t;\xi_*))|\to\infty$
gdy $t\uparrow T_*$. Nie istnieje zatem przedłużenie
rozwiązania jako funkcji $C^1$ na otoczenie
$(T_*,X(T_*;\xi_*))$.
\end{twierdzenie}
\begin{proof}
Dla $t<T_*$ mamy
$X_\xi(t;\xi)=1+tq(\xi)\geq1+tm>0$.
Odwzorowanie $X_t$ jest więc ściśle rosnące.
Z ograniczoności $a\circ u_0$ wynika
$X_t(\xi)-\xi=O(1)$ dla $\xi\to\pm\infty$,
zatem $X_t$ jest na $\R$. Twierdzenie o funkcji
odwrotnej daje odwrotność klasy $C^1$.
Z równości $u(t,X(t;\xi))=u_0(\xi)$ otrzymujemy
$u_x(t,X)X_\xi=u_0'(\xi)$, czyli
\eqref{eq:pde2-gradient-formula}.

Jeżeli $q(\xi_*)=m<0$, to
$a'(u_0(\xi_*))u_0'(\xi_*)\ne0$; w szczególności
$u_0'(\xi_*)\ne0$. Mianownik
\eqref{eq:pde2-gradient-formula} przy $\xi_*$
dąży do zera od strony dodatniej, a licznik jest niezerowy.
Zatem moduł pochodnej dąży do nieskończoności,
co wyklucza przedłużenie klasy $C^1$.
\end{proof}

\begin{przyklad}[Burgers: dokładny czas powstania osobliwości]
\label{ex:pde2-burgers}
Dla równania $u_t+u u_x=0$ mamy $a(u)=u$.
Wybierzmy $u_0(x)=-\tanh x$. Wówczas
$q(x)=u_0'(x)=-1/\cosh^2x$ osiąga minimum $-1$
w $x=0$, zaś $a\circ u_0$ jest ograniczona.
Twierdzenie \ref{thm:pde2-breaking} daje $T_*=1$,
$X(t;\xi)=\xi-t\tanh\xi$ oraz
$u_x(t,0)=-1/(1-t)$ dla $t<1$.
Charakterystyki istnieją jako krzywe gładkie także
później, ale nie mogą już wyznaczać globalnego
klasycznego rozwiązania w tym samym schemacie.
\end{przyklad}

\begin{center}
\begin{tikzpicture}[>=Stealth,scale=.86]
 \draw[->] (-3.4,0)--(3.6,0) node[right] {$x$};
 \draw[->] (-3.1,-.12)--(-3.1,3.75) node[above] {$t$};
 \foreach \xi in {-2.4,-1.6,-.8,-.4,.4,.8,1.6,2.4} {
  \pgfmathsetmacro{\dd}{-(exp(2*\xi)-1)/(exp(2*\xi)+1)}
  \draw[blue!65!black,thick] (\xi,0)--({\xi+0.95*\dd},2.95);
 }
 \draw[densely dashed] (-3.1,3.1)--(3.35,3.1);
 \node[anchor=west] at (3.55,3.08) {$t\approx T_*$};
 \node[below] at (0,0) {$0$};
 \node[align=left,anchor=west] at (3.6,1.45) {zbliżanie się\\charakterystyk};
\end{tikzpicture}
\par\smallskip\small
Dla malejących danych Burgersa charakterystyki
zbliżają się do siebie. Rysunek ma charakter schematyczny;
czas i nachylenia przeskalowano niezależnie.
\end{center}

\begin{uwaga}[Eksperyment obliczeniowy]
W rozdziale~\ref{chap:pde-code} zamieszczono kod Python,
który dla przykładu Burgersa z danymi $u_0=-\tanh x$
rysuje charakterystyki i profile rozwiązań dla
$t\in\{0,0.5,0.9,0.99\}$. Program korzysta z odwracania
monotonicznej mapy $\xi\mapsto X(t;\xi)$ przez
interpolację. Porównanie z dokładnym wzorem
\eqref{eq:pde2-gradient-formula} pozwala sprawdzić,
że $u_x(t,0)=-1/(1-t)$, podczas gdy $|u(t,x)|\leq1$.
Nie jest to algorytm kontynuacji rozwiązania przez szok:
dla $t\geq1$ mapa etykiet nie jest już globalnie
odwracalna w klasie $C^1$.
\end{uwaga}

\begin{uwaga}[Znaczenie czasu $T_*$]
Twierdzenie \ref{thm:pde2-breaking} dotyczy dokładnie
zadanego równania $u_t+a(u)u_x=0$, bez źródła,
i wskazanej klasy danych. Nie jest ogólnym twierdzeniem
o czasie istnienia dowolnych równań quasiliniowych.
Gdy $q(\xi)\geq0$ dla wszystkich $\xi$, charakterystyki
nie przecinają się dla $t\geq0$, o ile
$a\circ u_0$ pozostaje ograniczona; w takim przypadku
wzór daje globalne rozwiązanie klasy $C^1$.
\end{uwaga}

\section{Od przecinających się charakterystyk do fal uderzeniowych}
\label{sec:pde2-shocks}

Równanie $u_t+a(u)u_x=0$ można, gdy $a=f'$,
zapisać w postaci prawa zachowania
\begin{equation}\label{eq:pde2-conservation}
 u_t+\partial_x f(u)=0.
\end{equation}
Po pojawieniu się nieciągłości wzór z pochodnymi
klasycznymi przestaje mieć sens. Potrzebne jest
sformułowanie całkowe, które później zinterpretujemy
w języku dystrybucji.

\begin{definicja}[Rozwiązanie słabe prawa zachowania]
\label{def:pde2-weak}
Niech $f\in C^1(\R)$, $u_0\in L^1_{\mathrm{loc}}(\R)$,
a $u\in L^1_{\mathrm{loc}}([0,T)\times\R)$, przy czym
$f(u)$ jest lokalnie całkowalna. Funkcję $u$ nazywamy
\emph{rozwiązaniem słabym} zagadnienia
\eqref{eq:pde2-conservation} z danymi $u_0$, jeżeli
dla każdej $\varphi\in C_c^\infty([0,T)\times\R)$
zachodzi
\begin{equation}\label{eq:pde2-weak-form}
 \int_0^T\!\int_\R
 \bigl(u(t,x)\varphi_t(t,x)+f(u(t,x))\varphi_x(t,x)\bigr)
 \,\mathrm d\lambda_1(x)\,\mathrm dt
 +\int_\R u_0(x)\varphi(0,x)\,\mathrm d\lambda_1(x)=0.
\end{equation}
Przestrzeń $C_c^\infty([0,T)\times\R)$ oznacza
restrykcje funkcji gładkich w pewnym otoczeniu
$[0,T)\times\R$, których nośniki są zwarte i nie dochodzą
do czasu $T$.
\end{definicja}

\begin{stwierdzenie}[Warunek Rankine'a--Hugoniota]
\label{prop:pde2-rh}
Niech $f\in C^1(\R)$ i $s\in C^1((0,T))$.
Załóżmy, że $u$ jest klasy $C^1$ po obu stronach
krzywej $x=s(t)$, z jednostronnymi przedłużeniami
klasy $C^1$ do krzywej, i posiada ślady
$u_L(t)$ i $u_R(t)$ na tej krzywej oraz spełnia
\eqref{eq:pde2-conservation} klasycznie poza nią.
Wówczas $u$ spełnia równanie słabo w otoczeniu
krzywej wtedy i tylko wtedy, gdy
\begin{equation}\label{eq:pde2-rh}
 s'(t)\bigl(u_R(t)-u_L(t)\bigr)
 =f(u_R(t))-f(u_L(t))
\end{equation}
dla każdego $t\in(0,T)$, przy ciągłości śladów.
\end{stwierdzenie}
\begin{proof}
Wybierzmy funkcję testową $\varphi$ o nośniku
w wąskim otoczeniu rozważanego fragmentu krzywej.
Rozdzielmy całkę w słabym sformułowaniu na obszary
$x<s(t)$ oraz $x>s(t)$. Całkowanie przez części
w obu obszarach, z wykorzystaniem równania klasycznego,
usuwa wyrazy objętościowe. Pozostają tylko wyrazy
na krzywej. Ich suma wynosi
\[
 \int_0^T\bigl(s'(t)[u](t)-[f(u)](t)\bigr)
 \varphi(t,s(t))\,\mathrm dt,
\]
gdzie $[u]=u_R-u_L$ oraz $[f(u)]=f(u_R)-f(u_L)$.
Ślad funkcji testowej na krzywej można zadawać dowolnie
bez zakładania gładkości samej krzywej ponad klasę $C^1$.
Niech $\eta\in C_c^\infty((0,T))$ i niech
$J\Subset(0,T)$ zawiera jej nośnik. Zwartość $s(J)$
pozwala wybrać $\chi\in C_c^\infty(\R)$, równą $1$
w otoczeniu $s(J)$. Wtedy
$\varphi(t,x)=\eta(t)\chi(x)$ jest klasy $C_c^\infty$
i spełnia $\varphi(t,s(t))=\eta(t)$ dla $t\in J$.
Zatem dla ciągłej funkcji
$b(t)=s'(t)[u](t)-[f(u)](t)$ zachodzi
$\int_0^T b(t)\eta(t)\,\mathrm dt=0$ dla każdego takiego $\eta$.
Gdyby $b(t_0)\ne0$, ciągłość dawałaby przedział, na którym
$b$ ma stały, niezerowy znak; test nieujemny o nośniku w tym
przedziale prowadziłby do sprzeczności. Stąd $b=0$
w każdym punkcie $(0,T)$.
Warunek rozstrzyga wyłącznie poprawność konserwatywną;
do jednoznacznego wyboru skoku potrzebna jest jeszcze
nierówność entropijna.
\end{proof}

\begin{przyklad}[Skok Burgersa]
Dla $f(u)=u^2/2$ i stałych stanów po obu stronach
fali $u_L\ne u_R$ warunek \eqref{eq:pde2-rh} daje
$\sigma=s'=(u_L+u_R)/2$.
Przy $u_L>u_R$ zachodzi $u_L>\sigma>u_R$:
charakterystyki z obu stron wpadają w falę uderzeniową.
Jeżeli $u_L<u_R$, taka nieciągłość spełnia wprawdzie
warunek Rankine'a--Hugoniota, ale nie jest dopuszczalnym
rozwiązaniem entropijnym.
\end{przyklad}

\begin{przyklad}[Fala rozrzedzenia]
\label{ex:pde2-rarefaction}
Dla równania Burgersa z danymi Riemanna
$u_0(x)=u_L$ dla $x<0$ oraz $u_0(x)=u_R$ dla $x>0$,
przy $u_L<u_R$, funkcja
\[
 u(t,x)=
 \begin{cases}
 u_L,&x\leq u_Lt,\\
 x/t,&u_Lt<x<u_Rt,\\
 u_R,&x\geq u_Rt
 \end{cases}
 \qquad(t>0)
\]
jest ciągła na każdym przekroju $t>0$ i spełnia
równanie Burgersa klasycznie w każdym z trzech
otwartych obszarów. Bezpośrednio sprawdzamy, że
w środkowej części $u_t=-x/t^2$, $u_x=1/t$,
zatem $u_t+uu_x=0$. Ciągłość na liniach
$x=u_Lt$ i $x=u_Rt$ sprawia, że nie występują tam
wkłady skokowe; dla testów zwartych w czasoprzestrzeni
całkowanie przez części potwierdza
\eqref{eq:pde2-weak-form}. W miarę $t\downarrow0$
rozwiązanie przyjmuje dane $u_0$ w $L^1_{\mathrm{loc}}$.
To rozwiązanie jest entropijne; precyzyjną definicję
warunków entropijnych podamy w rozdziale o prawach zachowania.
\end{przyklad}

\begin{uwaga}[Dlaczego samo rozwiązanie słabe nie wystarcza?]
Dla danych z przykładu \ref{ex:pde2-rarefaction}
można też skonstruować skok poruszający się z prędkością
$(u_L+u_R)/2$. Spełnia on ten sam całkowy warunek
\eqref{eq:pde2-weak-form}, lecz narusza warunek entropijny.
Dla tego samego zagadnienia mamy więc co najmniej
dwa różne rozwiązania słabe. Selekcja rozwiązania
fizycznego wymaga dodatkowej informacji.
\end{uwaga}

\section{Równania w pełni nieliniowe i układ Charpita}
\label{sec:pde2-charpit}

Dla ogólnego równania pierwszego rzędu
\begin{equation}\label{eq:pde2-general-F}
 F(x,u(x),\nabla u(x))=0,
\end{equation}
gdzie $x\in\Omega\subset\R^n$,
$F\colon\Omega\times\R\times\R^n\to\R$,
rozpatrujemy dodatkowe zmienne $p\in\R^n$
reprezentujące gradient. Pojawia się
\emph{układ charakterystyczny Charpita}:
\begin{equation}\label{eq:pde2-charpit-system}
 \dot X=F_p(X,U,P),\qquad
 \dot U=P\cdot F_p(X,U,P),\qquad
 \dot P=-F_x(X,U,P)-F_u(X,U,P)P.
\end{equation}
Tutaj $F_x,F_p$ są gradientami po odpowiednich
zmiennych, a $F_u$ jest pochodną skalarną.

\begin{lemat}[Zachowanie równania i warunku styczności]
\label{lem:pde2-contact}
Niech $F\in C^2$ i niech $(X,U,P)$ będzie rozwiązaniem
\eqref{eq:pde2-charpit-system} na przedziale $I$.
Wtedy $s\mapsto F(X(s),U(s),P(s))$ jest stała.
Jeśli rozwiązania zależą klasy $C^1$ od parametru
$y\in V\subset\R^{n-1}$ i $F(X,U,P)=0$, to
\begin{equation}\label{eq:pde2-contact-ode}
 \frac{\mathrm d}{\mathrm ds}
 \bigl(U_{y_j}-P\cdot X_{y_j}\bigr)
 =-F_u(X,U,P)\bigl(U_{y_j}-P\cdot X_{y_j}\bigr)
\end{equation}
dla każdego $j=1,\ldots,n-1$.
\end{lemat}
\begin{proof}
Reguła łańcuchowa daje
\[
 \frac{\mathrm d}{\mathrm ds}F(X,U,P)
 =F_x\cdot F_p+F_u(P\cdot F_p)
 +F_p\cdot(-F_x-F_uP)=0.
\]
Różniczkując $R_j=U_{y_j}-P\cdot X_{y_j}$ po $s$,
wykorzystując \eqref{eq:pde2-charpit-system}, otrzymujemy
\[
 \dot R_j=P_{y_j}\cdot F_p+(F_x+F_uP)\cdot X_{y_j}.
\]
Ponieważ $F(X,U,P)=0$ dla każdego $(s,y)$,
mamy także
$F_x\cdot X_{y_j}+F_u U_{y_j}+F_p\cdot P_{y_j}=0$.
Po podstawieniu do wzoru na $\dot R_j$ dostajemy
$\dot R_j=-F_u(U_{y_j}-P\cdot X_{y_j})$, czyli
\eqref{eq:pde2-contact-ode}.
\end{proof}

\begin{twierdzenie}[Lokalna konstrukcja Charpita]
\label{thm:pde2-charpit}
Niech $n\geq2$, $\Omega\subset\R^n$ będzie otwarta,
$F\in C^2(\Omega\times\R\times\R^n)$,
$V\subset\R^{n-1}$ otwarty oraz
$\gamma\in C^2(V;\Omega)$ regularną parametryzacją
hiperpowierzchni. Niech $g\in C^2(V)$ i
$p_0\in C^1(V;\R^n)$ spełniają dla wszystkich $y\in V$
\begin{equation}\label{eq:pde2-charpit-initial}
 F(\gamma(y),g(y),p_0(y))=0,\qquad
 p_0(y)\cdot\partial_{y_j}\gamma(y)=\partial_{y_j}g(y)
 \quad(j=1,\ldots,n-1).
\end{equation}
Jeżeli w $y_0\in V$ wektor
$F_p(\gamma(y_0),g(y_0),p_0(y_0))$
nie należy do $T_{\gamma(y_0)}\Sigma$,
gdzie $\Sigma=\gamma(V)$ lokalnie, to istnieje
otoczenie $U$ punktu $\gamma(y_0)$ i rozwiązanie
$u\in C^2(U)$ równania \eqref{eq:pde2-general-F}
przyjmujące dane $u(\gamma(y))=g(y)$ w pobliżu $y_0$.
Jego gradient na powierzchni danych jest równy $p_0$.
W klasie rozwiązań $C^2$, które realizują te same
dane $g,p_0$, konstrukcja jest lokalnie jednoznaczna.
\end{twierdzenie}
\begin{proof}
Rozwiążmy układ \eqref{eq:pde2-charpit-system}
z danymi $(X,U,P)(0,y)=(\gamma(y),g(y),p_0(y))$.
Zależność rozwiązań ODE od parametrów daje funkcje
klasy $C^1$ względem $(s,y)$ w pobliżu $(0,y_0)$.
Lemat \ref{lem:pde2-contact} i pierwszy warunek
\eqref{eq:pde2-charpit-initial} pokazują, że
$F(X,U,P)=0$; drugi warunek i równanie
\eqref{eq:pde2-contact-ode} dają
$U_{y_j}=P\cdot X_{y_j}$ przez jednoznaczność
skalarnego równania liniowego. Ponadto z samego
układu wynika $U_s=P\cdot X_s$.

W punkcie $(0,y_0)$ kolumny macierzy $D_{(s,y)}X$
są równe $F_p$ i $\partial_{y_j}\gamma$;
warunek poprzeczności zapewnia, że macierz ta jest
odwracalna. Z twierdzenia o funkcji odwrotnej $X$
stanowi lokalny układ współrzędnych w $\Omega$.
Definiujemy $u(X(s,y))=U(s,y)$. Wspomniane wcześniej
równości są równoważne tożsamości
$\mathrm dU=P\cdot\mathrm dX$, zatem
$\nabla u(X(s,y))=P(s,y)$. Ponieważ $P$ i lokalna
odwrotność $X^{-1}$ są klasy $C^1$, gradient $\nabla u$
jest klasy $C^1$ na $U$, a więc $u\in C^2(U)$.
Ponieważ $F(X,U,P)=0$, otrzymujemy równanie
\eqref{eq:pde2-general-F}.

Jeśli $u\in C^2$ jest innym rozwiązaniem z tym samym
gradientem początkowym, to wzdłuż odpowiednich
charakterystyk jego wartości i pochodne spełniają
\eqref{eq:pde2-charpit-system}: wystarczy zróżniczkować
$F(x,u,\nabla u)=0$ i skorzystać z symetrii hesjanu.
Jednoznaczność ODE implikuje zgodność obu rozwiązań
na wspólnym otoczeniu punktu początkowego.
\end{proof}

\begin{przyklad}[Eikonal i wybór gałęzi]
Niech $F(x,u,p)=|p|^2-1$ w $\R^2$.
Na prostej $\Sigma=\{(x,0):x\in\R\}$ zadamy
$u(x,0)=0$. Zgodność gradientu stycznego wymusza
$p_x=0$, a równanie wymusza $p_y=\pm1$.
Są zatem dwie dopuszczalne gałęzie danych gradientu,
które prowadzą do rozwiązań $u(x,y)=y$ i $u(x,y)=-y$.
Oba są niecharakterystyczne dla odpowiedniego wektora
$F_p=2p$. Samo zadanie wartości na prostej nie określa
więc jednoznacznie rozwiązania w pełni nieliniowego.
\end{przyklad}

\section{Równanie Hamiltona--Jacobiego i charakterystyki Hamiltona}
\label{sec:pde2-hj}

Ważną szczególną postacią równania pierwszego rzędu jest
\begin{equation}\label{eq:pde2-hj}
 u_t+H(x,\nabla_xu)=0,
 \qquad u(0,x)=u_0(x).
\end{equation}
W tej sekcji $H\in C^2(\R^d\times\R^d)$ nie zależy
jawnie od czasu ani od wartości $u$.

\begin{stwierdzenie}[Metoda charakterystyk Hamiltona]
\label{prop:pde2-hamilton}
Niech $u_0\in C^2(\R^d)$, $H\in C^2(\R^d\times\R^d)$
i niech $\xi_0\in\R^d$. Rozwiążmy lokalnie układ
\begin{equation}\label{eq:pde2-hamilton-odes}
 \dot X=H_p(X,P),\qquad
 \dot P=-H_x(X,P),\qquad
 \dot U=P\cdot H_p(X,P)-H(X,P),
\end{equation}
z warunkami $X(0;\xi)=\xi$,
$P(0;\xi)=\nabla u_0(\xi)$ i $U(0;\xi)=u_0(\xi)$.
Dla dostatecznie małych $|t|$ i $\xi$ bliskich $\xi_0$
mapa $\xi\mapsto X(t;\xi)$ jest lokalnie odwracalna.
W tych współrzędnych
$u(t,X(t;\xi))=U(t;\xi)$ definiuje rozwiązanie
klasy $C^1$ równania \eqref{eq:pde2-hj}, a ponadto
$\nabla_xu(t,X(t;\xi))=P(t;\xi)$.
\end{stwierdzenie}
\begin{proof}
Istnienie ODE i lokalna odwracalność wynikają jak
w twierdzeniu \ref{thm:pde2-quasilinear}, ponieważ
$D_\xi X(0;\xi)=I$. Ustalmy współrzędną $\xi_j$
i oznaczmy $R_j=U_{\xi_j}-P\cdot X_{\xi_j}$.
Korzystając z \eqref{eq:pde2-hamilton-odes} oraz
reguły łańcuchowej, obliczamy
\begin{align*}
 \dot R_j&=\partial_{\xi_j}(P\cdot H_p-H)
 -\dot P\cdot X_{\xi_j}-P\cdot\partial_{\xi_j}H_p\\
 &=P_{\xi_j}\cdot H_p-H_x\cdot X_{\xi_j}
 -H_p\cdot P_{\xi_j}+H_x\cdot X_{\xi_j}=0.
\end{align*}
Początkowo
$R_j(0,\xi)=\partial_{\xi_j}u_0-\nabla u_0\cdot e_j=0$,
zatem $R_j=0$. Wynika stąd
$\nabla_xu(t,X(t;\xi))=P(t;\xi)$. Z kolei pochodna
$u(t,X(t;\xi))=U(t;\xi)$ po $t$ daje
$u_t+\nabla_xu\cdot H_p=P\cdot H_p-H$, a więc
$u_t+H=0$, jak należało wykazać.
\end{proof}

\begin{przyklad}[Hamiltonian kinetyczny i powstawanie kaustyki]
\label{ex:pde2-hj-blowup}
Niech $d=1$ i $H(p)=p^2/2$. Wtedy
$\dot P=0$, $\dot X=P$ oraz $\dot U=P^2/2$.
W konsekwencji
\[
 X(t;\xi)=\xi+t u_0'(\xi),\qquad
 P(t;\xi)=u_0'(\xi),\qquad
 U(t;\xi)=u_0(\xi)+\frac t2|u_0'(\xi)|^2.
\]
Dla $u_0(x)=-x^2/2$ otrzymujemy
$X(t;\xi)=(1-t)\xi$. Jeśli $0\leq t<1$, to
\[
 u(t,x)=-\frac{x^2}{2(1-t)},\qquad
 u_x(t,x)=-\frac{x}{1-t}.
\]
W chwili $t=1$ mapa $\xi\mapsto X(t;\xi)$ traci
odwracalność, a dla ustalonego $x\ne0$ gradient
rozwiązania staje się nieograniczony, gdy $t\uparrow1$.
Zjawisko utraty odwracalności rzutu
$(X,P)\mapsto X$ nazywa się w tym kontekście
powstawaniem \emph{kaustyki}. Rozwiązania lepkościowe
umożliwią później kontynuację odpowiednich zadań
Hamiltona--Jacobiego poza pierwszą kaustyką.
\end{przyklad}

\section{Podsumowanie: kiedy metoda charakterystyk działa?}
\label{sec:pde2-summary}

Metoda charakterystyk wykorzystuje zależność
\eqref{eq:pde2-chain}, ale jej wynik zależy od struktury
równania:
\begin{enumerate}
\item Dla transportu liniowego z wystarczająco regularnym
polem prędkości charakterystyki tworzą przepływ
dyfeomorfizmów i zapewniają globalne rozwiązanie klasyczne
na skończonym przedziale czasu.
\item Dla równań quasiliniowych lokalne rozwiązanie istnieje,
dopóki układ charakterystyczny daje odwracalną mapę
początkowych etykiet $\xi$ na położenia $x$.
\item Dla praw zachowania po utracie gładkości naturalne
stają się rozwiązania słabe. Sama słaba postać równania
nie wystarcza zwykle do jednoznaczności; potrzebne są
warunki entropijne.
\item Dla równań w pełni nieliniowych trzeba śledzić
również gradient $p=\nabla u$. Warunek poprzeczności
$F_p$ względem powierzchni danych jest kluczowy,
a rozwiązanie może istnieć jedynie lokalnie nawet
przy bardzo gładkich danych.
\end{enumerate}

W rozdziale następnym przejdziemy do równania Laplace'a,
teorii funkcji harmonicznych i zagadnień Dirichleta.

\chapter[Równania Laplace'a i Poissona]{Równania Laplace'a i Poissona: funkcje harmoniczne i potencjał}
\label{chap:pde-laplace}

W tym rozdziale badamy dwa równania na zbiorze otwartym $\Omega\subset\R^n$:
$\Delta u=0$ oraz $-\Delta u=f$. Pierwsze opisuje między innymi
stacjonarną temperaturę i potencjały bez źródeł; drugie uwzględnia
źródła. W przeciwieństwie do równania transportu nie wystarcza
znajomość danych na dowolnej hiperpowierzchni. Naturalnym zadaniem
jest zwykle problem Dirichleta na \emph{całym} brzegu.

Przez $\lambda_n$ oznaczamy miarę Lebesgue'a w $\R^n$, przez
$\sigma$ miarę powierzchniową na gładkiej hiperpowierzchni,
a przez $\omega_n=\sigma(S^{n-1})$ pole jednostkowej sfery.
W szczególności $\lambda_n(B_r(0))=\omega_n r^n/n$.
Znak laplasjanu ustalamy jako $\Delta=\sum_{j=1}^n\partial_{x_j}^2$;
rozwiązanie fundamentalne będziemy normalizować przez $-\Delta\Phi=\delta_0$.

\section{Własność wartości średniej}
\label{sec:pde3-mean}

\begin{definicja}[Funkcja harmoniczna]
Niech $\Omega\subset\R^n$ będzie niepustym zbiorem otwartym.
Funkcję rzeczywistą $u\in C^2(\Omega)$ nazywamy \emph{harmoniczną}
na $\Omega$, jeśli $\Delta u=0$ w każdym punkcie $\Omega$.
Funkcję $u\in C^2(\Omega)$ nazywamy \emph{subharmoniczną},
gdy $\Delta u\geq0$, oraz \emph{superharmoniczną},
gdy $\Delta u\leq0$. Później rozszerzymy te pojęcia
na funkcje bez klasycznych drugich pochodnych.
\end{definicja}

Dla $x_0\in\Omega$ i $0<r<\operatorname{dist}(x_0,\partial\Omega)$
zapisujemy średnią sferyczną
\begin{equation}\label{eq:pde3-sphere-mean}
 M_u(x_0,r)=\frac{1}{\omega_n r^{n-1}}
 \int_{\partial B_r(x_0)}u(y)\,\mathrm d\sigma(y)
 =\frac1{\omega_n}\int_{S^{n-1}}u(x_0+r\theta)\,
 \mathrm d\sigma(\theta).
\end{equation}
Jeśli $\Omega=\R^n$, przyjmujemy $\operatorname{dist}(x_0,\partial\Omega)=+\infty$.

\begin{twierdzenie}[Własność wartości średniej]
\label{thm:pde3-mean}
Niech $n\geq2$, $\Omega\subset\R^n$ będzie otwarta,
$u\in C^2(\Omega)$ i $\Delta u=0$ w $\Omega$.
Dla każdego $x_0\in\Omega$ oraz $r>0$ takich, że
$\overline{B_r(x_0)}\subset\Omega$, zachodzi
\begin{equation}\label{eq:pde3-both-means}
 u(x_0)=\frac{1}{\omega_n r^{n-1}}
 \int_{\partial B_r(x_0)}u(y)\,\mathrm d\sigma(y)
 =\frac{1}{\lambda_n(B_r(x_0))}
 \int_{B_r(x_0)}u(y)\,\mathrm d\lambda_n(y).
\end{equation}
\end{twierdzenie}
\begin{proof}
Rozważmy $M(r)=M_u(x_0,r)$. Różniczkując drugą
postać we wzorze~\eqref{eq:pde3-sphere-mean} i używając
wzoru Gaussa, otrzymujemy
\[
 M'(r)=\frac{1}{\omega_n r^{n-1}}
 \int_{\partial B_r(x_0)}\partial_\nu u(y)\,
 \mathrm d\sigma(y)
 =\frac{1}{\omega_n r^{n-1}}
 \int_{B_r(x_0)}\Delta u(y)\,\mathrm d\lambda_n(y)=0.
\]
Zatem $M$ jest stała. Z ciągłości $u$ mamy
$\lim_{r\downarrow0}M(r)=u(x_0)$, co daje średnią sferyczną.
Całkując tę tożsamość w zmiennej promieniowej z wagą
$\omega_n s^{n-1}$ od $0$ do $r$ i korzystając ze
współrzędnych biegunowych, otrzymujemy średnią po kuli.
\end{proof}

\begin{stwierdzenie}[Odwrotność własności średniej]
\label{prop:pde3-converse}
Niech $n\geq2$, $\Omega\subset\R^n$ będzie otwarta,
a $u\in C^2(\Omega)$. Jeśli dla każdego $x\in\Omega$
i wszystkich dostatecznie małych $r>0$ spełniona jest
pierwsza równość w~\eqref{eq:pde3-both-means}, to
$\Delta u=0$ w $\Omega$. Ten sam wniosek zachodzi przy
założeniu własności średniej po kulach.
\end{stwierdzenie}
\begin{proof}
Wzór Taylora w punkcie $x$ daje jednostajnie względem
$\theta\in S^{n-1}$:
$u(x+r\theta)=u(x)+r\nabla u(x)\cdot\theta+
\frac{r^2}{2}\theta^{\mathsf T}D^2u(x)\theta+o(r^2)$.
Przez symetrię sfery całki z $\theta_i$ znikają,
a $\omega_n^{-1}\int_{S^{n-1}}\theta_i\theta_j\,
\mathrm d\sigma(\theta)=\delta_{ij}/n$.
Stąd $M_u(x,r)=u(x)+\frac{r^2}{2n}\Delta u(x)+o(r^2)$.
Jeśli $M_u(x,r)=u(x)$, dzielimy przez $r^2$ i przechodzimy
do granicy. Wersja dla kul daje współczynnik
$r^2/(2(n+2))$ zamiast $r^2/(2n)$, a argument jest taki sam.
\end{proof}

\begin{uwaga}[Dlaczego średnia po kulach nie jest przypadkowa?]
Wartość funkcji harmonicznej w środku jest dokładnie
równowagą jej wartości wokół środka. Dla funkcji
subharmonicznej o klasie $C^2$ rachunek z dowodu pokazuje,
że $M_u(x,r)$ jest niemalejąca jako funkcja $r$, więc
$u(x)\leq M_u(x,r)$. Dla funkcji superharmonicznej
nierówność ma kierunek przeciwny.
\end{uwaga}

\begin{twierdzenie}[Własność wartości średniej bez założenia $C^2$]
\label{thm:pde-v31-continuous-mean}
Niech $n\ge2$, $\Omega\subset\R^n$ będzie otwarta,
a $u\in C(\Omega)$. Załóżmy, że dla każdego
$x\in\Omega$ i każdego $r>0$ z
$\overline{B_r(x)}\subset\Omega$ zachodzi
\[
 u(x)=\frac1{\omega_nr^{n-1}}
 \int_{\partial B_r(x)}u(y)\,\mathrm d\sigma(y).
\]
Wówczas $u\in C^\infty(\Omega)$ i $\Delta u=0$ w $\Omega$.
\end{twierdzenie}
\begin{proof}
Ustalmy kulę $\overline{B_{2R}(x_0)}\subset\Omega$.
Wybierzmy nieujemną gładką funkcję $\rho$ o nośniku
w $B_1(0)$ i całce $\int_{\R^n}\rho(z)\,
\mathrm d\lambda_n(z)=1$. Dla $0<\eps<R/2$ definiujemy
lokalnie
$u_\eps(x)=\int_{\R^n}\rho_\eps(z)u(x-z)\,
\mathrm d\lambda_n(z)$ na $B_{3R/2}(x_0)$,
gdzie $\rho_\eps(z)=\eps^{-n}\rho(z/\eps)$.
Funkcja $u_\eps$ jest gładka. Dla $x\in B_R(x_0)$
i dostatecznie małego $r>0$ zamiana kolejności całek
i założona własność średniej dla punktów $x-z$
prowadzą do
\[
 \frac1{\omega_nr^{n-1}}
 \int_{\partial B_r(x)}u_\eps(y)\,\mathrm d\sigma(y)
 =\int_{\R^n}\rho_\eps(z)u(x-z)\,
 \mathrm d\lambda_n(z)=u_\eps(x).
\]
Użycie twierdzenia Fubiniego jest dopuszczalne,
ponieważ całkujemy funkcję ciągłą po zbiorach zwartych.
Z odwrotności własności średniej dla funkcji klasy $C^2$
(stwierdzenie~\ref{prop:pde3-converse}) wynika
$\Delta u_\eps=0$ na $B_R(x_0)$.

Z ciągłości $u$ wynika $u_\eps\to u$ jednostajnie na
$\overline{B_R(x_0)}$. Stosujemy wzór Poissona
(twierdzenie~\ref{thm:pde3-poisson}) do każdej funkcji
$u_\eps$ na kuli $B_R(x_0)$:
\[
 u_\eps(x)=\int_{\partial B_R(x_0)}P_R(x-x_0,\zeta-x_0)
 u_\eps(\zeta)\,\mathrm d\sigma(\zeta),
 \qquad x\in B_R(x_0).
\]
Dodatniość i normalizacja jądra umożliwiają przejście
do granicy w powyższej całce. W konsekwencji $u$ jest
dokładnie rozwiązaniem Poissona z ciągłymi danymi
$u|_{\partial B_R(x_0)}$, a więc jest harmoniczna
i gładka na $B_R(x_0)$. Dowolność $x_0$ kończy dowód.
\end{proof}

\begin{uwaga}
Kluczowe jest, że własność średniej zakładamy dla
\emph{wszystkich} kul o domknięciu zawartym w $\Omega$.
Nie należy bez dodatkowego argumentu zastępować tego
warunku wymaganiem dotyczącym jednego promienia
zależnego od środka kuli.
\end{uwaga}

\section{Mocna zasada maksimum, jednoznaczność i twierdzenie Liouville'a}
\label{sec:pde3-max}

Słabą zasadę maksimum udowodniliśmy już w twierdzeniu
\ref{thm:laplace-max}. Własność wartości średniej daje
twierdzenie silniejsze, które nie wymaga ograniczoności obszaru.

\begin{twierdzenie}[Mocna zasada maksimum dla funkcji harmonicznych]
\label{thm:pde3-strong-max}
Niech $\Omega\subset\R^n$ będzie obszarem, $n\geq2$,
a $u\in C^2(\Omega)$ funkcją harmoniczną. Jeśli
$u$ osiąga maksimum lokalne albo minimum lokalne w pewnym
punkcie $x_0\in\Omega$, to jest stała na $\Omega$.
\end{twierdzenie}
\begin{proof}
Niech $M=u(x_0)$ będzie maksimum lokalnym. Dla pewnego
$r_0>0$ funkcja $u$ nie przekracza $M$ na $B_{r_0}(x_0)$.
Dla każdego $0<r<r_0$ jej średnia po sferze
$\partial B_r(x_0)$ równa się $M$. Ponieważ $u$ jest
ciągła i nie przekracza $M$, musi być równa $M$ na
każdej z tych sfer, a więc na całej kuli $B_{r_0}(x_0)$.

Aby rozciągnąć stałość na cały obszar,
korzystamy z twierdzenia o analityczności funkcji
harmonicznych (twierdzenie~\ref{thm:pde3-analytic}),
udowodnionego niżej niezależnie od obecnego wyniku.
Dla rzeczywiście analitycznej funkcji na spójnym zbiorze
otwartym równość z funkcją stałą na niepustym otwartym
podzbiorze wymusza taką samą równość wszędzie:
można to wykazać, przechodząc wzdłuż drogi przez
skończenie wiele nakładających się kul, na których
zbieżne szeregi Taylora wyznaczają funkcję jednoznacznie.
Dokładniej, jeśli analityczna funkcja jest zerowa na
niepustej części pierwszej kuli, wszystkie jej pochodne
w punktach tej części są zerowe, a jej szereg Taylora
przedstawia funkcję zerową w otoczeniu takiego punktu.
Przenosząc ten argument po nakładających się kulach,
uzyskujemy zerowanie się na całym spójnym obszarze.
Wskazane późniejsze twierdzenie o analityczności
korzysta tylko z reprezentacji Poissona i nie odwołuje
się do mocnej zasady maksimum.
Zatem $u\equiv M$ na $\Omega$. Przypadek minimum
sprowadza się do maksimum funkcji $-u$.
\end{proof}

\begin{stwierdzenie}[Oszacowanie pochodnych w środku kuli]
\label{prop:pde3-gradient}
Niech $n\geq2$, $R>0$, $x_0\in\R^n$,
a $u\in C^2(B_R(x_0))\cap C(\overline{B_R(x_0)})$
będzie harmoniczna. Wówczas dla każdego
$j=1,\dots,n$ zachodzi
\begin{equation}\label{eq:pde3-derivative-bound}
 |\partial_{x_j}u(x_0)|\leq \frac nR
 \sup_{y\in B_R(x_0)}|u(y)|.
\end{equation}
\end{stwierdzenie}
\begin{proof}
Pochodna $\partial_{x_j}u$ jest harmoniczna, gdy
$u$ ma trzecie pochodne, których jeszcze nie wiemy.
Dlatego nie zakładamy tego: różniczkujemy zamiast tego
średnią objętościową z ruchomym środkiem
$u(x)=\lambda_n(B_R)^{-1}\int_{B_R(0)}u(x+z)\,
\mathrm d\lambda_n(z)$, początkowo dla $x$ bliskich $x_0$
i promienia trochę mniejszego od $R$. Zatem
$\partial_{x_j}u(x_0)=\lambda_n(B_\rho)^{-1}
\int_{B_\rho(x_0)}\partial_{x_j}u(y)\,
\mathrm d\lambda_n(y)$ dla każdego $\rho<R$.
Twierdzenie Gaussa daje
\[
 \partial_{x_j}u(x_0)=\frac{1}{\lambda_n(B_\rho)}
 \int_{\partial B_\rho(x_0)}u(y)\nu_j(y)\,
 \mathrm d\sigma(y).
\]
Z $|\nu_j|\leq1$ i
$\sigma(\partial B_\rho)/\lambda_n(B_\rho)=n/\rho$
wynika oszacowanie z $\rho$ w mianowniku.
Przechodząc do granicy $\rho\uparrow R$ dostajemy
\eqref{eq:pde3-derivative-bound}.
\end{proof}

\begin{wniosek}[Twierdzenie Liouville'a]
\label{cor:pde3-liouville}
Niech $n\geq2$ i $u\in C^2(\R^n)$ będzie harmoniczna
oraz ograniczona. Wtedy $u$ jest stała.
\end{wniosek}
\begin{proof}
Dla każdego $x_0\in\R^n$ i $R>0$ stosujemy
\eqref{eq:pde3-derivative-bound}. Ograniczoność $u$
pozwala przejść z $R$ do nieskończoności, więc
$\nabla u(x_0)=0$. Ponieważ $x_0$ był dowolny,
a $\R^n$ jest spójna, funkcja $u$ jest stała.
\end{proof}

\begin{kontrprzyklad}[Dlaczego ograniczoność jest ważna?]
Funkcja $u(x)=x_1$ jest harmoniczna na $\R^n$,
lecz nie jest stała. Twierdzenie Liouville'a nie
obowiązuje bez ograniczoności (ani analogicznej
kontroli wzrostu).
\end{kontrprzyklad}

\section{Rozwiązanie fundamentalne laplasjanu}
\label{sec:pde3-fundamental}

Rozwiązaniem fundamentalnym operatora nazywa się funkcję lub dystrybucję,
która reaguje na punktowe źródło. Wprowadzimy ją najpierw
bez formalnej teorii dystrybucji: zapis
$-\Delta\Phi=\delta_0$ będzie oznaczał precyzyjną
tożsamość całkową, gdzie $\delta_0$ jest funkcjonałem
$\varphi\mapsto\varphi(0)$ na funkcjach testowych.
Tutaj $C_c^\infty(\R^n)$ oznacza przestrzeń funkcji
nieskończenie różniczkowalnych o zwartym nośniku, a
$L^1_{\mathrm{loc}}(\R^n)$ klasę funkcji całkowalnych
na każdym zwartym podzbiorze $\R^n$.

\begin{twierdzenie}[Rozwiązanie fundamentalne $-\Delta$]
\label{thm:pde3-phi}
Niech $n\geq2$ i $\omega_n=\sigma(S^{n-1})$. Dla
$x\in\R^n\setminus\{0\}$ połóżmy
\begin{equation}\label{eq:pde3-fundamental}
 \Phi_n(x)=
 \begin{cases}
 \displaystyle\frac{1}{(n-2)\omega_n}|x|^{2-n},&n\geq3,\\[4pt]
 \displaystyle-\frac1{2\pi}\log|x|,&n=2.
 \end{cases}
\end{equation}
Wówczas $\Phi_n\in L^1_{\mathrm{loc}}(\R^n)$,
$\Delta\Phi_n=0$ na $\R^n\setminus\{0\}$ oraz
\begin{equation}\label{eq:pde3-distdelta}
 \int_{\R^n}\Phi_n(x)(-\Delta\varphi(x))\,
 \mathrm d\lambda_n(x)=\varphi(0)
 \quad\text{dla każdej }\varphi\in C_c^\infty(\R^n).
\end{equation}
\end{twierdzenie}
\begin{proof}
Dla radialnej funkcji $w(x)=h(|x|)$ i $x\ne0$
obliczamy $\Delta w=h''(r)+(n-1)h'(r)/r$.
Rozwiązanie $h''+(n-1)h'/r=0$ daje $h'(r)=Cr^{1-n}$,
co prowadzi do potęg $r^{2-n}$ albo logarytmu.
Wybór stałej w~\eqref{eq:pde3-fundamental} zapewnia
$\partial_r\Phi_n(r)=-1/(\omega_n r^{n-1})$
($\omega_2=2\pi$). Osobliwość jest całkowalna lokalnie,
gdyż iloczyn $|\Phi_n(r)|r^{n-1}$ jest całkowalny w zerze.

Niech $\varphi$ ma nośnik w $B_R(0)$ i $0<\eps<R$.
Stosujemy drugą tożsamość Greena na pierścieniu
$A_{\eps,R}=B_R(0)\setminus\overline{B_\eps(0)}$.
Na zewnętrznym brzegu wszystkie składniki zawierające
$\varphi$ znikają. Na wewnętrznym brzegu normalna
zewnętrzna dla pierścienia jest skierowana do środka,
a zatem $\partial_\nu\Phi_n=1/(\omega_n\eps^{n-1})$.
Otrzymujemy
\[
 \int_{A_{\eps,R}}\Phi_n(y)(-\Delta\varphi(y))\,
 \mathrm d\lambda_n(y)
 =\int_{\partial B_\eps(0)}
 \bigl(\varphi(y)\,\partial_\nu\Phi_n(y)
       -\Phi_n(y)\,\partial_\nu\varphi(y)\bigr)\,\mathrm d\sigma(y).
\]
Pierwszy wyraz dąży do $\varphi(0)$; drugi znika,
ponieważ $\eps^{n-1}|\Phi_n(\eps)|\to0$.
Lokalna całkowalność pozwala przejść do granicy
również po lewej stronie. Otrzymujemy~\eqref{eq:pde3-distdelta}.
\end{proof}

\begin{uwaga}[Wymiar jeden]
Dla $n=1$ odpowiednikiem rozwiązania fundamentalnego
jest $\Phi_1(x)=-|x|/2$, ponieważ druga pochodna
$|x|$ w sensie dystrybucyjnym jest równa $2\delta_0$.
Przypadek jednowymiarowy jest w istocie teorią ODE,
dlatego w dalszych konstrukcjach potencjału przyjmujemy $n\geq2$.
\end{uwaga}

\begin{twierdzenie}[Potencjał Newtona i równanie Poissona]
\label{thm:pde3-newton}
Niech $n\geq2$ i $f\in C_c^\infty(\R^n)$. Określmy
\begin{equation}\label{eq:pde3-newton}
 (\Phi_n*f)(x)=\int_{\R^n}\Phi_n(x-y)f(y)\,
 \mathrm d\lambda_n(y).
\end{equation}
Wówczas $u=\Phi_n*f\in C^\infty(\R^n)$ i
$-\Delta u=f$ na $\R^n$.
\end{twierdzenie}
\begin{proof}
Po podstawieniu $z=x-y$ możemy zapisać
$u(x)=\int_{\R^n}\Phi_n(z)f(x-z)\,\mathrm d\lambda_n(z)$.
Dla $x$ z ustalonego zbioru zwartego funkcje
$z\mapsto f(x-z)$ mają nośniki zawarte we wspólnym
zbiorze zwartym. Dzięki lokalnej całkowalności
$\Phi_n$ można różniczkować całkę dowolnie wiele razy
przenosząc pochodne na $f$. Stąd $u\in C^\infty$.
Dla ustalonego $x$ wzór~\eqref{eq:pde3-distdelta}
zastosowany do funkcji testowej $z\mapsto f(x-z)$ daje
$\int_{\R^n}\Phi_n(z)(-\Delta f)(x-z)\,\mathrm d\lambda_n(z)=f(x)$,
czyli $-\Delta u(x)=f(x)$.
\end{proof}

\begin{uwaga}[Rozwiązanie szczególne a niejednoznaczność]
Jeżeli $u$ i $v$ są rozwiązaniami $-\Delta u=f$
na $\R^n$, ich różnica jest harmoniczna. Zatem potencjał
Newtona jest rozwiązaniem szczególnym; warunki wzrostu
w nieskończoności są potrzebne do rozstrzygnięcia
jednoznaczności. Przykładowo $u+x_1$ jest innym rozwiązaniem.
\end{uwaga}

\section{Funkcje Greena i tożsamości Greena}
\label{sec:pde3-green}

\begin{stwierdzenie}[Dwie tożsamości Greena]
\label{prop:pde3-green-identities}
Niech $\Omega\subset\R^n$ będzie obszarem ograniczonym
o brzegu klasy $C^1$ i niech $u,v\in C^2(\overline\Omega)$.
Wtedy
\begin{align}
 \int_\Omega u(x)\Delta v(x)\,\mathrm d\lambda_n(x)
 &=-\int_\Omega\nabla u(x)\cdot\nabla v(x)\,\mathrm d\lambda_n(x)
 +\int_{\partial\Omega}u(x)\,\partial_\nu v(x)\,\mathrm d\sigma(x),
 \label{eq:pde3-green1}\\
 \int_\Omega(u(x)\Delta v(x)-v(x)\Delta u(x))\,\mathrm d\lambda_n(x)
 &=\int_{\partial\Omega}(u(x)\partial_\nu v(x)-v(x)\partial_\nu u(x))
 \,\mathrm d\sigma(x).\label{eq:pde3-green2}
\end{align}
\end{stwierdzenie}
\begin{proof}
Dla pierwszej tożsamości stosujemy wzór Gaussa do
pola wektorowego $u\nabla v$ i korzystamy z
$\diver(u\nabla v)=\nabla u\cdot\nabla v+u\Delta v$.
Drugą otrzymujemy odejmując pierwszą tożsamość
z parami $(u,v)$ i $(v,u)$.
\end{proof}

\begin{definicja}[Funkcja Greena dla $-\Delta$]
\label{def:pde3-green}
Niech $\Omega\subset\R^n$ będzie obszarem ograniczonym,
$n\geq2$. Dla ustalonego $x\in\Omega$ funkcją Greena
(z warunkiem Dirichleta) nazywamy funkcję
$G_\Omega(x,\cdot)$ taką, że:
\begin{enumerate}
\item jest lokalnie całkowalna w $\Omega$ i klasy
$C^2$ na $\Omega\setminus\{x\}$;
\item $-\Delta_yG_\Omega(x,y)=\delta_x$ w sensie
równości całkowej~\eqref{eq:pde3-distdelta}, odpowiednio
przesuniętej do punktu $x$;
\item $G_\Omega(x,y)=0$ dla $y\in\partial\Omega$,
przy założeniu ciągłego przedłużenia z dala od $x$.
\end{enumerate}
W zastosowaniach do całek brzegowych wymagamy również
odpowiedniej gładkości funkcji Greena względem $y$ aż do brzegu.
\end{definicja}

\begin{twierdzenie}[Wzór reprezentacyjny Greena]
\label{thm:pde3-representation}
Niech $n\geq2$, $\Omega\subset\R^n$ będzie obszarem ograniczonym o brzegu klasy $C^2$, $f\in C(\overline\Omega)$
i $u\in C^2(\overline\Omega)$ spełnia $-\Delta u=f$
w $\Omega$. Załóżmy, że dla każdego $x\in\Omega$
istnieje funkcja Greena $G_\Omega(x,y)=\Phi_n(x-y)-h_x(y)$,
gdzie $h_x$ jest harmoniczna i klasy $C^2$ na pewnym
otoczeniu $\overline\Omega$ (w zmiennej $y$).
Wówczas dla $x\in\Omega$
\begin{equation}\label{eq:pde3-green-representation}
 u(x)=\int_\Omega G_\Omega(x,y)f(y)\,\mathrm d\lambda_n(y)
 -\int_{\partial\Omega}u(y)\partial_{\nu_y}G_\Omega(x,y)
 \,\mathrm d\sigma(y).
\end{equation}
\end{twierdzenie}
\begin{proof}
Ustalmy $x\in\Omega$ i wybierzmy
$\eps<\operatorname{dist}(x,\partial\Omega)$.
Druga tożsamość Greena dla $u$ i $G_\Omega(x,\cdot)$
na $\Omega_\eps=\Omega\setminus\overline{B_\eps(x)}$
daje, ponieważ $\Delta G_\Omega=0$ poza $x$,
\[
 \int_{\Omega_\eps}G_\Omega(x,y)f(y)\,\mathrm d\lambda_n(y)
 =\int_{\partial\Omega_\eps}
 \bigl(u(y)\partial_{\nu_y}G_\Omega(x,y)
 -G_\Omega(x,y)\partial_{\nu_y}u(y)\bigr)
 \,\mathrm d\sigma(y).
\]
Na zewnętrznym brzegu $G_\Omega=0$. Na wewnętrznym
brzegu część pochodząca od $h_x$ znika w granicy,
a dla $\Phi_n(x-y)$ pierwszy składnik dąży do $u(x)$,
podczas gdy drugi dąży do zera, jak w dowodzie
twierdzenia~\ref{thm:pde3-phi}. Po przejściu z
$\eps\downarrow0$ otrzymujemy
$\int_\Omega G_\Omega f=\int_{\partial\Omega}
u\,\partial_\nu G_\Omega+u(x)$,
co po przeniesieniu składnika brzegowego daje
\eqref{eq:pde3-green-representation}.
\end{proof}

\section{Jawne jądro Poissona i zagadnienie Dirichleta na kuli}
\label{sec:pde3-poisson-kernel}

W tym podrozdziale $n\geq2$, $R>0$,
$B_R=B_R(0)$ i $\zeta\in\partial B_R$.
Dla $y\in B_R\setminus\{0\}$ wprowadzamy punkt odwrócony
$y^*=R^2y/|y|^2$, leżący poza kulą. Zachodzi tożsamość
\begin{equation}\label{eq:pde3-kelvin-identity}
 |z-y^*|=\frac R{|y|}|z-y| \qquad (|z|=R).
\end{equation}

\begin{stwierdzenie}[Funkcja Greena dla kuli]
\label{prop:pde3-ball-green}
Dla $n\geq3$ i $x,y\in B_R$, $x\ne y$, $y\ne0$, funkcja
\begin{equation}\label{eq:pde3-ball-green}
 G_R(x,y)=\Phi_n(x-y)-\Bigl(\frac R{|y|}\Bigr)^{n-2}
 \Phi_n(x-y^*)
\end{equation}
jest funkcją Greena dla kuli w zmiennej $x$.
Dla $n=2$ odpowiednia funkcja ma postać
\[
 G_R(x,y)=\frac1{2\pi}\log\frac{|y|\,|x-y^*|}{R\,|x-y|}.
\]
W obu przypadkach formuły przedłużają się w punkcie $y=0$
przez granicę (dla $n\geq3$:
$G_R(x,0)=\Phi_n(x)-\Phi_n(R)$; dla $n=2$:
$G_R(x,0)=(2\pi)^{-1}\log(R/|x|)$).
Ponadto $G_R(x,y)=G_R(y,x)$.
\end{stwierdzenie}
\begin{proof}
Przy ustalonym $y$ pierwszy składnik jest rozwiązaniem
fundamentalnym o biegunie $y$, a drugi jest harmoniczny
w $x\in B_R$, bo jego biegun $y^*$ leży poza
$\overline{B_R}$. Na brzegu $|x|=R$ składniki znoszą
się na mocy~\eqref{eq:pde3-kelvin-identity};
dla $n=2$ zachodzi analogiczne znoszenie logarytmów.
Granice dla $y\to0$ wynikają z asymptotyki
$|x-y^*|\sim R^2/|y|$.

Dla symetrii można bezpośrednio przekształcić odległości:
\[
 |y|^2|x-y^*|^2=R^4-2R^2 x\cdot y+|x|^2|y|^2
 =|x|^2|y-x^*|^2.
\]
W połączeniu z symetrią $\Phi_n(x-y)$ i odpowiednimi
potęgami lub logarytmami daje to $G_R(x,y)=G_R(y,x)$.
\end{proof}

\begin{center}
\begin{tikzpicture}[scale=.87,>=Stealth]
 \fill[blue!5] (0,0) circle (2);
 \draw[thick] (0,0) circle (2);
 \draw[densely dashed,gray!65] (0,0)--(4.25,0);
 \fill (0,0) circle (1.6pt);
 \node[below left] at (0,0) {$0$};
 \fill[blue!75!black] (1.1,0) circle (2.5pt);
 \node[above] at (1.1,0) {$y$};
 \fill[red!75!black] (3.63,0) circle (2.5pt);
 \node[above] at (3.63,0) {$y^*$};
 \draw[<->,thin] (0,-.3)--(2,-.3);
 \node[below] at (1,-.3) {$R$};
 \node[align=left,anchor=west] at (4.2,1.0) {$|y|\,|y^*|=R^2$\\punkt odwrócony leży poza kulą};
\end{tikzpicture}
\par\smallskip\small
Inwersja w sferze $\partial B_R$: punkt $y$ wewnątrz
odpowiada punktowi $y^*$ na zewnątrz. Pozwala to
zbudować harmoniczną poprawkę rozwiązania fundamentalnego.
\end{center}

\begin{twierdzenie}[Jądro Poissona dla kuli]
\label{thm:pde3-poisson}
Niech $n\geq2$, $R>0$ i $g\in C(\partial B_R)$.
Dla $x\in B_R$ zdefiniujmy
\begin{equation}\label{eq:pde3-poisson}
 P_R(x,\zeta)=\frac{R^2-|x|^2}{\omega_n R|x-\zeta|^n},
 \qquad
 u(x)=\int_{\partial B_R}P_R(x,\zeta)g(\zeta)
 \,\mathrm d\sigma(\zeta).
\end{equation}
Wówczas $u\in C^\infty(B_R)\cap C(\overline{B_R})$,
$\Delta u=0$ w $B_R$ i $u|_{\partial B_R}=g$.
Jest to jedyne rozwiązanie z klasy
$C^2(B_R)\cap C(\overline{B_R})$ tego problemu Dirichleta.
Ponadto $P_R(x,\zeta)>0$ oraz
$\int_{\partial B_R}P_R(x,\zeta)\,\mathrm d\sigma(\zeta)=1$.
\end{twierdzenie}
\begin{proof}
Najpierw wykażemy związek z normalną pochodną funkcji
Greena. Dla $y\in B_R\setminus\{0\}$ i $|\zeta|=R$
korzystamy z relacji
$|\zeta-y^*|=(R/|y|)|\zeta-y|$.
Ponieważ $\nabla\Phi_n(z)=-z/(\omega_n|z|^n)$
(dla wszystkich $n\geq2$), dla $n\geq3$ otrzymujemy
\[
\begin{aligned}
-\partial_{\nu_\zeta}G_R(\zeta,y)
&=\frac{(\zeta-y)\cdot\zeta}{\omega_nR|\zeta-y|^n}
 -\Bigl(\frac R{|y|}\Bigr)^{n-2}
 \frac{(\zeta-y^*)\cdot\zeta}{\omega_nR|\zeta-y^*|^n}\\
&=\frac{(\zeta-y)\cdot\zeta-
 (|y|/R)^2(\zeta-y^*)\cdot\zeta}
 {\omega_nR|\zeta-y|^n}
 =\frac{R^2-|y|^2}{\omega_nR|\zeta-y|^n}.
\end{aligned}
\]
W przypadku $n=2$ różniczkujemy analogicznie
logarytmiczną postać $G_R$, otrzymując ten sam wzór.
W punkcie $y=0$ tożsamość wynika przez ciągłość.
Korzystając z symetrii $G_R(x,y)=G_R(y,x)$,
uzyskujemy więc
$P_R(x,\zeta)=-\partial_{\nu_\zeta}G_R(x,\zeta)$.
Uzasadnienie harmoniczności wymaga zachowania kolejności
operacji. Ustalmy zwarty zbiór $K\Subset B_R$ oraz punkt
$\zeta\in\partial B_R$. Dla $y$ w dostatecznie małym
otoczeniu $\zeta$ po wewnętrznej stronie brzegu,
oddzielonym od $K$, funkcja $x\mapsto G_R(x,y)$ jest
harmoniczna na otoczeniu $K$. Jawny wzór
\eqref{eq:pde3-ball-green} jest gładki względem
parametru $y$ w takim otoczeniu i pozwala najpierw
zróżniczkować po $y$ w kierunku normalnym, a dopiero
potem przejść do $y\to\zeta$ od wewnątrz.
Ponieważ pochodne względem $x$ i $y$ komutują,
otrzymujemy
$\Delta_xP_R(x,\zeta)=-\partial_{\nu_y}
\Delta_xG_R(x,y)|_{y=\zeta}=0$ dla $x\in K$.
Tu zapis wartości na brzegu oznacza granicę
\emph{pochodnej}, a nie pochodną funkcji
$G_R(x,\zeta)\equiv0$. Z nierówności
$|x-\zeta|\geq R-|x|>0$ wynika, że pochodne
względem $x$ można przenosić pod całkę na zwartych
podzbiorach kuli. Stąd $u$ jest gładka i harmoniczna.

Uzasadnijmy normalizację jądra bez powoływania się na
istnienie szukanego rozwiązania Dirichleta. Stosujemy
wzór reprezentacyjny~\eqref{eq:pde3-green-representation}
do funkcji stałej $1$, dla której $-\Delta1=0$.
Poprawka harmoniczna w funkcji Greena dla kuli
jest gładka na otoczeniu domkniętej kuli przy ustalonym
biegunie $x$, więc założenia wzoru są spełnione.
Wynika stąd $1=\int_{\partial B_R}P_R(x,\zeta)\,
\mathrm d\sigma(\zeta)$. Dodatniość jest jawna.

Ustalmy $\zeta_0\in\partial B_R$ i $\eps>0$.
Z jednostajnej ciągłości $g$ wybierzmy $\delta>0$
tak, że $|g(\zeta)-g(\zeta_0)|<\eps$, gdy
$|\zeta-\zeta_0|<\delta$. Posługując się dodatniością
i normalizacją $P_R$ dostajemy
\[
 |u(x)-g(\zeta_0)|\leq\eps+
 2\|g\|_\infty
 \int_{\substack{\zeta\in\partial B_R\\
 |\zeta-\zeta_0|\geq\delta}}P_R(x,\zeta)\,
 \mathrm d\sigma(\zeta).
\]
Gdy $x\to\zeta_0$ od wewnątrz, na zbiorze
całkowania mamy $|x-\zeta|\geq\delta/2$ dla
$x$ dostatecznie bliskich $\zeta_0$, podczas gdy
$R^2-|x|^2\to0$. Druga całka dąży więc do zera.
Otrzymujemy $u(x)\to g(\zeta_0)$, co daje ciągłe
przedłużenie na całą domkniętą kulę.
Jednoznaczność wynika z zasady maksimum
(twierdzenie~\ref{thm:laplace-max}).
\end{proof}

\begin{przyklad}[Równanie Dirichleta na dysku]
Niech $n=2$, $R=1$, a $g(\cos\theta,\sin\theta)=\cos(3\theta)$.
Funkcja $u(r\cos\theta,r\sin\theta)=r^3\cos(3\theta)$
ma $\Delta u=0$ i przyjmuje dane $g$ na brzegu.
Z jednoznaczności we wcześniejszym twierdzeniu jest
ona dokładnie całką Poissona z tych danych.
\end{przyklad}

\begin{stwierdzenie}[Nierówność Harnacka w kuli]
\label{prop:pde3-harnack}
Niech $n\geq2$, $R>0$, $u\in C^2(B_R)\cap C(\overline{B_R})$
będzie harmoniczna i nieujemna. Dla każdego $0\leq r<R$
i $|x|\leq r$ zachodzi
\begin{equation}\label{eq:pde3-harnack}
 \frac{R^{n-2}(R-r)}{(R+r)^{n-1}}u(0)
 \leq u(x)\leq
 \frac{R^{n-2}(R+r)}{(R-r)^{n-1}}u(0).
\end{equation}
\end{stwierdzenie}
\begin{proof}
Na podstawie~\eqref{eq:pde3-poisson} i $u|_{\partial B_R}\geq0$
możemy porównać jądro w punkcie $x$ i w punkcie $0$.
Ponieważ $P_R(0,\zeta)=1/(\omega_n R^{n-1})$,
\[
 \frac{P_R(x,\zeta)}{P_R(0,\zeta)}
 =\frac{R^{n-2}(R^2-|x|^2)}{|x-\zeta|^n}.
\]
Używając $R-|x|\leq|x-\zeta|\leq R+|x|$,
otrzymujemy
\[
 \frac{R^{n-2}(R-|x|)}{(R+|x|)^{n-1}}
 \leq\frac{P_R(x,\zeta)}{P_R(0,\zeta)}
 \leq\frac{R^{n-2}(R+|x|)}{(R-|x|)^{n-1}}.
\]
Obie funkcje graniczne są odpowiednio malejąca i
rosnąca względem $|x|$. Podstawiamy $|x|\leq r$,
mnożymy nierówności przez nieujemne $P_R(0,\zeta)u(\zeta)$
i całkujemy po $\partial B_R$. Otrzymujemy
\eqref{eq:pde3-harnack}.
\end{proof}

\begin{uwaga}[Nierówność Harnacka dla kul wewnętrznych]
Dla dodatniej funkcji harmonicznej na dowolnym obszarze
nierówność Harnacka obowiązuje na każdej kuli, której
domknięcie zawiera się w obszarze (po translacji).
Łącząc skończenie wiele takich kul, otrzymujemy
oszacowania na dowolnym zwartym podzbiorze spójnego obszaru. Nie zakłada się przy tym istnienia
ciągłego przedłużenia aż do \emph{brzegu całego obszaru}.
\end{uwaga}

\section{Problem Poissona i bariery porównawcze}
\label{sec:pde3-poisson-equation}

\begin{twierdzenie}[Problem Poissona na kuli]
\label{thm:pde3-ball-poisson}
Niech $n\geq2$, $R>0$, $f\in C_c^\infty(B_R)$,
a $g\in C(\partial B_R)$. Istnieje dokładnie jedno
$u\in C^2(B_R)\cap C(\overline{B_R})$ takie, że
$-\Delta u=f$ w $B_R$ i $u=g$ na $\partial B_R$.
Ma ono postać
\begin{equation}\label{eq:pde3-poisson-solution}
 u(x)=\int_{B_R}G_R(x,y)f(y)\,\mathrm d\lambda_n(y)
 +\int_{\partial B_R}P_R(x,\zeta)g(\zeta)
 \,\mathrm d\sigma(\zeta).
\end{equation}
\end{twierdzenie}
\begin{proof}
Przedłużmy $f$ przez zero poza $B_R$ i połóżmy
$w=\Phi_n*f$. Na mocy twierdzenia~\ref{thm:pde3-newton}
$w$ jest gładka i spełnia $-\Delta w=f$ w $\R^n$.
Wzór Poissona zastosowany do danych $g-w|_{\partial B_R}$
wyznacza harmoniczną funkcję $h$, ciągłą na domkniętej
kuli. Funkcja $u=w+h$ spełnia równanie i zadany
warunek brzegowy; jednoznaczność wynika z zasady maksimum.

Niech $H_R(x,y)=\Phi_n(x-y)-G_R(x,y)$.
Dla $y$ w zwartym nośniku $f$ funkcja $H_R(\cdot,y)$
jest harmoniczna w kuli, gładka aż do jej brzegu.
Zatem
$h_f(x)=\int_{B_R}H_R(x,y)f(y)\,\mathrm d\lambda_n(y)$
jest harmoniczna i ciągła na domkniętej kuli.
Wobec tego $v=w-h_f=\int_{B_R}G_R(x,y)f(y)\,
\mathrm d\lambda_n(y)$ spełnia $-\Delta v=f$.
Ponieważ $G_R(x,y)=0$ dla $x\in\partial B_R$,
a $y$ przebiega tylko zwarty nośnik $f\Subset B_R$,
mamy $v|_{\partial B_R}=0$. Stąd $v$ oraz całka
Poissona z danych $g$ dają rozwiązanie o zadanym
brzegu. Jego jednoznaczność dowodzi
\eqref{eq:pde3-poisson-solution}.
\end{proof}

\begin{stwierdzenie}[Oszacowanie przez barierę kwadratową]
\label{prop:pde3-barrier}
Niech $\Omega\subset B_R(0)$ będzie obszarem ograniczonym,
$n\geq2$, $u\in C^2(\Omega)\cap C(\overline\Omega)$
rozwiązaniem $-\Delta u=f$ w $\Omega$,
gdzie $f\in C(\Omega)$ jest ograniczona. Wówczas
\begin{equation}\label{eq:pde3-barrier}
 \sup_{x\in\Omega}|u(x)|\leq
 \sup_{x\in\partial\Omega}|u(x)|+
 \frac{R^2}{2n}\sup_{x\in\Omega}|f(x)|.
\end{equation}
\end{stwierdzenie}
\begin{proof}
Niech $M=\sup_\Omega|f|$ i
$b(x)=(R^2-|x|^2)/(2n)$. Mamy $b\geq0$ na $\Omega$
oraz $-\Delta b=1$. Stąd
$\Delta(u-Mb)=M-f\geq0$. Słaba zasada maksimum
daje $u-Mb\leq\sup_{\partial\Omega}u$, czyli
$u\leq\sup_{\partial\Omega}|u|+Mb$.
Ten sam argument zastosowany do $-u$ prowadzi do
$-u\leq\sup_{\partial\Omega}|u|+Mb$.
Ponieważ $0\leq b\leq R^2/(2n)$, otrzymujemy tezę.
\end{proof}

\begin{przyklad}[Wymuszenie stałe]
Niech $f\equiv1$ w $B_R(0)$, a dane Dirichleta będą równe $0$.
Wtedy $u(x)=(R^2-|x|^2)/(2n)$ jest dokładnym
rozwiązaniem $-\Delta u=1$. Oszacowanie
\eqref{eq:pde3-barrier} jest w tym przykładzie ostre.
\end{przyklad}

\section{Harmoniczność, analityczność i funkcje holomorficzne}
\label{sec:pde3-analytic}

\begin{twierdzenie}[Analityczność funkcji harmonicznych]
\label{thm:pde3-analytic}
Niech $n\geq2$, $\Omega\subset\R^n$ będzie otwarta,
a $u\in C^2(\Omega)$ harmoniczna. Wtedy $u$ jest
rzeczywiście analityczna w $\Omega$, czyli w otoczeniu
każdego punktu jest sumą zbieżnego szeregu potęgowego.
W szczególności $u\in C^\infty(\Omega)$.
\end{twierdzenie}
\begin{proof}
Ustalmy $x_0\in\Omega$ i kulę
$\overline{B_R(x_0)}\subset\Omega$. Z twierdzenia
\ref{thm:pde3-poisson} mamy w $B_R(x_0)$
\[
 u(x)=\int_{\partial B_R(x_0)}
 \frac{R^2-|x-x_0|^2}{\omega_n R|x-\zeta|^n}
 u(\zeta)\,\mathrm d\sigma(\zeta).
\]
Dla $x$ w dostatecznie małej kuli wokół $x_0$
wyrażenie $|x-\zeta|^{-n}$ jest rzeczywiście
analityczne w $x$: to funkcja analityczna dodatniego
wielomianu $|x-\zeta|^2$, którego wartość jest
jednostajnie oddzielona od zera, złożona z funkcją
$s\mapsto s^{-n/2}$. Szeregi potęgowe otrzymane
przez rozwinięcie dwumianowe mają wspólny dodatni
promień zbieżności i zbiegają jednostajnie względem
$\zeta\in\partial B_R(x_0)$ na mniejszej kuli.
Można więc całkować je wyraz po wyrazie. Otrzymany
szereg potęgowy przedstawia $u$ w otoczeniu $x_0$.
\end{proof}

\begin{stwierdzenie}[Harmoniczna sprzężona w wymiarze dwa]
\label{prop:pde3-conjugate}
Niech $\Omega\subset\R^2$ będzie obszarem jednospójnym, a $u\in C^2(\Omega)$
harmoniczna. Istnieje $v\in C^2(\Omega)$, jednoznaczna
z dokładnością do stałej addytywnej, taka, że
$F(x+iy)=u(x,y)+iv(x,y)$ jest holomorficzna w $\Omega$
traktowanej jako podzbiór $\C$.
\end{stwierdzenie}
\begin{proof}
Rozważmy formę $\alpha=-u_y\,\mathrm dx+u_x\,\mathrm dy$.
Ponieważ $u\in C^2$, zachodzi
$\mathrm d\alpha=(u_{xx}+u_{yy})\,\mathrm dx\wedge\mathrm dy=0$.
Na obszarze jednospójnym każda zamknięta forma
$1$ stopnia jest dokładna, zatem istnieje $v$ o
$\mathrm dv=\alpha$. Oznacza to $v_x=-u_y$ i $v_y=u_x$,
czyli równania Cauchy'ego--Riemanna. Ich spełnienie przy
regularności $C^1$ zapewnia holomorficzność $F$.
Jeśli $v_1,v_2$ mają te same pochodne, ich różnica
jest stała na obszarze.
\end{proof}

\begin{kontrprzyklad}[Przeszkoda topologiczna dla sprzężenia]
Na pierścieniu $\Omega=\{z\in\C:1<|z|<2\}$
funkcja $u(z)=\log|z|$ jest harmoniczna. Formalne
równania Cauchy'ego--Riemanna wymuszałyby
$\mathrm dv=-y\,(x^2+y^2)^{-1}\mathrm dx+
x\,(x^2+y^2)^{-1}\mathrm dy$. Całka tej formy po okręgu
$|z|=3/2$ wynosi $2\pi$, zatem forma nie jest dokładna.
Nie istnieje globalna jednowartościowa harmoniczna
sprzężona do $u$ na całym pierścieniu, choć istnieją
lokalne sprzężenia.
\end{kontrprzyklad}

\section{Wygładzanie harmoniczne na dysku: rozwiązanie przez szeregi Fouriera}
\label{sec:pde3-fourier-disk}

Jawna postać jądra Poissona ma szczególnie przejrzystą
interpretację na jednostkowym dysku. Niech $g\in C^2(S^1)$,
a jego rzeczywisty szereg Fouriera ma postać
$g(e^{i\theta})=a_0+\sum_{k=1}^\infty
(a_k\cos k\theta+b_k\sin k\theta)$.
Ponieważ $g\in C^2$, dwukrotne całkowanie przez części
we wzorze na współczynniki Fouriera daje
$|a_k|+|b_k|\leq C_g/k^2$. Szereg jest więc
bezwzględnie i jednostajnie zbieżny.

\begin{stwierdzenie}[Rozwinięcie harmoniczne na dysku]
\label{prop:pde3-fourier-disk}
Dla opisanych danych $g$ jedynym rozwiązaniem
$\Delta u=0$ w $B_1(0)\subset\R^2$,
$u|_{S^1}=g$, klasy $C^2(B_1)\cap C(\overline{B_1})$,
jest funkcja
\begin{equation}\label{eq:pde3-fourier-disk}
 u(r\cos\theta,r\sin\theta)=a_0+
 \sum_{k=1}^\infty r^k
 (a_k\cos k\theta+b_k\sin k\theta),\quad 0\leq r<1.
\end{equation}
\end{stwierdzenie}
\begin{proof}
W biegunowych współrzędnych na przebitym dysku
$\Delta=\partial_{rr}+r^{-1}\partial_r+r^{-2}\partial_{\theta\theta}$.
Dla każdego $k\geq1$ bezpośrednio sprawdzamy, że
$r^k\cos k\theta$ i $r^k\sin k\theta$ są harmoniczne.
Na każdym zwartym podzbiorze dysku szereg w
\eqref{eq:pde3-fourier-disk} można różniczkować
dowolnie wiele razy, dzięki czynnikowi $r^k$.
W zerze regularność wynika z tego, że
$r^k\cos k\theta=\operatorname{Re}(x+iy)^k$
(i analogicznie dla części urojonej).
Zbieżność jednostajna na domkniętym dysku wynika
z oszacowania $\sum_k(|a_k|+|b_k|)<\infty$.
Dlatego funkcja przedłuża się ciągle do brzegu,
gdzie przyjmuje wartości $g$. Jednoznaczność
wynika z zasady maksimum.
\end{proof}

Współczynnik wysokiej częstości $k$ zostaje wewnątrz
dysku pomnożony przez $r^k$. Dla każdego ustalonego
$r<1$ czynnik ten maleje wykładniczo wraz z $k$.
Mamy więc bardzo mocne \emph{wygładzanie w głąb obszaru}.
Odwrotna operacja, polegająca na odtwarzaniu danych
brzegowych z profilu na okręgu $|z|=r$, musi
dzielić współczynniki przez $r^k$ i jest niestabilna.
To geometryczny krewny kontrprzykładu Hadamarda
z rozdziału~\ref{chap:pde-intro}.

\begin{center}
\begin{tikzpicture}[>=Stealth,scale=1]
  \draw[->] (0,0)--(5.6,0) node[right] {$k$};
  \draw[->] (0,0)--(0,3.3) node[above] {$r^k$};
  \draw[densely dotted] (0,2.9)--(5.25,2.9);
  \node[left] at (0,2.9) {$1$};
  \draw[blue!70!black,thick,domain=0:5.1,samples=100]
    plot (\x,{2.9*exp(-.25*\x)});
  \draw[red!70!black,thick,domain=0:5.1,samples=100]
    plot (\x,{2.9*exp(-.9*\x)});
  \node[blue!70!black,anchor=west] at (3.5,1.5) {$r=e^{-1/4}$};
  \node[red!70!black,anchor=west] at (3.05,.70) {$r=e^{-9/10}$};
  \node[below] at (2.8,0) {częstość $k$};
\end{tikzpicture}
\par\smallskip\small
Tłumienie modów Fouriera w rozwiązaniu harmonicznym
na dysku: im dalej od brzegu, tym szybsze zanikanie
wysokich częstości.
\end{center}

\section{Zagadnienia brzegowe na bardziej ogólnych obszarach}
\label{sec:pde3-general-domain}

Jawny wzór Poissona jest szczególnym przywilejem kuli.
Dla ogólnego obszaru funkcja Greena może nie mieć
prostej postaci, a istnienie ciągłego rozwiązania Dirichleta
zależy od geometrii brzegu. Aby rozdzielić twierdzenia
udowodnione od materiału wymagającego późniejszych narzędzi,
formułujemy poniższy wynik jako zapowiedź bez dowodu.

\begin{uwaga}[Zapowiedź twierdzenia o ogólnych obszarach]
\label{prop:pde3-future}
Niech $n\geq2$ i niech $\Omega\subset\R^n$ będzie obszarem ograniczonym o brzegu klasy $C^2$.
Wtedy dla każdego $g\in C(\partial\Omega)$ istnieje
dokładnie jedna funkcja
$u\in C^2(\Omega)\cap C(\overline\Omega)$ spełniająca
$\Delta u=0$ w $\Omega$ i $u|_{\partial\Omega}=g$.
Jest to klasyczny wynik teorii potencjału. Nie traktujemy go
tu jako wyniku udowodnionego: dowód wymaga metody Perrona
i barier brzegowych lub odpowiednich twierdzeń o regularności
eliptycznej, które wprowadzimy w dalszej części skryptu.
\end{uwaga}

Warto odróżnić dwie kwestie. \emph{Jednoznaczność}
na każdym obszarze ograniczonym wynika z zasady maksimum,
niezależnie od regularności jej brzegu. Natomiast
\emph{istnienie} dla wszystkich ciągłych danych brzegowych
nie jest prawdziwe na każdym obszarze: punkt brzegowy
może być nieregularny w sensie teorii potencjału.
To zapowiedź kryterium Wienera i pojęcia pojemności,
które pojawi się w dalszej części skryptu.

\section{Podsumowanie i mapa dalszych związków}
\label{sec:pde3-summary}

Poznaliśmy dwa wzajemnie uzupełniające się spojrzenia
na równanie Laplace'a. Własność wartości średniej,
zasady maksimum, nierówność Harnacka i analityczność
ujawniają \emph{sztywność} funkcji harmonicznych.
Z kolei rozwiązanie fundamentalne, potencjał Newtona,
funkcja Greena i jądro Poissona dają \emph{konstrukcje}
rozwiązań, a nie tylko ich jakościowe własności.

Zasadnicze pytanie pozostaje otwarte w pełnej ogólności:
jak rozwiązywać $-\Delta u=f$ na obszarze, dla której
nie znamy jawnego jądra? Odpowiedź dadzą później
przestrzenie Sobolewa, metoda wariacyjna i teoria
regularności eliptycznej. W następnym rozdziale
porównamy te własności z równaniem ciepła i jego
charakterystycznym mechanizmem wygładzania.

% Rozdział 4: równanie ciepła.
\chapter[Równanie ciepła]{Równanie ciepła: jądro Gaussa, maksimum, wygładzanie i ewolucja}
\label{chap:pde-heat}

W poprzednim rozdziale warunek $\Delta u=0$ wymuszał silną
sztywność funkcji. Teraz badamy równanie, w którym laplasjan
steruje zmianą w czasie:
\begin{equation}\label{eq:pde4-heat}
 \partial_tu(t,x)-\kappa\Delta_xu(t,x)=f(t,x),\qquad
 t>0,\quad x\in\R^n,\quad \kappa>0.
\end{equation}
Gdy $f=0$, równanie opisuje swobodną dyfuzję. Fundamentalne pytania
są odmienne niż dla równania Laplace'a: jak odzyskać rozwiązanie
z profilu początkowego, jakie informacje zostają zachowane i dlaczego
każdy dodatni czas poprawia regularność? Wyprowadzimy rozwiązanie
całego zagadnienia Cauchy'ego, zasady maksimum, formułę Duhamela,
bilanse energii i przykłady pokazujące, że bez określenia klasy
rozwiązań nawet jednoznaczność może zawieść.

Przez $\mathrm d\lambda_n(x)$ oznaczamy miarę Lebesgue'a
w $\R^n$. Symbole $L^p(\R^n)$ i $\|\cdot\|_p$ oznaczają
przestrzenie funkcji całkowalnych z potęgą $p$ i ich normy,
a $\mathrm{BUC}(\R^n)$ --- przestrzeń funkcji ograniczonych
i jednostajnie ciągłych z normą supremum. Dla $p=\infty$
normą jest kres istotny. W dalszych twierdzeniach
\emph{silna ciągłość} oznacza zbieżność w normie danej
przestrzeni, a nie tylko punktową ani słabą. Będziemy korzystać z elementarnych
twierdzeń o całkowaniu: Fubiniego--Tonellego, zbieżności
zdominowanej, nierówności Minkowskiego i ciągłości przesunięć
w $L^p$; w rozdziałach o analizie funkcjonalnej wrócimy
do ich szerszego kontekstu.

\section{Jądro Gaussa: skąd bierze się rozwiązanie?}
\label{sec:pde4-gauss}

\begin{definicja}[Jądro cieplne]
\label{def:pde4-kernel}
Niech $n\geq1$ i $\kappa>0$. Dla $t>0$ oraz $x\in\R^n$
definiujemy
\begin{equation}\label{eq:pde4-kernel}
 \Gamma_\kappa(t,x)=\frac{1}{(4\pi\kappa t)^{n/2}}
 \exp\!\left(-\frac{|x|^2}{4\kappa t}\right).
\end{equation}
Zapis $\Gamma_\kappa(t)*g$ oznacza splot w zmiennej przestrzennej:
$(\Gamma_\kappa(t)*g)(x)=\int_{\R^n}\Gamma_\kappa(t,x-y)g(y)
\,\mathrm d\lambda_n(y)$, jeżeli całka ma sens.
\end{definicja}

\begin{twierdzenie}[Równanie, normalizacja i skalowanie]
\label{thm:pde4-gauss-basic}
Dla każdego $n\geq1$, $\kappa>0$, $t>0$ zachodzi:
$\Gamma_\kappa(t,x)>0$,
$(\partial_t-\kappa\Delta_x)\Gamma_\kappa(t,x)=0$
dla $x\in\R^n$, oraz
\begin{equation}\label{eq:pde4-normalize}
 \int_{\R^n}\Gamma_\kappa(t,x)\,\mathrm d\lambda_n(x)=1.
\end{equation}
Ponadto $\Gamma_\kappa(t,x)=t^{-n/2}\Gamma_\kappa(1,x/\sqrt t)$.
\end{twierdzenie}
\begin{proof}
Rachunek bezpośredni daje
$\partial_t\Gamma_\kappa=\Gamma_\kappa
(-n/(2t)+|x|^2/(4\kappa t^2))$.
Z drugiej strony
$\partial_{x_j}^2\Gamma_\kappa=\Gamma_\kappa
(x_j^2/(4\kappa^2t^2)-1/(2\kappa t))$.
Po zsumowaniu po $j$ otrzymujemy równanie.
W całce stosujemy podstawienie $z=x/\sqrt{4\kappa t}$
oraz elementarną tożsamość
$\int_{\R^n}e^{-|z|^2}\,\mathrm d\lambda_n(z)=\pi^{n/2}$;
tę ostatnią otrzymuje się przez pomnożenie całek
jednowymiarowych, z których kwadrat pierwszej
oblicza się we współrzędnych biegunowych w $\R^2$.
Skalowanie jest natychmiastowe ze wzoru.
\end{proof}

\begin{uwaga}[Skalowanie paraboliczne]
Jeśli $u$ rozwiązuje jednorodne równanie ciepła,
to $u_\rho(t,x)=u(\rho^2t,\rho x)$ również je rozwiązuje
dla każdej $\rho>0$. Czas skaluje się więc jak kwadrat
odległości, w odróżnieniu od skali falowej $t\sim|x|$.
Przy wyprowadzeniu jądra można poszukiwać funkcji
podobnej do siebie w różnych skalach,
$\Gamma(t,x)=t^{-n/2}F(x/\sqrt t)$;
profil Gaussa wynika z uzyskanego równania dla $F$.
\end{uwaga}

\begin{center}
\begin{tikzpicture}[>=Stealth,x=1.15cm,y=3.8cm]
 \draw[->] (-3.2,0)--(3.55,0) node[right] {$x$};
 \draw[->] (0,0)--(0,1.04) node[above] {$\Gamma_1(t,x)$};
 \draw[blue!70!black,thick,domain=-3:3,samples=110]
   plot (\x,{1/sqrt(4*pi*.12)*exp(-\x*\x/(4*.12))});
 \draw[red!70!black,thick,domain=-3:3,samples=110]
   plot (\x,{1/sqrt(4*pi*.45)*exp(-\x*\x/(4*.45))});
 \draw[green!45!black,thick,domain=-3:3,samples=110]
   plot (\x,{1/sqrt(4*pi*1.1)*exp(-\x*\x/(4*1.1))});
 \node[blue!70!black,anchor=west] at (1.5,.79) {$t=0{,}12$};
 \node[red!70!black,anchor=west] at (1.5,.68) {$t=0{,}45$};
 \node[green!45!black,anchor=west] at (1.5,.57) {$t=1{,}1$};
\end{tikzpicture}
\par\smallskip\small
Jądro cieplne w wymiarze jeden ($\kappa=1$): profil
staje się niższy i szerszy, lecz całkowita masa pozostaje równa $1$.
\end{center}

\section{Własność półgrupy i przybliżona jedność}
\label{sec:pde4-semigroup}

\begin{lemat}[Składanie jąder Gaussa]
\label{lem:pde4-convolution}
Dla $n\geq1$, $\kappa>0$ oraz $s,t>0$ zachodzi
$\Gamma_\kappa(t)*\Gamma_\kappa(s)=\Gamma_\kappa(t+s)$.
\end{lemat}
\begin{proof}
W wykładniku pod całką używamy tożsamości
\[
 \frac{|x-y|^2}{t}+\frac{|y|^2}{s}
 =\frac{t+s}{ts}\left|y-\frac{s}{t+s}x\right|^2
 +\frac{|x|^2}{t+s}.
\]
Po wyłączeniu czynnika zależnego od $x$ i obliczeniu
pozostałej całki Gaussowskiej otrzymujemy
$\Gamma_\kappa(t+s,x)$; normalizacja stałej wynika
z~\eqref{eq:pde4-normalize} albo z jawnego rachunku.
Wszystkie użyte całki są skończone i nieujemne,
więc przestawienie kolejności nie stanowi problemu.
\end{proof}

\begin{twierdzenie}[Przybliżona jedność]
\label{thm:pde4-approx-identity}
Niech $n\geq1$, $\kappa>0$. Dla $g\in\mathrm{BUC}(\R^n)$
funkcje $P_tg:=\Gamma_\kappa(t)*g$ spełniają
$\|P_tg-g\|_\infty\to0$ przy $t\downarrow0$.
Dla $1\leq p<\infty$ i $g\in L^p(\R^n)$ zachodzi
$\|P_tg-g\|_p\to0$ przy $t\downarrow0$.
Dla każdej $g\in L^p$ ($1\leq p\leq\infty$),
$\|P_tg\|_p\leq\|g\|_p$ dla $t>0$.
\end{twierdzenie}
\begin{proof}
Dla funkcji jednostajnie ciągłej i ograniczonej, dzięki
normalizacji jądra,
\[
 |P_tg(x)-g(x)|\leq
 \int_{\R^n}\Gamma_\kappa(t,y)|g(x-y)-g(x)|
 \,\mathrm d\lambda_n(y).
\]
Ustalmy $\eps>0$ i dobierzmy $\delta>0$ tak, aby
$|g(x-y)-g(x)|<\eps$ dla $|y|<\delta$, jednostajnie w $x$.
Całka po tej części jest nie większa niż $\eps$,
a po pozostałej nie przekracza
$2\|g\|_\infty\int_{|y|\geq\delta}\Gamma_\kappa(t,y)
\,\mathrm d\lambda_n(y)$. Po podstawieniu
$y=\sqrt t\,z$ ostatnia całka dąży do $0$.

Dla $1\leq p<\infty$ wykorzystujemy tę samą tożsamość,
teraz z normą $L^p$. Nierówność Minkowskiego daje
\[
 \|P_tg-g\|_p\leq
 \int_{\R^n}\Gamma_\kappa(t,y)
 \|g(\cdot-y)-g\|_p\,\mathrm d\lambda_n(y).
\]
Przesunięcia są ciągłe w $L^p$: najpierw sprawdza się
to bezpośrednio dla $C_c^\infty(\R^n)$, a następnie
przybliża dowolną funkcję $L^p$ przez funkcje tej klasy.
Oszacowanie $\|g(\cdot-y)-g\|_p\leq2\|g\|_p$
i skupianie się masy $\Gamma_\kappa(t,\cdot)$ przy zerze
kończą dowód. Nierówność kontrakcji otrzymujemy z
nierówności Minkowskiego i normowania jądra
($p=\infty$ bezpośrednio przez oszacowanie supremum).
\end{proof}

Z lematu~\ref{lem:pde4-convolution} i twierdzenia
\ref{thm:pde4-approx-identity} wynika
$P_{t+s}=P_tP_s$, a dla $1\leq p<\infty$
odwzorowanie $t\mapsto P_tg$ jest ciągłe w $L^p$
na $[0,\infty)$. Na $L^\infty$ taka zbieżność w normie
nie musi zachodzić; odpowiednią przestrzenią dla
ciągłości normowej w chwili $0$ jest $\mathrm{BUC}(\R^n)$.

\section{Zagadnienie Cauchy'ego na całej przestrzeni}
\label{sec:pde4-cauchy}

\begin{twierdzenie}[Jawne rozwiązanie z danych początkowych]
\label{thm:pde4-cauchy}
Niech $n\geq1$, $\kappa>0$ i $g\in\mathrm{BUC}(\R^n)$.
Definiujemy dla $t>0$,
\begin{equation}\label{eq:pde4-cauchy}
 u(t,x)=\int_{\R^n}\Gamma_\kappa(t,x-y)g(y)
 \,\mathrm d\lambda_n(y),\qquad u(0,x)=g(x).
\end{equation}
Wówczas $u\in C([0,\infty);\mathrm{BUC}(\R^n))$,
$u\in C^\infty((0,\infty)\times\R^n)$,
$u_t-\kappa\Delta_xu=0$ dla $t>0$ oraz
$\sup_{t\geq0}\|u(t,\cdot)\|_\infty\leq\|g\|_\infty$.
Ponadto $g\geq0$ implikuje $u\geq0$.
Jest to jedyne rozwiązanie z klasy
$C^{1,2}((0,T]\times\R^n)\cap C([0,T]\times\R^n)$,
ograniczone na $[0,T]\times\R^n$, o tych danych,
dla każdego ustalonego $T>0$.
\end{twierdzenie}
\begin{proof}
Dla $t$ z przedziału zwartego zawartego w $(0,\infty)$
każda pochodna przestrzenna i czasowa jądra
$\Gamma_\kappa(t,x-y)$ jest wielomianem w $x-y$
pomnożonym przez funkcję Gaussa, z oszacowaniem
całkowalnym w $y$ niezależnym od $x$ na zwartych
zbiorach. Ponieważ $g$ jest ograniczona, możemy
różniczkować pod całką dowolną liczbę razy.
Dla równania ciepła przenosimy operator
$\partial_t-\kappa\Delta_x$ pod całkę i stosujemy
twierdzenie~\ref{thm:pde4-gauss-basic}.
Ciągłość w normie w chwili $0$ wynika z twierdzenia
\ref{thm:pde4-approx-identity}; w pozostałych chwilach
z własności półgrupy i kontrakcji.
Dodatniość i oszacowanie supremum są następstwem
$\Gamma_\kappa\geq0$ oraz jego normalizacji.
Jednoznaczność udowodnimy poniżej niezależnie
w twierdzeniu~\ref{thm:pde4-global-uniqueness},
korzystając z zasady maksimum, bez powoływania się
na wzór jawny.
\end{proof}

\begin{wniosek}[Natychmiastowe rozchodzenie się zaburzenia]
\label{cor:pde4-positive}
Niech $n\geq1$, $\kappa>0$, a
$g\in C_c(\R^n)$ będzie nieujemna i nierówna identycznie
zeru. Wówczas rozwiązanie~\eqref{eq:pde4-cauchy}
ma własność $u(t,x)>0$ dla każdego $t>0$, $x\in\R^n$.
\end{wniosek}
\begin{proof}
Istnieje kula $B\subset\R^n$ i $c>0$ takie, że $g\geq c$
na $B$. Ponieważ $\Gamma_\kappa(t,x-y)>0$ dla
każdego $y$, mamy
$u(t,x)\geq c\int_B\Gamma_\kappa(t,x-y)
\,\mathrm d\lambda_n(y)>0$.
\end{proof}

\begin{stwierdzenie}[Zachowanie masy i tłumienie norm]
\label{prop:pde4-mass}
Niech $n\geq1$, $\kappa>0$, $t>0$.
Dla $g\in L^1(\R^n)$ zachodzi
\begin{equation}\label{eq:pde4-mass}
 \int_{\R^n}(P_tg)(x)\,\mathrm d\lambda_n(x)
 =\int_{\R^n}g(x)\,\mathrm d\lambda_n(x).
\end{equation}
Dla $g\in L^1(\R^n)$ zachodzi również
\begin{equation}\label{eq:pde4-L1Linf}
 \|P_tg\|_\infty\leq (4\pi\kappa t)^{-n/2}\|g\|_1,
 \qquad
 \|P_tg\|_2\leq (8\pi\kappa t)^{-n/4}\|g\|_1.
\end{equation}
\end{stwierdzenie}
\begin{proof}
Ponieważ $\int_{\R^n}\int_{\R^n}\Gamma_\kappa(t,x-y)|g(y)|
\,\mathrm d\lambda_n(y)\,\mathrm d\lambda_n(x)=\|g\|_1$,
można użyć twierdzenia Fubiniego i
\eqref{eq:pde4-normalize}, co daje zachowanie masy.
Pierwsze oszacowanie jest bezpośrednie z
$\|\Gamma_\kappa(t,\cdot)\|_\infty=(4\pi\kappa t)^{-n/2}$.
Drugie wynika z nierówności Minkowskiego dla całki
z funkcji $g(y)\Gamma_\kappa(t,\cdot-y)$ oraz
$\|\Gamma_\kappa(t,\cdot)\|_2=(8\pi\kappa t)^{-n/4}$.
\end{proof}

\section{Paraboliczna zasada maksimum i jednoznaczność}
\label{sec:pde4-maximum}

Dla obszaru ograniczonego $\Omega\subset\R^n$ i $T>0$
wprowadzamy cylinder $Q_T=(0,T)\times\Omega$ oraz jego
\emph{brzeg paraboliczny}
$\partial_pQ_T=(\{0\}\times\overline\Omega)
\cup([0,T]\times\partial\Omega)$.
To właśnie tam musimy kontrolować rozwiązanie;
na górnej podstawie $\{T\}\times\Omega$ nie zadaje się
niezależnych danych w zagadnieniu postępowym.

\begin{twierdzenie}[Słaba paraboliczna zasada maksimum]
\label{thm:pde4-max}
Niech $\Omega\subset\R^n$ będzie obszarem ograniczonym, $T>0$, $\kappa>0$, a
$u\in C([0,T]\times\overline\Omega)\cap
C^{1,2}((0,T]\times\Omega)$ spełnia
$u_t-\kappa\Delta_xu\leq0$ w $\Omega\times(0,T]$.
Wówczas
\begin{equation}\label{eq:pde4-max}
 \max_{[0,T]\times\overline\Omega}u
 \leq \max_{\partial_pQ_T}u.
\end{equation}
W konsekwencji dwie funkcje tej klasy rozwiązujące
to samo jednorodne równanie, o tych samych danych na
$\partial_pQ_T$, są równe.
\end{twierdzenie}
\begin{proof}
Dla $\eps>0$ połóżmy $v(t,x)=u(t,x)-\eps t$.
Wówczas $v_t-\kappa\Delta v\leq-\eps<0$.
Jeśli maksimum $v$ występuje poza $\partial_pQ_T$,
to istnieje punkt $(t_0,x_0)$, gdzie $t_0>0$,
$x_0\in\Omega$, taki że $v$ osiąga tam maksimum
na całym cylindrze. W tym punkcie
$\nabla_xv=0$, $\Delta_xv\leq0$ i $v_t\geq0$;
przy $t_0=T$ ostatnia nierówność dotyczy pochodnej
lewostronnej. W rezultacie
$v_t-\kappa\Delta v\geq0$, sprzeczność.
Zatem $\max v\leq\max_{\partial_pQ_T}v
\leq\max_{\partial_pQ_T}u$. Ponieważ
$u=v+\eps t\leq v+\eps T$, przechodzimy
z $\eps\downarrow0$ i dostajemy tezę.
Równość dwóch rozwiązań wynika z zastosowania
zasady maksimum do ich różnicy i do różnicy
z przeciwnym znakiem.
\end{proof}

\begin{twierdzenie}[Jednoznaczność w klasie rozwiązań ograniczonych]
\label{thm:pde4-global-uniqueness}
Niech $n\geq1$, $\kappa>0$, $T>0$. Jeśli
$u,v\in C([0,T]\times\R^n)\cap
C^{1,2}((0,T]\times\R^n)$ są ograniczone
na $[0,T]\times\R^n$, rozwiązują
$w_t-\kappa\Delta_xw=0$ (osobno dla $w=u$ i $w=v$)
i spełniają $u(0,x)=v(0,x)$ dla wszystkich
$x\in\R^n$, to $u=v$ na $[0,T]\times\R^n$.
\end{twierdzenie}
\begin{proof}
Niech $w=u-v$ i $M=\sup_{[0,T]\times\R^n}|w|$.
Dla $\eps>0$ rozważmy
$b_\eps(t,x)=\eps(|x|^2+2n\kappa t)$.
Jest to dokładne rozwiązanie równania jednorodnego,
ponieważ $\partial_tb_\eps=2n\kappa\eps=
\kappa\Delta b_\eps$.
Wybierzmy $R>0$ tak duże, żeby $\eps R^2>M$.
Na części bocznej $[0,T]\times\partial B_R$
mamy $w-b_\eps\leq M-\eps R^2<0$,
a na podstawie $\{0\}\times B_R$
mamy $w-b_\eps=-\eps|x|^2\leq0$.
Z zasady maksimum zastosowanej do $w-b_\eps$
na $[0,T]\times B_R$ wynika
$w(t,x)\leq\eps(|x|^2+2n\kappa t)$ dla $|x|<R$.
Ustalając $(t,x)$ i przechodząc z $\eps\downarrow0$
(dla każdego $\eps$ wybieramy nowe $R>|x|$),
dostajemy $w(t,x)\leq0$. Ten sam argument dla
$-w$ dowodzi $w=0$.
\end{proof}

\begin{uwaga}[Dlaczego trzeba podać klasę rozwiązań?]
Zagadnienie Cauchy'ego na przestrzeni nieograniczonej
może mieć rozwiązania nieograniczone względem $x$,
które nie są jednoznaczne nawet przy danych zerowych.
Nie przeczy to powyższemu twierdzeniu: z jego założeń
wynika ograniczoność funkcji na całym cylindrze.
Jawny kontrprzykład podamy pod koniec rozdziału.
\end{uwaga}

\section{Wygładzanie i analityczność przestrzenna}
\label{sec:pde4-smoothing}

\begin{twierdzenie}[Oszacowania pochodnych]
\label{thm:pde4-smoothing}
Niech $n\geq1$, $\kappa>0$, $g\in L^\infty(\R^n)$,
a $u(t,x)=P_tg(x)$ dla $t>0$.
Dla każdego wielowskaźnika $\alpha\in\N_0^n$
istnieje stała $C_{\alpha,n}>0$, niezależna od
$g,t,\kappa$, taka że
\begin{equation}\label{eq:pde4-derivative}
 \|\partial_x^\alpha u(t,\cdot)\|_\infty
 \leq C_{\alpha,n}(\kappa t)^{-|\alpha|/2}\|g\|_\infty.
\end{equation}
Co więcej, dla każdego ustalonego $t>0$ funkcja
$x\mapsto u(t,x)$ przedłuża się do funkcji holomorficznej
na $\C^n$ (jako funkcja zespolonych zmiennych $z_1,\ldots,z_n$).
\end{twierdzenie}
\begin{proof}
Każda pochodna $\partial_x^\alpha\Gamma_\kappa(t,x)$
jest iloczynem skalującego czynnika
$(\kappa t)^{-(n+|\alpha|)/2}$ i wielomianu w
$x/\sqrt{\kappa t}$ razy funkcję Gaussa.
Stąd
$\|\partial_x^\alpha\Gamma_\kappa(t,\cdot)\|_1
=C_{\alpha,n}(\kappa t)^{-|\alpha|/2}$.
Różniczkując splot i stosując oszacowanie
$L^1$--$L^\infty$, otrzymujemy~\eqref{eq:pde4-derivative}.

Dla $z=x+i\eta\in\C^n$ określamy
\[
 U_t(z)=\frac1{(4\pi\kappa t)^{n/2}}
 \int_{\R^n}
 \exp\!\left(-\frac{\sum_{j=1}^n(z_j-y_j)^2}{4\kappa t}\right)
 g(y)\,\mathrm d\lambda_n(y).
\]
Moduł całkowanego wyrażenia jest równy
$\exp(|\eta|^2/(4\kappa t))
\exp(-|x-y|^2/(4\kappa t))|g(y)|$
razy dodatni czynnik normalizujący.
Na każdym zwartym zbiorze zmiennej zespolonej $z$
funkcję tę ogranicza jedna całkowalna w $y$ funkcja
Gaussa (po ewentualnym zwiększeniu stałej i osłabieniu
wykładnika). Całka definiuje zatem funkcję
holomorficzną przez różniczkowanie pod całką, a jej
restrykcja do $\R^n$ jest równa $u(t,\cdot)$.
\end{proof}

\begin{przyklad}[Wygasanie wysokich częstości]
Niech $\xi\in\R^n$ i
$g(x)=\cos(\xi\cdot x)$. Wówczas
$u(t,x)=e^{-\kappa|\xi|^2t}\cos(\xi\cdot x)$
rozwiązuje równanie ciepła i spełnia dane początkowe.
Z twierdzenia o jednoznaczności w klasie ograniczonej
musi być równe $P_tg$. Wysokie częstości zanikają
z szybkością wykładniczą względem $|\xi|^2$.
\end{przyklad}

\begin{lemat}[Generator półgrupy cieplnej na gładkich danych]
\label{lem:pde-v31-heat-generator}
Niech $n\ge1$, $\kappa>0$ oraz
$h\in C^2(\R^n)$, przy czym $h$ i wszystkie pochodne
przestrzenne rzędu co najwyżej dwa są ograniczone
i jednostajnie ciągłe na $\R^n$.
Wtedy dla $t>0$ zachodzi
\begin{equation}\label{eq:pde-v31-heat-generator}
 P_th-h=\kappa\int_0^t P_s(\Delta h)\,\mathrm ds,
 \qquad
 \lim_{t\downarrow0}\left\|\frac{P_th-h}{t}
                 -\kappa\Delta h\right\|_\infty=0.
\end{equation}
\end{lemat}
\begin{proof}
Dla $t>0$ różniczkowanie jądra Gaussa i całkowanie
przez części względem zmiennej przestrzennej dają
$\partial_tP_th=\kappa P_t(\Delta h)$. Wyrazy
brzegowe znikają: można najpierw całkować po kuli,
a następnie puścić jej promień do nieskończoności,
wykorzystując ograniczoność $h,\nabla h$ i gaussowski
zanik jądra oraz jego pochodnych. Równość pochodnych
można również uzyskać z tożsamości
$\Delta_x\Gamma(t,x-y)=\Delta_y\Gamma(t,x-y)$.
Dla $0<\delta<t$ całkujemy równanie w czasie od
$\delta$ do $t$:
$P_th-P_\delta h=\kappa\int_\delta^tP_s(\Delta h)\,\mathrm ds$.
Z twierdzenia~\ref{thm:pde4-approx-identity},
zastosowanego do $h$ i $\Delta h$, wolno przejść
z $\delta\downarrow0$ w normie supremum.
Otrzymujemy pierwszą równość.
Dzieląc ją przez $t$ i odejmując $\kappa\Delta h$,
dostajemy
\[
 \left\|\frac{P_th-h}{t}-\kappa\Delta h\right\|_\infty
 \le\frac\kappa t\int_0^t\|P_s\Delta h-\Delta h\|_\infty
 \,\mathrm ds\longrightarrow0.
\]
Zbieżność ostatniej całki wynika z silnej ciągłości
półgrupy na $\mathrm{BUC}(\R^n)$.
\end{proof}

\section{Wymuszenie: zasada Duhamela}
\label{sec:pde4-duhamel}

\begin{twierdzenie}[Wzór Duhamela]
\label{thm:pde4-duhamel}
Niech $T>0$, $n\geq1$, $\kappa>0$, $g\in C_c^\infty(\R^n)$,
a $f\in C^\infty([0,T]\times\R^n)$ ma własność:
wszystkie jego pochodne względem $x$ i $t$ są ograniczone,
a nośniki $f(t,\cdot)$ są zawarte we wspólnym zbiorze zwartym
$K\subset\R^n$ dla $0\leq t\leq T$.
Wtedy funkcja
\begin{equation}\label{eq:pde4-duhamel}
 u(t,x)=P_tg(x)+\int_0^t(P_{t-s}f(s,\cdot))(x)
 \,\mathrm ds,\quad 0\leq t\leq T,
\end{equation}
przy konwencji $P_0h=h$, jest klasy $C^{1,2}$
na $[0,T]\times\R^n$, spełnia
$u_t-\kappa\Delta_xu=f$ dla $0<t\leq T$ i $u(0)=g$.
Ponadto
\begin{equation}\label{eq:pde4-duhamel-bound}
 \|u(t,\cdot)\|_\infty\leq\|g\|_\infty+
 \int_0^t\|f(s,\cdot)\|_\infty\,\mathrm ds.
\end{equation}
\end{twierdzenie}
\begin{proof}
Z regularności i ograniczoności pochodnych danych
wynika, że wszystkie przestrzenne pochodne $P_rg$
i $P_rf(s)$ można przenosić na $g$ lub $f(s)$.
Ponadto $\partial_rP_rh=\kappa P_r\Delta h$
dla $r>0$, a pochodna prawostronna w $r=0$
istnieje w normie supremum na mocy
lematu~\ref{lem:pde-v31-heat-generator}, ponieważ
$h$ i jego pochodne spełniają jego założenia.
Dla $v(t)=\int_0^tP_{t-s}f(s)\,\mathrm ds$
możemy zatem różniczkować po $t$:
\[
 \partial_tv(t)=f(t)+\kappa\int_0^t
 P_{t-s}\Delta f(s)\,\mathrm ds
 =f(t)+\kappa\Delta v(t).
\]
Ten sam rachunek z $P_tg$ daje
$\partial_tP_tg=\kappa\Delta P_tg$.
Jednostajna zbieżność przy $t\downarrow0$ i założenia
na pochodne danych gwarantują regularność aż do chwili
początkowej. Oszacowanie~\eqref{eq:pde4-duhamel-bound}
wynika z kontrakcji $P_t$ i nierówności trójkąta.
\end{proof}

\begin{przyklad}[Źródło stałe w czasie]
Jeśli $f(t,x)=h(x)$ z $h\in C_c^\infty(\R^n)$ oraz
$g=0$, to $u(t,x)=\int_0^tP_sh(x)\,\mathrm ds$.
Różniczkując dostajemy $u_t=P_th$,
a z tożsamości $\partial_sP_sh=\kappa\Delta P_sh$
wynika $\kappa\Delta u=P_th-h$, zatem
$u_t-\kappa\Delta u=h$.
\end{przyklad}

\section{Dyfuzja w obszarze ograniczonym: energia i widmo}
\label{sec:pde4-bounded}

Na obszarze ograniczonym wzór splotowy nie uwzględnia
brzegu. Naturalne warunki Dirichleta określają
temperaturę na brzegu, a Neumanna --- strumień ciepła.
W tej sekcji wyprowadzimy podstawowe wnioski
\emph{przy założeniu istnienia dostatecznie gładkiego
rozwiązania}. Ogólną teorię istnienia oprzemy później
na przestrzeniach Sobolewa i półgrupach operatorów.

\begin{twierdzenie}[Bilans energii]
\label{thm:pde4-energy}
Niech $\Omega\subset\R^n$ będzie obszarem ograniczonym
o brzegu klasy $C^1$, $T>0$, $\kappa>0$.
Niech $u\in C^{1,2}([0,T]\times\overline\Omega)$ spełnia
$u_t-\kappa\Delta_xu=0$ oraz $u=0$
na $[0,T]\times\partial\Omega$.
Wówczas dla $0\leq t\leq T$
\begin{equation}\label{eq:pde4-energy}
 \frac12\int_\Omega |u(t,x)|^2\,\mathrm d\lambda_n(x)
 +\kappa\int_0^t\!\int_\Omega|\nabla_xu(s,x)|^2
 \,\mathrm d\lambda_n(x)\,\mathrm ds
 =\frac12\int_\Omega|u(0,x)|^2\,\mathrm d\lambda_n(x).
\end{equation}
W szczególności $t\mapsto\|u(t,\cdot)\|_2$ jest
nierosnąca. W powyższej klasie rozwiązanie
z zadanymi danymi początkowymi i zerowym warunkiem
Dirichleta jest jednoznaczne.
\end{twierdzenie}
\begin{proof}
Mnożymy równanie przez $u$, całkujemy po $\Omega$
i używamy pierwszej tożsamości Greena
\eqref{eq:pde3-green1}:
\[
\begin{aligned}
 \frac12\frac{\mathrm d}{\mathrm dt}
 \int_\Omega u(t,x)^2\,\mathrm d\lambda_n(x)
 &=\kappa\int_\Omega u(t,x)\Delta_x u(t,x)\,\mathrm d\lambda_n(x)\\
 &=-\kappa\int_\Omega|\nabla_x u(t,x)|^2\,\mathrm d\lambda_n(x)\\
 &\quad+\kappa\int_{\partial\Omega}u(t,\zeta)\partial_\nu u(t,\zeta)
 \,\mathrm d\sigma(\zeta).
\end{aligned}
\]
Całka brzegowa znika. Całkujemy względem czasu
od $0$ do $t$, otrzymując~\eqref{eq:pde4-energy}.
Dla dwóch rozwiązań bierzemy ich różnicę;
przy zerowych danych początkowych energia jest zerowa,
więc różnica znika w każdym punkcie przez ciągłość.
\end{proof}

\begin{lemat}[Elementarna nierówność Poincarégo w obszarze ograniczonym]
\label{lem:pde4-poincare}
Niech $L>0$ i $\Omega\subset(-L,L)^n$ będzie obszarem ograniczonym o brzegu klasy $C^1$. Jeśli
$v\in C^1(\overline\Omega)$ i $v|_{\partial\Omega}=0$,
to
\begin{equation}\label{eq:pde4-poincare}
 \int_\Omega |v(x)|^2\,\mathrm d\lambda_n(x)
 \leq4L^2\int_\Omega|\nabla v(x)|^2\,\mathrm d\lambda_n(x).
\end{equation}
\end{lemat}
\begin{proof}
Przedłużamy $v$ przez zero do sześcianu
$Q=(-L,L)^n$. Z powodu zerowych wartości na brzegu
funkcja przedłużona jest ciągła; na niemal każdej
prostej równoległej do osi $x_1$ jest absolutnie ciągła,
a jej pochodna jest równa $\partial_{x_1}v$
w $\Omega$ i $0$ poza $\Omega$. Ten fakt wynika
z podziału prostej na składowe otwarte przecięcia
z $\Omega$ i z zerowania się $v$ na końcach każdej
składowej. Dla niemal wszystkich $x=(x_1,x')\in Q$
mamy więc
$|v(x)|=|\int_{-L}^{x_1}\partial_{x_1}v(s,x')\,
\mathrm ds|\leq\sqrt{2L}
(\int_{-L}^L|\partial_{x_1}v(s,x')|^2\,\mathrm ds)^{1/2}$.
Całkując tę nierówność kwadratową po $Q$,
otrzymujemy $\int_Q|v|^2\leq(2L)^2
\int_Q|\partial_{x_1}v|^2$, co daje tezę.
\end{proof}

\begin{wniosek}[Wykładniczy zanik energii przy warunku Dirichleta]
\label{cor:pde4-exponential}
Przy założeniach twierdzenia~\ref{thm:pde4-energy}
załóżmy dodatkowo $\Omega\subset(-L,L)^n$.
Wówczas
\begin{equation}\label{eq:pde4-exponential}
 \|u(t,\cdot)\|_2\leq
 \exp\!\left(-\frac{\kappa t}{4L^2}\right)
 \|u(0,\cdot)\|_2\qquad(0\leq t\leq T).
\end{equation}
\end{wniosek}
\begin{proof}
Niech $E(t)=\frac12\|u(t,\cdot)\|_2^2$.
Z bilansu energii i~\eqref{eq:pde4-poincare} mamy
$E'(t)=-\kappa\|\nabla u(t,\cdot)\|_2^2
\leq-\kappa(4L^2)^{-1}\|u(t,\cdot)\|_2^2
=-\kappa E(t)/(2L^2)$.
Nierówność Grönwalla daje
$E(t)\leq E(0)e^{-\kappa t/(2L^2)}$;
po spierwiastkowaniu otrzymujemy tezę.
\end{proof}

\begin{przyklad}[Mody własne na odcinku]
Niech $\Omega=(0,L)$, $n=1$, $\kappa>0$.
Dla dowolnego $m\in\N$ funkcja
$u_m(t,x)=e^{-\kappa(m\pi/L)^2t}\sin(m\pi x/L)$
rozwiązuje równanie ciepła, znika na $x=0,L$
i przyjmuje w chwili $0$ dane
$g_m(x)=\sin(m\pi x/L)$.
Skończona kombinacja liniowa takich modów również
jest rozwiązaniem. Ogólne rozwinięcia w szereg własnych
funkcji laplasjanu omówimy po wprowadzeniu teorii Fouriera.
\end{przyklad}

\section{Zachowanie dla dużych czasów i nieodwracalność}
\label{sec:pde4-large-time}

\begin{twierdzenie}[Zanik na całej przestrzeni]
\label{thm:pde4-large-time}
Niech $n\geq1$ i $\kappa>0$.
Jeśli $g\in L^1(\R^n)$, to $\|P_tg\|_\infty\to0$
przy $t\to\infty$.
Jeśli $g\in L^2(\R^n)$, to $\|P_tg\|_2\to0$
przy $t\to\infty$.
\end{twierdzenie}
\begin{proof}
Pierwsze stwierdzenie wynika bezpośrednio z
\eqref{eq:pde4-L1Linf}.
Jeśli dodatkowo $g\in L^1\cap L^2$, drugie oszacowanie
z~\eqref{eq:pde4-L1Linf} daje $\|P_tg\|_2\to0$.
Dla dowolnego $g\in L^2$ wybieramy ciąg
gładkich funkcji o zwartym nośniku $g_j\to g$ w $L^2$.
Dzięki kontrakcji $P_t$ w $L^2$ mamy
$\|P_tg\|_2\leq\|g-g_j\|_2+\|P_tg_j\|_2$.
Najpierw ustalamy duże $j$, a następnie
przechodzimy z $t\to\infty$. Otrzymujemy tezę.
\end{proof}

\begin{kontrprzyklad}[Wsteczne równanie ciepła jest niestabilne]
\label{ex:pde4-backward}
Niech $\Omega=(0,\pi)$, $\kappa=1$ i $T>0$.
Dla $m\in\N$ funkcja
$u_m(t,x)=e^{-m^2t}\sin(mx)$ jest rozwiązaniem
z jednorodnymi warunkami Dirichleta.
W chwili $0$ zachodzi $\|u_m(0,\cdot)\|_\infty=1$,
a w chwili $T$ mamy $\|u_m(T,\cdot)\|_\infty=e^{-m^2T}\to0$.
Zatem odwrotne odwzorowanie, które miałoby odzyskać
dane początkowe z końcowego profilu $u(T,\cdot)$,
nie jest ciągłe w normie supremum na swojej dziedzinie.
Odwrotność nie jest operatorem ograniczonym również
w normie $L^2$.
\end{kontrprzyklad}

\section{Nieoczekiwana niejednoznaczność: przykład Tychonoffa}
\label{sec:pde4-tychonoff}

Powyższa jednoznaczność wymagała ograniczoności
rozwiązania w całym cylindrze czasoprzestrzennym.
Pokażemy teraz, że nie jest to wyłącznie techniczny
warunek dowodu.

\begin{twierdzenie}[Istnienie niezerowego rozwiązania o zerowych danych]
\label{thm:pde4-tychonoff}
Istnieje niezerowa funkcja
$u\in C^\infty([0,\infty)\times\R)$
spełniająca $u_t=u_{xx}$ dla $t>0$, $x\in\R$,
przy czym $u(0,x)=0$ dla wszystkich $x\in\R$.
Funkcja ta nie może być ograniczona na żadnym
cylindrze $[0,T]\times\R$ z $T>0$.
\end{twierdzenie}
\begin{proof}
Niech $h(t)=e^{-1/t^2}$ dla $t>0$ i $h(0)=0$.
Kładziemy
\begin{equation}\label{eq:pde4-tychonoff}
 u(t,x)=\sum_{k=0}^\infty\frac{h^{(k)}(t)}{(2k)!}
 x^{2k}\quad(t>0),\qquad u(0,x)=0.
\end{equation}
Uzasadnimy najpierw zbieżność i możliwość
różniczkowania. Dla $t>0$ i $z\in\C$ z
$|z-t|\leq t/4$ mamy $\operatorname{Re}(1/z^2)
\geq c/t^2$ z pewną stałą $c>0$, niezależną od $t$.
Ponieważ $z\mapsto e^{-1/z^2}$ jest holomorficzna
poza zerem, wzór Cauchy'ego daje
\begin{equation}\label{eq:pde4-tychonoff-est}
 |h^{(k)}(t)|\leq k!\left(\frac4t\right)^k
 e^{-c/t^2},\qquad k\geq0.
\end{equation}
Ponieważ $(2k)!\geq(k!)^2$, dla $|x|\leq A$
otrzymujemy
\[
 |u(t,x)|\leq e^{-c/t^2}
 \sum_{k=0}^\infty\frac{(4A^2/t)^k}{k!}
 =\exp\!\left(-\frac c{t^2}+\frac{4A^2}t\right).
\]
Dowodzi to jednostajnej zbieżności do zera na
$|x|\leq A$ przy $t\downarrow0$.

Uzasadnijmy dokładniej oszacowanie wszystkich pochodnych.
Dla ustalonych $p,q\in\N_0$ i $2k\ge q$ mamy
\[
 \partial_t^p\partial_x^q
 \left(\frac{h^{(k)}(t)x^{2k}}{(2k)!}\right)
 =\frac{h^{(k+p)}(t)x^{2k-q}}{(2k-q)!}.
\]
Wybierzmy $B=\max\{1,A\}$ dla $|x|\le A$.
Oszacowanie~\eqref{eq:pde4-tychonoff-est} daje
\[
 \left|\frac{h^{(k+p)}(t)x^{2k-q}}{(2k-q)!}\right|
 \le e^{-c/t^2}\left(\frac4t\right)^p
 \frac{(k+p)!}{(2k-q)!}
 \left(\frac{4B^2}{t}\right)^k.
\]
Dla każdego ustalonego $p,q$ zachodzi
$(k+p)!/(2k-q)!\le C_{p,q}(1+k)^{p+q}/k!$
przy $2k\ge q$: istotnie, $(2k)!\ge(k!)^2$,
więc $(2k-q)!\ge(k!)^2/(2k)^q$ dla $k\ge1$,
a ponadto $(k+p)!\le C_p(1+k)^p k!$.
Dla skończenie wielu małych indeksów zwiększamy stałą.
Ponieważ dla $m\in\N_0$ szereg
$\sum_{k\ge0}(1+k)^m z^k/k!$ jest równy
$Q_m(z)e^z$ z wielomianem $Q_m$ stopnia $m$
(wynika to indukcyjnie przez działanie $z\partial_z$),
cały szereg pochodnych jest ograniczony na
$0<t\le1$, $|x|\le A$ przez
\[
 C_{p,q,A}t^{-p}(1+t^{-1})^{p+q}
 \exp\left(-\frac c{t^2}+\frac{4B^2}{t}\right).
\]
Wyrażenie to dąży do zera szybciej niż dowolna potęga
$t$ przy $t\downarrow0$. Wszystkie szeregi pochodnych
są lokalnie jednostajnie zbieżne dla $t>0$ oraz mają
jednostajnie zerowe granice w chwili $0$.
Kolejne zastosowanie twierdzenia o różniczkowaniu granicy
jednostajnej pokazuje, że przedłużenie $u(0,x)=0$
jest rzeczywiście klasy $C^\infty$ na $[0,\infty)\times\R$.
Na każdym zwartym przedziale $t\in[a,b]\subset(0,\infty)$
powyższe oszacowania zapewniają zbieżność jednostajną
wszystkich szeregów pochodnych na zwartych przedziałach
$x$, zatem wolno różniczkować wyraz po wyrazie.
Mamy wtedy
\[
 u_t(t,x)=\sum_{k=0}^\infty\frac{h^{(k+1)}(t)}{(2k)!}x^{2k}
 =u_{xx}(t,x).
\]
Wreszcie $u(t,0)=h(t)>0$ dla każdego $t>0$,
więc nie jest to rozwiązanie zerowe.
Gdyby było ograniczone na pewnym cylindrze $[0,T]\times\R$,
twierdzenie~\ref{thm:pde4-global-uniqueness}
wymusiłoby $u=0$, co jest sprzeczne z $u(t,0)>0$.
\end{proof}

\begin{uwaga}[Co ten kontrprzykład rzeczywiście oznacza?]
Niejednoznaczność nie pojawia się w naturalnych klasach
rozwiązań o kontrolowanym wzroście, na przykład w klasie
ograniczonej na skończonych cylindrach. Przykład Tychonoffa
pokazuje tylko, że po usunięciu wszelkiej kontroli
zachowania w nieskończoności sam warunek $u(0,\cdot)=0$
nie gwarantuje jednoznaczności.
\end{uwaga}

\section{Podsumowanie i związek z kolejnym rozdziałem}
\label{sec:pde4-summary}

Kluczowe mechanizmy równania ciepła są już widoczne
w~\eqref{eq:pde4-kernel}: dodatniość jądra wyjaśnia
zasadę zachowania porządku i natychmiastową propagację,
jego normalizacja --- zachowanie masy, a skalowanie
Gaussowskie --- coraz większe rozmywanie profilu.
Szybkość zanikania wysokich częstości prowadzi do
wygładzania, ale równocześnie czyni odtwarzanie
przeszłości niestabilnym.
W następnym rozdziale porównamy te fakty z równaniem
falowym, gdzie propagacja ma skończoną prędkość,
a klasyczna energia ma inny charakter.


% Rozdział 5. Równanie falowe.
\chapter[Równanie falowe]{Równanie falowe: charakterystyki, energia i propagacja}
\label{chap:pde-wave}

Równanie ciepła wygładza dane i przenosi niezerowe zaburzenie
na dowolną odległość w każdym dodatnim czasie. Równanie falowe
zachowuje się inaczej: zaburzenie nie przekracza określonej
prędkości, a oscylacje nie są samoczynnie tłumione. W tym
rozdziale wyprowadzimy jawne wzory, udowodnimy jednoznaczność
na podstawie energii i rozróżnimy słabą oraz silną zasadę
Huygensa.

Ustalmy $n\geq1$, $c>0$ i oznaczmy
$\square_c u=\partial_t^2u-c^2\Delta_xu$. Zagadnienie
Cauchy'ego ma postać
\begin{equation}\label{eq:pde5-wave}
 \square_cu=0\quad\text{dla }t>0,\ x\in\R^n,\qquad
 u(0,x)=g(x),\quad u_t(0,x)=h(x).
\end{equation}
Dwa warunki początkowe są potrzebne, ponieważ równanie
jest drugiego rzędu względem czasu. Przez
$\mathrm d\sigma$ oznaczamy standardową miarę
powierzchniową na gładkiej hiperpowierzchni, a przez
$\mathrm d\lambda_n(x)$ miarę Lebesgue'a. Wszystkie
rachunki z całkowaniem i różniczkowaniem wykonujemy
najpierw dla danych dostatecznie gładkich; słabe
rozwiązania wprowadzimy dopiero w kolejnych częściach.

\section{Równanie w jednym wymiarze: rozkład na fale biegnące}
\label{sec:pde5-dalembert}

Rozważmy $u_{tt}-c^2u_{xx}=0$ dla $(t,x)\in(0,\infty)\times\R$.
Operator rozkłada się na czynniki
$\partial_t^2-c^2\partial_x^2
=(\partial_t-c\partial_x)(\partial_t+c\partial_x)$.
Wprowadzając $\xi=x-ct$ i $\eta=x+ct$ otrzymujemy
$\partial_\xi\partial_\eta u=0$. Lokalnie każda
$C^2$-funkcja spełniająca to równanie ma więc postać
$u(t,x)=F(x-ct)+G(x+ct)$.

W rozdziale wprowadzającym otrzymaliśmy już wzór d'Alemberta
(twierdzenie~\ref{thm:dalembert}). Wracamy do niego tutaj,
aby powiązać reprezentację z charakterystykami, energią i
obszarem zależności. Powtórzony krótki dowód ma charakter
samodzielnego przypomnienia, nie stanowi odrębnego twierdzenia
o szerszych założeniach.

\begin{twierdzenie}[Wzór d'Alemberta]\label{thm:pde5-dalembert}
Niech $c>0$, $g\in C^2(\R)$ i $h\in C^1(\R)$.
Istnieje dokładnie jedno rozwiązanie
$u\in C^2([0,\infty)\times\R)$ zagadnienia
\eqref{eq:pde5-wave} dla $n=1$. Jest nim
\begin{equation}\label{eq:pde5-dalembert}
 u(t,x)=\frac{g(x-ct)+g(x+ct)}2
 +\frac1{2c}\int_{x-ct}^{x+ct}h(y)\,\mathrm d\lambda_1(y).
\end{equation}
\end{twierdzenie}
\begin{proof}
Warunki początkowe dla $F$ i $G$ dają
$F(x)+G(x)=g(x)$ oraz
$-cF'(x)+cG'(x)=h(x)$. Zatem
$F'=(g'-h/c)/2$ i $G'=(g'+h/c)/2$.
Po scałkowaniu i zsumowaniu otrzymujemy
\eqref{eq:pde5-dalembert} (stałe całkowania
znoszą się, ponieważ $F+G=g$).
Wzór definiuje funkcję $C^2$, a bezpośrednie
różniczkowanie po $t$ i $x$ daje równanie oraz
$u(0,x)=g(x)$, $u_t(0,x)=h(x)$.
Dla jednoznaczności niech $w$ będzie różnicą
dwóch rozwiązań. Współrzędne $\xi,\eta$ pokazują,
że $w=F(\xi)+G(\eta)$, a zerowe dane początkowe
wymuszają $F+G=0$ i $F'=G'=0$; stąd $w=0$.
\end{proof}

\begin{uwaga}[Dokładny obszar zależności]
Ze wzoru wynika, że wartość $u(t,x)$ zależy jedynie
od $g$ w dwóch punktach $x-ct,x+ct$ i od $h$ na
odcinku $[x-ct,x+ct]$. W szczególności dane z punktów
poza tym odcinkiem nie wpływają na $u(t,x)$.
Nie oznacza to, że w wymiarze jeden sygnał jest
odczuwany wyłącznie na czole fali: człon z $h$ całkuje
po całym wnętrzu stożka wstecznego.
\end{uwaga}

\begin{przyklad}[Dwa przesuwające się impulsy]
Jeżeli $h=0$ i $g\in C_c^2(\R)$, to
$u(t,x)=\frac12g(x-ct)+\frac12g(x+ct)$.
Początkowy impuls rozdziela się na dwa profile
biegnące w przeciwnych kierunkach bez zmiany kształtu.
Jeżeli natomiast $g=0$ i $h\geq0$ ma zwarty nośnik
oraz dodatnią całkę, to dla ustalonego $x$ i dużych $t$
cały nośnik $h$ leży w $[x-ct,x+ct]$, a
$u(t,x)=(2c)^{-1}\int_\R h(y)\,\mathrm d\lambda_1(y)>0$.
Ta druga sytuacja pokazuje, dlaczego silna zasada
Huygensa nie zachodzi w wymiarze jeden.
\end{przyklad}

\begin{center}
\begin{tikzpicture}[scale=.88,>=Stealth]
  \fill[blue!8] (-2.7,0)--(2.7,0)--(0,3.35)--cycle;
  \draw[->] (-4,0)--(4,0) node[right] {$x$};
  \draw[->] (0,-.25)--(0,3.9) node[above] {$t$};
  \draw[blue!65!black,thick] (-2.7,0)--(0,3.35)--(2.7,0);
  \fill (0,3.35) circle (1.8pt) node[above right] {$(t_0,x_0)$};
  \draw[<->,blue!50!black] (-2.7,-.3)--(2.7,-.3)
    node[midway,below] {$[x_0-ct_0,x_0+ct_0]$};
  \node[blue!65!black] at (0,1.05) {stożek wsteczny};
  \node at (3.4,3.2) {$x_0=0$};
\end{tikzpicture}
\par\smallskip\small
Obszar zależności w jednym wymiarze; osie i wierzchołek
narysowano dla $x_0=0$, a szerokość stożka odpowiada $c t_0$.
\end{center}

\section{Energia i jednoznaczność w wielu wymiarach}
\label{sec:pde5-energy}

\begin{definicja}[Energia falowa]
Dla $u\in C^2([0,T]\times\R^n)$ takiej, że
$u_t(t,\cdot)$ i $\nabla_xu(t,\cdot)$ mają zwarty nośnik,
definiujemy
\begin{equation}\label{eq:pde5-energy}
 E_u(t)=\frac12\int_{\R^n}
 \bigl(|u_t(t,x)|^2+c^2|\nabla_xu(t,x)|^2\bigr)
 \,\mathrm d\lambda_n(x).
\end{equation}
\end{definicja}

\begin{twierdzenie}[Zachowanie energii]\label{thm:pde5-energy}
Niech $c,T>0$, $n\ge1$, a $u\in C^2([0,T]\times\R^n)$
spełnia $\square_cu=0$. Załóżmy, że istnieje
zwarty zbiór $K\subset\R^n$, poza którym
$u_t(t,\cdot)=\nabla_xu(t,\cdot)=0$
dla każdego $t\in[0,T]$. Wówczas $E_u(t)=E_u(0)$
dla $0\leq t\leq T$.
\end{twierdzenie}
\begin{proof}
Różniczkując pod całką i całkując przez części w
$\R^n$, bez wyrazów brzegowych dzięki wspólnemu
zwartemu nośnikowi pochodnych, otrzymujemy
\[
 \frac{\mathrm d}{\mathrm dt}E_u(t)
 =\int_{\R^n}u_t(u_{tt}-c^2\Delta_xu)
 \,\mathrm d\lambda_n(x)=0.
\]
Różniczkowanie i całkowanie przez części są
uzasadnione gładkością i wspólnym zwartym nośnikiem.
\end{proof}

Najważniejszy wynik nie wymaga jednak założeń o zachowaniu
rozwiązania przy $|x|\to\infty$. Zamiast całej przestrzeni
rozważymy kurczące się kule.

\begin{twierdzenie}[Lokalna energia i skończona prędkość propagacji]
\label{thm:pde5-cone}
Niech $n\ge1$, $c,T>0$, $x_0\in\R^n$ oraz
$u\in C^2([0,T]\times\R^n)$ spełnia $\square_cu=0$.
Jeżeli $u(0,x)=u_t(0,x)=0$ dla wszystkich
$x\in\overline B(x_0,cT)$, to $u(T,x_0)=0$.
Co więcej, w stożku
$\{(t,x):0\le t<T,\ |x-x_0|<c(T-t)\}$
mamy $u(t,x)=0$.
\end{twierdzenie}
\begin{proof}
Dla $0\le t<T$ połóżmy $r(t)=c(T-t)$,
$B_t=B(x_0,r(t))$ i
$e=\frac12(u_t^2+c^2|\nabla_xu|^2)$. Definiujemy
$E(t)=\int_{B_t}e(t,x)\,\mathrm d\lambda_n(x)$.
Wzór na pochodną całki po ruchomej kuli otrzymujemy
bezpośrednio po zapisaniu jej we współrzędnych biegunowych:
$E(t)=\int_0^{r(t)}\!\int_{S^{n-1}}
e(t,x_0+\rho\theta)\rho^{n-1}\,\mathrm d\sigma(\theta)\,\mathrm d\lambda_1(\rho)$.
Różniczkowanie po czasie daje całkę z $e_t$ po $B_t$
oraz człon $r'(t)\int_{\partial B_t}e\,\mathrm d\sigma$.
Ponieważ $e_t=c^2\diver(u_t\nabla_xu)$,
po zastosowaniu twierdzenia Gaussa otrzymujemy
\[
\begin{aligned}
 E'(t)&=c^2\int_{\partial B_t}u_t\partial_\nu u
 \,\mathrm d\sigma-c\int_{\partial B_t}e\,\mathrm d\sigma\\
 &=-\frac c2\int_{\partial B_t}
 \bigl((u_t-c\partial_\nu u)^2+c^2|\nabla_{\tan}u|^2\bigr)
 \,\mathrm d\sigma\leq0.
\end{aligned}
\]
Tutaj $\nu$ jest zewnętrzną normalną jednostkową,
$\partial_\nu u=\nabla_xu\cdot\nu$, zaś
$\nabla_{\tan}u=\nabla_xu-(\partial_\nu u)\nu$.
Dla $n=1$ rozumiemy całkę brzegową jako sumę po
dwóch końcach przedziału, a składnik styczny jest
zerowy. Ponieważ $E(0)=0$ i $E(t)\geq0$, mamy
$E(t)=0$ dla wszystkich $t<T$. Z ciągłości
wynika $u_t=0$ i $\nabla_xu=0$ w całym otwartym
stożku. Każdy jego punkt można połączyć pionowym
odcinkiem z podstawą, gdzie $u=0$;
stąd $u=0$ w stożku. Ciągłość aż do wierzchołka
daje $u(T,x_0)=0$.
\end{proof}


\begin{twierdzenie}[Lokalne oszacowanie energii dla wymuszonego równania falowego]
\label{thm:pde5-cone-forcing}
Niech $n\ge1$, $c,T>0$, $x_0\in\R^n$, a
$u\in C^2([0,T]\times\R^n)$ spełnia
$u_{tt}-c^2\Delta_xu=f$ dla ciągłej funkcji
$f:[0,T]\times\R^n\to\R$.
Dla $0\le t<T$ oznaczmy $B_t=B(x_0,c(T-t))$ i
\[
 E(t)=\frac12\int_{B_t}
 \bigl(|u_t(t,x)|^2+c^2|\nabla_xu(t,x)|^2\bigr)
 \,\mathrm d\lambda_n(x).
\]
Wtedy dla każdego $0\le t<T$ zachodzi
\begin{equation}\label{eq:pde5-forced-cone-bound}
 \sqrt{E(t)}\le\sqrt{E(0)}+
 \frac1{\sqrt2}\int_0^t
 \|f(s,\cdot)\|_{L^2(B_s)}\,\mathrm d\lambda_1(s).
\end{equation}
W szczególności jeśli początkowe dane $u,u_t$ znikają
na $\overline B(x_0,cT)$, zaś $f=0$ w stożku
$\{(t,x):0\le t<T,\ x\in B_t\}$, to $u$ znika
w tym stożku oraz w jego wierzchołku $(T,x_0)$.
\end{twierdzenie}
\begin{proof}
Z twierdzenia o pochodnej całki po ruchomej kuli,
przeprowadzonego w dowodzie twierdzenia~\ref{thm:pde5-cone},
oraz z tożsamości
$\partial_t e=c^2\diver(u_t\nabla_xu)+u_tf$
otrzymujemy
\[
 E'(t)=-\frac c2\int_{\partial B_t}
 \bigl((u_t-c\partial_\nu u)^2+
 c^2|\nabla_{\tan}u|^2\bigr)\,\mathrm d\sigma
 +\int_{B_t}u_t(t,x)f(t,x)\,\mathrm d\lambda_n(x).
\]
Pierwszy składnik jest niedodatni, a nierówność
Cauchy'ego--Schwarza i definicja $E$ dają
$E'(t)\le\sqrt{2E(t)}\|f(t,\cdot)\|_{L^2(B_t)}$.
Nie wolno bez uzasadnienia dzielić przez $\sqrt{E(t)}$,
ponieważ energia może znikać. Dla $\delta>0$ piszemy
\[
 \frac{\mathrm d}{\mathrm dt}\sqrt{E(t)+\delta}
 \le\frac1{\sqrt2}\|f(t,\cdot)\|_{L^2(B_t)}.
\]
Całkujemy na $[0,t]$ i przechodzimy z $\delta\downarrow0$,
co daje \eqref{eq:pde5-forced-cone-bound}. Gdy prawa strona
jest zerowa, mamy $u_t=\nabla_xu=0$ w otwartym stożku.
Wartość $u$ jest wówczas stała na każdym pionowym odcinku
łączącym punkt stożka z jego podstawą, gdzie z założenia
wynosi zero. Ciągłość daje znikanie także w wierzchołku.
\end{proof}

\begin{wniosek}[Jednoznaczność bez ograniczeń wzrostu]
\label{cor:pde5-uniqueness}
Niech $n\ge1$, $c,T>0$ i niech
$u,v\in C^2([0,T]\times\R^n)$ rozwiązują to samo
jednorodne równanie falowe i mają takie same dane
początkowe w całej $\R^n$. Wówczas $u=v$.
\end{wniosek}
\begin{proof}
Do $w=u-v$ stosujemy twierdzenie~\ref{thm:pde5-cone}
dla dowolnego $(T',x_0)$, $0<T'\le T$;
dane $w$ są zerowe na dowolnej kuli wymaganej
w twierdzeniu, zatem $w(T',x_0)=0$.
\end{proof}

\begin{uwaga}[Kontrast z równaniem ciepła]
W rozdziale~\ref{chap:pde-heat} skonstruowaliśmy
niezerowe, gładkie rozwiązanie Tychonoffa o zerowych
danych początkowych, co wymagało porzucenia ograniczeń
wzrostu. Dla klasycznego równania falowego na całej
$\R^n$ taki fenomen nie występuje w klasie $C^2$:
lokalna tożsamość energetyczna zapewnia jednoznaczność
bez jakichkolwiek warunków przy $|x|\to\infty$.
\end{uwaga}

\section{Wzór Kirchhoffa w trzech wymiarach}
\label{sec:pde5-kirchhoff}

Dla $v\in C^2(\R^3)$, $x\in\R^3$ i $r>0$
wprowadzamy średnią sferyczną
\begin{equation}\label{eq:pde5-sphere}
 M_rv(x)=\frac1{4\pi}\int_{S^2}v(x+r\omega)
 \,\mathrm d\sigma_{S^2}(\omega),
\end{equation}
gdzie $S^2=\{\omega\in\R^3:|\omega|=1\}$,
a $\mathrm d\sigma_{S^2}$ jest miarą powierzchniową
o całkowitej masie $4\pi$. Definiujemy $M_0v(x)=v(x)$.

\begin{lemat}[Równanie dla średnich sferycznych]
\label{lem:pde5-darboux}
Jeżeli $v\in C^2(\R^3)$, to dla $r>0$
\begin{equation}\label{eq:pde5-darboux}
 \partial_r^2 M_rv(x)+\frac2r\partial_rM_rv(x)
 =M_r(\Delta v)(x).
\end{equation}
\end{lemat}
\begin{proof}
Dla ustalonego $x$ niech $V(r,\omega)=v(x+r\omega)$.
Rozkład laplasjanu na składową radialną i sferyczną daje
$\Delta v(x+r\omega)=\partial_r^2V+(2/r)\partial_rV
+r^{-2}\Delta_{S^2}V$. Całka po sferze z
$\Delta_{S^2}V$ jest zerowa, ponieważ sfera nie
ma brzegu (można to sprawdzić przez lokalną
formułę dywergencji na sferze). Całkujemy otrzymaną
tożsamość po $S^2$ i dzielimy przez $4\pi$.
\end{proof}

\begin{twierdzenie}[Wzór Kirchhoffa]
\label{thm:pde5-kirchhoff}
Niech $c>0$, $g\in C^3(\R^3)$ oraz
$h\in C^2(\R^3)$. Wówczas istnieje dokładnie
jedno rozwiązanie $u\in C^2([0,\infty)\times\R^3)$
zagadnienia~\eqref{eq:pde5-wave}, dane wzorem
\begin{equation}\label{eq:pde5-kirchhoff}
 u(t,x)=\frac{\mathrm d}{\mathrm dt}
 \bigl(t M_{ct}g(x)\bigr)+t M_{ct}h(x),\qquad t\ge0.
\end{equation}
\end{twierdzenie}
\begin{proof}
Najpierw załóżmy, że $u$ jest rozwiązaniem $C^2$.
Dla ustalonego $x$ kładziemy
$W(t,r)=rM_ru(t,\cdot)(x)$, $r\ge0$. Z lematu
\ref{lem:pde5-darboux} otrzymujemy dla $r>0$
$W_{tt}-c^2W_{rr}=0$, gdyż
$\partial_r^2(rM_ru)=rM_r\Delta u$.
Mamy $W(t,0)=0$,
$W(0,r)=rM_rg(x)$ i $W_t(0,r)=rM_rh(x)$.
Przedłużając dane nieparzyście na $r<0$
i stosując wzór d'Alemberta z twierdzenia
\ref{thm:pde5-dalembert}, otrzymujemy
$W_r(t,0)=\frac{\mathrm d}{\mathrm dr}(rM_rg)|_{r=ct}
 +c^{-1}(rM_rh)|_{r=ct}$.
Ponieważ $W_r(t,0)=u(t,x)$, dostajemy
\eqref{eq:pde5-kirchhoff}.

Dla istnienia zdefiniujmy dla $v\in C^2(\R^3)$
$V_v(t,x)=tM_{ct}v(x)$. Z tożsamości Darboux
otrzymujemy dla $t>0$ bezpośrednio
\[
 (\partial_t^2-c^2\Delta_x)V_v
 =t c^2\left(\partial_r^2M_rv+\frac2r\partial_rM_rv
              -M_r\Delta v\right)\bigg|_{r=ct}=0.
\]
Dla $g\in C^3(\R^3)$ otrzymujemy $\partial_tV_g\in C^2$,
a dla $h\in C^2(\R^3)$ także $V_h\in C^2$.
Aby uzasadnić, że $\square_c(\partial_tV_g)=0$
przy tej skończonej regularności, aproksymujemy $g$
lokalnie w $C^3$ przez funkcje gładkie; skończony
zasięg całek sferycznych daje lokalną zbieżność
$\partial_tV_{g_j}\to\partial_tV_g$ w $C^2$.
Dla danych gładkich można skomutować $\square_c$
z $\partial_t$, a przejście do granicy dowodzi żądanej równości.
Ponadto $M_0v=v$ i $\partial_rM_rv|_{r=0}=0$,
ponieważ średnia z $\omega$ na $S^2$ jest zerowa.
Zatem $u=\partial_tV_g+V_h$ spełnia
$u(0,x)=g(x)$ i $u_t(0,x)=h(x)$.
Jednoznaczność wynika z wniosku~\ref{cor:pde5-uniqueness}.
\end{proof}

\begin{uwaga}[Jakie dane widzi punkt w trzech wymiarach?]
Dla ustalonego $(t,x)$ wartość rozwiązania
zależy od $g,h$ (i pochodnych $g$) na sferze
$\{y\in\R^3:|y-x|=ct\}$. Dane z wnętrza kuli,
ale z dala od tej sfery, nie mają wpływu na
$u(t,x)$. Jest to silna zasada Huygensa dla
równania falowego w wymiarze trzy.
\end{uwaga}

\section{Dwa wymiary i zasada Huygensa}
\label{sec:pde5-huygens}

\begin{twierdzenie}[Wzór Poissona w wymiarze dwa]
\label{thm:pde5-2d}
Niech $c>0$, $g\in C_c^3(\R^2)$, $h\in C_c^2(\R^2)$.
Jedyne rozwiązanie $C^2$ równania falowego w
$\R^2$ z tymi danymi jest dane wzorem
\begin{equation}\label{eq:pde5-2d}
\begin{aligned}
 u(t,x)={}&\frac{\mathrm d}{\mathrm dt}
 \left[\frac1{2\pi c}\int_{|y-x|<ct}
 \frac{g(y)}{\sqrt{c^2t^2-|y-x|^2}}\,
 \mathrm d\lambda_2(y)\right]\\
 &+\frac1{2\pi c}\int_{|y-x|<ct}
 \frac{h(y)}{\sqrt{c^2t^2-|y-x|^2}}\,
 \mathrm d\lambda_2(y),\qquad t>0.
\end{aligned}
\end{equation}
\end{twierdzenie}
\begin{proof}
Rozszerzamy $g,h$ na $\R^3$ tak, aby były niezależne
od trzeciej zmiennej: $\widetilde g(y,z)=g(y)$,
$\widetilde h(y,z)=h(y)$. Wzór Kirchhoffa definiuje
wtedy rozwiązanie trójwymiarowe niezależne od $z$
(jednoznaczność ze wniosku~\ref{cor:pde5-uniqueness}).
Dla $r>0$ parametryzujemy dwie półsfery
$|y-x|^2+z^2=r^2$ przez
$z=\pm\sqrt{r^2-|y-x|^2}$. Element pola na każdej
półsferze wynosi
$r(r^2-|y-x|^2)^{-1/2}\,\mathrm d\lambda_2(y)$.
Dlatego dla $\varphi$ niezależnej od $z$ mamy
\[
 M_r\varphi(x,0)=\frac1{2\pi r}
 \int_{|y-x|<r}\frac{\varphi(y)}{\sqrt{r^2-|y-x|^2}}
 \,\mathrm d\lambda_2(y).
\]
Wstawiając $r=ct$ do wzoru Kirchhoffa,
otrzymujemy dokładnie \eqref{eq:pde5-2d}.
\end{proof}

Zwróćmy uwagę na różnicę geometryczną: w wymiarze
trzy całkowanie odbywa się po \emph{sferze},
a w wymiarze dwa po \emph{całym dysku} z osobliwą,
ale całkowalną wagą. Dla $g=0$ i nieujemnej,
niezerowej funkcji $h\in C_c^2(\R^2)$ druga całka
jest ściśle dodatnia, gdy dysk $B(x,ct)$ obejmuje
nośnik $h$ (przy $t$ odpowiednio dużym).
Oznacza to ogon fali, którego nie ma w wersji
trójwymiarowej.

\begin{uwaga}[Silna i słaba zasada Huygensa]
Słaba zasada oznacza, że wpływ danych początkowych
nie może wyjść poza domknięty stożek $|x-y|\le ct$.
W przedstawionym modelu jest to wniosek z
\ref{thm:pde5-cone}. Silna zasada oznacza zależność
od danych wyłącznie na czole stożka, czyli na
$|x-y|=ct$. Dla stałych współczynników w przestrzeni
Euklidesowej silna zasada zachodzi w wymiarach
nieparzystych $n\ge3$, lecz zawodzi w wymiarach
parzystych i w wymiarze jeden. Uogólnienia
na wszystkie wymiary wymagają dalszych rachunków
z rozwiązaniami fundamentalnymi; na tym etapie
udowodniliśmy przypadki $n=1,2,3$.
\end{uwaga}

\section{Zewnętrzne wymuszenie i wzór Duhamela dla fal}
\label{sec:pde5-forcing}

\begin{twierdzenie}[Wzór dla niejednorodnego równania w jednym wymiarze]
\label{thm:pde5-forcing}
Niech $c,T>0$, $g\in C^2(\R)$, $h\in C^1(\R)$,
$f\in C^1([0,T]\times\R)$ i załóżmy, że wszystkie
te funkcje oraz pochodne występujące we wzorze są
ciągłe (nie wymagamy ograniczoności w nieskończoności).
Wówczas funkcja
\begin{equation}\label{eq:pde5-forcing}
\begin{aligned}
 u(t,x)={}&\frac{g(x-ct)+g(x+ct)}2
 +\frac1{2c}\int_{x-ct}^{x+ct}h(y)\,\mathrm d\lambda_1(y)\\
 &+\frac1{2c}\int_0^t\int_{x-c(t-s)}^{x+c(t-s)}
 f(s,y)\,\mathrm d\lambda_1(y)\,\mathrm ds
\end{aligned}
\end{equation}
rozwiązuje $u_{tt}-c^2u_{xx}=f$, $u(0)=g$,
$u_t(0)=h$ na $[0,T]\times\R$. Jest jedyna
w klasie rozwiązań $C^2$.
\end{twierdzenie}
\begin{proof}
Jednorodna część rozwiązuje problem bez wymuszenia.
Oznaczmy trzeci składnik przez $v(t,x)$.
Różniczkując dwukrotnie z wykorzystaniem
reguły Leibniza otrzymujemy
\[
 v_t(t,x)=\frac12\int_0^t\bigl[
 f(s,x+c(t-s))+f(s,x-c(t-s))\bigr] \,\mathrm ds
\]
oraz
\[
 v_{tt}(t,x)=f(t,x)+\frac c2\int_0^t
 \bigl[f_x(s,x+c(t-s))-f_x(s,x-c(t-s))\bigr]
 \,\mathrm ds.
\]
Jednocześnie
$v_{xx}(t,x)=\frac1{2c}\int_0^t
 \bigl[f_x(s,x+c(t-s))-f_x(s,x-c(t-s))\bigr]
 \,\mathrm ds$.
Stąd $v_{tt}-c^2v_{xx}=f$ i
$v(0,x)=v_t(0,x)=0$. Jednoznaczność wynika
z wniosku~\ref{cor:pde5-uniqueness} zastosowanego
do różnicy rozwiązań.
\end{proof}

\section{Warunki brzegowe: struna na odcinku}
\label{sec:pde5-string}

Na odcinku $\Omega=(0,L)$ rozważmy
$u_{tt}-c^2u_{xx}=0$ z warunkiem
$u(t,0)=u(t,L)=0$. Zagadnienie można rozwiązać
metodą rozdzielenia zmiennych. Dla $m\in\N$
połóżmy $\omega_m=cm\pi/L$; wtedy
\[
 u_m(t,x)=\bigl(a_m\cos(\omega_mt)
 +b_m\sin(\omega_mt)\bigr)\sin(m\pi x/L)
\]
jest rozwiązaniem klasycznym. Jeśli zadane dane
początkowe są skończonymi kombinacjami funkcji
$\sin(m\pi x/L)$, dobieramy
$a_m$ jako współczynniki danych $g$,
a $b_m$ jako współczynniki $h$ podzielone przez
$\omega_m$.

\begin{przyklad}[Brak tłumienia modów]
Dla $g(x)=\sin(m\pi x/L)$ i $h=0$ dostajemy
$u(t,x)=\cos(\omega_mt)\sin(m\pi x/L)$.
Amplituda nie maleje monotonicznie, podczas gdy
odpowiedni mod równania ciepła zanika jak
$e^{-\kappa(m\pi/L)^2t}$. Właśnie dlatego
regularizujący charakter półgrupy cieplnej
nie przenosi się na grupę ewolucji falowej.
\end{przyklad}

\begin{stwierdzenie}[Energia struny]
\label{prop:pde5-stringenergy}
Jeśli $u\in C^2([0,T]\times[0,L])$ spełnia
$u_{tt}=c^2u_{xx}$ dla $0<t<T$, $0<x<L$,
a także $u(t,0)=u(t,L)=0$, to
\[
 E(t)=\frac12\int_0^L(u_t^2+c^2u_x^2)\,
 \mathrm d\lambda_1(x)
\]
jest stała na $[0,T]$.
\end{stwierdzenie}
\begin{proof}
Z warunku Dirichleta i regularności mamy
$u_t(t,0)=u_t(t,L)=0$ dla $0<t<T$. Zatem
$E'(t)=c^2[u_tu_x]_{x=0}^{x=L}=0$.
Ciągłość energii daje równość także w chwilach
$t=0,T$.
\end{proof}

\section{Dlaczego schemat numeryczny musi respektować stożek?}
\label{sec:pde5-numerics}

Na regularnej siatce $x_j=j\Delta x$, $t_m=m\Delta t$
przybliżamy drugie pochodne różnicami centralnymi:
\begin{equation}\label{eq:pde5-leapfrog}
 \frac{U_j^{m+1}-2U_j^m+U_j^{m-1}}{(\Delta t)^2}
 =c^2\frac{U_{j+1}^m-2U_j^m+U_{j-1}^m}{(\Delta x)^2}.
\end{equation}
Wprowadzamy liczbę Couranta $\lambda=c\Delta t/\Delta x$.
Dla modu $U_j^m=e^{i(kj\Delta x-\omega m\Delta t)}$
otrzymujemy relację dyspersji
\begin{equation}\label{eq:pde5-cfl}
 \sin^2\!\left(\frac{\omega\Delta t}2\right)
 =\lambda^2\sin^2\!\left(\frac{k\Delta x}2\right).
\end{equation}
Jeżeli $\lambda\le1$, istnieją rzeczywiste
częstotliwości $\omega$ dla wszystkich częstości
siatkowych. Gdy $\lambda>1$, istnieją częstości,
dla których prawa strona przekracza $1$, co
prowadzi do wykładniczego wzrostu rozwiązań
rekurencji. Jest to konieczny warunek stabilności
metody centralnej; w tym konkretnym jednowymiarowym
schemacie z prawidłowo zadanymi danymi
$\lambda\le1$ jest również klasycznym warunkiem
stabilności w odpowiedniej normie energetycznej
(z dodatkową ostrożnością dla częstości granicznych
przy $\lambda=1$). Pełny dowód numeryczny
pozostawimy do działu metod różnicowych.

\begin{uwaga}[Granica stabilności przy $\lambda=1$]
Dla częstotliwości siatkowej $k\Delta x=\pi$ amplituda
modu spełnia rekurencję $A_{m+1}+2A_m+A_{m-1}=0$.
Jej ogólne rozwiązanie to $(\alpha+\beta m)(-1)^m$:
podwójny pierwiastek $-1$ dopuszcza więc liniowy wzrost
przy dowolnie zadanych dwóch pierwszych warstwach.
Dla prędkości początkowej równej zeru i poprawnego
pierwszego kroku $A_1=-A_0$, skąd $\beta=0$.
Analiza stabilności w przypadku granicznym musi uwzględniać
przestrzeń danych i sposób inicjalizacji schematu;
proste kryterium rzeczywistości $\omega$ nie wystarcza.
\end{uwaga}

\begin{uwaga}[Zasięg numeryczny a zasięg fizyczny]
Schemat \eqref{eq:pde5-leapfrog} korzysta z
sąsiadów $j-1,j,j+1$ w poprzedniej warstwie
czasowej. Jeden krok może więc przenieść informację
o odległość $\Delta x$, podczas gdy fizyczna fala
przebywa $c\Delta t$. Warunek $c\Delta t\le\Delta x$
pozwala fizycznemu stożkowi zmieścić się w
numerycznym obszarze zależności. To tłumaczy
geometrię warunku CFL, ale nie zastępuje
dowodu stabilności.
\end{uwaga}

\section{Podsumowanie}
\label{sec:pde5-summary}

Dla równania falowego rozwiązanie w punkcie jest
wyznaczone przez dane w skończonym obszarze
początkowym; wynika to z kurczącej się energii
stożkowej. Równania ciepła i falowe są pod tym
względem fundamentalnie odmienne. Wzory d'Alemberta,
Kirchhoffa i Poissona pokazują również, że w samym
równaniu falowym silna zasada Huygensa zależy
od parzystości wymiaru. Następny rozdział poświęcimy
systematycznemu wykorzystaniu transformaty i
szeregów Fouriera, łącząc otrzymane tu konstrukcje
z metodami spektralnymi.

% ROZDZIAŁ 6: do wstawienia przed rozdziałem z kodem numerycznym.
\chapter[Metody Fouriera i rozdzielenie zmiennych]{Szeregi Fouriera, rozdzielenie zmiennych i zagadnienia Sturma--Liouville'a}
\label{chap:pde-fourier}

W rozdziałach~\ref{chap:pde-heat} i~\ref{chap:pde-wave} uzyskaliśmy rozwiązania na
całej przestrzeni i w prostych obszarach. Na przedziale lub prostokącie
bardziej naturalne jest korzystanie z \emph{modów własnych} operatora
przestrzennego. Metoda rozdzielenia zmiennych będzie skuteczna tylko wtedy,
gdy uzasadnimy, że funkcje własne tworzą dostatecznie bogatą rodzinę.
Najpierw przeprowadzimy odpowiedni rachunek dla szeregów Fouriera,
a następnie zbadamy jego wersję dla regularnych zagadnień
Sturma--Liouville'a. Formalny rachunek dystrybucji przedstawimy w rozdziale~\ref{chap:pde-distributions},
a transformatę Fouriera na $\R^n$ i dystrybucje temperowane ---
w rozdziale~\ref{chap:pde-fourier-Rd}, po wprowadzeniu
przestrzeni Schwartza. Rozróżnienie szeregów Fouriera na torusie
i transformaty na całej przestrzeni jest istotne także dla normalizacji.

\section{Szereg Fouriera: współczynniki, ortogonalność i zbieżność}
\label{sec:pde6-fourier}

Przez $\mathbb T=\R/(2\pi\mathbb Z)$ rozumiemy okrąg z
współrzędną $x\in[-\pi,\pi]$ i utożsamionymi końcami.
Dla $f\in L^1(-\pi,\pi)$ definiujemy
$\widehat f(k)=(2\pi)^{-1}\int_{-\pi}^{\pi}f(x)e^{-ikx}\,\mathrm d\lambda_1(x)$,
$k\in\mathbb Z$, oraz sumę częściową
$S_Nf(x)=\sum_{|k|\le N}\widehat f(k)e^{ikx}$.
W tym rozdziale $L^2(\mathbb T)$ oznacza zespoloną przestrzeń Hilberta
z iloczynem skalarnym
$\langle f,g\rangle=(2\pi)^{-1}\int_{-\pi}^{\pi}f(x)\overline{g(x)}\,\mathrm d\lambda_1(x)$.
Pojęcie całki Lebesgue'a i podstawowe własności $L^p$ traktujemy
jako znane; szerszą teorię funkcjonalną rozwiniemy w dalszych tomach.
Zwracamy uwagę na normalizację: w tym rozdziale norma $L^2(\T)$
pochodzi od miary $\mathrm d\lambda_1(x)/(2\pi)$.
W dalszych rozdziałach o PDE na torusie stosujemy zwykle
nieznormalizowaną miarę $\mathrm d\lambda_d(x)$; wzory
na wartości własne i współczynniki Fouriera pozostają zgodne,
ale normy $L^2$ różnią się odpowiednim czynnikiem.

\begin{lemat}[Ortogonalność i nierówność Bessela]
\label{lem:pde6-bessel}
Niech $f\in L^2(\mathbb T)$ oraz $N\ge0$. Wówczas
$\langle e^{ikx},e^{i\ell x}\rangle=\delta_{k\ell}$,
$\|S_Nf\|_2^2=\sum_{|k|\le N}|\widehat f(k)|^2$ i
$\sum_{|k|\le N}|\widehat f(k)|^2\le\|f\|_2^2$.
\end{lemat}
\begin{proof}
Pierwsza równość wynika z całki z $e^{i(k-\ell)x}$.
Reszta pierwszej tożsamości to twierdzenie Pitagorasa dla skończonego
układu ortonormalnego. Ponadto $f-S_Nf$ jest prostopadłe do każdego
$e^{ikx}$, $|k|\le N$, toteż
$\|f\|_2^2=\|S_Nf\|_2^2+\|f-S_Nf\|_2^2$.
\end{proof}

Nie można jednak wywnioskować z ortogonalności, że $S_Nf\to f$.
Do tego potrzebujemy zupełności układu wykładniczego.

\begin{lemat}[Jądro Fejéra]
\label{lem:pde6-fejer}
Dla $N\ge0$ niech
\begin{equation}\label{eq:pde6-fejer}
 F_N(t)=\frac1{N+1}\left|\sum_{j=0}^Ne^{ijt}\right|^2
 =\frac1{N+1}\left(\frac{\sin((N+1)t/2)}{\sin(t/2)}\right)^2
\end{equation}
(z wartością $F_N(0)=N+1$). Wówczas $F_N\ge0$,
$(2\pi)^{-1}\int_{-\pi}^{\pi}F_N(t)\,\mathrm d\lambda_1(t)=1$ oraz
dla każdego $\delta\in(0,\pi)$
$\int_{\delta\le|t|\le\pi}F_N(t)\,\mathrm d\lambda_1(t)\to0$.
\end{lemat}
\begin{proof}
Rozwinięcie kwadratu modułu daje
$F_N(t)=\sum_{|k|\le N}(1-|k|/(N+1))e^{ikt}$.
Wyraz stały jest równy $1$, skąd normalizacja.
Gdy $\delta\le|t|\le\pi$, mianownik drugiego wzoru
jest ograniczony od zera, a licznik przez $1$;
stąd $F_N(t)\le[(N+1)\sin^2(\delta/2)]^{-1}$.
\end{proof}

\begin{twierdzenie}[Fejér i zupełność układu trygonometrycznego]
\label{thm:pde6-fejer}
Niech $f\in C(\mathbb T)$. Średnie Cesàra
$\sigma_N f=(N+1)^{-1}\sum_{j=0}^N S_jf$
zbieżne są jednostajnie do $f$ na $\mathbb T$.
W konsekwencji wielomiany trygonometryczne są gęste
w $C(\mathbb T)$ w normie supremum i w $L^2(\mathbb T)$
w normie $L^2$.
\end{twierdzenie}
\begin{proof}
Z lematu~\ref{lem:pde6-fejer} wynika
$\sigma_N f(x)=(2\pi)^{-1}\int_{-\pi}^{\pi}
F_N(t)f(x-t)\,\mathrm d\lambda_1(t)$.
Odejmujemy $f(x)$ i dzielimy całkę na $|t|<\delta$
oraz $\delta\le|t|\le\pi$. Z jednostajnej ciągłości $f$
możemy dobrać $\delta$ tak, aby pierwszy składnik
miał moduł co najwyżej $\varepsilon$ dla wszystkich $x,N$.
Drugi ma moduł nie większy niż
$2\|f\|_\infty(2\pi)^{-1}\int_{|t|\ge\delta}F_N(t)
\,\mathrm d\lambda_1(t)$, co dąży do zera.
Gęstość w $L^2$ wynika dodatkowo z gęstości
funkcji ciągłych w $L^2$ na zwartym przedziale,
którą można udowodnić aproksymacją funkcji prostych
przez funkcje ciągłe (dla zbiorów mierzalnych używamy
regularności miary Lebesgue'a).
\end{proof}

\begin{twierdzenie}[Zbieżność w $L^2$ i tożsamość Parsevala]
\label{thm:pde6-parseval}
Dla dowolnego $f\in L^2(\mathbb T)$ zachodzą
$S_Nf\to f$ w normie $L^2$ oraz
\begin{equation}\label{eq:pde6-parseval}
 \frac1{2\pi}\int_{-\pi}^{\pi}|f(x)|^2\,\mathrm d\lambda_1(x)
 =\sum_{k\in\mathbb Z}|\widehat f(k)|^2.
\end{equation}
\end{twierdzenie}
\begin{proof}
Dla $\varepsilon>0$ wybieramy wielomian trygonometryczny $p$
z $\|f-p\|_2<\varepsilon$. Dla $N$ większych niż stopień
$p$ mamy $S_Np=p$. Ponieważ $S_N$ jest rzutem ortogonalnym,
$\|S_N(f-p)\|_2\le\|f-p\|_2$; zatem
$\|f-S_Nf\|_2\le2\varepsilon$.
Przechodząc do granicy w tożsamości Pitagorasa z
lematu~\ref{lem:pde6-bessel}, otrzymujemy Parsevala.
\end{proof}

\begin{uwaga}[Jakiego rodzaju zbieżność została udowodniona?]
Zbieżność $L^2$ nie oznacza zbieżności jednostajnej ani
punktowej. Dla funkcji ciągłych jednostajnie zbieżne są
\emph{średnie Fejéra}, niekoniecznie zwykłe sumy $S_N$.
W dalszej części do różniczkowania szeregów użyjemy
gładkich danych, dla których zachodzi dodatkowe szybkie
zanikanie współczynników.
\end{uwaga}

\begin{lemat}[Zanikanie współczynników gładkiej funkcji]
\label{lem:pde6-decay}
Jeśli $f\in C^m(\mathbb T)$, $m\ge1$, to dla $k\ne0$
$|\widehat f(k)|\le\|f^{(m)}\|_{L^1(-\pi,\pi)}/(2\pi|k|^m)$.
Jeśli $f\in C^\infty(\mathbb T)$, to szereg Fouriera wraz
ze wszystkimi szeregami otrzymanymi przez różniczkowanie
zbiega jednostajnie.
\end{lemat}
\begin{proof}
$m$-krotnie całkujemy przez części; wyrazy brzegowe
znikają dzięki zgodności pochodnych na końcach przedziału.
Dla $f\in C^\infty$ i ustalonego $r\ge0$ wybieramy
$m\ge r+2$ i szacujemy szereg
$\sum_{k\ne0}|k|^r|\widehat f(k)|$ przez
stałą pomnożoną przez $\sum_{k\ne0}|k|^{r-m}<\infty$.
\end{proof}

\section{Zagadnienie własne laplasjanu na odcinku}
\label{sec:pde6-dirichlet}

Dla $L>0$ rozważmy
\begin{equation}\label{eq:pde6-eigen}
 -X''(x)=\lambda X(x),\quad 0<x<L,\qquad X(0)=X(L)=0.
\end{equation}
\begin{twierdzenie}[Pełne spektrum Dirichleta w jednym wymiarze]
\label{thm:pde6-spectrum}
Dla $L>0$ zagadnienie~\eqref{eq:pde6-eigen} ma
nietrywialne rozwiązanie $X\in C^2([0,L])$
wtedy i tylko wtedy, gdy
$\lambda=\lambda_n=(n\pi/L)^2$ dla pewnego $n\in\N$.
Jego przestrzeń rozwiązań jest rozpięta przez
$X_n(x)=\sin(n\pi x/L)$. Ponadto funkcje
$e_n(x)=\sqrt{2/L}\sin(n\pi x/L)$, $n\ge1$,
tworzą bazę ortonormalną w $L^2(0,L)$.
\end{twierdzenie}
\begin{proof}
Mnożąc~\eqref{eq:pde6-eigen} przez $\overline X$
i całkując przez części otrzymujemy
$\lambda\int_0^L|X(x)|^2\,\mathrm d\lambda_1(x)
=\int_0^L|X'(x)|^2\,\mathrm d\lambda_1(x)$.
Dla nietrywialnego $X$ mamy $\lambda>0$ i jest rzeczywista.
Ogólne rozwiązanie ODE to
$A\sin(\sqrt\lambda x)+B\cos(\sqrt\lambda x)$.
Warunek w $0$ daje $B=0$, a warunek w $L$ wymusza
$\sqrt\lambda L\in\pi\N$. Ortogonalność i normowanie
sprawdza się całkowaniem elementarnym.
Dla zupełności weźmy $f\in L^2(0,L)$ i przedłużmy
ją nieparzyście do $\widetilde f\in L^2(-L,L)$,
a następnie $2L$-okresowo. Przeskalowana wersja
twierdzenia~\ref{thm:pde6-parseval} daje zbieżność
szeregu Fouriera $\widetilde f$ w $L^2(-L,L)$.
Przez nieparzystość wszystkie współczynniki cosinusowe
znikają; po ograniczeniu do $(0,L)$ pozostają wyłącznie
mody $\sin(n\pi x/L)$.
\end{proof}

\begin{wniosek}[Nierówność Poincarégo dla funkcji gładkich]
\label{cor:pde6-poincare}
Jeżeli $v\in C^1([0,L])$, $v(0)=v(L)=0$, to
\begin{equation}\label{eq:pde6-poincare}
 \int_0^L|v(x)|^2\,\mathrm d\lambda_1(x)
 \le \left(\frac L\pi\right)^2
 \int_0^L|v'(x)|^2\,\mathrm d\lambda_1(x).
\end{equation}
Stała $(L/\pi)^2$ jest optymalna.
\end{wniosek}
\begin{proof}
Niech $v_n=\langle v,e_n\rangle_{L^2(0,L)}$.
Całkowanie przez części daje
$\langle v',\sqrt{2/L}\cos(n\pi x/L)\rangle
=(n\pi/L)v_n$.
Rodzina kosinusów dla $n\ge1$ jest ortonormalna;
przez nierówność Bessela
$\sum_{n\ge1}(n\pi/L)^2|v_n|^2\le\|v'\|_2^2$.
Z Parsevala dla $v$ dostajemy~\eqref{eq:pde6-poincare}.
Równość zachodzi dla $v(x)=\sin(\pi x/L)$.
\end{proof}

\begin{uwaga}[Po co zachowywać stałe w oszacowaniach?]
Nierówność~\eqref{eq:pde6-poincare} wzmacnia
oszacowanie elementarne uzyskane w rozdziale~\ref{chap:pde-heat}.
Najmniejsza wartość własna $\lambda_1=\pi^2/L^2$
wyznacza dokładną szybkość zaniku pierwszego modu
równania ciepła i najniższą częstość drgań struny.
\end{uwaga}

\section{Rozdzielenie zmiennych: równanie ciepła}
\label{sec:pde6-heat}

Niech $\kappa,L>0$, a $u_t=\kappa u_{xx}$ dla $t>0$,
$x\in(0,L)$, z $u(t,0)=u(t,L)=0$ i $u(0,x)=g(x)$.
Podstawienie $u(t,x)=T(t)X(x)$ dla niezerowego iloczynu
daje $T'/(\kappa T)=X''/X=-\lambda$.
Dla $\lambda=\lambda_n$ otrzymujemy $T(t)=e^{-\kappa\lambda_nt}$.


\begin{lemat}[Jądro Gaussa i kontrakcja na odcinku]
\label{lem:pde6-dirichlet-heat}
Niech $L,\kappa>0$ i niech
$g\in C([0,L])$ spełnia $g(0)=g(L)=0$.
Oznaczmy przez $\widetilde g$ jej nieparzyste
$2L$-okresowe przedłużenie na $\R$. Dla $t>0$ połóżmy
\[
 (Q_tg)(x)=\int_{\R}(4\pi\kappa t)^{-1/2}
 e^{-z^2/(4\kappa t)}\widetilde g(x-z)
 \,\mathrm d\lambda_1(z),\qquad x\in[0,L].
\]
Wówczas $Q_tg$ znika dla $x=0,L$,
$\|Q_tg\|_{C([0,L])}\le\|g\|_{C([0,L])}$ oraz
$\|Q_tg-g\|_{C([0,L])}\to0$ dla $t\downarrow0$.
Ponadto współczynnik sinusowy $Q_tg$ odpowiadający
$\sin(n\pi x/L)$ jest równy
$e^{-\kappa(n\pi/L)^2t}b_n$, gdzie
$b_n=(2/L)\int_0^Lg(x)\sin(n\pi x/L)\,\mathrm d\lambda_1(x)$.
\end{lemat}
\begin{proof}
Przedłużenie jest ciągłe i jednostajnie ciągłe na $\R$,
a jego norma supremum jest równa normie $g$.
Jądro Gaussa jest dodatnie i ma całkę $1$, co daje
kontrakcję. Ponieważ splot funkcji parzystej z funkcją
nieparzystą jest nieparzysty, a splot zachowuje okres
$2L$, otrzymana funkcja znika dla $x=0$ i $x=L$.
Dla dowolnego $\delta>0$ rozdzielamy całkę splotową
na $|z|<\delta$ i $|z|\ge\delta$. Na pierwszej części
różnica $|\widetilde g(x-z)-\widetilde g(x)|$ jest
jednostajnie mała z jednostajnej ciągłości, a na drugiej
nie przekracza $2\|g\|_\infty$ i masa Gaussa dąży do
zera przy $t\downarrow0$. Daje to zbieżność jednostajną.
Aby znaleźć współczynniki, zapisujemy splot jako
splot okresowy z jądrem
$\sum_{k\in\mathbb Z}G_t(\cdot+2kL)$.
Szereg jąder zbiega jednostajnie na okresie dla $t>0$;
można zatem zastosować twierdzenie Fubiniego.
Dla częstości $n\pi/L$ współczynnik tego jądra
jest równy całce Gaussowskiej
$\int_{\R}G_t(z)e^{-in\pi z/L}\,\mathrm d\lambda_1(z)
=e^{-\kappa(n\pi/L)^2t}$.
Ponieważ $\widetilde g$ jest nieparzysta, współczynniki
cosinusowe znikają, co daje tezę.
\end{proof}

\begin{twierdzenie}[Rozwiązanie przez szereg sinusowy]
\label{thm:pde6-heat}
Niech $\kappa,L>0$ i $g\in C([0,L])$ spełnia
$g(0)=g(L)=0$. Definiujemy
$b_n=(2/L)\int_0^Lg(x)\sin(n\pi x/L)\,\mathrm d\lambda_1(x)$.
Wówczas
\begin{equation}\label{eq:pde6-heat}
 u(t,x)=\sum_{n=1}^\infty b_n e^{-\kappa\lambda_nt}
 \sin(n\pi x/L),\qquad t>0,\;x\in[0,L],
\end{equation}
jest klasycznym rozwiązaniem równania ciepła dla $t>0$,
ciągłym na $[0,\infty)\times[0,L]$ i przyjmuje dane $g$
w sposób jednostajny, gdy $t\downarrow0$. Rozwiązanie to
jest jedyne w klasie funkcji ciągłych na każdym
$[0,T]\times[0,L]$, klasy $C^{1,2}$ we wnętrzu
i spełniających dane i warunki Dirichleta.
\end{twierdzenie}
\begin{proof}
Z ograniczoności $g$ mamy $|b_n|\le2\|g\|_\infty$.
Dla $t\ge\delta>0$ szereg i wszystkie szeregi
z pochodnymi czasowymi i przestrzennymi zbieżne są
jednostajnie, bo dla dowolnych $r,s\ge0$
$\sum_{n\ge1}n^{r+2s}e^{-\kappa(n\pi/L)^2\delta}<\infty$.
Równanie i warunki brzegowe sprawdzamy wyraz po wyrazie.
Zagadnienie zbieżności w chwili $t=0$ wymaga osobnego kroku.
Dla wielomianu sinusowego $p$ rozwiązanie
$P_tp=\sum p_ne^{-\kappa\lambda_nt}\sin(n\pi x/L)$
zbiega jednostajnie do $p$. Każdą funkcję ciągłą na
$[0,L]$, zerującą się na końcach, można jednostajnie
aproksymować wielomianami sinusowymi: wystarczy
przedłużyć ją nieparzyście na $[-L,L]$, a potem
$2L$-okresowo, i zastosować
\ref{thm:pde6-fejer}.
Z lematu~\ref{lem:pde6-dirichlet-heat} wynika,
że operator splotu Gaussa po nieparzystym przedłużeniu
jest kontrakcją w normie supremum i ma dokładnie
te same współczynniki sinusowe co szereg w tezie.
Dla $t>0$ oba wyrażenia są funkcjami ciągłymi,
a zupełność układu sinusowego
(twierdzenie~\ref{thm:pde6-spectrum}) zapewnia ich równość.
W ten sposób kontrakcja została wykazana bez zakładania
z góry istnienia abstrakcyjnej półgrupy Dirichleta.
Stąd $\|P_tg-g\|_\infty\le
2\|g-p\|_\infty+\|P_tp-p\|_\infty\to 2\|g-p\|_\infty$,
a następnie wybieramy dowolnie dokładny $p$.
Jednoznaczność wynika z parabolicznej zasady maksimum
z rozdziału~\ref{chap:pde-heat}, zastosowanej do różnicy.
\end{proof}

\begin{wniosek}[Ostry wykładnik zaniku w normie $L^2$]
\label{cor:pde6-decay}
Jeżeli $g\in C([0,L])$ i $g(0)=g(L)=0$, to dla
rozwiązania~\eqref{eq:pde6-heat}
\begin{equation}\label{eq:pde6-decay}
 \|u(t,\cdot)\|_{L^2(0,L)}
 \le e^{-\kappa\pi^2t/L^2}\|g\|_{L^2(0,L)},\qquad t\ge0.
\end{equation}
Stała w wykładniku jest optymalna.
\end{wniosek}
\begin{proof}
Dla $t>0$ stosujemy ortogonalność sinusów i Parsevala:
\[
\|u(t)\|_2^2=\frac L2\sum_{n\ge1}|b_n|^2e^{-2\kappa\lambda_nt}
\le e^{-2\kappa\lambda_1t}\|g\|_2^2.
\]
Dla $t=0$ jest to równość. Przy $g(x)=\sin(\pi x/L)$
zachodzi równość dla każdego $t$.
\end{proof}


\begin{stwierdzenie}[Rozwiązanie cieplne dla danych $L^2$]
\label{prop:pde6-L2-semigroup}
Niech $\kappa,L>0$, $g\in L^2(0,L)$ i niech
$g_n=\langle g,e_n\rangle_{L^2(0,L)}$ dla ortonormalnych
funkcji $e_n$ z twierdzenia~\ref{thm:pde6-spectrum}.
Szereg $u(t)=\sum_{n\ge1}e^{-\kappa\lambda_nt}g_ne_n$
wyznacza funkcję $u\in C([0,\infty);L^2(0,L))$,
przy czym $u(0)=g$, a dla każdego $t>0$ otrzymujemy
funkcję gładką względem $x\in[0,L]$, klasyczną
w $t>0$ i spełniającą zerowe warunki Dirichleta.
Zachodzi także $\|u(t)\|_2\le e^{-\kappa\lambda_1t}\|g\|_2$.
\end{stwierdzenie}
\begin{proof}
Z twierdzenia Parsevala mamy $\sum_{n\ge1}|g_n|^2<\infty$.
Dla $t\ge0$ szereg zbiega w $L^2$, a
\[
 \|u(t)-g\|_2^2=\sum_{n\ge1}
 |e^{-\kappa\lambda_nt}-1|^2|g_n|^2\longrightarrow0
 \quad(t\downarrow0)
\]
na mocy zbieżności zdominowanej dla szeregów.
Ten sam argument w dowolnej chwili $t_0\ge0$
uzasadnia ciągłość $L^2$ w czasie.
Dla $t\ge\delta>0$ nierówność Cauchy'ego--Schwarza
zapewnia jednostajną zbieżność szeregów pochodnych:
\[
 \sum_{n\ge1}|g_n|n^k e^{-\kappa\lambda_n t}
 \le\Bigl(\sum_{n\ge1}|g_n|^2\Bigr)^{1/2}
 \Bigl(\sum_{n\ge1}n^{2k}e^{-2\kappa\lambda_n\delta}\Bigr)^{1/2}
 <\infty.
\]
Uwzględnienie dodatkowych potęg $\lambda_n$
pozwala również różniczkować po czasie. Stąd równanie
klasyczne i warunki brzegowe. Ostatnie oszacowanie
wynika wprost z $\lambda_n\ge\lambda_1$ i Parsevala.
\end{proof}

\section{Rozdzielenie zmiennych: równanie falowe i energia modów}
\label{sec:pde6-wave}

Rozważmy $u_{tt}=c^2u_{xx}$ na $(0,L)$,
z zerowymi warunkami Dirichleta oraz danymi
$u(0,x)=g(x)$, $u_t(0,x)=h(x)$.
Dla $X_n(x)=\sin(n\pi x/L)$ otrzymujemy równanie
$T_n''+c^2\lambda_nT_n=0$.

\begin{twierdzenie}[Struna: reprezentacja modalna]
\label{thm:pde6-wave}
Niech $c,L>0$ i $g,h\in C_c^\infty(0,L)$.
Niech $a_n=(2/L)\int_0^Lg(x)\sin(n\pi x/L)
\,\mathrm d\lambda_1(x)$ i
$b_n=(2/L)\int_0^Lh(x)\sin(n\pi x/L)\,\mathrm d\lambda_1(x)$.
Wówczas funkcja
\begin{equation}\label{eq:pde6-wave}
 u(t,x)=\sum_{n=1}^\infty
 \left[a_n\cos(c\sqrt{\lambda_n}t)
 +\frac{b_n}{c\sqrt{\lambda_n}}
 \sin(c\sqrt{\lambda_n}t)\right]\sin(n\pi x/L)
\end{equation}
jest jedynym rozwiązaniem klasy $C^2(\R\times[0,L])$
powyższego zagadnienia wśród rozwiązań klasycznych.
Jego energia falowa
$E(t)=\frac12\int_0^L(u_t(t,x)^2+c^2u_x(t,x)^2)
\,\mathrm d\lambda_1(x)$ jest stała.
\end{twierdzenie}
\begin{proof}
Nieparzyste $2L$-okresowe przedłużenia $g,h$
są klasy $C^\infty$, ponieważ obie funkcje
zerują się w otoczeniach końców przedziału.
Całkowanie przez części dowolnie wiele razy
pokazuje, że $|a_n|+|b_n|\le C_mn^{-m}$
dla każdego $m\ge0$. Na dowolnym zwartym przedziale
czasowym można więc różniczkować szereg
\eqref{eq:pde6-wave} dowolnie wiele razy;
uzyskujemy równanie, warunki początkowe i brzegowe.
Jednoznaczność wynika z zachowania energii
na odcinku (rozdział~\ref{chap:pde-wave}): różnica
rozwiązań ma energię zerową, więc $w_t=w_x=0$;
z zerowych warunków brzegowych wynika $w=0$.
Zachowanie energii można też otrzymać, sumując
energię modów za pomocą ortogonalności
$(\sin(n\pi x/L))_{n\ge1}$ i
$(\cos(n\pi x/L))_{n\ge1}$.
\end{proof}

\begin{przyklad}[Wygładzanie a brak wygładzania]
\label{ex:pde6-smoothing}
Dla równania ciepła amplituda modu $n$ zanika jak
$e^{-\kappa(n\pi/L)^2t}$. Dla równania falowego
mod ma amplitudę oscylującą i nie znika
z upływem czasu. Na przykład dla
$g(x)=\sin(n\pi x/L)$, $h=0$ otrzymujemy odpowiednio
$e^{-\kappa\lambda_nt}g(x)$ i
$\cos(c\sqrt{\lambda_n}t)g(x)$.
Ten kontrast jest spektralnym wyjaśnieniem
wygładzania parabolicznego.
\end{przyklad}

\section{Niejednorodne równania i wzór modalny Duhamela}
\label{sec:pde6-forcing}

Rozważmy $u_t-\kappa u_{xx}=f(t,x)$ na $(0,L)$
z jednorodnymi warunkami Dirichleta oraz danymi $g$.
Dla $f\in C^\infty([0,T]\times[0,L])$ z nośnikiem
przestrzennym zawartym w pewnym stałym zwartym
podzbiorze $(0,L)$, a także $g\in C_c^\infty(0,L)$,
możemy rozwijać wszystko w szeregi sinusowe.
Niech
$f_n(t)=(2/L)\int_0^Lf(t,x)\sin(n\pi x/L)
\,\mathrm d\lambda_1(x)$.
Mnożąc równanie przez odpowiedni mod i całkując
po $x$ otrzymujemy
$b_n'(t)+\kappa\lambda_nb_n(t)=f_n(t)$,
$\ b_n(0)=g_n$.
Po zastosowaniu czynnika całkującego dostajemy
\begin{equation}\label{eq:pde6-duhamel}
 b_n(t)=e^{-\kappa\lambda_nt}g_n
 +\int_0^t e^{-\kappa\lambda_n(t-s)}f_n(s)
 \,\mathrm d\lambda_1(s).
\end{equation}
W konsekwencji $u(t,x)=\sum_{n\ge1}b_n(t)
\sin(n\pi x/L)$. W podanych założeniach wszystkie
pochodne potrzebne do równania można sumować
jednostajnie na $[0,T]\times[0,L]$:
wielokrotne całkowanie przez części daje
$\sup_{t\in[0,T]}|f_n^{(j)}(t)|\le C_{m,j}n^{-m}$
i analogicznie dla $g_n$. Zatem wzór nie jest
jedynie rachunkiem formalnym.

\section{Regularne zagadnienie Sturma--Liouville'a}
\label{sec:pde6-sturm}

Przy rozdzieleniu zmiennych w równaniach ze
zmiennymi współczynnikami otrzymujemy często operator
$\mathcal L y=-(py')'+qy$. Zdefiniujmy dokładnie
zagadnienie i przestrzeń, w której badamy funkcje własne.

\begin{definicja}[Regularne zagadnienie Dirichleta]
\label{def:pde6-sturm}
Niech $L>0$, $p\in C^1([0,L];\R)$,
$q,w\in C([0,L];\R)$, $p(x)>0$, $w(x)>0$
i $q(x)\ge0$ na $[0,L]$. Zagadnieniem
Sturma--Liouville'a z warunkami Dirichleta nazywamy
\begin{equation}\label{eq:pde6-sturm}
 -(p(x)y'(x))'+q(x)y(x)=\lambda w(x)y(x),
 \quad y(0)=y(L)=0.
\end{equation}
Rozwiązaniem jest $y\in C^2([0,L])$ spełniające
równanie punktowo. Pracujemy w przestrzeni
$H=L^2((0,L),w(x)\,\mathrm d\lambda_1(x))$ z
iloczynem skalarnym $\langle f,g\rangle_H
=\int_0^Lf(x)\overline{g(x)}w(x)\,\mathrm d\lambda_1(x)$.
\end{definicja}

\begin{lemat}[Tożsamość Lagrange'a i rzeczywistość widma]
\label{lem:pde6-lagrange}
Niech $y,z\in C^2([0,L])$. Wtedy
\begin{equation}\label{eq:pde6-lagrange}
 \int_0^L\bigl((\mathcal Ly)(x)\overline{z(x)}
 -y(x)\overline{(\mathcal Lz)(x)}\bigr)\,\mathrm d\lambda_1(x)
 =\left[p\left(y\overline{z'}-y'\overline z\right)\right]_0^L.
\end{equation}
Jeżeli $y\not\equiv0$ jest funkcją własną
z jednorodnymi warunkami Dirichleta, to $\lambda>0$
jest rzeczywista. Funkcje własne dla różnych
wartości własnych są ortogonalne w $H$.
\end{lemat}
\begin{proof}
Pierwszą równość otrzymujemy przez całkowanie
przez części. Dla $z=y$ i jednorodnych warunków
brzegowych mamy
$\lambda\int_0^Lw(x)|y(x)|^2\,\mathrm d\lambda_1(x)
=\int_0^L(p(x)|y'(x)|^2+q(x)|y(x)|^2)\,\mathrm d\lambda_1(x)>0$.
Dla pary funkcji własnych $y,z$ o wartościach
$\lambda,\mu$ stosujemy tożsamość Lagrange'a;
ze znikania brzegów i rzeczywistości liczb własnych
wynika $(\lambda-\mu)\langle y,z\rangle_H=0$.
\end{proof}

\begin{twierdzenie}[Dyskretność i zupełność widma]
\label{thm:pde6-sturm}
Przy założeniach definicji~\ref{def:pde6-sturm}
istnieje rosnący do nieskończoności ciąg
$0<\lambda_1<\lambda_2<\cdots$ i odpowiednie
rzeczywiste funkcje własne $y_n\in C^2([0,L])$,
które po unormowaniu tworzą bazę ortonormalną $H$.
Każda wartość własna jest prosta.
\end{twierdzenie}
\begin{proof}
Najpierw konstruujemy operator zwarty, do którego
zastosujemy \emph{twierdzenie spektralne dla zwartych
operatorów samosprzężonych w przestrzeni Hilberta}:
każdy taki operator ma bazę ortonormalną złożoną
z wektorów własnych i wektorów należących do jądra.
Twierdzenie to traktujemy jako jawnie wskazany
wynik analizy funkcjonalnej, a nie dowód ukryty
we własnościach szeregów Fouriera.

Równanie $\mathcal Ly=f$ z danymi Dirichleta ma
dokładnie jedno rozwiązanie dla $f\in C([0,L])$.
Rzeczywiście, jednorodność i $q\ge0$ gwarantują
trywialność jądra przez energię. Standardowa
teoria liniowego ODE z dwoma warunkami brzegowymi
daje wtedy istnienie: wybieramy dowolne rozwiązanie
szczególne z warunkiem $y(0)=0$ i korygujemy
je wielokrotnością rozwiązania jednorodnego
z $y(0)=0$, $y'(0)=1$. To ostatnie nie może
zerować się w $L$, bo mielibyśmy nietrywialne
rozwiązanie jednorodne z dwoma zerowymi brzegami.

Niech $G(x,s)$ oznacza funkcję Greena dla $\mathcal L$
z warunkami Dirichleta, skonstruowaną z dwóch
rozwiązań jednorodnych: pierwszego o $y_1(0)=0$,
$y_1'(0)=1$ i drugiego o $y_2(L)=0$, $y_2'(L)=-1$.
Z równości $[p(y_1'y_2-y_1y_2')]'=0$ wynika,
że $C=p(s)(y_1'(s)y_2(s)-y_1(s)y_2'(s))$
jest stałą różną od zera, oraz
\begin{equation}\label{eq:pde6-green}
 G(x,s)=\frac{y_1(\min\{x,s\})y_2(\max\{x,s\})}{C}.
\end{equation}
Funkcja $G$ jest ciągła i symetryczna na
$[0,L]^2$ oraz spełnia
$\mathcal L_xG(x,s)=\delta_s$ w sensie skoku
$-p(s)(\partial_xG(s^+,s)-\partial_xG(s^-,s))=1$.
Definiujemy $Tf(x)=\int_0^LG(x,s)w(s)f(s)
\,\mathrm d\lambda_1(s)$ dla $f\in H$.
Jądro $G$ jest ciągłe, a $w$ dodatnia i ciągła;
zatem $T$ jest zwartym operatorem na $H$:
aproksymując jednostajnie jądro
$G(x,s)$ sumami rozdzielonych funkcji $a_j(x)b_j(s)$,
uzyskujemy operatory skończonego rzędu
zbieżne do $T$ w normie operatorowej.
Takie przybliżenie można skonstruować jawnie:
wybieramy podział $0=x_0<\cdots<x_N=L$ i liniowe
funkcje kapeluszowe $\theta_j$ tworzące rozkład jedności.
Interpolant
$G_N(x,s)=\sum_{i,j=0}^N G(x_i,x_j)\theta_i(x)\theta_j(s)$
jest sumą skończenie wielu iloczynów funkcji jednej zmiennej.
Z jednostajnej ciągłości $G$ na $[0,L]^2$ wynika
$\|G_N-G\|_\infty\to0$, a z nierówności
Cauchy'ego--Schwarza i ograniczoności dodatniej wagi $w$
wynika $\|T_N-T\|_{H\to H}\le
L\,\|w\|_\infty\,\|G_N-G\|_\infty\to0$.
Zatem zwartość nie wymaga tu nieudowodnionego twierdzenia
o aproksymacji ciągłych jąder.
Symetria jądra daje samosprzężoność.
Dla $f\in H$ funkcja $y=Tf$ jest klasy $C^1([0,L])$,
a $py'$ jest absolutnie ciągła. Istotnie, rozpisując całkę
z jądrem~\eqref{eq:pde6-green} na przedziały $[0,x]$
i $[x,L]$, można ją zróżniczkować po $x$; składniki
pochodzące z ruchomego punktu podziału znoszą się.
Ponowne zróżniczkowanie w sensie prawie wszędzie
oraz warunek skoku jądra dają
$-(py')'+qy=wf$ prawie wszędzie. Nie twierdzimy
jeszcze, że $y$ jest klasy $C^2$ dla dowolnego $f\in H$.
Całkowanie przez części z użyciem absolutnej ciągłości
$py'$ zapewnia
$\langle Tf,f\rangle_H=
\int_0^L(p(x)|y'(x)|^2+q(x)|y(x)|^2)\,\mathrm d\lambda_1(x)\ge0$;
dla $f\ne0$ nierówność jest ostra. Operator $T$
ma więc trywialne jądro, a z samosprzężoności
wynika $\overline{\operatorname{ran}T}=H$.
Twierdzenie spektralne daje bazę funkcji własnych
$Ty_n=\mu_n y_n$ z $\mu_n>0$ i $\mu_n\downarrow0$.
Ponieważ $\mathcal L(Ty_n)=wy_n$ prawie wszędzie,
otrzymujemy $\mathcal Ly_n=\mu_n^{-1}wy_n$ w tym
samym sensie. Funkcja $y_n=\mu_n^{-1}Ty_n$ jest ciągła,
a więc prawa strona równania jest ciągła. Ponieważ
$p\in C^1$ jest dodatnia, równanie daje najpierw
$py_n'\in C^1$, a następnie $y_n\in C^2([0,L])$.
Równanie zachodzi wtedy punktowo. Nie może być skończenie
wielu niezerowych wartości własnych, bo $T$ jest
iniektywny na nieskończeniewymiarowym $H$.
Wreszcie prostota wynika z tego, że dla
ustalonej $\lambda$ dwa rozwiązania warunku
$y(0)=0$ są proporcjonalne (są wyznaczone
przez wartość $y'(0)$).
\end{proof}

\begin{uwaga}[Co jest tu nowe?]
Dla stałych współczynników mieliśmy jawne funkcje
$\sin(n\pi x/L)$; przy zmiennych współczynnikach
możemy zwykle nie znać ich w postaci elementarnej,
ale nadal mamy dyskretne wartości własne oraz
zupełny układ modów. To właśnie uzasadnia metodę
rozdzielenia zmiennych również bez jawnych wzorów.
Twierdzenie o oscylacji Sturma, mówiące o liczbie
zer kolejnych funkcji własnych, omówimy osobno,
gdy będziemy dysponować szerszą teorią ODE i
zagadnień spektralnych.
\end{uwaga}

\section{Wyższe wymiary: prostokąt i geometria widma}
\label{sec:pde6-rectangle}

Dla prostokąta $\Omega=(0,L_1)\times(0,L_2)$
problem Dirichleta $-\Delta X=\lambda X$ rozdziela się
na dwa jednowymiarowe zagadnienia własne.
Otrzymujemy
\begin{equation}\label{eq:pde6-rectangle}
 X_{mn}(x,y)=\sin\left(\frac{m\pi x}{L_1}\right)
 \sin\left(\frac{n\pi y}{L_2}\right),\quad
 \lambda_{mn}=\pi^2\left(\frac{m^2}{L_1^2}
 +\frac{n^2}{L_2^2}\right),\quad m,n\ge1.
\end{equation}
Rodzina $\{2(L_1L_2)^{-1/2}X_{mn}\}_{m,n\ge1}$
jest bazą ortonormalną $L^2(\Omega)$. Rzeczywiście,
ortogonalność wynika z rozdzielenia całek, a
zupełność z zastosowania jednowymiarowego rozwinięcia
w zmiennej $x$, a potem $y$ (twierdzenie Fubiniego
i zbieżność w $L^2$ pozwalają na ten krok).
Niekoniecznie wszystkie wartości $\lambda_{mn}$ są
różne: symetria kwadratu $L_1=L_2$ powoduje
na przykład $\lambda_{12}=\lambda_{21}$.
Wartości własne mogą mieć krotność większą niż jeden,
mimo że na przedziale każda jest prosta.

\section{Dlaczego nie można bezwarunkowo odwrócić ewolucji cieplnej?}
\label{sec:pde6-backward}

Niech $\kappa,L,T>0$. Rozważmy operator
$P_T:g\mapsto u(T,\cdot)$ dla rozwiązania równania
ciepła z jednorodnymi warunkami Dirichleta.
Na modzie $e_n$ działa on przez mnożnik
$e^{-\kappa\lambda_nT}$. Stąd
$\|P_Te_n\|_2=e^{-\kappa\lambda_nT}\to0$,
podczas gdy $\|e_n\|_2=1$.
Jeżeli odwrotność $P_T^{-1}$ na jego obrazie byłaby
ograniczona w normie $L^2$, istniałaby stała $C>0$
z $1=\|e_n\|_2\le C\|P_Te_n\|_2$, sprzeczność.
Nie ma więc ciągłego operatora odtwarzającego
dowolne dane początkowe z dokładnych danych końcowych
w normie $L^2$. To spektralna wersja
niestabilności Hadamarda z rozdziału~\ref{chap:pde-heat}.

\section{Podsumowanie i dalsza perspektywa}
\label{sec:pde6-summary}

Dzięki średnim Fejéra uzyskaliśmy zupełność układu
trygonometrycznego i twierdzenie Parsevala.
Na ich podstawie skonstruowaliśmy jawne rozwiązania
ciepła i fal na odcinku oraz ostre oszacowanie zaniku.
Metoda ta przechodzi na zmienne współczynniki przez
regularną teorię Sturma--Liouville'a i na prostokąty
przez iloczyny tensorowe. W kolejnych rozdziałach
rozszerzymy ją na rozwiązania słabe, operatory
samosprzężone, metody wariacyjne i transformatę
Fouriera na $\R^n$.


\chapter[Dystrybucje i słabe pochodne]{Dystrybucje, słabe pochodne i rozwiązania fundamentalne}
\label{chap:pde-distributions}

W poprzednich rozdziałach rozwiązanie klasyczne musiało mieć tyle pochodnych,
ile wymagało równanie. W praktyce taki warunek jest często zbyt silny:
profil fali może mieć skok, dane początkowe mogą być tylko całkowalne,
a rozwiązanie eliptyczne może być najpierw otrzymane przez minimalizację
energii. Potrzebujemy rachunku różniczkowego, w którym można różniczkować
nawet funkcję skokową. Nie chodzi o arbitralne przypisywanie jej wartości
pochodnej w punkcie skoku, lecz o zastąpienie równań punktowych przez
tożsamości całkowe z funkcjami testowymi.

\section{Przestrzeń funkcji testowych}
\label{sec:pde7-tests}

\begin{definicja}[Funkcje testowe i ich zbieżność]
\label{def:pde7-test}
Niech $d\ge1$ oraz $\Omega\subset\R^d$ będzie niepustym zbiorem otwartym.
Przez $\mathcal D(\Omega)=C_c^\infty(\Omega;\C)$ oznaczamy przestrzeń
zespolonych funkcji nieskończenie różniczkowalnych o zwartym nośniku
zawartym w $\Omega$. Dla $K\Subset\Omega$ definiujemy
$\mathcal D_K(\Omega)=\{\varphi\in\mathcal D(\Omega):\supp\varphi\subset K\}$
oraz półnormy
$p_{K,m}(\varphi)=\max_{|\alpha|\le m}\sup_{x\in K}|\partial^\alpha\varphi(x)|$,
$m\in\N_0$. Mówimy, że $\varphi_j\to0$ w $\mathcal D(\Omega)$,
jeżeli istnieje jeden zbiór zwarty $K\Subset\Omega$ obejmujący nośniki
wszystkich $\varphi_j$ od pewnego indeksu oraz
$p_{K,m}(\varphi_j)\to0$ dla każdego $m\in\N_0$.
Topologia $\mathcal D(\Omega)$ jest standardową topologią granicy
induktywnej przestrzeni $\mathcal D_K(\Omega)$; powyższy opis
określa w szczególności jej zbieżne ciągi.
\end{definicja}

Notacja $K\Subset\Omega$ oznacza, że $K$ jest zwarty i $K\subset\Omega$.
Warto odróżnić tę sytuację od żądania $\varphi\in C^\infty(\overline\Omega)$:
funkcja testowa znika na pewnym otoczeniu brzegu $\Omega$.
Istnieje wiele funkcji testowych. Dla przykładu funkcja
\[
 \eta(x)=\begin{cases}
 \exp\!\bigl(-1/(1-|x|^2)\bigr),& |x|<1,\\
 0,& |x|\ge1
 \end{cases}
\]
jest klasy $C_c^\infty(\R^d)$ i jest dodatnia w $B_1(0)$.
Podczas weryfikacji gładkości na sferze wykorzystujemy fakt, że
$e^{-1/s}$ wraz ze wszystkimi pochodnymi dąży do zera przy $s\downarrow0$.

\begin{uwaga}[Dlaczego potrzebny jest wspólny nośnik?]
Przesunięcia $\varphi_j(x)=\eta(x-j)$ na $\R$ mają wszystkie pochodne
jednostajnie ograniczone, a dla każdego ustalonego $x$ od pewnego miejsca
znikają. Mimo to nie dążą do zera w $\mathcal D(\R)$, ponieważ ich
nośniki uciekają do nieskończoności. To rozróżnienie okaże się niezbędne
dla definicji dystrybucji.
\end{uwaga}

\section{Dystrybucje jako ciągłe funkcjonały liniowe}
\label{sec:pde7-def}

\begin{definicja}[Dystrybucja]
\label{def:pde7-distribution}
Dla niepustego otwartego $\Omega\subset\R^d$ dystrybucją na $\Omega$
nazywamy ciągły funkcjonał \emph{liniowy} $T:\mathcal D(\Omega)\to\C$.
Zbiór wszystkich takich funkcjonałów oznaczamy przez $\mathcal D'(\Omega)$.
Wartość funkcjonału zapisujemy jako $\langle T,\varphi\rangle$;
nie jest to na ogół iloczyn skalarny, a w szczególności
zapis jest liniowy względem $\varphi$, bez sprzężenia zespolonego.
\end{definicja}

Związek z definicją topologiczną można wyrazić za pomocą lokalnych
oszacowań. Jest to pierwszy ważny przykład tego, że element
przestrzeni dualnej do przestrzeni funkcji można opisywać
przez oszacowania na jej półnormach.

\begin{twierdzenie}[Lokalne oszacowanie dystrybucji]
\label{thm:pde7-local-order}
Niech $\Omega\subset\R^d$ będzie otwarta, a $T$ liniowym funkcjonałem
na $\mathcal D(\Omega)$. Wówczas $T\in\mathcal D'(\Omega)$ wtedy i tylko
wtedy, gdy dla każdego $K\Subset\Omega$ istnieją $m_K\in\N_0$
oraz $C_K>0$ takie, że
\begin{equation}\label{eq:pde7-estimate}
 |\langle T,\varphi\rangle|\le C_K
 \max_{|\alpha|\le m_K}\sup_{x\in K}|\partial^\alpha\varphi(x)|,
 \qquad \varphi\in\mathcal D_K(\Omega).
\end{equation}
\end{twierdzenie}
\begin{proof}
Jeżeli $T$ jest ciągły, to jego ograniczenie do przestrzeni
$\mathcal D_K(\Omega)$ jest ciągłym funkcjonałem liniowym na przestrzeni
Frécheta opisanej przez rosnącą rodzinę półnorm $p_{K,m}$.
Ciągłość w zerze oznacza, że dla pewnych $m$ i $\delta>0$ mamy
$|T(\varphi)|<1$, gdy $p_{K,m}(\varphi)<\delta$. Skalując
$\varphi$ przez dodatnią liczbę (a przypadek $p_{K,m}(\varphi)=0$
rozpatrując oddzielnie), otrzymujemy
$|T(\varphi)|\le (2/\delta)p_{K,m}(\varphi)$.
Odwrotnie, jeśli zachodzi \eqref{eq:pde7-estimate}, to $T$ jest
ciągły na każdym $\mathcal D_K(\Omega)$, a więc ciągły dla
opisanej topologii granicy induktywnej.
\end{proof}

Liczba $m_K$ może zależeć od $K$. Mówimy wtedy o \emph{lokalnie
skończonym rzędzie} dystrybucji. Nie należy bez dodatkowych założeń
wnioskować, że jedna liczba $m$ działa na wszystkich zwartych
podzbiorach obszaru, \emph{nawet gdy jest on ograniczony}.
Istotnie, na $\Omega=(0,1)$ połóżmy
$T=\sum_{k=2}^\infty 2^{-k^2}\partial^k\delta_{1/k}$.
Każda funkcja testowa ma nośnik oddzielony od zera, więc
na jej nośniku występuje jedynie skończenie wiele
składników. Suma określa dystrybucję w $\mathcal D'(0,1)$.
Nie istnieje jednak wspólny skończony rząd $m$:
w otoczeniu dowolnego punktu $1/k$ nieprzecinającym
pozostałych punktów nośnika dystrybucji mamy
$T=2^{-k^2}\partial^k\delta_{1/k}$, którego rząd wynosi $k$.
To pokazuje, że ograniczoność obszaru nie wystarcza
do globalnego oszacowania rzędu dystrybucji.

\begin{przyklad}[Dystrybucja regularna i miara punktowa]
\label{ex:pde7-regular}
Jeśli $f\in L^1_{\mathrm{loc}}(\Omega)$, to wzór
\begin{equation}\label{eq:pde7-regular}
 \langle T_f,\varphi\rangle
 =\int_\Omega f(x)\varphi(x)\,\mathrm d\lambda_d(x)
\end{equation}
definiuje dystrybucję, ponieważ dla $\supp\varphi\subset K$ zachodzi
$|T_f(\varphi)|\le\|\varphi\|_\infty\int_K|f(x)|\,\mathrm d\lambda_d(x)$.
Dla $a\in\Omega$ wzór $\langle\delta_a,\varphi\rangle=\varphi(a)$
definiuje \emph{deltę Diraca}, również dystrybucję rzędu zerowego
na każdym zwartym podzbiorze $\Omega$.
\end{przyklad}

\begin{stwierdzenie}[Iniektywność odwzorowania $f\mapsto T_f$]
\label{prop:pde7-injective}
Jeżeli $f,g\in L^1_{\mathrm{loc}}(\Omega)$ i $T_f=T_g$, to $f=g$
prawie wszędzie w $\Omega$ względem miary Lebesgue'a.
\end{stwierdzenie}
\begin{proof}
Wystarczy rozważyć $h=f-g$. Wybierzmy $a\in\Omega$
i $r>0$ tak, aby $B_{2r}(a)\Subset\Omega$.
Ustalmy nieujemną funkcję $\rho\in C_c^\infty(B_1(0))$
o całce $1$ i połóżmy
$\rho_\eps(x)=\eps^{-d}\rho(x/\eps)$ dla $0<\eps<r$.
Dla każdego $x\in B_r(a)$ funkcja
$y\mapsto\rho_\eps(x-y)$ jest testowa i ma nośnik
w $B_{2r}(a)$. Ponieważ $T_h=0$, mamy
$(h*\rho_\eps)(x)=0$ dla wszystkich takich $x$.
Aby przejść do granicy, przedłużmy $h$ po obcięciu do
$B_{2r}(a)$ przez zero i oznaczmy otrzymaną funkcję
przez $H\in L^1(\R^d)$. Z nierówności Minkowskiego
wynika
\[
 \|H*\rho_\eps-H\|_{L^1(\R^d)}
 \le\int_{B_1(0)}\rho(z)
 \|H(\cdot-\eps z)-H\|_{L^1(\R^d)}
 \,\mathrm d\lambda_d(z)\longrightarrow0.
\]
Istotnie, przesunięcia są ciągłe w $L^1$: najpierw
sprawdza się to dla funkcji ciągłych o zwartym nośniku,
a następnie przez ich gęstość w $L^1$; pod całką
stosuje się ograniczenie $2\|H\|_1$.
Na $B_r(a)$ sploty $H*\rho_\eps$ są zerowe, stąd
$h=0$ prawie wszędzie na tej kuli. Pokrywając $\Omega$
przeliczalnie takimi kulami, otrzymujemy tezę.
Ogólną teorię przybliżonych jedynek rozwiniemy później
w sekcji~\ref{sec:pde7-mollifiers}; nie jest ona
przesłanką niniejszego dowodu.
\end{proof}

\begin{uwaga}[Dlaczego delta nie jest zwykłą funkcją?]
Załóżmy, że $\delta_0=T_f$ dla pewnego
$f\in L^1_{\mathrm{loc}}(\R^d)$. Wybierzmy funkcję
$\chi\in\mathcal D(\R^d)$, $0\le\chi\le1$,
$\chi(0)=1$, $\supp\chi\subset B_1(0)$, i połóżmy
$\chi_\varepsilon(x)=\chi(x/\varepsilon)$.
Wtedy $\langle\delta_0,\chi_\varepsilon\rangle=1$,
natomiast
$|\langle T_f,\chi_\varepsilon\rangle|
\le\int_{B_\varepsilon}|f(x)|\,\mathrm d\lambda_d(x)\to0$
przez absolutną ciągłość całki Lebesgue'a. Sprzeczność.
\end{uwaga}

\section{Pochodne dystrybucyjne i skoki}
\label{sec:pde7-derivatives}

\begin{definicja}[Pochodna dystrybucyjna]
\label{def:pde7-derivative}
Dla $T\in\mathcal D'(\Omega)$ i wielowskaźnika
$\alpha\in\N_0^d$ określamy $\partial^\alpha T\in\mathcal D'(\Omega)$
wzorem
\begin{equation}\label{eq:pde7-distributional-derivative}
 \langle\partial^\alpha T,\varphi\rangle
 =(-1)^{|\alpha|}\langle T,\partial^\alpha\varphi\rangle,
 \qquad\varphi\in\mathcal D(\Omega).
\end{equation}
\end{definicja}

Definicja ma sens, gdyż różniczkowanie funkcji testowej jest
ciągłym operatorem $\mathcal D(\Omega)\to\mathcal D(\Omega)$.
Dzięki temu \emph{każda} dystrybucja ma pochodne dowolnego rzędu,
a mieszane pochodne komutują.

\begin{stwierdzenie}[Zgodność z klasycznym różniczkowaniem]
\label{prop:pde7-classical}
Jeśli $u\in C^1(\Omega)$, to $\partial_jT_u=T_{\partial_j u}$.
Jeśli $T\in\mathcal D'(\Omega)$ i $a\in C^\infty(\Omega)$,
to iloczyn $aT$, dany wzorem
$\langle aT,\varphi\rangle=\langle T,a\varphi\rangle$,
jest dystrybucją i zachodzi reguła Leibniza
$\partial_j(aT)=(\partial_ja)T+a\partial_jT$.
\end{stwierdzenie}
\begin{proof}
Dla funkcji $C^1$ wzór wynika z całkowania przez części,
bez składnika brzegowego, bo $\varphi$ ma zwarty nośnik w $\Omega$.
Dla iloczynu i testu $\varphi$ liczymy
$\langle\partial_j(aT),\varphi\rangle
=-\langle T,\partial_j(a\varphi)\rangle
=-\langle T,(\partial_ja)\varphi+a\partial_j\varphi\rangle$,
co po przegrupowaniu jest żądaną równością.
\end{proof}

\begin{przyklad}[Funkcja Heaviside'a]
\label{ex:pde7-heaviside}
Niech $H(x)=0$ dla $x<0$ i $H(x)=1$ dla $x>0$;
wartość w zerze nie ma znaczenia. Dla $\varphi\in\mathcal D(\R)$
\[
 \langle H',\varphi\rangle
 =-\int_0^\infty\varphi'(x)\,\mathrm d\lambda_1(x)
 =\varphi(0)=\langle\delta_0,\varphi\rangle.
\]
Zatem $H'=\delta_0$ w $\mathcal D'(\R)$.
Ponadto dla $u(x)=|x|$ zachodzi
$u'=\operatorname{sgn}(x)$ i $u''=2\delta_0$.
\end{przyklad}

\begin{przyklad}[Pochodna delty]
Dystrybucja $\delta'_0$ nie jest miarą, a jej działanie opisuje
$\langle\delta'_0,\varphi\rangle=-\varphi'(0)$.
Ogólniej $\langle\partial^\alpha\delta_a,\varphi\rangle
=(-1)^{|\alpha|}\partial^\alpha\varphi(a)$.
Jest to dystrybucja rzędu co najwyżej $|\alpha|$ na zwartych
podzbiorach zawierających punkt $a$.
\end{przyklad}

\begin{twierdzenie}[Wzór skoku przez hiperpowierzchnię]
\label{thm:pde7-jump}
Niech $\Omega\subset\R^d$ będzie otwarta, a
$S\subset\Omega$ zorientowaną hiperpowierzchnią klasy $C^1$,
która lokalnie rozdziela $\Omega$ na strony $\Omega_-$ i $\Omega_+$.
Niech $\nu$ będzie jednostkowym wektorem normalnym skierowanym
od $\Omega_-$ do $\Omega_+$. Załóżmy, że funkcja $u$ jest klasy
$C^1$ po obu stronach $S$, ma jednostronne ślady $u_\pm$ klasy $C^1$
aż do $S$ i jest lokalnie całkowalna. Wtedy, w sensie dystrybucji,
\begin{equation}\label{eq:pde7-jump}
 \partial_jT_u=T_{(\partial_ju)_{\mathrm{reg}}}
 +(u_+-u_-)\nu_j\delta_S,
\end{equation}
gdzie $(\partial_ju)_{\mathrm{reg}}$ oznacza klasyczną pochodną
poza $S$, natomiast
$\langle b\delta_S,\varphi\rangle
=\int_S b(s)\varphi(s)\,\mathrm d\sigma(s)$
dla ciągłej funkcji $b$ na $S$ (całka jest lokalna, gdy $S$ jest
niezwarta).
\end{twierdzenie}
\begin{proof}
Ustalmy $\varphi\in\mathcal D(\Omega)$ i wybierzmy otoczenie
jej nośnika, w którym można zastosować wzór Gaussa po obu stronach
$S$ (w ogólnym przypadku używamy skończonego podziału jedności).
Całkując $-u\partial_j\varphi$ po $\Omega_+$ i $\Omega_-$,
otrzymujemy całki objętościowe z klasycznych pochodnych i
składniki brzegowe $-u_+\varphi(-\nu_j)$ oraz
$-u_-\varphi\nu_j$. Ich suma to
$(u_+-u_-)\nu_j\varphi$, a wkłady z pozostałych granic
znikają dzięki nośnikowi $\varphi$. Daje to
\eqref{eq:pde7-jump}.
\end{proof}

\begin{wniosek}[Skok strumienia i równanie ciągłości]
Niech $F\in L^1_{\mathrm{loc}}(\Omega;\R^d)$ będzie polem
klasy $C^1$ po stronach hiperpowierzchni $S$ jak wyżej.
Wówczas
$\operatorname{div}T_F=T_{(\operatorname{div}F)_{\mathrm{reg}}}
+[F]\cdot\nu\,\delta_S$, gdzie $[F]=F_+-F_-$.
Stąd równanie dywergencji $\operatorname{div}F=0$
w sensie dystrybucyjnym wymaga zanikania skoku normalnego
strumienia. To jest ogólna postać mechanizmu stojącego za warunkiem
Rankine'a--Hugoniota z rozdziału~\ref{chap:pde-characteristics}.
\end{wniosek}

\section{Nośnik, lokalność i zbieżność dystrybucji}
\label{sec:pde7-support}

\begin{definicja}[Nośnik dystrybucji]
Dla $T\in\mathcal D'(\Omega)$ nośnik $\supp T$ jest dopełnieniem
największego otwartego zbioru $U\subset\Omega$ takiego, że
$\langle T,\varphi\rangle=0$ dla wszystkich
$\varphi\in\mathcal D(U)$.
Mówimy, że ciąg $T_n\in\mathcal D'(\Omega)$ zbiega do
$T\in\mathcal D'(\Omega)$ w sensie dystrybucji, jeśli
$\langle T_n,\varphi\rangle\to\langle T,\varphi\rangle$
dla każdego $\varphi\in\mathcal D(\Omega)$.
\end{definicja}

Dla funkcji lokalnie całkowalnej nośnik $T_f$ jest domknięciem
zbioru, na którym $f$ nie znika w sensie lokalnym prawie wszędzie.
Mamy $\supp(\partial^\alpha T)\subset\supp T$,
ponieważ test i wszystkie jego pochodne są zerowe
poza nośnikiem testu.

\begin{stwierdzenie}[Różniczkowanie a granica dystrybucyjna]
\label{prop:pde7-derivative-limit}
Niech $\Omega\subset\R^d$ będzie otwarta i
$T_n\to T$ w $\mathcal D'(\Omega)$. Dla każdego
wielowskaźnika $\alpha$ zachodzi
$\partial^\alpha T_n\to\partial^\alpha T$ w $\mathcal D'(\Omega)$.
\end{stwierdzenie}
\begin{proof}
Dla każdego $\varphi\in\mathcal D(\Omega)$ mamy
$\langle\partial^\alpha T_n,\varphi\rangle
=(-1)^{|\alpha|}\langle T_n,\partial^\alpha\varphi\rangle$.
Ponieważ $\partial^\alpha\varphi$ jest funkcją testową,
przechodzimy do granicy po prawej stronie.
\end{proof}

\begin{przyklad}[Gładkie funkcje zbieżne do delty]
Jeśli $\rho\in\mathcal D(\R^d)$ spełnia $\rho\ge0$ i
$\int_{\R^d}\rho(x)\,\mathrm d\lambda_d(x)=1$, a
$\rho_\varepsilon(x)=\varepsilon^{-d}\rho(x/\varepsilon)$,
to $T_{\rho_\varepsilon}\to\delta_0$ w $\mathcal D'(\R^d)$.
Rzeczywiście,
$\langle T_{\rho_\varepsilon},\varphi\rangle
=\int_{\R^d}\rho(y)\varphi(\varepsilon y)\,\mathrm d\lambda_d(y)
\to\varphi(0)$ przez zbieżność jednostajną na nośniku $\rho$.
\end{przyklad}

\section{Wygładzanie i przybliżone jedynki}
\label{sec:pde7-mollifiers}

Wybierzmy odtąd $\rho\in C_c^\infty(B_1(0))$,
$\rho\ge0$, o całce równej $1$, i definiujmy
$\rho_\varepsilon(x)=\varepsilon^{-d}\rho(x/\varepsilon)$.
Dla $u\in L^1_{\mathrm{loc}}(\Omega)$ i punktów
$x\in\Omega_\varepsilon
=\{x\in\Omega:\operatorname{dist}(x,\R^d\setminus\Omega)>\varepsilon\}$
wprowadzamy
$u_\varepsilon(x)=\int_\Omega u(y)\rho_\varepsilon(x-y)
\,\mathrm d\lambda_d(y)$. Dla dystrybucji stosujemy analogiczną
definicję $T_\varepsilon(x)=\langle T(y),\rho_\varepsilon(x-y)\rangle$.
Zapis $T(y)$ w nawiasie jest tylko sposobem wskazania zmiennej
funkcji testowej; nie oznacza wartości punktowej dystrybucji.
Splot jest tu \emph{lokalny}: dla danego $x$ musimy mieć
$\operatorname{supp}(\rho_\varepsilon(x-\cdot))\Subset\Omega$.
Nie wolno automatycznie rozciągać takiej definicji na całą
$\R^d$ bez przedłużenia $T$ lub założeń o nośniku.

\begin{twierdzenie}[Wygładzanie dystrybucji]
\label{thm:pde7-mollify}
Niech $\Omega\subset\R^d$ będzie otwarta, $T\in\mathcal D'(\Omega)$
i $\rho_\varepsilon$ jak wyżej. Wówczas
$T_\varepsilon\in C^\infty(\Omega_\varepsilon)$ oraz
\begin{equation}\label{eq:pde7-deriv-mollify}
 \partial^\alpha T_\varepsilon
 =(\partial^\alpha T)*\rho_\varepsilon
 =T*(\partial^\alpha\rho_\varepsilon)
 \quad\text{na }\Omega_\varepsilon.
\end{equation}
Ponadto $T_\varepsilon\to T$ w $\mathcal D'(\Omega)$,
rozumiejąc każde $T_\varepsilon$ na zwartych podzbiorach,
na których jest określone, dla dostatecznie małych $\varepsilon$.
\end{twierdzenie}
\begin{proof}
Dla ustalonego $x\in\Omega_\varepsilon$ przesunięta funkcja
$y\mapsto\rho_\varepsilon(x-y)$ jest testowa i jej nośnik
leży w $\Omega$. Z ciągłości $T$ oraz gładkiej zależności
funkcji testowej od $x$ wynikają gładkość $T_\varepsilon$
i wzór przez różniczkowanie po parametrze.
Tożsamość z $\partial^\alpha T$ wynika z definicji pochodnej
po scałkowaniu przez części w zmiennej testowej.
Dla $\varphi\in\mathcal D(\Omega)$ i małego $\varepsilon$
zamiana działania $T$ z całką daje
$\langle T_\varepsilon,\varphi\rangle
=\langle T,\check\rho_\varepsilon*\varphi\rangle$,
gdzie $\check\rho(x)=\rho(-x)$.
Na jednym zwartym podzbiorze $\Omega$ mamy
$\check\rho_\varepsilon*\varphi\to\varphi$
jednostajnie ze wszystkimi pochodnymi; wynika to
z jednostajnej ciągłości każdej pochodnej $\varphi$.
Z definicji ciągłości $T$ otrzymujemy zbieżność dystrybucyjną.
\end{proof}

\begin{twierdzenie}[Aproksymacja w $L^p_{\mathrm{loc}}$]
\label{thm:pde7-approx-id}
Niech $\Omega\subset\R^d$ będzie otwarta,
$1\le p<\infty$ i $u\in L^p_{\mathrm{loc}}(\Omega)$.
Wówczas $u_\varepsilon\to u$ w $L^p(K)$ dla każdego
$K\Subset\Omega$.
\end{twierdzenie}
\begin{proof}
Dla $K\Subset\Omega$ wybierzmy otwarty zbiór
$V$ o $K\subset V\Subset\Omega$ i funkcję odcinającą
$\chi\in C_c^\infty(\Omega)$ równą $1$ na $V$.
Funkcja $v=\chi u$ przedłużona zerem na $\R^d$
należy do $L^p(\R^d)$ i $u_\varepsilon=v*\rho_\varepsilon$
na $K$ dla małego $\varepsilon$. Z nierówności Minkowskiego
otrzymujemy
\[
 \|v*\rho_\varepsilon-v\|_{L^p(\R^d)}
 \le\int_{\R^d}\rho(z)
 \|v(\cdot-\varepsilon z)-v\|_{L^p(\R^d)}
 \,\mathrm d\lambda_d(z).
\]
Ciągłość przesunięć w $L^p(\R^d)$ daje zbieżność do zera
jednostajnie dla $z$ z ustalonego zwartego nośnika $\rho$.
Ciągłość przesunięć można wykazać najpierw dla
$C_c^\infty(\R^d)$, a następnie przez gęstość tych funkcji
w $L^p$ i izometryczność przesunięcia. To kończy dowód.
\end{proof}

\begin{uwaga}[Nie mylić sposobów zbieżności]
Dla $p=\infty$ na ogół nie zachodzi zbieżność w normie
$L^\infty$: wygładzenie skoku nie może aproksymować go
jednostajnie w otoczeniu skoku. Dla funkcji ciągłych
przybliżona jedynka daje jednak zbieżność jednostajną
na zwartych podzbiorach.
\end{uwaga}

\begin{center}
\begin{tikzpicture}[scale=.86,>=Stealth]
 \draw[->] (-4.3,0)--(4.7,0) node[below] {$x$};
 \draw[->] (0,-.1)--(0,3.5) node[left] {$\rho_\varepsilon(x)$};
 \draw[thick,blue!65!black,domain=-2.2:2.2,samples=80,smooth]
   plot (\x,{1.05*exp(-3*(\x)^2)});
 \draw[thick,red!70!black,domain=-1:1,samples=80,smooth]
   plot (\x,{2.7*exp(-13*(\x)^2)});
 \node[blue!65!black,anchor=west] at (2.05,1.05) {$\varepsilon$ większe};
 \node[red!70!black,anchor=west] at (.75,2.65) {$\varepsilon$ mniejsze};
 \node[below] at (0,-.1) {$0$};
\end{tikzpicture}
\par\smallskip\small
Poglądowe profile przybliżonej jedynki: masa pozostaje równa $1$,
choć wysokość rośnie, a nośnik koncentruje się wokół zera.
Rysunek ilustruje zjawisko, nie przedstawia znormalizowanej
funkcji o zadanym wzorze analitycznym.
\end{center}

\section{Słabe pochodne i pierwsze przestrzenie Sobolewa}
\label{sec:pde7-weak}

\begin{definicja}[Pochodna słaba]
\label{def:pde7-weak}
Niech $\Omega\subset\R^d$ będzie otwarta,
$u\in L^1_{\mathrm{loc}}(\Omega)$ i $\alpha\in\N_0^d$.
Funkcję $v\in L^1_{\mathrm{loc}}(\Omega)$ nazywamy
\emph{słabą pochodną} $\partial^\alpha u$, jeżeli
\begin{equation}\label{eq:pde7-weak}
 \int_\Omega u(x)\partial^\alpha\varphi(x)\,\mathrm d\lambda_d(x)
 =(-1)^{|\alpha|}\int_\Omega v(x)\varphi(x)\,
 \mathrm d\lambda_d(x)
 \quad (\varphi\in C_c^\infty(\Omega)).
\end{equation}
Słaba pochodna jest jednoznaczna prawie wszędzie,
bo $T_v=\partial^\alpha T_u$ i działa
stwierdzenie~\ref{prop:pde7-injective}.
\end{definicja}

Pochodna dystrybucyjna istnieje zawsze, ale słaba
pochodna w powyższym znaczeniu istnieje tylko wtedy,
gdy odpowiadająca jej dystrybucja jest regularna.
To rozróżnienie jest istotne: pochodna $H$ jest deltą,
która nie pochodzi od funkcji lokalnie całkowalnej.

\begin{definicja}[Przestrzeń $W^{1,p}$]
Niech $\Omega\subset\R^d$ będzie otwarta i $1\le p\le\infty$.
Przez $W^{1,p}(\Omega)$ oznaczamy zbiór tych
$u\in L^p(\Omega)$, dla których wszystkie pochodne
dystrybucyjne $\partial_jT_u$ są reprezentowane przez
funkcje z $L^p(\Omega)$. Przestrzeń wyposażamy w normę
$\|u\|_{W^{1,p}}=\|u\|_{L^p}+\sum_{j=1}^d
\|\partial_ju\|_{L^p}$. Pełną teorię przestrzeni Sobolewa,
włącznie ze zupełnością i twierdzeniami o zanurzeniach,
rozwiniemy w rozdziale 9.
\end{definicja}

\begin{uwaga}[Jak ten skrót ma się do dalszych rozdziałów?]
W tym rozdziale przestrzeń $W^{1,p}(\Omega)$ pojawia się
jedynie jako pierwszy przykład funkcji o słabych pochodnych.
Rozdział~\ref{chap:pde-sobolev} zbuduje ogólną teorię
$W^{k,p}(\Omega)$ na obszarach, wraz z zupełnością,
aproksymacją i przedłużaniem. Z kolei
rozdział~\ref{chap:pde-fourier-Rd} wprowadza
przestrzenie $H^s(\R^d)$ definiowane przez transformatę Fouriera.
Dla całkowitych rzędów na całej przestrzeni zachodzi
$H^m(\R^d)=W^{m,2}(\R^d)$, lecz na obszarach są to różne
konstrukcje i nie należy utożsamiać ich definicji.
\end{uwaga}

\begin{przyklad}[Funkcja z załamaniem i funkcja ze skokiem]
Na $(-1,1)$ funkcja $u(x)=|x|$ należy do
$W^{1,p}((-1,1))$ dla każdego $1\le p\le\infty$:
jej słabą pochodną jest $\operatorname{sgn}(x)$.
Natomiast funkcja Heaviside'a ograniczona do $(-1,1)$
nie należy do $W^{1,1}((-1,1))$, ponieważ jej
pochodna dystrybucyjna jest deltą.
Obie funkcje są niegładkie klasycznie, lecz tylko jedna
ma całkowalną słabą pochodną.
\end{przyklad}

\begin{twierdzenie}[Pochodna zerowa oznacza stałość w wymiarze jeden]
\label{thm:pde7-constant}
Niech $I\subset\R$ będzie przedziałem otwartym,
$T\in\mathcal D'(I)$ i $T'=0$. Istnieje dokładnie jedna
stała $c\in\C$ taka, że $T=T_c$ na $I$.
\end{twierdzenie}
\begin{proof}
Wybierzmy $\eta\in\mathcal D(I)$ o
$\int_I\eta(x)\,\mathrm d\lambda_1(x)=1$ i połóżmy
$c=\langle T,\eta\rangle$.
Dla dowolnej $\varphi\in\mathcal D(I)$ zdefiniujmy
$\psi=\varphi-\bigl(\int_I\varphi(x)\,\mathrm d\lambda_1(x)\bigr)\eta$.
Wtedy $\int_I\psi(x)\,\mathrm d\lambda_1(x)=0$. Ponieważ $I$ jest
przedziałem, funkcja $\Psi(x)=\int_{-\infty}^x\psi(t)
\,\mathrm d\lambda_1(t)$, po przedłużeniu $\psi$ zerem,
ma zwarty nośnik zawarty w $I$ i jest klasy $C^\infty$.
Mamy $\Psi'=\psi$, a więc
$\langle T,\psi\rangle=\langle T,\Psi'\rangle
=-\langle T',\Psi\rangle=0$. Otrzymujemy
$\langle T,\varphi\rangle=c\int_I\varphi(x)\,\mathrm d\lambda_1(x)$.
Jednoznaczność $c$ wynika z wyboru testu o niezerowej całce.
\end{proof}

\section{Rozwiązania fundamentalne operatorów różniczkowych}
\label{sec:pde7-fundamental}

\begin{definicja}[Rozwiązanie fundamentalne]
Niech $P(\partial)=\sum_{|\alpha|\le m}a_\alpha\partial^\alpha$
będzie operatorem różniczkowym o stałych współczynnikach
na $\R^d$, gdzie $a_\alpha\in\C$.
Dystrybucję $E\in\mathcal D'(\R^d)$ nazywamy
\emph{rozwiązaniem fundamentalnym} $P$, gdy
$P(\partial)E=\delta_0$ w $\mathcal D'(\R^d)$.
\end{definicja}

Jeśli $f\in C_c^\infty(\R^d)$, możemy zdefiniować
$(E*f)(x)=\langle E(y),f(x-y)\rangle$.
Funkcja ta jest gładka i spełnia
$P(\partial)(E*f)=(P(\partial)E)*f=\delta_0*f=f$.
Stąd rozwiązanie fundamentalne daje narzędzie rozwiązywania
niejednorodnego równania na całej przestrzeni.
Nie twierdzimy na tym etapie, że $E*f$ opisuje wszystkie
rozwiązania ani że bez dodatkowych warunków można w ten sposób
rozwiązywać dowolne zadanie brzegowe.


\begin{lemat}[Różniczkowanie splotu z dystrybucją]
\label{lem:pde7-convolution-derivatives}
Niech $d\ge1$, $T\in\mathcal D'(\R^d)$ i
$\psi\in C_c^\infty(\R^d)$. Funkcja
$u(x)=\langle T(y),\psi(x-y)\rangle$ należy do
$C^\infty(\R^d)$ i dla każdego wielowskaźnika
$\alpha\in\N_0^d$ zachodzi
\[
 \partial^\alpha(T*\psi)
 =T*(\partial^\alpha\psi)=(\partial^\alpha T)*\psi.
\]
W szczególności dla operatora o stałych współczynnikach
$P(\partial)$ i rozwiązania fundamentalnego $E$ mamy
$P(\partial)(E*\psi)=\psi$.
\end{lemat}
\begin{proof}
Dla $x$ w ustalonym zbiorze zwartym nośniki wszystkich
funkcji $y\mapsto\psi(x-y)$ są zawarte w jednym zwartym
zbiorze. Odwzorowanie $x\mapsto\psi(x-\cdot)$ jest więc
gładkie jako odwzorowanie w przestrzeń funkcji testowych:
każda pochodna względem parametru $x$ powstaje przez
różniczkowanie $\psi$. Ciągłość funkcjonału $T$ daje
$\partial_x^\alpha u(x)=\langle T(y),
\partial_x^\alpha\psi(x-y)\rangle$.
Dla pochodnej pierwszego rzędu korzystamy z definicji
pochodnej dystrybucyjnej:
\[
 \langle\partial_{y_j}T(y),\psi(x-y)\rangle
 =-\langle T(y),\partial_{y_j}\psi(x-y)\rangle
 =\langle T(y),\partial_{x_j}\psi(x-y)\rangle.
\]
Iteracja daje wszystkie wielowskaźniki. Na koniec
$P(\partial)(E*\psi)=(P(\partial)E)*\psi
=\delta_0*\psi=\psi$, gdzie ostatnia równość
wynika bezpośrednio z działania delty na test.
\end{proof}

\begin{twierdzenie}[Rozwiązanie fundamentalne operatora $-\Delta$]
\label{thm:pde7-laplace}
Niech $d\ge1$ i niech $\omega_d$ oznacza pole powierzchni
sfery jednostkowej $S^{d-1}\subset\R^d$ dla $d\ge2$.
Funkcja lokalnie całkowalna
\begin{equation}\label{eq:pde7-laplace-fundamental}
 \Phi_d(x)=\begin{cases}
 -|x|/2,&d=1,\\
 (2\pi)^{-1}\log(1/|x|),&d=2,\\
 ((d-2)\omega_d)^{-1}|x|^{2-d},&d\ge3,
 \end{cases}
\end{equation}
spełnia $-\Delta T_{\Phi_d}=\delta_0$ na $\R^d$.
\end{twierdzenie}
\begin{proof}
W wymiarze jeden stwierdzenie wynika z
$\partial_x^2|x|=2\delta_0$.
Dla $d\ge2$ funkcja $\Phi_d$ jest harmoniczna poza zerem
i spełnia $\partial_r\Phi_d(r)=-1/(\omega_dr^{d-1})$.
Dla $\varphi\in\mathcal D(\R^d)$ wybierzmy $R>0$
tak, że $\supp\varphi\Subset B_R$, i stosujmy drugą
tożsamość Greena na pierścieniu $B_R\setminus\overline{B_\varepsilon}$:
\[
 -\int_{B_R\setminus B_\varepsilon}\Phi_d(x)\Delta\varphi(x)
 \,\mathrm d\lambda_d(x)
 =\int_{\partial B_\varepsilon}\varphi(x)
 \frac{1}{\omega_d\varepsilon^{d-1}}\,\mathrm d\sigma(x)
 -\int_{\partial B_\varepsilon}\Phi_d(x)
 \partial_\nu\varphi(x)\,\mathrm d\sigma(x).
\]
Tu $\nu$ oznacza normalną zewnętrzną \emph{pierścienia},
a więc na wewnętrznym brzegu $\nu=-x/|x|$.
Pierwsza całka dąży do $\varphi(0)$; druga ma moduł
$O(\varepsilon)$ dla $d\ge3$ oraz
$O(\varepsilon|\log\varepsilon|)$ dla $d=2$.
Całka objętościowa na $B_\varepsilon$ jest znikoma,
ponieważ $\Phi_d\in L^1_{\mathrm{loc}}$.
Przechodząc do granicy, otrzymujemy
$\langle-\Delta T_{\Phi_d},\varphi\rangle=\varphi(0)$.
\end{proof}

\begin{uwaga}[Zgodność z wcześniejszą teorią potencjału]
W rozdziale~\ref{chap:pde-laplace} otrzymaliśmy te same
funkcje $\Phi_d$ w klasycznej teorii potencjału.
Tutaj wykazaliśmy dokładniejszą tożsamość:
nawet jeśli $\Phi_d$ nie ma klasycznej drugiej pochodnej
w zerze, jej laplasjan dystrybucyjny jest dokładnie
miarą punktową z odpowiednim znakiem.
\end{uwaga}

\section{Równanie ciepła i delta w chwili początkowej}
\label{sec:pde7-heat}

\begin{twierdzenie}[Jądro cieplne jako rozwiązanie fundamentalne]
\label{thm:pde7-heat}
Niech $d\ge1$ i $\kappa>0$. Na $\R^{1+d}$ definiujemy
funkcję lokalnie całkowalną
\[
 E(t,x)=\begin{cases}
 (4\pi\kappa t)^{-d/2}
 \exp\!\bigl(-|x|^2/(4\kappa t)\bigr),&t>0,\\
 0,&t\le0.
 \end{cases}
\]
Wtedy dla $L=\partial_t-\kappa\Delta_x$ zachodzi
$L T_E=\delta_{(0,0)}$ w $\mathcal D'(\R^{1+d})$.
\end{twierdzenie}
\begin{proof}
Oznaczmy przez $\Gamma_\kappa(t,x)$ dodatnią część
funkcji $E$. Dla $t>0$ mamy
$(\partial_t-\kappa\Delta_x)\Gamma_\kappa=0$ i
$\int_{\R^d}\Gamma_\kappa(t,x)\,\mathrm d\lambda_d(x)=1$.
Niech $\varphi\in\mathcal D(\R^{1+d})$. Z definicji
pochodnej dystrybucyjnej i dwukrotnego całkowania
przez części względem $x$ otrzymujemy
\[
 \langle LT_E,\varphi\rangle
 =\lim_{\varepsilon\downarrow0}
 \int_\varepsilon^\infty\int_{\R^d}
 \Gamma_\kappa(t,x)
 [-\partial_t\varphi(t,x)-\kappa\Delta_x\varphi(t,x)]
 \,\mathrm d\lambda_d(x)\,\mathrm d\lambda_1(t).
\]
Na obszarze $t>\varepsilon$ używamy równania
spełnianego przez $\Gamma_\kappa$ i całkowania przez części:
całka równa się
$\int_{\R^d}\Gamma_\kappa(\varepsilon,x)
\varphi(\varepsilon,x)\,\mathrm d\lambda_d(x)$.
Po podstawieniu $x=\sqrt{4\kappa\varepsilon}\,z$
dostajemy całkę Gaussowską z
$\varphi(\varepsilon,\sqrt{4\kappa\varepsilon}\,z)$,
która dąży do $\varphi(0,0)$ na mocy zbieżności
zdominowanej. To oznacza dokładnie
$LT_E=\delta_{(0,0)}$.
\end{proof}

\begin{uwaga}[Nie jest to rozwiązanie klasyczne w $t=0$]
Funkcja $E$ jest klasycznie gładka dla $t>0$, ale
jej przedłużenie przez zero dla $t\le0$ nie
spełnia równania jednorodnego na całej czasoprzestrzeni.
Osobliwość w chwili $t=0$ koduje punktowe źródło.
To rozróżnienie będzie ważne dla wzoru Duhamela,
źródeł impulsowych i równań ewolucyjnych.
\end{uwaga}

\section{Lemat Weyla: rozwiązanie dystrybucyjne jest gładkie}
\label{sec:pde7-weyl}

\begin{twierdzenie}[Lemat Weyla]
\label{thm:pde7-weyl}
Niech $d\ge2$, $\Omega\subset\R^d$ będzie otwarta,
a $T\in\mathcal D'(\Omega)$ spełnia $\Delta T=0$.
Wówczas istnieje dokładnie jedna harmoniczna funkcja
$h\in C^\infty(\Omega)$ taka, że $T=T_h$.
W szczególności $h$ jest rzeczywiście analityczna.
\end{twierdzenie}
\begin{proof}
Ustalmy $a\in\Omega$ i $r>0$ takie, że
$\overline{B_{3r}(a)}\subset\Omega$. Dla
$0<\varepsilon<r$ połóżmy $u_\varepsilon=T*\rho_\varepsilon$
na $B_{2r}(a)$. Ze wzoru \eqref{eq:pde7-deriv-mollify}
wynika $\Delta u_\varepsilon=0$.
Wybierzmy radialną $\eta\in C_c^\infty(B_r(0))$,
$\eta\ge0$, o całce $1$. Dla $x\in B_r(a)$
własność wartości średniej dla funkcji harmonicznych
i rozkład całki na sfery dają
\[
 u_\varepsilon(x)=\int_{\R^d}
 u_\varepsilon(y)\eta(x-y)\,\mathrm d\lambda_d(y).
\]
Prawa strona jest $u_\varepsilon*\eta$ na $B_r(a)$.
Wiemy, że $u_\varepsilon\to T$ w $\mathcal D'(\Omega)$.
Przejście do granicy uzasadnimy wprost na funkcjach testowych.
Niech $\varphi\in\mathcal D(B_r(a))$. Ponieważ
$\eta$ ma nośnik w $B_r(0)$, funkcja
$\check\eta*\varphi$ ma nośnik w $B_{2r}(a)$,
gdzie zdefiniowano $u_\varepsilon$; przyjmujemy
$\check\eta(x)=\eta(-x)=\eta(x)$. Z twierdzenia Fubiniego
i definicji splotu funkcji gładkich dostajemy
\[
 \langle T,\varphi\rangle
 =\lim_{\varepsilon\downarrow0}
 \int_{B_r(a)}u_\varepsilon(x)\varphi(x)
 \,\mathrm d\lambda_d(x)
 =\lim_{\varepsilon\downarrow0}
 \int_{B_{2r}(a)}u_\varepsilon(y)
 (\check\eta*\varphi)(y)\,\mathrm d\lambda_d(y)
 =\langle T,\check\eta*\varphi\rangle.
\]
Ostatnia liczba jest, z definicji lokalnego splotu,
równa $\langle T*\eta,\varphi\rangle$.
Pokazaliśmy więc $T=T*\eta$ na $B_r(a)$
bez nieuzasadnionego zamieniania granic w rodzinie
przesuniętych testów. Prawa strona jest klasy $C^\infty$,
więc $T$ ma gładką
reprezentację w otoczeniu $a$. Ponieważ $a$ jest dowolny,
otrzymujemy funkcję $h\in C^\infty(\Omega)$;
iniektywność regularnego osadzenia zapewnia zgodność
lokalnych reprezentacji. Z $\Delta T_h=0$ i
iniektywności wynika $\Delta h=0$ punktowo.
Analityczność wykazano już w
\ref{thm:pde3-analytic}.
\end{proof}

\section{Ograniczenia rachunku dystrybucji}
\label{sec:pde7-limits}

Pochodne i mnożenie przez funkcje gładkie są dobrze określone
na $\mathcal D'(\Omega)$, ale \emph{nie istnieje kanoniczny iloczyn
dowolnych dwóch dystrybucji}, który jednocześnie zachowuje
wszystkie naturalne własności mnożenia funkcji i reguły
różniczkowania. Zapis $\delta_0^2$ nie ma więc standardowego
znaczenia w teorii Schwartza. Analogicznie wyrażenie $u\,\partial_xu$
dla nieciągłego $u$ nie musi być określone dystrybucyjnie.

Dla równań zachowania rozwiązaniem jest najczęściej
\emph{postać dywergencyjna}, np.
$\partial_tu+\partial_x(u^2/2)=0$:
jeśli $u\in L^2_{\mathrm{loc}}$, to $u^2/2\in L^1_{\mathrm{loc}}$
i oba składniki są dystrybucjami. Nie oznacza to jeszcze
jednoznaczności rozwiązań; do niej potrzebne bywają
warunki entropijne omówione wcześniej i rozwinięte
w rozdziale o prawach zachowania.

\begin{przyklad}[Skąd bierze się problem z $\delta^2$?]
Niech $\rho_\varepsilon(x)=\varepsilon^{-1}\rho(x/\varepsilon)$
na $\R$, gdzie $\rho\in C_c^\infty(\R)$, $\rho\ge0$,
$\int_\R\rho(x)\,\mathrm d\lambda_1(x)=1$ i $\rho\not\equiv0$.
Mamy $\rho_\varepsilon\to\delta_0$, lecz dla testu
$\varphi$ z $\varphi(0)>0$
\[
 \int_\R \rho_\varepsilon(x)^2\varphi(x)
 \,\mathrm d\lambda_1(x)
 =\varepsilon^{-1}\int_\R \rho(y)^2\varphi(\varepsilon y)
 \,\mathrm d\lambda_1(y)\longrightarrow+\infty.
\]
Zatem kwadraty tej naturalnej aproksymacji delty
nie mają granicy w $\mathcal D'(\R)$.
Nie dowodzi to samo przez się żadnego twierdzenia o
niemożliwości wszystkich możliwych rozszerzeń iloczynu,
ale jasno pokazuje, dlaczego naiwny iloczyn nie działa.
\end{przyklad}

\section[Podsumowanie: przygotowanie do analizy Fouriera]{Podsumowanie i przygotowanie do analizy Fouriera na $\R^d$}
Dystrybucje pozwalają różniczkować obiekty niegładkie
bez tracenia zgodności z klasyczną analizą. Przybliżone jedynki
łączą rachunek dystrybucyjny z aproksymacją w $L^p$.
Rozwiązania fundamentalne zamieniają liniowe PDE o stałych
współczynnikach na operacje splotu, a lemat Weyla pokazuje,
że niektóre równania narzucają regularność nawet wtedy,
gdy nie zakładamy jej na początku. W następnym rozdziale
zbadamy transformatę Fouriera na przestrzeni Schwartza
i dystrybucjach temperowanych; to umożliwi m.in. opis
symboli operatorów i jąder potencjałowych w języku mnożników.



\chapter[Transformata Fouriera i dystrybucje temperowane]{Transformata Fouriera, przestrzeń Schwartza i dystrybucje temperowane}
\label{chap:pde-fourier-Rd}

Rozdział~\ref{chap:pde-distributions} pozwolił różniczkować funkcje
niegładkie i rozwiązywać równania o stałych współczynnikach za pomocą
rozwiązań fundamentalnych. Transformata Fouriera realizuje komplementarny
pomysł: zamienia pochodne na mnożenie przez wielomiany. Na całej
$\R^d$ nie możemy jednak bez zastrzeżeń używać klasycznej transformaty
całkowej dla wszystkich dystrybucji. Potrzebujemy więc przestrzeni
Schwartza $\mathcal S(\R^d)$ i jej przestrzeni dualnej $\mathcal S'(\R^d)$.

W całym rozdziale $d\ge1$, całki na $\R^d$ są liczone względem
miary Lebesgue'a $\lambda_d$, a wielowskaźniki należą do $\N_0^d$.
Przyjmujemy \emph{nieunitarną} konwencję
\begin{equation}\label{eq:pde8-convention}
 \widehat f(\xi)=\mathcal F f(\xi)
 =\int_{\R^d}e^{-ix\cdot\xi}f(x)\,\mathrm d\lambda_d(x),\qquad
 \mathcal F^{-1}g(x)=(2\pi)^{-d}
 \int_{\R^d}e^{ix\cdot\xi}g(\xi)\,\mathrm d\lambda_d(\xi).
\end{equation}
Współczynniki $(2\pi)^{-d}$ zależą od konwencji i nie będą
pomijane. W szczególności nie będziemy bez przeliczenia używać
wzorów wyprowadzonych dla transformaty unitarnej.
Nie należy też utożsamiać tej transformaty całkowej na $\R^d$
z dyskretnymi współczynnikami szeregów Fouriera z
rozdziału~\ref{chap:pde-fourier}: te ostatnie są uśredniane
przez czynnik $(2\pi)^{-d}$.

\section{Przestrzeń Schwartza: po co wymagać szybkiego zaniku?}
\label{sec:pde8-schwartz}

\begin{definicja}[Przestrzeń Schwartza]
\label{def:pde8-schwartz}
Przestrzenią Schwartza na $\R^d$ nazywamy
\[
 \mathcal S(\R^d)=\bigl\{f\in C^\infty(\R^d;\C):
 p_{\alpha,\beta}(f):=\sup_{x\in\R^d}
 |x^\alpha\partial^\beta f(x)|<\infty
 \ \text{dla wszystkich }\alpha,\beta\in\N_0^d\bigr\}.
\]
Wyposażamy ją w topologię wyznaczoną przez przeliczalną rodzinę
półnorm $p_{\alpha,\beta}$. Zbieżność $f_n\to f$ w $\mathcal S$
oznacza zbieżność do zera wszystkich półnorm $p_{\alpha,\beta}(f_n-f)$.
\end{definicja}

Warunek ten mówi, że funkcja i wszystkie jej pochodne zanikają szybciej
niż dowolna odwrotność wielomianu. Funkcja $e^{-|x|^2}$ należy do
$\mathcal S$, podobnie każda funkcja z $C_c^\infty$; natomiast
$(1+|x|^2)^{-1}$ na ogół nie należy do $\mathcal S$.

\begin{lemat}[Równoważne półnormy i operacje]
\label{lem:pde8-seminorms}
Dla $N,m\in\N_0$ połóżmy
$q_{N,m}(f)=\max_{|\beta|\le m}\sup_x(1+|x|)^N|\partial^\beta f(x)|$.
Rodzina $q_{N,m}$ określa tę samą topologię co $p_{\alpha,\beta}$.
Operator $f\mapsto x_jf$ oraz $f\mapsto\partial_j f$ jest ciągły
$\mathcal S\to\mathcal S$. Ponadto $\mathcal S(\R^d)\subset L^1\cap L^2$.
\end{lemat}
\begin{proof}
Każdy monomian $|x^\alpha|$ jest nie większy niż
$(1+|x|)^{|\alpha|}$, a $(1+|x|)^N$ można oszacować sumą
skończonej liczby $|x^\alpha|$ dla $|\alpha|\le N$, z odpowiednią
stałą zależną od $N,d$. Stąd równoważność rodzin półnorm.
Reguła Leibniza pokazuje, że półnorma pochodnej $x_jf$ jest
kontrolowana przez skończoną sumę półnorm $f$; dla pochodnej
$\partial_jf$ jest to natychmiastowe. Wybierając $N>d$, mamy
$|f(x)|\le q_{N,0}(f)(1+|x|)^{-N}$, co daje całkowalność.
Dla $L^2$ wystarcza $2N>d$.
\end{proof}

\begin{twierdzenie}[Gęstość funkcji o zwartym nośniku]
\label{thm:pde8-density}
Przestrzeń $C_c^\infty(\R^d)$ jest gęsta w $\mathcal S(\R^d)$.
Jest też gęsta w $L^p(\R^d)$ dla $1\le p<\infty$.
\end{twierdzenie}
\begin{proof}
Weźmy $\chi\in C_c^\infty(\R^d)$ taką, że
$\chi=1$ na $B_1(0)$ i $\chi=0$ poza $B_2(0)$.
Dla $f\in\mathcal S$ definiujemy $f_R(x)=\chi(x/R)f(x)$.
Różniczkując $(1-\chi(x/R))f(x)$, dostajemy wyrazy
zawierające $R^{-|\gamma|}\partial^\gamma\chi(x/R)
\partial^{\beta-\gamma}f(x)$, o nośnikach w $|x|\ge R$.
Każdy taki wyraz, pomnożony przez $x^\alpha$, ma supremum
dążące do zera: wybieramy półnormę $f$ z potęgą
$(1+|x|)^{|\alpha|+1}$ (lub większą) i korzystamy z
$|x|\ge R$. Zatem $f_R\to f$ w $\mathcal S$.
Gęstość w $L^p$ uzyskujemy najpierw przez obcinanie funkcji
$L^p$, potem wygładzanie przybliżoną jedynką z
rozdziału~\ref{chap:pde-distributions}; obie aproksymacje
są zbieżne w normie $L^p$.
\end{proof}

\section{Transformata Fouriera na przestrzeni Schwartza}
\label{sec:pde8-transform}

\begin{twierdzenie}[Różniczkowanie i mnożenie]
\label{thm:pde8-rules}
Jeżeli $f\in\mathcal S(\R^d)$, to $\widehat f\in\mathcal S(\R^d)$
oraz dla każdego $\alpha\in\N_0^d$
\begin{equation}\label{eq:pde8-rules}
 \mathcal F(\partial^\alpha f)(\xi)=(i\xi)^\alpha\widehat f(\xi),
 \qquad \mathcal F(x^\alpha f)(\xi)=i^{|\alpha|}
 \partial_\xi^\alpha\widehat f(\xi).
\end{equation}
Ponadto $\mathcal F:\mathcal S\to\mathcal S$ jest ciągła.
\end{twierdzenie}
\begin{proof}
Ponieważ $f$ wraz z pochodnymi zanika szybko, całkowanie przez
części nie wytwarza składników brzegowych. Dla pierwszej tożsamości:
$\widehat{\partial_jf}=i\xi_j\widehat f$.
Drugą otrzymujemy przez różniczkowanie pod całką:
$\partial_{\xi_j}\widehat f=\int_{\R^d}(-ix_j)e^{-ix\cdot\xi}f(x)
\,\mathrm d\lambda_d(x)$.
Powtarzając te operacje otrzymujemy
$\xi^\alpha\partial_\xi^\beta\widehat f$
jako kombinację transformat pochodnych funkcji
$x^\beta f$. Każdą taką transformatę szacujemy przez
normę $L^1$ odpowiedniej funkcji, a tę przez skończoną sumę
półnorm Schwartza za pomocą lematu~\ref{lem:pde8-seminorms}.
Daje to jednocześnie szybki zanik $\widehat f$ i ciągłość
odwzorowania $\mathcal F$ w topologii $\mathcal S$.
\end{proof}

\begin{stwierdzenie}[Lemat Riemanna--Lebesgue'a]
\label{prop:pde8-RL}
Jeśli $f\in L^1(\R^d)$, to $\widehat f$ jest ograniczona,
ciągła i $\widehat f(\xi)\to0$ przy $|\xi|\to\infty$.
\end{stwierdzenie}
\begin{proof}
Oszacowanie $|\widehat f(\xi)|\le\|f\|_{L^1}$ jest natychmiastowe.
Ciągłość wynika ze zbieżności zdominowanej. Dla
$g\in C_c^\infty$ poprzednie twierdzenie daje zanik
$\widehat g$. Ponieważ $C_c^\infty$ jest gęsta w $L^1$,
a $\|\widehat f-\widehat g\|_\infty\le\|f-g\|_1$,
przenosimy zanik na $f$.
\end{proof}

\section{Transformata funkcji Gaussa i wzór odwrócenia}
\label{sec:pde8-inversion}

\begin{lemat}[Transformata funkcji Gaussa]
\label{lem:pde8-gauss}
Dla $a>0$ i $x,\xi\in\R^d$ zachodzi
\begin{equation}\label{eq:pde8-gauss}
 \mathcal F(e^{-a|x|^2})(\xi)
 =\left(\frac\pi a\right)^{d/2}e^{-|\xi|^2/(4a)}.
\end{equation}
\end{lemat}
\begin{proof}
W wymiarze jeden funkcja
$I(\xi)=\int_\R e^{-ax^2}e^{-ix\xi}\,\mathrm d\lambda_1(x)$
spełnia po różniczkowaniu pod całką i całkowaniu przez części
$I'(\xi)=-\xi I(\xi)/(2a)$.
Ponieważ $I(0)=\sqrt{\pi/a}$ (klasyczna całka Gaussa),
otrzymujemy $I(\xi)=\sqrt{\pi/a}\,e^{-\xi^2/(4a)}$.
W $d$ wymiarach stosujemy twierdzenie Fubiniego.
\end{proof}

\begin{twierdzenie}[Odwrócenie transformaty]
\label{thm:pde8-inversion}
Dla $f\in\mathcal S(\R^d)$ zachodzi
\[
 f(x)=(2\pi)^{-d}\int_{\R^d}e^{ix\cdot\xi}
 \widehat f(\xi)\,\mathrm d\lambda_d(\xi),\qquad x\in\R^d.
\]
Zatem $\mathcal F$ jest topologicznym automorfizmem
$\mathcal S(\R^d)$, którego odwrotność podano w
\eqref{eq:pde8-convention}.
\end{twierdzenie}
\begin{proof}
Niech $\varepsilon>0$. Rozważmy funkcję
\[
 I_\varepsilon(x)=(2\pi)^{-d}\int_{\R^d}
 e^{ix\cdot\xi}e^{-\varepsilon|\xi|^2/2}
 \widehat f(\xi)\,\mathrm d\lambda_d(\xi).
\]
Funkcja podcałkowa jest całkowalna, więc wolno podstawić
wzór na $\widehat f$ i zamienić kolejność całkowania.
Na mocy lematu~\ref{lem:pde8-gauss} dostajemy
$I_\varepsilon=f*G_\varepsilon$, gdzie
$G_\varepsilon(z)=(2\pi\varepsilon)^{-d/2}
 e^{-|z|^2/(2\varepsilon)}$.
Ponieważ $G_\varepsilon$ jest przybliżoną jedynką,
$f*G_\varepsilon(x)\to f(x)$ (dla każdej ustalonej $x$,
a nawet jednostajnie w $x$, dzięki jednostajnej ciągłości $f$).
Z drugiej strony $\widehat f\in L^1$, zatem zdominowana
zbieżność przy $\varepsilon\downarrow0$ daje
$I_\varepsilon(x)\to(2\pi)^{-d}
\int_{\R^d}e^{ix\cdot\xi}\widehat f(\xi)\,\mathrm d\lambda_d(\xi)$.
To dowodzi wzoru. Równoważne wzory dla transformaty
odwrotnej wynikają przez podstawienie $\xi\mapsto-\xi$.
Ciągłość obu operatorów w $\mathcal S$ wynika z
oszacowań dowodu twierdzenia~\ref{thm:pde8-rules}.
\end{proof}

\begin{wniosek}[Wzory na splot i iloczyn]
\label{cor:pde8-convolution}
Dla $f,g\in\mathcal S(\R^d)$ mamy
\[
 \widehat{f*g}=\widehat f\,\widehat g,\qquad
 \widehat{fg}=(2\pi)^{-d}(\widehat f*\widehat g).
\]
\end{wniosek}
\begin{proof}
Pierwsza równość wynika z podstawienia $z=x-y$ oraz Fubiniego.
Druga: wstawiamy wzór odwrócenia dla $g$ do całki
definiującej $\widehat{fg}$, a następnie stosujemy Fubiniego.
Wszystkie użyte funkcje i ich transformaty należą do $L^1$.
\end{proof}

\section[Twierdzenie Plancherela i przestrzeń L2]{Twierdzenie Plancherela i przestrzeń $L^2$}
\label{sec:pde8-plancherel}

\begin{twierdzenie}[Tożsamość Parsevala--Plancherela]
\label{thm:pde8-plancherel}
Dla $f,g\in\mathcal S(\R^d)$ zachodzi
\begin{equation}\label{eq:pde8-parseval}
 \int_{\R^d}f(x)\overline{g(x)}\,\mathrm d\lambda_d(x)
 =(2\pi)^{-d}\int_{\R^d}\widehat f(\xi)
 \overline{\widehat g(\xi)}\,\mathrm d\lambda_d(\xi).
\end{equation}
Transformata Fouriera przedłuża się w sposób jednoznaczny do
przekształcenia ciągłego $\mathcal F:L^2(\R^d)\to L^2(\R^d)$,
dla którego $\|\widehat f\|_2=(2\pi)^{d/2}\|f\|_2$.
Przekształcenie $(2\pi)^{-d/2}\mathcal F$ jest unitarne
na $L^2(\R^d)$.
\end{twierdzenie}
\begin{proof}
Ze wzoru odwrócenia dla $g$ mamy
$\overline{g(x)}=(2\pi)^{-d}\int_{\R^d}e^{-ix\cdot\xi}
\overline{\widehat g(\xi)}\,\mathrm d\lambda_d(\xi)$.
Mnożymy przez $f(x)$, całkujemy po $x$ i stosujemy
Fubiniego (iloczyn norm $L^1$ jest skończony), co daje
\eqref{eq:pde8-parseval}. Dla $f=g$ otrzymujemy równość norm.
Gęstość $\mathcal S$ w $L^2$ pozwala przedłużyć
$(2\pi)^{-d/2}\mathcal F$ do izometrii na całym $L^2$.
Jej obraz jest domknięty, ponieważ $L^2$ jest zupełna;
zarazem zawiera $\mathcal F(\mathcal S)=\mathcal S$, zbiór gęsty.
Zatem obraz jest całym $L^2$, czyli przedłużenie jest unitarne.
\end{proof}

\begin{uwaga}[Właściwe znaczenie całki dla $L^2$]
Jeśli $f\in L^2\setminus L^1$, zapis $\widehat f$ oznacza
powyższe przedłużenie operatora, a nie automatycznie
zbieżną całkę z \eqref{eq:pde8-convention}.
Na przykład funkcja $f(x)=(1+|x|)^{-3/4}$ na $\R$ należy
do $L^2$, lecz nie do $L^1$.
\end{uwaga}

\section[Dystrybucje temperowane]{Dystrybucje temperowane: przestrzeń dualna do $\mathcal S$}
\label{sec:pde8-tempered}

\begin{definicja}[Dystrybucja temperowana]
\label{def:pde8-tempered}
Przestrzenią dystrybucji temperowanych $\mathcal S'(\R^d)$
nazywamy przestrzeń dualną, złożoną ze wszystkich ciągłych
funkcjonałów liniowych $T:\mathcal S(\R^d)\to\C$.
Zbieżność $T_n\to T$ w sensie dystrybucji temperowanych
oznacza zbieżność $\langle T_n,\varphi\rangle\to
\langle T,\varphi\rangle$ dla każdej $\varphi\in\mathcal S$.
\end{definicja}

\begin{twierdzenie}[Globalne oszacowanie dystrybucji temperowanej]
\label{thm:pde8-growth}
Funkcjonał liniowy $T$ na $\mathcal S(\R^d)$ jest ciągły
dokładnie wtedy, gdy istnieją $C>0$ i $N,m\in\N_0$ takie, że
\begin{equation}\label{eq:pde8-growth}
 |\langle T,\varphi\rangle|\le C\,q_{N,m}(\varphi)
 \qquad(\varphi\in\mathcal S(\R^d)).
\end{equation}
\end{twierdzenie}
\begin{proof}
Ciągłość w zerze oznacza ograniczenie $|T(\varphi)|<1$
na pewnym otoczeniu zdefiniowanym przez skończenie wiele
półnorm $q_{N_j,m_j}$. Ponieważ rodzina $q_{N,m}$ jest
monotoniczna względem obu indeksów, jedno $q_{N,m}$
kontroluje je wszystkie. Skalowanie daje żądane oszacowanie.
Odwrotna implikacja jest bezpośrednia z definicji topologii.
\end{proof}

\begin{przyklad}[Obiekty temperowane i nietemperowane]
\begin{enumerate}
\item Każda funkcja $f$ spełniająca $|f(x)|\le C(1+|x|)^M$
(nawet tylko lokalnie całkowalna, z oszacowaniem p.w.)
określa element $\mathcal S'$ przez całkowanie. Należą tu
wielomiany, stałe i funkcje ograniczone.
\item Każda delta $\delta_a$ i jej pochodne należą do $\mathcal S'$.
\item Każda dystrybucja o zwartym nośniku należy do $\mathcal S'$:
po odcięciu testu funkcją równą $1$ w otoczeniu nośnika
lokalne oszacowanie z rozdziału~\ref{chap:pde-distributions}
przechodzi w globalne oszacowanie \eqref{eq:pde8-growth}.
\item Dystrybucja regularna wyznaczona przez $e^{|x|^2}$
na $\R^d$ \emph{nie} jest temperowana. Gdyby przedłużała
się do ciągłego funkcjonału na $\mathcal S$, testy
$\chi(x/R)e^{-|x|^2/2}$ dążyłyby w $\mathcal S$ do
$e^{-|x|^2/2}$, ale ich wartości na tym funkcjonale
rosłyby do $+\infty$. Sprzeczność z ciągłością.
\end{enumerate}
\end{przyklad}

\begin{uwaga}[Relacja między $\mathcal D'$ a $\mathcal S'$]
Ograniczenie funkcjonału z $\mathcal S'$ do
$C_c^\infty(\R^d)$ jest dystrybucją. Jest to odwzorowanie
iniektywne, bo $C_c^\infty$ jest gęsta w $\mathcal S$.
Nie każda dystrybucja na $\R^d$ jest temperowana.
W szczególności zapis $\mathcal F T$ dla dowolnego
$T\in\mathcal D'$ wymaga dodatkowych założeń i
nie ma automatycznie sensu w opisanym rachunku.
\end{uwaga}

\section{Transformata Fouriera dystrybucji temperowanych}
\label{sec:pde8-fourier-distributions}

\begin{definicja}[Transformata na $\mathcal S'$]
\label{def:pde8-distFT}
Dla $T\in\mathcal S'(\R^d)$ definiujemy $\widehat T\in\mathcal S'(\R^d)$
przez \emph{transpozycję} operatora $\mathcal F$:
\[
 \langle\widehat T,\varphi\rangle
 =\langle T,\widehat\varphi\rangle,
 \qquad\varphi\in\mathcal S(\R^d).
\]
Analogicznie definiujemy $\mathcal F^{-1}T$ przez
$\langle\mathcal F^{-1}T,\varphi\rangle
=\langle T,\mathcal F^{-1}\varphi\rangle$.
\end{definicja}

\begin{twierdzenie}[Rachunek na $\mathcal S'$]
\label{thm:pde8-dist-rules}
Transformata $\mathcal F:\mathcal S'\to\mathcal S'$ jest
liniowym automorfizmem, którego odwrotnością jest
$\mathcal F^{-1}$. Dla $T\in\mathcal S'(\R^d)$ i
$\alpha\in\N_0^d$ zachodzi
\[
 \mathcal F(\partial^\alpha T)=(i\xi)^\alpha\widehat T,
 \qquad \mathcal F(x^\alpha T)=i^{|\alpha|}
 \partial_\xi^\alpha\widehat T.
\]
\end{twierdzenie}
\begin{proof}
Ciągłość $\mathcal F:\mathcal S\to\mathcal S$ oraz
jej odwrotności zapewnia, że transpozycje działają na
przestrzeni dualnej. Kompozycja transpozycji odpowiada
odwróceniu kolejności operatorów, więc
$\mathcal F^{-1}\mathcal F T=T$ i $\mathcal F\mathcal F^{-1}T=T$.
Dla pierwszej równości i testu $\varphi$ liczymy
\[
 \langle\widehat{\partial_jT},\varphi\rangle
 =-\langle T,\partial_j\widehat\varphi\rangle
 =\langle T,\widehat{i\xi_j\varphi}\rangle
 =\langle i\xi_j\widehat T,\varphi\rangle.
\]
Dla drugiej korzystamy z $\widehat{\partial_{\xi_j}\varphi}
=ix_j\widehat\varphi$ (w nazwach zmiennych odpowiednich
obszarów) i z definicji pochodnej dystrybucyjnej;
pozostałe przypadki otrzymujemy przez iterację.
\end{proof}

\begin{przyklad}[Cztery transformaty podstawowe]
\label{ex:pde8-four}
W $\mathcal S'(\R^d)$ zachodzą wzory
\[
 \widehat{\delta_0}=1,\quad
 \widehat 1=(2\pi)^d\delta_0,\quad
 \widehat{\delta_a}(\xi)=e^{-ia\cdot\xi},\quad
 \mathcal F(e^{ia\cdot x})=(2\pi)^d\delta_a.
\]
Na przykład $\langle\widehat{\delta_0},\varphi\rangle
=\widehat\varphi(0)=\int_{\R^d}\varphi(x)\,\mathrm d\lambda_d(x)$,
a $\widehat1=(2\pi)^d\delta_0$ wynika z odwrócenia
transformaty i symetrii przestrzeni Schwartza.
Równości po prawej należy rozumieć \emph{dystrybucyjnie},
a nie jako zwykłe całki bezwzględnie zbieżne.
\end{przyklad}

\begin{twierdzenie}[Splot dystrybucji temperowanej z funkcją Schwartza]
\label{thm:pde8-convol-temp}
Niech $T\in\mathcal S'(\R^d)$ i $\varphi\in\mathcal S(\R^d)$.
Funkcja $(T*\varphi)(x):=\langle T(y),\varphi(x-y)\rangle$
jest gładka, ma wzrost co najwyżej wielomianowy
(wraz z każdą ustaloną pochodną) i wyznacza dystrybucję
temperowaną. Ponadto
$\mathcal F(T*\varphi)=\widehat T\,\widehat\varphi$ w $\mathcal S'$.
\end{twierdzenie}
\begin{proof}
Dla $x\in\R^d$ funkcja $y\mapsto\varphi(x-y)$ należy do
$\mathcal S$. Różniczkowanie względem $x$ odbywa się
w topologii $\mathcal S$, a więc jest dozwolone po
zastosowaniu $T$. Z \eqref{eq:pde8-growth} oraz nierówności
$(1+|y|)^N\le(1+|x|)^N(1+|x-y|)^N$ wynika
$|\partial_x^\alpha(T*\varphi)(x)|\le C_\alpha(1+|x|)^N$.
Dla dowodu tożsamości Fouriera ustalmy
$\psi\in\mathcal S(\R^d)$. Z definicji transformaty
dystrybucyjnej i splotu otrzymujemy
\[
 \langle\mathcal F(T*\varphi),\psi\rangle
 =\left\langle T(y),\int_{\R^d}
 \varphi(x-y)\widehat\psi(x)\,\mathrm d\lambda_d(x)
 \right\rangle.
\]
Całka wewnętrzna jest funkcją Schwartza zmiennej $y$.
Istotnie, przed całkowaniem jej półnormy Schwartza
szacujemy za pomocą skończonej sumy półnorm $\varphi$
pomnożonych przez potęgi $(1+|x|)$, całkowalne
przeciwko $|\widehat\psi(x)|\,\mathrm d\lambda_d(x)$.
Uzasadnia to zamianę całkowania z działaniem ciągłego
funkcjonału $T$, np. przez przybliżenie całki sumami
Riemanna w topologii $\mathcal S$.
Po podstawieniu $z=x-y$ i użyciu wzoru
$\widehat\psi(x)=\int_{\R^d}e^{-ix\cdot\xi}\psi(\xi)
\,\mathrm d\lambda_d(\xi)$ liczymy, stosując Fubiniego,
\[
 \int_{\R^d}\varphi(x-y)\widehat\psi(x)
 \,\mathrm d\lambda_d(x)
 =\int_{\R^d}e^{-iy\cdot\xi}\widehat\varphi(\xi)
 \psi(\xi)\,\mathrm d\lambda_d(\xi)
 =\mathcal F(\widehat\varphi\,\psi)(y).
\]
Stąd $\langle\mathcal F(T*\varphi),\psi\rangle
=\langle T,\mathcal F(\widehat\varphi\psi)\rangle
=\langle\widehat T,\widehat\varphi\psi\rangle$,
czyli $\mathcal F(T*\varphi)=\widehat T\,\widehat\varphi$.
W szczególności nie pojawia się dodatkowy czynnik
$(2\pi)^{-d}$: w naszej konwencji dotyczy on
transformaty \emph{iloczynu} funkcji, a nie splotu.
\end{proof}

\section{Mnożniki Fouriera: kiedy operator jest ograniczony?}
\label{sec:pde8-multipliers}

\begin{twierdzenie}[Mnożnik w $L^2$]
\label{thm:pde8-multiplier}
Niech $m:\R^d\to\C$ będzie funkcją mierzalną i
istotnie ograniczoną. Dla $f\in L^2(\R^d)$ istnieje dokładnie
jedna funkcja $T_mf\in L^2(\R^d)$ spełniająca
$\widehat{T_mf}=m\widehat f$ w sensie $L^2$.
Operator $T_m$ jest liniowy i ograniczony oraz
$\|T_m\|_{L^2\to L^2}=\|m\|_{L^\infty}$.
\end{twierdzenie}
\begin{proof}
Mnożenie przez $m$ jest ograniczone na $L^2$, a transformata
Fouriera jest unitarnym operatorem z dokładnością do
stałego czynnika, co daje istnienie, jednoznaczność i
oszacowanie z góry. Dla dolnego oszacowania bierzemy
$\widehat f$ proporcjonalną do funkcji charakterystycznej
zbioru dodatniej i skończonej miary, na którym
$|m|>\|m\|_\infty-\varepsilon$. Taki zbiór istnieje
(po przecięciu z dostatecznie dużą kulą), a po
$\varepsilon\downarrow0$ otrzymujemy równość norm.
\end{proof}

\begin{przyklad}[Transformata Riesza w $L^2$]
Dla $1\le j\le d$ połóżmy
$m_j(\xi)=-i\xi_j/|\xi|$ dla $\xi\ne0$ i $m_j(0)=0$.
Poprzednie twierdzenie definiuje ograniczony operator
$R_j=T_{m_j}:L^2\to L^2$, tzw. transformatę Riesza,
z $\|R_j\|\le1$. Wartość mnożnika w pojedynczym punkcie
nie zmienia operatora $L^2$. Ograniczoność w $L^p$ dla
$p\ne2$ jest istotnie trudniejszym rezultatem teorii
Calderóna--Zygmunda i nie wynika z samego warunku
$m_j\in L^\infty$.
\end{przyklad}

\section{Zastosowanie I: równanie ciepła i równanie falowe}
\label{sec:pde8-evolution}

\begin{twierdzenie}[Rozwiązania ewolucyjne poprzez transformatę]
\label{thm:pde8-evolution}
Niech $\kappa,c>0$ i $f,g\in\mathcal S(\R^d)$.
Wtedy funkcje
\begin{align*}
 u(t,x)&=\mathcal F^{-1}
 \bigl(e^{-\kappa t|\xi|^2}\widehat f(\xi)\bigr)(x),&&t\ge0,\\
 v(t,x)&=\mathcal F^{-1}\biggl(
 \cos(ct|\xi|)\widehat f(\xi)
 +\frac{\sin(ct|\xi|)}{c|\xi|}\widehat g(\xi)
 \biggr)(x),&&t\in\R,
\end{align*}
gdzie iloraz w $\xi=0$ definiujemy przez wartość graniczną $t$,
są gładkimi rozwiązaniami odpowiednio zagadnień
$u_t-\kappa\Delta u=0$, $u(0,\cdot)=f$ oraz
$v_{tt}-c^2\Delta v=0$, $v(0,\cdot)=f$, $v_t(0,\cdot)=g$.
W każdym ustalonym czasie funkcje przestrzenne należą do $\mathcal S$.
\end{twierdzenie}
\begin{proof}
Z reguły różniczkowania otrzymujemy po transformacji
$\partial_t\widehat u=-\kappa|\xi|^2\widehat u$ oraz
$\partial_{tt}\widehat v=-c^2|\xi|^2\widehat v$.
Są to dla każdego $\xi$ zwyczajne równania różniczkowe
z wskazanymi danymi początkowymi. Ich rozwiązania mają
właśnie zapis z tezy. Mnożniki są gładkie w $\xi$:
zwłaszcza $\sin(ct|\xi|)/(c|\xi|)$ ma rozwinięcie w szereg
potęg $|\xi|^2$ w otoczeniu zera.
Dla $t$ w dowolnym ograniczonym przedziale ich
pochodne po $t$ i $\xi$ mają co najwyżej wielomianowy wzrost,
stąd mnożą $\mathcal S$ w siebie. Uzasadnia to wszystkie
różniczkowania i stosowanie odwrotnej transformaty.
Wartości początkowe otrzymujemy przez podstawienie $t=0$.
\end{proof}

\begin{uwaga}[Dwie różne skale propagacji]
Równanie ciepła działa w częstotliwościach przez
$m_t(\xi)=e^{-\kappa t|\xi|^2}$: dla $t>0$ tłumi wysokie
częstotliwości bardzo szybko. Równanie falowe używa
$\cos(ct|\xi|)$ i $\sin(ct|\xi|)/(c|\xi|)$; nie ma
analogicznego wygładzania wszystkich wysokich częstotliwości.
To jest spektralna wersja zjawisk z rozdziałów
\ref{chap:pde-heat} i~\ref{chap:pde-wave}.
\end{uwaga}

\section[Potencjały Bessela i przestrzenie Sobolewa]{Potencjały Bessela i pierwsze przestrzenie $H^s$}
\label{sec:pde8-bessel}

\begin{definicja}[Przestrzeń Sobolewa typu Bessela]
\label{def:pde8-Hs}
Dla dowolnego $s\in\R$ definiujemy
\[
 H^s(\R^d)=\bigl\{u\in\mathcal S'(\R^d):
 (1+|\xi|^2)^{s/2}\widehat u\in L^2(\R^d)\bigr\}.
\]
Zapis przynależności do $L^2$ oznacza tutaj, że odpowiednia
dystrybucja jest reprezentowana przez funkcję kwadratowalną.
Normę tej przestrzeni definiujemy wzorem
$\|u\|_{H^s}^2=(2\pi)^{-d}\int_{\R^d}
(1+|\xi|^2)^s|\widehat u(\xi)|^2\,\mathrm d\lambda_d(\xi)$.
W tym wzorze $\widehat u$ jest reprezentowana przez funkcję
lokalnie $L^2$: otrzymujemy ją mnożąc funkcję $L^2$
przez $(1+|\xi|^2)^{-s/2}$.
Operator $J^s=(I-\Delta)^{s/2}$ określamy na $\mathcal S'$
przez mnożnik $(1+|\xi|^2)^{s/2}$.
Dla ujemnego $s$ element $H^s$ może być autentyczną
dystrybucją, a nie funkcją lokalnie całkowalną;
wymóg $L^2$ odnosi się do ważonej transformaty,
nie do samej dystrybucji $u$.
\end{definicja}

Dla $m\in\N_0$ zapis $W^{m,2}(\R^d)$ oznacza przestrzeń
funkcji $u\in L^2(\R^d)$, których wszystkie pochodne słabe
$\partial^\alpha u$, $|\alpha|\le m$, należą do $L^2$.

\begin{uwaga}[Szeregi Fouriera na torusie a transformata Fouriera na $\R^d$]
Rozdział~\ref{chap:pde-fourier} pracował na torusie lub
na odcinku po okresowym przedłużeniu i korzystał z rozwinięć
w szeregi Fouriera, czyli dyskretnego spektrum częstości.
Tutaj pracujemy na całej przestrzeni $\R^d$, więc częstość
$\xi$ zmienia się ciągle i pojawia się całka po $\R^d$.
Obie teorie są zgodne pojęciowo, ale różnią się nośnikiem,
normalizacją oraz tym, że w części okresowej integrujemy na
$\T^d$ względem ustalonej miary, a tutaj na $\R^d$ względem
miary Lebesgue'a.
\end{uwaga}

\begin{twierdzenie}[Podstawowe własności $H^s$]
\label{thm:pde8-Hs}
Dla wszystkich $s,t\in\R$, $s\ge t$, zachodzą:
\begin{enumerate}
\item $H^s(\R^d)$ jest przestrzenią Hilberta, a
 $J^s:H^s\to L^2$ jest izometrią na całą przestrzeń $L^2$.
\item $H^s\hookrightarrow H^t$ ciągle i
 $\|u\|_{H^t}\le\|u\|_{H^s}$.
\item Dla $\alpha\in\N_0^d$, $\partial^\alpha:H^s\to
 H^{s-|\alpha|}$ jest ograniczony.
\item Dla $m\in\N_0$ zachodzi $H^m(\R^d)=W^{m,2}(\R^d)$
 ze wzajemnie równoważnymi normami.
\end{enumerate}
\end{twierdzenie}
\begin{proof}
Pierwsza teza wynika z utożsamienia $H^s$ poprzez
$J^s$ z $L^2$ i odwracalności mnożnika
$(1+|\xi|^2)^{s/2}$ na $\mathcal S'$. Pozostałe
oszacowania opierają się na nierównościach
$(1+|\xi|^2)^t\le(1+|\xi|^2)^s$ i
$|\xi^\alpha|^2\le(1+|\xi|^2)^{|\alpha|}$.
Gdy $m$ jest całkowite nieujemne, wielomian
$(1+|\xi|^2)^m$ jest równoważny
$\sum_{|\alpha|\le m}|\xi^\alpha|^2$,
ze stałymi zależnymi tylko od $m,d$. Dzięki
Plancherelowi warunek $u\in H^m$ jest równoważny
istnieniu wszystkich pochodnych dystrybucyjnych
$\partial^\alpha u\in L^2$ dla $|\alpha|\le m$.
W przypadku $m=0$ warunek ten mówi po prostu $u\in L^2$.
\end{proof}

\begin{wniosek}[Kiedy delta należy do ujemnej przestrzeni Sobolewa?]
\label{cor:pde8-delta-Hs}
Dla $s\in\R$ zachodzi
\[
 \delta_0\in H^s(\R^d)\quad\Longleftrightarrow\quad
 s<-\frac d2.
\]
\end{wniosek}
\begin{proof}
Ponieważ $\widehat{\delta_0}=1$, warunkiem jest
$\int_{\R^d}(1+|\xi|^2)^s\,\mathrm d\lambda_d(\xi)<\infty$.
Całka po kuli jednostkowej jest skończona. Poza nią,
przechodząc do współrzędnych biegunowych, porównujemy
całkę z $\int_1^\infty r^{2s+d-1}\,\mathrm d\lambda_1(r)$,
która jest skończona dokładnie wtedy, gdy $2s+d<0$.
\end{proof}

\section{Zastosowanie II: eliptyczne równanie z parametrem}
\label{sec:pde8-elliptic}

\begin{twierdzenie}[Równanie Yukawy--Helmholtza]
\label{thm:pde8-helmholtz}
Niech $a>0$ i $f\in L^2(\R^d)$. Istnieje dokładnie jedno
$u\in H^2(\R^d)$ takie, że $(-\Delta+a^2)u=f$
w sensie dystrybucji. Zachodzi
\[
 \widehat u(\xi)=\frac{\widehat f(\xi)}{|\xi|^2+a^2},
 \qquad \|u\|_{H^2}\le C_a\|f\|_{L^2},
\]
gdzie $H^2$ zdefiniowano w~\ref{def:pde8-Hs},
a $C_a=\max\{1,a^{-2}\}$ dla przyjętej tam normy.
\end{twierdzenie}
\begin{proof}
Funkcja $(|\xi|^2+a^2)^{-1}$ jest ograniczona.
Ponadto $(1+|\xi|^2)/(|\xi|^2+a^2)\le\max\{1,a^{-2}\}$.
Zatem zadany wzór wyznacza $u\in H^2$ i podane
oszacowanie przez twierdzenie Plancherela. Po
pomnożeniu przez $|\xi|^2+a^2$ otrzymujemy równanie.
Jeżeli $(-\Delta+a^2)u=0$ dla $u\in H^2$, to
$(|\xi|^2+a^2)\widehat u=0$ prawie wszędzie,
a więc $u=0$. To daje jednoznaczność.
\end{proof}

\begin{uwaga}[Dlaczego przypadek $a=0$ jest inny?]
Dla czystego równania Poissona mnożnik formalny
$|\xi|^{-2}$ ma osobliwość w zerze. Nie można więc
bez dodatkowej analizy definiować dla dowolnego
$f\in L^2$ rozwiązania $u\in H^2$ przez ten wzór.
Dokładniej, dla $1\le d\le4$ i $f\in\mathcal S(\R^d)$
z $\widehat f(0)\ne0$ iloraz
$|\xi|^{-2}\widehat f(\xi)$ nie jest funkcją $L^2$
w pobliżu zera, gdyż jego kwadrat daje po przejściu
do współrzędnych biegunowych całkę porównywalną z
$\int_0^1r^{d-5}\,\mathrm d\lambda_1(r)=\infty$.
W wymiarach $d\ge5$ ta konkretna przeszkoda
niskoczęstotliwościowa znika dla takiego gładkiego $f$;
nie zmienia to faktu, że dzielenie przez $|\xi|^2$
nie jest ograniczonym mnożnikiem $L^2\to H^2$
bez dodatkowych warunków na dane. Dodatkowo funkcje harmoniczne tworzą jądro
operatora $-\Delta$ w szerszych klasach rozwiązań.
Nie oznacza to braku rozwiązań dystrybucyjnych;
rozdział~\ref{chap:pde-distributions} konstruuje je
z użyciem funkcji fundamentalnej.
\end{uwaga}

\section{Zasada nieoznaczoności jako konsekwencja Plancherela}
\label{sec:pde8-uncertainty}

\begin{twierdzenie}[Nierówność Heisenberga w jednej współrzędnej]
\label{thm:pde8-uncertainty}
Jeśli $0\ne f\in\mathcal S(\R^d)$ i $1\le j\le d$, to
\[
 \|x_jf\|_{L^2}\,\|\xi_j\widehat f\|_{L^2}
 \ge\frac12\,\|f\|_{L^2}\,\|\widehat f\|_{L^2}.
\]
\end{twierdzenie}
\begin{proof}
Całkując pochodną $\partial_j(x_j|f|^2)$ i używając
zaniku funkcji Schwartza, dostajemy
$\|f\|_2^2=-2\operatorname{Re}\int_{\R^d}x_j\overline f\,
\partial_jf\,\mathrm d\lambda_d(x)$.
Nierówność Cauchy'ego--Schwarza daje
$\|x_jf\|_2\|\partial_jf\|_2\ge\tfrac12\|f\|_2^2$.
Z twierdzenia Plancherela i
$\widehat{\partial_jf}=i\xi_j\widehat f$ wynika
$\|\partial_jf\|_2=(2\pi)^{-d/2}\|\xi_j\widehat f\|_2$
oraz $\|f\|_2=(2\pi)^{-d/2}\|\widehat f\|_2$.
Po podstawieniu otrzymujemy tezę.
\end{proof}

\section{Podsumowanie i zakres dalszych zagadnień}
Transformata Fouriera jest automorfizmem przestrzeni Schwartza i
jej przestrzeni dualnej. Jest podstawowym narzędziem badania wielu
liniowych PDE o stałych współczynnikach, lecz sama nie zapewnia
rozwiązywalności ani jednoznaczności w dowolnej klasie funkcji;
nie rozwiązuje też automatycznie problemów brzegowych ani równań
ze zmiennymi współczynnikami. Mnożniki ograniczone działają na $L^2$,
lecz dla $L^p$ z $p\ne2$ potrzebna jest głębsza teoria.
W następnych rozdziałach zbudujemy systematycznie
przestrzenie Sobolewa na obszarach i ich twierdzenia o zanurzeniach.
W ten sposób domkniemy łańcuch pojęciowy rozpoczęty w
rozdziale~\ref{chap:pde-distributions}: od słabych pochodnych,
przez narzędzia Fouriera na całej przestrzeni, do funkcjonałowej
teorii $W^{k,p}(\Omega)$ potrzebnej dla metod wariacyjnych i
zagadnień brzegowych.



\chapter[Przestrzenie Sobolewa]{Przestrzenie Sobolewa: słabe pochodne, aproksymacja i przedłużanie}
\label{chap:pde-sobolev}

Funkcja $u(x)=|x|$ nie jest różniczkowalna w zerze, lecz jej pochodna
w sensie dystrybucji jest funkcją $\operatorname{sgn}x$. Z drugiej strony
pochodna funkcji skokowej jest deltą Diraca, a nie funkcją lokalnie
całkowalną. Ta różnica motywuje przestrzenie Sobolewa: dopuszczamy
niegładkie funkcje, ale wymagamy, aby ich pochodne dystrybucyjne
były reprezentowane przez funkcje z określonych klas całkowalności.

Niech w całym rozdziale $d,k\in\N$, $\Omega\subset\R^d$ będzie
niepustym zbiorem otwartym, $1\le p\le\infty$; chyba że wyraźnie
zaznaczono inaczej, funkcje są rzeczywiste. Pochodne i równości funkcji
rozumiemy prawie wszędzie. Wszystkie całki na $\Omega$ są liczone
względem miary Lebesgue'a $\lambda_d$. Definicje dystrybucji i
pochodnych słabych znajdują się w rozdziale~\ref{chap:pde-distributions},
a przestrzenie $H^s(\R^d)$ w~\ref{sec:pde8-bessel}.

\section{Definicje i pierwszy słownik}
\label{sec:pde9-def}

\begin{definicja}[Przestrzenie $W^{k,p}$]
\label{def:pde9-Wkp}
Dla $k\in\N_0$ i $1\le p\le\infty$ definiujemy
\[
 W^{k,p}(\Omega)=\{u\in L^p(\Omega):
  \partial^\alpha u\in L^p(\Omega)
  \text{ dla każdego }\alpha\in\N_0^d,\ |\alpha|\le k\},
\]
gdzie $\partial^\alpha u$ oznacza pochodną dystrybucyjną.
Wyrażenie $\partial^\alpha u\in L^p$ oznacza, że dystrybucja ta jest
reprezentowana przez funkcję z $L^p$. Normę określamy wzorami
\begin{align*}
 \|u\|_{W^{k,p}(\Omega)}
  &=\left(\sum_{|\alpha|\le k}\|\partial^\alpha u\|_{L^p(\Omega)}^p
    \right)^{1/p},&&1\le p<\infty,\\
 \|u\|_{W^{k,\infty}(\Omega)}
  &=\max_{|\alpha|\le k}\|\partial^\alpha u\|_{L^\infty(\Omega)}.
\end{align*}
Dla $k=0$ otrzymujemy $W^{0,p}=L^p$. Piszemy $H^k(\Omega)=W^{k,2}(\Omega)$,
a w szczególności $H^1(\Omega)=W^{1,2}(\Omega)$.
\end{definicja}

\begin{definicja}[Przestrzeń lokalna i zerowy brzeg]
\label{def:pde9-local-zero}
Niech $k\in\N_0$, $1\le p\le\infty$.
Funkcja $u$ należy do $W^{k,p}_{\mathrm{loc}}(\Omega)$, jeśli
$u\in W^{k,p}(U)$ dla każdego otwartego $U$ takiego, że
$\overline U\subset\Omega$ i $\overline U$ jest zwarta.
Dla $1\le p<\infty$ definiujemy
\[
 W^{k,p}_0(\Omega)=\overline{C_c^\infty(\Omega)}^{\,W^{k,p}(\Omega)}.
\]
Indeks $0$ oznacza w tej chwili \emph{definicję przez domknięcie},
a nie jeszcze twierdzenie o istnieniu śladu na brzegu.
\end{definicja}

\begin{uwaga}[Wybór normy]
Dla skończenie wielu pochodnych różne zwyczajowe normy, np.
$\sum_{|\alpha|\le k}\|\partial^\alpha u\|_p$, są równoważne.
Przy $p=2$ norma $H^1(\Omega)$ spełnia
$\|u\|_{H^1}^2=\|u\|_2^2+\|\nabla u\|_2^2$.
Zgodnie z twierdzeniem~\ref{thm:pde8-Hs}, na całej $\R^d$
przestrzenie $H^k$ zdefiniowane przez słabe pochodne pokrywają się
z przestrzeniami Bessela z rozdziału~\ref{chap:pde-fourier-Rd};
normy są równoważne, a dla $k=1$ są równe.
\end{uwaga}

\begin{przyklad}[Załamanie nie jest skokiem]
Niech $I=(-1,1)$. Funkcja $u(x)=|x|$ należy do $W^{1,p}(I)$
dla każdego $1\le p\le\infty$, ponieważ
$u'=\operatorname{sgn}x$ prawie wszędzie i pochodna ta jest
jednocześnie pochodną dystrybucyjną. Natomiast funkcja
$H=\chi_{(0,1)}$ nie należy do $W^{1,p}(I)$ dla żadnego
$p\ge1$: jej pochodna dystrybucyjna to $\delta_0$, która nie jest
reprezentowana przez funkcję z $L^p(I)$.
\end{przyklad}

\begin{przyklad}[Osobliwość potęgowa]
\label{ex:pde9-power}
Niech $d\ge2$, $0<\alpha<d$, a $u(x)=\chi(x)|x|^{-\alpha}$,
gdzie $\chi\in C_c^\infty(\R^d)$ jest równa $1$ w pewnym otoczeniu
zera. Dla $1\le p<\infty$ zachodzi
\[
 u\in L^p(\R^d)\iff \alpha p<d,
 \qquad u\in W^{1,p}(\R^d)\iff (\alpha+1)p<d.
\]
W drugim warunku zakładamy, że $\alpha>0$.
W istocie w pobliżu zera $|\nabla u|\asymp |x|^{-\alpha-1}$.
Jeśli $(\alpha+1)p<d$, całkowanie przez części poza kulą
$B_\varepsilon$ daje po przejściu do granicy słabą pochodną:
wyraz brzegowy jest $O(\varepsilon^{d-1-\alpha})$ i zanika,
gdyż $\alpha<d-1$. Gdy warunek zawodzi, klasyczna pochodna
na $\R^d\setminus\{0\}$ nie należy do $L^p$ w pobliżu zera,
więc nie może istnieć pochodna słaba w $L^p$.
\end{przyklad}

\section{Zupełność i zamkniętość operatorów pochodnych}
\label{sec:pde9-complete}

\begin{twierdzenie}[Przestrzeń Sobolewa jest Banacha]
\label{thm:pde9-banach}
Dla każdego otwartego $\Omega\subset\R^d$, $k\in\N_0$ i
$1\le p\le\infty$ przestrzeń $W^{k,p}(\Omega)$ jest przestrzenią
Banacha. W szczególności $H^k(\Omega)$ jest przestrzenią Hilberta
z iloczynem skalarnym
$\langle u,v\rangle_{H^k}=\sum_{|\alpha|\le k}
\int_\Omega\partial^\alpha u\,\overline{\partial^\alpha v}
\,\mathrm d\lambda_d$.
\end{twierdzenie}
\begin{proof}
Niech $(u_n)$ będzie ciągiem Cauchy'ego w $W^{k,p}(\Omega)$.
Dla każdego wielowskaźnika $|\alpha|\le k$ ciąg
$(\partial^\alpha u_n)$ jest Cauchy'ego w $L^p$, zatem istnieje
$f_\alpha\in L^p(\Omega)$ takie, że
$\partial^\alpha u_n\to f_\alpha$ w $L^p$. Oznaczmy $u=f_0$.
Dla dowolnego $\varphi\in C_c^\infty(\Omega)$ mamy
\[
 \int_\Omega u_n\partial^\alpha\varphi\,\mathrm d\lambda_d
 =(-1)^{|\alpha|}\int_\Omega
 (\partial^\alpha u_n)\varphi\,\mathrm d\lambda_d.
\]
Obie całki przechodzą do granicy: dla $p<\infty$ wynika to z
nierówności Höldera (dla $p=1$ funkcje testowe są ograniczone),
a dla $p=\infty$ z faktu, że testy i ich pochodne są w $L^1$.
Otrzymujemy $\partial^\alpha u=f_\alpha$ w $\mathcal D'(\Omega)$.
A zatem $u\in W^{k,p}$, a zbieżność wszystkich pochodnych
w $L^p$ daje $u_n\to u$ w normie Sobolewa.
Dla $p=2$ wskazany iloczyn skalarny jest dodatnio określony,
a uzyskana zupełność dowodzi własności Hilberta.
\end{proof}

\begin{wniosek}[Zamkniętość pochodnej]
\label{cor:pde9-closed}
Jeśli $u_n\to u$ w $L^p(\Omega)$ oraz
$\partial^\alpha u_n\to g$ w $L^p(\Omega)$, gdzie
$u_n\in W^{|\alpha|,p}(\Omega)$, to $\partial^\alpha u=g$
w sensie dystrybucji. Nie zakładamy tu zbieżności pozostałych
pochodnych tego samego rzędu.
\end{wniosek}
\begin{proof}
Wystarczy przetestować definicję pochodnej dystrybucyjnej,
jak w dowodzie twierdzenia~\ref{thm:pde9-banach}.
\end{proof}

\begin{uwaga}[Dlaczego nie należy oczekiwać podobnego wyniku w $C^0$?]
Norma $L^p$ kontroluje funkcję jedynie w sensie całkowym;
nic nie wymusza istnienia wartości punktowych.
Dopiero twierdzenia o zanurzeniach Sobolewa, omówione w następnym
rozdziale, wskażą warunki, w których klasy $W^{k,p}$ mają
reprezentantów ciągłych.
\end{uwaga}

\section{Wymiar jeden: absolutna ciągłość}
\label{sec:pde9-oneD}

\begin{twierdzenie}[Charakterystyka $W^{1,p}(a,b)$]
\label{thm:pde9-ac}
Niech $-\infty<a<b<\infty$ i $1\le p\le\infty$.
Klasa $u\in L^p(a,b)$ należy do $W^{1,p}(a,b)$ wtedy i tylko
wtedy, gdy ma reprezentanta absolutnie ciągłego
$\widetilde u:[a,b]\to\R$, którego pochodna klasyczna
istnieje prawie wszędzie i należy do $L^p(a,b)$.
Dla takiego reprezentanta
\[
 \widetilde u(y)-\widetilde u(x)
 =\int_x^y u'(t)\,\mathrm d\lambda_1(t)
 \qquad(a\le x\le y\le b).
\]
\end{twierdzenie}
\begin{proof}
Załóżmy $u\in W^{1,p}(a,b)$ i niech $g$ będzie jego
słabą pochodną. Ponieważ $(a,b)$ jest ograniczony,
$g\in L^1(a,b)$. Ustalmy $x_0\in(a,b)$ i zdefiniujmy
$U(x)=\int_{x_0}^xg(t)\,\mathrm d\lambda_1(t)$.
Funkcja $U$ jest absolutnie ciągła i $U'=g$ prawie wszędzie.
Dystrybucyjna pochodna $u-U$ jest zerowa. Z twierdzenia
o dystrybucjach o zerowej pochodnej
(\ref{chap:pde-distributions}) wynika, że $u-U$ jest
prawie wszędzie stała. Wobec tego $U+C$ jest żądanym
reprezentantem. Odwrotnie, dla funkcji absolutnie ciągłej
z pochodną w $L^p$ zwykłe całkowanie przez części przeciw
funkcji testowej pokazuje, że pochodna klasyczna jest również
pochodną słabą.
\end{proof}

\begin{wniosek}[Oszacowanie punktowe i ślady w jednym wymiarze]
\label{cor:pde9-trace-oneD}
Dla $I=(a,b)$, $L=b-a$ i $u\in W^{1,p}(I)$, $1\le p<\infty$,
jedyny ciągły reprezentant spełnia
\[
 \|u\|_{L^\infty(I)}\le
 L^{-1/p}\|u\|_{L^p(I)}+L^{1-1/p}\|u'\|_{L^p(I)}.
\]
Dla $p=\infty$ zachodzi $\|u\|_\infty\le\|u\|_{W^{1,\infty}}$.
W szczególności odwzorowania $u\mapsto u(a)$ i $u\mapsto u(b)$
są ciągłymi funkcjonałami liniowymi na $W^{1,p}(I)$.
\end{wniosek}
\begin{proof}
Istnieje $x_0\in(a,b)$ takie, że
$|u(x_0)|\le L^{-1/p}\|u\|_p$ (w przeciwnym przypadku
całka z $|u|^p$ byłaby zbyt duża). Zastosowanie wzoru
z twierdzenia~\ref{thm:pde9-ac} oraz nierówności Höldera
daje dla wszystkich $x\in[a,b]$ podane oszacowanie.
Przy $p=\infty$ jest ono oczywiste z definicji normy.
\end{proof}

\begin{przyklad}[Funkcja Cantora a pochodna słaba]
Klasyczna funkcja Cantora jest ciągła, niemalejąca i ma
pochodną klasyczną równą zeru prawie wszędzie, mimo że nie jest
stała. Nie należy do $W^{1,1}(0,1)$, gdyż nie jest absolutnie
ciągła: jej pochodna \emph{dystrybucyjna} jest niezerową
miarą singularną. Zatem pochodna istniejąca prawie wszędzie
nie musi być pochodną słabą.
\end{przyklad}

\section{Sploty i aproksymacja wewnętrzna}
\label{sec:pde9-mollify}

Wybierzmy nieujemną funkcję $\rho\in C_c^\infty(B_1(0))$,
$\int_{\R^d}\rho\,\mathrm d\lambda_d=1$, i niech
$\rho_\varepsilon(x)=\varepsilon^{-d}\rho(x/\varepsilon)$.
Dla zwartego $K\subset\Omega$ zapisujemy
$\operatorname{dist}(K,\partial\Omega)>0$; jeśli $\Omega=\R^d$,
przyjmujemy, że odległość ta jest nieskończona.

\begin{twierdzenie}[Lokalne wygładzanie Sobolewa]
\label{thm:pde9-mollify}
Niech $1\le p<\infty$, $k\in\N_0$ oraz $u\in W^{k,p}_{\mathrm{loc}}(\Omega)$.
Jeżeli $K\subset\Omega$ jest zwarty, to dla dostatecznie
małych $\varepsilon>0$ splot
$u_\varepsilon=\rho_\varepsilon*u$ jest gładki w otoczeniu $K$
i zachodzi
\[
 \|u_\varepsilon-u\|_{W^{k,p}(K)}\longrightarrow0
 \qquad(\varepsilon\downarrow0),
\]
gdzie norma na $K$ oznacza sumę norm $L^p(K)$ pochodnych.
Ponadto na $K$ obowiązuje
$\partial^\alpha u_\varepsilon
 =\rho_\varepsilon*(\partial^\alpha u)$ dla $|\alpha|\le k$.
\end{twierdzenie}
\begin{proof}
Wybierzmy zbiory otwarte $K\subset U$, $\overline U\subset V$,
$\overline V\subset\Omega$, przy czym $\overline V$ jest zwarta,
oraz $\chi\in C_c^\infty(V)$ równą $1$ na $U$.
Po przedłużeniu $\chi u$ przez zero poza $V$ dostajemy
funkcję z $W^{k,p}(\R^d)$: reguła Leibniza wynika z
rachunku dystrybucyjnego, ponieważ $\chi$ jest gładka.
Gdy $\varepsilon$ jest mniejsze od odległości $K$ od
$\Omega\setminus U$, sploty $\rho_\varepsilon*u$ i
$\rho_\varepsilon*(\chi u)$ zgadzają się na $K$.
Dla każdej pochodnej $|\alpha|\le k$ rachunek dystrybucyjny
daje
$\partial^\alpha(\rho_\varepsilon*(\chi u))
 =\rho_\varepsilon*\partial^\alpha(\chi u)$.
Przybliżone jedynki są zbieżne w $L^p(\R^d)$ dla
$1\le p<\infty$ na mocy
rozdziału~\ref{chap:pde-distributions}. Ograniczając normy do $K$
i korzystając z $\chi=1$ w otoczeniu $K$, otrzymujemy tezę.
\end{proof}

\begin{uwaga}[Dlaczego nie $L^\infty$?]
Konkluzja normowa na ogół zawodzi dla $p=\infty$:
pochodna $u(x)=|x|$ na $(-1,1)$ jest funkcją znak, której
ciągłe wygładzenia nie zbliżają się w normie $L^\infty$
w otoczeniu zera. Nadal otrzymujemy lokalną zbieżność
w $W^{k,q}$ dla każdego skończonego $q$ na zbiorach zwartych,
gdy $u\in W^{k,\infty}_{\mathrm{loc}}$.
\end{uwaga}

\begin{center}
\begin{tikzpicture}[scale=0.92]
 \draw[thick,rounded corners=12pt,fill=blue!5,draw=blue!55]
  (-4.1,-1.45) -- (3.9,-1.45) -- (4.2,1.25) -- (-3.7,1.40) -- cycle;
 \draw[thick,rounded corners=10pt,fill=green!9,draw=green!48!black]
  (-2.55,-.85) rectangle (2.65,.85);
 \draw[thick,rounded corners=9pt,fill=orange!11,draw=orange!75!black]
  (-1.35,-.52) rectangle (1.25,.52);
 \node at (-3.42,0.98) {$\Omega$};
 \node at (2.1,.59) {$U$};
 \node at (0,0) {$K$};
 \draw[-{Stealth[length=2.3mm]},thick] (0.85,-.10)--(2.1,-.10)
  node[midway,below] {$\varepsilon$};
\end{tikzpicture}

{\small\emph{Lokalność wygładzania.} Zwarty $K$ jest oddalony od
brzegu obszaru; promień splotu wybieramy mniejszy od odległości
do obszaru, na którym odcinamy funkcję.}
\end{center}

\section{Twierdzenie Meyersa--Serrina i gęstość funkcji gładkich}
\label{sec:pde9-density}

\begin{twierdzenie}[Meyers--Serrin]
\label{thm:pde9-MeyersSerrin}
Niech $\Omega\subset\R^d$ będzie otwarta, $k\in\N_0$
i $1\le p<\infty$. Wówczas $C^\infty(\Omega)\cap W^{k,p}(\Omega)$
jest gęsta w $W^{k,p}(\Omega)$.
Nie zakłada się regularności brzegu $\partial\Omega$.
\end{twierdzenie}
\begin{proof}
Wybierzmy lokalnie skończone pokrycie $\Omega$ otwartymi
zbiorami $U_j$ o zwartych domknięciach wewnątrz $\Omega$ oraz
gładki podział jedności $\{\psi_j\}$ z $\supp\psi_j\subset U_j$.
Można wybrać dodatkowo lokalnie skończone zbiory otwarte
$V_j$ takie, że $\supp\psi_j\subset V_j$, $\overline{V_j}\subset\Omega$.
Dla $u\in W^{k,p}(\Omega)$ każda funkcja $u_j=\psi_j u$
ma zwarty nośnik w $\Omega$ i należy do $W^{k,p}(\Omega)$.
Z twierdzenia~\ref{thm:pde9-mollify} oraz argumentu o nośniku
możemy wybrać $\varepsilon_j>0$ tak małe, by
$g_j=\rho_{\varepsilon_j}*u_j$ należało do
$C_c^\infty(V_j)$ i
$\|g_j-u_j\|_{W^{k,p}(\Omega)}<\varepsilon2^{-j}$.
Suma $g=\sum_{j=1}^\infty g_j$ jest lokalnie skończona,
a więc należy do $C^\infty(\Omega)$. Ponadto szereg
$\sum_j(g_j-u_j)$ jest bezwzględnie zbieżny w przestrzeni
Banacha $W^{k,p}(\Omega)$. Ponieważ
$\sum_j u_j=u$ w sensie dystrybucyjnym i punktowo
prawie wszędzie, w tej przestrzeni mamy
$g-u=\sum_j(g_j-u_j)$. Zatem
$\|g-u\|_{W^{k,p}}\le\sum_j\varepsilon2^{-j}\le\varepsilon$.
Użyliśmy jedynie standardowej konstrukcji lokalnie
skończonego podziału jedności dla zbioru otwartego w $\R^d$,
nie wymagającej żadnych założeń o jego brzegu.
\end{proof}

\begin{uwaga}[Gładkość nie oznacza zerowania na brzegu]
Twierdzenie~\ref{thm:pde9-MeyersSerrin} \emph{nie} głosi,
że $C_c^\infty(\Omega)$ jest gęsta w $W^{k,p}(\Omega)$.
W przeciwieństwie do $C^\infty(\Omega)$ funkcje o zwartym
nośniku wewnątrz obszaru muszą zanikać przy brzegu.
\end{uwaga}

\begin{przyklad}[Stała na przedziale nie należy do $W^{1,p}_0$]
\label{ex:pde9-constant}
Niech $\Omega=(0,1)$, $1\le p<\infty$, oraz $u\equiv1$.
Funkcja $u\in W^{1,p}(0,1)$, ale
$u\notin W^{1,p}_0(0,1)$. Gdyby istniał ciąg
$\varphi_n\in C_c^\infty(0,1)$ zbieżny do $u$
w $W^{1,p}$, to z oszacowania
\ref{cor:pde9-trace-oneD} wynikałoby
$\varphi_n(0)\to u(0)=1$. Jest to niemożliwe, bo
$\varphi_n(0)=0$ dla wszystkich $n$.
\end{przyklad}

\section[Przestrzeń Sobolewa z zerowym śladem i przedłużenie przez zero]{Przestrzeń $W^{k,p}_0$ i przedłużenie przez zero}
\label{sec:pde9-zero}

\begin{twierdzenie}[Przedłużenie przez zero]
\label{thm:pde9-zero-extension}
Niech $\Omega\subset\R^d$ będzie otwarta, $k\in\N_0$,
$1\le p<\infty$ i $u\in W^{k,p}_0(\Omega)$. Funkcja
\[
 \widetilde u(x)=\begin{cases}u(x),&x\in\Omega,\\0,&x\notin\Omega\end{cases}
\]
należy do $W^{k,p}(\R^d)$, a
$\|\widetilde u\|_{W^{k,p}(\R^d)}=\|u\|_{W^{k,p}(\Omega)}$.
Wszystkie pochodne słabe $\partial^\alpha\widetilde u$
są przedłużeniami przez zero pochodnych $\partial^\alpha u$.
\end{twierdzenie}
\begin{proof}
Z definicji $W^{k,p}_0$ istnieją
$\varphi_n\in C_c^\infty(\Omega)$ z
$\varphi_n\to u$ w $W^{k,p}(\Omega)$.
Ich przedłużenia przez zero do $\R^d$ są gładkie i mają
normy Sobolewa równe normom na $\Omega$.
Tworzą zatem ciąg Cauchy'ego w $W^{k,p}(\R^d)$.
Niech $v$ będzie jego granicą. Zbieżność w $L^p(\R^d)$
pokazuje, że $v=u$ prawie wszędzie w $\Omega$ i $v=0$
prawie wszędzie poza $\Omega$. Zbieżność pochodnych
identyfikuje je z przedłużeniami przez zero pochodnych $u$.
Równość norm wynika przez przejście do granicy.
\end{proof}

\begin{uwaga}[Nie każdą funkcję z przestrzeni Sobolewa można tak przedłużyć]
Dla $u\equiv1$ na $(0,1)$ przedłużenie przez zero ma skoki
w punktach $0,1$, więc jego pochodną dystrybucyjną
jest $\delta_0-\delta_1$, a nie funkcja z $L^p(\R)$.
Zatem założenia w twierdzeniu~\ref{thm:pde9-zero-extension}
nie wolno zastępować samym $u\in W^{1,p}(\Omega)$.
\end{uwaga}

\section{Nierówność Poincarégo na kostce}
\label{sec:pde9-poincare}

\begin{twierdzenie}[Poincaré--Wirtinger na kostce]
\label{thm:pde9-cube-Poincare}
Niech $Q=\prod_{j=1}^d(a_j,a_j+L)$, $L>0$,
$1\le p<\infty$. Dla $u\in W^{1,p}(Q)$ połóżmy
$u_Q=\lambda_d(Q)^{-1}\int_Qu\,\mathrm d\lambda_d$.
Wówczas
\[
 \|u-u_Q\|_{L^p(Q)}\le L\sum_{j=1}^d
 \|\partial_j u\|_{L^p(Q)}.
\]
W szczególności prawa strona znika dokładnie wtedy,
gdy $u$ jest prawie wszędzie stała na $Q$.
\end{twierdzenie}
\begin{proof}
Najpierw niech $u$ będzie gładka w otoczeniu $\overline Q$.
W jednym wymiarze dla $h\in C^1([a,a+L])$
\[
 \left|h(x)-\frac1L\int_a^{a+L}h(t)\,\mathrm d\lambda_1(t)\right|
 \le\int_a^{a+L}|h'(t)|\,\mathrm d\lambda_1(t).
\]
Biorąc normę $L^p$ względem $x$ i stosując Höldera,
otrzymujemy oszacowanie przez $L\|h'\|_{L^p(a,a+L)}$.
Niech $A_j$ oznacza operator uśrednienia względem pierwszych
$j$ współrzędnych, przy czym $A_0u=u$ oraz $A_du=u_Q$.
Z nierówności jednowymiarowej i twierdzenia Fubiniego
wynika $\|A_{j-1}u-A_ju\|_p\le L\|\partial_j A_{j-1}u\|_p$.
Różniczkowanie w $x_j$ komutuje z $A_{j-1}$,
a $A_{j-1}$ jest kontrakcją w $L^p(Q)$ na mocy Jensena.
Zatem $\|A_{j-1}u-A_ju\|_p\le L\|\partial_j u\|_p$.
Sumując po $j$ i stosując nierówność trójkąta,
otrzymujemy tezę dla gładkich $u$.
Dla ogólnego $u\in W^{1,p}(Q)$ rozważamy kostki wewnętrzne
$Q_\delta=\prod_j(a_j+\delta,a_j+L-\delta)$, $0<\delta<L/2$.
Na każdej $Q_\delta$ twierdzenie o lokalnych splotach
pozwala przybliżyć $u$ w $W^{1,p}(Q_\delta)$ przez funkcje
gładkie w otoczeniu $\overline{Q_\delta}$.
Wynik stosujemy do $Q_\delta$ z długością krawędzi
$L-2\delta$. Gdy $\delta\downarrow0$, całki i średnie
po $Q_\delta$ zbiegają do odpowiednich wyrażeń po $Q$;
ponieważ $u,\nabla u\in L^p(Q)$, zbieżność w normach $L^p$
wynika z absolutnej ciągłości całki. Otrzymujemy tezę z
nie większą stałą $L$.
Jeśli $\nabla u=0$, to $u=u_Q$ prawie wszędzie.
\end{proof}

\begin{wniosek}[Wersja z zerowym brzegiem]
\label{cor:pde9-Poincare-zero}
Dla $u\in W^{1,p}_0(Q)$ na kostce z poprzedniego twierdzenia
zachodzi $\|u\|_{L^p(Q)}\le L\|\partial_1u\|_{L^p(Q)}$.
\end{wniosek}
\begin{proof}
Dla $u\in C_c^\infty(Q)$ ustalamy pozostałe współrzędne
i korzystamy ze wzoru
$u(x)=\int_{a_1}^{x_1}\partial_1u(t,x_2,\ldots,x_d)
\,\mathrm d\lambda_1(t)$. Hölder i Fubini dają oszacowanie
po scałkowaniu po $Q$. Przenosimy je na domknięcie
$C_c^\infty(Q)$ w $W^{1,p}$.
\end{proof}

\section{Reguły rachunkowe w przestrzeniach Sobolewa}
\label{sec:pde9-calculus}

\begin{twierdzenie}[Mnożenie przez funkcję gładką i iloczyn]
\label{thm:pde9-products}
Niech $\Omega\subset\R^d$ będzie otwarta, $1\le p<\infty$
i $u\in W^{1,p}(\Omega)$.
Jeśli $\phi\in C^1(\Omega)$ oraz $\phi,\nabla\phi\in L^\infty(\Omega)$,
to $\phi u\in W^{1,p}(\Omega)$ i
$\partial_j(\phi u)=\phi\partial_j u+u\partial_j\phi$.
Ponadto jeśli $u\in W^{1,p}(\Omega)$ oraz
$v\in W^{1,q}(\Omega)$, gdzie $1/p+1/q=1$,
to $uv\in W^{1,1}(\Omega)$ i
$\partial_j(uv)=v\partial_j u+u\partial_j v$.
\end{twierdzenie}
\begin{proof}
Pierwsza tożsamość wynika bezpośrednio z
$\partial_j(\phi u)$ w sensie dystrybucyjnym; oba składniki
prawej strony są w $L^p$.
Dla iloczynu wybieramy lokalne wygładzenia $u_n,v_n$.
Na każdym zwartym podzbiorze $\Omega$ odpowiednie
pochodne są zbieżne w $L^p$ i $L^q$; nierówność Höldera
zapewnia w $L^1$ zbieżność $u_nv_n\to uv$ i
$v_n\partial_j u_n+u_n\partial_jv_n
\to v\partial_j u+u\partial_jv$.
Wzór dla gładkich funkcji przechodzi do granicy
na testach o zwartym nośniku. Globalnie wszystkie
wyrazy po prawej stronie oraz $uv$ należą do $L^1(\Omega)$
na mocy Höldera. Dla $p=1,q=\infty$ stosujemy
lokalne wygładzenia i lokalną zbieżność w $L^1$
oraz jednostajne ograniczenie funkcji $v_n,\nabla v_n$;
aby otrzymać zbieżność iloczynów wystarczy również
zbieżność prawie wszędzie po przejściu do podciągu,
a następnie zbieżność zdominowana na nośniku testu.
\end{proof}

\begin{twierdzenie}[Reguła łańcuchowa dla funkcji Lipschitzowskich]
\label{thm:pde9-chain}
Niech $1\le p<\infty$, $u\in W^{1,p}(\Omega)$ i
$F\in C^1(\R)$ z $F(0)=0$ oraz $\|F'\|_\infty<\infty$.
Wtedy $F\circ u\in W^{1,p}(\Omega)$ i
\[
 \partial_j(F(u))=F'(u)\,\partial_j u
 \quad\text{prawie wszędzie w }\Omega.
\]
Jeżeli $\Omega$ ma skończoną miarę, warunek $F(0)=0$
można pominąć.
\end{twierdzenie}
\begin{proof}
Z nierówności $|F(t)-F(0)|\le\|F'\|_\infty|t|$
wynika $F(u)\in L^p$.
Weźmy lokalne wygładzenia $u_n\to u$ w $W^{1,p}$.
Po przejściu do podciągu $u_n\to u$ i
$\nabla u_n\to\nabla u$ prawie wszędzie.
Wówczas $F(u_n)\to F(u)$ w $L^p$ z oszacowania
Lipschitzowskiego. Dla pochodnych rozkładamy
\[
 F'(u_n)\partial_j u_n-F'(u)\partial_j u
 =F'(u_n)(\partial_j u_n-\partial_j u)
 +(F'(u_n)-F'(u))\partial_j u.
\]
Pierwszy składnik zbiega do zera w $L^p$ z ograniczoności
$F'$, drugi z twierdzenia o zbieżności zdominowanej.
Przechodząc do granicy dystrybucyjnie, uzyskujemy tezę.
\end{proof}

\begin{lemat}[Gradient znika na zbiorach poziomicowych]
\label{lem:pde9-levelsets}
Jeśli $u\in W^{1,p}(\Omega)$, $1\le p<\infty$,
i $c\in\R$, to $\nabla u=0$ prawie wszędzie na
$\{x\in\Omega:u(x)=c\}$.
\end{lemat}
\begin{proof}
Wybierzmy ciąg funkcji $h_n\in C^\infty(\R)$ taki, że
$0\le h_n\le1$, $h_n(t)=0$ dla $|t-c|\le1/n$
oraz $h_n(t)=1$ dla $|t-c|\ge2/n$.
Niech $F_n(t)=\int_0^t h_n(s)\,\mathrm d\lambda_1(s)$.
Wtedy $F_n(0)=0$, $\|F_n'\|_\infty\le1$
oraz $|F_n(t)-t|\le4/n$ dla każdego $t\in\R$.
Ponadto $|F_n(t)|\le|t|$, więc z twierdzenia o
zbieżności zdominowanej otrzymujemy
$F_n(u)\to u$ w $L^p(\Omega)$, także gdy $\Omega$
ma nieskończoną miarę. Reguła łańcuchowa daje
$\nabla(F_n(u))=h_n(u)\nabla u$.
Ponieważ $h_n(u)\to\chi_{\{u\ne c\}}$ punktowo,
a $|h_n(u)\nabla u|\le|\nabla u|$, znów zbieżność
zdominowana zapewnia zbieżność w $L^p$ do
$\chi_{\{u\ne c\}}\nabla u$.
Korzystając z zamkniętości pochodnej słabej
(wniosek~\ref{cor:pde9-closed}), dostajemy
$\nabla u=\chi_{\{u\ne c\}}\nabla u$,
czyli tezę.
\end{proof}

\begin{wniosek}[Obcięcia i części dodatnie]
\label{cor:pde9-truncation}
Niech $u\in W^{1,p}(\Omega)$, $1\le p<\infty$.
Wtedy $u^+=\max(u,0)$ i $u^-=\max(-u,0)$ należą do
$W^{1,p}(\Omega)$, a prawie wszędzie
\[
 \nabla u^+=\chi_{\{u>0\}}\nabla u,
 \qquad \nabla u^-=-\chi_{\{u<0\}}\nabla u.
\]
Dla $M>0$ obcięcie $T_M(u)=\max(-M,\min(u,M))$
również należy do $W^{1,p}(\Omega)$.
\end{wniosek}
\begin{proof}
Przybliżamy funkcję $t\mapsto t^+$ przez gładkie
funkcje Lipschitzowskie $F_n$ z $F_n(0)=0$
i $0\le F'_n\le1$, które zbiegają punktowo do
typowej pochodnej części dodatniej poza $t=0$.
Poprzednie twierdzenie daje regułę łańcuchową.
Lemat~\ref{lem:pde9-levelsets} zapewnia, że
$\nabla u=0$ prawie wszędzie na $\{u=0\}$.
Pozwala to przejść do granicy również na zbiorze poziomicowym.
Dla $u^-$ i $T_M$ postępujemy analogicznie.
\end{proof}

\section{Przedłużanie z półprzestrzeni i rola brzegu}
\label{sec:pde9-extension}

\begin{twierdzenie}[Odbicie Sobolewa, rząd pierwszy]
\label{thm:pde9-reflection}
Niech $\R^d_+=\{(x',t)\in\R^{d-1}\times\R:t>0\}$,
$1\le p<\infty$. Istnieje liniowy, ograniczony operator
$E:W^{1,p}(\R^d_+)\to W^{1,p}(\R^d)$ taki, że
$(Eu)|_{\R^d_+}=u$. Można przyjąć
\[
 (Eu)(x',t)=\begin{cases}u(x',t),&t>0,\\
 u(x',-t),&t<0.\end{cases}
\]
Wówczas $\|Eu\|_{W^{1,p}(\R^d)}=2^{1/p}
\|u\|_{W^{1,p}(\R^d_+)}$.
\end{twierdzenie}
\begin{proof}
Najpierw dla funkcji gładkich aż do brzegu i
mających odpowiednio całkowalne pochodne zwykłe
różniczkowanie odbicia daje
$\partial_t Eu(x',t)=\operatorname{sgn}(t)
(\partial_tu)(x',|t|)$, a pochodne styczne odbijają się
bez zmiany znaku. Ponieważ wartości z dwóch stron
$t=0$ są takie same, nie pojawia się składnik delta.
W konsekwencji $\|Eu\|_{W^{1,p}}=2^{1/p}\|u\|_{W^{1,p}}$.

Dla ogólnego $u\in W^{1,p}(\R^d_+)$ aproksymujemy
funkcjami gładkimi do brzegu: przesuwamy najpierw $u$
w kierunku dodatnim $t$, czyli $u_h(x',t)=u(x',t+h)$,
$h>0$, co zbiega w $W^{1,p}(\R^d_+)$ przy $h\downarrow0$.
Funkcja $u_h$ jest określona także dla $t>-h$;
można ją więc wygładzić splotem z jądrem o promieniu
mniejszym niż $h/2$, uzyskując gładkie aproksymacje
do brzegu. Równość norm dla odbicia pokazuje, że ich
przedłużenia tworzą ciąg Cauchy'ego w $W^{1,p}(\R^d)$.
Jego granica jest przedłużeniem $u$, a równość norm
przechodzi do granicy.
\end{proof}

\begin{stwierdzenie}[Odbicia wyższego rzędu na półprzestrzeni]
\label{prop:pde9-highreflection}
Niech $k\in\N$, $1\le p<\infty$. Istnieje ograniczony,
liniowy operator $E_k:W^{k,p}(\R^d_+)\to W^{k,p}(\R^d)$,
który jest prawym odwrotnym do ograniczenia na $\R^d_+$.
\end{stwierdzenie}
\begin{proof}
Wybieramy liczby $a_1,\ldots,a_k$ spełniające
$\sum_{j=1}^k a_j(-j)^m=1$ dla $m=0,\ldots,k-1$.
Ten układ ma dokładnie jedno rozwiązanie, ponieważ jego
macierz jest macierzą Vandermonde'a o różnych węzłach
$-1,\ldots,-k$. Definiujemy $E_ku=u$ dla $t>0$, a dla $t<0$
\[
 (E_ku)(x',t)=\sum_{j=1}^k a_j u(x',-jt).
\]
Dla gładkiego $u$ ślady wszystkich pochodnych normalnych
rzędu $m<k$ po obu stronach $t=0$ są zgodne właśnie dzięki
układowi na $a_j$. Nie powstają więc dystrybucyjne
składniki singularne w pochodnych rzędu co najwyżej $k$.
Normy tych pochodnych po dolnej stronie granicy są
ograniczone przez skończoną sumę wyrażeń
$|a_j|j^{m-1/p}\|\partial_t^m\partial_{x'}^\beta u\|_p$,
przy $m+|\beta|\le k$, po zmianie zmiennej $s=-jt$.
Ponownie aproksymacja przez przesunięcie w głąb
i wygładzanie pozwalają przenieść oszacowanie na wszystkie
$u\in W^{k,p}(\R^d_+)$.
\end{proof}

\begin{uwaga}[Obszary Lipschitzowskie]
Jeśli ograniczony obszar $\Omega$ ma brzeg lokalnie
będący wykresem funkcji Lipschitzowskiej, można prostować
brzeg dwulipschitzowskimi zmianami współrzędnych,
zastosować odbicie pierwszego rzędu i skleić lokalne
przedłużenia podziałem jedności. Otrzymujemy ograniczony
operator $W^{1,p}(\Omega)\to W^{1,p}(\R^d)$.
Przypadek $W^{k,p}$ dla $k>1$ na dowolnym obszarze Lipschitzowskim wymaga bardziej subtelnej konstrukcji:
proste prostowanie Lipschitzowskie nie zachowuje ogólnie
drugich pochodnych w $L^p$. Pełną teorię przedłużania
i śladów rozwiniemy w następnym rozdziale.
\end{uwaga}

\begin{lemat}[Gęstość funkcji gładkich aż do płaskiego brzegu]
\label{lem:pde9-density-halfspace-slab}
Niech $1\le p<\infty$ oraz $d\ge1$. Dla
$G=\R^{d-1}\times(0,\infty)$ albo
$G=\R^{d-1}\times(0,1)$ każdą funkcję $u\in W^{1,p}(G)$
można przybliżyć w normie $W^{1,p}(G)$ restrykcjami funkcji
z $C^\infty(\R^d)$. Można przy tym wybrać przybliżenia
z $C_c^\infty(\R^d)$.
\end{lemat}
\begin{proof}
Dla półprzestrzeni używamy ograniczonego przedłużenia
przez odbicie z twierdzenia~\ref{thm:pde9-reflection}.
Dla warstwy przedłużamy najpierw przez odbicie w obu
płaszczyznach brzegowych: na $-1<t<0$ przyjmujemy
$u(x',-t)$, na $0<t<1$ funkcję $u(x',t)$, a na
$1<t<2$ funkcję $u(x',2-t)$. Zastosowanie argumentu
z dowodu twierdzenia~\ref{thm:pde9-reflection} osobno
w pobliżu $t=0$ i $t=1$ pokazuje, że sklejenie
należy do $W^{1,p}(\R^{d-1}\times(-1,2))$:
na prawie każdym pionowym odcinku funkcja ma reprezentanta
absolutnie ciągłego (twierdzenie~\ref{thm:pde9-ac}),
więc wartości z obu stron płaszczyzn odbicia są zgodne
i nie powstają dystrybucyjne składniki brzegowe.
Normy funkcji i pochodnych na częściach odbitych
równe odpowiednim normom na warstwie wyjściowej.
Mnożymy otrzymaną funkcję przez gładkie odcięcie
$\eta(t)$ równe $1$ na $[0,1]$ i o nośniku w $(-1,2)$,
a następnie przedłużamy wynik przez zero do $\R^d$.
Powstaje ograniczone liniowe przedłużenie
$E:W^{1,p}(G)\to W^{1,p}(\R^d)$ również dla warstwy.

Niech $v=Eu$. Najpierw mnożymy $v$ przez odcięcie
$\chi_R\in C_c^\infty(\R^d)$, równe $1$ na $B_R$,
o $|\nabla\chi_R|\le C/R$ i nośniku w $B_{2R}$.
Z nierówności Leibniza i zbieżności ogonów całek $L^p$
wynika $\chi_Rv\to v$ w $W^{1,p}(\R^d)$.
Dla ustalonego $R$ splatamy $\chi_Rv$ z przybliżoną
jedynką $\rho_\varepsilon$. Zbieżność splotów dla
funkcji i ich słabych pochodnych w $L^p$ daje
$\rho_\varepsilon*(\chi_Rv)\to\chi_Rv$
w $W^{1,p}(\R^d)$. Otrzymane funkcje są w
$C_c^\infty(\R^d)$, a ograniczenie ich do $G$
nie zwiększa normy. Wybierając kolejno $R\to\infty$
i $\varepsilon\downarrow0$, otrzymujemy tezę.
\end{proof}


\begin{uwaga}[Aproksymacja gładka nie oznacza zerowego śladu]
\label{rem:pde-v42-density}
Twierdzenie Meyersa--Serrina~\ref{thm:pde9-MeyersSerrin}
nie orzeka, że $C_c^\infty(\Omega)$ jest gęste w całej
przestrzeni $W^{1,p}(\Omega)$. Na przykład funkcja
$u\equiv1$ na ograniczonym obszarze Lipschitzowskim
należy do $W^{1,p}(\Omega)$, ale jej ślad jest równy $1$,
a więc nie należy do $W_0^{1,p}(\Omega)$ na mocy
twierdzenia~\ref{thm:pde10-trace-zero}.
Są to trzy różne twierdzenia: aproksymacja wewnętrzna,
gęstość gładkich restrykcji aż do brzegu oraz
charakterystyka jądra śladu. Ten ostatni rezultat
zostanie udowodniony w rozdziale~\ref{chap:pde-embeddings}.
\end{uwaga}

\section{Ślady na płaskim brzegu: pierwszy wynik}
\label{sec:pde9-traces}

\begin{twierdzenie}[Oszacowanie śladu na warstwie]
\label{thm:pde9-trace-slab}
Niech $S=\R^{d-1}\times(0,1)$ i $1\le p<\infty$.
Istnieje stała $C_p>0$ i jednoznaczny ograniczony
operator liniowy
$\operatorname{Tr}_0:W^{1,p}(S)\to L^p(\R^{d-1})$
taki, że dla $u\in C^\infty(\overline S)\cap W^{1,p}(S)$
z regularnym śladem jest nim
$\operatorname{Tr}_0u(x')=u(x',0)$. Zachodzi
\[
 \|\operatorname{Tr}_0u\|_{L^p(\R^{d-1})}
 \le C_p\bigl(\|u\|_{L^p(S)}+\|\partial_tu\|_{L^p(S)}\bigr).
\]
\end{twierdzenie}
\begin{proof}
Dla funkcji gładkiej $h:[0,1]\to\R$ i $s\in(0,1)$
mamy $|h(0)|\le|h(s)|+\int_0^1|h'(t)|\,
\mathrm d\lambda_1(t)$. Całkując po $s$ i używając Höldera,
otrzymujemy $|h(0)|\le\|h\|_{L^p(0,1)}+
\|h'\|_{L^p(0,1)}$, stąd
$|h(0)|^p\le2^{p-1}(\|h\|_p^p+\|h'\|_p^p)$.
Stosujemy to do $h(t)=u(x',t)$ i całkujemy po $x'$.
Daje to tezę dla funkcji gładkich. Gęstość restrykcji funkcji gładkich aż do obu brzegów
w $W^{1,p}(S)$ wynika z osobno udowodnionego
lematu~\ref{lem:pde9-density-halfspace-slab}. Oszacowanie pozwala
przedłużyć operator śladu przez ciągłość na całą
przestrzeń. Jednoznaczność jest konsekwencją gęstości.
\end{proof}

\section{Perspektywa wariacyjna: pierwsze słabe rozwiązanie}
\label{sec:pde9-weakPDE}

\begin{przyklad}[Równanie Poissona na przedziale]
\label{ex:pde9-dirichlet}
Niech $f\in L^2(0,1)$ i rozważmy
$-u''=f$ w $(0,1)$ z warunkami $u(0)=u(1)=0$.
Jeśli $u$ jest gładkie, mnożymy równanie przez
$\varphi\in C_c^\infty(0,1)$ i całkujemy przez części:
\[
 \int_0^1u'(x)\varphi'(x)\,\mathrm d\lambda_1(x)
 =\int_0^1 f(x)\varphi(x)\,\mathrm d\lambda_1(x).
\]
Równanie to zachowuje sens również wtedy, gdy szukamy
$u\in H_0^1(0,1)$: wystarcza żądać tożsamości dla każdego
$\varphi\in H_0^1(0,1)$. Jest to \emph{słabe zagadnienie Dirichleta}.

Dla $f\equiv1$ rozwiązaniem jest $u(x)=x(1-x)/2$.
Dla dowolnego $f\in L^2(0,1)$ można też podać jawny wzór
\[
 u(x)=\int_0^1G(x,y)f(y)\,\mathrm d\lambda_1(y),\qquad
 G(x,y)=\begin{cases}y(1-x),&y\le x,\\x(1-y),&x<y.\end{cases}
\]
Ponieważ $G(0,y)=G(1,y)=0$ i
$-\partial_x^2G(\cdot,y)=\delta_y$ w $\mathcal D'(0,1)$,
forma jest rozwiązaniem dystrybucyjnym.
Dokładniej, dla każdego $f\in L^2$ można dwukrotnie
całkować równanie w jednym wymiarze i uzyskać
$u\in H^2(0,1)\cap H_0^1(0,1)$, a zatem także tożsamość słabą.
Jeśli $f\equiv0$, wstawiając do niej $\varphi=u$,
otrzymujemy $\int_0^1|u'|^2\,\mathrm d\lambda_1=0$.
Twierdzenie~\ref{thm:pde9-ac} i zerowe wartości brzegowe
pokazują $u=0$, co dowodzi jednoznaczności w klasie $H_0^1$.
\end{przyklad}

\section{Skalowanie, krytyczne wykładniki i zadania na przyszłość}
\label{sec:pde9-scaling}

Niech $\phi\in C_c^\infty(\R^d)$, $\phi\not\equiv0$,
oraz dla $\varepsilon>0$ zdefiniujmy
$u_\varepsilon(x)=\varepsilon^{-a}\phi(x/\varepsilon)$.
Bezpośrednia zmiana zmiennych pokazuje
\[
 \|u_\varepsilon\|_{L^q}
 =\varepsilon^{d/q-a}\|\phi\|_{L^q},\qquad
 \|\nabla u_\varepsilon\|_{L^p}
 =\varepsilon^{d/p-a-1}\|\nabla\phi\|_{L^p}
\]
dla $1\le p,q<\infty$. Jeżeli chcemy, aby nierówność
$\|u\|_{L^q}\le C\|\nabla u\|_{L^p}$ na całej $\R^d$
była niezmiennicza względem skalowania, musimy mieć
$d/q=d/p-1$. Dla $1\le p<d$ dostajemy wykładnik
\[
 p^*=\frac{dp}{d-p}.
\]
Nie dowodzi to jeszcze zanurzenia Sobolewa
$W^{1,p}\hookrightarrow L^{p^*}$, ale wyjaśnia,
dlaczego dokładnie ten wykładnik pojawia się w twierdzeniu.
Konstrukcje skupiające masę w małej kuli będą również
źródłem kontrprzykładów dla zbyt silnych zanurzeń.

\section{Podsumowanie i zależności logiczne}
Wprowadziliśmy przestrzenie $W^{k,p}$ na dowolnych obszarach,
wykazaliśmy ich zupełność, lokalne wygładzanie i
teorię Meyersa--Serrina oraz przedłużenia z półprzestrzeni.
Oddzieliliśmy $W^{k,p}(\Omega)$ od $W^{k,p}_0(\Omega)$,
który koduje zerowy warunek brzegowy, a także udowodniliśmy
proste nierówności Poincarégo i śladu. W rozdziale 10
wykorzystamy te narzędzia do ogólnych twierdzeń o
zanurzeniach, zwartości Rellicha--Kondrachova i
pełniejszej teorii śladów na brzegach Lipschitzowskich.


% Rozdział 10. Samodzielna część wykładu; włączana przed dodatkiem z kodem.
\chapter[Zanurzenia i zwartość Sobolewa]{Zanurzenia Sobolewa, nierówność Morreya, zwartość i ślady}
\label{chap:pde-embeddings}

Przestrzeń $W^{1,p}(\Omega)$ została zdefiniowana przez całkowalność
funkcji i jej słabych pochodnych. Nie wynika z tej definicji wprost,
czy funkcja jest ciągła, czy można ją przybliżać jednostajnie
ani czy ciąg ograniczony w $W^{1,p}$ ma podciąg zbieżny w $L^p$.
Odpowiedzi zależą przede wszystkim od relacji $p$ do wymiaru $d$.
Będziemy odróżniać \emph{zanurzenie ciągłe}, oznaczane
$X\hookrightarrow Y$, od \emph{zanurzenia zwartego}, oznaczanego
$X\hookrightarrow\hookrightarrow Y$. Pierwsze znaczy, że
$\|u\|_Y\le C\|u\|_X$; drugie, że każdy ciąg ograniczony w $X$
ma podciąg zbieżny w normie $Y$.

W rozdziale $d\in\N$ i $1\le p\le\infty$. Nie zakładamy automatycznie,
że $\Omega$ ma regularny brzeg. Kiedy korzystamy z obszaru Lipschitzowskiego, podajemy to wyraźnie. Całki na $\R^d$ i na
$\Omega$ są całkami względem $\lambda_d$; na brzegu pojawia się
miara powierzchniowa $\mathcal H^{d-1}$.

\section{Nierówność geometryczna potrzebna do dowodu Sobolewa}
\label{sec:pde10-LW}

Zaczniemy od rachunkowego lematu, który pozwala udowodnić
nierówność Sobolewa bez odwoływania się do twierdzenia
izoperymetrycznego ani do teorii potencjałów Riesza.
Dla $x=(x_1,\ldots,x_d)$ przez $\widehat x_j$ oznaczamy
wektor powstały przez opuszczenie współrzędnej $x_j$.

\begin{lemat}[Nierówność Loomisa--Whitneya w postaci całkowej]
\label{lem:pde10-LW}
Niech $d\ge2$ i $F_j:\R^{d-1}\to[0,\infty]$ będą mierzalne,
$j=1,\ldots,d$. Wówczas
\[
 \int_{\R^d}\prod_{j=1}^d F_j(\widehat x_j)^{1/(d-1)}
 \,\mathrm d\lambda_d(x)
 \le \prod_{j=1}^d
 \left(\int_{\R^{d-1}}F_j\,\mathrm d\lambda_{d-1}\right)^{1/(d-1)}.
\]
Jeżeli niektóre całki po prawej są nieskończone, nierówność
rozumiemy w naturalnym sensie po zastosowaniu obcięć;
nie posługujemy się nieokreślonym iloczynem $0\cdot\infty$.
\end{lemat}
\begin{proof}
Wystarczy najpierw rozważyć nieujemne funkcje ograniczone
o zwartych nośnikach. Dla $d=2$ nierówność jest równością:
$F_1$ zależy tylko od $x_2$, a $F_2$ tylko od $x_1$;
twierdzenie Fubiniego kończy dowód.

Załóżmy indukcyjnie tezę w wymiarze $d-1$, gdzie $d\ge3$.
Całkujemy w lewej stronie po $x_d$. Czynnik
$F_d(x_1,\ldots,x_{d-1})$ nie zależy od $x_d$.
Nierówność H\"oldera zastosowana do pozostałych $d-1$
czynników, każdy z wykładnikiem $d-1$, daje
\[
 \int_{\R}\prod_{j<d}F_j(\widehat x_j)^{1/(d-1)}
 \,\mathrm d\lambda_1(x_d)
 \le\prod_{j<d}G_j(\widehat x^{\prime}_j)^{1/(d-1)},
\]
gdzie $x'=(x_1,\ldots,x_{d-1})$, a $G_j$ powstaje przez
scałkowanie $F_j$ względem $x_d$; nie zależy więc od $x_j$.
Oznaczmy $H(x')=\prod_{j<d}G_j(\widehat x'_j)^{1/(d-2)}$.
Zapiszmy dla skrócenia wzorów
$I_d=\int_{\R^{d-1}}F_d(x')\,\mathrm d\lambda_{d-1}(x')$,
$J=\int_{\R^{d-1}}H(x')\,\mathrm d\lambda_{d-1}(x')$
oraz, dla $j<d$,
$I_j=\int_{\R^{d-2}}G_j(z)\,\mathrm d\lambda_{d-2}(z)$.
Po użyciu H\"oldera na $\R^{d-1}$ z wykładnikami
$d-1$ oraz $(d-1)/(d-2)$ otrzymujemy
\begin{align*}
 &\int_{\R^{d-1}}F_d(x')^{1/(d-1)}
 \prod_{j<d}G_j(\widehat x'_j)^{1/(d-1)}
 \,\mathrm d\lambda_{d-1}(x')\\
 &\qquad\le I_d^{1/(d-1)}J^{(d-2)/(d-1)}
 \le I_d^{1/(d-1)}\prod_{j<d} I_j^{1/(d-1)}.
\end{align*}
Druga nierówność jest dokładnie założeniem indukcyjnym.
Na mocy twierdzenia Fubiniego
$I_j=\int_{\R^{d-1}}F_j(z)\,\mathrm d\lambda_{d-1}(z)$,
co kończy dowód. Dla dowolnych nieujemnych $F_j$ stosujemy
monotoniczną aproksymację funkcjami ograniczonymi o nośnikach
o skończonej mierze i twierdzenie o zbieżności monotonicznej.
\end{proof}

\begin{twierdzenie}[Nierówność Sobolewa dla $p=1$]
\label{thm:pde10-Sob1}
Dla $d\ge2$ i każdej $u\in C_c^1(\R^d)$ zachodzi
\[
 \|u\|_{L^{d/(d-1)}(\R^d)}
 \le\prod_{j=1}^d\|\partial_j u\|_{L^1(\R^d)}^{1/d}
 \le \frac1d\sum_{j=1}^d\|\partial_j u\|_{L^1(\R^d)}.
\]
W szczególności $\|u\|_{L^{d/(d-1)}}\le C_d\|\nabla u\|_{L^1}$.
\end{twierdzenie}
\begin{proof}
Dla ustalonego $j$ i każdego $x\in\R^d$, całkując
$\partial_j u$ wzdłuż prostej równoległej do osi $x_j$,
uzyskujemy
\[
 |u(x)|\le A_j(\widehat x_j),\qquad
 A_j(\widehat x_j)=\int_{\R}
 |\partial_j u(x_1,\ldots,t,\ldots,x_d)|\,\mathrm d\lambda_1(t).
\]
Mnożąc $d$ nierówności i podnosząc do potęgi $1/(d-1)$,
dostajemy $|u(x)|^{d/(d-1)}\le\prod_j
A_j(\widehat x_j)^{1/(d-1)}$. Lemat~\ref{lem:pde10-LW}
daje zatem
\[
 \|u\|_{d/(d-1)}^{d/(d-1)}
 \le \prod_{j=1}^d\|\partial_j u\|_1^{1/(d-1)}.
\]
Po podniesieniu do potęgi $(d-1)/d$ otrzymujemy pierwszą
nierówność. Druga jest nierównością pomiędzy średnią
geometryczną i arytmetyczną. Wszystkie całki są skończone
z uwagi na zwarty nośnik $u$.
\end{proof}

\section{Ostre wykładniki zanurzenia Sobolewa}
\label{sec:pde10-Sobolev}

\begin{twierdzenie}[Nierówność Sobolewa na całej przestrzeni]
\label{thm:pde10-Sobp}
Niech $d\ge2$ i $1\le p<d$. Przyjmijmy
$p^*=dp/(d-p)$. Istnieje stała $C(d,p)>0$ taka, że
\[
 \|u\|_{L^{p^*}(\R^d)}
 \le C(d,p)\|\nabla u\|_{L^p(\R^d)}
 \quad\text{dla każdego }u\in C_c^\infty(\R^d).
\]
W konsekwencji identyczne oszacowanie obowiązuje
(z ewentualnie inną stałą) dla $u\in W^{1,p}(\R^d)$.
\end{twierdzenie}
\begin{proof}
Dla $p=1$ jest to twierdzenie~\ref{thm:pde10-Sob1}.
Niech $1<p<d$ i $u\in C_c^\infty(\R^d)$. Wybierzmy
$\gamma=p(d-1)/(d-p)>1$. Stosujemy twierdzenie
\ref{thm:pde10-Sob1} do funkcji
$v=|u|^\gamma$; aby uniknąć niepotrzebnego założenia
o gładkości w zerach $u$, można najpierw zastąpić
$|u|^\gamma$ gładką funkcją
$(u^2+\delta^2)^{\gamma/2}-\delta^\gamma$
i przejść do granicy przy $\delta\downarrow0$.
Reguła łańcuchowa oraz nierówność H\"oldera dają
\begin{align*}
 \|u\|_{L^{\gamma d/(d-1)}}^\gamma
 &\le C_d\gamma\int_{\R^d}|u|^{\gamma-1}|\nabla u|
 \,\mathrm d\lambda_d\\
 &\le C_d\gamma
 \|u\|_{L^{(\gamma-1)p'}}^{\gamma-1}
 \|\nabla u\|_{L^p},\qquad p'=\frac p{p-1}.
\end{align*}
Z wyboru $\gamma$ wynika jednocześnie
$\gamma d/(d-1)=p^*$ i $(\gamma-1)p'=p^*$.
Jeśli $u\not\equiv0$, dzielimy obie strony przez
$\|u\|_{p^*}^{\gamma-1}$; przypadek $u=0$ jest oczywisty.
To daje nierówność dla funkcji gładkich.

Jeśli $u\in W^{1,p}(\R^d)$, wybieramy
$u_n\in C_c^\infty(\R^d)$ zbieżne do $u$ w $W^{1,p}$.
Istnienie takiego ciągu wynika z odcięcia $u$ na coraz
większych kulach, a następnie ze splotów; przy $p<\infty$
oba błędy znikają (reguła Leibniza i zbieżność w $L^p$).
Z udowodnionego oszacowania wynika, że $(u_n)$ jest
ciągiem Cauchy'ego również w $L^{p^*}$. Jego granica
w $L^{p^*}$ zgadza się prawie wszędzie z $u$, gdyż
$u_n\to u$ jednocześnie w $L^p$ i lokalnie w mierze.
Przechodzimy do granicy w nierówności.
\end{proof}

\begin{uwaga}[Dlaczego wykładnik jest krytyczny?]
\label{rem:pde10-critical}
Dla $\varphi\in C_c^\infty(B_1(0))$, $\varphi\ne0$, oraz
$x_0\in\R^d$ niech
$u_\varepsilon(x)=\varepsilon^{1-d/p}
\varphi((x-x_0)/\varepsilon)$. Wówczas
\[
 \|\nabla u_\varepsilon\|_p=\|\nabla\varphi\|_p,
 \qquad \|u_\varepsilon\|_q
 =\varepsilon^{1-d/p+d/q}\|\varphi\|_q.
\]
Dla $q=p^*$ drugi wykładnik jest równy zeru;
dla $q>p^*$ jest ujemny. Nie może więc istnieć
oszacowanie $\|u\|_q\le C\|u\|_{W^{1,p}}$ dla
wszystkich takich funkcji, nawet na dowolnym obszarze zawierającym kulę wokół $x_0$. Ten sam przykład pokaże
brak zwartości w wykładniku krytycznym.
\end{uwaga}

\begin{definicja}[Obszar Lipschitzowski]
\label{def:pde10-Lipschitz}
Ograniczony otwarty zbiór $\Omega\subset\R^d$ nazywamy
\emph{obszarem Lipschitzowskim}, jeżeli jest spójny, a jego
brzeg można pokryć skończoną liczbą otoczeń, w których
po odpowiednim sztywnym ruchu współrzędnych ma on postać
$\{(x',x_d):x_d=\phi(x')\}$ z $\phi$ Lipschitzowską,
zaś $\Omega$ leży lokalnie po jednej stronie tego wykresu.
Przyjmujemy, że otoczenia są mniejsze niż wykresowe
cylindry, aby dostępne były dodatnie odległości między
nośnikami funkcji odcinających a ich bocznymi ścianami.
\end{definicja}

\begin{lemat}[Przedłużanie pierwszego rzędu]
\label{lem:pde10-extension}
Niech $\Omega\subset\R^d$ będzie obszarem ograniczonym Lipschitzowskim i $1\le p\le\infty$. Istnieje liniowy operator
$E:W^{1,p}(\Omega)\to W^{1,p}(\R^d)$ oraz stała
$C=C(\Omega,p)>0$ takie, że $(Eu)|_\Omega=u$ i
$\|Eu\|_{W^{1,p}(\R^d)}\le C\|u\|_{W^{1,p}(\Omega)}$.
Można dodatkowo wybrać operator o obrazie w funkcjach
wspartych w jednej ustalonej kuli, niezależnej od $u$.
\end{lemat}
\begin{proof}
Przedstawimy konstrukcję i wyjaśnimy osobno przejście
od funkcji gładkich do Sobolewa. Dla pojedynczego wykresu
$x_d=\phi(x')$ prostujemy brzeg odwzorowaniem
$\Phi(x',t)=(x',t+\phi(x'))$. Funkcja $\Phi$ i jej odwrotność
są Lipschitzowskie; ich jakobian jest prawie wszędzie
równy $1$, a macierze pochodnych są jednostajnie ograniczone
przez stałą zależną od $\operatorname{Lip}\phi$.
Dla gładkiego $u$ reguła łańcuchowa pokazuje, że
$u\circ\Phi\in W^{1,p}$ na odpowiedniej półwarstwie,
a normy po zmianie zmiennych są kontrolowane przez normy
$u$ i $\nabla u$. Dla ogólnego $u$ stosujemy lokalne
przybliżenie gładkie z twierdzenia~\ref{thm:pde9-MeyersSerrin};
po wybraniu mniejszego cylindra oraz odcięcia otrzymujemy
ograniczony operator zmiany współrzędnych i jego odwrotność.
Dla $p=\infty$ tę samą własność uzyskujemy bez aproksymacji
w normie $W^{1,\infty}$: funkcje $W^{1,\infty}$ mają
lokalnie Lipschitzowskich reprezentantów, a ich złożenie
z odwzorowaniem dwulipschitzowskim zachowuje oszacowanie
na słaby gradient (można to także wykazać z twierdzenia
o absolutnej ciągłości na prawie wszystkich prostych).

Na wyprostowanej warstwie stosujemy parzyste odbicie
$w(x',t)\mapsto w(x',|t|)$, którego oszacowanie dla
$1\le p<\infty$ wykazano w twierdzeniu
\ref{thm:pde9-reflection}; dla $p=\infty$ identyczny rachunek
słabych pochodnych działa po lokalnej aproksymacji w
$W^{1,q}$ i przejściu do granicy słabo-* w $L^\infty$.
Mnożymy odbity wynik przez odcięcie o nośniku wewnątrz
większego cylindra i wracamy do oryginalnych współrzędnych.

Wybieramy skończony układ takich cylindrów oraz jeden
zbiór ściśle wewnętrzny, które pokrywają $\overline\Omega$.
Podział jedności $\chi_0+\cdots+\chi_N=1$ w otoczeniu
$\overline\Omega$ pozwala osobno przedłużyć funkcje
$\chi_j u$. Część wewnętrzna ma przedłużenie przez zero;
części brzegowe przedłużamy metodą opisaną wyżej.
Suma jest równa $u$ na $\Omega$, a ograniczoność liczby
składników i ich operatorów daje żądaną nierówność.
Końcowe pomnożenie przez funkcję równą $1$ na
$\overline\Omega$ ogranicza nośnik przedłużenia
bez naruszenia własności $(Eu)|_\Omega=u$.
\end{proof}

\begin{wniosek}[Gęstość funkcji gładkich aż do brzegu]
\label{cor:pde10-density-boundary}
Niech $\Omega\subset\R^d$ będzie ograniczonym obszarem
Lipschitzowskim, $1\le p<\infty$. Wówczas zbiór
$\{\varphi|_\Omega:\varphi\in C_c^\infty(\R^d)\}$
jest gęsty w $W^{1,p}(\Omega)$. W szczególności jest
gęsty zbiór funkcji gładkich w otoczeniu $\overline\Omega$.
\end{wniosek}
\begin{proof}
Dla $u\in W^{1,p}(\Omega)$ przedłużenie
$v=Eu\in W^{1,p}(\R^d)$ istnieje na mocy
lematu~\ref{lem:pde10-extension}. Wybierzmy funkcje
$\chi_R$ równe $1$ w $B_R$, o nośniku w $B_{2R}$
i $\|\nabla\chi_R\|_\infty\le C/R$.
Wtedy $\chi_Rv\to v$ w $W^{1,p}(\R^d)$:
dla składników bez pochodnych korzystamy ze zbieżności
ogonów całek, a dla gradientów dodatkowy składnik
$v\nabla\chi_R$ ma normę najwyżej $C\|v\|_p/R$.
Dla ustalonego $R$ regularizujemy $\chi_Rv$ splotem
z $\rho_\varepsilon$. Ponieważ $p<\infty$,
$\rho_\varepsilon*(\chi_Rv)\to\chi_Rv$ w $W^{1,p}(\R^d)$.
Funkcje te należą do $C_c^\infty(\R^d)$.
Ograniczenie do $\Omega$ nie zwiększa normy, więc
odpowiedni ciąg restrykcji zbiega do $u$ w $W^{1,p}(\Omega)$.
\end{proof}

\begin{twierdzenie}[Zanurzenia na obszarze ograniczonym Lipschitzowskim]
\label{thm:pde10-domain-embed}
Niech $d\ge2$ i $\Omega\subset\R^d$ będzie obszarem ograniczonym Lipschitzowskim. Wówczas zachodzą zanurzenia ciągłe:
\begin{enumerate}[label=\textnormal{(\roman*)}]
\item gdy $1\le p<d$, $W^{1,p}(\Omega)\hookrightarrow
 L^q(\Omega)$ dla $1\le q\le p^*=dp/(d-p)$;
\item gdy $p=d$, $W^{1,d}(\Omega)\hookrightarrow L^q(\Omega)$
dla każdego skończonego $q\ge1$;
\item gdy $p>d$, $W^{1,p}(\Omega)\hookrightarrow
 C^{0,\alpha}(\overline\Omega)$ dla $\alpha=1-d/p$.
Dla $p=\infty$ przyjmujemy $\alpha=1$.
\end{enumerate}
W części (iii) przestrzeń $C^{0,\alpha}(\overline\Omega)$
oznacza klasy z reprezentantem ciągłym o skończonej normie
$\|u\|_\infty+[u]_{C^{0,\alpha}}$.
\end{twierdzenie}
\begin{proof}
Część (i) wynika z lematu~\ref{lem:pde10-extension},
twierdzenia~\ref{thm:pde10-Sobp} i ograniczoności miary
$\Omega$: dla $q\le p^*$ nierówność H\"oldera daje
$\|u\|_q\le|\Omega|^{1/q-1/p^*}\|u\|_{p^*}$.

Gdy $p=d$, ustalmy dowolne skończone $q>d$.
Możemy dobrać $p_0\in(1,d)$ tak, że
$p_0^*=dp_0/(d-p_0)\ge q$. Przedłużenie $Eu$ ma
zwarty nośnik zawarty w ustalonej kuli $B_R$.
Ponieważ $p_0<d$, H\"older na $B_R$ daje
$\|Eu\|_{W^{1,p_0}}\le C\|Eu\|_{W^{1,d}}$.
Część (i), zastosowana na całej przestrzeni do $Eu$,
zapewnia $\|u\|_q\le C_q\|u\|_{W^{1,d}}$.
Dla $q\le d$ wystarcza H\"older na $\Omega$.
Część (iii) wynika z nierówności Morreya,
której pełny dowód podajemy w następnej sekcji.
\end{proof}

\section{Nierówność Morreya: dowód z całkowania po odcinkach}
\label{sec:pde10-Morrey}

\begin{lemat}[Oszacowanie różnicy punktu i średniej]
\label{lem:pde10-point-average}
Niech $d\ge2$, $B=B_R(a)\subset\R^d$, a $f\in C^1(\overline B)$.
Dla każdego $x\in B$ zachodzi
\[
 |f(x)-f_B|\le C_d\int_B
 \frac{|\nabla f(z)|}{|x-z|^{d-1}}\,\mathrm d\lambda_d(z),
 \quad f_B=\frac{1}{|B|}\int_Bf\,\mathrm d\lambda_d.
\]
Stała $C_d$ zależy tylko od wymiaru.
\end{lemat}
\begin{proof}
Z wypukłości kuli wynika, że odcinek od $x$ do każdego
$y\in B$ leży w $B$. Podstawowe twierdzenie rachunku
daje
$|f(x)-f(y)|\le|x-y|\int_0^1
|\nabla f(x+t(y-x))|\,\mathrm d\lambda_1(t)$.
Uśredniamy po $y\in B$ i zamieniamy całki nieujemne
na mocy Tonellego. Przy ustalonym $t$ podstawiamy
$z=x+t(y-x)$, zatem $\mathrm d\lambda_d(y)=
t^{-d}\,\mathrm d\lambda_d(z)$ i
$|x-y|=|x-z|/t$. Jeśli $y\in B$, to
$z\in B$ oraz $|x-z|\le2Rt$. W konsekwencji
\begin{align*}
 |f(x)-f_B|
 &\le\frac1{|B|}\int_B|\nabla f(z)|\,|x-z|
 \int_{|x-z|/(2R)}^1 t^{-d-1}\,\mathrm d\lambda_1(t)
 \,\mathrm d\lambda_d(z)\\
 &\le \frac{(2R)^d}{d|B|}
 \int_B\frac{|\nabla f(z)|}{|x-z|^{d-1}}
 \,\mathrm d\lambda_d(z).
\end{align*}
Ponieważ $|B|=\omega_dR^d$, współczynnik jest zależny
tylko od $d$. Ewentualny punkt $z=x$ nie wpływa na całkę.
\end{proof}

\begin{twierdzenie}[Morrey na $\R^d$]
\label{thm:pde10-Morrey}
Niech $d\ge2$ i $d<p<\infty$. Każda klasa
$u\in W^{1,p}(\R^d)$ ma jednoznacznego ciągłego
reprezentanta $\widetilde u$ i istnieje stała $C(d,p)$ taka, że
\[
 |\widetilde u(x)-\widetilde u(y)|
 \le C(d,p)|x-y|^{1-d/p}\|\nabla u\|_{L^p(\R^d)}
 \quad(x,y\in\R^d).
\]
Dodatkowo $\|\widetilde u\|_\infty\le
 C(d,p)\|u\|_{W^{1,p}(\R^d)}$.
Dla $p=\infty$ istnieje reprezentant globalnie Lipschitzowski
z $\operatorname{Lip}\widetilde u\le\|\nabla u\|_\infty$.
\end{twierdzenie}
\begin{proof}
Najpierw rozważmy $u\in C_c^\infty(\R^d)$. Ustalmy
$x\ne y$ i oznaczmy $r=|x-y|$.
Wybieramy kulę $B=B_{2r}(x)$, która zawiera oba punkty.
Z lematu~\ref{lem:pde10-point-average} otrzymujemy
$|u(x)-u(y)|\le|u(x)-u_B|+|u(y)-u_B|$.
Każdy składnik szacujemy nierównością H\"oldera.
Dla $z_0=x$ lub $z_0=y$ liczymy
\[
 \left(\int_B |z-z_0|^{-(d-1)p'}
 \,\mathrm d\lambda_d(z)\right)^{1/p'}
 \le C(d,p)r^{1-d/p}.
\]
Istotnie, $B\subset B_{4r}(z_0)$, a w biegunowych
całka jest ograniczona przez
$C_d\int_0^{4r}s^{d-1-(d-1)p'}\,\mathrm d\lambda_1(s)$.
Wykładnik przy zerze jest większy od $-1$ dokładnie
gdy $p>d$. Otrzymujemy oszacowanie H\"olderowskie.

Aby kontrolować wartość w jednym punkcie,
uśredniamy $u$ na kuli jednostkowej $B_1(x)$.
Z powyższego oszacowania przy $r\le2$ dostajemy
$|u(x)-u_{B_1(x)}|\le C\|\nabla u\|_p$,
a $|u_{B_1(x)}|\le |B_1|^{-1/p}\|u\|_p$.
To zapewnia oszacowanie $\|u\|_\infty$.

Dla ogólnego $u\in W^{1,p}(\R^d)$ wybieramy
gładkie funkcje o zwartym nośniku $u_n\to u$ w $W^{1,p}$.
Dwa udowodnione oszacowania zastosowane do $u_n-u_m$
pokazują, że $(u_n)$ jest ciągiem Cauchy'ego w normie
$C_b(\R^d)$ i w seminormie H\"oldera wykładnika
$1-d/p$. Granica jednostajna jest ciągła i równa $u$
prawie wszędzie, bo istnieje podciąg $u_n$ zbieżny do $u$
prawie wszędzie. Przechodzimy do granicy w oszacowaniu.

Dla $p=\infty$ dowód dla funkcji gładkich wynika z całkowania
$\nabla u$ po odcinku; przejście do funkcji $W^{1,\infty}$
można wykonać przez splot lokalny i granicę jednostajną
na kulach (oszacowanie Lipschitzowskie nie zależy od promienia
splotu), następnie przez wyczerpanie $\R^d$.
\end{proof}

\begin{wniosek}[Morrey na obszarze i wyższe rzędy]
\label{cor:pde10-Morrey-domain}
Jeśli $\Omega\subset\R^d$ jest obszarem ograniczonym Lipschitzowskim i $d<p\le\infty$, to
$W^{1,p}(\Omega)\hookrightarrow C^{0,1-d/p}(\overline\Omega)$,
z naturalną interpretacją wykładnika $1$ dla $p=\infty$.
Jeśli $k\ge2$ i $u\in W^{k,p}(\Omega)$, $d<p\le\infty$,
to pochodne słabe $\partial^\alpha u$ dla $|\alpha|\le k-1$
mają ciągłe reprezentanty H\"olderowskie wykładnika
$1-d/p$; stąd $u$ ma reprezentanta w $C^{k-1,1-d/p}$.
\end{wniosek}
\begin{proof}
Dla $k=1$ wystarczy ograniczony operator przedłużania
z lematu~\ref{lem:pde10-extension} oraz twierdzenie
\ref{thm:pde10-Morrey}. Dla $k\ge2$ każdy słaby
$\partial^\alpha u$, $|\alpha|\le k-1$, należy do
$W^{1,p}(\Omega)$, więc stosujemy przypadek $k=1$.
Aby utożsamić te ciągłe funkcje z pochodnymi klasycznymi,
można użyć lokalnej tożsamości całkowania wzdłuż prostych
dla przybliżeń gładkich i przejść do granicy jednostajnej.
\end{proof}

\begin{uwaga}[Wymiar $d=1$]
Na ograniczonym przedziale wyniki są silniejsze:
$W^{1,1}(a,b)$ ma reprezentant absolutnie ciągły,
a dla $p>1$ nierówność H\"oldera i twierdzenie
\ref{thm:pde9-ac} dają
$|u(x)-u(y)|\le |x-y|^{1-1/p}\|u'\|_{L^p(a,b)}$.
Nie wolno utożsamiać tego przypadku z uogólnieniem wzoru
$p^*=dp/(d-p)$ na $d=1$.
\end{uwaga}

\section{Zwartość: od średnich po kostkach do twierdzenia Rellicha}
\label{sec:pde10-Rellich}

\begin{lemat}[Aproksymacja funkcjami stałymi na kostkach]
\label{lem:pde10-cube-approx}
Niech $Q\subset\R^d$ będzie kostką o krawędzi $L$, a
$1\le p<\infty$. Dla każdego $N\in\N$ podzielmy $Q$
na $N^d$ równych kostek $Q_j$ o krawędzi $h=L/N$,
a dla $u\in L^1(Q)$ połóżmy
$P_Nu=\sum_j u_{Q_j}\chi_{Q_j}$, gdzie
$u_{Q_j}=|Q_j|^{-1}\int_{Q_j}u\,\mathrm d\lambda_d$.
Dla wszystkich $u\in W^{1,p}(Q)$ zachodzi
\[
 \|u-P_Nu\|_{L^p(Q)}\le dh\|\nabla u\|_{L^p(Q)}.
\]
\end{lemat}
\begin{proof}
Stosujemy nierówność Poincar\'ego--Wirtingera
z twierdzenia~\ref{thm:pde9-cube-Poincare} osobno na
każdej $Q_j$:
$\|u-u_{Q_j}\|_{L^p(Q_j)}\le
 h\sum_{i=1}^d\|\partial_i u\|_{L^p(Q_j)}$.
Ponieważ $(\sum_i a_i)^p\le d^{p-1}\sum_i a_i^p$
dla $a_i\ge0$, po podniesieniu do potęgi $p$
i zsumowaniu po rozłącznych kostkach otrzymujemy
$\|u-P_Nu\|_p^p\le d^{p-1}h^p\sum_i\|\partial_i u\|_p^p$.
Po wyciągnięciu pierwiastka mamy w istocie stałą
$d^{1-1/p}$; oszacowanie przez $d$ jest wystarczające.
\end{proof}

\begin{twierdzenie}[Rellich--Kondrachov: podstawowa zwartość]
\label{thm:pde10-Rellich-p}
Niech $\Omega\subset\R^d$ będzie obszarem ograniczonym Lipschitzowskim i $1\le p<\infty$. Wtedy zanurzenie
$W^{1,p}(\Omega)\hookrightarrow\hookrightarrow L^p(\Omega)$
jest zwarte.
\end{twierdzenie}
\begin{proof}
Niech $(u_n)$ będzie ograniczony w $W^{1,p}(\Omega)$.
Przedłużamy każdą funkcję operatorem
$E$ z lematu~\ref{lem:pde10-extension}. Wszystkie
$Eu_n$ mają nośniki zawarte w jednej ustalonej kostce $Q$
i są ograniczone w $W^{1,p}(Q)$ przez stałą $M$.

Ustalmy $\varepsilon>0$. Z lematu~\ref{lem:pde10-cube-approx}
wybieramy $N$ tak duże, aby
$\|Eu_n-P_NEu_n\|_{L^p(Q)}<\varepsilon/3$
jednocześnie dla wszystkich $n$.
Funkcje $P_NEu_n$ należą do skończeniewymiarowej przestrzeni
funkcji stałych na kostkach $Q_j$, a ich normy są
jednostajnie ograniczone przez $M$ (Jensen daje
$\|P_Nf\|_p\le\|f\|_p$). Zbiór ograniczony w przestrzeni
skończeniewymiarowej jest całkowicie ograniczony,
więc ma skończoną $\varepsilon/3$-sieć.
Z nierówności trójkąta wynika, że zbiór $\{Eu_n\}$
ma skończoną $\varepsilon$-sieć w $L^p(Q)$.

Dla pełnej ścisłości, wybierając kolejno sieci
dla $\varepsilon=2^{-m}$ i zagnieżdżone podciągi,
konstrukcja przekątniowa daje podciąg $(Eu_{n_k})$
Cauchy'ego w $L^p(Q)$; zupełność $L^p$
zapewnia jego zbieżność. Ograniczenie do $\Omega$
nie zwiększa normy, więc $u_{n_k}$ zbiega w
$L^p(\Omega)$. Nie wykorzystaliśmy
refleksywności $W^{1,p}$, dlatego dowód działa również
przy $p=1$.
\end{proof}

\begin{lemat}[Interpolacja elementarna norm $L^p$]
\label{lem:pde10-interpolation}
Niech $1\le a<q<b<\infty$, a $f\in L^a(\Omega)\cap L^b(\Omega)$.
Jeśli $0<\theta<1$ spełnia $1/q=\theta/a+(1-\theta)/b$,
to
$\|f\|_{L^q}\le\|f\|_{L^a}^\theta
\|f\|_{L^b}^{1-\theta}$.
\end{lemat}
\begin{proof}
Piszemy $|f|^q=(|f|^a)^{q\theta/a}
(|f|^b)^{q(1-\theta)/b}$.
Dwa wykładniki sumują się do $1$.
Oznaczmy $J_r=\int_\Omega|f(x)|^r\,\mathrm d\lambda_d(x)$
dla $r\in\{a,q,b\}$. Nierówność H\"oldera z wykładnikami
$a/(q\theta)$ i $b/(q(1-\theta))$ daje
\[
 J_q\le J_a^{q\theta/a}J_b^{q(1-\theta)/b},
\]
co kończy dowód.
\end{proof}

\begin{twierdzenie}[Rellich--Kondrachov: wykładniki podkrytyczne]
\label{thm:pde10-Rellich-q}
Niech $\Omega\subset\R^d$ będzie obszarem ograniczonym Lipschitzowskim, $d\ge2$.
\begin{enumerate}[label=\textnormal{(\roman*)}]
\item Jeśli $1\le p<d$, to
$W^{1,p}(\Omega)\hookrightarrow\hookrightarrow L^q(\Omega)$
dla $1\le q<p^*=dp/(d-p)$.
\item Jeśli $p=d$, to to samo zanurzenie jest zwarte dla
każdego skończonego $q\ge1$.
\item Jeśli $d<p\le\infty$, to
$W^{1,p}(\Omega)\hookrightarrow\hookrightarrow
C^{0,\beta}(\overline\Omega)$
dla każdego $0\le\beta<1-d/p$; przy $p=\infty$ wolno
przyjąć dowolne $\beta<1$.
\end{enumerate}
\end{twierdzenie}
\begin{proof}
W (i) niech $(u_n)$ będzie ograniczony w $W^{1,p}$.
Z twierdzenia~\ref{thm:pde10-Rellich-p} wybieramy
podciąg Cauchy'ego w $L^p$. Twierdzenie
\ref{thm:pde10-domain-embed} daje jednostajne
ograniczenie norm $L^{p^*}$, zatem różnice
$u_n-u_m$ mają jednostajnie ograniczone normy $L^{p^*}$.
Dla $p<q<p^*$ lemat~\ref{lem:pde10-interpolation}
implikuje, że podciąg jest Cauchy'ego w $L^q$.
Dla $q\le p$ wystarcza nierówność H\"oldera na
zbiorze o skończonej mierze. Tak uzyskujemy (i).

W (ii) jeśli $q>p=d$, wybieramy ustalone skończone
$r>q$. Zanurzenie $W^{1,d}\hookrightarrow L^r$ daje
ograniczenie norm $L^r$, a interpolacja między $L^d$
i $L^r$ daje zbieżność w $L^q$. Dla $q\le d$
postępujemy jak wcześniej.

W (iii) nierówność Morreya daje jednostajne ograniczenie
norm $C^{0,\alpha}(\overline\Omega)$, gdzie
$\alpha=1-d/p>0$ (albo $\alpha=1$ dla $p=\infty$).
Rodzina jest więc jednostajnie ograniczona i jednakowo
ciągła na zwartym $\overline\Omega$.
Twierdzenie Arzel\`a--Ascoliego zapewnia podciąg zbieżny
jednostajnie do ciągłej funkcji $u$.
Jeśli $f_n=u_n-u$, to $\|f_n\|_\infty\to0$
i $[f_n]_{C^{0,\alpha}}\le C$.
Dla $0<\beta<\alpha$ i każdego $\delta>0$ mamy
\[
 [f_n]_{C^{0,\beta}}
 \le\max\{2\|f_n\|_\infty\delta^{-\beta},
 C\delta^{\alpha-\beta}\}.
\]
Najpierw wybieramy $\delta$ małe, by drugi wyraz
był mały, potem $n$ duże, by pierwszy też był mały.
To dowodzi zbieżności w $C^{0,\beta}$.
Dla $\beta=0$ oznacza ona zbieżność jednostajną.
\end{proof}

\begin{kontrprzyklad}[Brak zwartości dla wykładnika krytycznego]
\label{ex:pde10-noncompact-critical}
Niech $1\le p<d$, $x_0\in\Omega$ oraz
$\varphi\in C_c^\infty(B_1(0))$ będzie niezerowa.
Funkcje $u_\varepsilon(x)=\varepsilon^{1-d/p}
\varphi((x-x_0)/\varepsilon)$, dla dostatecznie małych
$\varepsilon>0$, mają jednostajnie ograniczone normy
$W^{1,p}(\Omega)$ i stałą dodatnią normę $L^{p^*}(\Omega)$.
Jednocześnie $u_\varepsilon\to0$ prawie wszędzie, gdy
$\varepsilon\downarrow0$. Gdyby istniał podciąg zbieżny
silnie w $L^{p^*}$, jego granica musiałaby być równa zeru,
co przeczyłoby stałości norm. Dlatego zanurzenie przy
$q=p^*$ nie jest zwarte.
\end{kontrprzyklad}

\begin{kontrprzyklad}[Brak zwartości na całej $\R^d$]
\label{ex:pde10-translation}
Niech $\varphi\in C_c^\infty(B_1(0))$ będzie niezerowa
i $u_n(x)=\varphi(x-3ne_1)$.
Wówczas $\|u_n\|_{W^{1,p}(\R^d)}$ jest stała dla
każdego $1\le p<\infty$, lecz dla $m\ne n$ nośniki są
rozłączne, więc
$\|u_n-u_m\|_p^p=2\|\varphi\|_p^p$.
Żaden podciąg nie może być Cauchy'ego w $L^p$.
Ograniczoność obszaru w twierdzeniu Rellicha nie jest
zatem przypadkowym warunkiem technicznym.
\end{kontrprzyklad}

\section[Przypadek graniczny w zanurzeniach Sobolewa]{Przypadek graniczny $p=d$: dlaczego nie dostajemy $L^\infty$?}
\label{sec:pde10-Moser}

\begin{kontrprzyklad}[Ciąg koncentracyjny typu Mosera]
\label{ex:pde10-Moser}
Niech $d\ge2$, $B=B_1(0)$ oraz $0<\varepsilon<e^{-1}$.
Połóżmy $L=\log(1/\varepsilon)$ i
\[
 u_\varepsilon(x)=
 \begin{cases}
 L^{1-1/d}, & |x|\le\varepsilon,\\
 L^{-1/d}\log(1/|x|), & \varepsilon<|x|<1,\\
 0,&|x|\ge1.
 \end{cases}
\]
Funkcja jest ciągła i należy do $W^{1,d}(\R^d)$.
W istocie dla $\varepsilon<|x|<1$ mamy
$|\nabla u_\varepsilon(x)|=L^{-1/d}/|x|$, skąd
\[
 \|\nabla u_\varepsilon\|_{L^d(\R^d)}^d
 =|S^{d-1}|L^{-1}\int_\varepsilon^1\frac{\mathrm dr}{r}
 =|S^{d-1}|.
\]
Natomiast $\|u_\varepsilon\|_{L^d}$ jest jednostajnie
ograniczona: na $\varepsilon<|x|<1$ wystarcza
$L^{-1}\int_0^\infty t^de^{-dt}\,\mathrm dt<\infty$,
a na kuli $B_\varepsilon$ wkład wynosi
$\omega_d\varepsilon^dL^{d-1}$ i pozostaje ograniczony.
Pomimo tego
$\|u_\varepsilon\|_\infty=L^{1-1/d}\to\infty$.
Nie istnieje więc ciągłe zanurzenie
$W^{1,d}(B)\hookrightarrow L^\infty(B)$.
\end{kontrprzyklad}

\begin{uwaga}[Ostrość wykładnika Morreya]
Niech $p>d$ i $\alpha=1-d/p$. Dla tej samej skalowanej
rodziny $u_\varepsilon(x)=\varepsilon^\alpha
\varphi((x-x_0)/\varepsilon)$ normy $W^{1,p}$ są
ograniczone, lecz dla $\beta>\alpha$ ilorazy H\"oldera
rosną jak $\varepsilon^{\alpha-\beta}$.
Nie można więc ogólnie zastąpić wykładnika Morreya
większym wykładnikiem, nawet na kuli.
\end{uwaga}

\section{Ślady na brzegu Lipschitzowskim}
\label{sec:pde10-traces}

\begin{definicja}[Miara i przestrzeń $L^p$ na brzegu]
Jeśli $\Omega\subset\R^d$ jest obszarem ograniczonym Lipschitzowskim, przyjmujemy na $\partial\Omega$ miarę
$\mathcal H^{d-1}|_{\partial\Omega}$, czyli ograniczenie
$(d-1)$-wymiarowej miary Hausdorffa. Na każdym lokalnym
wykresie $x_d=\phi(x')$ miara ta ma element
$\sqrt{1+|\nabla\phi(x')|^2}\,\mathrm d\lambda_{d-1}(x')$
prawie wszędzie. Przestrzeń $L^p(\partial\Omega)$
rozumiemy względem tej miary.
\end{definicja}

\begin{twierdzenie}[Twierdzenie o śladzie]
\label{thm:pde10-trace}
Niech $\Omega\subset\R^d$ będzie obszarem ograniczonym Lipschitzowskim, $1\le p<\infty$. Istnieje dokładnie jeden
ograniczony operator liniowy
$\operatorname{Tr}:W^{1,p}(\Omega)\to L^p(\partial\Omega)$,
który dla każdej $u\in C^\infty(\overline\Omega)$
(gładkiej w otoczeniu $\overline\Omega$) daje zwykłą
restrykcję $u|_{\partial\Omega}$. Ponadto
\[
 \|\operatorname{Tr}u\|_{L^p(\partial\Omega)}
 \le C(\Omega,p)\|u\|_{W^{1,p}(\Omega)}.
\]
\end{twierdzenie}
\begin{proof}
Najpierw udowodnimy oszacowanie na płaskim brzegu.
Dla $w\in C^1(\R^{d-1}\times[0,a])$ o odpowiedniej
całkowalności, $t\in(0,a)$, podstawowe twierdzenie
rachunku daje
$|w(x',0)|\le|w(x',t)|+
\int_0^a|\partial_tw(x',s)|\,\mathrm d\lambda_1(s)$.
Całkujemy nierówność po $t\in(0,a)$, dzielimy przez $a$
i stosujemy H\"oldera. Otrzymujemy punktowo
\[
 |w(x',0)|\le a^{-1/p}\|w(x',\cdot)\|_{L^p(0,a)}
 +a^{1-1/p}\|\partial_tw(x',\cdot)\|_{L^p(0,a)}.
\]
Po podniesieniu do potęgi $p$ i scałkowaniu po $x'$
dostajemy
\[
 \|w(\cdot,0)\|_{L^p(\R^{d-1})}^p
 \le 2^{p-1}\bigl(a^{-1}\|w\|_{L^p}^p
 +a^{p-1}\|\partial_tw\|_{L^p}^p\bigr).
\]
W lokalnym wykresie brzegu prostujemy go odwzorowaniem
$\Phi(x',t)=(x',t+\phi(x'))$. Funkcję $\chi u$,
gdzie $\chi$ jest odcięciem wewnątrz mniejszego cylindra,
przenosimy do płaskiej warstwy. Zmiana zmiennych
ma jakobian $1$, a $\partial_t[(\chi u)\circ\Phi]$
zgadza się z $\partial_{x_d}(\chi u)\circ\Phi$.
Na wykresie element powierzchniowy jest ograniczony
z góry przez $\sqrt{1+\operatorname{Lip}(\phi)^2}$
razy $\mathrm d\lambda_{d-1}$.
Otrzymujemy więc lokalne oszacowanie śladu przez
$\|\chi u\|_{W^{1,p}(\Omega)}$.
Skończone pokrycie brzegu i podział jedności pozwalają
zsumować te nierówności; stąd globalne oszacowanie
na gładkich $u$.

Aby zdefiniować ślad ogólnej funkcji, wybieramy
ciąg restrykcji funkcji $C_c^\infty(\R^d)$ zbieżny
w $W^{1,p}(\Omega)$, którego istnienie wykazano
we wniosku~\ref{cor:pde10-density-boundary}.
Otrzymane oszacowanie pokazuje, że restrykcje przybliżeń
tworzą ciąg Cauchy'ego w $L^p(\partial\Omega)$;
jego granica nie zależy od wyboru przybliżeń.
Przedłużenie operatora jest liniowe i ograniczone,
a gęstość funkcji gładkich aż do brzegu zapewnia
jednoznaczność.
\end{proof}

\begin{definicja}[Przestrzeń danych brzegowych $H^{1/2}(\partial\Omega)$]
\label{def:pde10-halftrace}
Niech $\Omega\subset\R^d$ będzie ograniczonym obszarem Lipschitzowskim,
a $\operatorname{Tr}:H^1(\Omega)\to L^2(\partial\Omega)$
operatorem z twierdzenia~\ref{thm:pde10-trace}. Definiujemy
\[
 H^{1/2}(\partial\Omega):=\operatorname{Tr}(H^1(\Omega)),\qquad
 \|g\|_{H^{1/2}(\partial\Omega)}
 :=\inf_{\substack{G\in H^1(\Omega)\\\operatorname{Tr}G=g}}
 \|G\|_{H^1(\Omega)}.
\]
Jest to definicja przez przestrzeń śladów, z normą ilorazową.
Na ograniczonym brzegu Lipschitzowskim jest ona równoważna
zwyczajowej definicji przez ułamkowe seminormy Gagliarda;
dowód tej identyfikacji wymaga oddzielnego twierdzenia o śladach.
\end{definicja}

\begin{lemat}[Liniowy operator podnoszenia śladu]
\label{lem:pde10-halftrace-rightinverse}
Dla każdego ograniczonego obszaru Lipschitzowskiego $\Omega$
istnieje liniowy, ograniczony operator
$E_{\operatorname{Tr}}:H^{1/2}(\partial\Omega)\to H^1(\Omega)$
taki, że $\operatorname{Tr}E_{\operatorname{Tr}}g=g$ oraz
$\|E_{\operatorname{Tr}}g\|_{H^1(\Omega)}
=\|g\|_{H^{1/2}(\partial\Omega)}$.
\end{lemat}
\begin{proof}
Niech $H=H^1(\Omega)$ i $N=\ker\operatorname{Tr}$.
Ponieważ $\operatorname{Tr}:H\to L^2(\partial\Omega)$
jest ciągły, $N$ jest domkniętą podprzestrzenią Hilberta.
Mamy więc rozkład ortogonalny $H=N\oplus N^\perp$.
Dla $g\in\operatorname{Tr}(H)$ wybieramy dowolny $G\in H$
o śladzie $g$ i definiujemy
$E_{\operatorname{Tr}}g=P_{N^\perp}G$.
Jeżeli $G_1$ i $G_2$ mają ten sam ślad, to
$G_1-G_2\in N$, a zatem ich rzuty na $N^\perp$ są równe.
Konstrukcja jest więc niezależna od wyboru $G$ i liniowa.
Ponadto $G-P_{N^\perp}G\in N$, skąd
$\operatorname{Tr}E_{\operatorname{Tr}}g=g$.
Każde inne podniesienie jest postaci
$E_{\operatorname{Tr}}g+n$ dla $n\in N$, więc
z twierdzenia Pitagorasa
$\|E_{\operatorname{Tr}}g+n\|_H^2
=\|E_{\operatorname{Tr}}g\|_H^2+\|n\|_H^2$.
To dowodzi równości norm i ograniczoności operatora.
W szczególności norma ilorazowa czyni
$H^{1/2}(\partial\Omega)$ przestrzenią Hilberta.
\end{proof}

\begin{twierdzenie}[Zerowy ślad i domknięcie funkcji testowych]
\label{thm:pde10-trace-zero}
Niech $\Omega\subset\R^d$ będzie obszarem ograniczonym Lipschitzowskim, $1\le p<\infty$. Wówczas
\[
 W_0^{1,p}(\Omega)=\ker\operatorname{Tr}.
\]
\end{twierdzenie}
\begin{proof}
Jeśli $u\in W_0^{1,p}(\Omega)$, to z definicji jest
granicą w $W^{1,p}$ funkcji $C_c^\infty(\Omega)$.
Te funkcje mają zerowy ślad, więc przez ciągłość
$\operatorname{Tr}u=0$.

Dla implikacji odwrotnej niech $\operatorname{Tr}u=0$.
Rozważmy lokalny fragment brzegu wyprostowany do
półwarstwy $t>0$. Z twierdzenia o absolutnej ciągłości
na prawie wszystkich pionowych odcinkach dla funkcji
$W^{1,p}$ wynika, że dla prawie wszystkich $x'$
wartość brzegowa $w(x',0)$ istnieje i pokrywa się
ze śladem (dowodzi się tego przez gęstość w normie
$W^{1,p}$ i jednowymiarowe oszacowanie z poprzedniego
dowodu). Jest zatem równa zero. Przedłużenie $w$
przez zero na $t<0$ ma słabą pochodną normalną
równą przedłużeniu $\partial_tw$ przez zero:
podczas całkowania przez części w $t$ składnik brzegowy
$w(x',0)$ znika. Pochodne styczne również
przedłużają się przez zero, gdyż granica całkowania
w tych współrzędnych jest niezależna od $x'$.

Wybierzmy funkcję odcinającą $\chi$ o nośniku w
mniejszym cylindrze oraz dodatni kierunek pionowy
skierowany do wnętrza obszaru. Przesuwamy
przedłużenie $\chi w$ o odległość $a\varepsilon$
do wnętrza, a następnie wygładzamy jądrem o promieniu
$\varepsilon$; wybieramy stałą $a>1$ tak, aby wygładzenie
miało nośnik po stronie $t>0$. Translacja jest ciągła
w $W^{1,p}(\R^d)$ dla $p<\infty$, a splot
z przybliżoną jedynką zbiega w tej samej normie.
Wszystkie uzyskane funkcje są zatem gładkie, wspierają
się we wnętrzu odpowiedniego fragmentu obszaru
i zbiegają do $\chi w$ w $W^{1,p}$.

Po powrocie do pierwotnych współrzędnych na wykresie
Lipschitzowskim gładkość nie musi być zachowana.
Dlatego konstrukcję wykonujemy geometrycznie przed
prostowaniem: przesuwamy fragment $\chi u$ w stałym
kierunku pionowym oryginalnego cylindra na odległość
$a\varepsilon$, gdzie $a>1+\operatorname{Lip}\phi$,
a następnie wygładzamy w oryginalnych współrzędnych
promieniem $\varepsilon$. Warunek na $a$ gwarantuje,
że całe kule używane w splocie pozostają po stronie
wewnętrznej grafu. Zbieżność w $W^{1,p}$ wynika
z własności translacji w $\R^d$ i ograniczoności
operatorów mnożenia przez odcięcia.
Części podziału jedności wspierające się ściśle
wewnątrz $\Omega$ wygładzamy bez przesunięć.
Sumując skończenie wiele aproksymacji, otrzymujemy
ciąg w $C_c^\infty(\Omega)$ zbieżny do $u$
w $W^{1,p}(\Omega)$. To kończy dowód.
\end{proof}

\begin{wniosek}[Nierówność Poincar\'ego na obszarze]
\label{cor:pde10-Poincare}
Niech $\Omega\subset\R^d$ będzie obszarem ograniczonym Lipschitzowskim, $1\le p<\infty$.
Istnieje stała $C=C(\Omega,p)$ taka, że
\[
 \|u-u_\Omega\|_{L^p(\Omega)}
 \le C\|\nabla u\|_{L^p(\Omega)},\qquad
 u_\Omega=|\Omega|^{-1}\int_\Omega u\,\mathrm d\lambda_d.
\]
Dla $u\in W_0^{1,p}(\Omega)$ istnieje również stała
$C'$ z $\|u\|_p\le C'\|\nabla u\|_p$.
\end{wniosek}
\begin{proof}
Gdyby pierwsza nierówność była fałszywa, istniałyby
$u_n\in W^{1,p}(\Omega)$ takie, że
$(u_n)_\Omega=0$, $\|u_n\|_p=1$ i
$\|\nabla u_n\|_p\to0$ (odejmujemy średnią i
normalizujemy przez normę $L^p$).
Ciąg jest ograniczony w $W^{1,p}$, więc twierdzenie
\ref{thm:pde10-Rellich-p} daje podciąg zbieżny w
$L^p$ do pewnego $u$ o normie $1$ i średniej $0$.
Z zamkniętości pochodnej dystrybucyjnej
(wniosek~\ref{cor:pde9-closed}) otrzymujemy
$\nabla u=0$. Na obszarze funkcja lokalnie
całkowalna o zerowych pochodnych dystrybucyjnych
jest stała prawie wszędzie; jej średnia jest zerowa,
więc $u=0$, sprzeczność.
Dla drugiej nierówności przedłużamy funkcję z
$W_0^{1,p}(\Omega)$ przez zero na kostkę $Q\supset\Omega$
i stosujemy wniosek~\ref{cor:pde9-Poincare-zero}.
\end{proof}

% Audyt rozdziału 10: jawny krok techniczny używany w dalszej części tomu.
\begin{wniosek}[Przedłużenie przez zero funkcji o zerowym śladzie]
\label{cor:pde10-zeroextension}
Niech $\Omega\subset\R^d$ będzie obszarem ograniczonym Lipschitzowskim, $1\le p<\infty$ i
$u\in W^{1,p}(\Omega)$ z $\operatorname{Tr}u=0$.
Niech $E_0u$ oznacza funkcję równą $u$ prawie wszędzie
na $\Omega$ i $0$ poza $\Omega$. Wówczas
$E_0u\in W^{1,p}(\R^d)$ oraz
\[ \partial_j(E_0u)=E_0(\partial_j u)
\quad\text{w }\mathcal D'(\R^d),\qquad j=1,\ldots,d. \]
W konsekwencji normy $L^p$ funkcji i jej gradientu
są zachowane przez przedłużenie.
\end{wniosek}
\begin{proof}
Twierdzenie~\ref{thm:pde10-trace-zero} pozwala wybrać
$u_n\in C_c^\infty(\Omega)$ zbieżne do $u$ w $W^{1,p}(\Omega)$.
Przedłużenia przez zero $E_0u_n$ należą do
$C_c^\infty(\R^d)$ i spełniają
$\partial_j(E_0u_n)=E_0(\partial_j u_n)$ punktowo.
Z definicji przedłużenia
$\|E_0u_n-E_0u_m\|_{L^p(\R^d)}=
\|u_n-u_m\|_{L^p(\Omega)}$;
taka sama równość zachodzi dla gradientów.
Ciąg $(E_0u_n)$ jest zatem Cauchy'ego w $W^{1,p}(\R^d)$
i zbiega do pewnej funkcji $w$. Zbieżność w $L^p$ pokazuje,
że $w=E_0u$ prawie wszędzie. Zamkniętość pochodnej słabej
daje $\partial_jw=E_0(\partial_j u)$.
Równości norm wynikają z równości funkcji prawie wszędzie.
\end{proof}

\begin{uwaga}[Zależności twierdzeń i brzegi nieregularne]
Powyższy dowód świadomie korzysta z ustalonej wcześniej
równości $W_0^{1,p}(\Omega)=\ker\operatorname{Tr}$;
nie może więc zastąpić jej dowodu bez błędnego koła.
Na obszarze o nieregularnym brzegu nie należy
identyfikować obu przestrzeni bez dodatkowych hipotez.
\end{uwaga}

\section{Schemat związków i ograniczenia twierdzeń}
\label{sec:pde10-summary}
\begin{center}
\begin{tikzpicture}[node distance=8mm and 10mm,>=Stealth,
 every node/.style={font=\small,align=center}]
 \node[draw,rounded corners,fill=blue!5,text width=3.3cm] (a)
  {$W^{1,p}(\Omega)$\\$\Omega$ Lipschitzowska};
 \node[draw,rounded corners,fill=green!8,text width=3cm,
 below left=of a] (b) {$p<d$\\$L^{p^*}$ ciągle\\$L^q$, $q<p^*$ zwarto};
 \node[draw,rounded corners,fill=orange!10,text width=3cm,
 below=of a] (c) {$p=d$\\$L^q$ dla $q<\infty$\\zwarto dla $q<\infty$};
 \node[draw,rounded corners,fill=cyan!8,text width=3cm,
 below right=of a] (d) {$p>d$\\$C^{0,1-d/p}$ ciągle\\$C^{0,\beta}$ zwarto};
 \draw[->] (a)--(b);
 \draw[->] (a)--(c);
 \draw[->] (a)--(d);
\end{tikzpicture}

{\small\emph{Rysunek:} Zachowanie funkcji Sobolewa zależy od relacji
wykładnika całkowalności $p$ do wymiaru $d$. Zwarty charakter
zanurzeń zanika przy niektórych wykładnikach granicznych.}
\end{center}

Kluczowa różnica między zanurzeniem ciągłym a zwartym
polega na tym, że oszacowanie normy nie zapewnia istnienia
silnie zbieżnego podciągu. Przykład skupiania masy
pokazuje utratę zwartości przy $q=p^*$, a translacje
pokazują jej utratę na nieograniczonej przestrzeni.
Te dwa mechanizmy będą ważne w metodach wariacyjnych.


% --- Rozdział 11. Metody wariacyjne ---
\chapter[Metody wariacyjne i twierdzenie Laxa--Milgrama]{Metody wariacyjne: słabe rozwiązania, twierdzenie Laxa--Milgrama i minimalizacja energii}
\label{chap:pde-variational}

Równania eliptyczne rzadko dają się rozwiązać za pomocą jawnych wzorów.
Wcześniej nauczyliśmy się jednak kontrolować słabe pochodne,
operator śladu i normy przestrzeni Sobolewa. Teraz pokażemy, że
te narzędzia pozwalają udowadniać istnienie oraz jednoznaczność
rozwiązań bez konstruowania ich punkt po punkcie.

Kluczowe będą dwa podejścia. Pierwsze polega na rozwiązywaniu
równania liniowego w przestrzeni Hilberta; jego podstawą jest
twierdzenie Laxa--Milgrama. Drugie polega na minimalizacji energii,
czasem również nieliniowej. Będziemy wyraźnie odróżniać
\emph{formę dwuliniową} od \emph{funkcjonału energii}: jeśli forma nie
jest symetryczna, to zwykły funkcjonał kwadratowy nie daje
na ogół wyjściowego równania.

W całym rozdziale przestrzenie Hilberta oraz formy dwuliniowe są
rzeczywiste. Całkujemy względem miary Lebesgue'a
$\lambda_d$ na $\Omega\subset\R^d$, zaś na brzegu Lipschitzowskim
względem $\mathcal H^{d-1}$. Zapisy $\int_\Omega f\,\mathrm d\lambda_d(x)$
i $\int_{\partial\Omega}g\,\mathrm d\mathcal H^{d-1}(x)$
określają za każdym razem także miarę całkowania.

\section{Przestrzeń energetyczna i sformułowanie słabe}
\label{sec:pde11-weak}

Niech $\Omega\subset\R^d$ będzie obszarem ograniczonym Lipschitzowskim. Kładziemy $V=H_0^1(\Omega)=W_0^{1,2}(\Omega)$.
Z nierówności Poincar\'ego (wniosek~\ref{cor:pde10-Poincare}) wynika, że
\[ (u,v)_V=\int_\Omega\nabla u\cdot\nabla v\,\mathrm d\lambda_d(x) \]
wyznacza iloczyn skalarny, którego norma jest równoważna normie
$H^1(\Omega)$. W istocie $\|u\|_{H^1}^2=\|u\|_2^2+
\|\nabla u\|_2^2\le(1+C_P^2)\|\nabla u\|_2^2$,
a odwrotna nierówność jest bezpośrednia. Ponieważ $H_0^1$ jest
domkniętą podprzestrzenią $H^1$, przestrzeń $(V,(\cdot,\cdot)_V)$
jest przestrzenią Hilberta.

\begin{definicja}[Przestrzeń dualna i słabe rozwiązanie]
\label{def:pde11-dual}
Niech $V$ będzie rzeczywistą przestrzenią Hilberta.
Jej \emph{przestrzenią dualną} $V^*$ nazywamy przestrzeń wszystkich
ciągłych funkcjonałów liniowych $F:V\to\R$, wyposażoną w normę
$\|F\|_{V^*}=\sup_{\|v\|_V\le1}|F(v)|$.
Jeżeli $a:V\times V\to\R$ jest formą dwuliniową i $F\in V^*$,
to \emph{słabym rozwiązaniem zagadnienia} jest element $u\in V$
spełniający $a(u,v)=F(v)$ dla każdego $v\in V$.
Dla $V=H_0^1(\Omega)$ oznaczamy też $V^*=H^{-1}(\Omega)$;
w tej notacji $H^{-1}(\Omega)$ oznacza przestrzeń dualną do
$H_0^1(\Omega)$, a nie przestrzeń wszystkich dystrybucji
z ujemną regularnością na całej $\R^d$.
\end{definicja}

\begin{przyklad}[Od równania Poissona do formy słabej]
\label{ex:pde11-poisson-weak}
Załóżmy chwilowo, że $u\in C^2(\overline\Omega)$,
$u|_{\partial\Omega}=0$, $f\in L^2(\Omega)$ oraz $-\Delta u=f$.
Dla $v\in C_c^\infty(\Omega)$ całkowanie przez części daje
$\int_\Omega\nabla u\cdot\nabla v\,\mathrm d\lambda_d(x)
=\int_\Omega fv\,\mathrm d\lambda_d(x)$.
Lewa strona jest ciągła względem $v$ w normie $H_0^1$,
ponieważ ogranicza ją $\|\nabla u\|_2\|\nabla v\|_2$.
Prawa strona również jest ciągła:
$|\int_\Omega fv\,\mathrm d\lambda_d|
\le C_P\|f\|_2\|\nabla v\|_2$.
Gęstość $C_c^\infty(\Omega)$ w $H_0^1(\Omega)$
pozwala rozszerzyć tożsamość na wszystkie $v\in V$.
\end{przyklad}

W odwrotną stronę słaba tożsamość daje równość
$-\Delta u=f$ w sensie dystrybucji, bo wolno wstawić
każdą funkcję testową $v\in C_c^\infty(\Omega)$.
Warunek $u|_{\partial\Omega}=0$ jest zakodowany w wyborze $V$;
na obszarze Lipschitzowskim jest równoważny
$\operatorname{Tr}u=0$ (twierdzenie~\ref{thm:pde10-trace-zero}).
Nie wynika stąd jeszcze, że $u\in H^2$ lub $u\in C^2$.
\emph{Istnienie słabe i regularność klasyczna to dwa osobne twierdzenia.}

\section{Dwa lematy o przestrzeniach Hilberta}
\label{sec:pde11-hilbert}

Najpierw udowodnimy fakty, od których zależy twierdzenie Laxa--Milgrama.
Nie zakładamy z góry, że każdy funkcjonał można zapisać
w postaci iloczynu skalarnego.

\begin{lemat}[Rzut na domknięty zbiór wypukły]
\label{lem:pde11-projection}
Niech $H$ będzie rzeczywistą przestrzenią Hilberta,
$K\subset H$ zbiorem niepustym, domkniętym i wypukłym,
a $x\in H$. Istnieje dokładnie jeden $y\in K$ taki, że
$\|x-y\|=\inf_{z\in K}\|x-z\|$. Ponadto
\[ (x-y,z-y)_H\le0\qquad\text{dla każdego }z\in K. \]
\end{lemat}
\begin{proof}
Oznaczmy $r=\inf_{z\in K}\|x-z\|$ i wybierzmy $y_n\in K$ z
$\|x-y_n\|\to r$. Ponieważ $K$ jest wypukły,
$(y_n+y_m)/2\in K$. Tożsamość równoległoboku daje
\begin{align*}
\|y_n-y_m\|^2
&=2\|x-y_n\|^2+2\|x-y_m\|^2
 -4\left\|x-\frac{y_n+y_m}{2}\right\|^2\\
&\le 2\|x-y_n\|^2+2\|x-y_m\|^2-4r^2\longrightarrow0.
\end{align*}
Zatem $(y_n)$ jest ciągiem Cauchy'ego. Zupełność i domkniętość
zapewniają $y\in K$ oraz $\|x-y\|=r$.
Gdyby $y'\in K$ było drugim minimizatorem, ten sam rachunek
z $y_n=y$, $y_m=y'$ dałby $\|y-y'\|=0$.

Dla dowolnego $z\in K$ i $0\le t\le1$ mamy
$y+t(z-y)\in K$. Z minimalności $y$ wynika
$\|x-y-t(z-y)\|^2\ge\|x-y\|^2$.
Rozwijając kwadrat, dzieląc przez $t>0$ i przechodząc
do granicy $t\downarrow0$, otrzymujemy
$-2(x-y,z-y)_H\ge0$, czyli tezę.
\end{proof}

\begin{twierdzenie}[Twierdzenie Riesza o reprezentacji]
\label{thm:pde11-riesz}
Niech $H$ będzie rzeczywistą przestrzenią Hilberta.
Dla każdego $F\in H^*$ istnieje dokładnie jeden $w_F\in H$ taki, że
$F(v)=(w_F,v)_H$ dla wszystkich $v\in H$.
Ponadto $\|w_F\|_H=\|F\|_{H^*}$.
\end{twierdzenie}
\begin{proof}
Dla $F=0$ wystarczy $w_F=0$.
Jeżeli $F\ne0$, to $K=\ker F$ jest domkniętą
podprzestrzenią właściwą. Wybierzmy $x\notin K$.
Z lematu~\ref{lem:pde11-projection} otrzymujemy rzut
$y\in K$ punktu $x$ i kładziemy $e=x-y\ne0$.
Z charakterystyki rzutu, podstawiając $z=y\pm k$,
$k\in K$, dostajemy $(e,k)_H=0$.
Dla dowolnego $v\in H$ element
$k=v-[F(v)/F(e)]e$ należy do $K$, ponieważ
$F(e)=F(x)\ne0$. Stąd
$0=(e,k)_H=(e,v)_H-[F(v)/F(e)]\|e\|_H^2$.
Otrzymujemy $F(v)=(w_F,v)_H$ dla
$w_F=F(e)e/\|e\|_H^2$.
Jeżeli dwa wektory reprezentują $F$, podstawienie ich różnicy
jako argumentu daje ich równość. Nierówność Cauchy'ego--Schwarza
daje $\|F\|\le\|w_F\|$; biorąc
$v=w_F/\|w_F\|$ w przypadku $w_F\ne0$,
uzyskujemy nierówność odwrotną.
\end{proof}

\section{Twierdzenie Laxa--Milgrama z pełnym dowodem}
\label{sec:pde11-lax}

\begin{definicja}[Ograniczoność i koercywność]
\label{def:pde11-coercivity}
Niech $H$ będzie rzeczywistą przestrzenią Hilberta.
Formę dwuliniową $a:H\times H\to\R$ nazywamy
\emph{ograniczoną}, jeżeli istnieje $M\ge0$ takie, że
$|a(u,v)|\le M\|u\|_H\|v\|_H$ dla wszystkich $u,v\in H$.
Nazywamy ją \emph{koercywną}, jeżeli istnieje $\alpha>0$,
dla którego $a(v,v)\ge\alpha\|v\|_H^2$ dla każdego $v\in H$.
Koercywność nie jest tym samym co symetria: forma
może być koercywna i jednocześnie niesymetryczna.
\end{definicja}

\begin{twierdzenie}[Twierdzenie Laxa--Milgrama]
\label{thm:pde11-LM}
Niech $H$ będzie rzeczywistą przestrzenią Hilberta,
$a:H\times H\to\R$ ograniczoną formą dwuliniową
ze stałą $M>0$, koercywną ze stałą $\alpha>0$,
a $F\in H^*$. Wtedy istnieje dokładnie jeden element $u\in H$
spełniający
\[ a(u,v)=F(v)\qquad\text{dla każdego }v\in H. \]
Ponadto $\|u\|_H\le\alpha^{-1}\|F\|_{H^*}$.
Nie wymaga się symetrii formy $a$.
\end{twierdzenie}
\begin{proof}
\textbf{Krok 1: operator reprezentujący formę.}
Dla ustalonego $u\in H$ odwzorowanie $v\mapsto a(u,v)$
jest funkcjonałem liniowym i ograniczonym o normie
co najwyżej $M\|u\|_H$. Z twierdzenia Riesza
(\ref{thm:pde11-riesz}) istnieje dokładnie jeden
$Tu\in H$ spełniający
$(Tu,v)_H=a(u,v)$ dla każdego $v\in H$.
Liniowość $a$ w pierwszej zmiennej i jednoznaczność
reprezentanta implikują liniowość $T$.
Równość norm w twierdzeniu Riesza daje
$\|Tu\|_H\le M\|u\|_H$, a zatem $T$ jest ograniczony.
Koercywność oznacza $(Tu,u)_H\ge\alpha\|u\|_H^2$.
Z kolei $F(v)=(f,v)_H$ dla jednoznacznie wyznaczonego $f\in H$.
Wystarczy zatem rozwiązać równanie $Tu=f$.

\textbf{Krok 2: konstrukcja kontrakcji.}
Dla parametru $t>0$ zdefiniujmy
$S_t(u)=u-t(Tu-f)$. Dla $u,w\in H$, oznaczając $z=u-w$,
obliczamy dokładnie
\begin{align*}
\|S_t(u)-S_t(w)\|_H^2
&=\|z-tTz\|_H^2\\
&=\|z\|_H^2-2t(Tz,z)_H+t^2\|Tz\|_H^2\\
&\le (1-2t\alpha+t^2M^2)\|z\|_H^2.
\end{align*}
Wybierzmy $0<t<2\alpha/M^2$.
Ponieważ $\alpha\le M$ (z ograniczoności i koercywności),
liczba $q=\sqrt{1-2t\alpha+t^2M^2}$
należy do $[0,1)$. Odwzorowanie $S_t$ jest więc kontrakcją.

\textbf{Krok 3: punkt stały bez pomijania granicy.}
Ustalmy dowolne $u_0\in H$ i zdefiniujmy
$u_{n+1}=S_t(u_n)$. Z kontrakcji wynika
$\|u_{n+1}-u_n\|_H\le q^n\|u_1-u_0\|_H$.
Dla $m>n$ nierówność trójkąta daje
$\|u_m-u_n\|_H\le q^n(1-q)^{-1}\|u_1-u_0\|_H$,
więc $(u_n)$ jest ciągiem Cauchy'ego.
Niech $u=\lim_nu_n$. Ciągłość $S_t$ zapewnia
$u=S_t(u)$, a więc $Tu=f$.
Jeżeli $Tu=f=Tw$, to z koercywności
$\alpha\|u-w\|_H^2\le(T(u-w),u-w)_H=0$;
stąd $u=w$.

\textbf{Krok 4: oszacowanie rozwiązania.}
Podstawiając $v=u$ w tożsamości słabej, otrzymujemy
$\alpha\|u\|_H^2\le a(u,u)=F(u)
\le\|F\|_{H^*}\|u\|_H$.
Dla $u\ne0$ dzielimy przez $\|u\|_H$;
dla $u=0$ oszacowanie jest oczywiste.
\end{proof}

\begin{wniosek}[Stabilność względem prawej strony]
\label{cor:pde11-stability}
Przy założeniach twierdzenia~\ref{thm:pde11-LM},
niech $u_i\in H$ będą rozwiązaniami odpowiadającymi
$F_i\in H^*$, $i=1,2$. Wówczas
$\|u_1-u_2\|_H\le\alpha^{-1}\|F_1-F_2\|_{H^*}$.
\end{wniosek}
\begin{proof}
Odejmujemy słabe równania: $a(u_1-u_2,v)=(F_1-F_2)(v)$.
Wybór $v=u_1-u_2$ i oszacowanie z końca dowodu
Laxa--Milgrama dają tezę. To właśnie ilościowa stabilność
w sensie Hadamarda.
\end{proof}

\begin{uwaga}[Metoda iteracyjna a dowód istnienia]
Punkt stały w dowodzie dostarcza jednocześnie algorytmu:
$u_{n+1}=u_n-t(Tu_n-f)$, o zbieżności geometrycznej
przy podanym wyborze $t$. Sam operator $T$ zależy jednak
od wybranego iloczynu skalarnego w $H$, toteż bez
skończeniewymiarowej reprezentacji nie jest to automatycznie
praktyczna metoda numeryczna.
\end{uwaga}

\section{Równanie eliptyczne z mierzalnymi współczynnikami}
\label{sec:pde11-elliptic}

\begin{twierdzenie}[Zagadnienie Dirichleta w postaci dywergencyjnej]
\label{thm:pde11-elliptic}
Niech $\Omega\subset\R^d$ będzie obszarem ograniczonym Lipschitzowskim. Załóżmy, że $A:\Omega\to\R^{d\times d}$
ma mierzalne, istotnie ograniczone współczynniki,
a $c\in L^\infty(\Omega)$ jest nieujemna prawie wszędzie.
Załóżmy ponadto, że dla pewnego $\nu>0$ oraz
prawie wszystkich $x\in\Omega$ zachodzi
$\xi^TA(x)\xi\ge\nu|\xi|^2$ dla każdego $\xi\in\R^d$.
Wówczas dla każdego $F\in H^{-1}(\Omega)$ istnieje
jednoznaczne $u\in H_0^1(\Omega)$ spełniające
\[
\int_\Omega\big(A\nabla u\cdot\nabla v+cuv\big)
\,\mathrm d\lambda_d(x)=F(v)
\quad\text{dla wszystkich }v\in H_0^1(\Omega).
\]
Przy normie $\|v\|_V=\|\nabla v\|_2$ zachodzi
$\|u\|_V\le\nu^{-1}\|F\|_{V^*}$.
Nie zakłada się symetrii macierzy $A(x)$.
\end{twierdzenie}
\begin{proof}
Na $V=H_0^1(\Omega)$ zdefiniujmy
$a(u,v)=\int_\Omega(A\nabla u\cdot\nabla v+cuv)
\,\mathrm d\lambda_d(x)$.
Oznaczmy $K=\operatorname*{ess\,sup}_{x\in\Omega}
\|A(x)\|_{\mathrm{op}}$ oraz $C=\|c\|_\infty$.
Z Cauchy'ego--Schwarza i nierówności Poincar\'ego
otrzymujemy
\begin{align*}
|a(u,v)|
&\le K\|\nabla u\|_2\|\nabla v\|_2
+C\|u\|_2\|v\|_2\\
&\le(K+CC_P^2)\|u\|_V\|v\|_V.
\end{align*}
Tak więc $a$ jest formą ograniczoną.
Ponieważ $c\ge0$, eliptyczność zapewnia
$a(v,v)\ge\nu\int_\Omega|\nabla v|^2
\,\mathrm d\lambda_d=\nu\|v\|_V^2$.
Wszystkie założenia twierdzenia~\ref{thm:pde11-LM} są spełnione,
więc otrzymujemy istnienie, jednoznaczność i oszacowanie.
\end{proof}

\begin{uwaga}[Znaczenie sformułowania dystrybucyjnego]
Dla $F_f(v)=\int_\Omega fv\,\mathrm d\lambda_d(x)$,
$f\in L^2(\Omega)$, zachodzi
$\|F_f\|_{V^*}\le C_P\|f\|_2$.
W tym przypadku rozwiązanie spełnia w $\mathcal D'(\Omega)$
$-\operatorname{div}(A\nabla u)+cu=f$.
Nie ma tu obowiązku definiowania klasycznej dywergencji
$A\nabla u$: przy mierzalnych współczynnikach produkt
$A\nabla u$ może nie mieć słabych pochodnych w $L^2$,
ale ma sens jako wektor z $L^2$, którego dywergencja
jest dystrybucją.
\end{uwaga}

\begin{przyklad}[Jawne rozwiązanie w jednym wymiarze]
\label{ex:pde11-1d}
Dla $\Omega=(0,1)$, $A=1$, $c=0$ i $f=1$
szukamy $u\in H_0^1(0,1)$ spełniającego
$\int_0^1u'(x)v'(x)\,\mathrm d\lambda_1(x)
=\int_0^1v(x)\,\mathrm d\lambda_1(x)$.
Funkcja $u(x)=x(1-x)/2$ spełnia $-u''=1$ i zerowe
warunki brzegowe; całkowanie przez części dowodzi
słabej tożsamości. Twierdzenie~\ref{thm:pde11-elliptic}
wykazuje, że nie istnieje żadne inne rozwiązanie w $H_0^1$.
\end{przyklad}

\begin{twierdzenie}[Koercywność przy wyrazach niższego rzędu]
\label{thm:pde11-lower-order}
Niech $\Omega\subset\R^d$ będzie ograniczonym obszarem
Lipschitzowskim, a $V=H_0^1(\Omega)$ z normą
$\|u\|_V=\|\nabla u\|_{L^2(\Omega)}$. Niech
$A\in L^\infty(\Omega;\R^{d\times d})$ spełnia
$\xi^TA(x)\xi\ge\nu|\xi|^2$ dla pewnego $\nu>0$,
wszystkich $\xi\in\R^d$ i prawie wszystkich $x\in\Omega$.
Niech $b\in L^\infty(\Omega;\R^d)$ i
$c\in L^\infty(\Omega;\R)$; oznaczmy
$c_-(x)=\max\{-c(x),0\}$ i przez $C_P$ dowolną
stałą w nierówności $\|u\|_2\le C_P\|\nabla u\|_2$
dla $u\in V$. Jeśli
\[
\alpha:=\nu-C_P\|b\|_\infty-C_P^2\|c_-\|_\infty>0,
\]
to forma
\[
 a(u,v)=\int_\Omega\bigl(A\nabla u\cdot\nabla v
 +(b\cdot\nabla u)v+cuv\bigr)\,\mathrm d\lambda_d(x)
\]
jest ograniczona i koercywna na $V$ ze stałą co
najmniej $\alpha$. Dla każdego $F\in V^*$ równanie
$a(u,v)=F(v)$, $v\in V$, ma dokładnie jedno rozwiązanie
i zachodzi $\|u\|_V\le\alpha^{-1}\|F\|_{V^*}$.
\end{twierdzenie}
\begin{proof}
Z nierówności Cauchy'ego--Schwarza, H\"oldera i
Poincar\'ego otrzymujemy
\[
 |a(u,v)|\le
 (\|A\|_\infty+C_P\|b\|_\infty
 +C_P^2\|c\|_\infty)\|u\|_V\|v\|_V.
\]
Z kolei dla $u\in V$ dodatnia część $c$ nie może
obniżać wartości $a(u,u)$, a pozostałe wyrazy dają
\begin{align*}
 a(u,u)&\ge\nu\|\nabla u\|_2^2
 -\|b\|_\infty\|\nabla u\|_2\|u\|_2
 -\|c_-\|_\infty\|u\|_2^2\\
 &\ge\alpha\|u\|_V^2.
\end{align*}
Twierdzenie Laxa--Milgrama~\ref{thm:pde11-LM}
kończy dowód. Warunek ma charakter wystarczający;
ujemna część potencjału lub dryft mogą być większe,
jeżeli znamy dodatkową strukturę formy.
\end{proof}

\begin{przyklad}[Dryft bez ograniczenia małości]
\label{ex:pde11-divfree-drift}
Niech $b\in W^{1,\infty}(\Omega;\R^d)$ spełnia
$\operatorname{div}b=0$ prawie wszędzie oraz
$A=I$, $c=0$. Dla każdego $u\in C_c^\infty(\Omega)$
całkowanie przez części daje
$\int_\Omega(b\cdot\nabla u)u\,\mathrm d\lambda_d
=\frac12\int_\Omega b\cdot\nabla(u^2)\,\mathrm d\lambda_d=0$.
Ograniczoność $b$ i gęstość funkcji testowych w
$H_0^1(\Omega)$ przenoszą tę tożsamość na dowolne
$u\in H_0^1(\Omega)$. Dlatego
$a(u,u)=\|\nabla u\|_2^2$ i forma jest koercywna
ze stałą $1$ bez założenia, że $\|b\|_\infty$ jest mała.
Sam dryft nie jest więc przeszkodą: znaczenie ma również
jego struktura i warunki brzegowe.
\end{przyklad}

\section{Niezerowe dane brzegowe: metoda podnoszenia}
\label{sec:pde11-lifting}

\begin{twierdzenie}[Warunek Dirichleta z danymi brzegowymi]
\label{thm:pde11-lifting}
Niech $\Omega,A,c,\nu$ spełniają założenia
twierdzenia~\ref{thm:pde11-elliptic}. Niech
$f\in H^{-1}(\Omega)$ oraz $g\in H^1(\Omega)$.
Istnieje dokładnie jedna funkcja $u\in H^1(\Omega)$
taka, że $u-g\in H_0^1(\Omega)$ oraz
\[
\int_\Omega(A\nabla u\cdot\nabla v+cuv)
\,\mathrm d\lambda_d(x)=f(v)
\quad(v\in H_0^1(\Omega)).
\]
Jej ślad jest równy $\operatorname{Tr}g$.
Istnieje stała $C$, zależna wyłącznie od
$\Omega,\nu,\|A\|_\infty,\|c\|_\infty$, taka, że
$\|u\|_{H^1}\le C(\|f\|_{H^{-1}}+\|g\|_{H^1})$.
\end{twierdzenie}
\begin{proof}
Podstawiamy $u=w+g$. Żądany warunek brzegowy jest równoważny
$w\in V=H_0^1(\Omega)$, a równanie przyjmuje postać
$a(w,v)=G(v)$, gdzie $G(v)=f(v)-a(g,v)$.
Ograniczoność współczynników, Cauchy--Schwarz i
Poincar\'e dla $v\in V$ zapewniają
$|a(g,v)|\le C_0\|g\|_{H^1}\|v\|_V$.
Wobec tego $G\in V^*$ i
$\|G\|_{V^*}\le\|f\|_{V^*}+C_0\|g\|_{H^1}$.
Twierdzenie Laxa--Milgrama daje jednoznaczne $w\in V$ oraz
$\|w\|_V\le\nu^{-1}\|G\|_{V^*}$.
Ponieważ $\|w\|_{H^1}\le(1+C_P^2)^{1/2}\|w\|_V$,
uzyskujemy żądane oszacowanie dla $u=w+g$.
Równość śladów wynika z
$\operatorname{Tr}(u-g)=0$ i twierdzenia
\ref{thm:pde10-trace-zero}.
\end{proof}

\begin{uwaga}[Dane na samym brzegu]
Dane $\varphi$ na $\partial\Omega$ można podać niezależnie,
ale nie każde $\varphi\in L^2(\partial\Omega)$ jest śladem
funkcji z $H^1(\Omega)$. W twierdzeniu zakładamy explicite
istnienie podniesienia $g\in H^1(\Omega)$.
Przestrzeń danych brzegowych $H^{1/2}(\partial\Omega)$
z normą ilorazową i jej liniowe podnoszenie zdefiniowano
już w~\ref{def:pde10-halftrace} i
\ref{lem:pde10-halftrace-rightinverse}. Równoważność tej
normy z normą Gagliarda pozostaje oddzielnym twierdzeniem.
\end{uwaga}


\begin{uwaga}[Wersja z danymi zadanymi na brzegu]
Twierdzenie~\ref{thm:pde11-lifting} sformułowano
z wybranym podniesieniem $g\in H^1(\Omega)$.
Z lematu~\ref{lem:pde10-halftrace-rightinverse} można
bez nowych argumentów otrzymać dane
$\varphi\in H^{1/2}(\partial\Omega)$ i oszacowanie
przez ich normę śladową. Uogólnienie to, także
z wyrazami niższego rzędu, zostaje jawnie sformułowane
w twierdzeniu~\ref{thm:pde13-nonhom}; nie powielamy
tutaj jego dowodu.
\end{uwaga}

\section{Zagadnienie Neumanna i warunek zgodności}
\label{sec:pde11-neumann}

Dla laplasjanu z warunkiem Neumanna nie można oczekiwać
jednoznaczności: dodanie stałej do rozwiązania nie zmienia
gradientu ani pochodnej normalnej. Warunek całkowy
pojawia się natomiast już w sformułowaniu słabym.

\begin{twierdzenie}[Słabe zagadnienie Neumanna]
\label{thm:pde11-neumann}
Niech $\Omega\subset\R^d$ będzie obszarem ograniczonym Lipschitzowskim, $f\in L^2(\Omega)$ oraz
$h\in L^2(\partial\Omega)$.
Rozważmy równanie słabe
\[
\int_\Omega\nabla u\cdot\nabla v\,\mathrm d\lambda_d(x)
=\int_\Omega fv\,\mathrm d\lambda_d(x)
+\int_{\partial\Omega}h\operatorname{Tr}v
\,\mathrm d\mathcal H^{d-1}(x)
\quad(v\in H^1(\Omega)).
\]
Istnieje rozwiązanie $u\in H^1(\Omega)$ wtedy i tylko wtedy, gdy
\[
\int_\Omega f\,\mathrm d\lambda_d(x)
+\int_{\partial\Omega}h\,\mathrm d\mathcal H^{d-1}(x)=0.
\]
Gdy ono istnieje, wszystkie rozwiązania różnią się o stałą.
Po narzuceniu normalizacji $\int_\Omega u\,\mathrm d\lambda_d=0$
rozwiązanie jest jednoznaczne.
\end{twierdzenie}
\begin{proof}
\textbf{Konieczność.} Funkcja $v\equiv1$ należy do $H^1(\Omega)$
i ma zerowy gradient, więc jej podstawienie daje warunek całkowy.

\textbf{Wystarczalność.} Niech
$V_\circ=\{v\in H^1(\Omega):\int_\Omega v\,\mathrm d\lambda_d=0\}$.
Jest to domknięta podprzestrzeń $H^1$, więc przestrzeń Hilberta.
Nierówność Poincar\'ego z odjęciem średniej
(wniosek~\ref{cor:pde10-Poincare}) pokazuje, że
$\|v\|_{V_\circ}=\|\nabla v\|_2$ jest na niej normą
równoważną normie $H^1$.
Forma $a(u,v)=\int_\Omega\nabla u\cdot\nabla v\,
\mathrm d\lambda_d$ jest ograniczona i koercywna w tej normie.
Prawa strona $F(v)$ jest liniowa i ciągła na $H^1(\Omega)$:
$|\int_\Omega fv|\le\|f\|_2\|v\|_2$,
a twierdzenie o śladzie~\ref{thm:pde10-trace} daje
$|\int_{\partial\Omega}h\operatorname{Tr}v|
\le C\|h\|_{L^2(\partial\Omega)}\|v\|_{H^1}$.
Z twierdzenia Laxa--Milgrama istnieje jednoznaczne $u\in V_\circ$
z $a(u,v)=F(v)$ dla wszystkich $v\in V_\circ$.
Dla dowolnego $v\in H^1(\Omega)$ rozkładamy
$v=(v-v_\Omega)+v_\Omega$.
Pierwszy składnik należy do $V_\circ$; dla stałej
$v_\Omega$ obie strony równania są równe zeru dzięki
warunkowi zgodności. Dostajemy równość dla każdego $v\in H^1$.

\textbf{Jednoznaczność modulo stałe.} Jeśli $u_1,u_2$
są rozwiązaniami, różnica $w=u_1-u_2$ spełnia
$\int_\Omega|\nabla w|^2\,\mathrm d\lambda_d=0$
po podstawieniu $v=w$. Wobec tego $\nabla w=0$ prawie wszędzie,
a spójność $\Omega$ implikuje, że $w$ jest stała prawie wszędzie.
\end{proof}

\section{Metoda bezpośrednia rachunku wariacyjnego}
\label{sec:pde11-direct}

Założenie, że rozwiązanie jest punktem minimum odpowiedniej energii,
pozwala badać również równania nieliniowe.
Do tego potrzebujemy precyzyjnego rozróżnienia zbieżności
silnej i słabej.

\begin{definicja}[Zbieżność słaba]
\label{def:pde11-weakconv}
W przestrzeni Hilberta $H$ piszemy $u_n\rightharpoonup u$,
jeżeli $(u_n,v)_H\to(u,v)_H$ dla każdego $v\in H$.
W przestrzeni Banacha $X$ oznacza to warunek
$F(u_n)\to F(u)$ dla każdego $F\in X^*$.
Funkcjonał $J:X\to(-\infty,+\infty]$ jest
\emph{słabo półciągły z dołu w sensie ciągów}, jeżeli
$u_n\rightharpoonup u$ implikuje
$J(u)\le\liminf_n J(u_n)$.
Jest \emph{koercywny}, jeżeli $J(u)\to+\infty$,
gdy $\|u\|_X\to\infty$.
\end{definicja}

\begin{lemat}[Słaba zwartość ograniczonych ciągów w Hilbercie]
\label{lem:pde11-weakcompact}
Każdy ciąg ograniczony w rzeczywistej przestrzeni Hilberta
$H$ zawiera podciąg słabo zbieżny.
\end{lemat}
\begin{proof}
Niech $(u_n)$ będzie ciągiem ograniczonym przez $C$.
Jego domknięta otoczka liniowa
$H_0=\overline{\operatorname{span}\{u_n:n\ge1\}}$
jest separowalną przestrzenią Hilberta.
Jeżeli $H_0$ jest nieskończeniewymiarowa,
wybieramy przeliczalną bazę ortonormalną $(e_j)$.
Każdy ciąg liczb $((u_n,e_j))_n$ jest ograniczony;
argument przekątniowy daje podciąg $(u_{n_k})$ taki, że
$(u_{n_k},e_j)\to a_j$ dla każdego $j$.
Z nierówności Bessela
$\sum_{j=1}^N|a_j|^2
=\lim_{k\to\infty}\sum_{j=1}^N|(u_{n_k},e_j)|^2\le C^2$
dla każdego $N$. Zatem $\sum_j|a_j|^2\le C^2$,
więc istnieje $u\in H_0$ o współrzędnych $(a_j)$.
Dla skończonych kombinacji wektorów $e_j$ zachodzi
$(u_{n_k},v)\to(u,v)$.
Dla ogólnego $v\in H_0$ przybliżamy go takimi kombinacjami:
$|(u_{n_k}-u,v-v_N)|\le 2C\|v-v_N\|$,
co pozwala najpierw przejść $k\to\infty$, a potem $N\to\infty$.
Dla $v\in H$ rozkład $v=P_{H_0}v+(v-P_{H_0}v)$
redukuje twierdzenie do $H_0$, bo $u_n,u\in H_0$.
W przypadku skończeniewymiarowym zwykła zwartość
ograniczonego zbioru dostarcza nawet zbieżności silnej.
\end{proof}

\begin{lemat}[Słaba półciągłość normy]
\label{lem:pde11-weaklsc}
Niech $H$ będzie rzeczywistą przestrzenią Hilberta.
Jeżeli $u_n\rightharpoonup u$ w $H$, to
$\|u\|_H\le\liminf_n\|u_n\|_H$.
Ponadto dla każdego $F\in H^*$ mamy $F(u_n)\to F(u)$.
\end{lemat}
\begin{proof}
Jeżeli $u=0$, nierówność jest oczywista.
W przeciwnym razie słaba zbieżność daje
$\|u\|_H^2=\lim_n(u_n,u)_H
\le(\liminf_n\|u_n\|_H)\|u\|_H$,
po czym dzielimy przez $\|u\|_H$.
Druga własność wynika z definicji i twierdzenia Riesza.
\end{proof}

\begin{twierdzenie}[Metoda bezpośrednia w przestrzeni Hilberta]
\label{thm:pde11-direct}
Niech $H$ będzie rzeczywistą przestrzenią Hilberta,
a $J:H\to\R$ funkcjonałem ograniczonym z dołu,
koercywnym i słabo półciągłym z dołu w sensie ciągów.
Wówczas istnieje $u\in H$ taki, że
$J(u)=\inf_{v\in H}J(v)$.
Jeżeli $J$ jest ponadto ściśle wypukły,
to minimizator jest jedyny.
\end{twierdzenie}
\begin{proof}
Niech $m=\inf_HJ\in\R$ i wybierzmy $u_n\in H$ z
$m\le J(u_n)<m+1/n$. Koercywność zapewnia
ograniczoność $(u_n)$: gdyby istniał podciąg
o normach dążących do nieskończoności,
to jego energie dążyłyby do $+\infty$, w sprzeczności z
$J(u_n)\to m$.
Na mocy lematu~\ref{lem:pde11-weakcompact}
istnieje podciąg $u_{n_k}\rightharpoonup u$.
Słaba półciągłość z dołu daje
$J(u)\le\liminf_kJ(u_{n_k})=m$.
Z definicji infimum również $m\le J(u)$;
zatem $J(u)=m$.
Jeśli istniałyby dwa różne minimizatory $u,w$,
ścisła wypukłość zapewniłaby
$J((u+w)/2)<(J(u)+J(w))/2=m$, sprzeczność.
\end{proof}

\begin{uwaga}[Dlaczego słaba półciągłość jest niezbędna?]
Sama koercywność nie wystarcza. Przykładowo na $H=\R$
zdefiniujmy $J(0)=1$, $J(x)=x^2$ dla $x\ne0$.
Funkcjonał jest ograniczony z dołu i koercywny,
lecz jego infimum wynosi $0$ i nie jest osiągane.
Nie jest też półciągły z dołu w punkcie $0$.
\end{uwaga}

\section{Energia kwadratowa i konieczność symetrii}
\label{sec:pde11-energy}

\begin{twierdzenie}[Zasada Dirichleta dla formy symetrycznej]
\label{thm:pde11-dirichlet}
Niech $H$ będzie rzeczywistą przestrzenią Hilberta,
$a:H\times H\to\R$ ograniczoną, \emph{symetryczną}
i koercywną formą dwuliniową, $F\in H^*$.
Zdefiniujmy $J:H\to\R$ wzorem
$J(v)=\frac12 a(v,v)-F(v)$.
Wtedy $J$ ma dokładnie jeden minimizator $u$, który jest
jednocześnie jedynym rozwiązaniem słabym
$a(u,v)=F(v)$ dla wszystkich $v\in H$.
Dla każdego $w\in H$ zachodzi tożsamość
\[ J(w)-J(u)=\frac12a(w-u,w-u)
\ge\frac\alpha2\|w-u\|_H^2. \]
\end{twierdzenie}
\begin{proof}
\textbf{Koercywność funkcjonału.} Z założeń otrzymujemy
\[ J(v)\ge\frac\alpha2\|v\|_H^2
 -\|F\|_{H^*}\|v\|_H\longrightarrow+\infty
 \quad\text{gdy }\|v\|_H\to\infty. \]

\textbf{Półciągłość z dołu.} Koercywna forma symetryczna
$a$ definiuje iloczyn skalarny $(v,w)_a=a(v,w)$.
Ograniczoność i koercywność dają równoważność norm:
$\sqrt\alpha\|v\|_H\le\|v\|_a\le\sqrt M\|v\|_H$.
Z twierdzenia Riesza w $(H,(\cdot,\cdot)_a)$ i
równoważności norm wynika, że zbieżność słaba w $H$
jest taka sama jak zbieżność słaba w nowym iloczynie.
Lemat~\ref{lem:pde11-weaklsc} daje więc
$\|v\|_a^2\le\liminf_n\|v_n\|_a^2$
dla $v_n\rightharpoonup v$; składnik $F(v_n)$ zbiega do $F(v)$.
Stąd $J$ jest słabo półciągły z dołu.
Twierdzenie~\ref{thm:pde11-direct} daje minimizator.

\textbf{Równanie Eulera.} Dla dowolnych $u,v\in H$
i $t\in\R$ wykorzystujemy symetrię:
$J(u+tv)=J(u)+t[a(u,v)-F(v)]+\frac{t^2}{2}a(v,v)$.
Jeżeli $u$ minimalizuje $J$, wielomian po prawej ma
minimum w $t=0$, więc jego współczynnik liniowy znika:
$a(u,v)=F(v)$. Odwrotnie, jeśli $u$ spełnia tę równość,
to po rozwinięciu $J(u+(w-u))$ otrzymujemy
$J(w)-J(u)=\frac12a(w-u,w-u)\ge0$.
Koercywność czyni nierówność ostrą dla $w\ne u$,
dając jednoznaczność. Ostatnia nierówność daje również
ilościową kontrolę błędu energetycznego.
\end{proof}

\begin{kontrprzyklad}[Dlaczego symetrii nie wolno pominąć]
\label{ex:pde11-nonsymmetric}
Na $H=\R^2$ przyjmijmy
$B=\begin{pmatrix}1&-1\\1&1\end{pmatrix}$,
$a(u,v)=(Bu)\cdot v$ oraz $F(v)=(1,0)\cdot v$.
Wtedy $a(u,u)=|u|^2$, więc $a$ jest koercywna i ograniczona.
Minimizer funkcjonału $J(u)=|u|^2/2-(1,0)\cdot u$
wynosi $u_J=(1,0)$. Tymczasem równanie
$a(u,v)=F(v)$ dla wszystkich $v$ jest równoważne
$Bu=(1,0)$ i ma rozwiązanie $u_*=\frac12(1,-1)$.
Zatem $u_J\ne u_*$. Część antysymetryczna formy znika
w $a(v,v)$, ale nie znika z samego równania słabego.
\end{kontrprzyklad}

\section[Operator p-Laplace'a]{Pierwszy problem nieliniowy: operator $p$-Laplace'a}
\label{sec:pde11-plaplace}

W części nieliniowej musimy użyć faktu, że dla
$1<p<\infty$ przestrzeń $W_0^{1,p}(\Omega)$ jest refleksywna.
Jest ona domkniętą podprzestrzenią produktu skończonej liczby
przestrzeni $L^p$, a przestrzenie $L^p$ są refleksywne w tym zakresie.
Będziemy korzystać ze standardowego twierdzenia, że każdy ciąg
ograniczony w przestrzeni refleksywnej ma słabo zbieżny podciąg
(twierdzenie Eberleina--\v Smuliana); dowód tego ogólnego wyniku
należy do analizy funkcjonalnej, nie do konstrukcji samego PDE.

\begin{twierdzenie}[Istnienie i jednoznaczność dla $p$-laplasjanu]
\label{thm:pde11-plaplace}
Niech $\Omega\subset\R^d$ będzie obszarem ograniczonym Lipschitzowskim, $1<p<\infty$, $p'=p/(p-1)$,
a $f\in L^{p'}(\Omega)$. Istnieje dokładnie jeden
$u\in W_0^{1,p}(\Omega)$ taki, że
\[
\int_\Omega |\nabla u|^{p-2}\nabla u\cdot\nabla v
\,\mathrm d\lambda_d(x)
=\int_\Omega fv\,\mathrm d\lambda_d(x)
\quad\text{dla każdego }v\in W_0^{1,p}(\Omega).
\]
W szczególności $u$ jest rozwiązaniem dystrybucyjnym
$-\operatorname{div}(|\nabla u|^{p-2}\nabla u)=f$.
Ponadto
$\|\nabla u\|_p^{p-1}\le C_P\|f\|_{p'}$.
\end{twierdzenie}
\begin{proof}
\textbf{Krok 1: energia i ograniczenie z dołu.}
Na $X=W_0^{1,p}(\Omega)$ przyjmijmy
$\|v\|_X=\|\nabla v\|_p$, co jest normą równoważną
$W^{1,p}$ dzięki nierówności Poincar\'ego.
Definiujemy
\[ J_p(v)=\frac1p\int_\Omega|\nabla v|^p
\,\mathrm d\lambda_d(x)-\int_\Omega fv
\,\mathrm d\lambda_d(x). \]
H\"older i Poincar\'e dają
$|\int_\Omega fv|\le\|f\|_{p'}\|v\|_p
\le C_P\|f\|_{p'}\|\nabla v\|_p$.
Zatem
$J_p(v)\ge\frac1p\|v\|_X^p-C_P\|f\|_{p'}\|v\|_X$,
co implikuje koercywność i ograniczenie z dołu,
ponieważ $p>1$.

\textbf{Krok 2: minimalizator.}
Niech $(v_n)$ będzie ciągiem minimalizującym.
Koercywność zapewnia ograniczoność w $X$,
a refleksywność pozwala wybrać podciąg
$v_{n_k}\rightharpoonup u$ w $X$.
Operator gradientu jest ograniczony z $X$ do $L^p(\Omega)^d$,
więc $\nabla v_{n_k}\rightharpoonup\nabla u$ w $L^p$.
Norma w przestrzeni Banacha jest słabo półciągła z dołu:
$\|z\|=\sup_{\|\ell\|_{X^*}\le1}\ell(z)$,
a przy zbieżności słabej każdy składnik supremum
przechodzi do granicy; stąd
$\|\nabla u\|_p^p\le\liminf_k\|\nabla v_{n_k}\|_p^p$.
Ponieważ funkcjonał $v\mapsto\int_\Omega fv$ jest
ciągły i liniowy na $X$, jego wartości na $v_{n_k}$
zbiegają do wartości na $u$. Zatem
$J_p(u)\le\liminf_kJ_p(v_{n_k})=\inf_XJ_p$.
Odwrotna nierówność wynika z definicji infimum,
więc $u$ minimalizuje $J_p$.

\textbf{Krok 3: wyprowadzenie słabego równania.}
Dla $u,v\in X$ rozważamy $\psi(t)=J_p(u+tv)$.
Funkcja $\xi\mapsto|\xi|^p/p$ jest klasy $C^1$ na $\R^d$
dla każdego $p>1$ i ma gradient
$|\xi|^{p-2}\xi$ (w zerze przyjmujemy wartość $0$).
W celu uzasadnienia zamiany granicy z całkowaniem
ustalmy $|t|\le1$. Twierdzenie o wartości średniej daje
\[
\left|\frac{|\nabla u+t\nabla v|^p-|\nabla u|^p}{pt}\right|
\le C_p\big(|\nabla u|^{p-1}|\nabla v|+|\nabla v|^p\big).
\]
Funkcja po prawej jest całkowalna: pierwszy składnik
należy do $L^1$ dzięki H\"olderowi z wykładnikami
$p'$ i $p$, a drugi dzięki $v\in W^{1,p}$.
Twierdzenie o zbieżności zdominowanej daje
\[ \psi'(0)=\int_\Omega |\nabla u|^{p-2}\nabla u
\cdot\nabla v\,\mathrm d\lambda_d(x)
-\int_\Omega fv\,\mathrm d\lambda_d(x). \]
Ponieważ $u$ jest minimizatorem, $\psi'(0)=0$
dla każdego $v\in X$, co daje słabe równanie.

\textbf{Krok 4: jednoznaczność.}
Funkcja $\xi\mapsto|\xi|^p/p$ jest ściśle wypukła dla
$p>1$. Jej gradient jest zatem ściśle monotoniczny:
$(|\xi|^{p-2}\xi-|\eta|^{p-2}\eta)\cdot(\xi-\eta)>0$
dla $\xi\ne\eta$. Jeśli $u_1,u_2$ są rozwiązaniami,
testujemy różnicę ich równań funkcją $v=u_1-u_2$.
Wtedy całka z nieujemnej funkcji jest równa $0$,
więc $\nabla u_1=\nabla u_2$ prawie wszędzie.
Nierówność Poincar\'ego na $X$ daje $u_1=u_2$.

\textbf{Krok 5: oszacowanie.}
W słabym równaniu podstawiamy $v=u$.
Dostajemy $\|\nabla u\|_p^p=\int_\Omega fu
\,\mathrm d\lambda_d\le C_P\|f\|_{p'}\|\nabla u\|_p$.
Dla $\|\nabla u\|_p>0$ dzielimy przez tę normę;
przypadek zerowy jest oczywisty.
\end{proof}

\begin{uwaga}[Dlaczego pominięto przypadek $p=1$?]
Przy $p=1$ energia $\int_\Omega|\nabla u|\,\mathrm d\lambda_d(x)$ nie jest ściśle wypukła,
a $W_0^{1,1}$ nie jest refleksywna. Nie można zatem
bezkrytycznie powtórzyć dowodu; pojawiają się przestrzenie
funkcji o ograniczonej wariacji ($BV$) oraz inna postać
równań Eulera--Lagrange'a.
\end{uwaga}

\section{Metoda Galerkina i pierwsze oszacowanie błędu}
\label{sec:pde11-galerkin}

\begin{twierdzenie}[Lemat C\'ei]
\label{thm:pde11-cea}
Niech $H$ będzie rzeczywistą przestrzenią Hilberta,
$a:H\times H\to\R$ formą ograniczoną ze stałą $M$
i koercywną ze stałą $\alpha>0$, a $F\in H^*$.
Niech $H_h\subset H$ będzie skończeniewymiarową
podprzestrzenią. Jeśli $u\in H$ spełnia
$a(u,v)=F(v)$ dla wszystkich $v\in H$, a $u_h\in H_h$
spełnia $a(u_h,v_h)=F(v_h)$ dla wszystkich $v_h\in H_h$,
to $u_h$ istnieje i jest jednoznaczne oraz
\[
\|u-u_h\|_H\le\frac M\alpha
\inf_{w_h\in H_h}\|u-w_h\|_H.
\]
Gdy $a$ jest symetryczna, $u_h$ jest rzutem ortogonalnym
$u$ na $H_h$ w iloczynie skalarnym $a$, a więc
$\|u-u_h\|_a=\inf_{w_h\in H_h}\|u-w_h\|_a$.
\end{twierdzenie}
\begin{proof}
Ograniczenie formy $a$ do $H_h\times H_h$ zachowuje
ograniczoność i koercywność. Z twierdzenia Laxa--Milgrama
na $H_h$ otrzymujemy dokładnie jedno $u_h$.
Odejmując równania, dostajemy \emph{ortogonalność Galerkina}:
$a(u-u_h,v_h)=0$ dla wszystkich $v_h\in H_h$.
Niech $e=u-u_h$ i $w_h\in H_h$ będzie dowolny.
Ponieważ $w_h-u_h\in H_h$, mamy
\begin{align*}
\alpha\|e\|_H^2
&\le a(e,e)=a(e,u-w_h)+a(e,w_h-u_h)\\
&=a(e,u-w_h)\le M\|e\|_H\|u-w_h\|_H.
\end{align*}
Po podzieleniu przez $\alpha\|e\|_H$ i wzięciu infimum
po $w_h$ otrzymujemy tezę; dla $e=0$ jest oczywista.
Przy symetrii $a$ jest iloczynem skalarnym.
Ortogonalność Galerkina staje się
$(u-u_h,v_h)_a=0$, więc twierdzenie Pitagorasa
daje $\|u-w_h\|_a^2=\|u-u_h\|_a^2+
\|u_h-w_h\|_a^2$ dla wszystkich $w_h\in H_h$,
z czego wynika minimalność błędu.
\end{proof}

\begin{wniosek}[Zbieżność przy zagęszczaniu przestrzeni]
\label{cor:pde11-galerkin-density}
Przy założeniach twierdzenia~\ref{thm:pde11-cea},
niech $H_n\subset H$ będą skończeniewymiarowymi
podprzestrzeniami takimi, że
$\inf_{w\in H_n}\|u-w\|_H\to0$ dla rozważanego rozwiązania $u$.
Wówczas $\|u-u_n\|_H\to0$, gdzie $u_n$ jest rozwiązaniem otrzymanym metodą Galerkina
w $H_n$.
\end{wniosek}
\begin{proof}
Prawa strona oszacowania C\'ei dąży do zera.
Nie jest potrzebne założenie, że $H_n\subset H_{n+1}$,
lecz sama gęstość \emph{sumy} podprzestrzeni bez
kontroli dla kolejnych $n$ nie wystarczyłaby.
\end{proof}

\section{Metoda elementów skończonych na odcinku}
\label{sec:pde11-fem}

Pokażemy, że teoria wariacyjna prowadzi do konkretnego układu
macierzowego. Wróćmy do $-u''=1$ na $(0,1)$ z
warunkami $u(0)=u(1)=0$. Dzielimy $[0,1]$ na
$N\ge2$ równych pododcinków długości $h=1/N$,
a węzły oznaczamy $x_i=ih$, $i=0,\ldots,N$.

Niech $V_h\subset H_0^1(0,1)$ będzie przestrzenią funkcji
ciągłych, odcinkami afinicznych, zerujących się na końcach.
Funkcje bazowe $\phi_i$, $i=1,\ldots,N-1$, spełniają
$\phi_i(x_j)=\delta_{ij}$. Przyjmując
$u_h=\sum_{i=1}^{N-1}U_i\phi_i$,
otrzymujemy układ $KU=b$, gdzie
\[
K_{ij}=\int_0^1\phi_i'\phi_j'\,\mathrm d\lambda_1,
\qquad b_i=\int_0^1\phi_i\,\mathrm d\lambda_1.
\]
Rachunek na dwóch przedziałach, na których każda
$\phi_i$ jest niezerowa, daje
$K_{ii}=2/h$, $K_{i,i+1}=K_{i+1,i}=-1/h$,
$K_{ij}=0$ dla $|i-j|>1$ oraz $b_i=h$.

\begin{stwierdzenie}[Rozwiązanie i dokładne błędy na siatce]
\label{prop:pde11-fem-errors}
Przy powyższych definicjach rozwiązaniem macierzowym jest
$U_i=u(x_i)=x_i(1-x_i)/2$, $i=1,\ldots,N-1$.
Jeśli $u_h$ jest funkcją odcinkami afiniczną o tych wartościach,
to
\[
\|u'-u_h'\|_{L^2(0,1)}=\frac h{\sqrt{12}},
\qquad
\|u-u_h\|_{L^2(0,1)}=\frac{h^2}{\sqrt{120}}.
\]
\end{stwierdzenie}
\begin{proof}
Niech $U_0=U_N=0$ oraz
$U_i=(ih)(1-ih)/2$. Ponieważ $U_i$ jest
kwadratowym wielomianem w $i$, mamy
$-U_{i-1}+2U_i-U_{i+1}=h^2$.
Po pomnożeniu przez $1/h$ dostajemy
$-U_{i-1}/h+2U_i/h-U_{i+1}/h=h=b_i$,
więc wektor $U$ rozwiązuje układ. Dodatnia określoność
macierzy $K$ wynika z
$U^TKU=\int_0^1|\sum_iU_i\phi_i'|^2\,
\mathrm d\lambda_1>0$ dla $U\ne0$;
zatem rozwiązanie jest jednoznaczne.

Na komórce $I_i=[x_i,x_{i+1}]$ wielomian $u$
ma pochodną $u'(x)=1/2-x$, a pochodna interpolantu
$u_h'$ jest stała i równa $u'(m_i)$,
gdzie $m_i=(x_i+x_{i+1})/2$.
W konsekwencji
\[
\int_{I_i}|u'-u_h'|^2\,\mathrm d\lambda_1(x)
=\int_{-h/2}^{h/2}s^2\,\mathrm d\lambda_1(s)=h^3/12.
\]
Ponieważ jest $N=1/h$ komórek, kwadrat całkowitego błędu
jest równy $h^2/12$.
Dalej na $I_i$ zachodzi dokładnie
$u(x)-u_h(x)=(x-x_i)(x_{i+1}-x)/2$.
Dlatego
\[
\int_{I_i}|u-u_h|^2\,\mathrm d\lambda_1(x)
=\frac14\int_0^h s^2(h-s)^2\,\mathrm d\lambda_1(s)
=\frac{h^5}{120}.
\]
Po zsumowaniu $N$ identycznych wkładów otrzymujemy
$\|u-u_h\|_2^2=h^4/120$, jak w tezie.
\end{proof}

\begin{center}
\begin{tikzpicture}[x=9cm,y=7cm,>=Stealth]
 \draw[->] (-.02,0)--(1.04,0) node[right] {$x$};
 \draw[->] (0,-.01)--(0,.17) node[above] {$u$};
 \draw[thick,blue,domain=0:1,samples=80] plot (\x,{\x*(1-\x)/2});
 \draw[orange!85!black,thick] (0,0)--(.25,.09375)--(.5,.125)--(.75,.09375)--(1,0);
 \foreach \x in {0,.25,.5,.75,1}{\fill[orange!85!black] (\x,{\x*(1-\x)/2}) circle[radius=.5pt];}
 \draw[blue,thick] (.13,-.033)--(.22,-.033);
 \node[black,anchor=west,font=\footnotesize] at (.23,-.033) {dokładne};
 \draw[orange!85!black,thick] (.56,-.033)--(.65,-.033);
 \node[black,anchor=west,font=\footnotesize] at (.66,-.033) {Galerkin};
\end{tikzpicture}

{\small\emph{Rysunek:} Dokładne rozwiązanie $-u''=1$ oraz
rozwiązanie Galerkina dla czterech elementów liniowych.}
\end{center}

\section{Zakres zastosowań i granice metody}
\label{sec:pde11-limits}

\begin{uwaga}[Koercywność jest warunkiem wystarczającym]
Twierdzenie Laxa--Milgrama nie obejmuje każdej formy eliptycznej.
Na przykład forma $a(u,v)=\int_\Omega\nabla u\cdot\nabla v
-\lambda\int_\Omega uv$ może przestać być koercywna
przy dużym $\lambda$. Nie oznacza to automatycznie braku
rozwiązania: zagadnienie wymaga analizy widma lub
alternatywy Fredholma. Zbadamy to w rozdziale 12.
\end{uwaga}

\begin{uwaga}[Od istnienia do regularności]
W dowodzie słabego zagadnienia Dirichleta nie założyliśmy
gładkości współczynników $A$ ani brzegu wyższej niż
Lipschitzowska. To wystarcza do istnienia rozwiązania
w $H_0^1$, ale nie gwarantuje jego regularności $H^2$.
W rozdziałach o regularności eliptycznej będziemy analizować,
które dodatkowe założenia pozwalają odzyskać drugie
pochodne i oszacowania wyższych rzędów.
\end{uwaga}

\begin{uwaga}[Co zostanie rozwinięte dalej]
Założenie $c\ge0$ i brak wyrazów pierwszego rzędu
upraszczają dowód. Dla członów dryftu potrzebne bywają
oszacowania perturbacyjne lub warunki dodatkowe.
W następnym rozdziale przejdziemy do operatorów nieograniczonych,
teorii spektralnej i półgrup; wcześniej otrzymane twierdzenia
pozwolą m.in. skonstruować odwrotność laplasjanu Dirichleta.
\end{uwaga}

% ==============================================================
% ROZDZIAŁ 12. OPERATORY NIEOGRANICZONE I METODY SPEKTRALNE
% ==============================================================
\chapter[Operatory nieograniczone i teoria spektralna]{Operatory nieograniczone, samosprzężoność, spektrum i półgrupy}
\label{chap:pde12}

W rozdziale~\ref{chap:pde-variational} zagadnienie Dirichleta rozwiązaliśmy
w przestrzeni energetycznej $H_0^1(\Omega)$. Teraz zapytamy, kiedy
formalny operator $-\Delta$ jest operatorem samosprzężonym w
$L^2(\Omega)$, jak wyglądają jego wartości własne i dlaczego
jego widmo steruje ewolucją cieplną i falową. Dalsze rachunki
zostaną zapisane najpierw dla rzeczywistych przestrzeni Hilberta.
W argumentach z operatorami unitarnymi i widmem zespolonym
korzystamy z ich zespolenia. W zespolonej przestrzeni Hilberta
iloczyn skalarny przyjmujemy liniowy w pierwszym argumencie.

\section{Dziedzina jest częścią operatora}
\label{sec:pde12-domain}

\begin{definicja}[Operator liniowy, domknięty i domykalny]
\label{def:pde12-operator}
Niech $H$ będzie przestrzenią Hilberta. Operatorem liniowym
$T$ w $H$ nazywamy parę składającą się z podprzestrzeni
$D(T)\subset H$, zwanej \emph{dziedziną}, oraz odwzorowania
liniowego $T:D(T)\to H$. Wykresem $T$ nazywamy
$\mathcal G(T)=\{(u,Tu):u\in D(T)\}\subset H\times H$.
Operator $T$ jest \emph{domknięty}, jeżeli $\mathcal G(T)$ jest
domknięty w $H\times H$. Jest \emph{domykalny}, jeżeli
domknięcie $\mathcal G(T)$ jest wykresem pewnego operatora;
operator ten oznaczamy przez $\overline T$.
\end{definicja}

\begin{lemat}[Kryterium domkniętości i norma grafowa]
\label{lem:pde12-graph}
Dla operatora liniowego $T:D(T)\subset H\to H$
następujące warunki są równoważne:
(i) $T$ jest domknięty; (ii) ilekroć $u_n\in D(T)$,
$u_n\to u$ i $Tu_n\to v$ w $H$, to $u\in D(T)$ oraz $Tu=v$;
(iii) przestrzeń $D(T)$ z normą
$\|u\|_{D(T)}=(\|u\|_H^2+\|Tu\|_H^2)^{1/2}$ jest zupełna.
\end{lemat}
\begin{proof}
Zbieżność $(u_n,Tu_n)\to(u,v)$ w iloczynie $H\times H$
jest równoważna dwóm zbieżnościom z (ii), więc (i) i (ii)
są równoważne. Odwzorowanie $u\mapsto(u,Tu)$ jest izometrią
$D(T)$ z normą grafową na $\mathcal G(T)$ z normą iloczynową.
Domknięta podprzestrzeń przestrzeni zupełnej jest zupełna.
Odwrotnie, podprzestrzeń zupełna w przestrzeni metrycznej jest
domknięta: granica ciągu jej elementów jest granicą ciągu
Cauchy'ego w tej podprzestrzeni. To daje (i)$\Leftrightarrow$(iii).
\end{proof}

\begin{przyklad}[Pochodna i dwie różne dziedziny]
\label{ex:pde12-derivative}
W $H=L^2(0,1;\C)$ rozważmy operator $D_0u=-iu'$ z dziedziną
$D(D_0)=H_0^1(0,1;\C)$. Jest on domknięty, bo norma grafowa
jest równoważna normie $H^1$ na $H_0^1$. Jest też symetryczny:
całkowanie przez części i zerowe ślady dają
$(D_0u,v)=(u,D_0v)$ dla $u,v\in D(D_0)$.
Nie jest jednak samosprzężony. W istocie
$D(D_0^*)=H^1(0,1;\C)$ i $D_0^*v=-iv'$.
Aby wykazać to ostatnie, dla $v\in D(D_0^*)$ podstawiamy
najpierw $u\in C_c^\infty(0,1)$: istnienie $g\in L^2$
takiego, że $\int_0^1(-iu')\overline v\,\mathrm d\lambda_1
=\int_0^1u\overline g\,\mathrm d\lambda_1$ oznacza
$v'\in L^2$ dystrybucyjnie, zatem $v\in H^1$ i $g=-iv'$.
Odwrotnie, dla każdego $v\in H^1$ przybliżamy
$u\in H_0^1$ funkcjami testowymi i całkujemy przez części:
wyraz brzegowy znika. W szczególności stała funkcja $1$
należy do $D(D_0^*)$, ale nie do $D(D_0)$.
\end{przyklad}

\begin{definicja}[Operator sprzężony i samosprzężoność]
\label{def:pde12-adjoint}
Niech $T:D(T)\subset H\to H$ ma gęstą dziedzinę.
Element $v\in H$ należy do $D(T^*)$ wtedy i tylko wtedy,
gdy istnieje $g\in H$ takie, że
$(Tu,v)_H=(u,g)_H$ dla wszystkich $u\in D(T)$.
Wówczas $T^*v=g$; gęstość $D(T)$ zapewnia
jednoznaczność $g$. Operator jest \emph{symetryczny}, gdy
$(Tu,v)=(u,Tv)$ dla wszystkich $u,v\in D(T)$,
czyli $T\subset T^*$. Jest \emph{samosprzężony}, gdy
$T=T^*$, łącznie z równością dziedzin.
\end{definicja}

\begin{uwaga}
Symetria nie gwarantuje samosprzężoności; przykład
\ref{ex:pde12-derivative} pokazuje to nawet dla zamkniętego
operatora różniczkowego pierwszego rzędu. Warunki brzegowe
mogą radykalnie zmienić dziedzinę operatora sprzężonego,
a więc także jego widmo i ewolucję.
\end{uwaga}

\section{Konstrukcja laplasjanu Dirichleta przez twierdzenie Laxa--Milgrama}
\label{sec:pde12-dirichlet}

Od tej chwili, aż do końca sekcji~\ref{sec:pde12-compact-spectrum},
niech $\Omega\subset\R^d$ będzie obszarem ograniczonym Lipschitzowskim, $H=L^2(\Omega;\R)$,
$V=H_0^1(\Omega;\R)$, a $V$ niech będzie wyposażone w normę
$\|u\|_V^2=\|\nabla u\|_2^2+\|u\|_2^2$.
Rozważmy formę $b(u,v)=\int_\Omega(\nabla u\cdot\nabla v+uv)
\,\mathrm d\lambda_d(x)$.

\begin{twierdzenie}[Rezolwenta z twierdzenia Laxa--Milgrama]
\label{thm:pde12-resolvent}
Dla każdego $f\in H$ istnieje dokładnie jedno $u=Bf\in V$
takie, że $b(u,v)=(f,v)_H$ dla wszystkich $v\in V$.
Operator $B:H\to H$ jest liniowy, ograniczony, zwarty,
samosprzężony, dodatni i różnowartościowy. Ponadto
$\|B\|_{H\to H}\le1$ oraz $\overline{\operatorname{Ran}B}=H$.
\end{twierdzenie}
\begin{proof}
\textbf{Istnienie i oszacowanie.} Forma $b$ jest
iloczynem skalarnym na $V$, a więc jest ciągła i koercywna
ze stałą $1$. Funkcjonał $v\mapsto(f,v)_H$ jest ciągły na $V$:
$|(f,v)_H|\le\|f\|_2\|v\|_V$. Z twierdzenia
Laxa--Milgrama~\ref{thm:pde11-LM} otrzymujemy jedyne $u$.
Przy $v=u$ mamy
$\|u\|_V^2=(f,u)_H\le\|f\|_2\|u\|_2\le\|f\|_2\|u\|_V$.
Stąd $\|u\|_V\le\|f\|_2$, a tym bardziej
$\|Bf\|_2\le\|f\|_2$. Liniowość $B$ wynika z
jednoznaczności rozwiązania.

\textbf{Zwarty charakter.} Jeśli $(f_n)$ jest ograniczony w $H$,
to $(Bf_n)$ jest ograniczony w $V$. Twierdzenie
Rellicha--Kondrachova~\ref{thm:pde10-Rellich-p} (w postaci
zanurzenia $H_0^1(\Omega)\hookrightarrow L^2(\Omega)$)
dostarcza podciągu zbieżnego w $H$. To znaczy, że $B$
jest zwarty.

\textbf{Samosprzężoność i dodatniość.} Dla $u=Bf$, $w=Bg$
możemy wybrać odpowiednio $v=w$ i $v=u$ w równaniach
wariacyjnych. Symetria $b$ daje
$(Bf,g)_H=b(u,w)=(f,Bg)_H$, zatem $B=B^*$.
Dalej $(Bf,f)_H=b(Bf,Bf)=\|Bf\|_V^2\ge0$.

\textbf{Jądro i gęstość obrazu.} Gdy $Bf=0$, jego równanie
wariacyjne daje $(f,v)_H=0$ dla wszystkich $v\in V$.
Ponieważ $C_c^\infty(\Omega)\subset V$ jest gęste w $H$,
to $f=0$. Wreszcie
$(\operatorname{Ran}B)^\perp=\ker B^*=\ker B=\{0\}$,
a zatem obraz $B$ jest gęsty.
\end{proof}


\begin{uwaga}[Z czego wynika zwartość rezolwenty?]
\label{rem:pde-v42-compact-chain}
Zależność logiczna w poprzednim dowodzie jest istotna:
koercywność formy i twierdzenie Laxa--Milgrama
(\ref{thm:pde11-LM}) zapewniają oszacowanie
$B:L^2(\Omega)\to H_0^1(\Omega)$, a dopiero
zwartość zanurzenia Rellicha--Kondrachova
(\ref{thm:pde10-Rellich-p}) zapewnia zwartość
$B:L^2(\Omega)\to L^2(\Omega)$. Samo istnienie
ograniczonej rezolwenty nie daje zwartości.
W szczególności ograniczoność obszaru i warunki brzegowe
odgrywają rolę przy późniejszym rozkładzie spektralnym.
\end{uwaga}

\begin{definicja}[Laplasjan Dirichleta w $L^2$]
\label{def:pde12-laplacian}
Operatorem Dirichleta $A=-\Delta_D$ nazywamy operator w $H$
o dziedzinie
\[
D(A)=\big\{u\in V:\ \exists g\in H\quad
\int_\Omega\nabla u\cdot\nabla v\,\mathrm d\lambda_d(x)
=(g,v)_H\ \text{dla każdego }v\in V\big\},
\]
przy czym $Au=g$. Element $g$ jest wyznaczony
jednoznacznie dzięki gęstości $V$ w $H$.
Na funkcjach testowych $Au=-\Delta u$ w sensie dystrybucji;
warunek $u\in H_0^1(\Omega)$ narzuca zerowy ślad.
\end{definicja}

\begin{twierdzenie}[Samosprzężoność laplasjanu Dirichleta]
\label{thm:pde12-A-selfadjoint}
Operator $A$ z definicji~\ref{def:pde12-laplacian} jest
gęsto określony, nieujemny i samosprzężony.
Ponadto $D(A)=\operatorname{Ran}B$, a na tej dziedzinie
$A=B^{-1}-I$, więc $(A+I)^{-1}=B$.
\end{twierdzenie}
\begin{proof}
Jeśli $u=Bf$, to
$\int_\Omega\nabla u\cdot\nabla v\,\mathrm d\lambda_d(x)
=(f-u,v)_H$, czyli $u\in D(A)$ i $Au=f-u$.
Odwrotnie, jeśli $u\in D(A)$ i $Au=g$, to
$b(u,v)=(g+u,v)_H$, a z jednoznaczności równania
dla $B$ wynika $u=B(g+u)$. Tak otrzymujemy
dziedzinę, wzór na $A$ i gęstość dziedziny.
Wybór $v=u$ w definicji daje
$(Au,u)_H=\|\nabla u\|_2^2\ge0$.

Aby wykazać samosprzężoność bez ukrywania dziedziny,
zauważmy, że $A$ jest symetryczny:
$(Au,v)_H=(u,Av)_H$ dla $u,v\in D(A)$
przez symetrię całki z gradientów. Jeśli $w\in D(A^*)$,
to dla każdego $u=Bf\in D(A)$ mamy
$(Au,w)_H=(u,A^*w)_H$, a więc
$(f-Bf,w)_H=(Bf,A^*w)_H=(f,BA^*w)_H$.
Samosprzężoność $B$ daje
$(f,w)_H=(f,Bw+BA^*w)_H$ dla wszystkich $f\in H$.
Zatem $w=B(w+A^*w)\in\operatorname{Ran}B=D(A)$.
Otrzymaliśmy $D(A^*)\subset D(A)$; inkluzja przeciwna
wynika z symetrii. Tym samym $A=A^*$.
\end{proof}

\begin{uwaga}[Dziedzina nie zawsze jest $H^2\cap H_0^1$]
\label{rem:pde12-corners}
Jeśli brzeg $\partial\Omega$ jest klasy $C^{1,1}$, klasyczna
regularność eliptyczna daje $D(A)=H^2(\Omega)\cap H_0^1(\Omega)$.
Dla ogólnego obszaru Lipschitzowskiego równość może zawieść.
Na przykład w płaskim wycinku o kącie wklęsłym
$\omega\in(\pi,2\pi)$,
funkcja $r^{\pi/\omega}\sin(\pi\theta/\omega)$ jest
harmoniczna poza wierzchołkiem, należy lokalnie do $H^1$,
ale nie do $H^2$. Mnożąc ją przez gładkie obcięcie radialne
równe $1$ przy wierzchołku i $0$ przy zewnętrznym łuku,
dostajemy $u\in D(A)$, lecz $u\notin H^2(\Omega)$:
$\Delta u$ ma nośnik z dala od wierzchołka i należy do $L^2$,
a drugie pochodne nie są całkowalne z kwadratem w pobliżu
wierzchołka. Dlatego nie wolno identyfikować operatora
z samym wyrażeniem $-\Delta$.
\end{uwaga}

\section{Twierdzenie spektralne dla zwartego operatora dodatniego}
\label{sec:pde12-compact-spectrum}

\begin{twierdzenie}[Rozkład spektralny zwartego operatora dodatniego]
\label{thm:pde12-compact-spectral}
Niech $H$ będzie nieskończeniewymiarową, rzeczywistą lub
zespoloną, separowalną przestrzenią Hilberta, a $K:H\to H$
operatorem zwartym, samosprzężonym i nieujemnym.
Wówczas istnieje ortonormalny układ wektorów własnych
$(e_j)$ odpowiadających dodatnim wartościom własnym
$\mu_j$, uporządkowanym nierosnąco i z uwzględnieniem
krotności. Jeśli wartości własnych jest nieskończenie wiele,
to $\mu_j\to0$. Ponadto
$H=\ker K\oplus\overline{\operatorname{span}}\{e_j\}$,
a dla $x\in H$ zachodzi
$Kx=\sum_j\mu_j(x,e_j)e_j$ zbieżnie w $H$.
Jeżeli $K$ jest różnowartościowy, $(e_j)$ jest bazą
ortonormalną całego $H$.
\end{twierdzenie}
\begin{proof}
\textbf{Krok 1: maksimum ilorazu Rayleigha.}
Dla $K\ne0$ mamy $m=\sup_{\|x\|=1}(Kx,x)=\|K\|>0$.
Aby uzasadnić tę równość, połóżmy
$m_0=\sup_{\|x\|=1}(Kx,x)\le\|K\|$.
Dodatniość daje nierówność Cauchy'ego--Schwarza
dla formy $(Kx,y)$:
$|(Kx,y)|^2\le(Kx,x)(Ky,y)$.
Dla wektorów jednostkowych $x,y$ otrzymujemy
$|(Kx,y)|\le m_0$; biorąc supremum po $y$,
dostajemy $\|Kx\|\le m_0$, a następnie
$\|K\|\le m_0$. Zatem $m=m_0=\|K\|$.
Wybierzmy $\|x_n\|=1$ z $(Kx_n,x_n)\to m$.
Z ograniczoności i słabej zwartości kul w przestrzeni Hilberta
możemy przejść do podciągu $x_n\rightharpoonup x$.
Zwarty operator przenosi zbieżność słabą ciągu w zbieżność
silną: każda podgranica $Kx_{n_k}$ równa się $Kx$,
gdyż $(Kx_{n_k},y)=(x_{n_k},Ky)\to(x,Ky)$.
Stąd $Kx_n\to Kx$ i $(Kx_n,x_n)\to(Kx,x)=m$.
Ponieważ $(Kx,x)\le m\|x\|^2$ i $\|x\|\le1$,
otrzymujemy $\|x\|=1$.

\textbf{Krok 2: maksimum daje wektor własny.}
Dla $y\perp x$ rozpatrujemy
$\varphi(t)=\frac{(K(x+ty),x+ty)}{\|x+ty\|^2}$.
Funkcja osiąga maksimum przy $t=0$, więc
$0=\varphi'(0)=2\operatorname{Re}(Kx,y)$.
W przypadku zespolonym powtarzamy argument dla $iy$
i otrzymujemy również $\operatorname{Im}(Kx,y)=0$.
Wynika stąd, że $Kx$ jest
wielokrotnością $x$, a z $(Kx,x)=m$ dostajemy $Kx=mx$.

\textbf{Krok 3: iteracja na ortogonalnym dopełnieniu.}
Przestrzeń $x^\perp$ jest niezmiennicza względem $K$,
bo $(Ky,x)=(y,Kx)=m(y,x)=0$. Ograniczenie $K$ do niej
pozostaje zwarte, samosprzężone i dodatnie. Powtarzamy
konstrukcję, dopóki ograniczenie nie stanie się zerowe.
Powstają ortonormalne $e_j$ i dodatnie nierosnące $\mu_j$.
Jeśli $\mu_j$ nie dążyłyby do zera, istniałoby $c>0$
i nieskończenie wiele indeksów z $\|Ke_j\|=\mu_j\ge c$.
Ponieważ wektory $Ke_j=\mu_je_j$ są parami ortogonalne,
nie mają podciągu Cauchy'ego, co przeczy zwartości $K$.

\textbf{Krok 4: zupełność na $(\ker K)^\perp$.}
Jeżeli $z\perp e_j$ dla wszystkich $j$, to $z$
należy do kolejnych ortogonalnych dopełnień, na których
norma ograniczenia $K$ jest co najwyżej $\mu_{j+1}$.
Zatem $\|Kz\|\le\mu_{j+1}\|z\|\to0$;
gdy konstrukcja skończyła się po skończonej liczbie kroków,
operator na ostatnim dopełnieniu jest z definicji zerowy.
W obu przypadkach $Kz=0$. To daje rozkład ortogonalny.
Dla skończonych sum częściowych wzór na $K$ wynika
z liniowości; przejście do granicy jest dozwolone
z ograniczoności $K$. Przy $\ker K=\{0\}$ układ
jest więc bazą całego $H$.
\end{proof}

\begin{twierdzenie}[Twierdzenie spektralne dla zwartego operatora samosprzężonego]
\label{thm:pde12-compact-general-spectral}
Niech $H$ będzie rzeczywistą lub zespoloną separowalną
przestrzenią Hilberta, a $K\in\mathcal L(H)$ operatorem
zwartym i samosprzężonym, niekoniecznie dodatnim.
Wówczas $H$ jest domkniętą sumą ortogonalną $\ker K$ oraz
podprzestrzeni własnych odpowiadających niezerowym
rzeczywistym wartościom własnym $\kappa_j$.
Każda niezerowa wartość własna ma skończoną krotność,
a jeśli jest ich nieskończenie wiele (licząc z krotnościami),
to $|\kappa_j|\to0$. Istnieje ortonormalny układ
$(e_j)$ tych wektorów własnych taki, że
\[
 Kx=\sum_j\kappa_j(x,e_j)_He_j\quad(x\in H),
\]
gdzie szereg jest zbieżny w $H$; jeżeli $\ker K=\{0\}$,
to $(e_j)$ jest bazą ortonormalną całego $H$.
\end{twierdzenie}
\begin{proof}
Operator $K^2$ jest zwarty, samosprzężony i nieujemny:
$K^2$ jest zwarty jako złożenie operatora zwartego
i ograniczonego, a $(K^2x,x)=\|Kx\|^2\ge0$.
Ponadto $\ker K^2=\ker K$. Twierdzenie
\ref{thm:pde12-compact-spectral} zastosowane do $K^2$
(dla przestrzeni skończeniewymiarowej używamy zwykłego
twierdzenia o diagonalizacji macierzy samosprzężonej)
daje rozkład ortogonalny
$H=\ker K\oplus\overline{\bigoplus_{\mu>0}E_\mu}$,
gdzie $E_\mu=\ker(K^2-\mu I)$, $\mu>0$.
Każde $E_\mu$ jest skończeniewymiarowe, gdyż nieskończony
układ ortonormalny w $E_\mu$ miałby obrazy przez $K^2$
bez podciągu zbieżnego, co przeczyłoby zwartości.

Ponieważ $K$ komutuje z $K^2$, każda przestrzeń $E_\mu$
jest niezmiennicza względem $K$. Ograniczenie
$K|_{E_\mu}$ jest macierzą samosprzężoną, można je więc
zdiagonalizować w bazie ortonormalnej $E_\mu$.
Jeżeli $Ke=\kappa e$ i $e\in E_\mu\setminus\{0\}$,
to $K^2e=\kappa^2e=\mu e$, zatem
$\kappa=\pm\sqrt\mu$; w przypadku zespolonym
$\kappa$ jest rzeczywista dzięki samosprzężoności.
Bazy wszystkich $E_\mu$ tworzą wraz z bazą jądra
ortogonalny rozkład $H$ na wektory własne $K$.
Niezerowe wartości własne $K^2$ mogą gromadzić się
tylko przy zerze, więc to samo dotyczy $|\kappa|$.
Dla sum skończonych podany wzór na $K$ wynika
z diagonalizacji. Dla dowolnego $x\in H$ przechodzimy
do granicy, ponieważ rzuty na skończone sumy podprzestrzeni
własnych zbiegają silnie do rzutu na $(\ker K)^\perp$,
a $K$ jest ograniczony i znika na swoim jądrze.
\end{proof}

\begin{wniosek}[Dyskretne widmo laplasjanu Dirichleta]
\label{cor:pde12-eigen}
Dla obszaru ograniczonego Lipschitzowskiego $\Omega$ operator
$A=-\Delta_D$ ma bazę ortonormalną $(\phi_j)_{j\ge1}$ w
$L^2(\Omega)$ złożoną z wektorów własnych:
$A\phi_j=\lambda_j\phi_j$, gdzie
$0<\lambda_1\le\lambda_2\le\cdots$, $\lambda_j\to\infty$.
Każda wartość własna ma skończoną krotność,
a w sensie słabym
$\int_\Omega\nabla\phi_j\cdot\nabla v\,\mathrm d\lambda_d(x)
=\lambda_j\int_\Omega\phi_jv\,\mathrm d\lambda_d(x)$
dla wszystkich $v\in H_0^1(\Omega)$.
\end{wniosek}
\begin{proof}
Dla operatora $B=(A+I)^{-1}$ z
\ref{thm:pde12-resolvent} jądro jest zerowe,
więc poprzednie twierdzenie daje bazę $(\phi_j)$
i liczby $\mu_j>0$ z $B\phi_j=\mu_j\phi_j$.
Ponieważ $B$ działa na nieskończeniewymiarowym $L^2$
i jest różnowartościowy, wartości tych jest nieskończenie
wiele, a $\mu_j\downarrow0$. Z równości
$A=B^{-1}-I$ otrzymujemy
$\lambda_j=\mu_j^{-1}-1$. Nieujemność $A$ daje
$\lambda_j\ge0$. Gdyby $\lambda_j=0$, to
$\|\nabla\phi_j\|_2^2=(A\phi_j,\phi_j)=0$.
Nierówność Poincar\'ego w $H_0^1(\Omega)$ daje
$\phi_j=0$, sprzeczność z normalizacją.
Zatem $\lambda_j>0$. Każdy niezerowy punkt widma
zwartego operatora ma skończoną krotność; w powyższym
dowodzie wynika to z niemożliwości istnienia nieskończonej
rodziny ortonormalnej o tym samym niezerowym obrazie normowym.
Tożsamość słaba wynika z definicji $A$.
\end{proof}

\section{Obszar formowa, zasada min--max i oszacowania widma}
\label{sec:pde12-minmax}

\begin{twierdzenie}[Rozwinięcie spektralne formy energii]
\label{thm:pde12-formdomain}
Niech $(\lambda_j,\phi_j)$ będą wartościami i funkcjami
własnymi z wniosku~\ref{cor:pde12-eigen}. Dla $u\in L^2(\Omega)$
niech $c_j=(u,\phi_j)_{L^2}$. Wówczas
\begin{align*}
u\in H_0^1(\Omega)
&\quad\Longleftrightarrow\quad
\sum_{j=1}^\infty\lambda_j|c_j|^2<\infty,\\
u\in D(A)
&\quad\Longleftrightarrow\quad
\sum_{j=1}^\infty\lambda_j^2|c_j|^2<\infty.
\end{align*}
Dla $u\in H_0^1(\Omega)$ zachodzi
$\|\nabla u\|_2^2=\sum_j\lambda_j|c_j|^2$,
a dla $u\in D(A)$ zachodzi
$Au=\sum_j\lambda_jc_j\phi_j$ zbieżnie w $L^2$.
\end{twierdzenie}
\begin{proof}
Ponieważ $\phi_j=\mu_j^{-1}B\phi_j$, mamy
$\phi_j\in H_0^1(\Omega)$. Z równania słabego
$b(\phi_i,\phi_j)=\mu_i^{-1}(\phi_i,\phi_j)_{L^2}$
wynika, że wektory $\sqrt{\mu_j}\phi_j$ tworzą
układ ortonormalny w przestrzeni $V=H_0^1(\Omega)$
z iloczynem skalarnym $b$. Układ ten jest zupełny w $V$:
jeśli $v\in V$ jest do niego ortogonalny, to
$0=b(v,\phi_j)=\mu_j^{-1}(v,\phi_j)_{L^2}$
dla wszystkich $j$; zupełność $(\phi_j)$ w $L^2$ daje $v=0$.
Z twierdzenia Parsevala w $V$ dostajemy
\[
\|v\|_V^2=\sum_j|b(v,\sqrt{\mu_j}\phi_j)|^2
=\sum_j\mu_j^{-1}|(v,\phi_j)_{L^2}|^2
=\sum_j(1+\lambda_j)|c_j|^2.
\]
Dla ogólnego $u\in L^2$ warunek po prawej gwarantuje,
że sumy częściowe $u_N=\sum_{j\le N}c_j\phi_j$
są ciągiem Cauchy'ego w $V$. Ich granica $v\in V$
równa się $u$, gdyż zbieżność w $V$ implikuje zbieżność
w $L^2$, a w $L^2$ mamy $u_N\to u$ z Parsevala.
Tak wykazaliśmy równoważność dla $H_0^1$ i równość energii
przez odjęcie $\|u\|_2^2$.

Jeśli $u\in D(A)$, to $Au\in L^2$, a dla każdego $j$
$(Au,\phi_j)=(u,A\phi_j)=\lambda_jc_j$.
Parseval daje sumowalność $\lambda_j^2|c_j|^2$ i wzór
na $Au$. Odwrotnie, jeśli ta suma jest skończona,
to $u\in H_0^1$ z pierwszej części twierdzenia.
Niech $g=\sum_j\lambda_jc_j\phi_j\in L^2$.
Dla skończonych sum funkcji własnych $v$ mamy
$\int_\Omega\nabla u\cdot\nabla v\,\mathrm d\lambda_d
=(g,v)_{L^2}$. Takie sumy są gęste w $V$ na mocy
poprzedniego akapitu; przez ciągłość obu stron otrzymujemy
równość dla wszystkich $v\in V$, więc $u\in D(A)$ i $Au=g$.
\end{proof}

\begin{twierdzenie}[Zasada Couranta--Fischera]
\label{thm:pde12-minmax}
Dla każdej liczby całkowitej $k\ge1$
\[
\lambda_k=
\min_{\substack{E\subset H_0^1(\Omega)\\\dim E=k}}
\ \max_{\substack{u\in E\\u\ne0}}
\frac{\displaystyle\int_\Omega|\nabla u|^2\,\mathrm d\lambda_d(x)}
{\displaystyle\int_\Omega|u|^2\,\mathrm d\lambda_d(x)}.
\]
Minimum jest osiągane dla $E=\operatorname{span}
\{\phi_1,\ldots,\phi_k\}$.
\end{twierdzenie}
\begin{proof}
Dla wskazanego $E$ i $u=\sum_{j\le k}c_j\phi_j$
twierdzenie~\ref{thm:pde12-formdomain} daje
$\|\nabla u\|_2^2=\sum_{j\le k}\lambda_j|c_j|^2
\le\lambda_k\sum_{j\le k}|c_j|^2=\lambda_k\|u\|_2^2$.
Wektor $u=\phi_k$ daje równość, więc maksimum na tym $E$
wynosi $\lambda_k$. Z drugiej strony weźmy dowolne
$k$-wymiarowe $E\subset H_0^1$. Układ $k-1$ równań
liniowych $(u,\phi_j)=0$, $j<k$, ma niezerowe rozwiązanie
$u\in E$ (wymiar jądra jest co najmniej $1$).
Dla takiego $u$ w rozwinięciu $L^2$ znikają wszystkie
współczynniki o indeksach mniejszych od $k$, wobec czego
$\|\nabla u\|_2^2=\sum_{j\ge k}\lambda_j|c_j|^2
\ge\lambda_k\sum_{j\ge k}|c_j|^2=\lambda_k\|u\|_2^2$.
Maksimum na dowolnym $E$ jest zatem co najmniej $\lambda_k$.
Łącząc oba oszacowania, otrzymujemy wzór.
\end{proof}

\begin{wniosek}[Pierwsza wartość własna, Poincaré i monotoniczność]
\label{cor:pde12-first}
Jeśli $\Omega\subset\R^d$ jest obszarem ograniczonym Lipschitzowskim, to
$\lambda_1=\inf_{u\in H_0^1(\Omega)\setminus\{0\}}
\|\nabla u\|_2^2/\|u\|_2^2$.
Najlepsza stała $C_P$ w nierówności
$\|u\|_2\le C_P\|\nabla u\|_2$ wynosi
$C_P=\lambda_1^{-1/2}$. Ponadto jeżeli
$\Omega_1\subset\Omega_2$ są ograniczonymi obszarami
Lipschitzowskimi, to
$\lambda_k(\Omega_2)\le\lambda_k(\Omega_1)$ dla każdego $k$.
\end{wniosek}
\begin{proof}
Pierwsze dwa stwierdzenia wynikają z
\ref{thm:pde12-minmax} dla $k=1$.
Dla drugiego rozważmy przedłużenie przez zero
$H_0^1(\Omega_1)\hookrightarrow H_0^1(\Omega_2)$;
na mocy twierdzenia o przedłużaniu przez zero przedłużenia
funkcji testowych, a następnie granic w $H^1$, zachowują
zarówno normę $L^2$, jak i normę gradientu.
Każda $k$-wymiarowa podprzestrzeń pierwszej przestrzeni
staje się więc dopuszczalną podprzestrzenią drugiej,
z tym samym maksimum ilorazu Rayleigha. Przyjmując minima
otrzymujemy podaną nierówność.
\end{proof}

\begin{przyklad}[Odcinek i prostokąt]
Na odcinku $(0,L)$ rozwiązania problemu
$-u''=\lambda u$, $u(0)=u(L)=0$ mają postać
$\phi_n(x)=\sqrt{2/L}\sin(n\pi x/L)$ oraz
$\lambda_n=(n\pi/L)^2$, $n\ge1$.
W prostokącie $(0,L_1)\times(0,L_2)$ funkcje własne są
iloczynami sinusów, z wartościami
$\lambda_{m,n}=\pi^2(m^2/L_1^2+n^2/L_2^2)$.
Wartości te należy uporządkować rosnąco z krotnościami;
na kwadracie występują wielokrotności wskutek symetrii
$\lambda_{m,n}=\lambda_{n,m}$.
\end{przyklad}

\section{Widmo, rezolwenta i alternatywa Fredholma}
\label{sec:pde12-fredholm}

\begin{definicja}[Zbiór rezolwenty i widmo]
\label{def:pde12-spectrum}
Dla domkniętego gęsto określonego operatora $T$ w
zespolonej przestrzeni Hilberta jego \emph{zbiorem rezolwenty}
$\rho(T)$ nazywamy zbiór liczb $z\in\C$, dla których
$T-zI:D(T)\to H$ jest bijekcją, a odwrotność
$(T-zI)^{-1}:H\to H$ jest ograniczona.
\emph{Widmo} to $\sigma(T)=\C\setminus\rho(T)$.
Liczba $z$ jest wartością własną, jeżeli
$0\ne u\in D(T)$ spełnia $Tu=zu$.
\end{definicja}

\begin{twierdzenie}[Alternatywa Fredholma dla laplasjanu]
\label{thm:pde12-fredholm}
Niech $A=-\Delta_D$ będzie operatorem na obszarze ograniczonym Lipschitzowskim $\Omega$, a $\zeta\in\C$.
Jeżeli $\zeta\notin\{\lambda_j:j\ge1\}$, to dla
każdego $f\in L^2(\Omega;\C)$ istnieje dokładnie jedno
$u\in D(A)$ spełniające $(A-\zeta I)u=f$, i
\[
\|u\|_2\le\frac{\|f\|_2}
{\operatorname{dist}(\zeta,\{\lambda_j\})}.
\]
Jeżeli $\zeta$ jest wartością własną, to równanie
$(A-\zeta I)u=f$ jest rozwiązywalne wtedy i tylko wtedy,
gdy $f\perp\ker(A-\zeta I)$; rozwiązanie jest wówczas
jednoznaczne z dokładnością do dodania elementu jądra.
Rezolwenta $(A-\zeta I)^{-1}$ w każdym punkcie
$\zeta\in\rho(A)$ jest zwarta, a
$\sigma(A)=\{\lambda_j:j\ge1\}$ jako zbiór (bez krotności).
\end{twierdzenie}
\begin{proof}
Po zespoleniu bazy własnej piszemy
$f=\sum_jf_j\phi_j$, a dla potencjalnego rozwiązania
$u=\sum_ju_j\phi_j$. Równanie po sparowaniu z
$\phi_j$ wymaga $(\lambda_j-\zeta)u_j=f_j$.
Jeżeli $\zeta$ nie jest wartością własną, to odległość
$\delta=\inf_j|\lambda_j-\zeta|$ jest dodatnia:
$\lambda_j\to\infty$, więc jedynie skończenie wiele
wartości własnych może leżeć w ustalonym otoczeniu
$\zeta$. Definiujemy $u_j=f_j/(\lambda_j-\zeta)$.
Z Parsevala
$\|u\|_2^2=\sum_j|f_j|^2/|\lambda_j-\zeta|^2
\le\delta^{-2}\|f\|_2^2$.
Dodatkowo liczby $\lambda_j/|\lambda_j-\zeta|$
są jednostajnie ograniczone: dla dużych $j$
wyrażenie dąży do $1$, a pierwszych wyrazów jest
skończenie wiele. Stąd $\sum\lambda_j^2|u_j|^2<\infty$,
więc $u\in D(A)$ na mocy~\ref{thm:pde12-formdomain},
i rzeczywiście rozwiązuje równanie. Każde inne rozwiązanie
ma te same współczynniki, co dowodzi jednoznaczności.

Jeśli $\zeta$ jest wartością własną, to dla indeksów
$J=\{j:\lambda_j=\zeta\}$ równanie daje warunek
$f_j=0$. Jest on dokładnie ortogonalnością do jądra.
Dla $j\notin J$ możemy określić
$u_j=f_j/(\lambda_j-\zeta)$, a współczynniki $u_j$
z $j\in J$ pozostają dowolne. Są tylko skończenie liczne;
pozostała część daje element $D(A)$ jak wyżej.

Na koniec na podprzestrzeni rozpiętej przez
$\phi_1,\ldots,\phi_N$ rezolwenta jest operatorem rzędu
skończonego. Różnica między nią a pełną rezolwentą
ma normę nie większą niż
$\sup_{j>N}|\lambda_j-\zeta|^{-1}\to0$.
Granica w normie operatorowej operatorów zwartych jest
zwarta. Pokazuje to zwartość rezolwenty i identyfikację widma.
\end{proof}

\begin{twierdzenie}[Alternatywa Fredholma dla operatorów zwartych]
\label{thm:pde12-general-fredholm}
Niech $H$ będzie zespoloną przestrzenią Hilberta,
a $K:H\to H$ operatorem zwartym. Wówczas
$T=I-K$ ma domknięty obraz oraz skończeniewymiarowe
jądra $\ker T$ i $\ker T^*$. Zachodzi
$\dim\ker T=\dim\ker T^*$, a dla $f\in H$
równanie $(I-K)u=f$ ma rozwiązanie wtedy i tylko wtedy,
gdy $f\perp\ker(I-K^*)$. Rozwiązanie jest jednoznaczne
z dokładnością do dodania elementu $\ker(I-K)$.
\end{twierdzenie}
\begin{proof}
\textbf{Krok 1: skończony wymiar jądra.}
Jeśli $\ker T$ byłoby nieskończeniewymiarowe,
wybralibyśmy w nim nieskończony układ ortonormalny
$(e_n)$. Równość $Te_n=0$ oznacza $Ke_n=e_n$.
Ponieważ $(e_n)$ nie ma podciągu zbieżnego, byłoby
to sprzeczne ze zwartością $K$. Skończoną wymiarowość
$\ker T^*$ otrzymamy w kroku 4.

\textbf{Krok 2: domkniętość obrazu.} Niech
$N=\ker T$ i $M=N^\perp$. Pokażemy, że istnieje $c>0$
takie, że $\|Tx\|\ge c\|x\|$ dla wszystkich $x\in M$.
W przeciwnym razie można wybrać $x_n\in M$,
$\|x_n\|=1$, $Tx_n\to0$. Zwarty operator $K$
daje podciąg $(Kx_{n_k})$ zbieżny w $H$.
Ponieważ $x_{n_k}=Tx_{n_k}+Kx_{n_k}$,
ten sam podciąg $x_{n_k}$ zbiega do pewnego $x\in M$
o normie $1$. Ciągłość $T$ daje $Tx=0$,
czyli $x\in M\cap N=\{0\}$; sprzeczność.
Teraz niech $Tx_n\to y$ w $H$.
Zastępując $x_n$ przez jego rzut na $M$, nie zmieniamy
$Tx_n$. Uzyskane $x_n\in M$ są Cauchy'ego,
bo $\|x_n-x_m\|\le c^{-1}\|Tx_n-Tx_m\|$.
Zatem $x_n\to x\in M$ i $Tx=y$, dowodząc domkniętości
$\operatorname{Ran}T$.

\textbf{Krok 3: równość wymiarów jąder.}
Pokażemy najpierw, że $K$ można aproksymować
w normie operatorowej operatorami skończonego rzędu.
Dla $0<\varepsilon<1$ zbiór $K(B_H)$ jest względnie
zwarty, więc można znaleźć skończony zbiór
$y_1,\dots,y_m$ w $H$ taki, że każdy $Kx$, $\|x\|\le1$,
leży w odległości mniejszej niż $\varepsilon$ od
jednego z $y_j$. Jeśli $W=\operatorname{span}\{y_j\}$,
a $P_W$ oznacza rzut ortogonalny na $W$, to
$F=P_WK$ jest skończonego rzędu i
$\|K-F\|\le\varepsilon$.
Zapiszmy $K=F+E$, gdzie $\|E\|<1$.
Szereg Neumanna dowodzi, że $I-E$ jest odwracalny,
a faktoryzacja
\[
I-K=(I-E)(I-G),\qquad G=(I-E)^{-1}F
\]
przedstawia $T$ jako iloczyn operatora odwracalnego
i operatora $I-G$ z $G$ skończonego rzędu.
Niech $W'=\operatorname{Ran}G$ (przestrzeń skończeniewymiarowa).
W rozkładzie $H=W'\oplus(W')^\perp$ operator $I-G$
ma postać blokową
\[
(I-G)(w,z)=((I-G|_{W'})w-Gz,\ z).
\]
Jego jądro to $\ker(I-G|_{W'})\oplus\{0\}$.
Obraz ma skończony kowymiar równy
$\dim W'-\dim\operatorname{Ran}(I-G|_{W'})$:
rzutowanie obrazu na $(W')^\perp$ jest surjekcją,
a przy ustalonym $z$ dopuszczalne pierwsze składowe
różnią się o element $\operatorname{Ran}(I-G|_{W'})$.
Z twierdzenia o wymiarze jądra i obrazu w $W'$
ten kowymiar równa się $\dim\ker(I-G|_{W'})$.
Mnożenie z lewej przez odwracalny $I-E$ nie zmienia
wymiaru jądra ani kowymiaru obrazu. Otrzymujemy
$\dim\ker T=\operatorname{codim}\operatorname{Ran}T$.

\textbf{Krok 4: warunek rozwiązywalności.}
Dla dowolnego ograniczonego operatora w Hilbercie
$(\operatorname{Ran}T)^\perp=\ker T^*$,
gdyż $(Tu,v)=0$ dla wszystkich $u$ jest dokładnie
warunkiem $T^*v=0$. Obraz $T$ jest domknięty
z kroku 2, stąd
$\operatorname{Ran}T=(\ker T^*)^\perp$.
To dowodzi warunku na $f$ i jednocześnie
$\dim\ker T=\dim\ker T^*$ z kroku 3.
Każde dwa rozwiązania różnią się o element jądra.
\end{proof}

\begin{uwaga}
W odróżnieniu od twierdzenia~\ref{thm:pde12-fredholm},
nie założyliśmy tutaj samosprzężoności $K$.
Dlatego w warunku rozwiązywalności pojawia się
$\ker(I-K^*)$, a niekoniecznie $\ker(I-K)$.
W terminologii teorii operatorów wynik stwierdza,
że $I-K$ jest operatorem Fredholma o indeksie $0$.
\end{uwaga}


\begin{kontrprzyklad}[Brak zwartości na całej przestrzeni]
\label{ex:pde12-unbounded-domain}
W $L^2(\R^d)$ operator $A=-\Delta$ ze standardową
dziedziną $H^2(\R^d)$ jest samosprzężony,
lecz $(A+I)^{-1}$ nie jest zwarty. Wystarczy wziąć
$0\ne f\in C_c^\infty(\R^d)$, przesunięcia
$f_n(x)=f(x-nLe_1)$ o rozłącznych nośnikach
(i odpowiednio duże $L$). Ciąg $(f_n)$ jest ograniczony
w $L^2$ i nie ma podciągu zbieżnego. Gdyby
$(A+I)^{-1}$ był zwarty, obrazy $u_n=(A+I)^{-1}f_n$
miałyby podciąg zbieżny w $L^2$; lecz $u_n$ są
przesunięciami ustalonej niezerowej funkcji $u$.
Przesunięcia $u$ dążą słabo do $0$, zachowując normę
$L^2$, i dlatego nie posiadają podciągu silnie zbieżnego.
Istotnie, słabą zbieżność najpierw sprawdzamy dla
funkcji testowych o zwartym nośniku, a potem korzystamy
z ich gęstości w $L^2$.
Na mocy transformaty Fouriera $A$ jest równoważny
operatorowi mnożenia przez $|\xi|^2$; jego widmo jest
ciągłe i równe $[0,\infty)$. Nie należy więc oczekiwać
sumowania dyskretnego po funkcjach własnych jak na
obszarze ograniczonym.
\end{kontrprzyklad}

\section{Półgrupa cieplna zbudowana z widma}
\label{sec:pde12-heat}

\begin{definicja}[Silnie ciągła półgrupa]
Niech $H$ będzie przestrzenią Banacha. Rodzinę
$S(t)\in\mathcal L(H)$, $t\ge0$, nazywamy
\emph{silnie ciągłą półgrupą}, gdy $S(0)=I$,
$S(t+s)=S(t)S(s)$ dla wszystkich $s,t\ge0$, a dla
każdego $u\in H$ zachodzi $\|S(t)u-u\|_H\to0$
przy $t\downarrow0$. Jej \emph{generator} $G$ jest
dany przez $Gu=\lim_{t\downarrow0}(S(t)u-u)/t$
na dziedzinie tych $u$, dla których granica istnieje w $H$.
\end{definicja}

\begin{twierdzenie}[Półgrupa generowana przez $-A$]
\label{thm:pde12-heat}
Niech $A=-\Delta_D$ dla obszaru ograniczonego Lipschitzowskiego $\Omega$, a $H=L^2(\Omega)$.
Dla $t\ge0$ określmy
\[
S(t)u=\sum_{j=1}^\infty e^{-\lambda_jt}(u,\phi_j)_H\phi_j.
\]
Wtedy $(S(t))_{t\ge0}$ jest silnie ciągłą półgrupą na $H$
o generatorze dokładnie $-A$ i
$\|S(t)\|_{H\to H}=e^{-\lambda_1t}$.
Dla $t>0$, $m\in\N$ i $u\in H$ zachodzi
$S(t)u\in D(A^m)$ oraz
\[
\|A^mS(t)u\|_H
\le\bigg(\frac{m}{e t}\bigg)^m\|u\|_H.
\]
Dla $u_0\in D(A)$ funkcja $u(t)=S(t)u_0$
należy do $C^1([0,\infty);H)\cap C([0,\infty);D(A))$
i spełnia $u'(t)+Au(t)=0$, $u(0)=u_0$.
Dla dowolnego $u_0\in H$ to samo równanie zachodzi dla
$t>0$, a $u(t)\to u_0$ w $H$ przy $t\downarrow0$.
\end{twierdzenie}
\begin{proof}
\textbf{Definicja i własność półgrupy.} Z Parsevala
$\|S(t)u\|_H^2=\sum_je^{-2\lambda_jt}|(u,\phi_j)|^2
\le e^{-2\lambda_1t}\|u\|_H^2$, więc szereg
określa ograniczony operator. Dla $u=\phi_1$ mamy
równość, co daje dokładną normę. Mnożenie współczynników
$e^{-\lambda_js}e^{-\lambda_jt}=e^{-\lambda_j(s+t)}$
pokazuje $S(t)S(s)=S(t+s)$, a $S(0)=I$.

\textbf{Ciągłość silna.} Dla $u\in H$ mamy
$\|S(t)u-u\|_H^2=\sum_j|e^{-\lambda_jt}-1|^2|u_j|^2$.
Dla każdego $j$ czynnik w module dąży do zera,
a jest ograniczony przez $1$. Ponieważ
$\sum_j|u_j|^2<\infty$, z twierdzenia o zbieżności
zdominowanej dla szeregu otrzymujemy zbieżność do zera.
Ciągłość w każdym $t\ge0$ wynika z półgrupy i oszacowania
$\|S(t)\|\le1$ (dla przyrostu dodatniego i ujemnego
stosujemy odpowiednią faktoryzację).

\textbf{Wygładzanie.} Dla $m\in\N$ i $t>0$
$\sup_{x\ge0}x^me^{-tx}=(m/(et))^m$;
pochodna funkcji $x\mapsto x^me^{-tx}$ znika przy $x=m/t$.
Wobec tego
$\sum_j\lambda_j^{2m}e^{-2\lambda_jt}|u_j|^2
\le(m/(et))^{2m}\|u\|_H^2$.
Z charakterystyki dziedzin potęg operatora (otrzymywanej
indukcyjnie z~\ref{thm:pde12-formdomain}) wynika
$S(t)u\in D(A^m)$ i podane oszacowanie.

\textbf{Generator.} Dla $u\in D(A)$ iloraz
$(e^{-\lambda_jt}-1)/t$ dąży do $-\lambda_j$, a
$|(e^{-\lambda_jt}-1)/t|\le\lambda_j$.
Sumowalność $\sum\lambda_j^2|u_j|^2$ pozwala
ponownie zastosować zbieżność zdominowaną i otrzymać
$\lim_{t\downarrow0}(S(t)u-u)/t=-Au$ w $H$.
Odwrotnie, gdy granica ilorazu istnieje w $H$ i równa się
$g$, jej $j$-ty współczynnik musi równać się
$-\lambda_ju_j$. Stąd z Parsevala
$\sum_j\lambda_j^2|u_j|^2=\|g\|_H^2<\infty$,
a więc $u\in D(A)$. To dowodzi równości dziedzin generatora.
Ciągłość $u(t)$ z wartościami w $D(A)$ dla $u_0\in D(A)$
wynika z rozpisania normy grafowej na współczynniki
$u_j$ i $\lambda_ju_j$. Dla $t>0$ różniczkowalność
wynika z wygładzania i własności półgrupy.
\end{proof}

\begin{twierdzenie}[Wzór Duhamela w $L^2$]
\label{thm:pde12-duhamel}
Niech $T>0$, $u_0\in H=L^2(\Omega)$ i
$f\in C([0,T];H)$. Funkcja
\[
u(t)=S(t)u_0+\int_0^tS(t-s)f(s)\,\mathrm ds,
\qquad 0\le t\le T,
\]
gdzie całkę rozumiemy w sensie Bochnera, jest ciągła
z wartościami w $H$ i spełnia oszacowanie
$\|u(t)\|_H\le e^{-\lambda_1t}\|u_0\|_H
+\int_0^te^{-\lambda_1(t-s)}\|f(s)\|_H\,\mathrm ds$.
Jest jedynym rozwiązaniem \emph{łagodnym} zagadnienia
$u'(t)+Au(t)=f(t)$, $u(0)=u_0$, zdefiniowanym
przez powyższy wzór całkowy. Jeśli dodatkowo
$u_0\in D(A)$ i $f\in C^1([0,T];H)$, to
$u\in C^1([0,T];H)\cap C([0,T];D(A))$
i równanie zachodzi w $H$ dla każdego $t\in[0,T]$.
\end{twierdzenie}
\begin{proof}
Mapowanie $(t,s)\mapsto S(t-s)f(s)$ jest silnie ciągłe
na trójkącie $0\le s\le t\le T$, ponieważ $f$ i $S$
są silnie ciągłe, a $\|S\|\le1$. Całka Bochnera
istnieje; jej normę szacujemy przez całkę norm.
Ciągłość $u$ wynika przez rozdzielenie różnicy całek
na część o wspólnym przedziale całkowania i krótki odcinek,
a następnie przez zbieżność zdominowaną.

Dla silniejszych danych, po zmianie zmiennej $r=t-s$,
możemy pisać $F(t)=\int_0^tS(r)f(t-r)\,\mathrm dr$.
Różniczkowanie w $H$ pod całką, dozwolone przez ciągłość
$f'$, daje
$F'(t)=S(t)f(0)+\int_0^tS(r)f'(t-r)\,\mathrm dr$.
Z drugiej strony korzystamy z formuły
$A S(r)=-(\mathrm d/\mathrm dr)S(r)$ dla $r>0$
i po całkowaniu przez części w zmiennej $r$, najpierw
na $[\varepsilon,t]$, a następnie przechodząc z
$\varepsilon\downarrow0$, otrzymujemy
$AF(t)=f(t)-S(t)f(0)-\int_0^tS(r)f'(t-r)\,\mathrm dr$.
Uzasadnienie przejścia do granicy w grafie $A$ jest
następujące: całki po odcinku $[\varepsilon,t]$ należą
do $D(A)$, ich obrazy przez $A$ zbiegają w $H$ do
prawej strony, a same całki zbiegają w $H$ do $F(t)$.
Domkniętość $A$ daje $F(t)\in D(A)$ i powyższą równość.
Dla $S(t)u_0$ mamy $\frac{\mathrm d}{\mathrm dt}S(t)u_0
=-AS(t)u_0$ z poprzedniego twierdzenia.
Zsumowanie obu rachunków daje $u'+Au=f$.

Dla dowolnego $u_0\in H$ i ciągłego $f$ wynik definiuje
rozwiązanie łagodne. Jednoznaczność w tej klasie wynika
bezpośrednio z wymogu spełniania wzoru całkowego.
Jeśli rozwiązanie jest klasyczne, podstawienie
$\frac{\mathrm d}{\mathrm ds}[S(t-s)u(s)]
=S(t-s)f(s)$ dla $0<s<t$, z domkniętością i
ciągłością, prowadzi do wzoru Duhamela; zatem rozwiązanie
łagodne zgadza się z klasycznym, ilekroć drugie istnieje.
\end{proof}

\begin{uwaga}[Regularność w czasie i na brzegu]
Oszacowanie $\|A^mS(t)\|\lesssim t^{-m}$ jest
oszacowaniem operatorowym. Nie daje automatycznie
$C^\infty(\overline\Omega)$ na dowolnym obszarze Lipschitzowskim: w narożach może zabraknąć regularności
brzegowej. Dla $t>0$ rozwiązanie jest jednak gładkie
wewnątrz $\Omega$ na mocy lokalnej regularności eliptycznej
i równania ciepła, przy odpowiedniej iteracji oszacowań.
\end{uwaga}

\section{Równanie falowe: propagacja i zachowanie energii}
\label{sec:pde12-wave}

\begin{twierdzenie}[Spektralne rozwiązanie równania falowego]
\label{thm:pde12-wave}
Niech $A=-\Delta_D$ jak wyżej, $c>0$,
$u_0\in H_0^1(\Omega)$ i $u_1\in L^2(\Omega)$.
Oznaczmy $a_j=(u_0,\phi_j)$, $b_j=(u_1,\phi_j)$.
Wtedy szereg
\[
u(t)=\sum_{j=1}^\infty
\left(a_j\cos(c\sqrt{\lambda_j}t)
+\frac{b_j}{c\sqrt{\lambda_j}}\sin(c\sqrt{\lambda_j}t)\right)\phi_j
\]
wyznacza jedyne rozwiązanie energetyczne równania
$u_{tt}+c^2Au=0$ z danymi $u(0)=u_0$, $u_t(0)=u_1$,
przy czym
$u\in C(\R;H_0^1(\Omega))\cap C^1(\R;L^2(\Omega))$.
Energia
\[
E(t)=\frac12\|u_t(t)\|_2^2+
\frac{c^2}{2}\|\nabla u(t)\|_2^2
\]
jest stała i równa $E(0)$.
\end{twierdzenie}
\begin{proof}
Z twierdzenia~\ref{thm:pde12-formdomain} mamy
$\sum_j\lambda_j|a_j|^2<\infty$ i
$\sum_j|b_j|^2<\infty$. Dla współczynnika
$u_j(t)=a_j\cos(c\sqrt{\lambda_j}t)
+b_j\sin(c\sqrt{\lambda_j}t)/(c\sqrt{\lambda_j})$
zachodzi
\[
|u_j'(t)|^2+c^2\lambda_j|u_j(t)|^2
=|b_j|^2+c^2\lambda_j|a_j|^2.
\]
To tożsamość Pitagorasa dla obrotu wektora
$(c\sqrt{\lambda_j}a_j,b_j)$ w płaszczyźnie.
Sumując po $j$, otrzymujemy równość energii dla sum
częściowych i jednolite w $t$ oszacowanie ogonów
w $H_0^1\times L^2$. Wynika z tego zbieżność
sum częściowych w
$C([-T,T];H_0^1)\times C([-T,T];L^2)$
dla każdego $T>0$. Funkcje częściowe spełniają równanie
drugiego rzędu. Przejście do granicy w słabym równaniu
odbywa się przez testowanie funkcjami $v\in H_0^1$
oraz funkcjami testowymi czasu:
$u_{tt}+c^2Au=0$ rozumiemy w $H^{-1}(\Omega)$.
Rzeczywiście, forma $\int_\Omega\nabla u(x)\cdot\nabla v(x)\,\mathrm d\lambda_d(x)$ jest
ciągła na $H_0^1$, więc równanie dla sum częściowych
przechodzi do granicy w $H^{-1}$.

Dla jednoznaczności niech $w$ będzie różnicą dwóch
rozwiązań energetycznych z zerowymi danymi. Standardowe
uśrednienie w czasie albo projekcja na $\phi_j$ redukuje
równanie do $w_j''+c^2\lambda_jw_j=0$ z zerowymi
danymi, a więc $w_j=0$ dla każdego $j$. Zupełność bazy
daje $w=0$. Równość energii dla limitu wynika z
silnej zbieżności obu jej składników.
\end{proof}

\begin{uwaga}[Czego analiza spektralna nie zastępuje?]
Rozwinięcie w funkcje własne pokazuje istnienie i
zachowanie globalnej energii, ale nie ujawnia bezpośrednio
geometrii stożków charakterystycznych. Skończona prędkość
propagacji, uzyskana w rozdziale~5 metodą energii lokalnej,
pozostaje niezależną, silniejszą lokalnie informacją.
\end{uwaga}

\section{Przepływ Schrödingera i związek z operatorem unitarnym}
\label{sec:pde12-schroedinger}

\begin{twierdzenie}[Ewolucja unitarna dla samosprzężonego $A$]
\label{thm:pde12-unitary}
Na zespoleniu $L^2(\Omega;\C)$ zdefiniujmy
$U(t)u=\sum_je^{-it\lambda_j}(u,\phi_j)\phi_j$,
$t\in\R$. Wtedy $(U(t))_{t\in\R}$ jest silnie
ciągłą grupą unitarną z generatorem $-iA$.
Dla $u_0\in D(A)$ funkcja $u(t)=U(t)u_0$
spełnia $iu_t=Au$ i $\|u(t)\|_2=\|u_0\|_2$.
\end{twierdzenie}
\begin{proof}
Z Parsevala i $|e^{-it\lambda_j}|=1$ otrzymujemy
$\|U(t)u\|_2=\|u\|_2$. Ponadto
$U(t)U(s)=U(t+s)$ oraz $U(t)^{-1}=U(-t)$,
więc operatory są unitarne. Silna ciągłość wynika z
$\sum_j|e^{-it\lambda_j}-1|^2|u_j|^2\to0$
przez zbieżność zdominowaną. Dla $u\in D(A)$ ilorazy
$\frac{e^{-it\lambda_j}-1}{t}$ są zdominowane przez
$\lambda_j$, więc różniczkowanie w $L^2$ daje
$\frac{\mathrm d}{\mathrm dt}U(t)u=-iAU(t)u$.
Odwrotnie, istnienie granicy ilorazu w $L^2$ wymusza
$\sum_j\lambda_j^2|u_j|^2<\infty$ z Parsevala,
a więc dziedziną generatora jest dokładnie $D(A)$.
\end{proof}

\begin{uwaga}
Dla ogólnego operatora samosprzężonego, niekoniecznie
zwartego rezolwentnie, odpowiednie twierdzenia o
półgrupach i grupach unitarnych otrzymuje się przez
rachunek spektralny. Twierdzenie Stone'a identyfikuje
generator każdej silnie ciągłej grupy unitarnej jako
operator postaci $-iA$ dla pewnego samosprzężonego $A$.
W tym rozdziale dowiedliśmy konstrukcji dla laplasjanu
z czysto punktowym widmem; pełna ogólność Stone'a
wymaga twierdzenia spektralnego dla operatorów
samosprzężonych, również o widmie ciągłym.
\end{uwaga}

\section{Podsumowanie: co wynika z założeń, a co nie?}
\label{sec:pde12-summary}

Na obszarze ograniczonym Lipschitzowskim
słaba forma Dirichleta definiuje samosprzężony
operator $A=-\Delta_D$. Zwarty charakter zanurzenia
$H_0^1\hookrightarrow L^2$ jest dokładnym powodem
zwartej rezolwenty, dyskretnego widma i ortonormalnej
bazy funkcji własnych. Stąd otrzymujemy alternatywę
Fredholma, rozwinięcia dla równania ciepła i falowego
oraz przepływ unitarny równania Schrödingera.

Nie dowiedliśmy tu ogólnego twierdzenia spektralnego dla
nieograniczonych operatorów samosprzężonych ani
pełnego twierdzenia Hille'a--Yosidy. W kolejnych rozdziałach
wykorzystamy jednak już udowodnione przypadki oraz
wyraźnie oddzielimy twierdzenia w pełnej ogólności od
wyników obowiązujących dla zwartej rezolwenty.

% Rozdział 13. Tekst przeznaczony do włączenia przed \appendix.
\chapter[Eliptyczne zagadnienia brzegowe]{Eliptyczne zagadnienia brzegowe: formy dywergencyjne, nierówność Gårdinga i alternatywa Fredholma}
\label{chap:pde13}

W rozdziałach~\ref{chap:pde-variational} i~\ref{chap:pde12} zbudowaliśmy
słabą teorię laplasjanu. Teraz współczynniki będą zależały od punktu,
możliwe będą człony pierwszego i zerowego rzędu, a forma nie musi być
symetryczna. Rozdzielimy trzy pytania: czy operator jest eliptyczny,
czy związana z nim forma jest koercywna, i czy odpowiednie zagadnienie
brzegowe ma dokładnie jedno rozwiązanie. Wykazanie samej eliptyczności
nie rozstrzyga ostatniego pytania.

\section{Postać dywergencyjna a niedywergencyjna}
\label{sec:pde13-operators}
Niech $d\geq1$, a $\Omega\subset\R^d$ będzie obszarem ograniczonym Lipschitzowskim. Wszystkie współczynniki w tym rozdziale są rzeczywiste;
analogiczne sformułowania zespolone wymagają form półtoraliniowych.
Dla funkcji $u$ rozważmy formalne wyrażenie różniczkowe
\begin{equation}\label{eq:pde13-div}
 Lu=-\sum_{i,j=1}^d\partial_i(a_{ij}\partial_j u)
       +\sum_{i=1}^d b_i\partial_i u+cu.
\end{equation}
Zakładamy $a_{ij},b_i,c\in L^\infty(\Omega)$; piszemy
$A(x)=(a_{ij}(x))$, $b(x)=(b_i(x))$. Ponadto istnieje $\lambda>0$
takie, że dla prawie każdego $x\in\Omega$ i każdego $\xi\in\R^d$
\begin{equation}\label{eq:pde13-elliptic}
 \xi^{\mathsf T}A(x)\xi\geq\lambda|\xi|^2.
\end{equation}
Warunek ten nazywamy \emph{jednostajną eliptycznością}.
W przypadku macierzy niesymetrycznej zależy on tylko od części
symetrycznej $(A+A^{\mathsf T})/2$; macierz $A$ nie musi być dodatnio
określona jako operator samosprzężony.

\begin{definicja}[Symbol główny i postać niedywergencyjna]
\label{def:pde13-symbol}
Symbolem głównym wyrażenia $L$ z~\eqref{eq:pde13-div} jest
$p_2(x,\xi)=\sum_{i,j}a_{ij}(x)\xi_i\xi_j$.
Formalny operator w \emph{postaci niedywergencyjnej} ma postać
$\widetilde Lu=-\sum_{i,j}a_{ij}\partial_{ij}u+
\sum_i\widetilde b_i\partial_i u+\widetilde c u$.
Dla współczynników $a_{ij}\in W^{1,\infty}(\Omega)$ i dostatecznie
gładkiego $u$ zachodzi tożsamość
$-\partial_i(a_{ij}\partial_j u)=-a_{ij}\partial_{ij}u
-(\partial_i a_{ij})\partial_j u$, z sumowaniem po indeksach.
Przy samych współczynnikach mierzalnych nie wolno używać tego
przekształcenia jako tożsamości operatorów na $H^1$.
\end{definicja}

\begin{przyklad}[Ta sama eliptyczność, inne pojęcie rozwiązania]
\label{ex:pde13-interface}
Na $\Omega=(-1,1)$ niech $a(x)=1$ dla $x<0$ i $a(x)=4$ dla $x>0$.
Funkcja $u(x)=x$ dla $x\le0$ oraz $u(x)=x/4$ dla $x\ge0$
należy do $H^1(\Omega)$ i spełnia $-(au')'=0$ dystrybucyjnie,
ponieważ strumień $au'=1$ jest stały. Jednocześnie $u'$ ma skok,
a $u''$ zawiera deltę Diraca. Iloczyn $a u''$ nie jest w tym
przypadku kanonicznie zdefiniowany, ponieważ $a$ jest nieciągła
w punkcie podparcia delty. Postać niedywergencyjna nie jest więc
jedynie innym zapisem tego samego zagadnienia.
\end{przyklad}

\section{Rozwiązanie słabe i funkcjonał prawej strony}
\label{sec:pde13-weak}
Przez $H^{-1}(\Omega)$ rozumiemy przestrzeń dualną
$(H_0^1(\Omega))^*$ (dla rzeczywistych przestrzeni Sobolewa).
Niech $V=H_0^1(\Omega)$, z normą
$\|u\|_V^2=\|u\|_{L^2}^2+\|\nabla u\|_{L^2}^2$.
Definiujemy formę dwuliniową
\begin{equation}\label{eq:pde13-form}
 a(u,v)=\int_\Omega\big(A\nabla u\cdot\nabla v
 +(b\cdot\nabla u)v+cuv\big)\,\mathrm d\lambda_d(x).
\end{equation}
Dzięki nierówności Höldera istnieje $M>0$ zależne od norm
współczynników takie, że
$|a(u,v)|\le M\|u\|_V\|v\|_V$ dla wszystkich $u,v\in V$.

\begin{definicja}[Zagadnienie Dirichleta w sensie słabym]
\label{def:pde13-weak}
Dla $F\in V^*=H^{-1}(\Omega)$ funkcja $u\in V$ jest
\emph{słabym rozwiązaniem} zagadnienia $Lu=F$ z jednorodnym
warunkiem Dirichleta, jeżeli
\begin{equation}\label{eq:pde13-weak}
 a(u,v)=\langle F,v\rangle_{V^*,V}\qquad(v\in V).
\end{equation}
Jeśli $f\in L^2(\Omega)$, to utożsamiamy je z elementem $V^*$
przez $\langle F,v\rangle=\int_\Omega f v\,\mathrm d\lambda_d(x)$.
Warunek brzegowy jest zakodowany w przynależności $u\in H_0^1$.
\end{definicja}

\begin{stwierdzenie}[Zgodność ze znaczeniem dystrybucyjnym]
\label{prop:pde13-distributional}
Jeżeli $u\in H_0^1(\Omega)$ spełnia~\eqref{eq:pde13-weak}, to
$Lu=F$ jako dystrybucje na $\Omega$. Odwrotnie, gdy
$u\in H_0^1(\Omega)$ spełnia $Lu=F$ dystrybucyjnie dla
$F\in H^{-1}(\Omega)$, to spełnia~\eqref{eq:pde13-weak}.
\end{stwierdzenie}
\begin{proof}
Dla $\varphi\in C_c^\infty(\Omega)$ definicja pochodnej
dystrybucyjnej daje
$\langle-\partial_i(a_{ij}\partial_j u),\varphi\rangle
=\int_\Omega a_{ij}(\partial_j u)(\partial_i\varphi)
\,\mathrm d\lambda_d(x)$; pozostałe składniki są lokalnie całkowalne.
To jest dokładnie równość~\eqref{eq:pde13-weak} dla $v=\varphi$.
W drugą stronę, funkcjonały $v\mapsto a(u,v)$ i $F$ są
ciągłe na $V$, a $C_c^\infty(\Omega)$ jest gęste w $V$ z
definicji $H_0^1$. Równość na funkcjach testowych rozszerza się
więc na całe $V$.
\end{proof}

\section{Nierówność Gårdinga i wymuszenie koercywności}
\label{sec:pde13-garding}
\begin{twierdzenie}[Elementarna nierówność Gårdinga]
\label{thm:pde13-garding}
Przy założeniach~\eqref{eq:pde13-elliptic} oraz
$a_{ij},b_i,c\in L^\infty(\Omega)$ forma~\eqref{eq:pde13-form}
spełnia
\begin{equation}\label{eq:pde13-garding}
 a(u,u)\geq\frac\lambda2\|\nabla u\|_2^2
 -\Gamma\|u\|_2^2,\qquad u\in V,
\end{equation}
gdzie można przyjąć
$\Gamma=\|b\|_\infty^2/(2\lambda)+\|c^-\|_\infty$,
$c^-=\max\{-c,0\}$. W szczególności dla
$\mu\geq\Gamma+\lambda/2$ forma
$a_\mu(u,v)=a(u,v)+\mu(u,v)_{L^2}$ jest koercywna na $V$:
$a_\mu(u,u)\geq(\lambda/2)\|u\|_V^2$.
\end{twierdzenie}
\begin{proof}
Z eliptyczności
$\int_\Omega A\nabla u\cdot\nabla u\,\mathrm d\lambda_d
\geq\lambda\|\nabla u\|_2^2$.
Dalej $cuv$ dla $u=v$ jest nie mniejsze niż
$-\|c^-\|_\infty\|u\|_2^2$.
Człon pierwszego rzędu szacujemy przez Höldera i
nierówność Younga $st\leq\eps s^2/2+t^2/(2\eps)$:
\[
 \left|\int_\Omega(b\cdot\nabla u)u\,\mathrm d\lambda_d\right|
 \leq \|b\|_\infty\|\nabla u\|_2\|u\|_2
 \leq\frac\lambda2\|\nabla u\|_2^2+
 \frac{\|b\|_\infty^2}{2\lambda}\|u\|_2^2.
\]
Dodanie trzech oszacowań daje~\eqref{eq:pde13-garding}.
Po dodaniu $\mu\|u\|_2^2$ otrzymujemy
$a_\mu(u,u)\geq(\lambda/2)\|\nabla u\|_2^2+
(\mu-\Gamma)\|u\|_2^2$, co jest co najmniej
$(\lambda/2)\|u\|_V^2$ przy podanym wyborze $\mu$.
\end{proof}

\begin{twierdzenie}[Istnienie dla przesuniętego operatora]
\label{thm:pde13-shift}
Przy założeniach twierdzenia~\ref{thm:pde13-garding} i dla
$\mu\geq\Gamma+\lambda/2$ dla każdego $F\in V^*$ istnieje
dokładnie jedno $u\in V$ spełniające
$a(u,v)+\mu(u,v)_{L^2}=\langle F,v\rangle$ dla wszystkich $v\in V$.
Ponadto $\|u\|_V\leq(2/\lambda)\|F\|_{V^*}$.
\end{twierdzenie}
\begin{proof}
Forma $a_\mu$ jest ciągła jako suma dwóch form ciągłych,
a z poprzedniego twierdzenia jest koercywna. Twierdzenie
Laxa--Milgrama~\ref{thm:pde11-LM} stosuje się także do form
niesymetrycznych. Podstawiając $v=u$ do równania słabego,
uzyskujemy $(\lambda/2)\|u\|_V^2\leq
\langle F,u\rangle\leq\|F\|_{V^*}\|u\|_V$.
Jeżeli $u=0$, oszacowanie jest trywialne; w przeciwnym razie
dzielimy przez $\|u\|_V$.
\end{proof}

\begin{uwaga}[Dlaczego samo przesunięcie nie rozwiązuje zagadnienia?]
Wynik dotyczy operatora $L+\mu I$, a niekoniecznie $L$.
Przykładowo na $(0,1)$ z warunkami Dirichleta operator
$L=-\mathrm d^2/\mathrm dx^2-\pi^2$ ma niezerowe jądro
rozpięte przez $\sin(\pi x)$; rozwiązanie równania $Lu=f$
nie istnieje dla wszystkich $f\in L^2$.
\end{uwaga}


\begin{uwaga}[Trzy różne wnioski z oszacowań formy]
\label{rem:pde-v42-garding}
Koercywność z twierdzenia~\ref{thm:pde11-lower-order}
pozwala bezpośrednio zastosować twierdzenie Laxa--Milgrama.
Sama nierówność Gårdinga~\eqref{eq:pde13-garding}
nie daje tego wniosku dla operatora $L$, lecz zapewnia
koercywność odpowiednio przesuniętego operatora $L+\mu I$.
Dopiero argument Fredholma w sekcji~\ref{sec:pde13-fredholm}
pozwoli wrócić do $L$, z możliwym warunkiem ortogonalności
do jądra operatora sprzężonego. Te trzech kroków
nie należy utożsamiać.
\end{uwaga}

\section{Dirichlet: jednoznaczność, dane brzegowe i stabilność}
\label{sec:pde13-dirichlet}
\begin{twierdzenie}[Koercywność bez przesunięcia]
\label{thm:pde13-coercive}
Załóżmy, że $b=0$ i $c\geq0$ prawie wszędzie, a macierz $A$
spełnia~\eqref{eq:pde13-elliptic}. Wówczas dla każdego $F\in V^*$
zagadnienie~\eqref{eq:pde13-weak} ma dokładnie jedno rozwiązanie,
a $\|u\|_V\leq(1+C_P^2)\|F\|_{V^*}/\lambda$, gdzie
$C_P$ jest stałą nierówności Poincarégo
$\|v\|_2\leq C_P\|\nabla v\|_2$ na $V$.
To samo zachodzi przy ogólnych $b,c$, jeśli forma $a$ jest
koercywna na $V$; przykładowo wystarcza
$\Gamma C_P^2<\lambda/2$ w oszacowaniu~\eqref{eq:pde13-garding}.
\end{twierdzenie}
\begin{proof}
Dla $b=0$, $c\geq0$ mamy
$a(v,v)\geq\lambda\|\nabla v\|_2^2
\geq\lambda(1+C_P^2)^{-1}\|v\|_V^2$.
Twierdzenie Laxa--Milgrama daje istnienie i jednoznaczność,
a testowanie rozwiązaniem i nierówność dualności dają oszacowanie.
W drugim przypadku~\eqref{eq:pde13-garding} wraz z nierównością
Poincarégo daje
$a(v,v)\geq(\lambda/2-\Gamma C_P^2)\|\nabla v\|_2^2$.
Jeśli współczynnik w nawiasie jest dodatni, otrzymujemy
koercywność po ponownym użyciu nierówności Poincarégo.
\end{proof}

\begin{twierdzenie}[Niejednorodny warunek Dirichleta]
\label{thm:pde13-nonhom}
Niech $A,b,c$ spełniają założenia sekcji~\ref{sec:pde13-operators},
a forma $a$ będzie koercywna: $a(v,v)\geq\alpha\|v\|_V^2$
dla $v\in V$, gdzie $\alpha>0$.
Dla każdego $F\in V^*$ i $g\in H^{1/2}(\partial\Omega)$
istnieje dokładnie jedna funkcja $u\in H^1(\Omega)$ o śladzie
$\operatorname{Tr}u=g$ taka, że
$a(u,v)=\langle F,v\rangle$ dla wszystkich $v\in V$.
Ponadto
\[
 \|u\|_{H^1(\Omega)}\leq C\big(\|F\|_{V^*}
 +\|g\|_{H^{1/2}(\partial\Omega)}\big),
\]
gdzie $C$ zależy od $\alpha$, ciągłości formy i operatora
przedłużania śladu, ale nie od $F,g$.
\end{twierdzenie}
\begin{proof}
Definicja~\ref{def:pde10-halftrace} i lemat~\ref{lem:pde10-halftrace-rightinverse}
zapewniają liniowy ciągły prawy odwrotny operatora śladu:
wybieramy $G\in H^1(\Omega)$ tak, że $\operatorname{Tr}G=g$ i
$\|G\|_{H^1}\leq C_{\operatorname{Tr}}\|g\|_{H^{1/2}}$.
Szukamy $u=G+w$, gdzie $w\in V$.
Równanie dla $w$ przyjmuje postać
$a(w,v)=\langle F,v\rangle-a(G,v)$.
Prawa strona jest ciągłym funkcjonałem na $V$ o normie
co najwyżej $\|F\|_{V^*}+M\|G\|_{H^1}$.
Z koercywności i twierdzenia Laxa--Milgrama istnieje
jedno $w$ i
$\|w\|_V\leq\alpha^{-1}
(\|F\|_{V^*}+M\|G\|_{H^1})$.
Dodajemy normę $G$ i otrzymujemy żądane oszacowanie.
Różnica dwóch rozwiązań ma zerowy ślad, należy więc do $V$
i spełnia równanie jednorodne; koercywność wymusza jej znikanie.
\end{proof}

\begin{twierdzenie}[Słaba zasada minimum bez członu transportowego]
\label{thm:pde13-positivity}
Niech $b=0$, $c\geq0$ i $A$ spełnia~\eqref{eq:pde13-elliptic};
nie wymagamy symetrii macierzy $A$. Niech $u\in H^1(\Omega)$
spełnia $a(u,v)=\langle F,v\rangle$ dla wszystkich $v\in V$,
przy czym $\operatorname{Tr}u\geq0$ na $\partial\Omega$
(w sensie prawie wszędzie) i
$\langle F,v\rangle\geq0$ dla wszystkich nieujemnych $v\in V$.
Wtedy $u\geq0$ prawie wszędzie w $\Omega$.
\end{twierdzenie}
\begin{proof}
Ujemna część $u^-=\max\{-u,0\}$ należy do $H^1(\Omega)$.
Reguła łańcuchowa dla funkcji Lipschitzowskich daje
$\nabla u^-=-\chi_{\{u<0\}}\nabla u$ prawie wszędzie.
Ponieważ ślad $u$ jest nieujemny, $\operatorname{Tr}(u^-)=0$;
z charakteryzacji jądra śladu wynika $u^-\in V$.
Testujemy równanie $v=u^-$. Na zbiorze $\{u<0\}$ mamy
$u=-u^-$, $\nabla u=-\nabla u^-$, a poza nim
$u^-=\nabla u^-=0$ prawie wszędzie. Zatem
\[
 0\leq\langle F,u^-\rangle=a(u,u^-)
 =-\int_\Omega \big(A\nabla u^-\cdot\nabla u^-+c(u^-)^2\big)
 \,\mathrm d\lambda_d(x)
 \leq-\lambda\|\nabla u^-\|_2^2.
\]
Otrzymujemy $\nabla u^-=0$, a nierówność Poincarégo na $V$
daje $u^-=0$ w $L^2$. To jest równoważne $u\geq0$.
\end{proof}

\begin{twierdzenie}[Zasada porównania z dryfem o kontrolowanej dywergencji]
\label{thm:pde13-drift-comparison}
Niech $\Omega\subset\R^d$ będzie ograniczonym obszarem
Lipschitzowskim, $A\in L^\infty(\Omega;\R^{d\times d})$
spełnia~\eqref{eq:pde13-elliptic}, $b\in W^{1,\infty}(\Omega;\R^d)$,
$c\in L^\infty(\Omega)$ oraz
$c-\frac12\diver b\ge0$ prawie wszędzie. Jeżeli
$u_1,u_2\in H^1(\Omega)$ mają ślady spełniające
$\operatorname{Tr}u_1\le\operatorname{Tr}u_2$ prawie wszędzie
na $\partial\Omega$ oraz
\[
 a(u_1,\varphi)\le a(u_2,\varphi)
 \quad\text{dla wszystkich }\varphi\in H_0^1(\Omega),\ \varphi\ge0,
\]
to $u_1\le u_2$ prawie wszędzie w $\Omega$.
W szczególności założenie $\diver b=0$, $c\ge0$ zapewnia
porównanie niezależnie od wielkości normy $b$.
\end{twierdzenie}
\begin{proof}
Kładziemy $w=u_1-u_2$ i $v=w^+$. Reguła łańcuchowa
i własność śladu dają $v\in H_0^1(\Omega)$ oraz
$\nabla v=\chi_{\{w>0\}}\nabla w$.
Dopuszczalnym testem jest $\varphi=v$.
Człon pierwszego rzędu obliczamy bez założenia małości $b$:
\[
 \int_\Omega (b\cdot\nabla w)v\,\mathrm d\lambda_d(x)
 =\frac12\int_\Omega b\cdot\nabla(v^2)\,\mathrm d\lambda_d(x)
 =-\frac12\int_\Omega(\diver b)v^2\,\mathrm d\lambda_d(x).
\]
Ostatnia równość wynika z aproksymacji $v$ funkcjami
$C_c^\infty(\Omega)$ w $H^1$; ponieważ $v\in H_0^1$,
nie występuje człon brzegowy. Wobec tego
\[
 0\ge a(w,v)=\int_\Omega\left(A\nabla v\cdot\nabla v+
 (c-\tfrac12\diver b)v^2\right)\,\mathrm d\lambda_d(x)
 \ge\lambda\|\nabla v\|_2^2.
\]
Nierówność Poincarégo daje $v=0$. Jest to argument
energetyczny, a nie twierdzenie o punktowej regularności
rozwiązań.
\end{proof}



\section{Alternatywa Fredholma dla eliptycznych zagadnień Dirichleta}
\label{sec:pde13-fredholm}
Dla współczynników dolnego rzędu bez warunku znakowego forma $a$
może nie być koercywna. Mimo to operator $L:V\to V^*$,
$\langle Lu,v\rangle=a(u,v)$, ma precyzyjną własność Fredholma.
Nie należy go mylić z realizacją $L$ w $L^2$, której dziedzina
wymaga dodatkowej definicji.

\begin{twierdzenie}[Alternatywa Fredholma dla $L:V\to V^*$]
\label{thm:pde13-fredholm}
Niech $\Omega\subset\R^d$ będzie obszarem ograniczonym Lipschitzowskim, a $A,b,c\in L^\infty$ niech spełniają
jednostajną eliptyczność~\eqref{eq:pde13-elliptic}.
Wówczas $L:V\to V^*$ jest operatorem Fredholma indeksu $0$:
ma domknięty obraz, skończeniewymiarowe jądro i skończony
kowymiar obrazu, równy wymiarowi jądra.
Definiując operator transponowany $L^{\mathrm t}:V\to V^*$ przez
$\langle L^{\mathrm t}z,v\rangle=a(v,z)$, mamy
$\dim\ker L=\dim\ker L^{\mathrm t}<\infty$.
Równanie $Lu=F$, $F\in V^*$, ma rozwiązanie wtedy i tylko wtedy, gdy
\begin{equation}\label{eq:pde13-compat}
 \langle F,z\rangle=0\qquad(z\in\ker L^{\mathrm t}).
\end{equation}
Rozwiązania różnią się o elementy $\ker L$. W szczególności
gdy $\ker L=\{0\}$, dla każdego $F\in V^*$ istnieje jedno rozwiązanie.
\end{twierdzenie}
\begin{proof}
\textbf{Krok 1: odwracalny operator przesunięty.}
Wybieramy $\mu\geq\Gamma+\lambda/2$.
Z twierdzenia~\ref{thm:pde13-shift} operator
$L_\mu=L+\mu J:V\to V^*$ jest izomorfizmem ograniczonym,
przy czym $J:V\to V^*$ definiujemy wzorem
$\langle Ju,v\rangle=\int_\Omega uv\,\mathrm d\lambda_d(x)$.

\textbf{Krok 2: zwartość zaburzenia.}
Zanurzenie $i:V\hookrightarrow L^2(\Omega)$ jest zwarte
na mocy twierdzenia Rellicha--Kondrachova. Operator
$j:L^2(\Omega)\to V^*$, $\langle jf,v\rangle=(f,v)_{L^2}$,
jest ograniczony, ponieważ $\|v\|_2\leq\|v\|_V$.
Zatem $J=j\circ i$ jest zwarty jako złożenie operatora
zwartego z ograniczonym. W konsekwencji operator
$K=\mu L_\mu^{-1}J:V\to V$ jest zwarty.

\textbf{Krok 3: sprowadzenie do alternatywy dla operatora zwartego.}
Mamy równość operatorów
$L=L_\mu(I-K)$; istotnie
$L_\mu K=\mu J$.
Twierdzenie~\ref{thm:pde12-general-fredholm} zostało
udowodnione dla zespolonych przestrzeni Hilberta;
jego dowód oparty na aproksymacji operatorami skończonego rzędu
pozostaje bez zmian prawdziwy w rzeczywistych przestrzeniach
Hilberta. Daje to, że $I-K:V\to V$ ma domknięty obraz,
skończeniewymiarowe jądro i kowymiar równy wymiarowi jądra.
Ponieważ $L_\mu$ jest izomorfizmem, własności te przechodzą
na $L:V\to V^*$, a zatem $\operatorname{ind}L=0$.

\textbf{Krok 4: dokładny warunek zgodności.}
Dla $z\in V$ warunek $z\in\ker L^{\mathrm t}$ oznacza
$a(v,z)=0$ dla wszystkich $v\in V$. To jest równoważne temu,
że ciągły funkcjonał na $V^*$ zadany przez
$F\mapsto\langle F,z\rangle$ znika na $\operatorname{Ran}L$.
Każdy ciągły funkcjonał na $V^*$ ma tę postać, ponieważ
$V$ jest refleksywna (jako przestrzeń Hilberta).
Przestrzeń funkcjonałów zerujących obraz $L$ ma zatem
wymiar $\dim\ker L^{\mathrm t}$.
Ponieważ obraz jest domknięty, twierdzenie Hahna--Banacha
lub identyfikacja przestrzeni ilorazowej daje
$\operatorname{Ran}L=\{F\in V^*: \langle F,z\rangle=0
\text{ dla każdego } z\in\ker L^{\mathrm t}\}$.
Kowymiar obrazu jest równy $\dim\ker L^{\mathrm t}$,
a z kroku 3 także $\dim\ker L$. Wreszcie różnica rozwiązań
jednego równania należy do jądra $L$.
\end{proof}

\begin{przyklad}[Zgodność przy rezonansie]
\label{ex:pde13-resonance}
Dla $\Omega=(0,1)$ i $L=-\mathrm d^2/\mathrm dx^2-\pi^2$
z warunkiem Dirichleta mamy
$\ker L=\ker L^{\mathrm t}=\operatorname{span}\{\sin(\pi x)\}$.
Równanie $-u''-\pi^2u=f$, $u(0)=u(1)=0$,
ma słabe rozwiązanie dokładnie wtedy, gdy
$\int_0^1 f(x)\sin(\pi x)\,\mathrm d\lambda_1(x)=0$.
Gdy ten warunek zachodzi, rozwiązanie jest wyznaczone
z dokładnością do dodania wielokrotności $\sin(\pi x)$.
W szczególności dla $f=\sin(\pi x)$ rozwiązanie nie istnieje.
\end{przyklad}

\section{Zagadnienie Neumanna: warunek zgodności i stałe}
\label{sec:pde13-neumann}
Dla $b=c=0$ i jednorodnej eliptyczności rozważmy
$-\diver(A\nabla u)=f$ wraz z zadanym strumieniem na brzegu.
Przy słabych współczynnikach nie musimy definiować strumienia
punktowo; dane brzegowe można wprowadzić przez funkcjonał śladu.

\begin{lemat}[Nierówność Poincarégo--Wirtingera]
\label{lem:pde13-wirtinger}
Jeżeli $\Omega$ jest obszarem ograniczonym Lipschitzowskim,
to istnieje $C_W>0$ takie, że dla każdego $u\in H^1(\Omega)$
\[
\|u-u_\Omega\|_{L^2(\Omega)}\le C_W\|\nabla u\|_{L^2(\Omega)},
\qquad u_\Omega=\frac1{\lambda_d(\Omega)}
\int_\Omega u\,\mathrm d\lambda_d(x).
\]
\end{lemat}
\begin{proof}
Gdyby twierdzenie było fałszywe, znaleźlibyśmy $v_n\in H^1$
o średniej zero takie, że $\|v_n\|_2=1$ i
$\|\nabla v_n\|_2\to0$ (odejmujemy średnią i normalizujemy).
Ciąg $v_n$ jest ograniczony w $H^1$, więc dzięki
zwartości $H^1(\Omega)\hookrightarrow L^2(\Omega)$
ma podciąg zbieżny silnie w $L^2$ do $v$. Po wybraniu
podciągu słabo zbieżnego w $H^1$ otrzymujemy
$\nabla v=0$ w sensie dystrybucyjnym.
Ponieważ $\Omega$ jest spójna, $v$ jest prawie wszędzie stała.
Ciągłość funkcjonału $w\mapsto\int_\Omega w$ oraz średnie
zerowe pokazują, że ta stała wynosi zero.
Ale silna zbieżność w $L^2$ wymusza $\|v\|_2=1$,
co jest sprzecznością.
\end{proof}

\begin{twierdzenie}[Neumann: istnienie modulo stałe]
\label{thm:pde13-neumann}
Niech $\Omega$ będzie obszarem ograniczonym Lipschitzowskim, a $A\in L^\infty(\Omega;\R^{d\times d})$
spełnia~\eqref{eq:pde13-elliptic}. Niech
$f\in L^2(\Omega)$ i $h\in L^2(\partial\Omega)$,
przy czym na brzegu całkujemy względem miary powierzchniowej
$\sigma=\mathcal H^{d-1}|_{\partial\Omega}$.
Istnieje $u\in H^1(\Omega)$ spełniające
\begin{equation}\label{eq:pde13-neumann}
\int_\Omega A\nabla u\cdot\nabla v\,\mathrm d\lambda_d(x)
=\int_\Omega fv\,\mathrm d\lambda_d(x)
+\int_{\partial\Omega} h\,\operatorname{Tr}v\,\mathrm d\sigma
\quad(v\in H^1(\Omega))
\end{equation}
wtedy i tylko wtedy, gdy
\begin{equation}\label{eq:pde13-neumann-comp}
\int_\Omega f\,\mathrm d\lambda_d(x)
+\int_{\partial\Omega}h\,\mathrm d\sigma=0.
\end{equation}
Rozwiązanie jest jednoznaczne modulo stałe. Istnieje dokładnie
jedno rozwiązanie o średniej zero; spełnia ono
$\|u\|_{H^1}\leq C(\|f\|_2+\|h\|_{L^2(\partial\Omega)})$.
\end{twierdzenie}
\begin{proof}
\textbf{Konieczność.} Podstawiamy $v=1$ do
\eqref{eq:pde13-neumann}; gradient stałej znika i
otrzymujemy dokładnie~\eqref{eq:pde13-neumann-comp}.

\textbf{Wystarczalność.} Rozważmy domkniętą podprzestrzeń
$V_\diamond=\{v\in H^1(\Omega):v_\Omega=0\}$.
Na niej norma $\|\nabla v\|_2$ jest równoważna normie
$H^1$ na mocy lematu~\ref{lem:pde13-wirtinger}.
Forma $B(u,v)=\int_\Omega A\nabla u\cdot\nabla v
\,\mathrm d\lambda_d(x)$ jest ciągła oraz koercywna
względem tej normy, ponieważ
$B(v,v)\geq\lambda\|\nabla v\|_2^2$.
Funkcjonał po prawej stronie~\eqref{eq:pde13-neumann}
jest ciągły na $V_\diamond$: przez nierówność Höldera
oraz ciągłość śladu jego moduł jest nie większy niż
$C(\|f\|_2+\|h\|_{L^2(\partial\Omega)})\|v\|_{H^1}$.
Z twierdzenia Laxa--Milgrama istnieje jedno $u\in V_\diamond$
spełniające równanie dla wszystkich $v\in V_\diamond$.
Dla dowolnego $v\in H^1$ piszemy $v=(v-v_\Omega)+v_\Omega$.
Pierwsza część jest w $V_\diamond$, a prawa strona na stałej
znika dzięki zgodności~\eqref{eq:pde13-neumann-comp}.
Równanie jest zatem prawdziwe dla całego $H^1$.
Oszacowanie wynika z Laxa--Milgrama i lematu Wirtingera.
Jeśli $u_1,u_2$ są rozwiązaniami, ich różnica $w$ spełnia
$B(w,v)=0$ dla wszystkich $v\in H^1$. Biorąc $v=w$,
otrzymujemy $\nabla w=0$, a spójność $\Omega$ implikuje,
że $w$ jest prawie wszędzie stała.
\end{proof}

\begin{uwaga}[Znak strumienia]
Dla gładkiego $u$ i $A$ całkowanie przez części daje
$\int_\Omega A\nabla u\cdot\nabla v
=\int_\Omega[-\diver(A\nabla u)]v
+\int_{\partial\Omega}(A\nabla u\cdot\nu)v$.
Dlatego $h$ w~\eqref{eq:pde13-neumann} oznacza
\emph{zewnętrzny} strumień $A\nabla u\cdot\nu$.
Przy danych szorstkich jest to jedynie interpretacja
funkcjonałowa; nie zakładamy istnienia punktowego śladu
wektora $A\nabla u$.
\end{uwaga}

\section{Robin: człon brzegowy przywracający jednoznaczność}
\label{sec:pde13-robin}
\begin{lemat}[Nierówność Poincarégo z wyrazem brzegowym]
\label{lem:pde13-boundary-poincare}
Na obszarze ograniczonym Lipschitzowskim istnieje
$C>0$ takie, że dla wszystkich $u\in H^1(\Omega)$
\[
\|u\|_{H^1}^2\le C\big(\|\nabla u\|_2^2+
 \|\operatorname{Tr}u\|_{L^2(\partial\Omega)}^2\big).
\]
\end{lemat}
\begin{proof}
Gdyby oszacowanie nie zachodziło, mielibyśmy ciąg
$\|u_n\|_{H^1}=1$ i $\|\nabla u_n\|_2^2+
\|\operatorname{Tr}u_n\|_{L^2(\partial\Omega)}^2\to0$.
Ze zwartości zanurzenia $H^1\hookrightarrow L^2$ i słabej
zwartości kuli w $H^1$ otrzymujemy podciąg słabo zbieżny
w $H^1$, silnie w $L^2$ do pewnego $u$.
Ponieważ $\nabla u_n\to0$ w $L^2$, mamy $\nabla u=0$,
a więc $u$ jest stała. Z ciągłości śladu $H^1\to
L^2(\partial\Omega)$ otrzymujemy słabą zbieżność śladów
do $\operatorname{Tr}u$, lecz ślady dążą silnie do zera.
Zatem stała $u$ ma zerowy ślad, czyli $u=0$.
Wówczas $\|u_n\|_2\to0$ oraz $\|\nabla u_n\|_2\to0$,
wbrew $\|u_n\|_{H^1}=1$.
\end{proof}

\begin{twierdzenie}[Zagadnienie Robina]
\label{thm:pde13-robin}
Niech $\Omega$ będzie obszarem ograniczonym Lipschitzowskim, $A\in L^\infty$ spełnia
\eqref{eq:pde13-elliptic}, a $\beta\in L^\infty(\partial\Omega)$
spełnia $\beta\geq\beta_0>0$ prawie wszędzie.
Dla każdego $f\in L^2(\Omega)$ i $h\in L^2(\partial\Omega)$
istnieje dokładnie jedno $u\in H^1(\Omega)$ spełniające
\begin{equation}\label{eq:pde13-robin}
\int_\Omega A\nabla u\cdot\nabla v\,\mathrm d\lambda_d(x)
+\int_{\partial\Omega}\beta\,\operatorname{Tr}u\,
\operatorname{Tr}v\,\mathrm d\sigma
=\int_\Omega fv\,\mathrm d\lambda_d(x)
+\int_{\partial\Omega}h\,\operatorname{Tr}v\,\mathrm d\sigma
\end{equation}
dla wszystkich $v\in H^1(\Omega)$, a ponadto
$\|u\|_{H^1}\leq C(\|f\|_2+\|h\|_{L^2(\partial\Omega)})$.
Formalnie warunek brzegowy ma postać
$A\nabla u\cdot\nu+\beta u=h$.
\end{twierdzenie}
\begin{proof}
Ciągłość formy w lewej stronie wynika z ograniczoności $A,\beta$
oraz ciągłości operatora śladu $H^1\to L^2(\partial\Omega)$.
Jej wartość dla $v=u$ jest nie mniejsza niż
$\lambda\|\nabla u\|_2^2+
\beta_0\|\operatorname{Tr}u\|_{L^2(\partial\Omega)}^2$,
a z lematu~\ref{lem:pde13-boundary-poincare} wynika koercywność
na całym $H^1$. Prawa strona jest ciągłym funkcjonałem
na $H^1$ przez Höldera i twierdzenie o śladzie.
Wszystkie założenia Laxa--Milgrama są więc spełnione;
otrzymujemy istnienie, jednoznaczność i oszacowanie.
\end{proof}

\section{Przykład obliczeniowy: strumień przez granicę ośrodków}
\label{sec:pde13-transmission}
Niech $\Omega=(-1,1)$, $f=1$ i
$a(x)=1$ dla $x<0$, $a(x)=4$ dla $x>0$.
Rozważamy $-(a u')'=1$ wraz z $u(-1)=u(1)=0$.
Na obu podprzedziałach $u''=-1/a$, a zatem
\begin{equation}\label{eq:pde13-exact}
 u(x)=\begin{cases}
 -x^2/2-3x/10+1/5,& -1\le x\le0,\\
 -x^2/8-3x/40+1/5,& 0\le x\le1.
 \end{cases}
\end{equation}
\begin{center}
\begin{tikzpicture}[x=3.1cm,y=0.7cm,>=Stealth]
  \fill[blue!8] (-1,0) rectangle (0,1.35);
  \fill[orange!13] (0,0) rectangle (1,1.35);
  \draw[thick] (-1,0) rectangle (1,1.35);
  \draw[thick,dashed] (0,0)--(0,1.35);
  \node at (-.52,.87) {$a_-=1$};
  \node at (.52,.87) {$a_+=4$};
  \node at (-.52,.38) {ośrodek I};
  \node at (.52,.38) {ośrodek II};
  \draw[->] (-1.15,-.24)--(1.18,-.24) node[right] {$x$};
  \foreach \x/\lab in {-1/-1,0/0,1/1}{
    \draw (\x,-.31)--(\x,-.17);\node[below] at (\x,-.31) {$\lab$};}
  \node[above] at (0,1.35) {interfejs};
\end{tikzpicture}

\small\emph{Dwa ośrodki o różnych współczynnikach: ciągłe $u$ i $a u'$, lecz skok $u'$.}
\end{center}

Wartość $u(0^-)=u(0^+)=1/5$ jest ciągła, natomiast
$u'(0^-)=-3/10$ i $u'(0^+)=-3/40$. Mimo skoku pochodnej
strumień $a(x)u'(x)=-x-3/10$ jest ciągły w $0$.

\begin{stwierdzenie}[Warunki transmisyjne z równania słabego]
\label{prop:pde13-transmission}
Niech $a$ jest dodatnią stałą $a_-$ na $(-1,0)$ i
$a_+$ na $(0,1)$, $f\in L^2(-1,1)$, a $u\in H^1(-1,1)$
ma na obu podprzedziałach regularność $H^2$.
Wówczas $-(au')'=f$ w $\mathcal D'(-1,1)$ wtedy i tylko wtedy,
gdy $-a_\pm u''=f$ prawie wszędzie po odpowiednich stronach,
a ślady strumienia spełniają
$a_-u'(0^-)=a_+u'(0^+)$.
\end{stwierdzenie}
\begin{proof}
Z ciągłości funkcji z $H^1(-1,1)$ w wymiarze jednym
wynika $u(0^-)=u(0^+)$. Testując równanie funkcjami o nośniku
w każdym z podprzedziałów, dostajemy tam równania klasyczne
prawie wszędzie. Dla dowolnej $\varphi\in C_c^\infty(-1,1)$
całkujemy osobno po obu przedziałach:
\[
\int_{-1}^1 au'\varphi'\,\mathrm d\lambda_1
=\int_{-1}^1 f\varphi\,\mathrm d\lambda_1
 +\big(a_-u'(0^-)-a_+u'(0^+)\big)\varphi(0).
\]
Zatem równość dystrybucyjna na całym przedziale jest
równoważna znikaniu współczynnika przy $\varphi(0)$.
W drugą stronę te same rachunki, przy zerowym skoku strumienia,
dają równanie globalne. Dla~\eqref{eq:pde13-exact} skok strumienia
wynosi zero, co potwierdza poprawność rozwiązania.
\end{proof}

\section{Zakres otrzymanych rezultatów}
\label{sec:pde13-summary}
Jednostajna eliptyczność zapewnia kontrolę części gradientowej,
lecz człony dolnego rzędu mogą zniszczyć koercywność.
Nierówność Gårdinga pozwala odzyskać ją po przesunięciu,
a zwartość zanurzenia Sobolewa prowadzi do alternatywy Fredholma
dla problemu pierwotnego. W zagadnieniu Neumanna przeszkodą
jest przestrzeń stałych, zaś dodatni człon brzegowy w problemie
Robina usuwa tę niejednoznaczność.

Wszystkie twierdzenia o istnieniu dotyczą rozwiązań słabych.
Nie twierdzimy, że sam warunek eliptyczności i ograniczoność
mierzalnych współczynników dają $H^2$ lub gładkość klasyczną.
Przykład transmisyjny pokazuje dokładnie, gdzie taka
regularność może zawieść. Regularność wewnętrzną i brzegową
rozwiniemy w kolejnych rozdziałach.

% Iteracja 14: wykład, do włączenia przed dodatkiem z kodem.
\chapter[Zasady maksimum i teoria potencjału]{Zasady maksimum, lemat Hopfa, nierówności Harnacka i podstawy teorii potencjału}
\label{chap:pde14}

W rozdziale~\ref{chap:pde13} uzyskaliśmy rozwiązania słabe
zagadnień eliptycznych metodami wariacyjnymi. Samo istnienie w
$H^1$ nie mówi jeszcze, jaki znak ma rozwiązanie, czy osiąga ono
maksimum wewnątrz obszaru ani jak wartości w różnych punktach
zależą od siebie. Odpowiedzi na te pytania dostarczają zasady
maksimum i nierówności Harnacka. Zaczniemy od współczynników
mierzalnych, dla których wystarczy argument energetyczny, a potem
przejdziemy do silniejszych twierdzeń dla funkcji harmonicznych.

\section{Słaba zasada maksimum dla rozwiązań słabych}
\label{sec:pde14-weak}
Niech $\Omega\subset\R^d$ będzie obszarem ograniczonym Lipschitzowskim. Załóżmy, że $A\in L^\infty(\Omega;\R^{d\times d})$
spełnia $\xi^{\mathsf T}A(x)\xi\ge\lambda|\xi|^2$
dla pewnego $\lambda>0$, prawie każdego $x\in\Omega$ i każdego
$\xi\in\R^d$. Nie zakładamy symetrii $A$.
Rozważamy operator $Lu=-\diver(A\nabla u)+cu$, gdzie
$c\in L^\infty(\Omega)$ i $c\ge0$ prawie wszędzie.

\begin{definicja}[Słabe rozwiązanie nierówności]
\label{def:pde14-subsolution}
Funkcja $u\in H^1(\Omega)$ jest \emph{słabym podrozwiązaniem}
$Lu\le0$ w $\Omega$, jeżeli
\[
 \int_\Omega\big(A\nabla u\cdot\nabla\varphi
 +cu\varphi\big)\,\mathrm d\lambda_d(x)\le0
 \quad\text{dla każdego }\varphi\in H_0^1(\Omega),\ \varphi\ge0.
\]
Analogicznie definiujemy nadrozwiązanie, odwracając nierówność.
Równość dla wszystkich funkcji testowych definiuje rozwiązanie słabe.
W definicji nie zakładamy istnienia klasycznych pochodnych drugiego rzędu.
\end{definicja}

\begin{lemat}[Obcięcie poziomowe]
\label{lem:pde14-trunc}
Jeżeli $u\in H^1(\Omega)$ i $k\in\R$, to
$w=(u-k)^+:=\max\{u-k,0\}\in H^1(\Omega)$ oraz
$\nabla w=\chi_{\{u>k\}}\nabla u$ prawie wszędzie.
Jeżeli ponadto $\operatorname{Tr}u\le k$ prawie wszędzie na
$\partial\Omega$, to $w\in H_0^1(\Omega)$.
\end{lemat}
\begin{proof}
Przybliżamy funkcję $t\mapsto t^+$ funkcjami klasy $C^1$
o jednostajnie ograniczonej pochodnej i stosujemy regułę
łańcuchową dla funkcji Sobolewa. Graniczne pochodne zbiegają
w $L^2$ dzięki zbieżności zdominowanej, po uprzednim
ustaleniu pochodnej prawie wszędzie poza zbiorem poziomicowym
$\{u=k\}$. Na tym zbiorze $\nabla u=0$ prawie wszędzie:
fakt ten udowodniono w rozdziale~\ref{chap:pde-sobolev}
(dowód dla obcięć funkcji Sobolewa). Otrzymujemy wskazany
wzór na gradient. Ciągłość operatora śladu oraz jego zgodność
ze złożeniem przez funkcję Lipschitzowską dają
$\operatorname{Tr}w=(\operatorname{Tr}u-k)^+=0$.
Ponieważ na obszarze Lipschitzowskim
$\ker\operatorname{Tr}=H_0^1(\Omega)$, kończy to dowód.
\end{proof}

Poniższy wynik uogólnia elementarną zasadę maksimum
z twierdzenia~\ref{thm:laplace-max}: tutaj współczynniki
są jedynie mierzalne, a równanie i warunek brzegowy
interpretujemy w sensie słabym. Hipotezy obu twierdzeń
są różne, więc nie utożsamiamy ich dowodów.

\begin{twierdzenie}[Słaba zasada maksimum]
\label{thm:pde14-weak-max}
Niech $u\in H^1(\Omega)$ będzie podrozwiązaniem $Lu\le0$,
przy powyższych założeniach o $A$ i $c$. Niech $M\ge0$ będzie
liczbą taką, że $\operatorname{Tr}u\le M$ prawie wszędzie na
$\partial\Omega$. Wówczas $u\le M$ prawie wszędzie w $\Omega$.
Jeżeli $c=0$, ten sam wniosek jest prawdziwy dla dowolnego
$M\in\R$ spełniającego warunek na ślad.
\end{twierdzenie}
\begin{proof}
Kładziemy $w=(u-M)^+$. Z lematu~\ref{lem:pde14-trunc}
$w\in H_0^1(\Omega)$ i możemy użyć go jako nieujemnej
funkcji testowej. Na zbiorze $\{u>M\}$ mamy
$\nabla w=\nabla u$ i $u=M+w$. Zatem
\[
 0\ge \int_\Omega(A\nabla u\cdot\nabla w+cuw)\,
 \mathrm d\lambda_d(x)
 \ge \lambda\|\nabla w\|_2^2
      +\int_{\{u>M\}}c(M+w)w\,\mathrm d\lambda_d(x).
\]
Ponieważ $c\ge0$, $M\ge0$ i $w\ge0$, ostatnia całka jest
nieujemna. Otrzymujemy $\nabla w=0$, a nierówność Poincarégo
dla $H_0^1(\Omega)$ daje $w=0$. Jeśli $c=0$, składnik
zerowego rzędu znika, więc znak $M$ nie odgrywa roli.
\end{proof}

\begin{wniosek}[Porównanie]
\label{cor:pde14-compare}
Niech $u,v\in H^1(\Omega)$, $c\ge0$, a
$L=-\diver(A\nabla\cdot)+c$. Jeżeli $Lu\le Lv$ słabo w
$\Omega$ i $\operatorname{Tr}u\le\operatorname{Tr}v$ prawie
wszędzie na $\partial\Omega$, to $u\le v$ prawie wszędzie
w $\Omega$.
\end{wniosek}
\begin{proof}
Różnica $w=u-v$ jest podrozwiązaniem $Lw\le0$ i ma
niedodatni ślad. Stosujemy twierdzenie~\ref{thm:pde14-weak-max}
z $M=0$. Nie wykorzystujemy symetrii macierzy $A$.
\end{proof}

\begin{uwaga}[Granice wniosku]
Przy członie $b\cdot\nabla u$ poprzedni test zawiera
$\int_\Omega(b\cdot\nabla w)w\,\mathrm d\lambda_d$,
którego znak nie jest określony. Klasyczne zasady maksimum
dla operatorów z dryfem wymagają dodatkowych argumentów;
nie wynikają bezpośrednio z powyższego dowodu. Nie wolno
również pomijać założenia $c\ge0$.
\end{uwaga}

\begin{kontrprzyklad}[Ujemny człon zerowego rzędu]
\label{ex:pde14-negative-c}
Na $\Omega=(0,1)$ funkcja $u(x)=\sin(\pi x)$ jest dodatnia
we wnętrzu i ma zerowy ślad, choć spełnia
$-u''-\pi^2u=0$. Gdyby dla $c=-\pi^2$ zachodziła zasada
maksimum z twierdzenia~\ref{thm:pde14-weak-max},
wynikałoby $u\le0$, co jest fałszem. Jest to jednocześnie
przykład rezonansu z rozdziału~\ref{chap:pde13}.
\end{kontrprzyklad}

\section{Funkcje harmoniczne i mocna zasada maksimum}
\label{sec:pde14-strong}
Funkcją harmoniczną w otwartym zbiorze $U\subset\R^d$
nazywamy funkcję $u\in C^2(U)$ spełniającą $\Delta u=0$.
Z lematu Weyla (rozdział~\ref{chap:pde-distributions}) wiemy,
że również rozwiązanie dystrybucyjne równania $\Delta u=0$
ma reprezentanta gładkiego; dlatego poniższe twierdzenia
stosują się także do takich rozwiązań.

\begin{twierdzenie}[Mocna zasada maksimum dla laplasjanu]
\label{thm:pde14-strong}
Niech $U\subset\R^d$ będzie obszarem, a $u\in C^2(U)$ niech spełnia $\Delta u=0$ w $U$.
Jeżeli $u$ osiąga maksimum lub minimum w pewnym punkcie
$x_0\in U$, to $u$ jest stała na $U$.
\end{twierdzenie}
\begin{proof}
Załóżmy, że $u(x_0)=M=\sup_U u$. Dla dowolnego
$r>0$ takiego, że $\overline{B_r(x_0)}\subset U$, własność
wartości średniej daje
$M=u(x_0)=\frac{1}{\sigma(\partial B_r(x_0))}\int_{\partial B_r(x_0)}u\,\mathrm d\sigma$.
Funkcja $M-u$ jest nieujemna i ciągła na sferze.
Jeżeli byłaby dodatnia w choć jednym jej punkcie, to
byłaby dodatnia na pewnym otwartym fragmencie sfery,
a średnia $u$ byłaby ściśle mniejsza od $M$.
Tak więc $u\equiv M$ na każdej takiej sferze, a zatem
również w całej małej kuli.
Zbiór $E=\{x\in U:u(x)=M\}$ jest niepusty i domknięty
w $U$ z ciągłości. Ten sam argument zastosowany do każdego
$x\in E$ pokazuje, że $E$ jest też otwarty w $U$.
Spójność $U$ wymusza $E=U$. Dla minimum stosujemy wynik do $-u$.
\end{proof}

\begin{wniosek}[Dodatniość i jednoznaczność]
\label{cor:pde14-positive}
Jeżeli $U$ jest spójna, a $u\ge0$ jest harmoniczna w $U$,
to $u\equiv0$ albo $u(x)>0$ dla wszystkich $x\in U$.
Jeżeli $U$ jest ponadto ograniczona oraz
$u,v\in C^2(U)\cap C(\overline U)$ są harmoniczne i mają
te same wartości na $\partial U$, to $u=v$ na $U$.
\end{wniosek}
\begin{proof}
Jeżeli $u(x_0)=0$ we wnętrzu, funkcja osiąga tam minimum,
więc jest stała, a z nieujemności wynika $u=0$.
Dla jednoznaczności $w=u-v$ jest harmoniczna i znika na
brzegu. Ciągłość na zwartym $\overline U$ daje istnienie
maksimum i minimum; gdyby którakolwiek ekstremalna
wartość była różna od zera, ekstremum osiągnięte byłoby
we wnętrzu, co w myśl twierdzenia~\ref{thm:pde14-strong}
wymusiłoby stałość, sprzeczną z zerowym śladem.
\end{proof}

\section{Lemat Hopfa: znak pochodnej normalnej}
\label{sec:pde14-hopf}
Słaba zasada maksimum mówi, gdzie może znajdować się
minimum. Lemat Hopfa wyjaśnia, jak funkcja zachowuje się
przy zbliżaniu do punktu minimum na brzegu.

\begin{definicja}[Warunek kuli wewnętrznej w punkcie]
Niech $U\subset\R^d$ będzie otwarta, a $x_0\in\partial U$.
Mówimy, że w $x_0$ zachodzi \emph{warunek kuli wewnętrznej},
jeżeli istnieją $z\in U$ i $R>0$ takie, że
$B_R(z)\subset U$ oraz $x_0\in\partial B_R(z)$.
Wektor $\nu=(x_0-z)/R$ jest zewnętrzną normalną do tej kuli.
\end{definicja}

\begin{twierdzenie}[Lemat Hopfa dla funkcji harmonicznych]
\label{thm:pde14-hopf}
Niech $U\subset\R^d$ będzie obszarem ograniczonym,
$x_0\in\partial U$ i niech w $x_0$ zachodzi warunek kuli
wewnętrznej $B_R(z)\subset U$. Załóżmy, że
$u\in C^2(U)\cap C(\overline U)$ jest harmoniczna, nie jest
stała, a $u(x_0)=\min_{\overline U}u$.
Wówczas dla kierunku $-\nu=(z-x_0)/R$ zachodzi
\[
 \liminf_{t\downarrow0}\frac{u(x_0-t\nu)-u(x_0)}{t}>0.
\]
Jeżeli dodatkowo $u$ ma w $x_0$ pochodną kierunkową w
kierunku $\nu$ (na przykład $u\in C^1(\overline U)$
w otoczeniu $x_0$), to $\partial_\nu u(x_0)<0$.
\end{twierdzenie}
\begin{proof}
Odejmując stałą, zakładamy $u(x_0)=0$ i $u\ge0$ na
$\overline U$. Z mocnej zasady minimum
$u>0$ wewnątrz $U$. W pierścieniu
$\mathcal A=\{x:R/2<|x-z|<R\}$ rozważmy funkcję
\[
 w(x)=e^{-\alpha|x-z|^2}-e^{-\alpha R^2},\qquad \alpha>0.
\]
Jest ona dodatnia wewnątrz pierścienia, a na jego zewnętrznej
sferze wynosi zero. Rachunek daje
$\Delta w=(4\alpha^2|x-z|^2-2d\alpha)e^{-\alpha|x-z|^2}$.
Jeśli $\alpha>2d/R^2$, to $\Delta w>0$ na
$\overline{\mathcal A}$. Z ciągłości i dodatniości $u$
na zwartej wewnętrznej sferze
$m:=\min_{|x-z|=R/2}u(x)>0$. Wybieramy
$0<\varepsilon\le m/\max_{|x-z|=R/2}w(x)$.
Wtedy $u-\varepsilon w\ge0$ na obu składowych brzegu
pierścienia: na wewnętrznej z wyboru $\varepsilon$, na
zewnętrznej z $u\ge0$ oraz $w=0$. Ponieważ
$\Delta(u-\varepsilon w)=-\varepsilon\Delta w<0$,
słaba zasada minimum dla funkcji superharmonicznych
daje $u\ge\varepsilon w$ na $\overline{\mathcal A}$.
Po podstawieniu $x=x_0-t\nu$ otrzymujemy
\[
 \frac{u(x_0-t\nu)}t\ge
 \varepsilon\frac{e^{-\alpha(R-t)^2}-e^{-\alpha R^2}}t
 \longrightarrow 2\varepsilon\alpha R e^{-\alpha R^2}>0.
\]
W przypadku istnienia pochodnej normalnej lewa strona
zbiega do $-\partial_\nu u(x_0)$, co daje tezę.
\end{proof}

\begin{przyklad}[Konieczność geometrii brzegu]
\label{ex:pde14-cusp}
W klinie płaskim $U=\{(r\cos\theta,r\sin\theta):
0<r<1,\ 0<\theta<\omega\}$, gdzie $0<\omega<\pi$,
funkcja $u(r,\theta)=r^{\pi/\omega}\sin(\pi\theta/\omega)$
jest dodatnia i harmoniczna we wnętrzu oraz znika na obu
ramionach klina. Przy wierzchołku $r=0$ ma minimum,
ale jej przyrost w każdym ustalonym kierunku wewnętrznym
jest rzędu $r^{\pi/\omega}=o(r)$. W tym punkcie
nie zachodzi warunek kuli wewnętrznej: nie wolno
stosować lematu Hopfa z dodatnią liniową dolną granicą.
\end{przyklad}

\section{Jądro Poissona i ilościowe porównywanie wartości}
\label{sec:pde14-poisson}
Oznaczmy przez $\omega_{d-1}$ pole powierzchni
jednostkowej sfery $S^{d-1}$ dla $d\ge2$. W kuli
$B_R=B_R(0)\subset\R^d$ jądro Poissona ma postać
\[
 P_R(x,y)=\frac{R^2-|x|^2}{\omega_{d-1}R|x-y|^d},
 \qquad x\in B_R,\quad y\in\partial B_R.
\]
Jest dodatnie i spełnia
$\int_{\partial B_R}P_R(x,y)\,\mathrm d\sigma(y)=1$.
Dla $g\in C(\partial B_R)$ funkcja
$u(x)=\int_{\partial B_R}P_R(x,y)g(y)\,\mathrm d\sigma(y)$
jest harmoniczna i przyjmuje dane brzegowe $g$ w sensie
ciągłego przedłużenia na $\overline{B_R}$; dowód tych
faktów przedstawiono w rozdziale~\ref{chap:pde-laplace}.

\begin{twierdzenie}[Nierówność Harnacka w kuli]
\label{thm:pde14-harnack}
Niech $d\ge2$, $0<r<R$, a $u\ge0$ będzie funkcją
harmoniczną w otoczeniu $\overline{B_R(0)}$.
Wtedy dla każdego $x\in\overline{B_r(0)}$ zachodzi
\begin{equation}\label{eq:pde14-harnack}
 \frac{R^2-r^2}{R^2}\left(\frac{R}{R+r}\right)^d u(0)
 \le u(x)\le
 \left(\frac{R}{R-r}\right)^d u(0).
\end{equation}
W szczególności $\sup_{B_r}u\le C_{d,r/R}\inf_{B_r}u$
przy stałej zależnej wyłącznie od $d$ i $r/R$.
\end{twierdzenie}
\begin{proof}
Przez dodatniość $u$ i $P_R$ wolno porównywać jądra
punktowo pod znakiem całki Poissona. Ponieważ
$R-r\le|x-y|\le R+r$ dla $|x|\le r$, $|y|=R$,
oraz $R^2-r^2\le R^2-|x|^2\le R^2$, otrzymujemy
\[
 \frac{R^2-r^2}{R^2}\left(\frac{R}{R+r}\right)^d
 \le\frac{P_R(x,y)}{P_R(0,y)}\le
 \left(\frac{R}{R-r}\right)^d.
\]
Całkujemy po $\partial B_R$ z nieujemną wagą $u(y)$.
Po lewej i prawej zostaje $u(0)$, a w środku $u(x)$.
Jeżeli $u(0)=0$, mocna zasada minimum daje $u\equiv0$;
w innym przypadku dolne ograniczenie w~\eqref{eq:pde14-harnack}
pozwala oszacować $u(0)$ przez $\inf_{B_r}u$, a górne
$\sup_{B_r}u$ przez $u(0)$. Dokładniej, można wziąć
$C_{d,r/R}=(R/(R-r))^d\cdot R^2/(R^2-r^2)
\cdot((R+r)/R)^d$.
\end{proof}

\begin{uwaga}[Usunięcie założenia o przedłużeniu]
Jeżeli $u\ge0$ jest harmoniczna jedynie w $B_R$, bez
ciągłości na $\partial B_R$, stosujemy twierdzenie w
kulach $B_\rho$ o $r<\rho<R$, a następnie przechodzimy
do granicy $\rho\uparrow R$. Stałe w oszacowaniach
są ciągłe względem promienia $\rho$, więc otrzymujemy
te same nierówności. Nie używamy nieistniejącego śladu
na $\partial B_R$.
\end{uwaga}

\begin{wniosek}[Nierówność Harnacka na zwartych podzbiorach]
\label{cor:pde14-harnack-chain}
Niech $U\subset\R^d$ będzie obszarem,
$K\subset U$ będzie zwarty i niepusty, a $u\ge0$ harmoniczna
w $U$. Istnieje stała $C=C(U,K,d)>0$, niezależna od $u$,
taka że $\sup_Ku\le C\inf_Ku$.
\end{wniosek}
\begin{proof}
Wybieramy punkt $x_*\in K$. Każdy $x\in K$ można
połączyć z $x_*$ łamaną o zwartym obrazie zawartym w $U$:
w otwartej, spójnej części $\R^d$ istnieją takie łamane.
Dla każdej z tych łamanych wybieramy skończony łańcuch
kul $B_{2\rho_j}(z_j)\Subset U$ z nakładającymi się
kulami $B_{\rho_j}(z_j)$. Nierówność Harnacka w jednej kuli
porównuje wszystkie wartości w jej mniejszej kuli,
a więc porównuje także punkty sąsiedniego ogniwa łańcucha.
Otrzymujemy oszacowanie $u(x)\le C_xu(x_*)$ oraz
$u(x_*)\le C_xu(x)$ w pewnym otoczeniu $x$.
Zwartość $K$ pozwala wybrać skończenie wiele takich
otoczeń i największą ze skończenie wielu stałych.
Dla wszystkich $x,y\in K$ wynika
$u(x)\le C_xC_yu(y)\le C\,u(y)$; przejście do kresów
daje tezę. Jeżeli $u\equiv0$, nierówność jest oczywista.
\end{proof}

\begin{wniosek}[Twierdzenie Liouville'a z ograniczeniem jednostronnym]
\label{cor:pde14-liouville}
Jeżeli $u$ jest harmoniczna na $\R^d$, $d\ge2$, i jest
ograniczona z dołu albo z góry, to $u$ jest stała.
\end{wniosek}
\begin{proof}
Po odjęciu dolnego ograniczenia zakładamy $u\ge0$.
Dla ustalonego $x\in\R^d$ wybieramy $R>|x|$ i stosujemy
\eqref{eq:pde14-harnack} z $r=|x|$ (dla $x=0$ teza jest
oczywista). Gdy $R\to\infty$, oba współczynniki przy $u(0)$
dążą do $1$. Z twierdzenia o trzech ciągach wynika
$u(x)=u(0)$. Gdy $u$ jest ograniczona z góry, stosujemy
powyższe do $-u$ po odjęciu stosownej stałej.
\end{proof}

\section{Oszacowania pochodnych i lokalne wygładzanie}
\label{sec:pde14-deriv}
\begin{twierdzenie}[Oszacowanie gradientu funkcji harmonicznej]
\label{thm:pde14-gradient}
Niech $d\ge2$, $R>0$, $x_0\in\R^d$ i niech $u$ będzie
harmoniczna w otoczeniu $\overline{B_R(x_0)}$.
Wtedy
\[
 |\nabla u(x_0)|\le \frac dR
 \sup_{B_R(x_0)}|u|.
\]
W szczególności, jeżeli $0<r<R$ i $u$ jest harmoniczna
w $B_R(x_0)$, to
$\sup_{B_r(x_0)}|\nabla u|\le
\frac{2d}{R-r}\sup_{B_R(x_0)}|u|$.
\end{twierdzenie}
\begin{proof}
Możemy przyjąć $x_0=0$. Różniczkujemy formułę Poissona
względem $x$ w punkcie $0$; operacja jest legalna, gdyż
jądro oraz jego pochodne są ciągłe na
$\overline{B_{R/2}}\times\partial B_R$, a sfera jest zwarta.
Ponieważ $P_R(0,y)=1/(\omega_{d-1}R^{d-1})$ i
$\nabla_x P_R(0,y)=d y/(\omega_{d-1}R^{d+1})$, mamy
\[
 \nabla u(0)=\frac{d}{\omega_{d-1}R^{d+1}}
 \int_{\partial B_R}y\,u(y)\,\mathrm d\sigma(y).
\]
Wykorzystując $|y|=R$ oraz
$\sigma(\partial B_R)=\omega_{d-1}R^{d-1}$,
otrzymujemy żądane ograniczenie. Dla $x\in B_r(x_0)$
możemy zastosować pierwszy wynik do kuli o środku $x$
i promieniu $(R-r)/2$, której domknięcie leży w $B_R(x_0)$.
Stąd druga nierówność. Ponieważ dowolną lokalną pochodną
funkcji harmonicznej można ponownie traktować jako
funkcję harmoniczną, iteracja argumentu prowadzi do
oszacowań wyższych pochodnych na mniejszych kulach.
\end{proof}

\section{Subharmoniczność i nierówność wartości średniej}
\label{sec:pde14-subharmonic}
\begin{definicja}[Funkcja subharmoniczna w sensie klasycznym]
Niech $U\subset\R^d$ będzie otwarta, a $u\in C^2(U)$.
Nazywamy $u$ \emph{subharmoniczną}, jeśli $\Delta u\ge0$
w $U$, zaś \emph{superharmoniczną}, jeśli $\Delta u\le0$.
Istnieje również ogólniejsza definicja dla funkcji
półciągłych i ich odpowiedników dystrybucyjnych; w tym
rozdziale podajemy twierdzenia i dowody w klasie $C^2$.
\end{definicja}

\begin{twierdzenie}[Nierówność wartości średniej]
\label{thm:pde14-sub-mean}
Jeżeli $u\in C^2(U)$ jest subharmoniczna i
$\overline{B_R(x_0)}\subset U$, to dla każdego $0<r\le R$
\[
 u(x_0)\le \frac{1}{\sigma(\partial B_r(x_0))}
 \int_{\partial B_r(x_0)}u(y)\,\mathrm d\sigma(y).
\]
Dla funkcji superharmonicznej nierówność ma przeciwny zwrot.
\end{twierdzenie}
\begin{proof}
Niech $m(r)=\frac{1}{\sigma(\partial B_r(x_0))}\int_{\partial B_r(x_0)}u\,\mathrm d\sigma$.
Po podstawieniu $y=x_0+r\theta$, $\theta\in S^{d-1}$,
i różniczkowaniu pod znakiem całki otrzymujemy
\[
 m'(r)=\frac1{\omega_{d-1}}
 \int_{S^{d-1}}\nabla u(x_0+r\theta)\cdot\theta\,
 \mathrm d\sigma(\theta)
 =\frac1{\omega_{d-1}r^{d-1}}
 \int_{B_r(x_0)}\Delta u(y)\,\mathrm d\lambda_d(y).
\]
Druga równość jest twierdzeniem Gaussa zastosowanym do
$\nabla u$. Dla $\Delta u\ge0$ pochodna $m'$ jest nieujemna.
Ponieważ $\lim_{r\downarrow0}m(r)=u(x_0)$,
otrzymujemy tezę. Dla $\Delta u\le0$ rozumowanie jest
identyczne z odwróconymi nierównościami.
\end{proof}

\begin{uwaga}[Wersja dystrybucyjna]
Dla $u\in L^1_{\mathrm{loc}}(U)$ warunek
$\Delta u\ge0$ w $\mathcal D'(U)$ oznacza
$\langle\Delta u,\varphi\rangle\ge0$ dla nieujemnych
$\varphi\in C_c^\infty(U)$. Po wygładzeniu przez splot
otrzymujemy gładkie funkcje subharmoniczne.
Do zdefiniowania wartości $u(x_0)$ i nierówności średniej
należy następnie wybrać odpowiedniego, półciągłego z góry
reprezentanta. Nie wolno bez tego kroku przypisywać
wartości punktowych dowolnej klasie $L^1_{\mathrm{loc}}$.
\end{uwaga}

\begin{twierdzenie}[Silna zasada minimum dla funkcji superharmonicznych]
\label{thm:pde14-strong-super}
Niech $U\subset\R^d$ będzie obszarem i niech
$u\in C^2(U)$ spełnia $\Delta u\le0$ w $U$.
Jeżeli $u$ osiąga minimum globalne względem $U$ w punkcie
$x_0\in U$, to jest stała na $U$.
\end{twierdzenie}
\begin{proof}
Niech $m=u(x_0)=\inf_Uu$. Z twierdzenia
\ref{thm:pde14-sub-mean}, zastosowanego do $-u$, wynika
$u(x_0)\ge\frac{1}{\sigma(\partial B_r(x_0))}\int_{\partial B_r(x_0)}u\,\mathrm d\sigma$
dla każdej kuli o domknięciu zawartym w $U$.
Ponieważ $u\ge m$, średnia wynosi jednocześnie co
najmniej $m$; ciągłość wymusza $u=m$ na każdej
z tych sfer, a więc w małej kuli o środku $x_0$.
Zbiór $\{x\in U:u(x)=m\}$ jest zatem otwarty, a dzięki
ciągłości również domknięty w spójnym $U$. Jest niepusty,
więc równa się $U$.
\end{proof}

\begin{wniosek}[Lemat Hopfa dla nierówności $-\Delta u\ge0$]
\label{cor:pde14-hopf-super}
Niech $U\subset\R^d$ będzie obszarem ograniczonym,
$u\in C^2(U)\cap C(\overline U)$, $\Delta u\le0$ w $U$.
Załóżmy, że $u$ nie jest stała,
$u(x_0)=\min_{\overline U}u$ dla $x_0\in\partial U$
i istnieje kula $B_R(z)\subset U$ styczna do brzegu w $x_0$.
Wówczas, przy $\nu=(x_0-z)/R$,
\[
 \liminf_{t\downarrow0}\frac{u(x_0-t\nu)-u(x_0)}t>0.
\]
Jeśli istnieje pochodna normalna w $x_0$, jest ona ujemna.
\end{wniosek}
\begin{proof}
Z twierdzenia~\ref{thm:pde14-strong-super} wynika,
że $u(x)>u(x_0)$ dla każdego $x\in U$. Odejmujemy
$u(x_0)$ i używamy tej samej bariery
$w(x)=e^{-\alpha|x-z|^2}-e^{-\alpha R^2}$,
co w dowodzie twierdzenia~\ref{thm:pde14-hopf}.
Na pierścieniu $R/2<|x-z|<R$ mamy $\Delta w>0$
dla dostatecznie dużego $\alpha$, a więc
$\Delta(u-\varepsilon w)<0$. Na obu sferach brzegowych
$u-\varepsilon w\ge0$ po wyborze małego
$\varepsilon>0$. Słaba zasada minimum dla funkcji
superharmonicznych, wynikająca z
\ref{thm:pde14-sub-mean}, daje $u\ge\varepsilon w$.
Wzdłuż promienia $x=x_0-t\nu$ otrzymujemy granicę
$2\varepsilon\alpha R e^{-\alpha R^2}>0$, tak jak
w dowodzie twierdzenia~\ref{thm:pde14-hopf}.
\end{proof}



\section{Nierówność Caccioppolego dla współczynników mierzalnych}
\label{sec:pde14-caccioppoli}
Przechodzimy ponownie do form dywergencyjnych, ale
bez założenia gładkości współczynników.

\begin{twierdzenie}[Oszacowanie wewnętrzne Caccioppolego]
\label{thm:pde14-caccioppoli}
Niech $\Omega\subset\R^d$ będzie otwarta,
$A\in L^\infty(\Omega;\R^{d\times d})$,
$\xi^{\mathsf T}A(x)\xi\ge\lambda|\xi|^2$ i
$|A(x)\xi|\le\Lambda|\xi|$ prawie wszędzie dla
$\lambda,\Lambda>0$. Jeżeli $u\in H^1_{\mathrm{loc}}(\Omega)$
spełnia $-\diver(A\nabla u)=0$ słabo, to dla każdej
$\eta\in C_c^\infty(\Omega)$ zachodzi
\begin{equation}\label{eq:pde14-caccioppoli}
 \int_\Omega\eta^2|\nabla u|^2\,\mathrm d\lambda_d
 \le \frac{4\Lambda^2}{\lambda^2}
 \int_\Omega u^2|\nabla\eta|^2\,\mathrm d\lambda_d.
\end{equation}
W szczególności, gdy $B_{2r}(x_0)\Subset\Omega$,
\[
 \int_{B_r(x_0)}|\nabla u|^2\,\mathrm d\lambda_d
 \le \frac{C\Lambda^2}{\lambda^2r^2}
 \int_{B_{2r}(x_0)}|u|^2\,\mathrm d\lambda_d,
\]
gdzie $C$ zależy wyłącznie od wyboru standardowego obcięcia.
\end{twierdzenie}
\begin{proof}
Ponieważ $\eta$ ma zwarty nośnik w $\Omega$,
możemy użyć $\varphi=\eta^2u\in H_0^1(\Omega)$ jako
funkcji testowej (w razie potrzeby aproksymując ją
gładkimi funkcjami zwartonośnymi). Z reguły iloczynu
$\nabla(\eta^2u)=\eta^2\nabla u+2\eta u\nabla\eta$.
Równanie słabe daje
\[
 \int_\Omega\eta^2 A\nabla u\cdot\nabla u\,\mathrm d\lambda_d
 =-2\int_\Omega\eta u A\nabla u\cdot\nabla\eta\,\mathrm d\lambda_d(x).
\]
Eliptyczność i ograniczoność macierzy prowadzą do
$\lambda X^2\le2\Lambda XY$, gdzie
\[
 X^2=\int_\Omega\eta^2|\nabla u|^2\,\mathrm d\lambda_d(x),
 \qquad
 Y^2=\int_\Omega u^2|\nabla\eta|^2\,\mathrm d\lambda_d(x).
\]
Gdy $X=0$, teza jest oczywista. W przeciwnym razie
$X\le 2\Lambda Y/\lambda$, co po podniesieniu do kwadratu
daje~\eqref{eq:pde14-caccioppoli}.
Wybierając funkcję $\eta$ równą $1$ na $B_r(x_0)$,
zerową poza $B_{2r}(x_0)$ i z $|\nabla\eta|\le C_0/r$,
otrzymujemy drugie oszacowanie.
\end{proof}

\begin{uwaga}[Dlaczego nie dowodzimy tu ogólnej nierówności Harnacka?]
Oszacowanie Caccioppolego jest punktem wyjścia do iteracji
De Giorgiego i Mosera. Pełna nierówność Harnacka dla
nieujemnych słabych rozwiązań przy mierzalnych współczynnikach
jest wynikiem znacznie głębszym niż nierówność dla funkcji
harmonicznych wyprowadzona z jądra Poissona. Przedstawimy
jej dowód w rozdziale poświęconym twierdzeniom
De Giorgiego--Nasha--Mosera; nie używamy jej tutaj jako
nieuzasadnionego lematu.
\end{uwaga}

\section{Funkcja Greena i porównanie dla problemu Dirichleta}
\label{sec:pde14-green}
Niech $d\ge2$ oraz $\Phi$ będzie rozwiązaniem fundamentalnym
operatora $-\Delta$ w $\R^d$:
\[
 \Phi(z)=\begin{cases}
 \dfrac{|z|^{2-d}}{(d-2)\omega_{d-1}},&d\ge3,\\[4pt]
 -\dfrac1{2\pi}\log|z|,&d=2.
 \end{cases}
\]
W obu przypadkach $-\Delta\Phi=\delta_0$ w sensie
dystrybucyjnym. Nie utożsamiamy jądra Greena z samym
rozwiązaniem fundamentalnym: warunek brzegowy wymaga
odjęcia części harmonicznej.

\begin{definicja}[Funkcja Greena Dirichleta]
Niech $\Omega\subset\R^d$ będzie obszarem ograniczonym,
w którym dla każdego $y\in\Omega$ istnieje funkcja
$h_y\in C(\overline\Omega)\cap C^2(\Omega)$ harmoniczna
w $\Omega$ i taka, że $h_y(x)=\Phi(x-y)$ dla
$x\in\partial\Omega$. Definiujemy
$G_\Omega(x,y)=\Phi(x-y)-h_y(x)$ dla $x\ne y$.
Wówczas $-\Delta_xG_\Omega(x,y)=\delta_y$ dystrybucyjnie,
a $G_\Omega(\cdot,y)=0$ na brzegu.
\end{definicja}

\begin{twierdzenie}[Dodatniość funkcji Greena]
\label{thm:pde14-green-pos}
Przy powyższych założeniach dla każdego $y\in\Omega$
$G_\Omega(x,y)>0$ dla $x\in\Omega\setminus\{y\}$.
\end{twierdzenie}
\begin{proof}
Dla małego $\rho>0$ kula $\overline{B_\rho(y)}$ mieści
się w $\Omega$. Ponieważ $h_y$ jest ograniczona na
$\overline\Omega$, podczas gdy $\Phi(x-y)\to+\infty$
gdy $x\to y$ (również przy $d=2$), istnieje $\rho_0>0$
takie, że $G_\Omega(x,y)>0$ dla $|x-y|=\rho\le\rho_0$.
Na zewnętrznej granicy $\partial\Omega$ funkcja $G$ znika.
W obszarze pierścieniowym
$\Omega\setminus\overline{B_\rho(y)}$ jest harmoniczna;
słaba zasada minimum daje $G\ge0$ wszędzie w tym obszarze. Mocna zasada minimum daje ścisłą dodatniość.
Istotnie, $\Omega\setminus\overline{B_\rho(y)}$ jest spójna:
dowolną łamaną w $\Omega$ przecinającą kulę można zmodyfikować
wzdłuż sfery o promieniu nieco większym niż $\rho$,
która nadal leży w $\Omega$. Dowolny $x\ne y$
możemy objąć takim pierścieniem, wybierając $\rho<|x-y|$.
\end{proof}

\begin{lemat}[Jawna funkcja Greena kuli i jej pochodna normalna]
\label{lem:pde14-green-ball}
Niech $d\ge2$, $R>0$ i $x\in B_R\setminus\{0\}$.
Zapiszmy $x^*=R^2x/|x|^2$, a więc $|x^*|>R$. Dla
$y\in B_R$, $y\ne x$, funkcja Greena ma postać
\[
 G_{B_R}(x,y)=\begin{cases}
 \Phi(x-y)-(R/|x|)^{d-2}\Phi(x^*-y),&d\ge3,\\[3pt]
 \dfrac1{2\pi}\log\left(\dfrac{|x|\,|x^*-y|}
 {R|x-y|}\right),&d=2.
 \end{cases}
\]
W punkcie $x=0$ rozumiemy ją przez granicę. Jądro jest
symetryczne: $G_{B_R}(x,y)=G_{B_R}(y,x)$. Ponadto
$G_{B_R}(x,y)=0$ dla $|y|=R$, a na brzegu zachodzi
$-\partial_{\nu_y}G_{B_R}(x,y)=P_R(x,y)$, gdzie
$\nu_y=y/R$.
\end{lemat}
\begin{proof}
Dla $|y|=R$ otrzymujemy tożsamość geometryczną
$|x^*-y|=(R/|x|)|x-y|$. Po jej podstawieniu oba wzory
na $G$ przyjmują wartość zero. W zmiennej $y$ pierwszy
składnik jest harmoniczny poza biegunem $x$, drugi zaś
jest harmoniczny w całej kuli, ponieważ jego biegun
$x^*$ leży na zewnątrz. W pobliżu $y=x$ pierwszy
składnik jest rozwiązaniem fundamentalnym $-\Delta$,
więc $-\Delta_yG=\delta_x$.
Symetria wynika z równości
\[
 |x|^2|x^*-y|^2=R^4-2R^2x\cdot y+|x|^2|y|^2,
\]
której prawa strona jest symetryczna w $x,y$.

Dla $d\ge3$ wykorzystujemy
$\nabla_y\Phi(x-y)=(x-y)/(\omega_{d-1}|x-y|^d)$.
Różniczkując na zewnętrznej sferze i używając ponownie
$|x^*-y|=(R/|x|)|x-y|$, otrzymujemy
\begin{align*}
 \partial_{\nu_y}G(x,y)
 &=\frac{x\cdot y/R-R}{\omega_{d-1}|x-y|^d}
 -\frac{(|x|/R)^2(x^*\cdot y/R-R)}
 {\omega_{d-1}|x-y|^d}\\
 &=\frac{|x|^2-R^2}{\omega_{d-1}R|x-y|^d}
 =-P_R(x,y).
\end{align*}
Dla $d=2$ pochodna logarytmu spełnia tę samą
formułę, gdyż $\omega_1=2\pi$ i wykładnik $d-2$ wynosi
zero; składnik stały względem $y$ nie wnosi pochodnej.
Przejście $x\to0$ daje pozostały przypadek.
\end{proof}

\begin{twierdzenie}[Reprezentacja Greena w kuli]
\label{thm:pde14-green-repr}
Niech $d\ge2$, $R>0$, $f\in C(\overline{B_R})$,
$g\in C^2(\partial B_R)$, a
$u\in C^2(B_R)\cap C^1(\overline{B_R})$ spełnia
$-\Delta u=f$ w $B_R$ i $u=g$ na $\partial B_R$.
Wówczas dla $x\in B_R$ zachodzi
\begin{equation}\label{eq:pde14-green-repr}
 u(x)=\int_{B_R}G_{B_R}(x,y)f(y)\,\mathrm d\lambda_d(y)
 +\int_{\partial B_R}P_R(x,y)g(y)\,\mathrm d\sigma(y).
\end{equation}
W szczególności dla $f\ge0$ i $g\ge0$ mamy $u\ge0$.
\end{twierdzenie}
\begin{proof}
Dla ustalonego $x$ stosujemy drugą tożsamość Greena
w $B_R\setminus\overline{B_\rho(x)}$ do pary
$u(y)$ i $G_{B_R}(x,y)$. W tym obszarze $G$ jest
harmoniczna i $\Delta u=-f$. Otrzymujemy
\[
 \int_{B_R\setminus B_\rho(x)}G_{B_R}(x,y)f(y)
 \,\mathrm d\lambda_d(y)
 =\int_{\partial B_R}g\,\partial_{\nu_y}G\,
 \mathrm d\sigma+I_\rho,
\]
gdzie $I_\rho$ to wyraz z małej sfery, z normalną
skierowaną do jej środka. Ponieważ
$G(x,y)=\Phi(x-y)+O(1)$ w pobliżu $y=x$ oraz
$\partial_{\nu}G=\partial_{\nu}\Phi+O(1)$,
bezpośrednie obliczenie strumienia $\Phi$ daje
$\lim_{\rho\downarrow0}I_\rho=u(x)$.
Ponadto całka objętościowa przechodzi do granicy:
osobliwość $\Phi$ jest lokalnie całkowalna dla $d\ge2$.
Tożsamość $-\partial_{\nu_y}G_{B_R}(x,y)=P_R(x,y)$
została jawnie wyprowadzona w lemacie~\ref{lem:pde14-green-ball}.
Po przeniesieniu wyrazów otrzymujemy~\eqref{eq:pde14-green-repr}.
Wniosek o znaku wynika z dodatniości obu jąder.
\end{proof}

\begin{uwaga}[Uwaga o znakach we wzorze Greena]
W powyższym dowodzie normalna na \emph{wewnętrznej}
sferze $\partial B_\rho(x)$ jest skierowana \emph{do środka}
wyciętej kuli. To właśnie dlatego wyraz z tej sfery
ma znak prowadzący do $+u(x)$ w reprezentacji.
Rozpisanie orientacji na każdej granicy chroni przed
częstą pomyłką znaku przy przejściu od $\Phi$ do $G$.
\end{uwaga}

\section{Rysunek: bariera przy brzegu i geometria nierówności}
\label{sec:pde14-figure}
\begin{center}
\begin{tikzpicture}[scale=0.88,>=Latex]
  % Dwie kule styczne wewnetrznie: x0 nalezy do obu brzegow.
  \fill[blue!5] (-0.6,0) circle (2.45);
  \draw[thick,blue!65!black] (-0.6,0) circle (2.45);
  \draw[thick,orange!90!black] (0,0) circle (1.85);
  \fill (0,0) circle(2pt) node[below left] {$z$};
  \fill (1.85,0) circle(2.5pt) node[above right] {$x_0$};
  \draw[dashed] (0,0)--(1.85,0) node[midway,above] {$R$};
  \draw[-{Latex[length=2.5mm]},very thick] (1.85,0) -- (2.65,0)
    node[right] {$\nu$};
  \draw[-{Latex[length=2.5mm]},very thick,red!75!black]
    (1.78,-0.28) -- (0.90,-0.28);
  \node[red!70!black,below] at (0.15,-0.65) {kierunek $-\nu$};
  \node[blue!70!black] at (-2.3,1.25) {$\Omega$};
  \node[orange!85!black] at (-0.7,1.20) {$B_R(z)$};
\end{tikzpicture}

{\small Geometria lematu Hopfa: kula wewnętrzna dotyka brzegu
$\partial\Omega$ w $x_0$, a bariera wykrywa dodatni przyrost
w kierunku do wnętrza $-\nu$.}
\end{center}

\section{Podsumowanie i przejście do teorii regularności}
\label{sec:pde14-summary}
Wykazaliśmy słabą zasadę maksimum przy mierzalnych
współczynnikach i $c\ge0$, nie wykorzystując gładkości
ani symetrii części głównej. Dla laplasjanu udowodniliśmy
mocną zasadę maksimum, lemat Hopfa, ilościową nierówność
Harnacka, oszacowania gradientu i własności subharmoniczne.
Nierówność Caccioppolego pozostaje prawdziwa dla słabych
rozwiązań z mierzalną macierzą współczynników, ale sama
nie zapewnia od razu ciągłości rozwiązań ani nierówności
Harnacka. Te głębsze zagadnienia prowadzą do
teorii De Giorgiego--Nasha--Mosera oraz regularności
eliptycznej w kolejnych rozdziałach.

% Iteracja 15: regularność eliptyczna, całość rozdziału.
\chapter[Regularność Sobolewa rozwiązań eliptycznych]{Regularność Sobolewa rozwiązań eliptycznych: ilorazy różnicowe, oszacowania wewnętrzne i brzegowe}
\label{chap:pde15}

W rozdziale~\ref{chap:pde13} skonstruowaliśmy rozwiązania słabe
$u\in H^1$. Teraz zbadamy, kiedy równanie pozwala odzyskać drugie
pochodne $u$ w $L^2$. Zasadnicza trudność jest logiczna: nie wolno
różniczkować równania, zanim nie udowodnimy, że odpowiednie pochodne
istnieją. Ominiemy ją dzięki ilorazom różnicowym, których zastosowanie
wymaga jedynie $H^1$. W końcowej części zobaczymy, dlaczego gładkość
współczynników i brzegu naprawdę jest potrzebna.

W całym rozdziale $d\ge1$, a obszarem w $\R^d$ nazywamy niepusty,
otwarty i spójny podzbiór $\R^d$. Dla otwartych zbiorów
$U,V$ zapis $U\Subset V$ oznacza, że $\overline U$ jest zwartym
podzbiorem $V$. Symbole $H^k(U)=W^{k,2}(U)$ oraz
$\|\cdot\|_{2,U}=\|\cdot\|_{L^2(U)}$ mają znaczenie ustalone
w rozdziale~\ref{chap:pde-sobolev}. Całki po $U$ zapisujemy względem
miary Lebesgue'a $\mathrm d\lambda_d(x)$.

\section{Eliptyczność a możliwość różniczkowania rozwiązań}
\label{sec:pde15-model}
Rozważamy równanie w postaci dywergencyjnej
\begin{equation}\label{eq:pde15-basic}
 -\diver(A\nabla u)=f\qquad\text{w }\Omega,
\end{equation}
gdzie $\Omega\subset\R^d$ jest obszarem, $A=(a_{ij})$ ma współczynniki
rzeczywiste, a $f\in L^2_{\mathrm{loc}}(\Omega)$.

\begin{definicja}[Jednostajna eliptyczność i rozwiązanie słabe]
\label{def:pde15-elliptic}
Mówimy, że $A\in L^\infty_{\mathrm{loc}}(\Omega;\R^{d\times d})$
jest jednostajnie eliptyczna na $U\subset\Omega$, jeżeli istnieją
$\lambda,\Lambda>0$ takie, że dla prawie każdego $x\in U$ i
każdego $\xi\in\R^d$ zachodzi
$\xi^{\mathsf T}A(x)\xi\ge\lambda|\xi|^2$ oraz
$|A(x)\xi|\le\Lambda|\xi|$. Nie wymagamy symetrii $A$.
Funkcja $u\in H^1_{\mathrm{loc}}(\Omega)$ jest rozwiązaniem słabym
równania~\eqref{eq:pde15-basic}, jeżeli dla każdego
$\varphi\in C_c^\infty(\Omega)$ zachodzi
\[
 \int_\Omega A\nabla u\cdot\nabla\varphi\,\mathrm d\lambda_d(x)
 =\int_\Omega f\varphi\,\mathrm d\lambda_d(x).
\]
Z gęstości funkcji testowych równość rozciąga się na funkcje
$\varphi\in H^1(\Omega)$ o zwartym nośniku w $\Omega$.
\end{definicja}

\begin{uwaga}[Dlaczego samo $A\in L^\infty$ nie wystarcza?]
Jeśli $u$ ma tylko jedną słabą pochodną, to iloczyn $A\nabla u$
należy jedynie do $L^2_{\mathrm{loc}}$. Wyrażenie
$\diver(A\nabla u)$ jest jednak poprawnie określoną dystrybucją.
Przejście do $\nabla^2u\in L^2$ wymaga dodatkowej kontroli zmian
współczynników $A$. W poniższym twierdzeniu wystarczy warunek
$A\in W^{1,\infty}_{\mathrm{loc}}$.
\end{uwaga}

\section{Ilorazy różnicowe i ich podstawowe własności}
\label{sec:pde15-differences}
Dla $1\le j\le d$ i $h\ne0$ wprowadzamy
$\delta_h^jv(x)=\bigl(v(x+he_j)-v(x)\bigr)/h$, o ile oba punkty
należą do zbioru, na którym określona jest $v$.

\begin{lemat}[Oszacowania i całkowanie przez części]
\label{lem:pde15-difference}
Niech $U\subset\R^d$ będzie otwarty, $K\Subset U$ zwarty,
a $|h|<\operatorname{dist}(K,\partial U)$.
Jeżeli $v\in H^1(U)$, to
$\|\delta_h^jv\|_{L^2(K)}\le\|\partial_jv\|_{L^2(U)}$.
Jeżeli ponadto $v,w$ mają nośniki zwarte w $U$ i dla małych
$|h|$ wszystkie przesunięte nośniki pozostają w $U$, to
\[
 \int_U(\delta_h^jv)w\,\mathrm d\lambda_d
 =-\int_Uv(\delta_{-h}^jw)\,\mathrm d\lambda_d.
\]
Dla $v\in H^1_{\mathrm{loc}}(U)$ zachodzi też
$\delta_h^j(\partial_kv)=\partial_k(\delta_h^jv)$
w części wspólnej odpowiednich otwartych zbiorów.
\end{lemat}
\begin{proof}
Dla funkcji gładkiej podstawowe twierdzenie rachunku całkowego
wzdłuż odcinka daje
$\delta_h^jv(x)=\int_0^1\partial_jv(x+the_j)\,\mathrm dt$.
Nierówność Jensena i całkowanie po $K$ prowadzą do
\[
 \int_K|\delta_h^jv|^2\,\mathrm d\lambda_d
 \le\int_0^1\int_K|\partial_jv(x+the_j)|^2
 \mathrm d\lambda_d(x)\,\mathrm dt
 \le\int_U|\partial_jv|^2\,\mathrm d\lambda_d.
\]
Dla ogólnej funkcji Sobolewa aproksymujemy ją gładko
na otoczeniu $K$ powiększonym o odcinki przesunięć;
ciągłość translacji w $L^2$ pozwala przejść do granicy.
Całkowanie przez części dla ilorazów różnicowych wynika
wprost z podstawienia $y=x+he_j$ w pierwszym składniku.
Tożsamość z pochodną słabą wynika z przemienności
translacji z różniczkowaniem dystrybucyjnym.
\end{proof}

\begin{lemat}[Kryterium istnienia słabej pochodnej]
\label{lem:pde15-criterion}
Niech $U\subset\R^d$ będzie otwarty i $v\in L^2_{\mathrm{loc}}(U)$.
Jeśli dla każdego $K\Subset U$ istnieją $C_K,h_K>0$ takie, że
$\|\delta_h^jv\|_{L^2(K)}\le C_K$ dla $0<|h|<h_K$, to
$\partial_jv\in L^2_{\mathrm{loc}}(U)$ i
$\|\partial_jv\|_{L^2(K)}\le C_{K'}$ dla $K\Subset\operatorname{int}K'\Subset U$.
\end{lemat}
\begin{proof}
Wybierzmy otwarty $V$ z $K\Subset V\Subset U$ i zwartym
$K'\Subset U$ zawierającym $\overline V$. Z ograniczoności w
$L^2(V)$ wybieramy ciąg $h_n\to0$ taki, że
$\delta_{h_n}^jv\rightharpoonup w$ słabo w $L^2(V)$.
Dla $\varphi\in C_c^\infty(V)$ mamy na mocy poprzedniego lematu
\[
 \int_V\delta_{h_n}^jv\,\varphi\,\mathrm d\lambda_d
 =-\int_Uv\,\delta_{-h_n}^j\varphi\,\mathrm d\lambda_d
 \longrightarrow-\int_Vv\,\partial_j\varphi\,\mathrm d\lambda_d.
\]
Lewa strona zbiega do $\int_Vw\varphi\,\mathrm d\lambda_d$.
Stąd $w=\partial_jv$ dystrybucyjnie w $V$.
Półciągłość normy przy słabej zbieżności daje żądane
oszacowanie. Jednoznaczność pochodnej dystrybucyjnej
pozwala skleić reprezentanty otrzymane na różnych $V$.
\end{proof}

\section[Wewnętrzna regularność Sobolewa drugiego rzędu]{Wewnętrzna regularność $H^2$}
\label{sec:pde15-h2}
\begin{twierdzenie}[Wewnętrzna regularność drugiego rzędu]
\label{thm:pde15-h2}
Niech $\Omega\subset\R^d$ będzie obszarem,
$A\in W^{1,\infty}_{\mathrm{loc}}(\Omega;\R^{d\times d})$ będzie
jednostajnie eliptyczna i ograniczona na zwartych podzbiorach,
$f\in L^2_{\mathrm{loc}}(\Omega)$, a $u\in H^1_{\mathrm{loc}}(\Omega)$
rozwiązaniem słabym równania~\eqref{eq:pde15-basic}.
Wówczas $u\in H^2_{\mathrm{loc}}(\Omega)$. Dla każdych
obszarów $U,V$ takich, że $U\Subset V\Subset\Omega$, zachodzi
\begin{equation}\label{eq:pde15-h2-est}
 \|u\|_{H^2(U)}\le C\bigl(\|f\|_{L^2(V)}
 +\|u\|_{H^1(V)}\bigr),
\end{equation}
gdzie $C$ zależy od $U,V,d$, stałych eliptyczności na $V$
oraz $\|A\|_{W^{1,\infty}(V)}$.
\end{twierdzenie}
\begin{proof}
Wybierzmy $\eta\in C_c^\infty(V)$, równą $1$ w otoczeniu
$\overline U$, i ustalmy kierunek $j$.
Dla małych $|h|$ nośniki $\eta$ i $\eta(\cdot-he_j)$ są
zawarte w $V$. Oznaczmy $w_h=\delta_h^ju$.
Odejmując słabą tożsamość dla $u(x+he_j)$ od odpowiedniej
tożsamości dla $u(x)$, uzyskujemy dla $\varphi$ o nośniku
wewnątrz $V$ równanie
\begin{equation}\label{eq:pde15-diff-eq}
 \int_V A(x)\nabla w_h\cdot\nabla\varphi\,\mathrm d\lambda_d
 =\int_V(\delta_h^jf)\varphi\,\mathrm d\lambda_d
 -\int_V(\delta_h^jA)(x)\nabla u(x+he_j)
 \cdot\nabla\varphi\,\mathrm d\lambda_d.
\end{equation}
Współczynniki $A\in W^{1,\infty}(V)$ mają lokalnie
Lipschitzowskie reprezentanty, więc
$\|\delta_h^jA\|_\infty\le\|\partial_jA\|_\infty$.
Używamy funkcji testowej $\varphi=\eta^2w_h$, która należy do
$H_0^1(V)$; jest dopuszczalna na mocy gęstości.
Eliptyczność i nierówność Younga dają
\begin{align*}
 \lambda\|\eta\nabla w_h\|_2^2
 &\le\left|\int_VA\nabla w_h\cdot(2\eta w_h\nabla\eta)
 \,\mathrm d\lambda_d\right|\\
 &\quad+\left|\int_V\delta_h^jf\,\eta^2w_h
 \,\mathrm d\lambda_d\right|
 +\|\delta_h^jA\|_\infty\|\nabla u(\cdot+he_j)\|_{2,V'}
 \|\nabla(\eta^2w_h)\|_2,
\end{align*}
gdzie $V'\Subset V$ wybieramy tak, by zawierał wszystkie
przesunięcia nośnika $\eta$.
Kluczowe jest to, że \emph{nie} szacujemy $\delta_h^jf$
przez $\|\delta_h^jf\|_2$, bo taki szacunek nie jest
jednostajny. Zamiast tego stosujemy dyskretne całkowanie
przez części i lemat~\ref{lem:pde15-difference}:
\[
 \left|\int_V\delta_h^jf\,\eta^2w_h\,\mathrm d\lambda_d\right|
 =\left|\int_V f\,\delta_{-h}^j(\eta^2w_h)
 \,\mathrm d\lambda_d\right|
 \le\|f\|_{2,V}\|\nabla(\eta^2w_h)\|_{2,V}.
\]
Ponadto
$\|\nabla(\eta^2w_h)\|_2
 \le C_\eta(\|\eta\nabla w_h\|_2+\|w_h\|_{2,\operatorname{supp}\eta})$.
Pierwszy lemat zapewnia
$\|w_h\|_{2,\operatorname{supp}\eta}\le C\|\nabla u\|_{2,V}$,
a analogicznie $\|\nabla u(\cdot+he_j)\|_{2,V'}$ można
oszacować normą na ustalonym większym zbiorze zawartym w $\Omega$.
Po użyciu nierówności Younga i pochłonięciu składnika z
$\|\eta\nabla w_h\|_2$ dostajemy
\[
 \|\eta\nabla w_h\|_2^2\le
 C\bigl(\|f\|_{2,V}^2+\|u\|_{H^1(V)}^2\bigr),
\]
z $C$ niezależną od $h$ (dobieramy $\eta$ i wszystkie
przesunięcia w obrębie $V$).
Na obszarze $U$ zachodzi więc jednostajne oszacowanie
$\delta_h^j\partial_ku=\partial_kw_h$ w $L^2(U)$.
Lemat~\ref{lem:pde15-criterion} daje
$\partial_j\partial_ku\in L^2(U)$ dla wszystkich $j,k$.
Przechodząc do granicy słabo w $L^2$ uzyskujemy
również~\eqref{eq:pde15-h2-est}. Ponieważ $U\Subset\Omega$
był dowolny, dowód jest zakończony.
\end{proof}

\begin{uwaga}[Zależność od $V$ w oszacowaniu]
Przy wyborze funkcji odcinającej wystarczy
$U\Subset V$; używane w dowodzie przesunięcia zachodzą
w skończenie wielu ustalonych podzbiorach $V$.
Stała w oszacowaniu~\eqref{eq:pde15-h2-est} rośnie przy
zbliżaniu $U$ do brzegu $V$, co odpowiada naturze lokalnego
oszacowania eliptycznego.
\end{uwaga}

\section{Oszacowanie Caccioppolego z prawą stroną}
\label{sec:pde15-caccioppoli}
\begin{stwierdzenie}[Wersja niejednorodna]
Niech $B_{2r}(x_0)\Subset\Omega$, $A\in L^\infty(B_{2r})$
spełnia warunki eliptyczności z definicji~\ref{def:pde15-elliptic},
$f\in L^2(B_{2r})$, a $u\in H^1(B_{2r})$ jest słabym
rozwiązaniem $-\diver(A\nabla u)=f$. Istnieje $C=C(d,\lambda,\Lambda)>0$
taka, że
\[
 \int_{B_r}|\nabla u|^2\,\mathrm d\lambda_d
 \le C\left(r^{-2}\int_{B_{2r}}|u|^2\,\mathrm d\lambda_d
 +r^2\int_{B_{2r}}|f|^2\,\mathrm d\lambda_d\right).
\]
\end{stwierdzenie}
\begin{proof}
Wybieramy $\eta\in C_c^\infty(B_{2r})$, równą $1$ na $B_r$
i spełniającą $|\nabla\eta|\le C_0/r$. W równaniu słabym
podstawiamy $\varphi=\eta^2u$:
\[
 \int_{B_{2r}}\eta^2A\nabla u\cdot\nabla u\,\mathrm d\lambda_d(x)
 =\int_{B_{2r}} f\eta^2u\,\mathrm d\lambda_d(x)
 -2\int_{B_{2r}}\eta u A\nabla u\cdot\nabla\eta\,\mathrm d\lambda_d(x).
\]
Pierwszy wyraz prawej strony szacujemy przez
$\|f\|_2\|\eta u\|_2$.
Ponieważ $\eta u\in H_0^1(B_{2r})$, nierówność Poincarégo
ze skalowaniem daje
$\|\eta u\|_2\le Cr\|\nabla(\eta u)\|_2$.
Następnie
$\|\nabla(\eta u)\|_2\le\|\eta\nabla u\|_2+
C_0r^{-1}\|u\|_2$.
Oba składniki prawej strony pochłaniamy nierównością Younga
z parametrem zależnym od $\lambda$; otrzymujemy
$\|\eta\nabla u\|_2^2\le C(r^{-2}\|u\|_2^2+r^2\|f\|_2^2)$.
Wybór $\eta\equiv1$ na $B_r$ kończy dowód.
\end{proof}

\section{Regularność wyższych rzędów i bootstrap}
\label{sec:pde15-bootstrap}
\begin{twierdzenie}[Wewnętrzny bootstrap Sobolewa]
\label{thm:pde15-bootstrap}
Niech $m\in\N_0$, $\Omega\subset\R^d$ będzie obszarem,
$A\in W^{m+1,\infty}_{\mathrm{loc}}(\Omega;\R^{d\times d})$
będzie jednostajnie eliptyczna, $f\in H^m_{\mathrm{loc}}(\Omega)$,
a $u\in H^1_{\mathrm{loc}}(\Omega)$ niech spełnia
$-\diver(A\nabla u)=f$ słabo. Wtedy
$u\in H^{m+2}_{\mathrm{loc}}(\Omega)$.
Dla $U\Subset V\Subset\Omega$ istnieje $C$ takie, że
$\|u\|_{H^{m+2}(U)}\le C(\|f\|_{H^m(V)}+
\|u\|_{H^1(V)})$; stała zależy ponadto od $m$, eliptyczności,
$U,V$ i norm $A$ w $W^{m+1,\infty}(V)$.
\end{twierdzenie}
\begin{proof}
Dla $m=0$ jest to twierdzenie~\ref{thm:pde15-h2}.
Załóżmy indukcyjnie, że przy odpowiednich założeniach
z $f\in H^{m-1}_{\mathrm{loc}}$ uzyskujemy
$u\in H^{m+1}_{\mathrm{loc}}$.
Dla każdego wielowskaźnika $\alpha$, $|\alpha|=m$,
funkcja $v_\alpha=\partial^\alpha u$ należy do
$H^1_{\mathrm{loc}}$, a po zróżniczkowaniu równania
w sensie dystrybucyjnym spełnia
\[
 -\diver(A\nabla v_\alpha)=\partial^\alpha f
 +\diver\!\left(\sum_{0<\beta\le\alpha}
 {\alpha\choose\beta}(\partial^\beta A)
 \nabla\partial^{\alpha-\beta}u\right).
\]
Ostatnie wyrażenie jest \emph{dywergencją pola $H^1_{\mathrm{loc}}$},
bo z założenia indukcyjnego dla $|\beta|\ge1$ mamy
$\nabla\partial^{\alpha-\beta}u\in H^{|\beta|}_{\mathrm{loc}}$
i $\partial^\beta A\in W^{1,\infty}_{\mathrm{loc}}$
dla $|\beta|\le m$. Zatem prawa strona należy do
$L^2_{\mathrm{loc}}$, również składnik $\partial^\alpha f$.
Możemy zastosować twierdzenie~\ref{thm:pde15-h2} do
$v_\alpha$ i otrzymać $v_\alpha\in H^2_{\mathrm{loc}}$.
Oznacza to $u\in H^{m+2}_{\mathrm{loc}}$.
Przy szacowaniu norm wybieramy ciąg zagnieżdżonych
zbiorów otwartych między $U$ i $V$ i stosujemy kolejne
oszacowania drugiego rzędu; liczba kroków jest skończona.
\end{proof}

\begin{wniosek}[Gładkość wewnętrzna]
Jeżeli $A\in C^\infty(\Omega;\R^{d\times d})$ jest lokalnie
jednostajnie eliptyczna, $f\in C^\infty(\Omega)$ i
$u\in H^1_{\mathrm{loc}}(\Omega)$ jest słabym rozwiązaniem,
to $u\in C^\infty(\Omega)$.
\end{wniosek}
\begin{proof}
Każda pochodna współczynników jest ograniczona na zbiorze
$V\Subset\Omega$. Poprzednie twierdzenie stosujemy dla
każdego $m$ i otrzymujemy $u\in H^k_{\mathrm{loc}}(\Omega)$
dla wszystkich $k$. Lokalne zanurzenia Sobolewa dają
$u\in C^\ell_{\mathrm{loc}}(\Omega)$ dla dowolnego $\ell$.
\end{proof}

\section{Regularność przy brzegu: metoda spłaszczania}
\label{sec:pde15-boundary}

\begin{lemat}[Forma dywergencyjna po spłaszczeniu brzegu]
\label{lem:pde15-flatten}
Niech $V,W\subset\R^d$ będą zbiorami otwartymi i
$\Phi:W\to V$ dyfeomorfizmem klasy $C^{1,1}$,
którego macierz Jacobiego $J(y)=D\Phi(y)$ oraz macierz $J(y)^{-1}$
są jednostajnie ograniczone. Przyjmijmy, że
$u\in H^1_{\mathrm{loc}}(V)$ jest słabym rozwiązaniem
$-\operatorname{div}(A\nabla u)=f$ na $V$, gdzie
$A\in W^{1,\infty}_{\mathrm{loc}}(V)$ i $f\in L^2_{\mathrm{loc}}(V)$.
Wtedy $v=u\circ\Phi$ spełnia na $W$ równanie słabe
$-\operatorname{div}(B\nabla v)=F$, gdzie
\[
 B(y)=|\det J(y)|\,J(y)^{-1}A(\Phi(y))J(y)^{-\mathsf T},
 \qquad F(y)=|\det J(y)|f(\Phi(y)).
\]
Jeżeli $A$ jest jednostajnie eliptyczna, a $J,J^{-1}$
mają jednostajne oszacowania na rozważanym zwartym
otoczeniu, to $B$ jest jednostajnie eliptyczna i
$B\in W^{1,\infty}_{\mathrm{loc}}(W)$.
\end{lemat}
\begin{proof}
Dla $\psi\in C_c^\infty(W)$ w równaniu dla $u$ wybieramy
$\varphi=\psi\circ\Phi^{-1}$ (po standardowej aproksymacji
funkcji klasy $C^{1,1}$ przez testy gładkie w $H^1$).
Z reguły łańcuchowej $\nabla_xu(\Phi(y))=J(y)^{-\mathsf T}\nabla_yv(y)$
i analogicznie
$\nabla_x\varphi(\Phi(y))=J(y)^{-\mathsf T}\nabla_y\psi(y)$.
Stosując wzór zmiany zmiennych w obu całkach dostajemy
\[
 \int_W B\nabla v\cdot\nabla\psi\,\mathrm d\lambda_d(y)
 =\int_W F\psi\,\mathrm d\lambda_d(y).
\]
Dla $\xi\in\R^d$ mamy
$\xi^{\mathsf T}B\xi=|\det J|
(J^{-\mathsf T}\xi)^{\mathsf T}A(J^{-\mathsf T}\xi)
\ge |\det J|\lambda|J^{-\mathsf T}\xi|^2$.
Na zwartym otoczeniu wyznacznik jest oddzielony od zera,
a najmniejsza wartość osobliwa $J^{-\mathsf T}$ jest
dodatnia; to daje nowe stałe eliptyczności.
Reguła łańcuchowa dla funkcji Lipschitzowskich i
$J\in W^{1,\infty}_{\mathrm{loc}}$ daje
$B\in W^{1,\infty}_{\mathrm{loc}}$.
\end{proof}

\begin{lemat}[Oszacowanie pochodnych stycznych w półkuli]
\label{lem:pde15-tangential}
Niech $d\ge2$, $B_2^+=\{x\in\R^d:|x|<2,\ x_d>0\}$,
$A\in W^{1,\infty}(B_2^+;\R^{d\times d})$ będzie
jednostajnie eliptyczna, $f\in L^2(B_2^+)$ oraz
$v\in H^1(B_2^+)$ niech spełnia
\[
 \int_{B_2^+}A\nabla v\cdot\nabla\varphi\,\mathrm d\lambda_d(x)
 =\int_{B_2^+}f\varphi\,\mathrm d\lambda_d(x)
 \quad(\varphi\in H_0^1(B_2^+)).
\]
Jeżeli ślad $v$ na płaskiej części brzegu $\{x_d=0,|x|<2\}$
wynosi zero, to dla każdego $j<d$
$\partial_j\nabla v\in L^2(B_1^+)$ i
\[
 \|\partial_j\nabla v\|_{L^2(B_1^+)}
 \le C\bigl(\|f\|_{L^2(B_2^+)}+\|v\|_{H^1(B_2^+)}\bigr),
\]
gdzie $C$ zależy od $d$, eliptyczności oraz
$\|A\|_{W^{1,\infty}(B_2^+)}$.
\end{lemat}
\begin{proof}
Ustalmy $j<d$. Wybieramy $\eta\in C_c^\infty(B_{3/2})$
równą $1$ na $B_1$, i dla małych $|h|$ kładziemy
$w_h=\delta_h^jv$. Ponieważ przesunięcia są styczne do
płaskiego brzegu, $w_h$ ma tam zerowy ślad.
W szczególności $\eta^2w_h\in H_0^1(B_2^+)$.
Ponadto $-\delta_{-h}^j(\eta^2w_h)$ również należy
do $H_0^1(B_2^+)$: jego nośnik omija część sferyczną,
a ślad na płaskiej jest zerowy.
Testujemy słabe równanie ostatnią funkcją i wykonujemy
dyskretne całkowanie przez części. Otrzymujemy
\[
 \int_{B_2^+}A\nabla w_h\cdot\nabla(\eta^2w_h)
 \,\mathrm d\lambda_d(x)
 =\int_{B_2^+}(\delta_h^jf)\eta^2w_h\,\mathrm d\lambda_d(x)
 -\int_{B_2^+}(\delta_h^jA)\nabla v(\cdot+he_j)
 \cdot\nabla(\eta^2w_h)\,\mathrm d\lambda_d(x).
\]
W pierwszym członie po prawej nie różniczkujemy $f$:
\[
 \left|\int_{B_2^+}(\delta_h^jf)\eta^2w_h
 \,\mathrm d\lambda_d\right|
 =\left|\int_{B_2^+}f\,\delta_{-h}^j(\eta^2w_h)
 \,\mathrm d\lambda_d\right|
 \le\|f\|_2\|\partial_j(\eta^2w_h)\|_2.
\]
Własności ilorazów różnicowych dają
$\|w_h\|_{L^2(\operatorname{supp}\eta)}\le\|\nabla v\|_2$
i $\|\delta_h^jA\|_\infty\le\|\partial_jA\|_\infty$.
Rozwijamy $\nabla(\eta^2w_h)=\eta^2\nabla w_h+
2\eta w_h\nabla\eta$, stosujemy eliptyczność do
pierwszego członu i nierówność Younga do pozostałych.
Po pochłonięciu otrzymujemy jednostajnie w $h$
\[
 \|\eta\nabla w_h\|_2^2
 \le C(\|f\|_2^2+\|v\|_{H^1(B_2^+)}^2).
\]
Kryterium z lematu~\ref{lem:pde15-criterion}, stosowane
lokalnie w półkuli i z oszacowaniem jednostajnym aż do
płaskiego brzegu, daje $\partial_j\partial_kv\in L^2(B_1^+)$
dla każdego $k$. Nie rozciągamy funkcji przez zero
w kierunku normalnym, ponieważ takie przesunięcie
nie zachowuje równania słabego.
\end{proof}



\begin{twierdzenie}[Globalne oszacowanie $H^2$ dla Dirichleta]
\label{thm:pde15-boundary}
Niech $\Omega\subset\R^d$ będzie ograniczonym obszarem
klasy $C^{1,1}$, $A\in W^{1,\infty}(\Omega;\R^{d\times d})$
będzie jednostajnie eliptyczna i ograniczona, a
$f\in L^2(\Omega)$. Jeżeli $u\in H_0^1(\Omega)$ spełnia
$-\diver(A\nabla u)=f$ słabo, to $u\in H^2(\Omega)$ oraz
\begin{equation}\label{eq:pde15-boundary}
 \|u\|_{H^2(\Omega)}\le C\bigl(\|f\|_{L^2(\Omega)}
 +\|u\|_{H^1(\Omega)}\bigr),
\end{equation}
gdzie $C$ zależy od $d,\Omega$, stałych eliptyczności
i $\|A\|_{W^{1,\infty}(\Omega)}$.
\end{twierdzenie}
\begin{proof}
Podzielimy dowód na kroki, rozdzielając dwa różne zagadnienia:
kontrolę pochodnych stycznych i odzyskanie pochodnej normalnej.

\emph{Krok 1: skończone pokrycie brzegu.}
Ponieważ $\partial\Omega$ jest zwarty i klasy $C^{1,1}$,
można pokryć go skończenie wieloma otoczeniami, w których
po obrocie i translacji $\Omega$ ma postać
$x_d>\gamma(x')$, gdzie $\gamma\in C^{1,1}$.
Odwzorowanie $y' =x'$, $y_d=x_d-\gamma(x')$
i jego odwrotność są klasy $C^{1,1}$, a ich pierwsze
pochodne są Lipschitzowskie.
Dokładny wzór na nową macierz $B$ i prawą stronę $F$
podaje lemat~\ref{lem:pde15-flatten}; w tej zmianie zmiennych
$\det D\Phi=1$. Otrzymujemy jednostajnie eliptyczne
współczynniki $B\in W^{1,\infty}$ na mniejszym otoczeniu.
Na płaskim brzegu $\{y_d=0\}$ funkcja $v$ ma ślad zerowy.

\emph{Krok 2: pochodne styczne.}
Ustalmy funkcję odcinającą $\eta$ w półkuli $B_R^+$.
Dla $j<d$ przesunięcia $y\mapsto y+he_j$ nie przecinają
płaskiego brzegu. Iloraz różnicowy $w_h=\delta_h^jv$
ma zatem zerowy ślad na tym brzegu. Dopuszczalny jest test
$\varphi=-\delta_{-h}^j(\eta^2w_h)$ w słabym równaniu;
w przypadku $\eta^2w_h$ o zwartym nośniku wewnątrz
półkuli z wyjątkiem płaskiej części brzegu test należy do
$H_0^1(B_R^+)$.
Stosujemy lemat~\ref{lem:pde15-tangential} po przeskalowaniu
półkuli i odcięciu równania w mniejszym otoczeniu.
Otrzymujemy $\partial_j\nabla v\in L^2$ lokalnie do
płaskiego brzegu dla $j<d$ i kontrolę tych pochodnych
przez normy $\|f\|_2+\|v\|_{H^1}$.
Lemat uzasadnia dokładnie, dlaczego wolno przenieść
iloraz różnicowy na test, bez zakładania pochodnych $f$.

\emph{Krok 3: druga pochodna normalna.}
Dla $d=1$ nie ma kierunków stycznych: z równania
$-(a v')'=F$ mamy $(a v')'=-F\in L^2$, a ponieważ
$a\in W^{1,\infty}$ i $a\ge\lambda>0$, wynika
$v''\in L^2$. Dalej rozważamy $d\ge2$.
Rozpiszmy równanie (w zmiennych $y$) jako
\[
 -a_{dd}\partial_{dd}v=f+
 \sum_{(i,j)\ne(d,d)}\partial_i(a_{ij}\partial_jv)
 +(\partial_da_{dd})\partial_dv.
\]
Wszystkie składniki prawej strony oprócz
$\partial_{dd}v$ są już w $L^2$: gdy $i<d$,
korzystamy z pochodnych stycznych $\partial_i\nabla v$,
a gdy $i=d$ i $j<d$, korzystamy z przemienności
słabych pochodnych mieszanych.
Ponadto $a_{dd}\ge\lambda>0$ prawie wszędzie.
Stąd $\partial_{dd}v\in L^2$, a więc $v\in H^2$
lokalnie aż do spłaszczonego brzegu.

\emph{Krok 4: sklejenie oszacowań.}
Wnętrze $\Omega$ pokrywamy skończenie wieloma mniejszymi
zbiorami, a brzeg powyższymi mapami. Wybieramy partycję jedności
i sumujemy lokalne oszacowania. Ponieważ zmiany zmiennych
są klasy $C^{1,1}$, przeniesienie drugich słabych pochodnych
z powrotem do współrzędnych $x$ zachowuje $H^2$.
Daje to~\eqref{eq:pde15-boundary}.
\end{proof}

\begin{uwaga}[Niejednorodne dane Dirichleta]
Jeśli $g$ jest śladem pewnej funkcji $G\in H^2(\Omega)$,
to dla $v=u-G\in H_0^1(\Omega)$ otrzymujemy
$-\diver(A\nabla v)=f+\diver(A\nabla G)\in L^2(\Omega)$.
Stosując twierdzenie~\ref{thm:pde15-boundary} do $v$,
otrzymujemy $u\in H^2(\Omega)$. Zamiast wymagania
$G\in H^2$ można sformułować warunek na $g$ w odpowiedniej
przestrzeni śladów wyższego rzędu; wymaga to osobnego
twierdzenia o operatorze śladu.
\end{uwaga}


\begin{uwaga}[Granica obowiązywania identyfikacji dziedziny]
Dla $A=I$ twierdzenie~\ref{thm:pde15-boundary}
wyjaśnia identyfikację $D(-\Delta_D)=H^2(\Omega)\cap
H_0^1(\Omega)$ na obszarze klasy $C^{1,1}$, zapowiedzianą
w uwadze~\ref{rem:pde12-corners}. Na ogólnym obszarze
Lipschitzowskim taka identyfikacja może być fałszywa;
nie wolno jej używać w argumentach spektralnych
bez odpowiednich dodatkowych założeń.
\end{uwaga}

\section{Dwa kontrprzykłady: współczynniki i geometria brzegu}
\label{sec:pde15-counterexamples}
\begin{kontrprzyklad}[Skok współczynnika: $H^1$ bez $H^2$]
\label{ex:pde15-jump}
Na $\Omega=(-1,1)$ połóżmy
\[
 a(x)=\begin{cases}1,&x<0,\\2,&x\ge0,\end{cases}
 \qquad
 u(x)=\begin{cases}x,&x\le0,\\x/2,&x>0.\end{cases}
\]
Funkcja $u$ jest Lipschitzowska, więc należy do $H^1(\Omega)$.
Mamy $a(x)u'(x)=1$ prawie wszędzie, zatem
$-(a u')'=0$ dystrybucyjnie.
Jednak $u'$ ma skok równy $-1/2$ w zerze, a więc
$u''=-\tfrac12\delta_0$ jako dystrybucja; stąd
$u\notin H^2(\Omega)$. Samo ograniczenie i eliptyczność
$A$ nie gwarantują drugich pochodnych w $L^2$.
\end{kontrprzyklad}

\begin{kontrprzyklad}[Narożnik wklęsły]
\label{ex:pde15-corner}
Niech $\Omega$ będzie ograniczonym obszarem płaskim,
który w sąsiedztwie zera jest sektorem o kącie rozwarcia
$3\pi/2$, np. obszarem w kształcie litery L.
Współrzędne biegunowe dobieramy tak, że w tym sektorze
$0<\theta<3\pi/2$. Funkcja
$w(r,\theta)=r^{2/3}\sin(2\theta/3)$ jest harmoniczna dla $r>0$
i znika na ramionach sektora. Wybieramy gładką radialną
funkcję $\chi$, równą $1$ w otoczeniu zera i zerową poza
małą kulą zawartą w sąsiedztwie narożnika. Wtedy
$u=\chi w\in H_0^1(\Omega)$ oraz
$f=-\Delta u\in L^2(\Omega)$: blisko zera mamy $f=0$,
a tam, gdzie $\nabla\chi\ne0$, wszystkie składniki są gładkie.
Tymczasem $|\nabla^2u|\asymp r^{-4/3}$ blisko zera na
podsektorze dodatniej miary kątowej i dlatego
\[
 \int_{\Omega\cap B_\eps(0)}|\nabla^2u|^2\,\mathrm d\lambda_2
 \ge c\int_0^\eps r^{-8/3}r\,\mathrm dr=+\infty.
\]
Zatem $u\notin H^2(\Omega)$ mimo gładkiego operatora
$-\Delta$ i prawej strony $f\in L^2$. Obszar Lipschitzowski
nie musi mieć globalnej regularności $H^2$.
\end{kontrprzyklad}

\section{Wnioski i zakres dalszych twierdzeń}
\label{sec:pde15-summary}
Regularność $H^2$ we wnętrzu wynika z eliptyczności i
Lipschitzowskich współczynników; nie trzeba zakładać
$C^2$-regularności rozwiązania. Regularność przy brzegu
wymaga natomiast dodatkowej geometrii brzegu i poprawnie
postawionych warunków brzegowych.
W kolejnym rozdziale przejdziemy od oszacowań $L^2$ do
regularności $W^{2,p}$ oraz oszacowań Schaudera.
Nie należy przy tym mylić regularności Sobolewa z analitycznością:
$C^\infty$-gładkość współczynników gwarantuje gładkość rozwiązania
w powyższej sytuacji, lecz sama nie gwarantuje analityczności.

\chapter[Oszacowania Calderóna--Zygmunda i Schaudera]{Regularność $W^{2,p}$, całki osobliwe Calderóna--Zygmunda i oszacowania Schaudera}
\label{chap:pde16}

Rozdział~\ref{chap:pde15} dostarczył oszacowań drugich pochodnych w $L^2$.
Nie jest to koniec regularności. Z równania $-\Delta u=f$ możemy
często wywnioskować $D^2u\in L^p$ dla każdego $1<p<\infty$, a przy
założeniu Hölderowskiej regularności $f$ nawet $D^2u\in C^{0,\alpha}$.
Przejścia te nie wynikają z różniczkowania funkcji $f$: mechanizm
opiera się na strukturze całek osobliwych i na wzajemnym znoszeniu
się ich dodatnich i ujemnych części. Właśnie dlatego końcowe
przypadki $p=1$, $p=\infty$ oraz $\alpha=1$ są subtelne.

Przez $\Omega$ rozumiemy obszar w $\R^d$; $U\Subset V$ oznacza
$\overline U\subset V$ i zwartość $\overline U$. W tym rozdziale
$d\ge2$, chyba że jawnie powiemy inaczej. Wszystkie całki po
$\R^d$ i jego podzbiorach zapisujemy względem $\mathrm d\lambda_d$.
Przyjmujemy konwencję Fouriera z~\eqref{eq:pde8-convention}.

\section[Dlaczego pojawiają się całki osobliwe?]{Dlaczego $D^2(-\Delta)^{-1}$ nie jest operatorem splotu z funkcją całkowalną?}
Niech $\omega_{d-1}=\mathcal H^{d-1}(S^{d-1})$ oznacza pole
jednostkowej sfery, a $\Gamma$ rozwiązanie fundamentalne $-\Delta$:
$-\Delta\Gamma=\delta_0$. Dla $d\ge3$ można wziąć
$\Gamma(x)=\bigl((d-2)\omega_{d-1}\bigr)^{-1}|x|^{2-d}$,
a dla $d=2$ funkcję $\Gamma(x)=-(2\pi)^{-1}\log|x|$.
Dla $x\ne0$ mamy
\begin{equation}\label{eq:pde16-kernel}
 K_{ij}(x):=\partial_{ij}\Gamma(x)
 =\frac{d x_ix_j-\delta_{ij}|x|^2}
 {\omega_{d-1}|x|^{d+2}},\qquad 1\le i,j\le d.
\end{equation}
Jądro jest jednorodne stopnia $-d$, spełnia
$|K_{ij}(x)|\le C|x|^{-d}$ i
$|\nabla K_{ij}(x)|\le C|x|^{-d-1}$. Ponadto
$\int_{S^{d-1}}K_{ij}(\theta)\,\mathrm d\mathcal H^{d-1}(\theta)=0$,
ponieważ średnia $\theta_i\theta_j$ na sferze wynosi
$\delta_{ij}/d$.

W pobliżu zera funkcja $|x|^{-d}$ nie jest całkowalna. Wobec tego
naiwny wzór $D_{ij}u=K_{ij}*f$ jest błędny nawet dla $f\in C_c^\infty$.

\begin{definicja}[Główna wartość całki]
Niech $f\in C_c^1(\R^d)$, a $K$ będzie jądrem klasy $C^1$ poza
zerem, jednorodnym stopnia $-d$ i o zerowej średniej sferycznej.
Definiujemy
\[
 T_Kf(x)=\operatorname{p.v.}\!\int_{\R^d}K(x-y)f(y)\,\mathrm d\lambda_d(y)
 :=\lim_{\varepsilon\downarrow0}
 \int_{|x-y|>\varepsilon}K(x-y)f(y)\,\mathrm d\lambda_d(y).
\]
Dla $f\in C_c^1$ granica istnieje: w kuli $B_1(x)$ odejmujemy
$f(x)$, a z warunku zerowej średniej sferycznej wynika, że
całka ze stałej po każdym pierścieniu jest równa zeru.
\end{definicja}

\begin{lemat}[Hessian rozwiązania fundamentalnego jako dystrybucja]
\label{lem:pde16-distr-kernel}
Dla każdego $1\le i,j\le d$ zachodzi w $\mathcal D'(\R^d)$
\begin{equation}\label{eq:pde16-distr-kernel}
 \partial_{ij}\Gamma
 =\operatorname{p.v.}K_{ij}-\frac{\delta_{ij}}d\delta_0.
\end{equation}
W szczególności dla $f\in C_c^\infty(\R^d)$ i $u=\Gamma*f$,
\begin{equation}\label{eq:pde16-repr}
 \partial_{ij}u=T_{ij}f-\frac{\delta_{ij}}d f,
 \qquad T_{ij}=T_{K_{ij}}.
\end{equation}
\end{lemat}
\begin{proof}
Testujemy dystrybucję $\partial_{ij}\Gamma$ na
$\varphi\in C_c^\infty$ i całkujemy dwukrotnie przez części na
$\{|x|>\varepsilon\}$. Po pierwszym całkowaniu otrzymujemy
wyraz z $\partial_j\Gamma$, po drugim całkę
$\int_{|x|>\varepsilon}K_{ij}\varphi$ i wyrazy na sferze
$|x|=\varepsilon$. Tam
$\partial_i\Gamma(x)=-x_i/(\omega_{d-1}|x|^d)$, zatem
jedyny wyraz brzegowy, który nie zanika w granicy, ma postać
$-\omega_{d-1}^{-1}\int_{S^{d-1}}\theta_i\theta_j
\varphi(\varepsilon\theta)\,\mathrm d\mathcal H^{d-1}(\theta)$.
Granica wynosi $-\delta_{ij}\varphi(0)/d$.
Pozostały wyraz brzegowy jest rzędu
$\varepsilon|\log\varepsilon|$ w wymiarze $2$, a w wyższych
wymiarach rzędu $\varepsilon$, ponieważ $\varphi$ i jej
pochodne są ograniczone. Zerowa średnia $K_{ij}$ i oszacowanie
$|\varphi(x)-\varphi(0)|\le C|x|$ uzasadniają przejście do
głównej wartości. To daje~\eqref{eq:pde16-distr-kernel}.
Splot z funkcją testową i przemienność splotu z pochodnymi
dają~\eqref{eq:pde16-repr}.
\end{proof}

\section[Oszacowanie Hilbertowskie i rozkład Calderóna--Zygmunda]{Oszacowanie $L^2$ i rozkład Calderóna--Zygmunda}
\begin{lemat}[Oszacowanie w $L^2$]
\label{lem:pde16-l2}
Dla $f\in C_c^\infty(\R^d)$ zachodzi
$\|T_{ij}f\|_2\le C_d\|f\|_2$. Operator $T_{ij}$
przedłuża się jednoznacznie do ograniczonego operatora na
$L^2(\R^d)$.
\end{lemat}
\begin{proof}
Z~\eqref{eq:pde16-repr}, $-\Delta u=f$ i transformaty Fouriera
wynika, dla $\xi\ne0$,
$\widehat{\partial_{ij}u}(\xi)=-(\xi_i\xi_j/|\xi|^2)\widehat f(\xi)$.
Zatem $T_{ij}$ ma mnożnik
$m_{ij}(\xi)=\delta_{ij}/d-\xi_i\xi_j/|\xi|^2$.
Jego moduł jest ograniczony przez $2$; twierdzenie Plancherela
daje oszacowanie. Użycie dystrybucji temperowanych pozwala
wykonać rachunek także przy $d=2$, gdy sama $\Gamma$ rośnie
logarytmicznie. Gęstość $C_c^\infty$ w $L^2$ kończy dowód.
\end{proof}

\begin{lemat}[Rozkład Calderóna--Zygmunda]
\label{lem:pde16-cz-decomp}
Niech $f\in L^1(\R^d)\cap L^2(\R^d)$ i $s>0$. Istnieje
co najwyżej przeliczalna rodzina parami rozłącznych kostek
dyadycznych $Q_k$ oraz rozkład $f=g+\sum_k b_k$ taki, że
\begin{enumerate}[label=(\roman*)]
\item $|g|\le 2^ds$ prawie wszędzie i $\|g\|_1\le\|f\|_1$;
\item $\operatorname{supp}b_k\subset Q_k$, $\int_{Q_k} b_k\,\mathrm d\lambda_d=0$,
 $\sum_k\|b_k\|_1\le2\|f\|_1$;
\item $\sum_k|Q_k|\le s^{-1}\|f\|_1$ oraz
 $\|g\|_2^2\le 2^ds\|f\|_1$.
\end{enumerate}
\end{lemat}
\begin{proof}
Spośród kostek jednej ustalonej siatki dyadycznej wybieramy
maksymalne względem inkluzji kostki, na których średnia
$|f|$ przekracza $s$. Nie istnieje nieskończony łańcuch
rosnących kostek złych: ich objętości rosną do nieskończoności,
a całka z $|f|$ pozostaje skończona. Maksymalne kostki są więc
rozłączne; ich ojcowie mają średnią nie większą niż $s$.
Stąd
$ s<|Q_k|^{-1}\int_{Q_k}|f|\le2^ds$.
Poza ich sumą mamy $|f|\le s$ prawie wszędzie, co wynika z
twierdzenia o różniczkowaniu całek Lebesgue'a dla zagnieżdżonych
kostek dyadycznych. Kładziemy
$g=f$ poza kostkami oraz $g=f_{Q_k}$ na każdej $Q_k$,
$ b_k=(f-f_{Q_k})\chi_{Q_k}$, gdzie
$f_{Q_k}=|Q_k|^{-1}\int_{Q_k}f$.
Teraz średnia każdej $b_k$ jest zerowa, a
$\|b_k\|_1\le\int_{Q_k}|f|+|Q_k||f_{Q_k}|
\le2\int_{Q_k}|f|$.
Zsumowanie nierówności $s|Q_k|<\int_{Q_k}|f|$
daje (iii) dla objętości. Ponieważ $|g|\le2^ds$ i
$\|g\|_1\le\|f\|_1$, także ostatni warunek (iii) jest
spełniony. Rozkład ma zbieżność w $L^1$ i $L^2$, gdyż
$\sum_k b_k=f-g$, a funkcje $b_k$ mają rozłączne nośniki.
\end{proof}

\begin{twierdzenie}[Słaby typ $(1,1)$ operatorów $T_{ij}$]
\label{thm:pde16-weak11}
Istnieje $C_d>0$ taka, że dla każdego
$f\in L^1(\R^d)\cap L^2(\R^d)$ i $t>0$ zachodzi
\begin{equation}\label{eq:pde16-weak11}
 \lambda_d\bigl(\{x:|T_{ij}f(x)|>t\}\bigr)
 \le \frac{C_d}{t}\,\|f\|_{L^1(\R^d)}.
\end{equation}
\end{twierdzenie}
\begin{proof}
Stosujemy lemat~\ref{lem:pde16-cz-decomp} z $s=t$.
Niech $Q_k^*$ oznacza kostkę o tym samym środku co $Q_k$,
lecz boku $3\sqrt d$ razy większym, oraz
$E=\bigcup_k Q_k^*$. Wtedy
$|E|\le C_d\|f\|_1/t$.
Dla części dobrej nierówność Czebyszewa i
lemat~\ref{lem:pde16-l2} dają
\[
 \lambda_d(\{|T_{ij}g|>t/2\})
 \le 4t^{-2}\|T_{ij}g\|_2^2
 \le C_dt^{-1}\|f\|_1.
\]
Dla $x\notin Q_k^*$, dzięki $\int_{Q_k} b_k\,\mathrm d\lambda_d=0$, mamy
$T_{ij}b_k(x)=\int_{Q_k}[K_{ij}(x-y)-K_{ij}(x-c_k)]
 b_k(y)\,\mathrm d\lambda_d(y)$, gdzie $c_k$ jest środkiem
$Q_k$. Z twierdzenia o wartości średniej i oszacowania
$|\nabla K_{ij}(z)|\le C_d|z|^{-d-1}$ otrzymujemy
\[
 |T_{ij}b_k(x)|\le C_d\int_{Q_k}
 \frac{|y-c_k|\,|b_k(y)|}{|x-c_k|^{d+1}}
 \mathrm d\lambda_d(y).
\]
Całkowanie po $\R^d\setminus Q_k^*$ w współrzędnych
biegunowych, zaczynając od promienia proporcjonalnego do
boku $Q_k$, prowadzi do
$\int_{\R^d\setminus Q_k^*}|T_{ij}b_k|
\,\mathrm d\lambda_d\le C_d\|b_k\|_1$.
Sumujemy po $k$; z twierdzenia Tonellego wynika
$\int_{\R^d\setminus E}|T_{ij}(\sum b_k)|
\,\mathrm d\lambda_d\le C_d\sum_k\|b_k\|_1
\le C_d\|f\|_1$.
Ostatnią nierówność można najpierw uzasadnić dla
skończonych sum, a następnie przejść do granicy, korzystając
ze zbieżności w $L^2$ i podciągu prawie wszędzie.
Nierówność Czebyszewa daje miarę zbioru, na którym
$|T_{ij}(\sum b_k)|>t/2$, poza $E$, nie większą niż
$C_d\|f\|_1/t$. Połączenie trzech oszacowań kończy dowód.
\end{proof}

\begin{lemat}[Interpolacja pomiędzy słabym $L^1$ i silnym $L^2$]
\label{lem:pde16-interp}
Niech $T$ będzie operatorem liniowym o słabym typie $(1,1)$
i ograniczonym na $L^2$ na $\R^d$. Wtedy dla każdego
$1<p<2$ istnieje $C_p$ takie, że
$\|Tf\|_p\le C_p\|f\|_p$, początkowo dla
$f\in L^1\cap L^2\cap L^p$, a następnie przez gęstość.
\end{lemat}
\begin{proof}
Dla $t>0$ rozkładamy $f=f_t^{>}+f_t^{\le}$,
gdzie $f_t^{>}=f\chi_{\{|f|>t\}}$.
Z liniowości $T$ i obu założeń
\[
 \lambda_d(\{|Tf|>2t\})
 \le \frac C{t}\int_{\{|f|>t\}}|f|\,\mathrm d\lambda_d
 +\frac C{t^2}\int_{\{|f|\le t\}}|f|^2\,\mathrm d\lambda_d.
\]
Korzystamy z tożsamości warstwowej
$\|Tf\|_p^p=p\int_0^\infty s^{p-1}
\lambda_d(\{|Tf|>s\})\,\mathrm ds$.
Po podstawieniu $s=2t$ i użyciu Tonellego pierwszy
składnik sprowadza się do
\[
 C\int_{\R^d} |f(x)|\left(\int_0^{|f(x)|}t^{p-2}
\,\mathrm dt\right)\mathrm d\lambda_d(x)
 =\frac C{p-1}\|f\|_p^p.
\]
Drugi jest równy co najwyżej
$C\int_{\R^d} |f(x)|^2\int_{|f(x)|}^\infty t^{p-3}
\,\mathrm dt\,\mathrm d\lambda_d(x)
 =C(2-p)^{-1}\|f\|_p^p$.
To uzasadnia oszacowanie z zależnością stałej od $p$.
\end{proof}

\begin{twierdzenie}[Oszacowanie Calderóna--Zygmunda dla laplasjanu]
\label{thm:pde16-cz}
Niech $d\ge2$ i $1<p<\infty$. Dla
$f\in C_c^\infty(\R^d)$ oraz $u=\Gamma*f$ zachodzi
\begin{equation}\label{eq:pde16-cz}
 \|D^2u\|_{L^p(\R^d)}\le C_{d,p}\|f\|_{L^p(\R^d)}.
\end{equation}
Oszacowanie przedłuża operator $f\mapsto D^2(\Gamma*f)$
do $L^p$; nie oznacza ono, że samo $u$ musi należeć
do $L^p(\R^d)$.
\end{twierdzenie}
\begin{proof}
Z twierdzenia~\ref{thm:pde16-weak11},
lematu~\ref{lem:pde16-l2} i
lematu~\ref{lem:pde16-interp} wynika ograniczoność
$T_{ij}$ w $L^p$ dla $1<p\le2$.
Dla $p>2$ wykorzystujemy dualność $L^p$ i $L^{p'}$:
sprzężony operator $T_{ij}^*$ ma jądro
$K_{ij}(-x)$, które spełnia takie same oszacowania i jest
również ograniczone na $L^{p'}$, bo $1<p'<2$.
Dla funkcji testowej $h$ o normie $L^{p'}$ nie większej niż
jeden mamy
$|\int_{\R^d}(T_{ij}f)h\,\mathrm d\lambda_d|=|\int_{\R^d} f(T_{ij}^*h)\,\mathrm d\lambda_d|
\le C\|f\|_p$; supremum po $h$ daje oszacowanie w $L^p$.
Wzór~\eqref{eq:pde16-repr} i nierówność trójkąta kończą
dowód po zsumowaniu po $i,j$.
\end{proof}

\begin{uwaga}[Dlaczego przypadki $1$ i $\infty$ są osobne?]
\label{rem:pde16-endpoints}
Twierdzenie zapewnia wszystkie wykładniki $1<p<\infty$,
lecz nie daje silnego oszacowania $L^1\to L^1$ ani
$L^\infty\to L^\infty$. Słabe oszacowanie
\eqref{eq:pde16-weak11} jest rzeczywistym wynikiem końcowym.
Na przykład w $d=2$ weźmy $f=\chi_{(0,1)^2}$ i
$u=\Gamma*f$. Dla $x=(-t,-t)$, $t\downarrow0$,
$\partial_{12}u(x)=\int_{(0,1)^2}K_{12}(x-y)
\,\mathrm d\lambda_2(y)$ rośnie co do modułu jak
$\log(1/t)$: na wycinku pierścienia pierwszej ćwiartki
$C t<|y|<1/2$ jądro jest dodatnie i rzędu $|y|^{-2}$.
Zatem ograniczona prawa strona nie gwarantuje ograniczoności
wszystkich drugich pochodnych. Wartość funkcji
$\partial_{12}u$ w punktach spoza nośnika $f$ jest
zwykłą całką, więc przykład nie wymaga interpretacji głównej
wartości w tych punktach.
\end{uwaga}

\section[Lokalna regularność Sobolewa i zamrażanie współczynników]{Lokalna regularność $W^{2,p}$ i zamrażanie współczynników}
\begin{lemat}[Lokalizacja oszacowania dla laplasjanu]
\label{lem:pde16-localization}
Niech $1<p<\infty$, $B_r\Subset B_R\subset\R^d$ i
$u\in W^{2,p}(B_R)$. Wówczas
\[
 \|D^2u\|_{L^p(B_r)}\le C\bigl(
 \|\Delta u\|_{L^p(B_R)}+
 (R-r)^{-2}\|u\|_{L^p(B_R)}+
 (R-r)^{-1}\|\nabla u\|_{L^p(B_R)}\bigr).
\]
\end{lemat}
\begin{proof}
Wybieramy $\eta\in C_c^\infty(B_R)$ równą jeden na $B_r$,
$|D^k\eta|\le C_k(R-r)^{-k}$ dla $k=1,2$, a następnie
przedłużamy $v=\eta u$ przez zero na $\R^d$.
Z równości $\Delta v=\eta\Delta u+2\nabla\eta\cdot\nabla u
+u\Delta\eta$ i~\eqref{eq:pde16-cz} dla $v$
wynika teza (przez aproksymację $W^{2,p}$ funkcjami gładkimi).
\end{proof}

\begin{lemat}[Interpolacja pierwszych pochodnych]
\label{lem:pde16-gradinterp}
Niech $1<p<\infty$, $B_r\Subset B_R$, $u\in W^{2,p}(B_R)$.
Dla każdego $\varepsilon>0$ zachodzi
\[
 \|\nabla u\|_{L^p(B_r)}\le
 \varepsilon\|D^2u\|_{L^p(B_R)}
 +C_{\varepsilon,r,R}\|u\|_{L^p(B_R)}.
\]
\end{lemat}
\begin{proof}
Na $\R^d$ dla $v\in C_c^\infty$ i dowolnego $h>0$
podstawowe twierdzenie rachunku całkowego zastosowane
do $t\mapsto v(x+te_j)$, a następnie do pochodnej tej
funkcji daje (po całkowaniu z odpowiednią wagą)
\[
 \partial_jv(x)=\frac{v(x+he_j)-v(x)}h
 -\frac1h\int_0^h(h-t)\partial_{jj}v(x+te_j)\,\mathrm dt.
\]
Nierówność Minkowskiego daje
$\|\partial_jv\|_p\le 2h^{-1}\|v\|_p+(h/2)\|\partial_{jj}v\|_p$.
Stosujemy teraz tę samą punktową tożsamość bezpośrednio
do $u$ dla $x\in B_r$ i $0<h<(R-r)/2$.
Wówczas cały odcinek $x+[0,h]e_j$ leży w $B_R$.
Po scałkowaniu nierówności na $B_r$ i zastosowaniu
Minkowskiego otrzymujemy
$\|\partial_ju\|_{L^p(B_r)}\le
2h^{-1}\|u\|_{L^p(B_R)}+(h/2)\|\partial_{jj}u\|_{L^p(B_R)}$.
Wybierając $h$ tak małe, aby $h/2\le\varepsilon/d$,
i sumując po $j$, uzyskujemy tezę.
Aproksymacja w $W^{2,p}$ usprawiedliwia użyty wzór
również dla funkcji niegładkich.
\end{proof}

\begin{twierdzenie}[Wewnętrzne oszacowanie dla współczynników ciągłych]
\label{thm:pde16-variable}
Niech $\Omega\subset\R^d$ będzie obszarem, $1<p<\infty$,
$a_{ij}\in C(\Omega)$ i niech symetryczna macierz
$A(x)=(a_{ij}(x))$ spełnia na każdym zbiorze względnie
zwartym warunki $\lambda|\xi|^2\le\xi^{\mathsf T}A(x)\xi
\le\Lambda|\xi|^2$.
Jeżeli $u\in W^{2,p}_{\mathrm{loc}}(\Omega)$ i
$a_{ij}\partial_{ij}u=f\in L^p_{\mathrm{loc}}(\Omega)$ prawie wszędzie,
to dla $U\Subset V\Subset\Omega$
\begin{equation}\label{eq:pde16-variable}
 \|u\|_{W^{2,p}(U)}\le C\bigl(
 \|f\|_{L^p(V)}+\|u\|_{L^p(V)}\bigr).
\end{equation}
Stała zależy od $d,p,U,V$, eliptyczności i modułu ciągłości $A$
na $V$. Twierdzenie jest tu \emph{oszacowaniem a priori}: samo
w sobie nie konstruuje rozwiązania $W^{2,p}$.
\end{twierdzenie}
\begin{proof}
Ustalmy $x_0\in V$ i zamroźmy współczynniki w $A_0=A(x_0)$.
Po liniowej zamianie zmiennych $z=A_0^{-1/2}(x-x_0)$
operator $a_{ij}(x_0)\partial_{ij}$ przechodzi w laplasjan.
Stałe tej zamiany są kontrolowane przez $\lambda,\Lambda$.
Dla $v\in W^{2,p}$ o zwartym nośniku w dostatecznie małej
kuli wokół $x_0$ z oszacowania Calderóna--Zygmunda
otrzymujemy
\[
 \|D^2v\|_p\le C\|A_0:D^2v\|_p
 \le C\|A(x):D^2v\|_p+C\|A-A_0\|_{\infty}\|D^2v\|_p.
\]
Z ciągłości $A$ możemy zmniejszyć promień kuli tak, aby
ostatni współczynnik nie przekraczał $1/2$ i pochłonąć
ten wyraz do lewej strony. Dla $v=\eta u$ obliczamy
$A:D^2(\eta u)=\eta f+2A\nabla\eta\cdot\nabla u
+u\,A:D^2\eta$.
Otrzymujemy więc lokalne oszacowanie $D^2u$ przez
$f$, $u$ i $\nabla u$ na trochę większej kuli. Stosujemy
lemat~\ref{lem:pde16-gradinterp}, wybierając $\varepsilon$
tak, by pochłonąć składnik z $D^2u$; aby uniknąć
błędnego koła, przeprowadzamy to najpierw na zagnieżdżonej
rodzinie kul i sumujemy oszacowania z geometrycznie
malejącymi wagami. Skończone pokrycie $\overline U$ przez
otwarte kule zawarte w $V$ oraz partycja jedności
kończą dowód~\eqref{eq:pde16-variable}.
\end{proof}

\begin{uwaga}[Postać dywergencyjna a niedywergencyjna]
Jeżeli $A\in W^{1,\infty}_{\mathrm{loc}}$ i
$-\diver(A\nabla u)=f$, to dopiero \emph{po} wykazaniu
$u\in H^2_{\mathrm{loc}}$ wolno napisać prawie wszędzie
$-a_{ij}\partial_{ij}u=f+(\partial_i a_{ij})\partial_ju$.
Aby przejść do $W^{2,p}$ trzeba niezależnie sprawdzić,
że prawa strona należy do $L^p$; samo $f\in L^p$
bez dostatecznej kontroli $\nabla u$ nie uzasadnia takiego
wniosku.
\end{uwaga}

\section{Przestrzenie Höldera i regularność w przestrzeniach Höldera}
\begin{definicja}[Norma Höldera]
\label{def:pde16-holder}
Dla $0<\alpha<1$ i ograniczonego zbioru $U\subset\R^d$
definiujemy
$[f]_{C^{0,\alpha}(U)}=
\sup_{x\ne y\in U}|f(x)-f(y)|/|x-y|^\alpha$.
Przestrzeń $C^{k,\alpha}(\overline U)$ składa się z funkcji,
których pochodne do rzędu $k$ przedłużają się ciągle na
$\overline U$, a pochodne rzędu $k$ mają skończoną półnormę
Höldera. Normą jest
$\|f\|_{C^{k,\alpha}(\overline U)}
=\sum_{|\beta|\le k}\|\partial^\beta f\|_{\infty,U}
+\sum_{|\beta|=k}[\partial^\beta f]_{C^{0,\alpha}(U)}$.
\end{definicja}

\begin{lemat}[Potencjał Newtona: oszacowanie Höldera]
\label{lem:pde16-newton-holder}
Niech $0<\alpha<1$, $f\in C_c^{0,\alpha}(\R^d)$,
$\operatorname{supp}f\subset B_R(0)$ i $u=\Gamma*f$.
Wtedy $u\in C^{2,\alpha}_{\mathrm{loc}}(\R^d)$ i dla
każdego $M>0$
\[
 \|D^2u\|_{C^{0,\alpha}(B_M)}
 \le C(d,\alpha,R,M)\|f\|_{C^{0,\alpha}(\R^d)}.
\]
W szczególności równanie $-\Delta u=f$ przy danych
Hölderowskich daje lokalnie drugie pochodne Hölderowskie.
\end{lemat}
\begin{proof}
Zaczynamy od $f\in C_c^\infty$; wynik ogólny dostaniemy
przez wygładzanie z jednostajnymi oszacowaniami.
Równość~\eqref{eq:pde16-repr} oraz zerowa średnia jądra
po sferach dają
\begin{equation}\label{eq:pde16-cancel}
 \partial_{ij}u(x)=-\frac{\delta_{ij}}d f(x)
 +\int_{|z|<1}K_{ij}(z)[f(x-z)-f(x)]\,\mathrm d\lambda_d(z)
 +\int_{|z|\ge1}K_{ij}(z)f(x-z)\,\mathrm d\lambda_d(z).
\end{equation}
Pierwsza całka jest absolutnie zbieżna, bo
$|K_{ij}(z)[f(x-z)-f(x)]|\le C[f]_\alpha|z|^{\alpha-d}$,
a $\int_0^1 r^{\alpha-1}\,\mathrm dr=1/\alpha$.
Druga również jest absolutnie zbieżna dla każdego $x$,
ponieważ $f$ ma zwarty nośnik.

Wykażemy oszacowanie różnicy tych dwóch całek w punktach
$x,x'$ z $|x-x'|=h<1/4$. W każdej całce przepisujemy
zmienną $y=x-z$ i rozdzielamy integrację na część
$|x-y|<2h$ (oraz analogiczną kulę wokół $x'$) i dopełnienie.
W części bliskiej odjęcie stałej $f(x)$ lub $f(x')$ daje
$C[f]_\alpha\int_0^{3h}r^{\alpha-1}\,\mathrm dr
\le C[f]_\alpha h^\alpha$.
W dopełnieniu odejmujemy te same stałe i korzystamy z
$|K_{ij}(x-y)-K_{ij}(x'-y)|\le
Ch|x-y|^{-d-1}$, ważnej dla $|x-y|>2h$.
Odpowiadający wkład z różnicy jąder ogranicza całka
\[
 C h[f]_\alpha\int_{2h}^1r^{\alpha-2}\,\mathrm dr
 \le C_\alpha[f]_\alpha h^\alpha.
\]
Pozostałe wyrazy pochodzące ze zmiany odjętych stałych,
ze sfer granicznych i z części $|z|\ge1$ są ograniczone
przez $C\|f\|_{C^{0,\alpha}}(h^\alpha+h)$;
na zewnętrznych częściach jądro i jego pochodna są
całkowalne po odpowiednich zwartych przekrojach nośnika.
W ten sposób
$|D_{ij}u(x)-D_{ij}u(x')|\le C h^\alpha\|f\|_{C^{0,\alpha}}$.
Dla $h\ge1/4$ teza wynika z oszacowania supremum w
\eqref{eq:pde16-cancel} i ograniczoności $x,x'$ do $B_M$.

Dla $f\in C_c^{0,\alpha}$ wybieramy wygładzenia $f_n$ o
wspólnym zwartym nośniku i jednostajnie ograniczonych
normach Höldera. Potencjały $u_n=\Gamma*f_n$ zbiega
lokalnie w $C^1$, zaś drugie pochodne są jednostajnie
ograniczone w $C^{0,\alpha}$; Arzelà--Ascoli daje lokalnie
jednostajną zbieżność podciągu $D^2u_n$ do funkcji ciągłych.
Granice te są dystrybucyjnymi drugimi pochodnymi $u$,
a półnormy Höldera przechodzą do granicy przez
nierówność punktową. Dostajemy $u\in C^{2,\alpha}_{\rm loc}$
i żądane oszacowanie.
\end{proof}

\begin{twierdzenie}[Wewnętrzne oszacowanie Schaudera: model Poissona]
\label{thm:pde16-schauder-poisson}
Niech $0<\alpha<1$, $B_r\Subset B_R\subset\R^d$, a
$u\in C^2(B_R)\cap C^0(\overline{B_R})$ spełnia
$-\Delta u=f$ w $B_R$ dla $f\in C^{0,\alpha}(\overline{B_R})$.
Wówczas $u\in C^{2,\alpha}(\overline{B_r})$ i
\begin{equation}\label{eq:pde16-schauder-int}
 \|u\|_{C^{2,\alpha}(\overline{B_r})}
 \le C\bigl(\|u\|_{C^0(\overline{B_R})}
 +\|f\|_{C^{0,\alpha}(\overline{B_R})}\bigr).
\end{equation}
Stała zależy od $d,\alpha,r,R$.
\end{twierdzenie}
\begin{proof}
Wybieramy $B_r\Subset B_\rho\Subset B_R$ i funkcję odcinającą
$\eta\in C_c^\infty(B_R)$ równą jeden na $B_\rho$.
Zastosujmy poprzedni lemat do przedłużonej przez zero
funkcji $F=\eta f$: funkcja $v=\Gamma*F$ ma drugie pochodne
Hölderowskie na $B_\rho$ z normą ograniczoną przez normę $f$.
W $B_\rho$ mamy $-\Delta(u-v)=0$. Funkcja $h=u-v$ jest
więc harmoniczna i gładka (rozdział~\ref{chap:pde-laplace},
także lemat Weyla). Z własności wartości średniej i
różniczkowania jądra Poissona na kulach o promieniu
mniejszym niż $\rho-r$ dostajemy dla $k\le3$
\[
 \|D^kh\|_{L^\infty(B_r)}
 \le C_{k,d}(\rho-r)^{-k}\|h\|_{L^\infty(B_\rho)}.
\]
W szczególności $[D^2h]_{C^{0,\alpha}(B_r)}$ ograniczamy
przez $\|D^3h\|_\infty\operatorname{diam}(B_r)^{1-\alpha}$.
Ponieważ $\|v\|_{L^\infty(B_\rho)}\le C\|f\|_\infty$
(z lokalnej całkowalności $\Gamma$ i zwartego nośnika $F$),
odtwarzamy wszystkie składniki normy po lewej
stronie~\eqref{eq:pde16-schauder-int}. W dowodzie nie
zakładamy z góry Hölderowskiej regularności $D^2u$.
\end{proof}

\begin{twierdzenie}[Pełny przypadek: skalarny współczynnik w postaci niedywergencyjnej]
\label{thm:pde16-scalar-schauder}
Niech $d\ge2$, $0<\alpha<1$, $\Omega\subset\R^d$ będzie obszarem,
$a,f\in C^{0,\alpha}_{\mathrm{loc}}(\Omega)$, przy czym
$a(x)\ge\lambda>0$ dla każdego $x\in\Omega$.
Jeżeli $u\in C^2(\Omega)$ spełnia $a(x)\Delta u=f$ w $\Omega$,
to $u\in C^{2,\alpha}_{\mathrm{loc}}(\Omega)$.
Dla $U\Subset V\Subset\Omega$ zachodzi
\[
\|u\|_{C^{2,\alpha}(\overline U)}
\le C\bigl(\|f\|_{C^{0,\alpha}(\overline V)}
+\|u\|_{C^0(\overline V)}\bigr),
\]
gdzie $C$ zależy tylko od $d,\alpha,U,V,\lambda$ i
$\|a\|_{C^{0,\alpha}(\overline V)}$.
\end{twierdzenie}
\begin{proof}
Na $\overline V$ funkcja $1/a$ należy do $C^{0,\alpha}$,
gdyż $|a(x)^{-1}-a(y)^{-1}|
\le\lambda^{-2}[a]_{C^{0,\alpha}}|x-y|^\alpha$.
W szczególności $g:=f/a\in C^{0,\alpha}(\overline V)$ i
\[
\|g\|_{C^{0,\alpha}(\overline V)}
\le\lambda^{-1}\|f\|_{C^{0,\alpha}(\overline V)}
+\lambda^{-2}\|f\|_\infty[a]_{C^{0,\alpha}(\overline V)}.
\]
Z równania $\Delta u=g$ wynika $-\Delta u=-g$.
Wybieramy skończone pokrycie $\overline U$ kulami,
których domknięcia leżą w $V$, i stosujemy na każdej z nich
udowodnione twierdzenie~\ref{thm:pde16-schauder-poisson}
(z prawą stroną $-g$). Sumując lokalne oszacowania
otrzymujemy tezę. Ten dowód nie potrzebuje
nieudowodnionego kroku regularyzacyjnego z ogólnego
przypadku macierzowych współczynników.
\end{proof}


\begin{lemat}[Oszacowanie Schaudera \emph{a priori} i lokalizacja]
\label{lem:pde16-schauder-apriori}
Niech $0<\alpha<1$, $B_r\Subset B_R\Subset\Omega$, a
$A=(a_{ij})$, $b=(b_i)$, $c$ należą do
$C^{0,\alpha}(\overline{B_R})$; załóżmy, że $A$ jest
symetryczna i jednostajnie eliptyczna.
Jeżeli $v\in C^{2,\alpha}(\overline{B_R})$, to
\begin{equation}\label{eq:pde16-schauder-apriori}
 \|v\|_{C^{2,\alpha}(\overline{B_r})}
 \le C\bigl(\|A:D^2v+b\cdot\nabla v+cv\|_{C^{0,\alpha}(\overline{B_R})}
 +\|v\|_{C^0(\overline{B_R})}\bigr),
\end{equation}
gdzie $C$ zależy od $d,\alpha,r,R$, eliptyczności i norm
Hölderowskich współczynników, ale nie od $v$.
\end{lemat}
\begin{proof}
Najpierw wyjaśniamy dwie operacje użyte do pochłaniania
wyrazów, aby nie ukrywać założenia o regularności.
Dla kuli $B_\rho$ wprowadzamy normę skalowaną
\[
 \|w\|_{k,\alpha;B_\rho}^{*}
 =\sum_{j=0}^{k}\rho^j\|D^jw\|_{\infty,B_\rho}
 +\rho^{k+\alpha}[D^kw]_{C^{0,\alpha}(B_\rho)}.
\]
Dla $w\in C_c^{2,\alpha}(B_\rho)$ i stałej dodatnio
określonej macierzy $A_0$ po liniowej zmianie zmiennych
$z=A_0^{-1/2}x$ lemat~\ref{lem:pde16-newton-holder}
(stosowany do $-A_0:D^2w$ o zwartym nośniku)
oraz oszacowanie supremum potencjału dają
\begin{equation}\label{eq:pde16-frozen}
 \rho^2\|D^2w\|_{\infty}
 +\rho^{2+\alpha}[D^2w]_{\alpha}
 \le C\rho^2\|A_0:D^2w\|_{0,\alpha;B_\rho}^{*}.
\end{equation}
Stała nie zależy od promienia: podstawiamy $x=\rho z$
i korzystamy z jednorodności operatora drugiego rzędu.

Ustalmy środek $x_0$ i zamroźmy $A_0=A(x_0)$.
Dla $w=\eta v$ ze standardowym obcięciem $\eta$ równym
$1$ na $B_{\rho/2}(x_0)$, zerowym poza $B_\rho(x_0)$,
$\|D^j\eta\|_\infty\le C_j\rho^{-j}$, mamy
\[
 A_0:D^2(\eta v)=\eta Lv+(A_0-A):D^2(\eta v)
 +\mathcal R_\eta,
 \qquad Lv=A:D^2v+b\cdot\nabla v+cv,
\]
gdzie $\mathcal R_\eta$ jest sumą wyrazów zawierających
$D\eta,D^2\eta,v,Dv,b,c$; każdy wynika z reguły
Leibniza. Na $B_\rho(x_0)$ zachodzi
\[
 \|A-A_0\|_{0,\alpha;B_\rho}^{*}
 \le C[A]_{C^{0,\alpha}(B_R)}\rho^\alpha.
\]
Stąd, korzystając z nierówności dla iloczynów w
$C^{0,\alpha}$ i~\eqref{eq:pde16-frozen}, otrzymujemy
na dowolnych dwóch współśrodkowych kulach o promieniach
$\rho/2$ i $\rho$ oszacowanie
\begin{align*}
 \rho^2\|D^2v\|_{0,\alpha;B_{\rho/2}}^{*}
 &\le C\bigl(\rho^2\|Lv\|_{0,\alpha;B_\rho}^{*}
       +\|v\|_{\infty,B_\rho}
       +\rho\|Dv\|_{0,\alpha;B_\rho}^{*}\bigr)\\
 &\quad+C\bigl(\rho^\alpha[A]_\alpha
       +\rho\|b\|_{0,\alpha;B_R}^{*}
       +\rho^2\|c\|_{0,\alpha;B_R}^{*}\bigr)
       \|v\|_{2,\alpha;B_\rho}^{*}.
\end{align*}
Tutaj po lewej przez $\rho^2\|D^2v\|_{0,\alpha}^{*}$
rozumiemy dokładnie
$\rho^2\|D^2v\|_\infty+\rho^{2+\alpha}[D^2v]_\alpha$;
analogicznie dla wyrazów z $Lv$ i $Dv$.
Nierówność interpolacyjna
\[
 \rho\|Dv\|_{0,\alpha;B_\rho}^{*}
 \le\varepsilon\|v\|_{2,\alpha;B_\rho}^{*}
 +C_\varepsilon\|v\|_{\infty,B_\rho}
\]
wynika przez przeskalowanie do kuli jednostkowej.
Dla funkcji na całej przestrzeni wzór Taylora wzdłuż
odcinka długości $h$ daje
$\|Dv\|_\infty\le 2h^{-1}\|v\|_\infty+
(h/2)\|D^2v\|_\infty$. Ponieważ $Dv$ jest funkcją
Lipschitzowską, rozdzielenie ilorazu Höldera na przypadki
$|x-y|\le t$ oraz $|x-y|>t$ daje
\[
 [Dv]_\alpha\le
 2^{1-\alpha}\|Dv\|_\infty^{1-\alpha}
 \|D^2v\|_\infty^\alpha.
\]
Podstawiając pierwszy szacunek, wybierając $h$ i $t$
i stosując nierówność Younga, otrzymujemy współczynnik
$\varepsilon$ przy najwyższej normie.
Na kulach korzystamy z tych samych wzorów na odcinkach
zawartych w nieco większej współśrodkowej kuli;
standardowe przedłużenie z kuli klasy $C^{2,\alpha}$
pozwala sformułować oszacowanie również na całej kuli.
Stałe po przeskalowaniu nie zależą od $\rho$.
Wybieramy promień $\rho_0>0$ i następnie
$\varepsilon>0$ na tyle małe, aby suma współczynników
przy najwyższych normach była dowolnie mała.

Dla ścisłego sklejenia nierówności używamy ważonej normy
w kuli $B_R$: niech
$\delta(x)=\frac18\min\{\rho_0,\operatorname{dist}(x,\partial B_R)\}$.
Definiujemy $\mathcal N(v)$ jako supremum
$\delta(x)^2|D^2v(x)|$ po $x\in B_R$ oraz supremum
\[
 \delta(x)^{2+\alpha}
 \frac{|D^2v(x)-D^2v(y)|}{|x-y|^\alpha},
 \qquad x,y\in B_R,\ 0<|x-y|<\delta(x).
\]
Występujące w oszacowaniu dwie współśrodkowe kule
wybieramy o promieniach porównywalnych z $\delta(x)$.
Na większej kuli wartości $\delta$ różnią się co najwyżej
o stały czynnik, więc wszystkie normy drugich pochodnych
po prawej kontroluje $C\mathcal N(v)$. Ponieważ współczynnik
pochłaniania możemy wybrać dowolnie mały, uzyskujemy
\[
 \mathcal N(v)\le\tfrac12\mathcal N(v)
 +C\bigl(\|Lv\|_{C^{0,\alpha}(\overline{B_R})}
       +\|v\|_{C^0(\overline{B_R})}\bigr).
\]
Wszystkie normy są skończone, gdyż przyjęliśmy
$v\in C^{2,\alpha}(\overline{B_R})$, zatem wolno pochłonąć
pierwszy składnik. Na $B_r$ funkcja $\delta$ jest oddzielona
od zera; dla par punktów oddalonych o więcej niż
$\inf_{B_r}\delta$ półnormę Höldera kontroluje
norma supremum drugich pochodnych. Otrzymujemy oszacowanie
$D^2v$ na $B_r$, a nierówności interpolacyjne kontrolują
$Dv$ i $v$. Jest to~\eqref{eq:pde16-schauder-apriori}.
\end{proof}

\begin{lemat}[Komutator wygładzania i mnożenia]
\label{lem:pde16-commutator}
Niech $0<\alpha<1$, $a\in C^{0,\alpha}(\R^d)$,
$g\in L^\infty(\R^d)$, a $\rho_\varepsilon$
będzie standardowym gładkim jądrem wygładzającym,
o nośniku w $B_\varepsilon(0)$. Dla
\[
 \mathcal C_\varepsilon[a,g](x)
 =a(x)(g*\rho_\varepsilon)(x)
 -(ag)*\rho_\varepsilon(x)
\]
zachodzi jednostajnie względem $0<\varepsilon\le1$
\[
 \|\mathcal C_\varepsilon[a,g]\|_{C^{0,\alpha}(\R^d)}
 \le C_{d,\alpha,\rho}[a]_{C^{0,\alpha}(\R^d)}\|g\|_\infty.
\]
Ponadto $\|\mathcal C_\varepsilon[a,g]\|_\infty
\le C\varepsilon^\alpha[a]_\alpha\|g\|_\infty$.
\end{lemat}
\begin{proof}
Z definicji
$\mathcal C_\varepsilon[a,g](x)
 =\int_{\R^d}(a(x)-a(y))g(y)\rho_\varepsilon(x-y)
\,\mathrm d\lambda_d(y)$, co daje oszacowanie supremum.
Niech $h=|x-x'|$. Jeśli $h\ge\varepsilon$,
to różnicę wartości komutatora szacujemy przez
$2C\varepsilon^\alpha\|g\|_\infty[a]_\alpha
\le2Ch^\alpha\|g\|_\infty[a]_\alpha$.
Gdy $0<h<\varepsilon$, piszemy dokładnie
\begin{align*}
 \mathcal C_\varepsilon(x)-\mathcal C_\varepsilon(x')
 &= (a(x)-a(x'))(g*\rho_\varepsilon)(x)\\
 &\quad+\int_{\R^d}(a(x')-a(y))g(y)
 [\rho_\varepsilon(x-y)-\rho_\varepsilon(x'-y)]
 \,\mathrm d\lambda_d(y).
\end{align*}
Pierwszy wyraz jest ograniczony przez
$[a]_\alpha\|g\|_\infty h^\alpha$.
W drugim $|x'-y|\le2\varepsilon$ na nośniku różnicy jąder,
a $\|\rho_\varepsilon(\cdot-x)-
\rho_\varepsilon(\cdot-x')\|_{L^1}\le
C h/\varepsilon$. Dlatego ten wyraz jest nie większy niż
$C[a]_\alpha\|g\|_\infty\varepsilon^{\alpha-1}h
\le C[a]_\alpha\|g\|_\infty h^\alpha$.
Połączenie przypadków kończy dowód.
\end{proof}

\begin{twierdzenie}[Oszacowania Schaudera dla zmiennych współczynników]
\label{thm:pde16-schauder-variable}
Niech $0<\alpha<1$, $\Omega\subset\R^d$ będzie obszarem,
$a_{ij},b_i,c,f\in C^{0,\alpha}_{\mathrm{loc}}(\Omega)$,
a macierz $A=(a_{ij})$ będzie symetryczna i lokalnie
jednostajnie eliptyczna. Jeżeli $u\in C^2(\Omega)$ spełnia
$a_{ij}\partial_{ij}u+b_i\partial_i u+cu=f$, to
$u\in C^{2,\alpha}_{\mathrm{loc}}(\Omega)$ oraz dla
$U\Subset V\Subset\Omega$
\[
 \|u\|_{C^{2,\alpha}(\overline U)}
 \le C\bigl(\|f\|_{C^{0,\alpha}(\overline V)}
 +\|u\|_{C^0(\overline V)}\bigr),
\]
gdzie $C$ zależy od $U,V,d,\alpha$ i lokalnych oszacowań
współczynników oraz eliptyczności.
\end{twierdzenie}
\begin{proof}
\emph{Krok 1: regularizacja bez zakładania tezy.}
Wybierzmy kule $B_r\Subset B_R\Subset W\Subset\Omega$.
Bierzemy $\chi\in C_c^\infty(W)$, równą $1$ w otoczeniu
$\overline{B_R}$, i kładziemy
$u_\varepsilon=(\chi u)*\rho_\varepsilon$.
Dla dostatecznie małych $\varepsilon$ funkcja $u_\varepsilon$
jest gładka w $B_R$, a $u_\varepsilon\to u$ w
$C^2(\overline{B_R})$. Nie wygładzamy współczynników:
na $B_R$ równanie dla $u_\varepsilon$ jest dokładnie
\begin{align*}
 A:D^2u_\varepsilon+b\cdot\nabla u_\varepsilon
   +cu_\varepsilon
 &= f*\rho_\varepsilon
 +\sum_{i,j}\mathcal C_\varepsilon[a_{ij},\partial_{ij}u]
 +\sum_i\mathcal C_\varepsilon[b_i,\partial_i u]
 +\mathcal C_\varepsilon[c,u].
\end{align*}
Wszystkie komutatory liczymy po lokalnym przedłużeniu
współczynników i pochodnych z $W$ na $\R^d$; wartości
w $B_R$ nie zależą od sposobu przedłużenia, jeśli
$\varepsilon$ jest mniejsze niż odstęp $B_R$ od
$\partial W$. Ponieważ $u\in C^2(W)$, jego pochodne
do rzędu drugiego są ograniczone na $\operatorname{supp}\chi$.
Lemat~\ref{lem:pde16-commutator} oraz zwykłe oszacowanie
splotu w $C^{0,\alpha}$ pokazują, że prawa strona tego
równania ma normę $C^{0,\alpha}(\overline{B_R})$
\emph{jednostajnie ograniczoną} względem $\varepsilon$.
Nie zakładamy Hölderowskiej ciągłości $D^2u$.

\emph{Krok 2: kontrola norm i granica.}
Ponieważ $u_\varepsilon\in C^{2,\alpha}(\overline{B_R})$,
możemy do niego zastosować lemat
\ref{lem:pde16-schauder-apriori}. Otrzymujemy
$\sup_\varepsilon\|u_\varepsilon\|_{C^{2,\alpha}
(\overline{B_r})}<\infty$. Dla każdych $x,y\in B_r$
przechodzimy w oszacowaniu
$|D^2u_\varepsilon(x)-D^2u_\varepsilon(y)|\le C|x-y|^\alpha$
do granicy, ponieważ $D^2u_\varepsilon\to D^2u$
jednostajnie. Wobec tego $u\in C^{2,\alpha}(\overline{B_r})$.
Dowolność kul daje $u\in C^{2,\alpha}_{\mathrm{loc}}(\Omega)$.

\emph{Krok 3: usunięcie zależności od $\|D^2u\|_\infty$.}
Powyższa jednostajna kontrola komutatorów służyła tylko
do wykazania przynależności do $C^{2,\alpha}$; jej stała
mogła zależeć od lokalnej normy $C^2$ rozwiązania.
Teraz stosujemy lemat~\ref{lem:pde16-schauder-apriori}
bezpośrednio do $u$ na każdej parze kul o domknięciu
zawartym w $V$, gdzie $Lu=f$. Otrzymujemy oszacowanie
zależne wyłącznie od normy $f$ w $C^{0,\alpha}(\overline V)$
i normy $u$ w $C^0(\overline V)$. Skończone pokrycie
$\overline U$ kulami i kontrola par punktów należących
do różnych kul (przez supremum $D^2u$) kończą dowód.
\end{proof}

\begin{uwaga}[Granice formułowanych twierdzeń]
\label{rem:pde16-boundary}
Globalne oszacowanie $C^{2,\alpha}(\overline\Omega)$ dla
zagadnienia Dirichleta wymaga na ogół brzegu klasy
$C^{2,\alpha}$ i danych brzegowych w $C^{2,\alpha}$.
W przypadku $W^{2,p}$ typowym założeniem na brzeg jest
$C^{1,1}$. Nie utożsamiamy tych dwóch teorii.
Dowód globalnego twierdzenia o oszacowaniach Schaudera wymaga dodatkowych
półkulowych oszacowań brzegowych, których nie zastępuje
wyłącznie rezultat wewnętrzny.
\end{uwaga}


\begin{uwaga}[Postać równania i przestrzeń rozwiązania]
\label{rem:pde-v42-two-regularity-classes}
Twierdzenie~\ref{thm:pde15-boundary} dotyczy słabych
rozwiązań równania w postaci dywergencyjnej
$-\diver(A\nabla u)=f$ przy $A\in W^{1,\infty}$.
Natomiast twierdzenie~\ref{thm:pde16-schauder-variable}
dotyczy rozwiązań klasy $C^2$ równania
w postaci niedywergencyjnej
$a_{ij}\partial_{ij}u+b_i\partial_i u+cu=f$,
przy współczynnikach H\"olderowskich.
Nie wolno przenosić tezy $C^{2,\alpha}$ na rozwiązania
samej postaci dywergencyjnej z jedynie
$C^{0,\alpha}$-współczynnikami; różniczkowanie
współczynników byłoby wtedy nieuzasadnione.
\end{uwaga}

\section{Kontrprzykłady i granice regularności}
\begin{kontrprzyklad}[Brak ciągłości $D^2u$ przy ograniczonej prawej stronie]
Na przedziale $(-1,1)$ niech $u(x)=-\tfrac12x|x|$.
Wtedy $u\in W^{2,p}(-1,1)$ dla każdego $1\le p<\infty$,
$-u''=\operatorname{sgn}(x)$ prawie wszędzie, ale $u''$
ma skok w zerze, a więc $u\notin C^2(-1,1)$.
W wyższych wymiarach analogiczny przykład daje funkcja
zależna wyłącznie od jednej współrzędnej.
\end{kontrprzyklad}

\begin{kontrprzyklad}[Ostrość wykładnika Höldera]
Dla $0<\alpha<1$ połóżmy na $(-1,1)$
$u(x)=-|x|^{2+\alpha}/((1+\alpha)(2+\alpha))$.
Mamy $-u''=|x|^\alpha$ dla $x\ne0$ i także w sensie
klasycznym w zerze. Zatem $u\in C^{2,\alpha}$,
ale $u\notin C^{2,\beta}$ dla żadnego $\beta>\alpha$,
ponieważ $|u''(h)-u''(0)|/|h|^\beta=h^{\alpha-\beta}\to\infty$.
\end{kontrprzyklad}

\section{Podsumowanie i miejsce w dalszej teorii}
W oszacowaniach $L^p$ decydującą rolę odgrywa znoszenie się
jądra o średniej zero, rozkład Calderóna--Zygmunda i dualność.
W teorii oszacowań Schaudera ten sam mechanizm pojawia się jako
odejmowanie $f(x)$ pod całką osobliwą. Nie wystarczy
jednak powiedzieć, że $f$ jest ograniczona: różnica
$f(x-z)-f(x)$ musi zanikać dostatecznie szybko.
Późniejsze twierdzenia De Giorgiego--Nasha--Mosera pokażą, że
niektóre oszacowania Hölderowskie można otrzymać także przy
samych mierzalnych współczynnikach, ale dla \emph{samego
rozwiązania}, a nie automatycznie dla wszystkich jego
drugich pochodnych.

% ------------------------------------------------------------------
% ROZDZIAŁ 17. Teoria De Giorgiego--Nasha--Mosera
% ------------------------------------------------------------------
\chapter[Regularność De Giorgiego--Nasha--Mosera]{Teoria De Giorgiego--Nasha--Mosera: regularność przy mierzalnych współczynnikach}
\label{chap:pde17}

W rozdziałach~\ref{chap:pde15}--\ref{chap:pde16} uzyskiwaliśmy regularność
pochodnych przy odpowiedniej gładkości współczynników. Teraz przyjmujemy
\emph{jedynie mierzalność} macierzy $A$. Mimo to rozwiązanie słabe
jednostajnie eliptycznego równania skalarnego ma lokalnie reprezentanta
Hölderowskiego. Jest to szczególny przypadek twierdzenia
De Giorgiego--Nasha--Mosera. Przedstawimy metodę iteracji Mosera,
od oszacowań energetycznych aż do nierówności Harnacka i regularności.
Nie będziemy utożsamiać tej teorii z wyższą regularnością $W^{2,p}$:
przy skokowych współczynnikach pochodna rozwiązania może mieć skok.

\begin{uwaga}[Mapa założeń: zasada maksimum, Schauder i Harnack]
\label{rem:pde-v43-regularity-map}
Klasyczne zasady maksimum z rozdziału~\ref{chap:pde14}
wykorzystują regularność klasyczną albo dodatkowe argumenty
aproksymacyjne. Oszacowania Schaudera z
rozdziału~\ref{chap:pde16} wymagają regularności H\"olderowskiej
współczynników i dotyczą pochodnych drugiego rzędu.
W niniejszym rozdziale pracujemy z równaniem \emph{skalarnym}
w postaci dywergencyjnej, z mierzalnym, jednostajnie
eliptycznym $A$; regularność H\"olderowską otrzymujemy
wyłącznie dla samego rozwiązania. Nie przenosimy automatycznie
nierówności Harnacka na układy ani na równania z niekontrolowanymi
wyrazami niższego rzędu.
\end{uwaga}

Przyjmujemy oznaczenia $u_+=\max\{u,0\}$,
$\operatorname*{ess\,osc}_E u=
\operatorname*{ess\,sup}_E u-\operatorname*{ess\,inf}_E u$ oraz
$\fint_E h\,\mathrm d\lambda_d=|E|^{-1}
\int_E h\,\mathrm d\lambda_d$ dla zbiorów o dodatniej skończonej mierze.

\section{Założenia, słabe nierówności i skalowanie}
\label{sec:pde17-assumptions}
\begin{definicja}[Jednostajna eliptyczność i rozwiązania słabe]
Niech $d\ge2$, $\Omega\subset\R^d$ będzie obszarem, a
$A\colon\Omega\to\R^{d\times d}$ mierzalną macierzą, taką że dla
prawie każdego $x\in\Omega$, $\xi,\zeta\in\R^d$ zachodzą
\begin{equation}\label{eq:pde17-ellipticity}
 \xi^{\mathsf T}A(x)\xi\ge\lambda|\xi|^2,
 \qquad |A(x)\xi\cdot\zeta|\le\Lambda|\xi||\zeta|,
 \qquad 0<\lambda\le\Lambda<\infty.
\end{equation}
Nie wymagamy symetrii $A$. Funkcję $u\in H^1_{\rm loc}(\Omega)$
nazywamy słabym rozwiązaniem równania
$-\operatorname{div}(A\nabla u)=0$, jeżeli
\begin{equation}\label{eq:pde17-weak}
 \int_\Omega A\nabla u\cdot\nabla\varphi\,\mathrm d\lambda_d=0
 \quad\text{dla każdej }\varphi\in C_c^\infty(\Omega).
\end{equation}
Słabym podrozwiązaniem (odpowiednio nadrozwiązaniem) nazywamy funkcję,
dla której lewa strona~\eqref{eq:pde17-weak} jest $\le0$
(odpowiednio $\ge0$) dla każdej nieujemnej funkcji testowej.
\end{definicja}

\begin{uwaga}[Dopuszczalne testy]
Liniowość i ciągłość formy na $H^1(K)$, gdzie $K\Subset\Omega$,
pozwalają zastąpić $C_c^\infty$ przez funkcje $H^1$ o zwartym
nośniku w $\Omega$. Testowanie obcięciami i potęgami wykonujemy
najpierw dla funkcji ograniczonych i Lipschitzowskich po obcięciu
wartości $u$, a następnie przechodzimy do granicy przez zbieżność
zdominowaną i słabą dolną półciągłość energii. Nie zakładamy, że
$u$ jest gładka, ani nie różniczkujemy współczynników $A$.
\end{uwaga}

Jeżeli $u$ rozwiązuje równanie w $B_{2R}(x_0)$, funkcja
$v(y)=u(x_0+Ry)$ rozwiązuje je w $B_2(0)$ z macierzą
$A_R(y)=A(x_0+Ry)$ i \emph{tymi samymi} stałymi $\lambda,\Lambda$.
Z tego powodu stałe we wszystkich lokalnych nierównościach zależą
od $d$ i $\Lambda/\lambda$, lecz nie od środka ani promienia kuli.

\section{Nierówność Caccioppolego dla obcięć De Giorgiego}
\label{sec:pde17-caccioppoli}
\begin{lemat}[Energia części dodatniej]
\label{lem:pde17-caccioppoli}
Niech $u\in H^1(B_R)$ będzie podrozwiązaniem w $B_R\subset\Omega$.
Dla $k\in\R$ i $\eta\in C_c^1(B_R)$, $0\le\eta\le1$, zachodzi
\begin{equation}\label{eq:pde17-cacc}
 \int_{B_R}\eta^2|\nabla(u-k)_+|^2\,\mathrm d\lambda_d
 \le \frac{4\Lambda^2}{\lambda^2}
 \int_{B_R}(u-k)_+^2|\nabla\eta|^2\,\mathrm d\lambda_d.
\end{equation}
Analogicznie, dla nadrozwiązania oszacowanie dotyczy $(k-u)_+$.
\end{lemat}
\begin{proof}
Oznaczmy $w=(u-k)_+$. Funkcja $w$ należy do $H^1$;
$\nabla w=\chi_{\{u>k\}}\nabla u$ prawie wszędzie,
co wynika z reguły łańcuchowej dla funkcji Lipschitzowskich.
Testujemy nierówność podrozwiązania funkcją $\varphi=\eta^2w\ge0$.
Ponieważ $\nabla\varphi=\eta^2\nabla w+2\eta w\nabla\eta$
na $\{u>k\}$, dostajemy
\[
 \lambda\int_{B_R}\eta^2|\nabla w|^2\,\mathrm d\lambda_d
 \le 2\Lambda\int_{B_R}\eta w|\nabla w||\nabla\eta|\,\mathrm d\lambda_d.
\]
Nierówność Younga
$2\Lambda ab\le(\lambda/2)a^2+(2\Lambda^2/\lambda)b^2$
daje tezę. Dla nadrozwiązań testujemy $-(k-u)_+$ w równoważnej
nierówności, uwzględniając zmianę znaku.
\end{proof}

\begin{uwaga}[Kontrast z metodą ilorazów różnicowych]
W poprzednich rozdziałach do regularności $H^2$ wykorzystywaliśmy
pochodne $A$. W~\eqref{eq:pde17-cacc} wystarcza informacja o
$A\in L^\infty$ i dodatniości części symetrycznej macierzy.
To właśnie pozwala prowadzić argument dla współczynników mierzalnych.
\end{uwaga}

\section{Lemat iteracyjny De Giorgiego i lokalne ograniczenie}
\label{sec:pde17-degiorgi}
\begin{lemat}[Zbieżność iteracji superliniowej]
\label{lem:pde17-iteration}
Niech $Y_n\ge0$, $a>1$, $b>1$, $\delta>0$ oraz
$Y_{n+1}\le a b^n Y_n^{1+\delta}$ dla $n\ge0$.
Jeżeli $Y_0\le a^{-1/\delta}b^{-1/\delta^2}$,
to $Y_n\to0$.
\end{lemat}
\begin{proof}
Niech $Z_n=a^{1/\delta}b^{n/\delta+1/\delta^2}Y_n$.
Po podstawieniu do rekurencji otrzymujemy
$Z_{n+1}\le Z_n^{1+\delta}$. Ponieważ $0\le Z_0\le1$,
ciąg $Z_n$ jest ograniczony przez $1$, a zatem
$Y_n\le a^{-1/\delta}b^{-n/\delta-1/\delta^2}\to0$.
\end{proof}

\begin{twierdzenie}[Oszacowanie $L^2\to L^\infty$ podrozwiązania]
\label{thm:pde17-l2linfty}
Niech $B_{2R}(x_0)\Subset\Omega$ i $u\in H^1(B_{2R}(x_0))$
będzie podrozwiązaniem równania z~\eqref{eq:pde17-ellipticity}.
Wtedy
\begin{equation}\label{eq:pde17-l2linfty}
 \operatorname*{ess\,sup}_{B_R(x_0)}u_+
 \le C\left(\fint_{B_{2R}(x_0)}u_+^2\,\mathrm d\lambda_d\right)^{1/2},
 \qquad C=C(d,\Lambda/\lambda).
\end{equation}
\end{twierdzenie}
\begin{proof}
Skalowanie pozwala przyjąć $R=1$, $x_0=0$. Ustalmy
$K>0$ i ciągi
$r_n=1+2^{-n}$, $k_n=K(1-2^{-n-1})$, $n\ge0$.
Niech $w_n=(u-k_n)_+$ oraz
$Y_n=\int_{B_{r_n}}w_n^2\,\mathrm d\lambda_d$.
Wybieramy $\eta_n=1$ na $B_{r_{n+1}}$, o nośniku w $B_{r_n}$,
z $|\nabla\eta_n|\le C2^n$.
Z lematu~\ref{lem:pde17-caccioppoli} oraz nierówności Sobolewa
na kuli (z przedłużeniem przez zero funkcji $\eta_nw_n$)
wynika dla pewnego $\chi>1$ zależnego wyłącznie od $d$
\begin{equation}\label{eq:pde17-sobstep}
 \left(\int_{B_{r_n}}|\eta_nw_n|^{2\chi}\,\mathrm d\lambda_d
 \right)^{1/\chi}\le C4^nY_n.
\end{equation}
Dla $d>2$ bierzemy $\chi=d/(d-2)$, a dla $d=2$
możemy wybrać $\chi=2$ (zanurzenie $H_0^1(B_2)\hookrightarrow L^4$).
Na zbiorze $E_{n+1}=B_{r_{n+1}}\cap\{u>k_{n+1}\}$ mamy
$w_n\ge k_{n+1}-k_n=K2^{-n-2}$, zatem
$|E_{n+1}|\le C4^nK^{-2}Y_n$. Nierówność Höldera i
\eqref{eq:pde17-sobstep} dają
\[
 Y_{n+1}\le
 \left(\int_{B_{r_n}}|\eta_nw_n|^{2\chi}\,\mathrm d\lambda_d\right)^{1/\chi}
 |E_{n+1}|^{1-1/\chi}
 \le Cb^nK^{-2\delta}Y_n^{1+\delta},
 \quad\delta=1-1/\chi>0,
\]
przy pewnym $b=b(d)>1$. Po podzieleniu przez $K^2$
otrzymujemy rekurencję dla $Z_n=Y_n/K^2$ postaci
z lematu~\ref{lem:pde17-iteration}. Wystarczy wziąć
$K=C_0\|u_+\|_{L^2(B_2)}$ z odpowiednio dużą stałą $C_0$.
Wtedy $Z_0$ jest dostatecznie małe. Otrzymujemy
$Y_n\to0$, czyli $(u-K)_+=0$ prawie wszędzie na $B_1$.
Po powrocie do skali $R$ i podzieleniu przez objętość kuli
otrzymujemy~\eqref{eq:pde17-l2linfty}.
\end{proof}

\begin{wniosek}[Lokalna ograniczoność rozwiązania]
Jeżeli $u$ jest rozwiązaniem słabym w $B_{2R}$, to
$u$ i $-u$ są podrozwiązaniami. Stąd
$\|u\|_{L^\infty(B_R)}\le C(\fint_{B_{2R}}|u|^2\,\mathrm d\lambda_d)^{1/2}$.
W szczególności każde rozwiązanie słabe ma lokalnie ograniczonego
reprezentanta w sensie prawie wszędzie.
\end{wniosek}

\section{Energia potęg i mechanizm iteracji Mosera}
\label{sec:pde17-moser}
W dalszej części zakładamy $v\ge\eps>0$. Dla nieujemnego
rozwiązania $u$ otrzymujemy to, biorąc $v=u+\eps$,
a po udowodnieniu oszacowań przechodzimy do granicy
$\eps\downarrow0$. Takie przesunięcie jest rozwiązaniem,
ponieważ równanie nie zawiera wyrazu zerowego rzędu.
Ponadto z lokalnej ograniczoności rozwiązania otrzymujemy
$v\in L^\infty_{\rm loc}$.\footnote{Jeżeli rozwiązanie jest
zdefiniowane tylko na kuli otwartej, obcięcie funkcji wartości
wykonuje się na mniejszej kuli zawierającej nośnik testu.
Wszystkie stałe w dowodach pozostają niezależne od $\eps$.}

\begin{lemat}[Oszacowanie dla potęg dodatnich i ujemnych]
\label{lem:pde17-powers}
Niech $v\in H^1(B_R)\cap L^\infty(B_R)$, $v\ge\eps>0$,
będzie słabym rozwiązaniem, $\beta\in(-\infty,0)\cup(1,\infty)$,
a $\eta\in C_c^1(B_R)$, $0\le\eta\le1$.
Wówczas
\begin{equation}\label{eq:pde17-powerenergy}
 \int_{B_R}|\nabla(\eta v^{\beta/2})|^2\,\mathrm d\lambda_d
 \le C\left(1+\frac{\beta^2}{(\beta-1)^2}\right)
 \int_{B_R}v^\beta|\nabla\eta|^2\,\mathrm d\lambda_d,
\end{equation}
gdzie $C$ zależy tylko od $\lambda,\Lambda$.
W szczególności stała jest ograniczona jednostajnie dla
$\beta\ge2$ lub $\beta<0$.
\end{lemat}
\begin{proof}
Testujemy równanie funkcją $\eta^2v^{\beta-1}$.
Jest to dopuszczalne, ponieważ $\eps\le v\le M$ na nośniku
$\eta$; mapy potęgowe mają tam ograniczone pochodne.
Różniczkując test i przenosząc składnik mieszany, uzyskujemy
\[
 |\beta-1|\lambda\int_{B_R}\eta^2v^{\beta-2}|\nabla v|^2
 \,\mathrm d\lambda_d
 \le2\Lambda\int_{B_R}\eta v^{\beta-1}|\nabla v||\nabla\eta|
 \,\mathrm d\lambda_d.
\]
Z nierówności Younga wynika
$\int_{B_R}\eta^2v^{\beta-2}|\nabla v|^2\le
C(\beta-1)^{-2}\int_{B_R} v^\beta|\nabla\eta|^2$.
Ponieważ $|\nabla(v^{\beta/2})|^2
=(\beta^2/4)v^{\beta-2}|\nabla v|^2$,
zaś $|\nabla(\eta v^{\beta/2})|^2$ jest co najwyżej
$2\eta^2|\nabla(v^{\beta/2})|^2+2v^\beta|\nabla\eta|^2$,
otrzymujemy~\eqref{eq:pde17-powerenergy}.
\end{proof}

\begin{lemat}[Iteracja $L^q\to L^\infty$ i potęg odwrotnych]
\label{lem:pde17-moser-iteration}
Niech $B_{2R}\Subset\Omega$, $v>0$ będzie rozwiązaniem słabym,
a $0<r<s\le2R$. Dla każdego $q>0$ istnieje
$C=C(d,\lambda,\Lambda,q)$ takie, że
\begin{align}
 \operatorname*{ess\,sup}_{B_r}v
 &\le C(s-r)^{-d/q}\|v\|_{L^q(B_s)},\label{eq:pde17-moser-up}\\
 \operatorname*{ess\,sup}_{B_r}v^{-1}
 &\le C(s-r)^{-d/q}\|v^{-1}\|_{L^q(B_s)}.
 \label{eq:pde17-moser-inv}
\end{align}
Wymagana jest lokalna ograniczoność $v$ oraz $v\ge\eps>0$;
dla $u\ge0$ stosujemy nierówności do $v=u+\eps$.
\end{lemat}
\begin{proof}
Przy $d>2$ bierzemy $\chi=d/(d-2)$, a przy $d=2$
$\chi=2$. Nierówność Sobolewa i~\eqref{eq:pde17-powerenergy}
dla $|\beta|\ge2$ dają dla $\rho<\sigma$:
\[
 \left(\int_{B_\rho}v^{\beta\chi}\,\mathrm d\lambda_d\right)^{1/\chi}
 \le\frac{C}{(\sigma-\rho)^2}
 \int_{B_\sigma}v^\beta\,\mathrm d\lambda_d.
\]
W przypadku $\beta>0$ po wyciągnięciu pierwiastka
stopnia $\beta$ jest to nierówność z normy $L^\beta$
do normy $L^{\beta\chi}$. W przypadku $\beta<0$
ten sam wzór, po zamianie $v$ na $v^{-1}$, jest
nierównością z normy $L^{|\beta|}$ do $L^{|\beta|\chi}$
\emph{dla odwrotności}. Biorąc $\beta_k=2\chi^k$
(odpowiednio $\beta_k=-2\chi^k$) oraz malejące promienie
$\rho_k=r+(s-r)2^{-k}$, $\rho_0=s$, $\rho_k\downarrow r$,
stosujemy nierówność kolejno do $B_{\rho_{k+1}}\subset
B_{\rho_k}$. Po podniesieniu do potęgi $1/|\beta_k|$
i pomnożeniu otrzymanych oszacowań, wykorzystaniu
$\sum_{k\ge0}\chi^{-k}=\chi/(\chi-1)$, a następnie
przeliczeniu objętości przy skalowaniu $B_{2R}$ na $B_2$,
otrzymujemy
\[
 \|v^{\pm1}\|_{L^\infty(B_r)}
 \le C(s-r)^{-d/2}\|v^{\pm1}\|_{L^2(B_s)}.
\]
Iloczyn stałych zbiega: jego logarytm jest ograniczony przez
$\sum_{k\ge0}C(1+k)/\chi^k<\infty$.

Dla $0<q<2$ interpolujemy
$\|w\|_{L^2(B_\sigma)}\le
\|w\|_{L^\infty(B_\sigma)}^{1-q/2}
\|w\|_{L^q(B_\sigma)}^{q/2}$,
gdzie $w=v$ lub $w=v^{-1}$. Aby uzasadnić absorpcję,
ustalmy $r_j=r+(s-r)(1-2^{-j})$, $j\ge0$, i
$M_j=\|w\|_{L^\infty(B_{r_j})}$ oraz
$Q=\|w\|_{L^q(B_s)}$. Stosując oszacowanie $L^2\to
L^\infty$ pomiędzy $B_{r_j}$ a $B_{r_{j+1}}$
otrzymujemy
\[
 M_j\le C2^{jd/2}(s-r)^{-d/2}M_{j+1}^{1-q/2}Q^{q/2}.
\]
Dla dowolnego $0<\delta<1$ nierówność Younga z wykładnikami
$2/(2-q)$ i $2/q$ daje
\[
 M_j\le\delta M_{j+1}
 +C_{d,q,\delta}2^{jd/q}(s-r)^{-d/q}Q.
\]
Wybieramy $\delta$ tak małe, aby $\delta2^{d/q}<1$,
i iterujemy. Sumujemy szereg geometryczny
$\sum_{j\ge0}(\delta2^{d/q})^j$. Pozostałość
$\delta^nM_n$ dąży do zera: dzięki lokalnej ograniczoności
$w$ jest ograniczona na $B_s\Subset\Omega$. Uzyskujemy
$M_0\le C(s-r)^{-d/q}Q$, czyli
\eqref{eq:pde17-moser-up}--\eqref{eq:pde17-moser-inv}.
Przypadek $q\ge2$ uzyskujemy, rozpoczynając tę samą iterację
bezpośrednio od $\beta_0=\pm q$, co daje właściwy wykładnik
$(s-r)^{-d/q}$.
\end{proof}


\begin{uwaga}[Jawny zakres stosowania potęg ujemnych]
\label{rem:pde17-negpower}
W oszacowaniu~\eqref{eq:pde17-moser-inv} nie wolno od razu
podstawić $u^{-1}$, jeśli nie wykazano dodatniej dolnej granicy
rozwiązania. Dla $u\ge0$ pracujemy z $v=u+\eps$;
wówczas $v^{-1}\in L^\infty$ na każdej kuli, a reguła
łańcuchowa jest dopuszczalna. Najpierw dowodzimy obu
oszacowań dla każdego $\eps>0$, przy tych samych stałych,
a dopiero na końcu przechodzimy do granicy monotonnej.
Nie otrzymujemy w ten sposób osobnego oszacowania
$\|u^{-1}\|_{L^q}$, które mogłoby być nieskończone.
To właśnie połączenie oszacowania dla $v^{-1}$
z mostem~\eqref{eq:pde17-bridge} usuwa tę trudność.
\end{uwaga}

\section{Logarytm rozwiązania i wykładnicza całkowalność}
\label{sec:pde17-log}
\begin{lemat}[Energia logarytmu]
\label{lem:pde17-log}
Niech $v\ge\eps>0$ będzie słabym dodatnim nadrozwiązaniem
w $B_{2r}(z)$. Wówczas
\begin{equation}\label{eq:pde17-log}
 \int_{B_r(z)}|\nabla\log v|^2\,\mathrm d\lambda_d
 \le C r^{d-2},\qquad C=C(d,\Lambda/\lambda).
\end{equation}
Dla każdej kuli $B_s(y)\subset B_r(z)$ otrzymujemy w konsekwencji
$\fint_{B_s(y)}|\log v-(\log v)_{B_s(y)}|\,\mathrm d\lambda_d
\le C(d,\Lambda/\lambda)$, o ile równanie obowiązuje
na kuli o podwójnym promieniu.
\end{lemat}
\begin{proof}
Testujemy nierówność nadrozwiązania nieujemną funkcją
$\varphi=\eta^2v^{-1}$. Uzyskujemy
\[
 \int_{B_{2r}(z)}\eta^2v^{-2}A\nabla v\cdot\nabla v\,\mathrm d\lambda_d
 \le 2\int_{B_{2r}(z)}\eta v^{-1}A\nabla v\cdot\nabla\eta\,\mathrm d\lambda_d.
\]
Eliptyczność i Young dają
$\int_{B_{2r}(z)}\eta^2|\nabla\log v|^2\le
C\int_{B_{2r}(z)}|\nabla\eta|^2$.
Wybieramy $\eta=1$ na $B_r(z)$ z nośnikiem w $B_{2r}(z)$,
$|\nabla\eta|\le C/r$. Wniosek o oscylacji średniej jest
konsekwencją nierówności Poincarégo
$\fint_B|g-g_B|\le Cs(\fint_B|\nabla g|^2)^{1/2}$
przy zastosowaniu oszacowania do odpowiedniej podwójnej kuli.
\end{proof}

Potrzebujemy teraz jednego lemmatu analizy harmonicznej.
Nie traktujemy go jako automatycznej własności funkcji $H^1$:
przydatna jest \emph{jednostajna} kontrola średnich oscylacji
na wszystkich podkulach.

\begin{lemat}[Lemat Johna--Nirenberga; lokalna postać]
\label{lem:pde17-jn}
Niech $B=B_R(z)\subset\R^d$, a
$g\in L^1(B_{3R/2}(z))$. Zakładamy, że
$\sup_{B'\subset B_{3R/2}(z)}\fint_{B'}
|g-g_{B'}|\,\mathrm d\lambda_d\le M$,
gdzie supremum biegnie po kulach zawartych w większej kuli.
Wtedy istnieją $c_d,C_d>0$, dla których
\begin{equation}\label{eq:pde17-jn}
 \fint_B\exp\left(\frac{c_d}{M}|g-g_B|\right)
 \,\mathrm d\lambda_d\le C_d
\end{equation}
(gdy $M=0$, funkcja jest stała prawie wszędzie).
\end{lemat}
\begin{proof}[Dowód metodą kostek zatrzymania]
Najpierw ustalmy kostkę $Q$ i załóżmy, że średnia oscylacja
graniczy się przez $M$ na wszystkich jej podkostkach.
Wybieramy $L=2^{d+2}M$. Maksymalne kostki dyadyczne
$Q_j\subset Q$, dla których
$\fint_{Q_j}|g-g_Q|>L$, są rozłączne. Ponieważ
$\int_Q|g-g_Q|\le M|Q|$, łączna miara tych kostek
jest nie większa od $(M/L)|Q|\le 2^{-d-2}|Q|$.
Maksymalność i porównanie z rodzicem dają
$|g_{Q_j}-g_Q|\le2^dL$. Poza ich sumą zachodzi
$|g-g_Q|\le L$ prawie wszędzie, dzięki różniczkowaniu
całek Lebesgue'a. W każdej złej kostce powtarzamy procedurę
z $g_{Q_j}$ zamiast $g_Q$. Po $k$ pokoleniach łączna
miara złych kostek nie przekracza $2^{-k(d+2)}|Q|$,
a na dopełnieniu mamy
$|g-g_Q|\le(1+k2^d)L$ prawie wszędzie.
Stąd $|\{|g-g_Q|>t\}|\le C|Q|e^{-ct/M}$.
Całkując funkcję rozkładu wykładniczego otrzymujemy
$\fint_Q e^{c'|g-g_Q|/M}\le C'$ dla dostatecznie małego
$c'>0$.
Przejście do kuli wymaga jeszcze uzasadnienia. Ustalmy
$B=B_R(z)$ i wybierzmy siatkę kostek o boku
$\ell=R/(100\sqrt d)$. Można wybrać skończenie wiele
jednakowych otwartych kostek $Q_1,\ldots,Q_N$ o środkach
na siatce z krokiem $\ell/2$, które pokrywają
$\overline B$, mają nakładające się wnętrza wzdłuż
łańcuchów sąsiedztwa, a ich trzykrotne powiększenia
zawierają się w $B_{3R/2}(z)$. Liczba $N$ zależy tylko od
$d$. (Można także użyć innej, stałej proporcji
$\ell/R$, zależnej wyłącznie od $d$.)

Na każdej podkostce $Q'\subset Q_j$ średnia oscylacja
$g$ nie przekracza $C_dM$. Rzeczywiście, otaczamy $Q'$
kulą $B'$ o promieniu co najwyżej $\sqrt d$ razy
bok $Q'$, nadal zawartą w $B_{3R/2}(z)$. Z nierówności
$\fint_{Q'}|g-g_{Q'}|\le
2\fint_{Q'}|g-g_{B'}|\le
2(|B'|/|Q'|)\fint_{B'}|g-g_{B'}|$
dostajemy żądane ograniczenie, ponieważ iloraz objętości
zależy tylko od $d$. Dyadyczny dowód dla kostek
zastosowany do każdego $Q_j$ zapewnia zatem
\[
 \fint_{Q_j}\exp\bigl(c_d|g-g_{Q_j}|/M\bigr)
 \,\mathrm d\lambda_d\le C_d.
\]
Jeżeli sąsiednie kostki $Q_i,Q_j$ mają część wspólną
o objętości co najmniej $c_d|Q_i|$, to
$|g_{Q_i}-g_{Q_j}|\le C_dM$: wystarczy całkować
$|g_{Q_i}-g_{Q_j}|\le
|g-g_{Q_i}|+|g-g_{Q_j}|$ po ich przecięciu,
a oba składniki oszacować przez oscylacje na kostkach.
Każdą parę kostek pokrycia łączy łańcuch co najwyżej
$N-1$ takich sąsiadów, stąd
$|g_{Q_j}-g_{Q_1}|\le C_dM$ dla wszystkich $j$.
Sumując wykładnicze oszacowania po pokryciu $B$ i
zmniejszając w razie potrzeby stałą $c_d$, otrzymujemy
$\fint_Be^{c_d|g-g_{Q_1}|/M}\le C_d$.
Z nierówności Jensena lub bezpośrednio z oszacowania
pierwszego momentu wynika
$|g_B-g_{Q_1}|\le C_dM$. Zastępując $g_{Q_1}$
przez $g_B$, kończymy dowód~\eqref{eq:pde17-jn}.
\end{proof}

\begin{wniosek}[Most między potęgami dodatnimi i ujemnymi]
\label{cor:pde17-bridge}
Jeżeli $v\ge\eps>0$ jest nadrozwiązaniem w $B_{4R}(x_0)$,
to istnieje $q_0=q_0(d,\Lambda/\lambda)>0$, dla którego
\begin{equation}\label{eq:pde17-bridge}
 \left(\fint_{B_{2R}}v^{q_0}\,\mathrm d\lambda_d\right)
 \left(\fint_{B_{2R}}v^{-q_0}\,\mathrm d\lambda_d\right)
 \le C(d,\Lambda/\lambda).
\end{equation}
\end{wniosek}
\begin{proof}
Bierzemy $g=\log v$. Z lematu~\ref{lem:pde17-log}
i nierówności Poincarégo wynika, że średnia oscylacja $g$
na każdej podkuli $B_{2R}$ jest jednostajnie ograniczona.
Uściślijmy użycie kul. Dla dowolnej
$B_\rho(y)\subset B_{3R}(x_0)$ i $\rho\le R$
kula $B_{2\rho}(y)$ mieści się w $B_{4R}(x_0)$,
więc testowanie logarytmu na podwójnej kuli daje
jednostajną kontrolę oscylacji. Jeśli $R<\rho\le3R$, kula $B_\rho(y)$ mieści się w
$B_{3R}(x_0)$. Oszacowanie energii na $B_{7R/2}(x_0)$
(z odcięciem wspartym w $B_{4R}(x_0)$) daje
$\int_{B_{3R}(x_0)}|\nabla g|^2\,\mathrm d\lambda_d(x)
\le C R^{d-2}$.
Nierówność Poincarégo na $B_\rho(y)$ i $\rho>R$
prowadzą zatem do
$\frac{1}{\lambda_d(B_\rho(y))}\int_{B_\rho(y)}|g(x)-g_{B_\rho(y)}|\,\mathrm d\lambda_d(x)
\le C\rho(R^{d-2}/\rho^d)^{1/2}\le C_d$.
W obu przypadkach
otrzymujemy ograniczenie średnich oscylacji zależne
tylko od $d$ i $\Lambda/\lambda$. Stosujemy zatem
zlokalizowaną wersję lematu Johna--Nirenberga z kulą
mniejszą $B_{2R}(x_0)$ i większą $B_{3R}(x_0)$.
W lemacie~\ref{lem:pde17-jn} dobieramy $q_0=c_d/M$.
Zachodzą tożsamości
$v^{\pm q_0}=\exp(\pm q_0(g-g_{B_{2R}}))
\exp(\pm q_0g_{B_{2R}})$.
Mnożąc całki i używając~\eqref{eq:pde17-jn}, usuwamy
przeciwne czynniki zależne od średniej $g$.
\end{proof}

\section{Nierówność Harnacka i regularność Höldera}
\label{sec:pde17-harnack}
\begin{twierdzenie}[Nierówność Harnacka dla słabych rozwiązań]
\label{thm:pde17-harnack}
Niech $d\ge2$, $A$ spełnia~\eqref{eq:pde17-ellipticity},
$B_{4R}(x_0)\Subset\Omega$, zaś $u\in H^1(B_{4R})$
będzie nieujemnym słabym rozwiązaniem
$-\operatorname{div}(A\nabla u)=0$.
Wówczas istnieje stała $C=C(d,\Lambda/\lambda)$ taka, że
\begin{equation}\label{eq:pde17-harnack}
 \operatorname*{ess\,sup}_{B_R(x_0)}u
 \le C\operatorname*{ess\,inf}_{B_R(x_0)}u.
\end{equation}
\end{twierdzenie}
\begin{proof}
Do $v=u+\eps$ stosujemy kolejno iterację
\eqref{eq:pde17-moser-up} dla potęgi dodatniej i
\eqref{eq:pde17-moser-inv} dla potęgi odwrotnej,
z wykładnikiem $q_0$ z~\eqref{eq:pde17-bridge}:
\[
 \operatorname*{ess\,sup}_{B_R}v
 \le C\left(\fint_{B_{2R}}v^{q_0}\,\mathrm d\lambda_d\right)^{1/q_0},
 \qquad
 \left(\fint_{B_{2R}}v^{-q_0}\,\mathrm d\lambda_d\right)^{-1/q_0}
 \le C\operatorname*{ess\,inf}_{B_R}v.
\]
Nierówność~\eqref{eq:pde17-bridge} ogranicza pierwszą średnią
przez odwrotność drugiej. Łącząc trzy oszacowania, otrzymujemy
$\operatorname*{ess\,sup}_{B_R}v\le
C\operatorname*{ess\,inf}_{B_R}v$ ze stałą niezależną od
$\eps$. Przejście $\eps\downarrow0$ kończy dowód.
\end{proof}

\begin{twierdzenie}[Twierdzenie De Giorgiego--Nasha--Mosera o regularności lokalnej]
\label{thm:pde17-holder}
Niech $d\ge2$, $\Omega\subset\R^d$ będzie obszarem,
macierz $A$ spełnia~\eqref{eq:pde17-ellipticity}, a
$u\in H^1_{\rm loc}(\Omega)$ niech będzie słabym rozwiązaniem
$-\operatorname{div}(A\nabla u)=0$.
Istnieje $\alpha=\alpha(d,\Lambda/\lambda)\in(0,1)$
taka, że $u$ ma reprezentanta należącego do $C^{0,\alpha}_{\rm loc}(\Omega)$.
Dla $B_{4R}(x_0)\Subset\Omega$ zachodzi oszacowanie
\begin{equation}\label{eq:pde17-holder}
 [u]_{C^{0,\alpha}(\overline{B_R(x_0)})}
 \le C R^{-\alpha}
 \left(\fint_{B_{4R}(x_0)}|u|^2\,\mathrm d\lambda_d\right)^{1/2},
\end{equation}
z $C$ zależną tylko od $d$ i $\Lambda/\lambda$.
\end{twierdzenie}
\begin{proof}
Najpierw uzasadnimy redukcję oscylacji. Oznaczmy
$m=\operatorname*{ess\,inf}_{B_{4r}}u$,
$M=\operatorname*{ess\,sup}_{B_{4r}}u$, dla kuli
$B_{4r}\Subset\Omega$. Z lokalnego ograniczenia
oba wyrazy są skończone. Funkcje $u-m$ i $M-u$ są
nieujemnymi rozwiązaniami w $B_{4r}$, a po zastosowaniu
nierówności Harnacka na $B_r$ (na kulach pośrednich
wewnątrz $B_{4r}$) otrzymujemy, oznaczając
$m_r=\operatorname*{ess\,inf}_{B_r}u$,
$M_r=\operatorname*{ess\,sup}_{B_r}u$,
\[
 M_r-m\le C(m_r-m),\qquad M-m_r\le C(M-M_r).
\]
Jeżeli $a=m_r-m$, $b=M-M_r$, $c=M_r-m_r$, to
$c\le(C-1)a$ i $c\le(C-1)b$. Ponieważ
$M-m=a+b+c$, mamy
\begin{equation}\label{eq:pde17-osc-decay}
 \operatorname*{ess\,osc}_{B_r}u
 \le\vartheta\operatorname*{ess\,osc}_{B_{4r}}u,
 \qquad \vartheta=\frac{C-1}{C+1}<1.
\end{equation}
Iterując dla promieni $4^{-k}R$ otrzymujemy
$\operatorname*{ess\,osc}_{B_\rho(x_0)}u
\le C(\rho/R)^\alpha
\operatorname*{ess\,osc}_{B_R(x_0)}u$ z
$\alpha=-\log_4\vartheta>0$. Stałą Harnacka możemy przyjąć ściśle większą niż $1$,
więc $0<\vartheta<1$; w razie potrzeby zmniejszamy wykładnik
do liczby z $(0,1)$.
Lokalna ograniczoność z~\eqref{eq:pde17-l2linfty}
kontroluje oscylację na kuli początkowej przez średnią $L^2$
na większej kuli.

Dla dowolnych $x,y\in B_R(x_0)$ bierzemy kulę o środku
w $x$ i promieniu porównywalnym z $|x-y|$, zawartą w
$B_{2R}(x_0)$. Z powyższego oszacowania oscylacji i
jednostajnej kontroli norm otrzymujemy
$|u(x)-u(y)|\le C R^{-\alpha}|x-y|^\alpha
(\fint_{B_{4R}}|u|^2)^{1/2}$ dla punktów Lebesgue'a.
Wartości na tych punktach przedłużają się jednoznacznie
na cały $B_R$ przez ciągłość, tworząc reprezentanta
Hölderowskiego. Tak otrzymujemy~\eqref{eq:pde17-holder}.
\end{proof}

\begin{wniosek}[Mocna zasada maksimum przy mierzalnych współczynnikach]
Jeżeli $u$ jest ciągłym reprezentantem słabego rozwiązania
w spójnym obszarze $\Omega$ i osiąga wewnętrzne maksimum
lub minimum, to jest stałe w $\Omega$.
\end{wniosek}
\begin{proof}
Dla maksimum $M$ funkcja $M-u\ge0$ w dostatecznie małej
kuli i przyjmuje wartość zero w środku. Z ciągłości i
nierówności Harnacka wynika, że znika na mniejszej kuli.
Zbiór punktów, w których $u=M$, jest zatem jednocześnie
otwarty i domknięty w $\Omega$. Spójność daje tezę;
minimum rozpatrujemy analogicznie.
\end{proof}

\section[Współczynniki skokowe i nieciągłość pochodnych]{Współczynniki skokowe: dlaczego nie otrzymujemy $C^1$?}
\label{sec:pde17-counterexample}
\begin{kontrprzyklad}[Brak różniczkowalności pomimo jednostajnej eliptyczności]
Dla $x\in(-1,1)$ przyjmijmy $a(x)=1$ gdy $x<0$ i
$a(x)=2$ gdy $x\ge0$. Funkcja
$u(x)=x$ dla $x<0$ oraz $u(x)=x/2$ dla $x\ge0$ należy do
$H^1(-1,1)\cap C^{0,1}([-1,1])$, ale nie do $C^1(-1,1)$.
Ponieważ $a(x)u'(x)=1$ prawie wszędzie, mamy
$-(au')'=0$ w sensie dystrybucji. W wymiarach $d\ge2$
można wziąć $A(x)=\operatorname{diag}(a(x_1),1,\ldots,1)$
i $u(x)=u(x_1)$.
\end{kontrprzyklad}

\begin{kontrprzyklad}[Ostrość wykładnika Höldera dla macierzy mierzalnych]
\label{ex:pde17-radial}
Niech $0<\gamma<1$ i $\Omega=B_1(0)\subset\R^2$. Dla $x\ne0$
niech $e_r=x/|x|$, $e_\theta=(-x_2,x_1)/|x|$ i
\[
 A(x)=e_r\otimes e_r+\gamma^2 e_\theta\otimes e_\theta,
 \qquad u(r,\theta)=r^\gamma\cos\theta.
\]
Definiujemy $A(0)=I$. Wtedy $\gamma^2I\le A\le I$,
$u\in H^1(B_1)$, a $u$ jest słabym rozwiązaniem
$-\operatorname{div}(A\nabla u)=0$ na całej kuli.
Rzeczywiście, dla $r>0$ rachunek biegunowy daje
\[
 \operatorname{div}(A\nabla u)=
 \frac1r\partial_r(r\partial_ru)
 +\frac{\gamma^2}{r^2}\partial_{\theta\theta}u
 =\gamma^2r^{\gamma-2}\cos\theta
 -\gamma^2r^{\gamma-2}\cos\theta=0.
\]
W pobliżu zera $|\nabla u|\le Cr^{\gamma-1}$, więc
$\int_{B_1}|\nabla u|^2\,\mathrm d\lambda_2<\infty$.
Po scałkowaniu przez części poza $B_\eps$ wyraz brzegowy
ma moduł nie większy niż $C\eps^\gamma\|\varphi\|_\infty$,
a zatem znika w granicy. Funkcja $u$ ma regularność
$C^{0,\gamma}$, ale nie należy do $C^{0,\beta}$ w żadnym
otoczeniu zera dla $\beta>\gamma$, ponieważ
$u(r,0)-u(0)=r^\gamma$.
\end{kontrprzyklad}

\section{Granice teorii i dalsze kierunki}
\label{sec:pde17-final}
Dowiedliśmy nierówności Harnacka oraz regularności Höldera
dla jednorodnych, \emph{skalarnych} równań dywergencyjnych
przy jedynie mierzalnych i jednostajnie eliptycznych
współczynnikach. Kluczowy mechanizm to oszacowanie energii,
zanurzenie Sobolewa, iteracja i kontrola $\log u$.
Dla $-\operatorname{div}(A\nabla u)=f$ prawa strona wnosi
nowe człony: typowym warunkiem regularności jest
$f\in L^q_{\rm loc}$ z $q>d/2$ (dla odpowiednio dobranej
klasy rozwiązań). W twierdzeniach dla układów równań
sytuacja może być zasadniczo gorsza. Rozdział kolejny
zajmie się teorią spektralną, ale osobno zachowamy miejsce
na pełne warianty niejednorodne, paraboliczne i kontrprzykłady.

% -----------------------------------------------------------------------
% Rozdział 18 — do włączenia PRZED \appendix w pliku wykładu
% -----------------------------------------------------------------------
\chapter[Widmo operatorów eliptycznych]{Widmo operatorów eliptycznych: zasada min--max, asymptotyka Weyla i oszacowania wartości własnych}
\label{chap:pde18}

\section{Cel i zakres rozdziału}
\label{sec:pde18-intro}
W rozdziale~\ref{chap:pde12} skonstruowaliśmy laplasjan Dirichleta
jako operator związany z domkniętą formą i uzyskaliśmy bazę
funkcji własnych. Teraz chcemy odpowiedzieć na pytania ilościowe:
jak duże są kolejne wartości własne, jak zależą od geometrii
obszaru i współczynników oraz co dokładnie oznacza asymptotyka
liczby modów o energii poniżej ustalonego progu? Podamy dowody
zasady min--max, oszacowania Li--Yau, prawa Weyla dla laplasjanu
Dirichleta i wzoru na wiodący człon śladu półgrupy cieplnej.

W całym rozdziale $d\ge1$, a $\Omega\subset\R^d$ oznacza
\emph{ograniczony obszar}. Zwartość zanurzenia
$H_0^1(\Omega)\hookrightarrow L^2(\Omega)$ wynika już z
ograniczoności obszaru i nie wymaga regularnego brzegu.
Silniejsze założenia geometryczne stawiamy jawnie tam,
gdzie potrzebujemy śladów lub twierdzeń o regularności.
Dla prawa Weyla wystarcza ograniczoność i warunek
$\lambda_d(\partial\Omega)=0$; w dowodzie zostanie
to wskazane osobno.

\begin{uwaga}[Od widma do ewolucji — zakres zastosowania]
\label{rem:pde-v43-spectral-heat}
Rozkład spektralny dotyczy operatora samosprzężonego,
niekoniecznie formalnego wyrażenia $-\operatorname{div}(A\nabla)+V$
na dowolnie przyjętej dziedzinie. W rozdziale~\ref{chap:pde20}
półgrupa $e^{-tL}$ i jej rozwinięcie w funkcje własne
są używane w przypadku \emph{autonomicznego}, dodatniego
operatora Dirichleta. Dla macierzy $A(t,x)$ zależnej od
czasu, zwłaszcza niesymetrycznej, nie wolno stosować
tego samego rozwinięcia do operatora chwilowego;
konstrukcja Galerkina wykorzysta wówczas bazę
\emph{pomocniczego} laplasjanu Dirichleta.
\end{uwaga}

\section{Od formy do operatora z dyskretnym widmem}
\label{sec:pde18-form}
\begin{definicja}[Symetryczna forma eliptyczna]
Niech $\Omega\subset\R^d$ będzie ograniczonym obszarem
Lipschitzowskim. Niech $A\in L^\infty(\Omega;\R^{d\times d})$
będzie symetryczna prawie wszędzie, a $V\in L^\infty(\Omega;\R)$.
Zakładamy, że istnieją $0<a_0\le a_1<\infty$ takie, iż
$ a_0|\xi|^2\le\xi^{\mathsf T}A(x)\xi\le a_1|\xi|^2$
dla prawie każdego $x\in\Omega$ i każdego $\xi\in\R^d$.
Na $H_0^1(\Omega)$ określamy formę
\begin{equation}\label{eq:pde18-form}
 q(u,v)=\int_\Omega \bigl(A\nabla u\cdot\nabla v+Vuv\bigr)
 \,\mathrm d\lambda_d(x),\qquad q[u]=q(u,u).
\end{equation}
Przez $L_{A,V}$ oznaczamy operator samosprzężony w
$L^2(\Omega)$ związany z tą formą: $u\in D(L_{A,V})$
wtedy i tylko wtedy, gdy $u\in H_0^1(\Omega)$ oraz istnieje
$f\in L^2(\Omega)$ takie, że $q(u,v)=(f,v)_{L^2}$
dla każdego $v\in H_0^1(\Omega)$; wówczas $L_{A,V}u=f$.
\end{definicja}

\begin{uwaga}[Istotność dziedziny operatora]
Symbol $-\operatorname{div}(A\nabla u)+Vu$ należy na początku
rozumieć dystrybucyjnie. Przy $A\in L^\infty$
nie można bez dowodu przyjąć $D(L_{A,V})=H^2(\Omega)\cap
H_0^1(\Omega)$. Definicja przez formę jest poprawna również dla
współczynników skokowych. W tym rozdziale \emph{obszar} to
otwarty zbiór spójny, a \emph{dziedzina operatora} to $D(L_{A,V})$.
\end{uwaga}

\begin{twierdzenie}[Samosprzężoność i zwarta rezolwenta]
\label{thm:pde18-compact}
Przy założeniach poprzedniej definicji operator $L_{A,V}$
jest samosprzężony i ograniczony z dołu. Jego widmo składa się
z rzeczywistych wartości własnych, liczonych z krotnościami,
$\lambda_1\le\lambda_2\le\cdots$, przy czym
$\lambda_k\to+\infty$. Istnieje ortonormalna baza
$L^2(\Omega)$ złożona z funkcji własnych $L_{A,V}$.
\end{twierdzenie}
\begin{proof}
Oznaczmy $M=\|V_-\|_{L^\infty}$, gdzie
$V_-=\max\{-V,0\}$. Przy $c=M+1$ forma
$q_c(u,v)=q(u,v)+c(u,v)_{L^2}$ spełnia
$q_c[u]\ge a_0\|\nabla u\|_2^2+\|u\|_2^2$.
Jest zatem koercywna na $H_0^1(\Omega)$, a norma
$q_c[u]^{1/2}$ jest równoważna normie $H_0^1$.
Dla każdego $f\in L^2$ twierdzenie Laxa--Milgrama
wyznacza jednoznacznie $u=T f\in H_0^1$ takie, że
$q_c(u,v)=(f,v)_{L^2}$ dla wszystkich $v\in H_0^1$.
Otrzymujemy ciągły operator $T:L^2\to H_0^1\hookrightarrow L^2$.
Ponieważ zanurzenie $H_0^1(\Omega)\hookrightarrow L^2(\Omega)$
jest zwarte, operator $T$ działający w $L^2$ jest zwarty.

Z symetrii formy wynika
$(Tf,g)_{L^2}=q_c(Tf,Tg)=(f,Tg)_{L^2}$;
ponadto $(Tf,f)_{L^2}=q_c[Tf]\ge0$ i $\ker T=\{0\}$.
Twierdzenie spektralne dla zwartych operatorów samosprzężonych
(rozdział~\ref{chap:pde12}) daje bazę ortonormalną
$(\varphi_k)$ z $T\varphi_k=\mu_k\varphi_k$,
$\mu_k>0$, $\mu_k\downarrow0$. Definiując
$L=T^{-1}-cI$ na $\operatorname{ran}T$,
otrzymujemy samosprzężony operator związany z formą oraz
$L\varphi_k=(\mu_k^{-1}-c)\varphi_k$.
Gdyby $\mu_k$ nie dążyły do zera, zwartość $T$ przeczyłaby
ortonormalności $(\varphi_k)$. Stąd $\lambda_k=\mu_k^{-1}-c$
dążą do $+\infty$ i tworzą całe widmo.
\end{proof}

\section{Zasada min--max i jej konsekwencje}
\label{sec:pde18-minmax}
\begin{twierdzenie}[Zasada min--max Couranta--Fischera, wersja formowa]
\label{thm:pde18-minmax}
Niech $L_{A,V}$ będzie operatorem z
Twierdzenia~\ref{thm:pde18-compact}, a $q$ jego formą.
Dla każdego $k\ge1$ zachodzi
\begin{equation}\label{eq:pde18-minmax}
 \lambda_k=\min_{\substack{E\subset H_0^1(\Omega)\\\dim E=k}}
 \ \max_{u\in E\setminus\{0\}}\frac{q[u]}{\|u\|_2^2}
 =\max_{\substack{F\subset L^2(\Omega)\\\dim F=k-1}}
 \ \inf_{\substack{u\in H_0^1(\Omega)\setminus\{0\}\\u\perp F}}
 \frac{q[u]}{\|u\|_2^2}.
\end{equation}
W drugim wyrażeniu dla $k=1$ przyjmujemy $F=\{0\}$.
\end{twierdzenie}
\begin{proof}
Niech $(\varphi_j)$ będzie ortonormalną bazą funkcji
własnych, a wartości własne niech będą uporządkowane rosnąco.
Dla $u=\sum_jc_j\varphi_j\in H_0^1(\Omega)$ zachodzi
$q[u]=\sum_j\lambda_j|c_j|^2$.
Równość wynika najpierw z tożsamości formowej na skończonych
sumach, a następnie z ich zbieżności w normie formy $q_c$:
$\sum_j(\lambda_j+c)|c_j|^2=q_c[u]<\infty$.

Jeśli $E$ ma wymiar $k$, odwzorowanie
$u\mapsto((u,\varphi_1),\ldots,(u,\varphi_{k-1}))$
z $E$ do $\R^{k-1}$ ma niezerowe jądro. Istnieje więc
$0\ne u\in E$ ortogonalne do $\varphi_1,\ldots,
\varphi_{k-1}$, a wtedy $q[u]\ge\lambda_k\|u\|_2^2$.
Przy $E_k=\operatorname{span}\{\varphi_1,\ldots,\varphi_k\}$
każdy niezerowy $u$ ma iloraz Rayleigha najwyżej
$\lambda_k$. To dowodzi pierwszej równości.

Dla dowolnego $F$ o wymiarze $k-1$ istnieje niezerowy
$u\in E_k\cap F^\perp$. Stąd
$\inf_{u\perp F}q[u]/\|u\|_2^2\le\lambda_k$.
Jeśli $F=\operatorname{span}\{\varphi_1,\ldots,
\varphi_{k-1}\}$, wspomniana suma spektralna daje
odwrotną nierówność. Otrzymujemy drugą równość.
\end{proof}

\begin{wniosek}[Porównanie współczynników i potencjałów]
\label{cor:pde18-compare}
Niech $L_1,L_2$ mają wspólną dziedzinę formową
$H_0^1(\Omega)$, a ich symetryczne formy spełniają
$q_1[u]\le q_2[u]$ dla każdego $u\in H_0^1(\Omega)$.
Wtedy $\lambda_k(L_1)\le\lambda_k(L_2)$ dla każdego $k$.
W szczególności, gdy $A$ spełnia warunki z
\ref{sec:pde18-form}, a $v_-\le V\le v_+$ prawie wszędzie,
to
\begin{equation}\label{eq:pde18-compare}
 a_0\lambda_k(-\Delta_D)+v_-
 \le\lambda_k(L_{A,V})\le
 a_1\lambda_k(-\Delta_D)+v_+.
\end{equation}
\end{wniosek}
\begin{proof}
Dla dowolnej $k$-wymiarowej podprzestrzeni $E$ maksimum
ilorazu Rayleigha $q_1$ nie przewyższa odpowiedniego maksimum
$q_2$. Bierzemy minimum po $E$ i stosujemy
\eqref{eq:pde18-minmax}. Nierówności
$a_0\|\nabla u\|_2^2+v_-\|u\|_2^2\le q[u]
\le a_1\|\nabla u\|_2^2+v_+\|u\|_2^2$
prowadzą do~\eqref{eq:pde18-compare}.
\end{proof}

\begin{wniosek}[Monotoniczność Dirichleta względem obszaru]
\label{cor:pde18-domain}
Jeżeli $\Omega_1\subset\Omega_2$ są ograniczonymi obszarami,
a $-\Delta_{D,\Omega_i}$ oznacza laplasjan Dirichleta
w $L^2(\Omega_i)$, to
$\lambda_k(-\Delta_{D,\Omega_2})\le
\lambda_k(-\Delta_{D,\Omega_1})$ dla każdego $k$.
\end{wniosek}
\begin{proof}
Przedłużenie przez zero zanurza izometrycznie
$H_0^1(\Omega_1)$ w $H_0^1(\Omega_2)$, zachowując
$\|u\|_2$ i $\|\nabla u\|_2$.
Minimum w~\eqref{eq:pde18-minmax} dla $\Omega_2$
biegnie zatem po większej rodzinie podprzestrzeni.
\end{proof}

\begin{przyklad}[Odcinek i prostokąt]
Na przedziale $(0,L)$ wartości własne laplasjanu Dirichleta
to $\lambda_n=(\pi n/L)^2$, a unormowane funkcje własne
$\varphi_n(x)=\sqrt{2/L}\sin(n\pi x/L)$.
Na prostokącie $\Omega=(0,L_1)\times(0,L_2)$ mamy
$\lambda_{m,n}=\pi^2(m^2/L_1^2+n^2/L_2^2)$,
$m,n\ge1$. Gdy $L_1=L_2$, wartości własne dla
$(m,n)$ i $(n,m)$ mogą być równe; na ogół widmo
nie jest więc proste.
\end{przyklad}


\begin{stwierdzenie}[Skalowanie widma przy dylatacji obszaru]
\label{prop:pde18-scale}
Niech $\Omega\subset\R^d$ będzie ograniczonym obszarem,
$a>0$, a $a\Omega=\{ax:x\in\Omega\}$.
Dla każdego $k\ge1$ zachodzi
$\lambda_k(-\Delta_{D,a\Omega})=a^{-2}\lambda_k(-\Delta_{D,\Omega})$.
W konsekwencji
$N_{D,a\Omega}(\Lambda)=N_{D,\Omega}(a^2\Lambda)$.
\end{stwierdzenie}
\begin{proof}
Dla $v\in H_0^1(\Omega)$ określamy
$Uv(x)=a^{-d/2}v(x/a)$ dla $x\in a\Omega$.
Zmiana zmiennych pokazuje, że $U:L^2(\Omega)\to
L^2(a\Omega)$ jest unitarne i przeprowadza $H_0^1(\Omega)$
na $H_0^1(a\Omega)$. Ponadto
$\|\nabla Uv\|_{L^2(a\Omega)}^2
=a^{-2}\|\nabla v\|_{L^2(\Omega)}^2$.
Zasada min--max zastosowana do $U(E)$ i $U^{-1}(E)$
daje pierwszą równość, a definicja funkcji zliczającej drugą.
\end{proof}

\section{Funkcja zliczająca i prawo Weyla}
\label{sec:pde18-weyl}
\begin{definicja}[Funkcja zliczająca]
Dla operatora z dyskretnym widmem określamy
$N_L(\Lambda)=\#\{k\in\N:\lambda_k(L)\le\Lambda\}$,
licząc krotności wartości własnych. Symbol $\omega_d$
oznacza objętość jednostkowej kuli w $\R^d$.
\end{definicja}

\begin{lemat}[Liczba modów na sześcianie]
\label{lem:pde18-cube}
Niech $Q=(0,h)^d$, $h>0$. Dla laplasjanu z warunkiem
Dirichleta albo Neumanna zachodzi
\begin{equation}\label{eq:pde18-cube}
 N_{D,Q}(\Lambda)=\frac{\omega_dh^d}{(2\pi)^d}
 \Lambda^{d/2}+O_{d,h}(1+\Lambda^{(d-1)/2}),
 \qquad \Lambda\to\infty,
\end{equation}
a taki sam człon główny ma $N_{N,Q}$.
\end{lemat}
\begin{proof}
Rozdzielenie zmiennych daje wartości własne
$\pi^2h^{-2}\sum_{j=1}^d n_j^2$, przy czym
$n\in\{1,2,\ldots\}^d$ dla Dirichleta i
$n\in\{0,1,\ldots\}^d$ dla Neumanna. Niech
$R=h\sqrt\Lambda/\pi$. Liczba takich $n$ to liczba
punktów kratowych w dodatnim ortancie kuli promienia $R$.
Każdemu punktowi kratowemu przyporządkowujemy jednostkowy
sześcian. Różnica między sumą objętości sześcianów i
objętością ortantu kuli znajduje się w warstwie o szerokości
co najwyżej $\sqrt d+1$ przy sferze i hiperpłaszczyznach
współrzędnych. Jej objętość jest
$O_d(1+R^{d-1})$. Objętość dodatniego ortantu wynosi
$2^{-d}\omega_dR^d$. Po podstawieniu $R=h\sqrt\Lambda/\pi$
uzyskujemy~\eqref{eq:pde18-cube}. Dopuszczenie zerowych
współrzędnych zmienia liczbę punktów o co najwyżej
$O_d(1+R^{d-1})$.
\end{proof}

\begin{lemat}[Rozdzielanie na sześciany i ograniczenie liczby modów]
\label{lem:pde18-bracket}
Niech $P$ będzie wnętrzem sumy skończenie wielu domkniętych
sześcianów $Q_1,\ldots,Q_m$ o parami rozłącznych wnętrzach,
pochodzących z tej samej siatki. Dla każdego $\Lambda\in\R$
\begin{equation}\label{eq:pde18-bracket}
 \sum_{j=1}^m N_{D,Q_j}(\Lambda)
 \le N_{D,P}(\Lambda)
 \le\sum_{j=1}^m N_{N,Q_j}(\Lambda).
\end{equation}
\end{lemat}
\begin{proof}
Funkcje z $\bigoplus_jH_0^1(Q_j)$, przedłużone przez zero,
należą do $H_0^1(P)$. Ich forma Dirichleta jest sumą form
z poszczególnych sześcianów, a normy $L^2$ również się sumują.
Zatem zasada min--max i uwzględnienie krotności dają
pierwszą nierówność. Z drugiej strony przyporządkowanie
$u\in H_0^1(P)$ rodzinie ograniczeń $(u|_{Q_j})_j$
zanurza izometrycznie jej przestrzeń formową w
$\bigoplus_jH^1(Q_j)$ (w normie $L^2$ oraz energii).
Rozszerzenie przestrzeni dopuszczalnych funkcji o niezależne
wartości na sztucznych granicach sześcianów może tylko
obniżyć wartości min--max. Daje to drugą nierówność,
ponieważ na każdej części odpowiada ona warunkowi Neumanna.
\end{proof}

\begin{twierdzenie}[Prawo Weyla dla laplasjanu Dirichleta]
\label{thm:pde18-weyl}
Niech $\Omega\subset\R^d$ będzie ograniczonym obszarem
o brzegu miary Lebesgue'a zero, tj.
$\lambda_d(\partial\Omega)=0$. Niech $-\Delta_D$
oznacza laplasjan Dirichleta w $L^2(\Omega)$.
Wówczas
\begin{equation}\label{eq:pde18-weyl}
 \lim_{\Lambda\to\infty}
 \frac{N_{-\Delta_D}(\Lambda)}{\Lambda^{d/2}}
 =\frac{\omega_d}{(2\pi)^d}|\Omega|.
\end{equation}
Równoważnie, gdy $k\to\infty$,
\begin{equation}\label{eq:pde18-lambda-weyl}
 \lambda_k(-\Delta_D)\sim
 (2\pi)^2\left(\frac{k}{\omega_d|\Omega|}\right)^{2/d}.
\end{equation}
\end{twierdzenie}
\begin{proof}
Ustalmy siatkę sześcianów o boku $h>0$ na dostatecznie
dużym sześcianie zawierającym $\overline\Omega$.
Przez $P_h^-$ oznaczmy wnętrze sumy wszystkich domkniętych
komórek zawartych w $\Omega$, a przez $P_h^+$ wnętrze sumy
wszystkich domkniętych komórek, których wnętrza przecinają
$\Omega$. Mamy $P_h^-\subset\Omega\subset P_h^+$.
Ostatnie zawieranie jest prawdziwe również dla punktów
$\Omega$ leżących na płaszczyznach siatki: ich małe otoczenia
należą do sumy komórek przecinających otwarty zbiór $\Omega$.

Z monotoniczności Dirichleta względem obszaru wynika
$N_{D,P_h^-}(\Lambda)\le N_{D,\Omega}(\Lambda)
\le N_{D,P_h^+}(\Lambda)$.
Stosujemy nierówności~\eqref{eq:pde18-bracket} do obu
wielościanów, a następnie~\eqref{eq:pde18-cube} do ich
skończenie wielu komórek. Dla każdego ustalonego $h$
otrzymujemy po podzieleniu przez $\Lambda^{d/2}$ i
przejściu $\Lambda\to\infty$:
\[
 \frac{\omega_d}{(2\pi)^d}|P_h^-|
 \le\liminf_{\Lambda\to\infty}\frac{N_{D,\Omega}(\Lambda)}{\Lambda^{d/2}}
 \le\limsup_{\Lambda\to\infty}\frac{N_{D,\Omega}(\Lambda)}{\Lambda^{d/2}}
 \le\frac{\omega_d}{(2\pi)^d}|P_h^+|.
\]
Każda komórka należąca do $P_h^+\setminus P_h^-$ jest
w odległości co najwyżej $C_dh$ od $\partial\Omega$.
Ponieważ $\partial\Omega$ jest zwartym zbiorem miary zero,
objętość jego $C_dh$-otoczenia dąży do zera wraz z $h\to0$.
Zatem $|P_h^-|\to|\Omega|$ oraz
$|P_h^+|\to|\Omega|$, co daje~\eqref{eq:pde18-weyl}.

Aby uzyskać~\eqref{eq:pde18-lambda-weyl}, oznaczmy
$C_\Omega=\omega_d|\Omega|/(2\pi)^d$. Dla każdego
$\eps>0$ i dużych $\Lambda$ mamy
$(C_\Omega-\eps)\Lambda^{d/2}\le N(\Lambda)
\le(C_\Omega+\eps)\Lambda^{d/2}$. Stosując te nierówności
do $\Lambda$ nieco mniejszych i nieco większych niż
$(k/C_\Omega)^{2/d}$, otrzymujemy ograniczenia na
$\lambda_k/(k/C_\Omega)^{2/d}$ zbieżne do $1$
przy $\eps\downarrow0$. Argument uwzględnia krotności,
gdyż korzysta wyłącznie z funkcji zliczającej.
\end{proof}

\begin{uwaga}[Co pozostaje trudne w ogólnej teorii Weyla?]
Powyższy dowód wykorzystuje wyłącznie geometrię siatki,
min--max i jawne widmo sześcianu. Dla operatorów z macierzą
$A(x)$ oczekiwany człon główny dla odpowiednio regularnych
współczynników to
$(2\pi)^{-d}\omega_d\Lambda^{d/2}
\int_\Omega(\det A(x))^{-1/2}\,\mathrm d\lambda_d(x)$.
Nie wyprowadzamy tu tej mocniejszej asymptotyki, ponieważ
wymaga ona dodatkowej lokalizacji w przestrzeni fazowej.
Natomiast porównanie~\eqref{eq:pde18-compare} natychmiast
daje ilościowe dwustronne ograniczenia $\lambda_k(L_{A,V})
\asymp k^{2/d}$ dla dostatecznie dużych $k$.
\end{uwaga}

\section{Nierówność Li--Yau: dolne oszacowanie sumy wartości własnych}
\label{sec:pde18-liyau}
\begin{lemat}[Zasada wypełniania najniższych częstotliwości]
\label{lem:pde18-bathtub}
Niech $F:\R^d\to[0,M]$ będzie mierzalna,
$0<\int_{\R^d}F\,\mathrm d\lambda_d=k<\infty$,
$M>0$, i niech $\rho=(k/(M\omega_d))^{1/d}$.
Wtedy
\begin{equation}\label{eq:pde18-bathtub}
 \int_{\R^d}|\xi|^2F(\xi)\,\mathrm d\lambda_d(\xi)
 \ge\frac{d}{d+2}k\rho^2.
\end{equation}
\end{lemat}
\begin{proof}
Oznaczmy $F_0=M\chi_{B_\rho(0)}$.
Na $B_\rho$ mamy $(|\xi|^2-\rho^2)(F-F_0)\ge0$,
ponieważ oba czynniki są niedodatnie; poza kulą oba są
nieujemne. Całkując i korzystając z równości
$\int_{\R^d}F\,\mathrm d\lambda_d=\int_{\R^d}F_0\,\mathrm d\lambda_d=k$, dostajemy
$\int_{\R^d}|\xi|^2F(\xi)\,\mathrm d\lambda_d(\xi)\ge\int_{\R^d}|\xi|^2F_0(\xi)\,\mathrm d\lambda_d(\xi)$.
Ostatnia całka wynosi
$M\,d\omega_d\rho^{d+2}/(d+2)
=\frac d{d+2}k\rho^2$, co dowodzi tezy.
\end{proof}

\begin{twierdzenie}[Nierówność Li--Yau dla laplasjanu Dirichleta]
\label{thm:pde18-liyau}
Niech $\Omega\subset\R^d$ będzie ograniczonym obszarem,
a $0<\lambda_1\le\lambda_2\le\cdots$ wartościami własnymi
laplasjanu Dirichleta. Dla każdego $k\ge1$ zachodzi
\begin{equation}\label{eq:pde18-liyau}
 \sum_{j=1}^k\lambda_j\ge
 \frac{d}{d+2}(2\pi)^2
 \left(\frac{1}{\omega_d|\Omega|}\right)^{2/d}
 k^{1+2/d}.
\end{equation}
W szczególności
$\lambda_k\ge\frac d{d+2}(2\pi)^2
(k/(\omega_d|\Omega|))^{2/d}$.
\end{twierdzenie}
\begin{proof}
Niech $\varphi_1,\ldots,\varphi_k$ będą ortonormalnymi
funkcjami własnymi w $L^2(\Omega)$ i przedłużmy je przez
zero do $\R^d$. Ponieważ należą do $H_0^1(\Omega)$,
ich przedłużenia należą do $H^1(\R^d)$. Używamy unitarnej
transformaty Fouriera
$\widehat f(\xi)=(2\pi)^{-d/2}
\int_{\R^d}f(x)e^{-ix\cdot\xi}\,\mathrm d\lambda_d(x)$.
Definiujemy $F(\xi)=\sum_{j=1}^k|\widehat\varphi_j(\xi)|^2$.

Z twierdzenia Plancherela wynika $\int_{\R^d}F(\xi)\,\mathrm d\lambda_d(\xi)=k$ oraz
\begin{equation}\label{eq:pde18-fft-energy}
 \int_{\R^d}|\xi|^2F(\xi)\,\mathrm d\lambda_d(\xi)
 =\sum_{j=1}^k\int_\Omega|\nabla\varphi_j|^2
 \,\mathrm d\lambda_d(x)=\sum_{j=1}^k\lambda_j.
\end{equation}
Dla każdego ustalonego $\xi$ funkcje
$\varphi_1,\ldots,\varphi_k$ są ortonormalne w $L^2(\Omega)$.
Zatem nierówność Bessela zastosowana do
$x\mapsto(2\pi)^{-d/2}e^{-ix\cdot\xi}$ daje
\[
 F(\xi)=\sum_{j=1}^k
 |(\varphi_j,(2\pi)^{-d/2}e^{-i\cdot\xi})_{L^2(\Omega)}|^2
 \le(2\pi)^{-d}|\Omega|=:M.
\]
Stosujemy Lemat~\ref{lem:pde18-bathtub} z tym $M$,
$\rho^d=k/(M\omega_d)=(2\pi)^dk/(\omega_d|\Omega|)$.
Wstawiając oszacowanie do~\eqref{eq:pde18-fft-energy},
otrzymujemy~\eqref{eq:pde18-liyau}. Ponieważ
$k\lambda_k\ge\sum_{j=1}^k\lambda_j$, otrzymujemy
podany dolny szacunek pojedynczej wartości własnej.
\end{proof}

\begin{uwaga}[Dlaczego współczynnik jest naturalny?]
Prawa strona~\eqref{eq:pde18-liyau} jest dokładnie całką
średniego kosztu $|\xi|^2$ przy równomiernym wypełnianiu
najniższych częstotliwości w przestrzeni fazowej. Jest to wariacyjna interpretacja problemu
optymalizacyjnego dla funkcji $F$ o zadanej masie i ograniczeniu
punktowym; dowód nie wymaga znajomości prawa Weyla.
\end{uwaga}


\begin{wniosek}[Oszacowanie funkcji zliczającej bez prawa Weyla]
\label{cor:pde18-count-liyau}
Niech $d\ge1$, a $\Omega\subset\R^d$ będzie ograniczonym
obszarem. Dla $\Lambda>0$ funkcja zliczająca laplasjanu
Dirichleta spełnia
\[
 N_{-\Delta_D}(\Lambda)\le
 \frac{\omega_d|\Omega|}{(2\pi)^d}
 \left(\frac{d+2}{d}\right)^{d/2}\Lambda^{d/2}.
\]
Wniosek nie wymaga założenia $\lambda_d(\partial\Omega)=0$.
\end{wniosek}
\begin{proof}
Niech $k=N_{-\Delta_D}(\Lambda)$. Gdy $k=0$, nie ma
czego dowodzić. W przeciwnym razie $\lambda_k\le\Lambda$.
Na mocy nierówności~\eqref{eq:pde18-liyau} mamy
$\Lambda\ge\lambda_k\ge\frac d{d+2}(2\pi)^2
(k/(\omega_d|\Omega|))^{2/d}$. Przekształcenie tego
oszacowania daje tezę. Zwróćmy uwagę, że jest to
\emph{nierówność dla każdej} $\Lambda>0$, a nie
asymptotyka z prawa Weyla.
\end{proof}

\section{Pierwsza wartość własna, dodatniość i prostota}
\label{sec:pde18-ground}
\begin{twierdzenie}[Prostota pierwszej wartości własnej laplasjanu]
\label{thm:pde18-simple}
Niech $\Omega\subset\R^d$ będzie ograniczonym, spójnym
obszarem klasy $C^{1,1}$. Dla laplasjanu Dirichleta
$\lambda_1>0$ jest prosta, a odpowiadającą jej funkcję
własną można wybrać ściśle dodatnią we wnętrzu $\Omega$.
\end{twierdzenie}
\begin{proof}
Z nierówności Poincarégo mamy
$\lambda_1=\min_{u\in H_0^1\setminus\{0\}}
\|\nabla u\|_2^2/\|u\|_2^2>0$. Niech $u$ będzie
rzeczywistą funkcją własną dla $\lambda_1$.
Ponieważ $|u|\in H_0^1(\Omega)$, a
$|\nabla|u||=|\nabla u|$ prawie wszędzie poza zerami,
$|u|$ osiąga ten sam minimalny iloraz Rayleigha. Jest więc
również funkcją własną. Wewnętrzna regularność eliptyczna
dla $-\Delta|u|=\lambda_1|u|$, wraz z bootstrapem Sobolewa,
daje ciągłego, a w istocie gładkiego reprezentanta.
Ponieważ $|u|\ge0$ oraz
$-\Delta|u|=\lambda_1|u|\ge0$, mocna zasada minimum
dla funkcji superharmonicznych zapewnia $|u|>0$ na $\Omega$.
Stąd $u$ nie może zmieniać znaku bez przyjęcia wartości zero,
więc każdy rzeczywisty wektor własny dla $\lambda_1$
ma stały, ściśle dodatni albo ujemny znak.

Przypuśćmy, że przestrzeń własna ma wymiar co najmniej dwa.
Można wtedy znaleźć niezerowe, rzeczywiste, ortogonalne
funkcje własne $u,v$ dla $\lambda_1$.
Po ewentualnej zmianie znaków obie byłyby dodatnie w $\Omega$,
co przeczyłoby $(u,v)_{L^2}=0$. Pierwsza wartość własna jest
więc prosta.
\end{proof}

\section{Warunek Neumanna i problem zerowej energii}
\label{sec:pde18-neumann}
\begin{twierdzenie}[Widmo Neumanna i pierwsza dodatnia wartość]
\label{thm:pde18-neumann}
Niech $\Omega\subset\R^d$ będzie ograniczonym, spójnym
obszarem Lipschitzowskim, a $-\Delta_N$ laplasjanem
związanym z formą $\int_\Omega|\nabla u|^2\,
\mathrm d\lambda_d$ o dziedzinie formowej $H^1(\Omega)$.
Jego widmo jest dyskretne i nieujemne:
$0=\mu_1<\mu_2\le\mu_3\le\cdots\to\infty$.
Pierwsza przestrzeń własna składa się z funkcji stałych.
Ponadto
\begin{equation}\label{eq:pde18-neumann}
 \mu_2=\min_{\substack{u\in H^1(\Omega)\setminus\{0\}\\
 \int_\Omega u\,\mathrm d\lambda_d=0}}
 \frac{\int_\Omega|\nabla u|^2\,\mathrm d\lambda_d}
 {\int_\Omega|u|^2\,\mathrm d\lambda_d}>0.
\end{equation}
\end{twierdzenie}
\begin{proof}
Zanurzenie $H^1(\Omega)\hookrightarrow L^2(\Omega)$ jest zwarte.
Forma $\int_\Omega(|\nabla u|^2+|u|^2)\,\mathrm d\lambda_d$ jest koercywna,
więc poprzednia konstrukcja rezolwenty działa po zastąpieniu
$H_0^1$ przez $H^1$. Funkcje stałe mają energię zero.
Jeżeli energia $u\in H^1(\Omega)$ znika, to $\nabla u=0$
prawie wszędzie; z połączenia spójności i własności Sobolewa
wynika, że $u$ jest stała prawie wszędzie.
Zatem zero jest prostą wartością własną.
Po odrzuceniu stałych pierwsza dodatnia wartość jest
minimum ilorazu Rayleigha na podprzestrzeni średniej zero;
dodatniość zapewnia nierówność Poincarégo--Wirtingera,
a istnienie minimum zwartość zanurzenia.
\end{proof}

\begin{uwaga}[Pułapka z monotonicznością]
Monotoniczność wartości własnych względem zawierania
obszarów udowodniliśmy dla warunku Dirichleta.
Nie wolno bez dodatkowych hipotez przenosić jej na
laplasjan Neumanna: przedłużenie przez zero nie działa
wtedy jako zanurzenie $H^1(\Omega_1)\to H^1(\Omega_2)$.
\end{uwaga}

\section{Ślad półgrupy cieplnej i odczytywanie wymiaru}
\label{sec:pde18-heat}
\begin{twierdzenie}[Wiodący człon śladu cieplnego]
\label{thm:pde18-heat}
Niech $\Omega$ spełnia założenia
Twierdzenia~\ref{thm:pde18-weyl}. Dla każdego $t>0$
operator $e^{t\Delta_D}$ jest klasy śladowej w $L^2(\Omega)$,
a przy $t\downarrow0$ zachodzi
\begin{equation}\label{eq:pde18-heat}
 \operatorname{Tr}(e^{t\Delta_D})=
 \sum_{k=1}^\infty e^{-t\lambda_k}
 \sim\frac{|\Omega|}{(4\pi t)^{d/2}}.
\end{equation}
\end{twierdzenie}
\begin{proof}
Prawo Weyla daje $N(\Lambda)\le C(1+\Lambda^{d/2})$.
Dla $t>0$ suma $\sum_ke^{-t\lambda_k}$ jest więc skończona.
Ponieważ $N(\Lambda)=\sum_k\chi_{[\lambda_k,\infty)}(\Lambda)$,
zasada Tonellego prowadzi do tożsamości
\begin{equation}\label{eq:pde18-heat-count}
 \sum_ke^{-t\lambda_k}=t\int_0^\infty e^{-t\Lambda}
 N(\Lambda)\,\mathrm d\Lambda
 =\int_0^\infty e^{-s}N(s/t)\,\mathrm ds.
\end{equation}
Oznaczmy $C_\Omega=\omega_d|\Omega|/(2\pi)^d$.
Z~\eqref{eq:pde18-weyl} wynika, że dla ustalonego $s>0$
$t^{d/2}N(s/t)\to C_\Omega s^{d/2}$.
Wskutek oszacowania liczby modów funkcje pod całką są
zdominowane przez $C e^{-s}(1+s^{d/2})$ dla $0<t\le1$.
Stosujemy zbieżność zdominowaną:
\[
 \lim_{t\downarrow0}t^{d/2}\sum_ke^{-t\lambda_k}
 =C_\Omega\int_0^\infty e^{-s}s^{d/2}\,\mathrm ds
 =C_\Omega\Gamma(1+d/2).
\]
Równość $\omega_d=\pi^{d/2}/\Gamma(1+d/2)$ daje
$C_\Omega\Gamma(1+d/2)=|\Omega|/(4\pi)^{d/2}$.
To dowodzi wzoru~\eqref{eq:pde18-heat}.
\end{proof}

\begin{uwaga}[Powiązanie z geometrią spektralną]
Z zachowania $\operatorname{Tr}(e^{t\Delta_D})\asymp t^{-d/2}$
można odczytać wymiar $d$, a z wiodącego współczynnika
objętość obszaru. Wyższe człony rozwinięcia dla gładkiego
brzegu zawierają jego pole i dane krzywiznowe, ale ich
wyznaczenie wymaga osobnych oszacowań jądra ciepła.
Nie wnioskujemy tu automatycznie z prawa Weyla o pełnym
rozwinięciu asymptotycznym w potęgach $t^{1/2}$.
\end{uwaga}


\begin{lemat}[Przejście abelowe od liczby modów do śladu]
\label{lem:pde18-abel}
Niech $(\nu_j)_{j\ge1}$ będzie niemalejącym ciągiem
liczb dodatnich, $\nu_j\to\infty$, oraz
$N(\Lambda)=\#\{j:\nu_j\le\Lambda\}$. Załóżmy, że dla
pewnych $a>0$, $C>0$, $0\le\gamma<a$ i $K>0$ zachodzi
$|N(\Lambda)-C\Lambda^a|\le K(1+\Lambda^\gamma)$
dla wszystkich $\Lambda\ge0$. Wtedy, gdy $t\downarrow0$,
\[
 \sum_{j\ge1}e^{-t\nu_j}
 =C\Gamma(a+1)t^{-a}+O(1+t^{-\gamma}).
\]
\end{lemat}
\begin{proof}
Ponieważ $N(\Lambda)\le C'(1+\Lambda^a)$, szereg
po lewej jest zbieżny dla $t>0$. Z twierdzenia Tonellego
oraz $e^{-t\nu_j}=t\int_{\nu_j}^\infty e^{-t\Lambda}
\,\mathrm d\Lambda$ otrzymujemy
$\sum_j e^{-t\nu_j}=t\int_0^\infty e^{-t\Lambda}
N(\Lambda)\,\mathrm d\Lambda$.
Po podstawieniu $s=t\Lambda$ wyraz główny ma wartość
$C t^{-a}\int_0^\infty e^{-s}s^a\,\mathrm ds
=C\Gamma(a+1)t^{-a}$. Moduł reszty nie przekracza
$K\int_0^\infty e^{-s}(1+t^{-\gamma}s^\gamma)
\,\mathrm ds=K(1+\Gamma(\gamma+1)t^{-\gamma})$.
\end{proof}

\begin{uwaga}[Różnica między wiodącym wyrazem a rozwinięciem]
Prawo Weyla~\eqref{eq:pde18-weyl} daje tylko $o(t^{-d/2})$
w~\eqref{eq:pde18-heat}. Oszacowanie liczby modów z jawną
resztą prowadzi, przez poprzedni lemat, do oszacowania
reszty śladu. Nie wolno jednak wyprowadzać z samego prawa
Weyla ani dwuczłonowej asymptotyki śladu, ani współczynnika
zależnego od pola brzegu.
\end{uwaga}

\section{Perturbacje widma i wzór Hellmanna--Feynmana}
\label{sec:pde18-perturb}
\begin{wniosek}[Stabilność względem potencjału]
Niech $V_1,V_2\in L^\infty(\Omega;\R)$, a macierz $A$
niech będzie ustalona, symetryczna i jednostajnie eliptyczna.
Wtedy
$|\lambda_k(L_{A,V_1})-\lambda_k(L_{A,V_2})|
\le\|V_1-V_2\|_{L^\infty(\Omega)}$ dla każdego $k\ge1$.
\end{wniosek}
\begin{proof}
Dla każdego $u$ zachodzi
$|q_{V_1}[u]-q_{V_2}[u]|
\le\|V_1-V_2\|_\infty\|u\|_2^2$.
Stosujemy dwukrotnie zasadę min--max.
\end{proof}

\begin{stwierdzenie}[Wzór Hellmanna--Feynmana przy różniczkowalnej gałęzi]
\label{prop:pde18-hf}
Niech $L_t=L_0+tW$, gdzie $L_0$ jest samosprzężonym
operatorem z dyskretnym widmem, a $W$ ograniczonym,
samosprzężonym operatorem na $L^2(\Omega)$.
Przyjmijmy, że na otwartym przedziale $I\subset\R$
można wybrać różniczkowalne w normie grafowej funkcje
$t\mapsto\lambda(t)\in\R$ i
$t\mapsto\varphi(t)\in D(L_0)$ takie, że
$L_t\varphi(t)=\lambda(t)\varphi(t)$ oraz
$\|\varphi(t)\|_2=1$. Wtedy
\begin{equation}\label{eq:pde18-hf}
 \lambda'(t)=(W\varphi(t),\varphi(t))_{L^2}.
\end{equation}
Jeżeli $W$ jest operatorem mnożenia przez funkcję $W(x)$,
to prawa strona ma postać
\[
\int_\Omega W(x)|\varphi(t,x)|^2\,\mathrm d\lambda_d(x).
\]
\end{stwierdzenie}
\begin{proof}
Różniczkujemy równość
$L_t\varphi(t)=\lambda(t)\varphi(t)$ w $L^2$.
Bierzemy iloczyn skalarny z $\varphi(t)$.
Samosprzężoność $L_t$ powoduje skasowanie składników
z $\varphi'(t)$, natomiast z warunku normalizacji
wynika $\|\varphi(t)\|_2=1$. Pozostaje~\eqref{eq:pde18-hf}.
\end{proof}

\begin{uwaga}[Dlaczego założenie różniczkowalności jest jawne?]
Wzór nie mówi sam przez się, że dowolna numerowana gałąź
$\lambda_k(t)$ jest różniczkowalna w punkcie zderzenia
wartości własnych. Dla prostych, izolowanych wartości własnych
regularność gałęzi wynika z teorii perturbacji, której pełnego
dowodu nie wprowadzamy w tej części skryptu.
\end{uwaga}

\section{Podsumowanie i mapa dalszych związków}
\label{sec:pde18-summary}
Z \eqref{eq:pde18-minmax}, \eqref{eq:pde18-weyl} i
\eqref{eq:pde18-liyau} otrzymujemy trzy wzajemnie
uzupełniające się opisy: porządek geometryczny wartości
własnych, asymptotyczną gęstość wysokich częstotliwości
i globalne, ilościowe dolne ograniczenie sumy wartości
własnych. Wzór~\eqref{eq:pde18-heat} pokazuje, że to samo
widmo kontroluje krótkoczasowe zachowanie półgrupy cieplnej.
W następnym rozdziale przejdziemy do pojemności
Sobolewa, osobliwości usuwalnych i własności subtelnych
zbiorów brzegowych.
% Rozdział 19: włączany przed \appendix do PDE_tom_I_wyklad.tex.
\chapter[Pojemność Sobolewa i osobliwości]{Pojemność Sobolewa, osobliwości usuwalne i regularność brzegu}
\label{chap:pde19}

\section{Dlaczego miara Lebesgue'a nie wystarcza?}
\label{sec:pde19-motivation}
W teorii słabych rozwiązań PDE rozróżniamy zbiory pomijalne w sensie
miary Lebesgue'a od zbiorów pomijalnych dla energii Dirichleta.
Przykładowo pojedynczy punkt w $\R^2$ ma miarę zero, a funkcja
$\log|x|$ jest harmoniczna poza tym punktem; jej gradient nie jest
jednak całkowalny z kwadratem w otoczeniu zera. Tymczasem funkcja
harmoniczna, która \emph{należy do $H^1$ także w pobliżu osobliwości},
może nie odczuwać usunięcia punktu. Sformułujemy tę zasadę przy
pomocy pojemności.

Zawsze, gdy piszemy $\Omega\subset\R^d$ jako \emph{obszar},
zakładamy, że jest on otwarty i spójny. W definicjach lokalnych
pojemności potrzebna jest jedynie otwartość. Normy, jeśli nie
zaznaczono inaczej, są rzeczywiste; argumenty dla funkcji zespolonych
prowadzi się osobno dla części rzeczywistej i urojonej.

\begin{uwaga}[Równość prawie wszędzie a quasi-ciągłość]
\label{rem:pde-v43-capacity-ae}
Element $H^1(\Omega)$ jest klasą funkcji równych prawie
wszędzie; w tym sensie nie ma zadanego punktowego śladu
na dowolnym zbiorze miary zero. Pojemność oparta na energii
pozwala rozróżniać zbiory, na których reprezentant
quasi-ciągły może mieć znaczenie. Nie wynika z tego, że
każdy zbiór miary Lebesgue'a zero ma zerową pojemność.
Relację z operatorem śladu na brzegu Lipschitzowskim
omówiono w rozdziale~\ref{chap:pde-embeddings}; są to
różne konstrukcje, wymagające odmiennych założeń.
\end{uwaga}

\section{Definicja pojemności i jej podstawowe własności}
\label{sec:pde19-capacity}
\begin{definicja}[Względna pojemność Dirichleta]
\label{def:pde19-cap}
Niech $\Omega\subset\R^d$ będzie \emph{ograniczonym zbiorem otwartym},
a $K\Subset\Omega$ zbiorem zwartym. Oznaczmy przez $\mathcal A_\Omega(K)$ zbiór tych $v\in H_0^1(\Omega)$,
dla których istnieje otwarty $U$ spełniający $K\subset U\Subset\Omega$
oraz $v\ge1$ prawie wszędzie na $U$. Definiujemy
\begin{equation}\label{eq:pde19-cap}
 \operatorname{cap}_{2,\Omega}(K)=
 \inf_{v\in\mathcal A_\Omega(K)}
 \int_\Omega|\nabla v|^2\,\mathrm d\lambda_d(x).
\end{equation}
Pojemność zbioru pustego wynosi zero. Warunek $v\ge1$ oznacza,
że rozważamy rzeczywistą nierówność prawie wszędzie, a nie
niezdefiniowane wartości punktowe funkcji Sobolewa.
\end{definicja}

\begin{uwaga}[Pojemność i norma $H^1$]
Ze względu na nierówność Poincarégo na ograniczonym zbiorze otwartym
$\Omega$ seminorma $\|\nabla v\|_2$ jest równoważna normie
$H_0^1(\Omega)$. Zatem zerowanie się pojemności z
\eqref{eq:pde19-cap} jest równoważne zerowaniu się pojemności
zdefiniowanej przez pełną normę $H^1$. Wartości liczbowe tych dwóch
pojemności są jednak różne. Wzory na pojemności kul poniżej dotyczą
\emph{energii gradientu}, a nie pełnej normy $H^1$.
\end{uwaga}

\begin{lemat}[Próby gładkie i obcinanie]
\label{lem:pde19-smooth}
W definicji~\ref{def:pde19-cap} można zastąpić $H_0^1(\Omega)$
przestrzenią $C_c^\infty(\Omega)$ bez zmiany wartości infimum.
Można ponadto wymagać $0\le v\le1$ prawie wszędzie, używając
wówczas funkcji Sobolewa.
\end{lemat}
\begin{proof}
Wystarczy przybliżyć dowolną funkcję dopuszczalną $v\in H_0^1$.
Obcięcie $T(v)=\min\{1,\max\{v,0\}\}$ zachowuje warunek
$T(v)=1$ na otoczeniu $K$, zaś reguła łańcuchowa dla funkcji
Lipschitzowskich daje $|\nabla T(v)|\le|\nabla v|$ prawie wszędzie.
Możemy więc założyć $0\le v\le1$.

Wybierzmy $K\Subset V\Subset U\Subset\Omega$, gdzie $v=1$
prawie wszędzie na $U$, oraz $\chi\in C_c^\infty(U)$ spełniającą
$\chi=1$ na otoczeniu $\overline V$. Niech $v_n\in C_c^\infty(\Omega)$
zbiega do $v$ w $H_0^1(\Omega)$, zaś $\rho_h$ będzie standardową
nieujemną aproksymacją jedynki. Przedłużając $v$ przez zero do
$\R^d$, kładziemy
$w_{n,h}=\chi(\rho_h*v)+(1-\chi)v_n$.
Oba składniki należą do $C_c^\infty(\Omega)$; dla dostatecznie
małego $h$ na $V$ mamy $\rho_h*v=1$, więc $w_{n,h}=1$ na otoczeniu
$K$. Gdy $h\downarrow0$ i $n\to\infty$, to $w_{n,h}\to v$ w $H^1$:
wynika to z ciągłości mnożenia przez $\chi$ i zbieżności splotu
w $H^1(\R^d)$. W szczególności energie gradientów zbiegają
do energii $v$. Jedna nierówność między infimami jest oczywista,
a druga wynika z tej aproksymacji.
\end{proof}

\begin{twierdzenie}[Monotoniczność i (skończona) podaddytywność]
\label{thm:pde19-basic}
Niech $\Omega\subset\R^d$ będzie ograniczonym zbiorem otwartym,
a $K,L\Subset\Omega$ zbiorami zwartymi. Wówczas:
\begin{enumerate}[label=\textup{(\roman*)}]
\item $0\le\operatorname{cap}_{2,\Omega}(K)<\infty$;
\item $K\subset L$ implikuje
$\operatorname{cap}_{2,\Omega}(K)\le\operatorname{cap}_{2,\Omega}(L)$;
\item $\operatorname{cap}_{2,\Omega}(K\cup L)\le
\operatorname{cap}_{2,\Omega}(K)+\operatorname{cap}_{2,\Omega}(L)$;
\item jeśli $K\Subset\Omega_1\subset\Omega_2$ i oba zbiory są
ograniczone i otwarte, to
$\operatorname{cap}_{2,\Omega_2}(K)\le
\operatorname{cap}_{2,\Omega_1}(K)$.
\end{enumerate}
Ponadto $\operatorname{cap}_{2,\Omega}(K)=0$ implikuje
$\lambda_d(K)=0$.
\end{twierdzenie}
\begin{proof}
Dla (i) wybieramy funkcję $\eta\in C_c^\infty(\Omega)$ równą
$1$ w otoczeniu $K$. Monotoniczność (ii) wynika z zawierania
zbiorów funkcji dopuszczalnych.

Niech $v,w\in H_0^1(\Omega)$ będą nieujemnymi funkcjami
dopuszczalnymi odpowiednio dla $K$ i $L$, z energiami nie większymi
od odpowiednich infimów plus $\eps$. Funkcja
$z=\max\{v,w\}\in H_0^1(\Omega)$ jest dopuszczalna dla $K\cup L$.
Ponieważ niemal wszędzie
\[
 |\nabla\max(v,w)|^2+|\nabla\min(v,w)|^2
 =|\nabla v|^2+|\nabla w|^2,
\]
energia $z$ nie przewyższa sumy obu energii. Granica
$\eps\downarrow0$ dowodzi (iii). W (iv) przedłużamy funkcję
$H_0^1(\Omega_1)$ przez zero do $\Omega_2$, co nie zmienia
energii ani warunku dopuszczalności.

Jeżeli $v$ jest dopuszczalna dla $K$, to
$\lambda_d(K)\le\int_\Omega|v|^2\,\mathrm d\lambda_d$.
Nierówność Poincarégo (obowiązująca dla $H_0^1$ na
ograniczonym zbiorze otwartym) daje
$\lambda_d(K)\le C_\Omega\int_\Omega|\nabla v|^2$.
Biorąc infimum, otrzymujemy ostatnią tezę.
\end{proof}

\begin{uwaga}[Miara zero nie wystarcza]
Implikacji z ostatniego twierdzenia nie można odwrócić.
W wymiarze jeden pojedynczy punkt ma dodatnią pojemność,
choć oczywiście ma miarę zero. W większych wymiarach
podobnie mogą zachowywać się odpowiednio duże fragmenty
hiperpłaszczyzn (widać to z twierdzenia o śladzie).
\end{uwaga}

\section{Dokładna pojemność kondensatora kulistego}
\label{sec:pde19-ball}
\begin{twierdzenie}[Pojemność współśrodkowych kul]
\label{thm:pde19-balls}
Niech $d\ge1$, $0<r<R$ i $K=\overline{B_r(0)}$.
Dla względnej pojemności w $B_R(0)$ zachodzi
\begin{equation}\label{eq:pde19-ball-exact}
 \operatorname{cap}_{2,B_R}(K)=
 \begin{cases}
 \displaystyle\frac{2}{R-r},&d=1,\\[5pt]
 \displaystyle\frac{2\pi}{\log(R/r)},&d=2,\\[6pt]
 \displaystyle\frac{(d-2)|S^{d-1}|}{r^{2-d}-R^{2-d}},&d\ge3,
 \end{cases}
\end{equation}
gdzie $|S^{d-1}|=d\omega_d$ jest polem sfery jednostkowej.
\end{twierdzenie}
\begin{proof}
Każda dopuszczalna funkcja $v\in H_0^1(B_R)$ jest co
najmniej równa $1$ w otoczeniu $\overline{B_r}$.
Dla prawie każdego kierunku $\vartheta\in S^{d-1}$
reprezentant $s\mapsto v(s\vartheta)$ jest absolutnie ciągły
na $[r,R]$ (z odpowiednimi śladami), a $v(r\vartheta)\ge1$,
$v(R\vartheta)=0$. Otrzymujemy
\[
 1\le\int_r^R|\partial_sv(s\vartheta)|\,\mathrm ds
 \le\left(\int_r^R|\partial_sv|^2s^{d-1}\,\mathrm ds\right)^{1/2}
 \left(\int_r^R s^{1-d}\,\mathrm ds\right)^{1/2}.
\]
Integrujemy po sferze i wykorzystujemy
$|\nabla v|\ge|\partial_sv|$. Energia całkowita jest
więc nie mniejsza od
$|S^{d-1}|/I_{r,R}$, gdzie
$I_{r,R}=\int_r^R s^{1-d}\,\mathrm ds$.

Aby wykazać osiągalność infimum, definiujemy $v=1$
na $\overline{B_r}$, $v=0$ poza $B_R$ i dla $r<s<R$
przyjmujemy funkcję radialną
\[
 v(s)=\frac{\displaystyle\int_s^R t^{1-d}\,\mathrm dt}
 {\displaystyle\int_r^R t^{1-d}\,\mathrm dt}.
\]
Należy ona do $H_0^1(B_R)$, a jej energia jest równa
$|S^{d-1}|/I_{r,R}$. Ponieważ w definicji dopuszczalności
wymagamy nierówności na \emph{otoczeniu} zbioru
$\overline{B_r}$, tę funkcję trzeba rozumieć jako granicę
funkcji skonstruowanych analogicznie z $r$ zastąpionym
przez $r+\eps$. Ich energie zbiegają do podanej liczby.
Dla $d=1$ rachunek wykonywany jest na dwóch odcinkach,
a $|S^0|=2$. Obliczenie $I_{r,R}$ daje trzy wzory
w~\eqref{eq:pde19-ball-exact}.
\end{proof}

\begin{wniosek}[Pojemność punktu]
\label{cor:pde19-point}
Jeżeli $\Omega\subset\R^d$ jest ograniczonym zbiorem otwartym,
a $x_0\in\Omega$, to
\begin{equation}\label{eq:pde19-point}
 \operatorname{cap}_{2,\Omega}(\{x_0\})=0\quad(d\ge2),
 \qquad
 \operatorname{cap}_{2,\Omega}(\{x_0\})>0\quad(d=1).
\end{equation}
\end{wniosek}
\begin{proof}
Wybierzmy $B_R(x_0)\Subset\Omega$. Z monotoniczności
względem zbioru i obszaru otrzymujemy
$\operatorname{cap}_{2,\Omega}(\{x_0\})\le
\operatorname{cap}_{2,B_R}(\overline{B_r(x_0)})$.
Prawa strona dąży do zera przy $r\downarrow0$ dla $d\ge2$;
w $d=2$ tempo jest logarytmiczne, a w $d\ge3$ rzędu
$r^{d-2}$. W jednym wymiarze nierówność
$|v(x_0)|\le\sqrt{|\Omega|}\,\|v'\|_{L^2(\Omega)}$
(dla funkcji o zerowym śladzie na końcach składowej
zawierającej $x_0$) pokazuje, że energia każdej funkcji
dopuszczalnej jest ograniczona od dołu przez dodatnią stałą.
\end{proof}

\section[Pojemność Sobolewa i wymiar krytyczny]{Pojemność $W^{1,p}$ i wymiar krytyczny}
\label{sec:pde19-p}
\begin{definicja}[Pojemność $p$-energetyczna]
Niech $1\le p<\infty$, $K\Subset\Omega$, gdzie $\Omega$ jest
ograniczonym zbiorem otwartym. Definiujemy
\[
 \operatorname{cap}_{p,\Omega}(K)=\inf\left\{
 \int_\Omega|\nabla v|^p\,\mathrm d\lambda_d(x):
 v\in W_0^{1,p}(\Omega),\ v\ge1\ \text{p.w. na otoczeniu }K\right\}.
\]
Dla $p=2$ jest to poprzednia definicja.
\end{definicja}
\begin{twierdzenie}[Punkt a wykładnik Sobolewa]
\label{thm:pde19-point-p}
Niech $d\ge2$, $1\le p<\infty$ i $x_0$ należy do ograniczonego
zbioru otwartego $\Omega\subset\R^d$. Wtedy
$\operatorname{cap}_{p,\Omega}(\{x_0\})=0$ dla $p\le d$,
natomiast $\operatorname{cap}_{p,\Omega}(\{x_0\})>0$
dla $p>d$. W przypadku $d=1$ pojemność punktu jest dodatnia
dla każdego $p\ge1$.
\end{twierdzenie}
\begin{proof}
Po translacji przyjmijmy $x_0=0$, wybierając
$B_R\Subset\Omega$. Dla $p<d$ wystarczy próbna funkcja
$v_r(x)=\eta(x/r)$ z $\eta\in C_c^\infty(B_2)$,
$\eta=1$ na $B_1$. Zmiana zmiennych daje
$\int_{\R^d}|\nabla v_r|^p\,\mathrm d\lambda_d(x) =r^{d-p}\int_{\R^d}|\nabla\eta|^p\,\mathrm d\lambda_d(x)\to0$.

Dla $p=d>1$ wybieramy $0<r<R$, a na pierścieniu
$r<|x|<R$ funkcję logarytmiczną z dowodu
Twierdzenia~\ref{thm:pde19-balls}. Jej energia $d$-ta wynosi
\[
 |S^{d-1}|\,\frac{\displaystyle\int_r^R s^{-1}\,\mathrm ds}
 {\bigl(\log(R/r)\bigr)^d}
 =\frac{|S^{d-1}|}{(\log(R/r))^{d-1}}\longrightarrow0.
\]
Dla $p>d$ nierówność Morreya i Poincarégo na kuli $B_R$
zawartej w $\Omega$ dają lokalne oszacowanie punktowe przez
normę $W^{1,p}$; aby kontrolować energię na całym zbiorze,
stosujemy je do przedłużenia przez zero funkcji z
$W_0^{1,p}(\Omega)$ na ustalony duży sześcian zawierający
$\Omega$. Otrzymujemy $|v(x_0)|\le C\|\nabla v\|_{L^p(\Omega)}$
dla ciągłego reprezentanta. Każda funkcja dopuszczalna
spełnia $v(x_0)\ge1$, co daje dodatnią dolną granicę energii.
W $d=1$ ten sam wniosek daje absolutna ciągłość i nierówność
Höldera, również dla $p=1$.
\end{proof}

\begin{uwaga}[Wymiar, miara i pojemność]
Wynik nie orzeka, że każdy zbiór wymiaru Hausdorffa $d-p$
ma zerową pojemność: przypadki graniczne są subtelne.
Zależności z miarą Hausdorffa wymagają osobnych twierdzeń,
a dla $p=d$ pojawiają się poprawki logarytmiczne.
\end{uwaga}

\section{Usuwalność osobliwości dla słabych rozwiązań}
\label{sec:pde19-removable}
\begin{twierdzenie}[Zbiór zerowej pojemności jest niewidoczny dla rozwiązań energetycznych]
\label{thm:pde19-remove}
Niech $\Omega\subset\R^d$ będzie ograniczonym zbiorem otwartym,
a $K\Subset\Omega$ zbiorem zwartym takim, że
$\operatorname{cap}_{2,\Omega}(K)=0$.
Niech $A\in L^\infty(\Omega;\R^{d\times d})$, przy czym
$\|A\|_{L^\infty}<\infty$; nie zakładamy tu symetrii.
Jeżeli $u\in H^1(\Omega)$ oraz
\begin{equation}\label{eq:pde19-off-K}
 \int_{\Omega\setminus K}A(x)\nabla u\cdot\nabla\varphi
 \,\mathrm d\lambda_d(x)=0
 \quad\text{dla każdego }\varphi\in C_c^\infty(\Omega\setminus K),
\end{equation}
to $u$ jest słabym rozwiązaniem $-\operatorname{div}(A\nabla u)=0$
na całym $\Omega$.
\end{twierdzenie}
\begin{proof}
Z Lematu~\ref{lem:pde19-smooth} i zerowej pojemności istnieją
$\eta_j\in C_c^\infty(\Omega)$, $\eta_j\ge1$ na otoczeniu
$K$, takie że $\|\nabla\eta_j\|_2\to0$. Możemy przyjąć
$0\le\eta_j\le1$ po obcięciu i odpowiedniej gładkiej
aproksymacji albo użyć gładkiej funkcji obcinającej
$\Theta$ o wartościach w $[0,1]$, równej $1$ dla
argumentów co najmniej $1$ oraz $0$ dla niedodatnich,
z ograniczoną pochodną. Wtedy
$\psi_j=\Theta(\eta_j)\in C_c^\infty(\Omega)$,
$\psi_j=1$ w otoczeniu $K$ i
$\|\psi_j\|_{H^1(\Omega)}\to0$: ostatni wniosek wynika
z nierówności Poincarégo i oszacowań
$|\Theta(s)|\le C|s|$, $|\Theta'(s)|\le C$.

Weźmy dowolną $\varphi\in C_c^\infty(\Omega)$.
Funkcja $\varphi_j=(1-\psi_j)\varphi$ jest gładka,
ma zwarty nośnik w $\Omega\setminus K$, więc z
\eqref{eq:pde19-off-K} mamy
\[
 0=\int_\Omega A\nabla u\cdot\nabla\varphi_j
 \mathrm d\lambda_d
 =\int_\Omega A\nabla u\cdot\nabla\varphi
 \mathrm d\lambda_d
 -\int_\Omega A\nabla u\cdot
 (\psi_j\nabla\varphi+\varphi\nabla\psi_j)\mathrm d\lambda_d.
\]
Dwa ostatnie składniki dążą do zera dzięki
nierówności Cauchy'ego--Schwarza, ograniczoności
$A,\varphi,\nabla\varphi$ i zbieżności
$\psi_j\to0$ w $H^1$. Otrzymujemy tożsamość słabą
na całym $\Omega$.
\end{proof}

\begin{wniosek}[Równość przestrzeni Dirichleta po usunięciu zbioru]
\label{cor:pde19-H0}
Przy założeniach Twierdzenia~\ref{thm:pde19-remove} zachodzi
\[
 H_0^1(\Omega\setminus K)=H_0^1(\Omega)
\]
po naturalnym utożsamieniu funkcji przez przedłużenie zerem
na $K$ (który ma miarę zero).
\end{wniosek}
\begin{proof}
Najpierw uzasadnijmy, że $K$ ma miarę Lebesgue'a zero.
Z definicji zerowej pojemności istnieją $v_j\in H_0^1(\Omega)$
takie, że $v_j\ge1$ na otoczeniu $K$ i
$\|\nabla v_j\|_{L^2(\Omega)}\to0$.
Nierówność Poincarégo daje $\|v_j\|_{L^2(\Omega)}\to0$,
natomiast $\lambda_d(K)\le\|v_j\|_{L^2(\Omega)}^2$;
stąd $\lambda_d(K)=0$.
Zawieranie od lewej do prawej wynika z przedłużenia przez zero.
Dla odwrotnego wystarczy gęstość na $C_c^\infty(\Omega)$.
Jeżeli $\varphi\in C_c^\infty(\Omega)$ i $\psi_j$ są jak
w poprzednim dowodzie, to
$\varphi_j=(1-\psi_j)\varphi\in C_c^\infty(\Omega\setminus K)$
i $\varphi_j\to\varphi$ w $H^1(\Omega)$. Zatem każda
funkcja z $H_0^1(\Omega)$ jest granicą funkcji testowych
unikających $K$.
\end{proof}

\begin{kontrprzyklad}[Dlaczego warunek $H^1$ jest konieczny?]
\label{ex:pde19-fundamental}
Dla $d\ge3$ funkcja $u(x)=|x|^{2-d}$, a dla $d=2$
$u(x)=\log|x|$, jest harmoniczna na $B_1\setminus\{0\}$,
ale nie należy do $H^1(B_1)$. Rzeczywiście,
$|\nabla u(x)|\asymp|x|^{1-d}$ (dla $d\ge2$), więc
\[
 \int_{B_\eps\setminus\{0\}}|\nabla u|^2
 \,\mathrm d\lambda_d(x)
 \asymp\int_0^\eps r^{1-d}\,\mathrm dr=+\infty.
\]
Po odpowiedniej normalizacji $-\Delta u$ jest deltą Diraca
w sensie dystrybucyjnym. Zerowa pojemność punktu nie
usuwa więc osobliwości przy nieograniczonej energii.
\end{kontrprzyklad}

\section{Funkcje harmoniczne ograniczone na obszarze z przekłuciem}
\label{sec:pde19-puncture}
\begin{twierdzenie}[Usuwalność ograniczonej osobliwości izolowanej]
\label{thm:pde19-bounded}
Niech $d\ge2$, $B_R(x_0)\subset\R^d$ będzie kulą,
a $u\in C^2(B_R(x_0)\setminus\{x_0\})$ funkcją harmoniczną.
Jeżeli $u$ jest ograniczona w otoczeniu $x_0$, to istnieje
jedyna funkcja harmoniczna $\widetilde u\in C^\infty(B_R(x_0))$
równa $u$ poza $x_0$.
\end{twierdzenie}
\begin{proof}
Po translacji przyjmijmy $x_0=0$. Wybierzmy
$\chi\in C_c^\infty(B_R)$ równą $1$ na $B_{R/2}$.
Ze wzoru na pojemność punktu wybieramy
$\psi_j\in C_c^\infty(B_{R/2})$ równe $1$ przy zerze,
$0\le\psi_j\le1$, z
$\|\nabla\psi_j\|_2\to0$ i $\psi_j(x)\to0$ dla $x\ne0$.
Taką rodzinę otrzymujemy, biorąc podpory w kulach
$B_{R_j}$, $R_j\downarrow0$, oraz promienie wewnętrzne
$r_j=R_j^2/R$ w $d=2$; dla $d\ge3$ wystarczą skalowane
funkcje odcięcia.

Niech $M=\|u\|_{L^\infty(B_{R/2}\setminus\{0\})}$.
W słabym równaniu harmonicznym poza zerem podstawiamy
test $\chi^2(1-\psi_j)^2u$ (dopuszczalny przez lokalną
aproksymację, ponieważ unika zera). Po rozwinięciu gradientu
i absorpcji wyrazów mieszanych otrzymujemy
\[
 \int_{B_R}\chi^2(1-\psi_j)^2|\nabla u|^2
 \,\mathrm d\lambda_d
 \le C M_0^2\int_{B_R}
 (|\nabla\chi|^2+\chi^2|\nabla\psi_j|^2)
 \,\mathrm d\lambda_d,
\]
gdzie $M_0$ jest skończonym kresem $|u|$ na nośniku $\chi$.
Stała po prawej nie zależy od $j$.
Lemat Fatou i $\psi_j\to0$ prawie wszędzie dają
$\nabla u\in L^2(B_{R/2})$.

Ponadto $u\in L^2(B_{R/2})$, a obcięcie testów
$\varphi(1-\psi_j)$ pokazuje, że pochodne dystrybucyjne
$u$ zgadzają się z jej klasycznymi pochodnymi poza zerem:
wyrazy $\int_{B_R} u\varphi\nabla\psi_j\,\mathrm d\lambda_d(x)$ dążą do zera,
gdyż $u$ jest ograniczona, a
$\|\nabla\psi_j\|_{L^1}\le|B_R|^{1/2}\|\nabla\psi_j\|_2\to0$.
Wynika stąd $u\in H^1(B_{R/2})$. Twierdzenie
\ref{thm:pde19-remove} i lemat Weyla z rozdziału~7
pokazują, że $u$ przedłuża się harmonicznie przez zero.
Argument na mniejszych i większych kulach daje przedłużenie
na całe $B_R$; jednoznaczność wynika z ciągłości.
\end{proof}

\begin{przyklad}[Punkt brzegowy, w którym dane Dirichleta giną]
\label{ex:pde19-irregular}
Niech $d\ge2$ i $\Omega=B_1(0)\setminus\{0\}$.
Na brzegu $\partial\Omega=S^{d-1}\cup\{0\}$ zadajmy
funkcję ciągłą $g$, równą $0$ na $S^{d-1}$ oraz $1$ w zerze.
Nie istnieje ograniczona funkcja harmoniczna na $\Omega$,
ciągła na $\overline\Omega$ i przyjmująca wartości $g$.
W istocie taka funkcja na mocy poprzedniego twierdzenia
przedłużyłaby się harmonicznie na $B_1$. Z zasady maksimum,
przy zerowych danych na $S^{d-1}$, byłaby wszędzie zerowa.
Przeczy to wartości $1$ narzuconej w zerze.
To przykład \emph{nieregularnego punktu brzegu}.
\end{przyklad}

\section{Pojemność, reprezentanci i zbiory wyjątkowe}
\label{sec:pde19-quasi}
\begin{definicja}[Pojemność zbioru otwartego i quasi-ciągłość]
Dla otwartego $U\subset\Omega$ definiujemy
$\operatorname{cap}_{2,\Omega}(U)$ jak w~\eqref{eq:pde19-cap},
ale z warunkiem $v\ge1$ prawie wszędzie na $U$;
nie wymagamy $U\Subset\Omega$ i dopuszczamy wartość $+\infty$.
Dla dowolnego $E\subset\Omega$ kładziemy
\[
 \operatorname{cap}_{2,\Omega}(E)
 =\inf\{\operatorname{cap}_{2,\Omega}(U):E\subset U\subset\Omega,
 \ U\text{ otwarty}\}.
\]
Funkcję $\widetilde u:\Omega\to\R$, określoną z dokładnością
do zbioru pojemności zero, nazywamy \emph{quasi-ciągłą}, jeśli
 dla każdego $\eps>0$ istnieje otwarty $G\subset\Omega$,
$\operatorname{cap}_{2,\Omega}(G)<\eps$, taki że
$\widetilde u|_{\Omega\setminus G}$ jest ciągła.
Własność zachodzącą poza zbiorem pojemności zero nazywamy
własnością zachodzącą \emph{quasi-wszędzie}.
\end{definicja}

\begin{lemat}[Podaddytywność przeliczalna pojemności]
\label{lem:pde19-countable}
Dla dowolnych $E_j\subset\Omega$, $j\ge1$, zachodzi
$\operatorname{cap}_{2,\Omega}(\bigcup_jE_j)
 \le\sum_j\operatorname{cap}_{2,\Omega}(E_j)$.
\end{lemat}
\begin{proof}
Najpierw rozważmy zbiory otwarte $U_j$ o skończonej sumie
pojemności. Wybieramy $v_j\in H_0^1(\Omega)$, $v_j\ge1$
prawie wszędzie na $U_j$, $0\le v_j\le1$, z
$\int_\Omega|\nabla v_j|^2\,\mathrm d\lambda_d(x)\le\operatorname{cap}_{2,\Omega}(U_j)
 +\eps 2^{-j}$. Dla $m\ge1$ kładziemy
$w_m=\max(v_1,\ldots,v_m)$. Z tożsamości gradientowej
użytej w dowodzie Twierdzenia~\ref{thm:pde19-basic}
wynika $\|\nabla w_m\|_2^2\le\sum_{j=1}^m
\|\nabla v_j\|_2^2$. Normy $H_0^1$ są więc jednostajnie
ograniczone. Ponadto $0\le w_m\le1$ i $w_m$ jest
niemalejący. Z monotonicznej zbieżności punktowej
i zdominowanej zbieżności w $L^2(\Omega)$ wynika, że
$w_m\to w$ silnie w $L^2$; słaba zwartość w $H_0^1$
i półciągłość normy dają $w\in H_0^1$ oraz
$\|\nabla w\|_2^2\le\sum_j\|\nabla v_j\|_2^2$.
Ponieważ $w\ge1$ prawie wszędzie na $\bigcup_jU_j$,
uzyskujemy żądaną nierówność dla zbiorów otwartych.
Z definicji zewnętrznej pojemności wybieramy
$U_j\supset E_j$ z
$\operatorname{cap}(U_j)\le\operatorname{cap}(E_j)+\eps2^{-j}$
i przechodzimy z $\eps\downarrow0$.
\end{proof}

\begin{twierdzenie}[Quasi-ciągły reprezentant funkcji Sobolewa]
\label{thm:pde19-quasi}
Niech $\Omega\subset\R^d$ będzie ograniczonym zbiorem otwartym.
Każda rzeczywista funkcja $u\in H_0^1(\Omega)$ ma
reprezentanta quasi-ciągłego $\widetilde u$, równego $u$
prawie wszędzie.
\end{twierdzenie}
\begin{proof}
Wybieramy $u_n\in C_c^\infty(\Omega)$ tak szybko zbieżne
do $u$ w $H_0^1(\Omega)$, aby
$\|u_{n+1}-u_n\|_{H^1}^2\le2^{-4n}$.
Niech $G_n=\{x\in\Omega:|u_{n+1}(x)-u_n(x)|>2^{-n}\}$.
Ponieważ $u_{n+1}-u_n\in C_c^\infty(\Omega)$,
funkcja $2^n|u_{n+1}-u_n|$ należy do $H_0^1$ i jest
co najmniej $1$ na $G_n$. Zatem
\[
 \operatorname{cap}_{2,\Omega}(G_n)
 \le2^{2n}\|\nabla(u_{n+1}-u_n)\|_2^2\le2^{-2n}.
\]
Dla $F_m=\bigcup_{n\ge m}G_n$ z
Lematu~\ref{lem:pde19-countable} mamy
$\operatorname{cap}(F_m)\le\sum_{n\ge m}2^{-2n}\to0$.
Na $\Omega\setminus F_m$ szereg
$\sum_{n\ge m}(u_{n+1}-u_n)$ jest jednostajnie i bezwzględnie
zbieżny, ponieważ każdy składnik ma tam moduł nie większy niż
$2^{-n}$. Granica $\widetilde u$ jest na tym zbiorze ciągła
jako jednostajna granica funkcji ciągłych.

Zbiór $F=\bigcap_m F_m$ ma pojemność zero, ponieważ
$F\subset F_m$ dla wszystkich $m$; poza $F$
ciąg $u_n$ zbiega punktowo. Przez zbieżność $L^2$ po
wybraniu podciągu otrzymujemy także $\widetilde u=u$
prawie wszędzie. Wartości na $F$ można zdefiniować
dowolnie. Dla każdego $\eps>0$ wybór dużego $m$ sprawia,
że $\operatorname{cap}(F_m)<\eps$, a ograniczenie
$\widetilde u$ do dopełnienia $F_m$ jest ciągłe.
\end{proof}

\begin{uwaga}[Po co quasi-ciągłość?]
Funkcję Sobolewa określamy zwykle jedynie prawie wszędzie,
więc zapis $u(x_0)$ nie ma sensu dla dowolnie wybranego
punktu. Twierdzenie~\ref{thm:pde19-quasi} zastępuje
nieistniejącą ciągłość klasyczną ciągłością poza zbiorem
\emph{małym energetycznie}. Jest to podstawa subtelnej
interpretacji warunków brzegowych oraz przeszkód w zagadnieniach
wariacyjnych. Nie utożsamiamy przy tym automatycznie
pojęcia quasi-ciągłości z regularnym śladem na
brzegu Lipschitzowskim.
\end{uwaga}

\section{Bariera zewnętrznej kuli i kryterium Wienera}
\label{sec:pde19-wiener}
\begin{definicja}[Bariera harmoniczna]
Niech $\Omega\subset\R^d$ będzie ograniczonym obszarem,
a $x_0\in\partial\Omega$. Lokalną barierą harmoniczną
w $x_0$ nazywamy nieujemną funkcję harmoniczną $b$
w $\Omega\cap B_\rho(x_0)$, dodatnią w jego wnętrzu,
która spełnia $b(x)\to0$ przy $x\to x_0$ od strony $\Omega$.
Do użycia w rozwiązywaniu zagadnienia Dirichleta potrzebna
jest również odpowiednia dodatnia dolna granica bariery
na sztucznej części brzegu lokalnej dziedziny.
\end{definicja}

\begin{stwierdzenie}[Jawna bariera przy zewnętrznej kuli]
\label{prop:pde19-exterior-ball}
Niech $d\ge2$, $\Omega\subset\R^d$ będzie ograniczonym
obszarem, a $x_0\in\partial\Omega$. Załóżmy, że istnieje
kula $B_r(y)\subset\R^d\setminus\overline\Omega$
taka, że $|x_0-y|=r$. Wtedy funkcja
\[
 b(x)=\begin{cases}
 r^{2-d}-|x-y|^{2-d},&d\ge3,\\[2pt]
 \log(|x-y|/r),&d=2
 \end{cases}
\]
jest dodatnia i harmoniczna w $\Omega$, a
$b(x)\to0$ przy $x\to x_0$ w $\Omega$.
\end{stwierdzenie}
\begin{proof}
Ponieważ kula $B_r(y)$ jest zawarta w dopełnieniu
$\overline\Omega$, każdy $x\in\Omega$ spełnia
$|x-y|>r$. Funkcja radialna $s\mapsto s^{2-d}$
(dla $d\ge3$) albo $s\mapsto\log s$ (dla $d=2$)
jest harmoniczna poza środkiem $y$. Wskazane kombinacje
są dodatnie dla $s>r$ i zerują się dla $s=r$.
Ciągłość funkcji radialnej daje granicę zero w $x_0$.
Na zwartej części brzegu lokalnego obszaru oddzielonej
od $x_0$ bariera ma dodatnie minimum, jeżeli jest ona
oddzielona również od sfery $\partial B_r(y)$.
\end{proof}

\begin{uwaga}[Warunek geometryczny a rozwiązywalność]
Istnienie bariery zewnętrznej kuli prowadzi do regularności
$x_0$ dla zagadnienia Dirichleta w sensie Perrona, po
zastosowaniu zasady porównania. Pełna konstrukcja Perrona
nie jest tutaj dowodzona. Warunek zewnętrznej kuli jest
wystarczający, ale niekonieczny.
\end{uwaga}

\begin{definicja}[Bezwzględna pojemność Newtonowska]
Dla $d\ge3$ i zwartego $K\subset\R^d$ określamy
\[
 \operatorname{Cap}_2(K)=\inf\left\{
 \int_{\R^d}|\nabla\varphi|^2\,\mathrm d\lambda_d(x):
 \varphi\in C_c^\infty(\R^d),\ \varphi\ge1
 \text{ na otoczeniu }K\right\}.
\]
W odróżnieniu od~\eqref{eq:pde19-cap} nie występuje tu
ustalony obszar zewnętrzny ani wyraz $L^2$.
\end{definicja}


\begin{lemat}[Porównanie pojemności względnej i bezwzględnej]
\label{lem:pde19-cap-comparison}
Niech $d\ge3$, $\Omega\subset\R^d$ będzie ograniczonym
zbiorem otwartym, a $D\Subset\Omega$ zwartym zbiorem.
Istnieje stała $C=C(d,D,\Omega)$ taka, że dla każdego
zwartego $K\subset D$
\[
 \operatorname{Cap}_2(K)\le\operatorname{cap}_{2,\Omega}(K)
 \le C\operatorname{Cap}_2(K).
\]
W szczególności zerowość tych dwóch pojemności jest
równoważna na ustalonym zwartym podzbiorze wnętrza $\Omega$.
\end{lemat}
\begin{proof}
W pierwszej nierówności stosujemy przedłużenie przez
zero funkcji dopuszczalnych. Dzięki
Lematowi~\ref{lem:pde19-smooth} możemy ograniczyć się
do prób gładkich. Dla drugiej nierówności wybierzmy
$\chi\in C_c^\infty(\Omega)$, $\chi=1$ na otoczeniu $D$.
Jeżeli $\varphi\in C_c^\infty(\R^d)$ jest dopuszczalna
dla $\operatorname{Cap}_2(K)$, to $\chi\varphi$ jest
dopuszczalna dla $\operatorname{cap}_{2,\Omega}(K)$.
Mamy $|\nabla(\chi\varphi)|^2\le
2\chi^2|\nabla\varphi|^2+2|\nabla\chi|^2|\varphi|^2$.
Z nierówności Sobolewa w $\R^d$ dla $d\ge3$ oraz
nierówności Höldera na zwartym nośniku $\nabla\chi$
wynika
\[
 \int_\Omega|\nabla(\chi\varphi)|^2
 \,\mathrm d\lambda_d(x)
 \le C\int_{\R^d}|\nabla\varphi|^2
 \,\mathrm d\lambda_d(x).
\]
Bierzemy infimum po $\varphi$.
\end{proof}

\begin{twierdzenie}[Kryterium Wienera, $d\ge3$ — wynik przywołany]
\label{thm:pde19-wiener}
Niech $d\ge3$, $\Omega\subset\R^d$ będzie ograniczonym
obszarem i $x_0\in\partial\Omega$. Dla $k$ dostatecznie
dużych niech
\[
 E_k=(\R^d\setminus\Omega)\cap
 \{x:2^{-k-1}\le|x-x_0|\le2^{-k}\}.
\]
Punkt $x_0$ jest regularny dla klasycznego zagadnienia
Dirichleta w sensie Perrona wtedy i tylko wtedy, gdy
\begin{equation}\label{eq:pde19-wiener}
 \sum_{k=k_0}^\infty
 2^{k(d-2)}\operatorname{Cap}_2(E_k)=+\infty.
\end{equation}
Tutaj regularność oznacza, że rozwiązanie Perrona
przyjmuje w granicy od strony $\Omega$ wartość
$f(x_0)$ dla każdej ciągłej funkcji brzegowej
$f\in C(\partial\Omega)$.
\end{twierdzenie}

\begin{uwaga}[Zakres dowodu kryterium]
Kryterium Wienera jest głębokim twierdzeniem teorii
potencjału: \emph{nie dowodzimy tu jego pełnej implikacji
w obu kierunkach}. Z poprzednich twierdzeń wykazaliśmy
natomiast jawnie wzór pojemności kul i nieregularność
izolowanego punktu brzegowego. Dla przekłutej kuli
z~\ref{ex:pde19-irregular} wszystkie dostatecznie małe
$E_k$ są puste, zatem szereg jest zbieżny. Dla dopełnienia
zawierającego stały stożek z wierzchołkiem w $x_0$
pojemności kolejnych fragmentów są porównywalne z
$2^{-k(d-2)}$, co daje rozbieżność. Kryterium w wymiarze
dwa przyjmuje postać logarytmiczną i nie jest tu
formułowane.
\end{uwaga}


\begin{stwierdzenie}[Zewnętrzny stożek a szereg Wienera]
\label{prop:pde19-cone-wiener}
Niech $d\ge3$, $\Omega\subset\R^d$ będzie ograniczonym obszarem
oraz $x_0\in\partial\Omega$. Załóżmy, że istnieją $r_0>0$
i niepusty, otwarty zbiór kierunków $\Sigma\subset S^{d-1}$
takie, że $x_0+r\theta\notin\Omega$ dla każdego
$0<r<r_0$ i $\theta\in\Sigma$.
Wtedy dla zbiorów $E_k$ z Twierdzenia~\ref{thm:pde19-wiener}
istnieje $c>0$ i $k_1$ takie, że
\[
  2^{k(d-2)}\operatorname{Cap}_2(E_k)\ge c
  \qquad(k\ge k_1).
\]
W szczególności szereg Wienera jest rozbieżny, a
$x_0$ jest regularnym punktem brzegowym, jeśli stosujemy
przywołane Twierdzenie~\ref{thm:pde19-wiener}.
\end{stwierdzenie}
\begin{proof}
Wybierzmy $\theta_0\in\Sigma$ i $\delta>0$ tak małe,
aby kierunki z kuli sferycznej wokół $\theta_0$ o promieniu
$2\delta$ należały do $\Sigma$, przy czym $\delta<1/16$.
Dla $h=2^{-k}<r_0$ weźmy
$z_k=x_0+3h\theta_0/4$.
Z elementarnej geometrii wynika, że dla dostatecznie dużych
$k$ kula $\overline{B_{\delta h}(z_k)}$ leży w stożku
zewnętrznym i w pierścieniu
$\{h/2\le|x-x_0|\le h\}$, a zatem w $E_k$.
Monotoniczność pojemności bezwzględnej daje
$\operatorname{Cap}_2(E_k)\ge
\operatorname{Cap}_2(\overline{B_{\delta h}(z_k)})$.
Przesunięcie nie zmienia pojemności, a z obliczenia energii
radialnej na pierścieniu i przejścia z jego promieniem
zewnętrznym do nieskończoności (dla $d\ge3$) mamy
$\operatorname{Cap}_2(\overline{B_\rho})
=(d-2)|S^{d-1}|\rho^{d-2}$.
Stąd dolne ograniczenie jest równe
$(d-2)|S^{d-1}|\delta^{d-2}h^{d-2}$, co po
pomnożeniu przez $h^{-(d-2)}$ dowodzi tezy.
\end{proof}

\begin{przyklad}[Przekłuta kula a kryterium Wienera]
\label{ex:pde19-wiener-puncture}
Dla $d\ge3$ oraz $\Omega=B_1(0)\setminus\{0\}$
w punkcie $x_0=0$ wszystkie dostatecznie małe pierścienie
$E_k$ są puste, gdyż dopełnienie obszaru zawiera
wewnątrz $B_1$ jedynie punkt $0$, niewchodzący do
żadnego z tych pierścieni. Suma z kryterium Wienera wynosi
zero, zatem $x_0$ jest nieregularny. Wynik zgadza się z
Przykładem~\ref{ex:pde19-irregular}.
\end{przyklad}



\begin{lemat}[Usuwalność ograniczonej osobliwości harmonicznej w wymiarze dwa]
\label{lem:pde-v5-removable-bounded-2d}
Niech $R>0$ i niech $h$ będzie rzeczywistą funkcją harmoniczną
na $B_R(0)\setminus\{0\}\subset\R^2$. Jeżeli $h$ jest ograniczona
w otoczeniu zera, to ma jednoznaczne przedłużenie harmoniczne
na całą kulę $B_R(0)$.
\end{lemat}
\begin{proof}
Wybierzmy $0<r_0<R$ i zapiszmy $x=(r\cos\theta,r\sin\theta)$.
Dla $0<r<r_0$ oznaczmy przez
\[
 c_n(r)=\frac1{2\pi}\int_{0}^{2\pi}
 h(r\cos\theta,r\sin\theta)e^{-in\theta}
 \,\mathrm d\lambda_1(\theta),\qquad n\in\mathbb Z,
\]
zespolone współczynniki Fouriera ograniczonej funkcji
$\theta\mapsto h(r\cos\theta,r\sin\theta)$.
Równanie Laplace'a we współrzędnych biegunowych daje
$r^2c_n''+rc_n'-n^2c_n=0$. Dla $n=0$ zatem
$c_0(r)=A_0+B_0\log r$, a dla $n\ne0$
$c_n(r)=A_nr^{|n|}+B_nr^{-|n|}$.
Ograniczoność $h$ oznacza $|c_n(r)|\le\|h\|_\infty$
dla dostatecznie małych $r$, więc wszystkie współczynniki
$B_n$ muszą znikać. Niech $H$ będzie harmonicznym
przedłużeniem Poissona funkcji $h|_{\partial B_{r_0}}$
do $B_{r_0}(0)$. Jego współczynniki Fouriera na
okręgu promienia $r$ wynoszą
$(r/r_0)^{|n|}c_n(r_0)$; dla $n=0$ są stałe.
Są zatem identyczne z $c_n(r)$ dla $0<r<r_0$.
Z jednoznaczności szeregów Fouriera dla funkcji ciągłych
na okręgu wynika $H=h$ w przekłutej kuli $B_{r_0}$.
Funkcja otrzymana przez dołączenie $H(0)$ jest więc
harmonicznym przedłużeniem $h$ na całą $B_R$.
Jednoznaczność wynika z ciągłości.
\end{proof}

\begin{stwierdzenie}[Nieregularność izolowanego punktu brzegu w wymiarze dwa]
\label{prop:pde-v5-punctured-disk}
Niech $\Omega=B_1(0)\setminus\{0\}\subset\R^2$.
Punkt $0\in\partial\Omega$ nie jest regularny dla
klasycznego zagadnienia Dirichleta w sensie Perrona.
Dokładniej, dla danych brzegowych
$f(0)=1$ oraz $f|_{\partial B_1}=0$ rozwiązanie
Perrona jest identycznie równe zeru w $\Omega$.
\end{stwierdzenie}
\begin{proof}
Funkcja $f$ jest ciągła na rozłącznej sumie
$\partial\Omega=\{0\}\cup\partial B_1$.
Z zasady porównania ze stałymi wynika, że rozwiązanie
Perrona $u$ jest ograniczone przez $0\le u\le1$;
standardowa konstrukcja Perrona zapewnia jego harmoniczność
wewnątrz $\Omega$.
Z lematu~\ref{lem:pde-v5-removable-bounded-2d} funkcja $u$
przedłuża się harmonicznie przez zero. Każdy punkt
zewnętrznej sfery $\partial B_1$ ma barierę pochodzącą
od kuli zewnętrznej, więc $u$ dąży tam do wartości $0$
(przywołujemy tu standardową własność rozwiązania Perrona
w punktach z barierą).
Zasada maksimum dla funkcji harmonicznych w $B_1$
implikuje $u\equiv0$. Stąd granica $u(x)$ przy $x\to0$
wynosi $0$, a nie $f(0)=1$, co dowodzi nieregularności.
Argument nie wymaga pełnego logarytmicznego kryterium Wienera.
\end{proof}

\begin{przyklad}[Rodzina punktów regularnych i nieregularnych
sterowana szeregiem]
\label{ex:pde19-wiener-balls}
Niech $d\ge3$, $e_1=(1,0,\ldots,0)$, $h_k=2^{-k}$,
$k\ge2$, oraz $0<\delta_k\le1/16$. Określmy
$z_k=\tfrac34h_ke_1$, $K_k=\overline{B_{\delta_kh_k}(z_k)}$
i
\[
 \Omega=B_1(0)\setminus\biggl(\{0\}\cup
 \bigcup_{k=2}^\infty K_k\biggr).
\]
Wówczas $\Omega$ jest ograniczonym, otwartym i spójnym
obszarem, a $0\in\partial\Omega$. Kryterium
Wienera~\ref{thm:pde19-wiener} daje dokładnie
\[
 0\text{ jest regularny}\quad\Longleftrightarrow\quad
 \sum_{k=2}^\infty\delta_k^{d-2}=+\infty.
\]
W szczególności, dla $\delta_k=\frac1{16}k^{-a}$,
$a\ge0$, punkt $0$ jest regularny wtedy i tylko wtedy,
gdy $a(d-2)\le1$.
\end{przyklad}
\begin{proof}
Kule $K_k$ są parami rozłączne. Ich środki leżą na
promieniach $3h_k/4$, a promienie nie przekraczają
$h_k/16$. Dlatego $K_k$ zawiera się w otwartym
pierścieniu $h_k/2<|x|<h_k$. Pozostałe kule nie przecinają
tego pierścienia; zbiór $E_k$ z definicji kryterium Wienera
jest więc równy $K_k$. Usunięty zbiór jest domknięty
w $B_1$, ponieważ jego jedynym punktem skupienia poza
kulami jest $0$; dołączono ten punkt. Spójność $\Omega$
można wykazać, łącząc dowolne dwa punkty łamanymi,
a następnie omijając skończenie wiele kul przecinających
zadaną łamaną; pozostałe kule leżą przy $0$, którego
łamana może unikać w dodatniej odległości.

Dla $d\ge3$ bezwzględna pojemność kuli ma wzór
$\operatorname{Cap}_2(K_k)=(d-2)|S^{d-1}|
(\delta_kh_k)^{d-2}$ (Twierdzenie~\ref{thm:pde19-balls}
z promieniem zewnętrznym dążącym do nieskończoności).
Stąd
\[
 2^{k(d-2)}\operatorname{Cap}_2(E_k)
 =(d-2)|S^{d-1}|\delta_k^{d-2}.
\]
Szereg Wienera jest zatem rozbieżny dokładnie wtedy,
gdy rozbieżny jest wskazany szereg liczbowy. Ostatnia
równoważność wynika z kryterium zbieżności szeregu
$\sum_{k\ge2}k^{-a(d-2)}$.
\end{proof}

\begin{uwaga}[Granica samodzielności argumentu]
Geometria kul, obliczenie pojemności i redukcja do szeregu
liczbowego zostały wyprowadzone powyżej. Równoważność
między rozbieżnością szeregu a regularnością punktu brzegu
jest zastosowaniem \emph{przywołanego}, a nie udowodnionego
tutaj kryterium Wienera.
\end{uwaga}

\section{Związki z dalszą regularnością i podsumowanie}
\label{sec:pde19-summary}
Pojemność rozróżnia zbiory, na których w naturalny sposób
może koncentrować się energia rozwiązań, od tych, które
są dla przestrzeni Sobolewa niewidoczne. Wykazaliśmy
trzy konkretne mechanizmy: znikanie pojemności punktów
w wymiarach $d\ge2$, przedłużanie równań dywergencyjnych
przez zbiory pojemności zero dla rozwiązań $H^1$ oraz
istnienie quasi-ciągłych reprezentantów funkcji Sobolewa.
Kryterium Wienera opisuje natomiast regularność warunku
brzegowego za pomocą rozłożenia pojemności na skalach.

W późniejszej teorii jedyności kontynuacji i analizie
mikrolokalnej pojawią się inne sposoby mierzenia wielkości
zbiorów wyjątkowych; nie są one równoważne pojemności
w dowolnych klasach współczynników.


% Rozdzial 20 -- samodzielny fragment do wlaczenia przed \appendix.
\chapter[Paraboliczne równania ewolucyjne]{Paraboliczne równania ewolucyjne: rozwiązania słabe, metoda Galerkina i regularność}
\label{chap:pde20}

\section{Od równania ciepła do współczynników mierzalnych}
W rozdziale~4 rozwiązywaliśmy równanie ciepła w przestrzeni całej
przy użyciu jądra Gaussa. Obecnie rozważymy znacznie ogólniejsze
zagadnienie na ograniczonym obszarze $\Omega\subset\R^d$:
\begin{equation}\label{eq:pde20-eq}
 \partial_t u-\operatorname{div}(A(t,x)\nabla u)=f
 \quad\text{w }(0,T)\times\Omega,\qquad
 u|_{\partial\Omega}=0,\qquad u(0)=u_0.
\end{equation}
Macierz $A$ może być tylko mierzalna zarówno w czasie, jak i w przestrzeni.
W takich warunkach nie należy oczekiwać klasycznych pochodnych drugiego
rzędu. W przeciwieństwie do zagadnienia eliptycznego rozwiązanie zależy
od czasu, ale nie zakładamy istnienia pochodnej $u_t$ o wartościach w $L^2(\Omega)$.

\begin{definicja}[Dane i eliptyczność paraboliczna]\label{def:pde20-data}
Niech $\Omega\subset\R^d$ będzie ograniczonym obszarem Lipschitzowskim,
$0<T<\infty$, $H=L^2(\Omega)$, $V=H_0^1(\Omega)$, $V^*=H^{-1}(\Omega)$.
Rozpatrujemy rzeczywistą macierz $A=(a_{ij})\in
L^\infty((0,T)\times\Omega;\R^{d\times d})$ taką, że dla pewnych
$0<\lambda\le\Lambda<\infty$, dla prawie wszystkich $(t,x)$ i każdego
$\xi\in\R^d$ zachodzi
\[
 \lambda|\xi|^2\le A(t,x)\xi\cdot\xi,
 \qquad |A(t,x)\xi|\le\Lambda|\xi|.
\]
Nie wymagamy symetrii macierzy $A$. Dla $v,w\in V$ określamy
$a(t;v,w)=\int_\Omega A(t,x)\nabla v(x)\cdot\nabla w(x)
\,\mathrm d\lambda_d(x)$. Przyjmujemy $u_0\in H$ oraz
$f\in L^2(0,T;V^*)$; dualność $V^*,V$ oznaczamy przez
$\langle\cdot,\cdot\rangle$.
\end{definicja}


\begin{center}
\begin{tikzpicture}[scale=0.9,>=Latex]
  \fill[blue!4] (0,0) rectangle (6,3.4);
  \draw[very thick,blue!45!black] (0,0) -- (0,3.4) -- (6,3.4) -- (6,0);
  \draw[very thick,teal!65!black] (0,0)--(6,0);
  \draw[->] (-0.4,0)--(6.8,0) node[right] {$x$};
  \draw[->] (0,-0.25)--(0,4.1) node[above] {$t$};
  \node[below] at (3,-0.07) {$u(0,x)=u_0(x)$};
  \node[rotate=90,above] at (0,1.8) {$u(t,0)=0$};
  \node[rotate=90,below] at (6,1.8) {$u(t,1)=0$};
  \node[above] at (3,3.4) {$t=T$ (wartość nie jest zadawana)};
  \node at (3,1.75) {$Q_T=(0,T)\times(0,1)$};
  \node[below] at (0,-0.4) {$0$};
  \node[below] at (6,-0.4) {$1$};
\end{tikzpicture}

\small Rysunek: jednowymiarowy przykład ($\Omega=(0,1)$). Podstawa i brzeg boczny walca parabolicznego.
Dane początkowe i warunek Dirichleta określa się na różnych częściach brzegu.
\end{center}

Normę w $V$ ustalamy jako $\|v\|_V=\|\nabla v\|_{L^2(\Omega)}$.
Dzięki nierówności Poincarégo jest ona równoważna zwykłej normie $H^1$.
Mamy ciągłe, gęste zanurzenia $V\hookrightarrow H\cong H^*\hookrightarrow V^*$.
W szczególności $\|h\|_{V^*}\le C_P\|h\|_H$ dla $h\in H$.

\begin{uwaga}[Trzy odrębne narzędzia: forma, baza i zwartość]
\label{rem:pde-v43-parabolic-basis}
W dowodzie istnienia wykorzystamy koercywność formy
(znaną z rozdziału~\ref{chap:pde-variational}), bazę własną
pomocniczego $-\Delta_D$ z rozdziału~\ref{chap:pde18}
oraz zwarte zanurzenie z rozdziału~\ref{chap:pde-embeddings}.
Baza laplasjanu nie diagonalizuje na ogół macierzy
$A(t,x)$; służy do zbudowania stabilnych rzutów
$P_n$ w $V$, $H$ i $V^*$. Dlatego ani niesymetryczność,
ani mierzalna zależność $A$ od czasu nie podważają
przedstawionej konstrukcji. Własności samosprzężonej
półgrupy $e^{-tL}$ w dalszej części rozdziału
odnoszą się natomiast do wyraźnie wskazanego przypadku
$L=-\Delta_D$ (lub innego stałego operatora dodatniego
samosprzężonego), nie do ogólnego $A(t,x)$.
\end{uwaga}

\begin{definicja}[Rozwiązanie słabe w przestrzeni energii]
\label{def:pde20-weak}
Rozwiązaniem słabym zagadnienia~\eqref{eq:pde20-eq} nazywamy funkcję
\[
 u\in L^2(0,T;V),\qquad u'\in L^2(0,T;V^*),
\]
której pochodną rozumiemy dystrybucyjnie w czasie oraz która spełnia
$\langle u'(t),v\rangle+a(t;u(t),v)=\langle f(t),v\rangle$
dla każdego $v\in V$ i prawie każdego $t\in(0,T)$, przy czym zbiór
wyjątkowy można wybrać niezależnie od $v$.
Warunek $u(0)=u_0$ dotyczy ciągłego reprezentanta $H$-wartościowego,
którego istnienie wykażemy poniżej.
\end{definicja}

\section{Ciągłość czasowa i tożsamość energii}
\begin{lemat}[Tożsamość energii w trójce Gelfanda]\label{lem:pde20-energy}
Niech $V=H_0^1(\Omega)$, $H=L^2(\Omega)$, a
$u\in L^2(0,T;V)$ ma pochodną $u'\in L^2(0,T;V^*)$.
Wówczas $u$ ma jednoznacznego reprezentanta należącego do
$C([0,T];H)$ i dla wszystkich $0\le s\le t\le T$ zachodzi
\begin{equation}\label{eq:pde20-energy-identity}
 \frac12\|u(t)\|_H^2-\frac12\|u(s)\|_H^2
 =\int_s^t\langle u'(r),u(r)\rangle\,\mathrm dr.
\end{equation}
\end{lemat}
\begin{proof}
Niech $\{e_j\}_{j\ge1}$ będzie ortonormalną bazą $H$ złożoną
z funkcji własnych dodatniego laplasjanu Dirichleta,
$-\Delta_D e_j=\mu_j e_j$, $\mu_j>0$.
Z twierdzenia spektralnego dla jego formy wynika, że
$\|v\|_V^2=\sum_{j\ge1}\mu_j|(v,e_j)_H|^2$ dla $v\in V$.
Rzut $P_n$ na $\operatorname{span}\{e_1,\ldots,e_n\}$ ma normę
co najwyżej $1$ zarówno w $V$, jak i w $V^*$ (w drugiej przestrzeni
stosujemy rozwinięcie z wagami $\mu_j^{-1}$).
Co więcej, $P_nu\to u$ w $L^2(0,T;V)$ i
$P_nu'\to u'$ w $L^2(0,T;V^*)$.

Współczynniki $t\mapsto (u(t),e_j)_H$ należą do $H^1(0,T)$,
ponieważ ich pochodnymi są $\langle u'(t),e_j\rangle$.
Dla skończonego $n$ zwykła reguła różniczkowania daje
\[
 \|P_nu(t)\|_H^2-\|P_nu(s)\|_H^2
 =2\int_s^t\langle P_nu'(r),P_nu(r)\rangle\,\mathrm dr.
\]
Wykażemy najpierw zbieżność $P_nu$ w $C([0,T];H)$.
Dla $w=P_nu-P_mu$ wybieramy punkt $s\in[0,T]$ taki, że
$\|w(s)\|_H^2\le T^{-1}\|w\|_{L^2(0,T;H)}^2$;
istnieje on, gdyż średnia nie może być mniejsza od wszystkich wartości.
Stosując poprzednią tożsamość na przedziale o końcach $s,t$
i nierówność Cauchy'ego--Schwarza, otrzymujemy
\begin{equation}\label{eq:pde20-cont-est}
 \sup_{0\le t\le T}\|w(t)\|_H^2
 \le T^{-1}\|w\|_{L^2(0,T;H)}^2
 +2\|w'\|_{L^2(0,T;V^*)}\|w\|_{L^2(0,T;V)}.
\end{equation}
Oba składniki po prawej dążą do zera dla $m,n\to\infty$.
Granica jest więc $H$-wartościowo ciągła i równa $u$
prawie wszędzie. Przechodząc do granicy w skończeniewymiarowej
tożsamości energii, korzystamy z jednostajnej zbieżności
wartości końcowych oraz zbieżności obu czynników w przestrzeniach
$L^2(V^*)$ i $L^2(V)$. To daje~\eqref{eq:pde20-energy-identity}.
Jednoznaczność reprezentanta wynika z ciągłości w $H$.
\end{proof}

\section{Istnienie i jednoznaczność: pełna konstrukcja Galerkina}
\begin{twierdzenie}[Paraboliczne zagadnienie Dirichleta]\label{thm:pde20-existence}
Niech $\Omega,T,A,f,u_0$ spełniają założenia
Definicji~\ref{def:pde20-data}. Wtedy istnieje dokładnie jedno rozwiązanie
słabe zagadnienia~\eqref{eq:pde20-eq} w sensie
Definicji~\ref{def:pde20-weak}. Spełnia ono oszacowanie
\begin{equation}\label{eq:pde20-apriori}
 \sup_{0\le t\le T}\|u(t)\|_H^2
 +\lambda\int_0^T\|u(t)\|_V^2\,\mathrm dt
 \le 2\|u_0\|_H^2
 +2\lambda^{-1}\int_0^T\|f(t)\|_{V^*}^2\,\mathrm dt.
\end{equation}
Stała w nierówności nie zależy od gładkości współczynników.
\end{twierdzenie}
\begin{proof}
\emph{Krok 1: równanie skończeniewymiarowe.}
Niech $V_n=\operatorname{span}\{e_1,\ldots,e_n\}$, gdzie $e_j$ są
funkcjami własnymi użytymi w Lemacie~\ref{lem:pde20-energy}.
Szukamy $u_n(t)=\sum_{j=1}^n c_j(t)e_j$, które spełnia
\[
 (u_n'(t),e_i)_H+a(t;u_n(t),e_i)
 =\langle f(t),e_i\rangle,\quad i=1,\ldots,n,
 \qquad u_n(0)=P_nu_0.
\]
Jest to liniowy układ $c'(t)+B_n(t)c(t)=F_n(t)$,
gdzie $B_n$ jest mierzalna i ograniczona na $[0,T]$, a
$F_n\in L^2(0,T;\R^n)$. Równanie całkowe
$c(t)=c(0)+\int_0^t(F_n(r)-B_n(r)c(r))\,\mathrm dr$
ma jednoznaczne absolutnie ciągłe rozwiązanie: na krótkim przedziale
operator po prawej jest kontrakcją w normie supremum, a skończoną
liczbę takich przedziałów sklejamy. Otrzymujemy rozwiązanie na $[0,T]$.

\emph{Krok 2: energia.}
Mnożymy $i$-te równanie przez $c_i(t)$ i sumujemy po $i$.
Otrzymujemy dla prawie każdego $t$
\[
 \frac12\frac{\mathrm d}{\mathrm dt}\|u_n\|_H^2
 +\lambda\|u_n\|_V^2
 \le\|f\|_{V^*}\|u_n\|_V
 \le\frac{\lambda}{2}\|u_n\|_V^2
 +\frac1{2\lambda}\|f\|_{V^*}^2.
\]
Całkując od $0$ do $t$ i korzystając z
$\|P_nu_0\|_H\le\|u_0\|_H$, dostajemy
\begin{equation}\label{eq:pde20-energy-n}
 \|u_n(t)\|_H^2+\lambda\int_0^t\|u_n(r)\|_V^2\,\mathrm dr
 \le\|u_0\|_H^2+\lambda^{-1}\|f\|_{L^2(0,T;V^*)}^2.
\end{equation}
To jednocześnie zapobiega utracie kontroli nad rozwiązaniami przy $n\to\infty$.

\emph{Krok 3: kontrola pochodnej czasowej.}
Dla każdego $v\in V$ zachodzi
$\langle u_n',v\rangle=\langle f,P_nv\rangle-a(t;u_n,P_nv)$,
ponieważ $u_n'\in V_n$ i $P_n$ jest ortogonalny w $H$.
Korzystamy z $\|P_nv\|_V\le\|v\|_V$ oraz ograniczoności formy:
\[
 \|u_n'(t)\|_{V^*}\le\|f(t)\|_{V^*}+\Lambda\|u_n(t)\|_V.
\]
Zatem $\{u_n'\}$ jest ograniczony w $L^2(0,T;V^*)$.

\emph{Krok 4: granica.}
Po wybraniu podciągu, dzięki refleksywności przestrzeni Hilberta,
$u_n\rightharpoonup u$ w $L^2(0,T;V)$,
$u_n'\rightharpoonup g$ w $L^2(0,T;V^*)$.
Ponadto $u_n$ jest ograniczony w $L^\infty(0,T;H)$.
Dla dowolnego $v\in V$ i $\eta\in C_c^\infty(0,T)$ mamy
\[
 \int_0^T\langle u_n'(t),v\rangle\eta(t)\,\mathrm dt
 =-\int_0^T(u_n(t),v)_H\eta'(t)\,\mathrm dt.
\]
Przechodząc do granicy, stwierdzamy, że $g=u'$.
Ustalmy $v=e_j$. Dla $n\ge j$ równanie Galerkina można
pomnożyć przez $\eta\in L^2(0,T)$; liniowość i ograniczoność
formy $w\mapsto a(t;w,e_j)$ pozwalają przejść do granicy
w całce po czasie. Otrzymujemy równanie słabe dla każdego $e_j$.
Przez gęstość ich rozpiętości w $V$ i ciągłość obu stron
równość rozszerza się na dowolny $v\in V$.

\emph{Krok 5: warunek początkowy.}
Dla $v\in V_n$, po scałkowaniu równania Galerkina od $0$ do $T$
i całkowaniu przez części z dowolną $\eta\in C^1([0,T])$,
$\eta(T)=0$, otrzymujemy
\[
 -\int_0^T(u_n(t),v)_H\eta'(t)\,\mathrm dt
 +\int_0^T a(t;u_n(t),v)\eta(t)\,\mathrm dt
 =\int_0^T\langle f(t),v\rangle\eta(t)\,\mathrm dt
 +(P_nu_0,v)_H\eta(0).
\]
Dla każdego ustalonego $v=e_j$ przechodzimy do granicy jak wyżej.
Odpowiadające równanie dla $u$ otrzymujemy z różniczkowania
jego ciągłego reprezentanta (Lemat~\ref{lem:pde20-energy});
porównanie wyrazów w czasie $0$ daje
$(u(0),e_j)_H=(u_0,e_j)_H$ dla każdego $j$.
Zatem $u(0)=u_0$ w $H$.

\emph{Krok 6: oszacowanie i jednoznaczność.}
Wstawiamy $v=u(t)$ do równania słabego (jest to dopuszczalne
w równości w $V^*$ dla prawie każdego $t$) i używamy
Lematu~\ref{lem:pde20-energy}. Dostajemy oszacowanie
\eqref{eq:pde20-apriori} dokładnie tak jak w Kroku 2.
Jeżeli $u,w$ są rozwiązaniami z tymi samymi danymi,
to $z=u-w$ spełnia równanie z zerową prawą stroną
i zerową wartością początkową. Tożsamość energii daje
$\|z(t)\|_H^2+2\lambda\int_0^t\|z\|_V^2\,\mathrm dr\le0$.
Stąd $z=0$.
\end{proof}

\begin{uwaga}[Zmienne w czasie współczynniki i brak symetrii]
W dowodzie nie różniczkowaliśmy $A$ względem czasu
ani nie odwoływaliśmy się do symetrii $a(t;\cdot,\cdot)$.
Jest to istotna przewaga słabego sformułowania nad próbami
obliczania $\Delta u$ punkt po punkcie.
\end{uwaga}

\section{Zależność od danych i paraboliczna zasada maksimum}
\begin{wniosek}[Stabilność rozwiązania]\label{cor:pde20-stability}
Jeżeli $u_i$ odpowiadają danym $(u_{0,i},f_i)$, $i=1,2$,
przy tej samej macierzy $A$, to dla $z=u_1-u_2$ mamy
\[
 \sup_{t\in[0,T]}\|z(t)\|_H^2
 +\lambda\|z\|_{L^2(0,T;V)}^2
 \le2\|u_{0,1}-u_{0,2}\|_H^2
 +2\lambda^{-1}\|f_1-f_2\|_{L^2(0,T;V^*)}^2.
\]
\end{wniosek}
\begin{proof}
Równanie na $z$ ma tę samą formę i odpowiadającą jej prawą
stronę $f_1-f_2$. Stosujemy Twierdzenie~\ref{thm:pde20-existence}.
\end{proof}

\begin{lemat}[Reguła obcięcia w czasie]\label{lem:pde20-trunc}
Jeżeli $u\in L^2(0,T;V)$ i $u'\in L^2(0,T;V^*)$, to
$u^-:=\max\{-u,0\}\in L^2(0,T;V)$ i funkcja
$t\mapsto\|u^-(t)\|_H^2$ jest absolutnie ciągła, przy czym
\[
 \frac12\frac{\mathrm d}{\mathrm dt}\|u^-(t)\|_H^2
 =-\langle u'(t),u^-(t)\rangle
 \quad\text{dla prawie wszystkich }t.
\]
\end{lemat}
\begin{proof}
Aby uniknąć nieuprawnionego różniczkowania funkcji ujemnej,
wybieramy funkcje wypukłe $\Phi_\eps\in C^2(\R)$ z
$\Phi_\eps(0)=\Phi_\eps'(0)=0$, tak że
$\Phi_\eps(s)\to(s^-)^2/2$,
$\Phi_\eps'(s)\to-s^-$, oraz $|\Phi_\eps'(s)|\le |s|$,
$|\Phi_\eps''(s)|\le C$.
Dla wygładzonych w czasie aproksymacji o wartościach w $V$
i dla funkcjonału $\int_\Omega\Phi_\eps(u)\,\mathrm d\lambda_d(x)$
reguła łańcuchowa daje
$\mathrm d/\mathrm dt\int_\Omega\Phi_\eps(u(t,x))\,\mathrm d\lambda_d(x)=
\langle u',\Phi_\eps'(u)\rangle$.
Ponieważ $\Phi_\eps'(0)=0$ i $|\Phi_\eps''|\le C$,
reguła łańcuchowa w przestrzeni Sobolewa daje
$\Phi_\eps'(v)\in V$ oraz
$\|\Phi_\eps'(v)\|_V\le C\|v\|_V$.
Uzasadnijmy jawnie ciągłość tej operacji w $V$:
przy $v_n\to v$ w $H_0^1$, po przejściu do podciągu
$v_n\to v$ i $\nabla v_n\to\nabla v$ prawie wszędzie.
Wtedy
$\nabla\Phi_\eps'(v_n)=\Phi_\eps''(v_n)\nabla v_n$
zbiega do $\Phi_\eps''(v)\nabla v$ w $L^2$,
bo pierwszy składnik różnicy kontroluje
$C\|\nabla v_n-\nabla v\|_2$, a w drugim
stosujemy zbieżność zdominowaną przy ustalonym
$\nabla v$.
To samo rozumowanie, najpierw na podciągu i następnie
dla całego ciągu, daje silną zbieżność po
wygładzeniu czasowym w $L^2(0,T;V)$.
Otrzymujemy więc regułę łańcuchową dla ustalonego
$\eps$ przez granicę w dualności $L^2(V^*)$--$L^2(V)$.
Przy $\eps\downarrow0$ zachodzi
$\Phi_\eps'(u)\to-u^-$ w $L^2(0,T;V)$:
poza poziomicą $\{u=0\}$ wynika to z reguły
łańcuchowej, a na $\{u=0\}$ mamy $\nabla u=0$
prawie wszędzie. Ponownie stosujemy zbieżność
zdominowaną i otrzymujemy żądaną tożsamość.
Otrzymujemy wskazaną tożsamość w sensie rozkładów w czasie,
a obie strony są całkowalne, więc w sensie funkcji
absolutnie ciągłych. (Jest to standardowy argument
reguły łańcuchowej dla Lipschitzowskiego obcięcia.)
\end{proof}

\begin{twierdzenie}[Zachowanie dodatniości]\label{thm:pde20-positivity}
Niech $u$ będzie rozwiązaniem z Twierdzenia~\ref{thm:pde20-existence}.
Jeżeli $u_0\ge0$ prawie wszędzie i $\langle f(t),v\rangle\ge0$
dla każdej nieujemnej $v\in V$ i prawie każdego $t$, to
$u(t,x)\ge0$ dla prawie każdego $x\in\Omega$ i każdego $t\in[0,T]$.
\end{twierdzenie}
\begin{proof}
Testujemy równanie przez $u^-(t)$. Z reguły łańcuchowej
$\langle u',u^-\rangle=-\frac12\mathrm d\|u^-\|_H^2/\mathrm dt$.
Na zbiorze $\{u<0\}$ mamy $\nabla u=-\nabla u^-$,
a na dopełnieniu $\nabla u^-=0$ prawie wszędzie. Zatem
$a(t;u,u^-)=-a(t;u^-,u^-)\le-\lambda\|u^-\|_V^2$.
W rezultacie
\[
 \frac12\frac{\mathrm d}{\mathrm dt}\|u^-(t)\|_H^2
 +\lambda\|u^-(t)\|_V^2\le0.
\]
Ponieważ $u^-(0)=0$, lewa strona po scałkowaniu daje
$\|u^-(t)\|_H^2=0$ dla każdego $t$.
\end{proof}

\begin{wniosek}[Zasada porównania]
Dla tej samej macierzy $A$ niech rozwiązania $u,v$ mają
odpowiednio dane $(u_0,f)$ i $(v_0,g)$.
Jeżeli $u_0\le v_0$ prawie wszędzie i
$\langle g(t)-f(t),\phi\rangle\ge0$ dla wszystkich
nieujemnych $\phi\in V$ i prawie wszystkich $t$,
to $u(t)\le v(t)$ prawie wszędzie dla każdego $t$.
\end{wniosek}
\begin{proof}
Stosujemy Twierdzenie~\ref{thm:pde20-positivity}
do różnicy $v-u$.
\end{proof}

\section{Wygładzanie i zanik energii dla laplasjanu Dirichleta}
Niech $L=-\Delta_D$ będzie dodatnim laplasjanem Dirichleta
w $H$. Jego wartości własne oznaczamy przez $\mu_j>0$, a funkcje
własne przez $e_j$. W tym podrozdziale zakładamy
$A=I$ i $f=0$, chyba że wyraźnie zaznaczymy inaczej.

\begin{twierdzenie}[Półgrupa cieplna i wygładzanie]\label{thm:pde20-semigroup}
Dla każdego $u_0\in H$ funkcja
$u(t)=e^{-tL}u_0=\sum_{j\ge1}e^{-\mu_jt}(u_0,e_j)_He_j$
jest jedynym rozwiązaniem słabym zagadnienia jednorodnego.
Dla każdego $s\ge0$ i $t>0$ zachodzi
\[
 \|L^se^{-tL}u_0\|_H
 \le C_s t^{-s}\|u_0\|_H,
 \qquad C_s=(s/\mathrm e)^s\ (s>0),\quad C_0=1.
\]
Ponadto $\|u(t)\|_H\le e^{-\mu_1t}\|u_0\|_H$.
\end{twierdzenie}
\begin{proof}
Dla skończonych sum własnych obliczamy
$u'(t)+Lu(t)=0$, a tożsamość energii jest bezpośrednia.
Dla ogólnego $u_0$ szereg zbiega w $H$ jednostajnie na
zwartych przedziałach czasu dzięki $|e^{-\mu_jt}|\le1$.
Mamy też
\[
 \int_0^T\|L^{1/2}e^{-tL}u_0\|_H^2\,\mathrm dt
 =\sum_j\mu_j|(u_0,e_j)_H|^2
 \int_0^T e^{-2\mu_jt}\,\mathrm dt
 \le\tfrac12\|u_0\|_H^2.
\]
Analogicznie $u'=-Lu$ należy do $L^2(0,T;V^*)$,
ponieważ $\|Lu\|_{V^*}=\|L^{1/2}u\|_H$.
Granica spełnia więc definicję rozwiązania słabego;
jednoznaczność wynika z Twierdzenia~\ref{thm:pde20-existence}.
Dla $s>0$ supremum funkcji $y\mapsto y^se^{-ty}$ na
$[0,\infty)$ wynosi $(s/(\mathrm et))^s$.
Przykładając to oszacowanie do każdego współczynnika szeregu,
dostajemy nierówność wygładzającą. Ponieważ $\mu_j\ge\mu_1$,
otrzymujemy także zanik wykładniczy.
\end{proof}


\begin{twierdzenie}[Wzór Duhamela dla wymuszenia w $L^2(H)$]
\label{thm:pde20-duhamel}
Niech $\Omega\subset\R^d$ będzie ograniczonym obszarem
Lipschitzowskim, $H=L^2(\Omega)$, $L=-\Delta_D$,
$u_0\in H$, $f\in L^2(0,T;H)$.
Jedynym rozwiązaniem słabym zagadnienia $u'+Lu=f$,
$u(0)=u_0$, jest funkcja
\[
 u(t)=e^{-tL}u_0+\int_0^t e^{-(t-s)L}f(s)\,\mathrm ds,
 \qquad 0\le t\le T,
\]
gdzie całka jest całką Bochnera w $H$.
\end{twierdzenie}
\begin{proof}
Operator $e^{-tL}$ jest kontrakcją w $H$. Dla niemal
każdego $t$ całka Bochnera istnieje, ponieważ
$\int_0^T\|f(s)\|_H\,\mathrm ds\le
\sqrt T\|f\|_{L^2(0,T;H)}$.
Ciągłość składnika całkowego w $H$ wynika z silnej
ciągłości półgrupy i zbieżności zdominowanej.

Niech $e_j$ będzie bazą funkcji własnych, $Le_j=\mu_je_j$.
Dla $u_n=P_nu$ i $f_n=P_nf$ wzór redukuje się do
$r_j'(t)+\mu_jr_j(t)=(f(t),e_j)_H$,
$r_j(0)=(u_0,e_j)_H$, którego rozwiązanie ma dokładnie
podaną postać. Dla różnic dwóch skończonych przybliżeń
stosujemy oszacowanie energii z
Twierdzenia~\ref{thm:pde20-existence}:
\[
 \sup_{0\le t\le T}\|u_n(t)-u_m(t)\|_H^2
 +\int_0^T\|u_n-u_m\|_V^2\,\mathrm dt
 \le C\bigl(\|P_nu_0-P_mu_0\|_H^2
 +\|P_nf-P_mf\|_{L^2(0,T;H)}^2\bigr),
\]
przy czym w prawej stronie korzystamy z ciągłego
zanurzenia $H\hookrightarrow V^*$. Prawa strona dąży do
zera. Jednocześnie z kontraktywności półgrupy wynika
jednostajna zbieżność wzoru całkowego w $H$.
Granica jest rozwiązaniem słabym, a jednoznaczność
daje Twierdzenie~\ref{thm:pde20-existence}.
\end{proof}

\begin{twierdzenie}[Wyższa regularność energii]\label{thm:pde20-max-reg}
Niech $L=-\Delta_D$, $u_0\in V=D(L^{1/2})$,
$f\in L^2(0,T;H)$. Wtedy rozwiązanie słabe równania
$u'+Lu=f$ spełnia
$u\in H^1(0,T;H)\cap L^2(0,T;D(L))$ oraz
\[
 \int_0^T\|u'(t)\|_H^2\,\mathrm dt
 +\sup_{t\in[0,T]}\|L^{1/2}u(t)\|_H^2
 \le2\|L^{1/2}u_0\|_H^2
 +2\int_0^T\|f(t)\|_H^2\,\mathrm dt.
\]
Dla obszarów klasy $C^{1,1}$ można ponadto utożsamić
$D(L)=H^2(\Omega)\cap H_0^1(\Omega)$; bez regularności brzegu
takiego utożsamienia nie zakładamy.
\end{twierdzenie}
\begin{proof}
Dla projekcji $P_n$ korzystamy z równania
$u_n'+Lu_n=P_nf$ z początkową daną $P_nu_0$.
Iloczyn skalarny z $u_n'$ daje
$\|u_n'\|_H^2+\frac12\frac{\mathrm d}{\mathrm dt}
\|L^{1/2}u_n\|_H^2=(P_nf,u_n')_H$.
Stosujemy nierówność
$|(P_nf,u_n')|\le(\|f\|_H^2+\|u_n'\|_H^2)/2$,
po czym całkujemy i otrzymujemy oszacowanie niezależne od $n$.
Ponieważ $Lu_n=P_nf-u_n'$, także $\{Lu_n\}$ jest ograniczone
w $L^2(0,T;H)$. Przechodzimy do słabej granicy i używamy
domkniętości operatora $L$ w przestrzeni Bochnera:
$u\in L^2(0,T;D(L))$, $Lu=f-u'$. Dolna półciągłość norm
oraz ciągłość czasowa $V$-wartościowego reprezentanta,
uzyskana tą samą argumentacją spektralną co w
Lemacie~\ref{lem:pde20-energy}, kończą dowód.
\end{proof}

\section{Zwartość czasowo-przestrzenna: lemat Aubina--Lionsa w naszym modelu}
\begin{twierdzenie}[Zwartość dla parabolicznych ciągów Galerkina]\label{thm:pde20-compact}
Niech $\Omega\subset\R^d$ będzie ograniczonym obszarem
Lipschitzowskim. Jeżeli ciąg $\{u_n\}$ jest ograniczony
w $L^2(0,T;V)$ i ciąg $\{u_n'\}$ jest ograniczony
w $L^2(0,T;V^*)$, to $\{u_n\}$ jest względnie zwarty
w $L^2(0,T;H)$.
\end{twierdzenie}
\begin{proof}
Niech $P_N$ oznacza stały rzut na pierwsze $N$ funkcji
własnych laplasjanu Dirichleta. Z rozwinięcia w bazie
$\{e_j\}$ dostajemy dla prawie każdego $t$
\[
 \|(I-P_N)u_n(t)\|_H^2
 \le\mu_{N+1}^{-1}\|u_n(t)\|_V^2.
\]
Całkowanie po czasie i wspólne ograniczenie norm $V$ dają
\begin{equation}\label{eq:pde20-tail}
 \sup_n\|(I-P_N)u_n\|_{L^2(0,T;H)}^2
 \le C\mu_{N+1}^{-1}\xrightarrow[N\to\infty]{}0.
\end{equation}
Dla ustalonego $N$ współczynniki
$c_{n,j}(t)=(u_n(t),e_j)_H$, $j\le N$, są ograniczone
w $L^2(0,T)$, a ich pochodne spełniają
$|c_{n,j}'(t)|=|\langle u_n'(t),e_j\rangle|
\le\|u_n'(t)\|_{V^*}\|e_j\|_V$.
Zatem dla każdego $j\le N$ ciąg $\{c_{n,j}\}_n$ jest
ograniczony w $H^1(0,T)$, a więc względnie zwarty
w $C([0,T])$: wynika to z oszacowania
$|c(t)-c(s)|\le\|c'\|_{L^2}|t-s|^{1/2}$
i z twierdzenia Arzelà--Ascoliego.
Przez wybór podciągu mamy zbieżność $P_Nu_n$
w $L^2(0,T;H)$ dla każdego $N$.
Wybór przekątny względem $N$, a potem oszacowanie
ogonów~\eqref{eq:pde20-tail}, pokazują, że wybrany podciąg
$\{u_{n_k}\}$ jest Cauchy'ego w $L^2(0,T;H)$.
\end{proof}


\begin{wniosek}[Szacunek przesunięć czasowych w normie $H$]
\label{cor:pde20-translations}
Niech $\Omega\subset\R^d$ będzie ograniczonym obszarem
Lipschitzowskim, $V=H_0^1(\Omega)$, $H=L^2(\Omega)$,
a $u\in L^2(0,T;V)$, $u'\in L^2(0,T;V^*)$.
Dla $0<h<T$ zachodzi
\[
 \int_0^{T-h}\|u(t+h)-u(t)\|_H^2\,\mathrm dt
 \le 2h\,\|u\|_{L^2(0,T;V)}
 \|u'\|_{L^2(0,T;V^*)}.
\]
W szczególności rodzina funkcji o jednostajnie
ograniczonych obu normach ma jednostajny moduł ciągłości
względem translacji czasu w $L^2(H)$.
\end{wniosek}
\begin{proof}
Z twierdzenia o ciągłości w trójce Gelfanda wynika
w $V^*$ równość
$u(t+h)-u(t)=\int_t^{t+h}u'(s)\,\mathrm ds$.
Dzięki nierówności Minkowskiego dla przesunięć oraz
przedłużeniu $u'$ przez zero poza $(0,T)$ mamy
\[
 \|u(\cdot+h)-u\|_{L^2(0,T-h;V^*)}
 \le h\|u'\|_{L^2(0,T;V^*)}.
\]
Z drugiej strony nierówność trójkąta daje
$\|u(\cdot+h)-u\|_{L^2(0,T-h;V)}
\le2\|u\|_{L^2(0,T;V)}$.
Z osadzeń $V\hookrightarrow H\hookrightarrow V^*$,
przy identyfikacji $H=H^*$, otrzymujemy
$\|w\|_H^2\le\|w\|_V\|w\|_{V^*}$.
Całkujemy tę nierówność dla $w(t)=u(t+h)-u(t)$,
a następnie stosujemy Cauchy'ego--Schwarza w czasie.
\end{proof}

\begin{uwaga}[Zastosowanie do nieliniowości]
Zbieżność słaba w $L^2(0,T;V)$ na ogół nie pozwala
przechodzić do granicy w wyrażeniu nieliniowym $F(u_n)$.
Zbieżność silna w $L^2(0,T;H)$ otrzymana powyżej
jest często brakującym elementem dowodów istnienia
w nieliniowych PDE. Ogólne twierdzenie Aubina--Lionsa
obejmuje więcej par przestrzeni i inne wykładniki;
tutaj udowodniliśmy samodzielnie wystarczający przypadek.
\end{uwaga}

\section{Nieliniowy kontrast: możliwość eksplozji}
\begin{twierdzenie}[Warunkowa eksplozja dla równania reakcji--dyfuzji]
\label{thm:pde20-blowup}
Niech $p>1$, $\Omega\subset\R^d$ będzie ograniczonym,
gładkim obszarem, a $\phi_1>0$ wewnątrz $\Omega$
będzie pierwszą funkcją własną laplasjanu Dirichleta,
$-\Delta\phi_1=\mu_1\phi_1$, z normalizacją
$\int_\Omega\phi_1\,\mathrm d\lambda_d(x)=1$.
Niech $u\ge0$ będzie dostatecznie regularnym rozwiązaniem
$u_t-\Delta u=u^p$ na $[0,T_*)\times\Omega$
z jednorodnym warunkiem Dirichleta. Jeżeli
$F(0):=\int_\Omega u(0,x)\phi_1(x)\,\mathrm d\lambda_d(x)
>\mu_1^{1/(p-1)}$, to nie istnieje globalne, nieujemne rozwiązanie klasyczne
na $[0,\infty)$ z funkcjonałem $F(t)$ skończonym w każdej chwili.
Dokładniej, jeśli $T_y<\infty$ jest czasem eksplozji rozwiązania
$y'=-\mu_1y+y^p$, $y(0)=F(0)$, to rozwiązanie $u$ nie może
pozostawać klasyczne na żadnym przedziale $[0,T]$ z $T\ge T_y$.
Jeśli istnieje aż do $T_y$, wówczas $F(t)\to\infty$ przy $t\uparrow T_y$.
\end{twierdzenie}
\begin{proof}
Mnożąc równanie przez $\phi_1$ i całkując po $\Omega$,
a następnie korzystając z całkowania przez części
i z warunków Dirichleta, otrzymujemy
\[
 F'(t)=-\mu_1F(t)+\int_\Omega u(t,x)^p\phi_1(x)
 \,\mathrm d\lambda_d(x).
\]
Miara $\phi_1(x)\,\mathrm d\lambda_d(x)$ jest probabilistyczna,
więc nierówność Jensena daje
$\int_\Omega u^p\phi_1\,\mathrm d\lambda_d\ge F(t)^p$.
Stąd $F'\ge-\mu_1F+F^p$.
Dla $F(0)>\mu_1^{1/(p-1)}$ rozwiązanie skalarnego ODE
$y'=-\mu_1y+y^p$ rośnie i osiąga nieskończoność
w skończonym czasie, ponieważ dla dużych $y$
mamy $y^p-\mu_1y\ge y^p/2$, a
$\int_{y_0}^\infty 2s^{-p}\,\mathrm ds<\infty$.
Z zasady porównania dla skalarnych ODE
$F(t)\ge y(t)$, dopóki oba wyrażenia są skończone.
To wyklucza globalne ograniczone rozwiązanie wskazanego typu.
\end{proof}

\section{Stabilność wstecznego Eulera i rozdzielenie czasu od przestrzeni}
Dla dyskretyzacji laplasjanu Dirichleta na siatce jednowymiarowej
niech $L_h$ będzie symetryczną dodatnio określoną macierzą
różnicową. Przy kroku czasowym $\tau>0$ schemat
wstecznego Eulera ma postać
\begin{equation}\label{eq:pde20-implicit}
 \frac{U^{m+1}-U^m}{\tau}+L_hU^{m+1}=F^{m+1}.
\end{equation}
\begin{stwierdzenie}[Bezwarunkowa stabilność energetyczna]
\label{prop:pde20-be}
Jeżeli $F^{m+1}=0$ dla wszystkich $m$, to rozwiązanie
schematu~\eqref{eq:pde20-implicit} spełnia
$\|U^{m+1}\|_h\le\|U^m\|_h$ niezależnie od $\tau/h^2$,
gdzie $\|U\|_h^2=h\sum_j|U_j|^2$.
\end{stwierdzenie}
\begin{proof}
Bierzemy dyskretny iloczyn skalarny równania przez $U^{m+1}$.
Tożsamość $2(a-b,a)=\|a\|^2-\|b\|^2+\|a-b\|^2$
daje
\[
 \|U^{m+1}\|_h^2-\|U^m\|_h^2
 +\|U^{m+1}-U^m\|_h^2
 +2\tau(L_hU^{m+1},U^{m+1})_h=0.
\]
Ostatnie dwa składniki są nieujemne; stąd teza.
Jest to twierdzenie o stabilności, nie samoistny dowód
zbieżności: potrzebna jest jeszcze kontrola błędu aproksymacji.
\end{proof}

\begin{uwaga}[Granice bieżącego rozdziału]
Dla współczynników zmiennych w czasie udowodniliśmy
istnienie i jednoznaczność rozwiązania w naturalnej klasie
$L^2(V)\cap H^1(V^*)$ bez gładkości $A$. Nie dowodziliśmy
tu pełnej parabolicznej teorii De Giorgiego--Nasha--Mosera,
teorii oszacowań Schaudera ani maksymalnej regularności
w klasach $L^p$ dla dowolnych operatorów niesymetrycznych.
Odpowiednie zagadnienia wymagają odrębnych założeń i metod.
\end{uwaga}

\chapter[Układy hiperboliczne i propagacja]{Równania hiperboliczne: układy symetryczno-hiperboliczne, energia i propagacja informacji}
\label{chap:pde21}

\section{Dlaczego potrzebujemy układów pierwszego rzędu?}
Rozdział~5 dotyczył klasycznego równania falowego, przede wszystkim
jego jawnych formuł. Obecnie rozważamy układy o zmiennych współczynnikach,
bez dostępnej formuły rozwiązania. Ich centralną własnością będzie możliwość
kontrolowania energii bez tracenia pochodnych. Pozwoli to oddzielić
\emph{hiperboliczność symbolu} od \emph{poprawności zagadnienia Cauchy'ego}.

\begin{uwaga}[Dwie zasady energetyczne: dyfuzja i propagacja]
\label{rem:pde-v43-energy-types}
Oszacowanie paraboliczne z
Twierdzenia~\ref{thm:pde20-existence} kontroluje także
całkę w czasie z normy przestrzennej $H_0^1$ dzięki
koercywności operatora dyfuzji. W układzie
symetryczno-hiperbolicznym energia kontroluje przede
wszystkim normę $L^2$ w chwili $t$, natomiast jej
lokalny bilans daje stożki zależności. Nie wolno
wywnioskować wygładzania przestrzennego z samej
nierówności energetycznej dla układu hiperbolicznego.
\end{uwaga}

Niech $d,m\ge1$ i $T>0$. Dla kolumny $u=(u_1,\ldots,u_m)^\mathsf T$
rozważmy na $[0,T]\times\R^d$ równanie
\begin{equation}\label{eq:pde21-main}
 \partial_tu+\sum_{j=1}^d A_j(t,x)\partial_{x_j}u+B(t,x)u=f(t,x),
 \qquad u(0,x)=u_0(x).
\end{equation}
Macierze $A_j,B$ mają rozmiar $m\times m$ i rzeczywiste współczynniki.
Dla funkcji wektorowych przyjmujemy $\|u\|_2^2=\int_{\R^d}|u(x)|^2\,\mathrm d\lambda_d(x)$.

\begin{definicja}[Układ symetryczno-hiperboliczny]\label{def:pde21-sym}
Układ~\eqref{eq:pde21-main} nazywamy \emph{symetryczno-hiperbolicznym},
jeżeli $A_j(t,x)=A_j(t,x)^\mathsf T$ dla wszystkich $j$ i wszystkich
$(t,x)$ (lub prawie wszędzie, gdy pracujemy ze współczynnikami mierzalnymi).
Jego symbolem głównym jest $P(t,x;\tau,\xi)=\tau I_m+\sum_j\xi_jA_j(t,x)$.
Kowektor $(\tau,\xi)\ne0$ jest charakterystyczny, jeżeli
$\det P(t,x;\tau,\xi)=0$.
\end{definicja}

\begin{stwierdzenie}[Rzeczywiste prędkości charakterystyczne]
\label{prop:pde21-real}
Jeżeli $A_j=A_j^\mathsf T$, to dla każdego $(t,x)$ oraz
$\xi\in\R^d\setminus\{0\}$ wszystkie pierwiastki wielomianu
$\tau\mapsto\det(\tau I_m+\sum_j\xi_jA_j)$ są rzeczywiste.
Płaszczyzna $t=\mathrm{const}$ jest niecharakterystyczna.
\end{stwierdzenie}
\begin{proof}
Macierz $A(\xi)=\sum_j\xi_jA_j$ jest rzeczywista i symetryczna.
Z twierdzenia spektralnego istnieje baza ortonormalna jej wektorów
własnych i rzeczywiste liczby $\lambda_1(\xi),\ldots,\lambda_m(\xi)$.
W tej bazie $\det(\tau I_m+A(\xi))=\prod_k(\tau+\lambda_k(\xi))$.
Dla konormalnej do płaszczyzny $t=\mathrm{const}$ mamy $(\tau,\xi)=(1,0)$,
a zatem $P=I_m$ jest odwracalna.
\end{proof}

\begin{uwaga}
Symetryczna hiperboliczność nie wymaga, aby wartości własne symbolu były
różne. Nie należy jej utożsamiać ze ścisłą hiperbolicznością. Co więcej,
sama rzeczywistość pierwiastków wielomianu charakterystycznego nie
zapewnia oszacowania energetycznego; zobaczymy to w
Kontrprzykładzie~\ref{ex:pde21-jordan}.
\end{uwaga}

\section{Tożsamość energii i zależność od danych}
\begin{lemat}[Lokalny bilans energii]\label{lem:pde21-local}
Niech $A_j\in C^1([0,T]\times\R^d;\R^{m\times m})$ będą symetryczne,
$B$ ciągła, a $u\in C^1([0,T]\times\R^d;\R^m)$ rozwiązuje
\eqref{eq:pde21-main}. Wówczas punktowo zachodzi
\begin{equation}\label{eq:pde21-flux}
 \partial_t|u|^2+\sum_{j=1}^d\partial_{x_j}(u^\mathsf TA_ju)
 =u^\mathsf T\!\left(\sum_{j=1}^d\partial_{x_j}A_j-B-B^\mathsf T\right)u
 +2u^\mathsf Tf.
\end{equation}
\end{lemat}
\begin{proof}
Mnożymy równanie przez $2u^\mathsf T$. Otrzymujemy
$2u^\mathsf Tu_t+2\sum_j u^\mathsf TA_ju_{x_j}+2u^\mathsf TBu=2u^\mathsf Tf$.
Pierwszy składnik jest $\partial_t|u|^2$. Symetria $A_j$ daje
$\partial_{x_j}(u^\mathsf TA_ju)=2u^\mathsf TA_ju_{x_j}
+u^\mathsf T(\partial_{x_j}A_j)u$.
Wreszcie $2u^\mathsf TBu=u^\mathsf T(B+B^\mathsf T)u$.
Przenosząc odpowiednie wyrazy, otrzymujemy~\eqref{eq:pde21-flux}.
\end{proof}

\begin{twierdzenie}[Oszacowanie energii na całej przestrzeni]\label{thm:pde21-energy}
Niech $A_j\in W^{1,\infty}((0,T)\times\R^d)$ będą symetryczne,
$B\in L^\infty((0,T)\times\R^d)$, a $u$ będzie gładkim rozwiązaniem
\eqref{eq:pde21-main} o dostatecznym zaniku przy $|x|\to\infty$.
Niech $M=\|\sum_j\partial_{x_j}A_j-B-B^\mathsf T\|_{L^\infty_{t,x};\mathrm{op}}$.
Wtedy dla każdego $t\in[0,T]$
\begin{equation}\label{eq:pde21-L2}
 \|u(t)\|_2\le e^{Mt/2}\left(\|u_0\|_2+
 \int_0^t\|f(r)\|_2\,\mathrm dr\right).
\end{equation}
To samo oszacowanie zachodzi dla różnicy dwóch rozwiązań, z prawą
stroną równą różnicy wymuszeń.
\end{twierdzenie}
\begin{proof}
Całkujemy~\eqref{eq:pde21-flux} względem $x$.
Wyraz dywergencji znika dla rozwiązań o zwartym nośniku; przy
założonym zaniku otrzymujemy to samo, stosując funkcję odcinającą
$\chi_R$ i przechodząc z $R\to\infty$.
Po oznaczeniu $E(t)=\|u(t)\|_2^2$ mamy
$E'(t)\le ME(t)+2\sqrt{E(t)}\|f(t)\|_2$.
Aby nie dzielić przez zero, definiujemy $y_\eps(t)=\sqrt{E(t)+\eps^2}$.
Wówczas $y_\eps'\le(M/2)y_\eps+\|f\|_2$ prawie wszędzie.
Mnożymy przez $e^{-Mt/2}$, całkujemy od $0$ do $t$, a potem
przechodzimy z $\eps\downarrow0$. Otrzymujemy~\eqref{eq:pde21-L2}.
Dla różnicy dwóch rozwiązań powtarzamy identyczny rachunek.
\end{proof}

\begin{wniosek}[Jednoznaczność]
W klasie rozwiązań spełniających oszacowanie energii~\eqref{eq:pde21-L2}
zagadnienie Cauchy'ego dla~\eqref{eq:pde21-main} ma co najwyżej jedno rozwiązanie.
\end{wniosek}
\begin{proof}
Dla różnicy dwóch rozwiązań z tymi samymi danymi mamy zerowe
dane początkowe i zerową prawą stronę. Stosujemy~\eqref{eq:pde21-L2}.
\end{proof}

\section{Stożki zależności i skończona prędkość propagacji}
\begin{definicja}[Maksymalna prędkość symbolu]
Dla symetrycznych $A_j$ definiujemy
$c_*:=\operatorname*{ess\,sup}_{(t,x)\in(0,T)\times\R^d}
\sup_{|\nu|=1}\|\sum_j\nu_jA_j(t,x)\|_{\mathrm{op}}$.
W dalszych rozważaniach zakładamy $c_*<\infty$.
\end{definicja}

\begin{twierdzenie}[Stożkowa zasada zależności]\label{thm:pde21-cone}
Niech $A_j\in C^1_b([0,T]\times\R^d)$ będą symetryczne,
$B\in C_b([0,T]\times\R^d)$ i niech $u$ będzie klasycznym
rozwiązaniem~\eqref{eq:pde21-main}. Ustalmy $x_0\in\R^d$,
$t_0\in(0,T]$, $R\ge0$ i $c\ge c_*$. Jeśli
$u(0,x)=0$ dla $|x-x_0|\le R+ct_0$ oraz
$f(t,x)=0$ dla $0\le t\le t_0$, $|x-x_0|\le R+c(t_0-t)$,
to $u(t_0,x)=0$ dla $|x-x_0|\le R$.
Dla $R=0$ wniosek dotyczy punktu $x_0$.
\end{twierdzenie}
\begin{proof}
Najpierw załóżmy $R>0$. Niech $r(t)=R+c(t_0-t)$
i $D(t)=B_{r(t)}(x_0)$. Funkcja
$E_D(t)=\int_{D(t)}|u(t,x)|^2\,\mathrm d\lambda_d(x)$
jest różniczkowalna. Wzór Reynoldsa dla poruszającej się kuli
daje
\[
 E_D'(t)=\int_{D(t)}\partial_t|u|^2\,\mathrm d\lambda_d
 +r'(t)\int_{\partial D(t)}|u|^2\,\mathrm d\mathcal H^{d-1}.
\]
Wstawiamy~\eqref{eq:pde21-flux} i twierdzenie Gaussa.
Na brzegu pojawia się wyraz
$-u^\mathsf TA_\nu u-c|u|^2$, gdzie
$A_\nu=\sum_j\nu_jA_j$ i $\nu$ jest normalną zewnętrzną.
Ponieważ $|u^\mathsf TA_\nu u|\le\|A_\nu\|_{\mathrm{op}}|u|^2$
oraz $c\ge c_*$, całka brzegowa jest niedodatnia.
Wnętrze daje oszacowanie $E_D'(t)\le M E_D(t)$, ponieważ
$f=0$ w rozważanym stożku. Z warunku $E_D(0)=0$
i lematu Gronwalla wynika $E_D(t)=0$ dla $0\le t\le t_0$.
Ciągłość $u$ pozwala przejść od zerowania prawie wszędzie
do zerowania punktowego. Dla $R=0$ i $c>0$ przeprowadzamy ten sam rachunek na $[0,t_0)$:
promień $r(t)=c(t_0-t)$ jest tam dodatni, a zerowa energia daje
$u(t,x_0)=0$ dla wszystkich $t<t_0$. Ciągłość w czasie pozwala
przejść do $t=t_0$. Jeżeli $c=0$, to również $c_*=0$, a więc
$A_j=0$ i równanie w każdym punkcie jest zwykłym równaniem
różniczkowym; wniosek wynika wprost z jednoznaczności ODE.
\end{proof}



\begin{center}
\begin{tikzpicture}[scale=.88,>=Latex]
\fill[blue!8] (0.8,0) -- (6.0,0) -- (3.4,3.55) -- cycle;
\draw[->] (0,0) -- (6.8,0) node[right] {$x$};
\draw[->] (0,0) -- (0,4.2) node[above] {$t$};
\draw[thick,blue!65!black] (.8,0)--(3.4,3.55)--(6,0);
\draw[dashed] (1.7,1.2)--(5.1,1.2);
\node[left] at (0,3.55) {$t_0$};
\node[below] at (3.4,0) {$x_0$};
\node[above right] at (3.4,3.55) {$(t_0,x_0)$};
\node at (3.4,.5) {dane i wymuszenie w stożku};
\node[right] at (5.1,1.2) {$t=t_1$};
\end{tikzpicture}

\small Rysunek: wsteczny stożek zależności w jednym wymiarze.
Gdy $R=0$, podstawa ma promień $ct_0$.
\end{center}

\section{Równanie falowe jako układ symetryczny}
Niech $c>0$ i $u_{tt}-c^2u_{xx}=0$. Określamy $v=u_t$,
$w=cu_x$ i $U=(v,w)^\mathsf T$. Wówczas
\begin{equation}\label{eq:pde21-wavefirst}
 U_t+\begin{pmatrix}0&-c\\-c&0\end{pmatrix}U_x=0,
 \qquad |U|^2=u_t^2+c^2u_x^2.
\end{equation}
Macierz główna jest symetryczna, z wartościami własnymi $\pm c$.
Są to prędkości przenoszenia dwóch niezależnych modów.

\begin{stwierdzenie}[Rozkład na fale podróżujące]
Niech $v,w$ będą gładkimi funkcjami rozwiązującymi
\eqref{eq:pde21-wavefirst} na $[0,T]\times\R$. Funkcje
$r_+=v-w$ i $r_-=v+w$ spełniają odpowiednio
$(\partial_t+c\partial_x)r_+=0$ oraz
$(\partial_t-c\partial_x)r_-=0$.
W konsekwencji $r_+(t,x)=r_+(0,x-ct)$ i
$r_-(t,x)=r_-(0,x+ct)$.
\end{stwierdzenie}
\begin{proof}
Równania składowych to $v_t=cw_x$ i $w_t=cv_x$.
Po odjęciu otrzymujemy $(v-w)_t=-c(v-w)_x$,
a po dodaniu $(v+w)_t=c(v+w)_x$. Pozostały wniosek jest
rozwiązaniem liniowego równania transportu wzdłuż charakterystyk.
\end{proof}

\section{Oszacowania wyższych pochodnych bez ich utraty}
\begin{lemat}[Komutator z pochodną wielokrotną]\label{lem:pde21-comm}
Niech $s\ge1$ będzie liczbą całkowitą,
$A\in W^{s,\infty}(\R^d;\R^{m\times m})$ oraz
$u\in H^s(\R^d;\R^m)$. Dla każdego wielowskaźnika
$\alpha$ z $|\alpha|\le s$ zachodzi
\[
 \|[\partial^\alpha,A]\partial_j u\|_2
 \le C_{d,s}\|A\|_{W^{s,\infty}}\|u\|_{H^s},
\]
gdzie komutator oznacza
$[\partial^\alpha,A]v=\partial^\alpha(Av)-A\partial^\alpha v$.
\end{lemat}
\begin{proof}
Dla gładkich funkcji reguła Leibniza daje
$[\partial^\alpha,A]\partial_ju=
\sum_{0<\beta\le\alpha}\binom{\alpha}{\beta}
(\partial^\beta A)\partial_j\partial^{\alpha-\beta}u$.
Dla każdego składnika mamy
$1+|\alpha|-|\beta|\le|\alpha|\le s$,
ponieważ $|\beta|\ge1$. Stosując oszacowanie
$\|Cv\|_2\le\|C\|_\infty\|v\|_2$, sumując skończoną
liczbę składników i następnie aproksymując dane funkcjami
gładkimi, otrzymujemy tezę.
\end{proof}

\begin{twierdzenie}[Energia w $H^s$]\label{thm:pde21-Hs}
Niech $s\ge1$ będzie liczbą całkowitą,
$A_j\in C([0,T];W^{s,\infty}(\R^d))$ będą symetryczne,
$B\in C([0,T];W^{s,\infty}(\R^d))$, a ponadto
$\partial_{x_j}A_j\in L^\infty$.
Jeżeli $u$ jest gładkim rozwiązaniem~\eqref{eq:pde21-main}, to
\begin{equation}\label{eq:pde21-Hs}
 \|u(t)\|_{H^s}\le e^{Ct}
 \left(\|u_0\|_{H^s}+\int_0^t\|f(r)\|_{H^s}\,\mathrm dr\right),
\end{equation}
gdzie $C$ zależy od $d,m,s$ i jednostajnych norm
$W^{s,\infty}$ współczynników, nie zaś od $u$ i $f$.
\end{twierdzenie}
\begin{proof}
Niech $|\alpha|\le s$. Po zastosowaniu $\partial^\alpha$
do~\eqref{eq:pde21-main} otrzymujemy
\[
 \partial_t(\partial^\alpha u)+\sum_jA_j\partial_j(\partial^\alpha u)
 +B\partial^\alpha u
 =\partial^\alpha f-\sum_j[\partial^\alpha,A_j]\partial_ju
 -[\partial^\alpha,B]u.
\]
Dla każdego $\alpha$ stosujemy oszacowanie
\eqref{eq:pde21-L2} (lub bezpośrednio różniczkową nierówność
na kwadrat normy), a komutatory ograniczamy Lematem~\ref{lem:pde21-comm}
i regułą Leibniza dla $B$. Sumując względem $|\alpha|\le s$,
dostajemy dla $E_s=\sum_{|\alpha|\le s}\|\partial^\alpha u\|_2^2$
nierówność $E_s'\le C_1 E_s+2\sqrt{E_s}\,\|f\|_{H^s}$.
Argument z $\sqrt{E_s+\eps^2}$ i Gronwallem, opisany w dowodzie
Twierdzenia~\ref{thm:pde21-energy}, prowadzi do~\eqref{eq:pde21-Hs}.
\end{proof}


\begin{lemat}[Przedłużenie tożsamości energii na rozwiązania słabe]
\label{lem:pde21-friedrichs-comm}
Niech $A_j\in W^{1,\infty}(\R^d;\R^{m\times m})$ będą symetryczne,
$B\in L^\infty$, wszystkie współczynniki będą mierzalne w czasie
z jednostajnymi podanymi ograniczeniami. Jeżeli
$v\in L^\infty(0,T;L^2)$ rozwiązuje dystrybucyjnie
$v_t+\sum_j A_j\partial_jv+Bv=F$ z $F\in L^1(0,T;L^2)$,
to ma reprezentanta z $C([0,T];L^2)$ i zachodzi dla niego
energetyczna tożsamość otrzymana przez scałkowanie
\eqref{eq:pde21-flux}. W szczególności spełnia
oszacowanie~\eqref{eq:pde21-L2} z dowolnej chwili początkowej.
\end{lemat}
\begin{proof}
Niech $J_\eps v=\rho_\eps*v$ będzie splotem tylko względem $x$,
$\rho\in C_c^\infty(\R^d)$, $\int_{\R^d}\rho(x)\,\mathrm d\lambda_d(x)=1$. Wygładzona funkcja
$v_\eps=J_\eps v$ spełnia równanie z prawą stroną
$J_\eps F+R_\eps$, gdzie $R_\eps$ jest sumą komutatorów
$[A_j,J_\eps]\partial_jv$ i $[B,J_\eps]v$.
Dla każdego ustalonego $t$ (pomijamy go w zapisie)
\begin{align*}
 [J_\eps,A]\partial_jv(x)
 &=\int_{\R^d}\partial_j\rho_\eps(z)
  (A(x-z)-A(x))(v(x-z)-v(x))\,\mathrm d\lambda_d(z)\\
 &\quad+\int_{\R^d}\rho_\eps(z)
 (\partial_j A)(x-z)(v(x)-v(x-z))\,\mathrm d\lambda_d(z).
\end{align*}
Wzór otrzymujemy dwukrotnie całkując przez części w sensie dystrybucji;
wyrazy z $v(x)$ w drugiej całce znoszą część pierwszej.
Ponieważ $|A(x-z)-A(x)|\le\|\nabla A\|_\infty|z|$,
norma $L^2$ prawej strony jest ograniczona przez
\[
 C\|\nabla A\|_\infty
 \sup_{|z|\le\eps R_\rho}\|v(\cdot-z)-v\|_2,
\]
gdzie $R_\rho$ jest promieniem nośnika $\rho$. Wyrażenie to
zbiega do zera dla każdej funkcji $v\in L^2$, na mocy
ciągłości translacji w $L^2$. Jest również zdominowane przez
$C\|v\|_2$, co pozwala zastosować zbieżność zdominowaną
w czasie. Podobnie $[B,J_\eps]v\to0$ w $L^1_tL^2_x$
(bez potrzeby różniczkowania $B$).
Dla ustalonego $\eps$ równanie wygładzone ma wszystkie wyrazy
w $L^1_tL^2_x$, zatem $v_\eps\in W^{1,1}(0,T;L^2)$
i można zastosować klasyczną tożsamość energii.
Po przejściu z $\eps\downarrow0$ komutatory znikają,
a pozostałe całki zbiegają dla prawie wszystkich końców przedziału.
Otrzymana absolutnie ciągła funkcja normy $\|v(t)\|_2^2$
ma ciągłe przedłużenie na cały przedział.
Ponieważ samo równanie zapewnia słabą ciągłość $v$ w $L^2$
(testujemy przeciw gładkim, gęstym funkcjom przestrzennym),
a ciągłość normy i słaba ciągłość w przestrzeni Hilberta
implikują ciągłość silną, dostajemy $v\in C([0,T];L^2)$.
Powtórzenie rachunku z $v_\eps$ oraz użycie Gronwalla
daje również oszacowanie energii.
\end{proof}

\section{Istnienie przez regularizację Friedricha}
Warto odróżnić metodę zastosowaną poniżej od rozdziału~20:
w równaniach parabolicznych otrzymywaliśmy zwartość z całki po czasie
normy $H^1$. Tutaj podstawowym narzędziem jest kontrola \emph{wszystkich}
pochodnych przestrzennych dzięki symetrii macierzy głównych.

\begin{twierdzenie}[Zagadnienie Cauchy'ego dla gładkich współczynników]
\label{thm:pde21-existence}
Niech $s>d/2+2$ będzie liczbą całkowitą. Załóżmy, że
$A_j=A_j^\mathsf T$ oraz
$A_j,B\in C([0,T];W^{s+1,\infty}(\R^d;\R^{m\times m}))$,
zaś $u_0\in H^s(\R^d;\R^m)$ i
$f\in C([0,T];H^s(\R^d;\R^m))$.
Wtedy istnieje dokładnie jedno rozwiązanie zagadnienia
\eqref{eq:pde21-main}, dla którego
\[
 u\in C([0,T];H^s(\R^d;\R^m))\cap
 C^1([0,T];H^{s-1}(\R^d;\R^m)).
\]
Spełnia ono~\eqref{eq:pde21-Hs}.
\end{twierdzenie}
\begin{proof}
\emph{Krok 1: regularizacja.}
Niech $P_N$ będzie ortogonalnym rzutem Fouriera na częstotliwości
$|\xi|\le N$. Operator $P_N$ jest kontrakcją w każdym $H^k$
i komutuje z pochodnymi. Szukamy rozwiązania
\begin{equation}\label{eq:pde21-friedrichs}
 u_N'+P_N\left(\sum_j A_j\partial_j u_N+Bu_N\right)=P_Nf,
 \quad u_N(0)=P_Nu_0,
 \quad P_Nu_N=u_N.
\end{equation}
Na podprzestrzeni $P_NH^s$ operator
$v\mapsto P_N(A_j\partial_jv+Bv)$ jest ograniczony, ponieważ
$\|\partial_jv\|_{H^s}\le N\|v\|_{H^s}$.
Równanie~\eqref{eq:pde21-friedrichs} jest więc zwykłym
liniowym równaniem różniczkowym w przestrzeni Banacha, o ciągłych
w czasie współczynnikach; metoda kolejnych przybliżeń daje
rozwiązanie na krótkim przedziale, a liniowość pozwala je
przedłużyć do $T$.

\emph{Krok 2: jednostajne oszacowania.}
Ponieważ $P_N$ jest samosprzężony i $P_N\partial^\alpha u_N=\partial^\alpha u_N$,
rachunek energii po zastosowaniu $\partial^\alpha$ wygląda tak
samo jak w dowodzie~\eqref{eq:pde21-Hs}. Istotnie
$\langle\partial^\alpha u_N,P_N(A_j\partial_j\partial^\alpha u_N)\rangle
=\langle\partial^\alpha u_N,A_j\partial_j\partial^\alpha u_N\rangle$.
Komutatory są oszacowane przez Lemat~\ref{lem:pde21-comm}.
Stąd $\sup_{N,t}\|u_N(t)\|_{H^s}\le C$.
Z~\eqref{eq:pde21-friedrichs} i ograniczoności mnożenia
przez $W^{s,\infty}$ dostajemy także
$\sup_{N,t}\|u_N'(t)\|_{H^{s-1}}\le C$.

\emph{Krok 3: zwartość lokalna.}
Na każdej kuli $B_R$ zanurzenie $H^s(B_R)\hookrightarrow
H^{s-1}(B_R)$ jest zwarte. Ponadto
$\|u_N(t)-u_N(r)\|_{H^{s-1}(B_R)}\le C|t-r|$.
Z twierdzenia Arzelà--Ascoliego w metrycznej przestrzeni
$H^{s-1}(B_R)$ otrzymujemy podciąg zbieżny w
$C([0,T];H^{s-1}(B_R))$. Wybierając podciąg przekątny dla
$R=1,2,\ldots$, dostajemy funkcję $u$, która jest granicą lokalną
oraz słabą-* granicą w $L^\infty(0,T;H^s(\R^d))$.

\emph{Krok 4: równanie i dane.}
Dla $\varphi\in C_c^\infty((0,T)\times\R^d;\R^m)$
mnożymy~\eqref{eq:pde21-friedrichs} przez $\varphi$
i całkujemy. Przenosimy $P_N$ na funkcję testową
$P_N\varphi$; ponieważ $P_N\varphi\to\varphi$ w każdym
$H^k$, przechodzimy do granicy w wyrazach liniowych,
korzystając z lokalnej silnej zbieżności oraz jednostajnych
oszacowań na całej przestrzeni. Otrzymujemy~\eqref{eq:pde21-main}
w sensie dystrybucyjnym. Zbieżność
$u_N(0)=P_Nu_0\to u_0$ w $H^{s-1}$ i lokalna ciągłość implikują
$u(0)=u_0$ jako element $H^{s-1}(\R^d)$.

\emph{Krok 5: energia, jednoznaczność i ciągłość silna.}
Dla każdego wielowskaźnika $|\alpha|\le s$ funkcja
$v=\partial^\alpha u$ należy do $L^\infty(0,T;L^2)$.
Po zróżniczkowaniu równania jej prawa strona jest sumą
$\partial^\alpha f$ i komutatorów z
Lematu~\ref{lem:pde21-comm}, a więc należy do $L^1(0,T;L^2)$.
Lemat~\ref{lem:pde21-friedrichs-comm} stosuje się zatem
osobno do wszystkich $v=\partial^\alpha u$; daje
$v\in C([0,T];L^2)$ i po zsumowaniu oszacowanie
~\eqref{eq:pde21-Hs}. Dzięki temu $u\in C([0,T];H^s)$
bez konieczności zakładania z góry ciągłości najwyższych pochodnych.
Równanie ma prawą stronę ciągłą z wartościami w $H^{s-1}$,
ponieważ mnożenie przez przyjęte współczynniki jest ciągłe
na $H^{s-1}$, zatem $u_t\in C([0,T];H^{s-1})$.
Ostatecznie $u\in C^1([0,T];H^{s-1})$.
Jednoznaczność wynika z Twierdzenia~\ref{thm:pde21-energy}
zastosowanego do różnicy dwóch rozwiązań.
\end{proof}

\begin{uwaga}[Wersja o mniejszej regularności]
Następnie można przez aproksymację danych w $L^2$
i oszacowanie~\eqref{eq:pde21-L2} otrzymać rozwiązania energetyczne
z ciągłością $C([0,T];L^2)$, ale wymaga to osobnego uzasadnienia
reguły energii dla granicznych rozwiązań. Twierdzenie powyżej
celowo sformułowaliśmy z mocniejszymi założeniami i bez ukrytej
utraty pochodnych.
\end{uwaga}

\section{Symetryzatory: gdy macierze nie są symetryczne}
\begin{definicja}[Układ symetryzowalny]\label{def:pde21-symmetrizer}
Niech $A_j$ będą rzeczywistymi macierzami $m\times m$.
Mówimy, że układ~\eqref{eq:pde21-main} ma \emph{symetryzator} $S(t,x)$,
jeżeli $S=S^\mathsf T$, istnieją stałe $0<\sigma\le\Sigma$ takie,
że $\sigma I_m\le S\le\Sigma I_m$, a macierze $SA_j$ są symetryczne.
Dla oszacowań energii zakładamy dodatkowo
$S\in W^{1,\infty}$ i $A_j,B\in W^{1,\infty}$.
\end{definicja}

\begin{twierdzenie}[Energia ważona symetryzatorem]
\label{thm:pde21-symmetrizer}
Niech założenia Definicji~\ref{def:pde21-symmetrizer} zachodzą
na $[0,T]\times\R^d$, a $u$ jest gładkim, zanikającym
w nieskończoności rozwiązaniem~\eqref{eq:pde21-main}.
Wówczas dla $E_S(t)=\int_{\R^d}u^\mathsf TS u\,\mathrm d\lambda_d$
 istnieje stała $C$, zależna tylko od
$\sigma,\Sigma$ i norm odpowiednich współczynników,
taka że
\[
 \sqrt{E_S(t)}\le e^{Ct}\left(\sqrt{E_S(0)}+
 \sqrt\Sigma\int_0^t\|f(r)\|_2\,\mathrm dr\right).
\]
\end{twierdzenie}
\begin{proof}
Różniczkując, otrzymujemy
$\partial_t(u^\mathsf TS u)=2u^\mathsf TS u_t+u^\mathsf TS_tu$.
Po wstawieniu równania i użyciu symetrii $SA_j$:
\[
 \partial_t(u^\mathsf TS u)+\sum_j\partial_j(u^\mathsf TSA_ju)
 =u^\mathsf T\Bigl(S_t+\sum_j\partial_j(SA_j)-SB-B^\mathsf TS\Bigr)u
 +2u^\mathsf TSf.
\]
Całkujemy po przestrzeni. Z $E_S\ge\sigma\|u\|_2^2$
wynika ograniczenie wyrazu kwadratowego przez $C E_S$,
a z nierówności Cauchy'ego--Schwarza w iloczynie skalarnym
wyznaczonym przez $S$ dostajemy
$\left|\int_{\R^d}u(t,x)^\mathsf TS(t,x)f(t,x)\,\mathrm d\lambda_d(x)\right|\le\sqrt{E_S(t)}\sqrt\Sigma\|f(t)\|_2$.
Argument z $\sqrt{E_S+\eps^2}$ i Gronwallem kończy dowód.
\end{proof}

\begin{kontrprzyklad}[Rzeczywiste charakterystyki bez stabilności]\label{ex:pde21-jordan}
Niech $U=(u,v)^\mathsf T$ i rozważmy na
$[0,T]\times\R$ układ $u_t=v_x$, $v_t=0$.
Macierz główna w zapisie $U_t+AU_x=0$ jest równa
$A=\begin{pmatrix}0&-1\\0&0\end{pmatrix}$.
Jej jedyną wartością własną jest $0$, lecz nie ma bazy
wektorów własnych. Wybierzmy niezerową funkcję
$\phi\in C_c^\infty(\R)$ i dane
$u(0,x)=0$, $v(0,x)=n^{-1}\phi(x)\sin(nx)$.
Równanie ma jawne rozwiązanie
\[
 v(t,x)=n^{-1}\phi(x)\sin(nx),\qquad
 u(t,x)=t\bigl(\phi(x)\cos(nx)+n^{-1}\phi'(x)\sin(nx)\bigr).
\]
Norma początkowa $\|(u(0),v(0))\|_{L^2(\R;\R^2)}$
nie przekracza $n^{-1}\|\phi\|_2$ i dąży do zera.
Z kolei dla każdego ustalonego $t>0$,
na mocy lematu Riemanna--Lebesgue'a,
$\|u(t)\|_2^2\to t^2\|\phi\|_2^2/2>0$:
całka $\int_{\R}|\phi(x)|^2\cos^2(nx)\,
\mathrm d\lambda_1(x)$ dąży do $\|\phi\|_2^2/2$,
a składniki z $n^{-1}\phi'$ znikają w $L^2$.
Nie istnieje zatem ograniczony operator ewolucji
$L^2(\R;\R^2)\to L^2(\R;\R^2)$ w żadnej chwili $t>0$.
Przykład na torusie powstaje prościej, biorąc $\phi\equiv1$,
lecz funkcji okresowych nie wolno traktować jako elementów
$L^2(\R)$.
\end{kontrprzyklad}

\section{Zmiennowspółczynnikowe równanie falowe}
Niech $G(x)=G(x)^\mathsf T$ będzie macierzą $d\times d$,
$G\in W^{1,\infty}(\R^d)$ oraz
$\lambda|\xi|^2\le\xi^\mathsf TG(x)\xi\le\Lambda|\xi|^2$.
Rozważamy
\begin{equation}\label{eq:pde21-wavevariable}
 u_{tt}-\operatorname{div}(G(x)\nabla u)=f.
\end{equation}
\begin{twierdzenie}[Energia równania falowego]\label{thm:pde21-wave-energy}
Dla dostatecznie gładkiego rozwiązania~\eqref{eq:pde21-wavevariable},
które wraz z pochodnymi dostatecznie szybko zanika, energia
\[
 \mathcal E(t)=\frac12\int_{\R^d}
 (|u_t|^2+\nabla u^\mathsf TG\nabla u)\,\mathrm d\lambda_d
\]
spełnia $\mathcal E'(t)=\int_{\R^d}f u_t\,\mathrm d\lambda_d$.
Jeżeli $f=0$, energia jest stała. Jeśli $f$ znika we wstecznym
stożku o prędkości $c\ge\sqrt\Lambda$, a dane $u_0,u_1$ znikają
na jego podstawie, to $u$ znika w całym wnętrzu stożka.
\end{twierdzenie}
\begin{proof}
Mnożymy równanie przez $u_t$ i całkujemy.
Ponieważ $G$ jest niezależna od $t$, całkowanie przez części daje
$\int_{\R^d}u_{tt}u_t\,\mathrm d\lambda_d(x)
=\frac12\frac{\mathrm d}{\mathrm dt}\|u_t\|_2^2$
oraz
\[
-\int_{\R^d}\operatorname{div}(G\nabla u)u_t\,\mathrm d\lambda_d(x)
=\int_{\R^d}G\nabla u\cdot\nabla u_t\,\mathrm d\lambda_d(x)
=\frac12\frac{\mathrm d}{\mathrm dt}
\int_{\R^d}\nabla u^\mathsf TG\nabla u\,\mathrm d\lambda_d(x).
\]
To daje tożsamość energii.

Dla skończonej prędkości propagacji używamy lokalnej gęstości
$e=(|u_t|^2+\nabla u^\mathsf TG\nabla u)/2$
i strumienia $J=-u_tG\nabla u$. Z równania otrzymujemy
$\partial_te+\operatorname{div}J=f u_t$.
Na brzegu ruchomej kuli mamy
$|J\cdot\nu|\le |u_t|\,|G\nabla u\cdot\nu|
\le\sqrt\Lambda |u_t|\sqrt{\nabla u^\mathsf TG\nabla u}
\le\sqrt\Lambda\,e$.
Zatem przy ruchu brzegu z szybkością $-c$ całka energii
w malejącej kuli jest nierosnąca, jeśli $f=0$ wewnątrz.
Zerowe dane $u_1$ i $\nabla u_0$ oznaczają zerową energię
na podstawie. Stąd $u_t=\nabla u=0$ we wnętrzu stożka;
ponieważ $u_0=0$ na podstawie, otrzymujemy tam $u=0$.
\end{proof}

\section{Warunki brzegowe: dlaczego nie można narzucić dowolnych danych?}
Na obszarze ograniczonym energia układu symetrycznego spełnia
prawo~\eqref{eq:pde21-flux}, ale po całkowaniu pojawia się
strumień przez brzeg $\partial\Omega$:
$-\int_{\partial\Omega}u^\mathsf TA_\nu u\,\mathrm d\mathcal H^{d-1}$.
W szczególności dla zagadnienia Cauchy'ego na półprostej
$x>0$ z macierzą diagonalną $A=\operatorname{diag}(a_1,\ldots,a_m)$
warunki brzegowe podaje się tylko na składowych z $a_k>0$,
które wchodzą do obszaru przez $x=0$; dla $a_k<0$
wartości brzegowe są wyznaczane przez wnętrze.

\begin{przyklad}[Dwa kierunki transportu na półprostej]
Dla $a>0$ rozwiązanie $u_t+au_x=0$ na $x>0$ wymaga
jednej wartości brzegowej $u(t,0)=g(t)$: charakterystyki
wychodzące z brzegu wchodzą do obszaru. Natomiast dla $a<0$
charakterystyki biegną ku brzegowi i wartość $u(t,0)$
wynika już z danej początkowej. Dodatkowe narzucenie
niezależnego $g$ może prowadzić do braku rozwiązania.
\end{przyklad}


\begin{twierdzenie}[Dysypatywne dane brzegowe na półprostej]
\label{thm:pde21-halfline-energy}
Niech $A=A^{\mathsf T}\in\R^{m\times m}$ będzie macierzą stałą,
$T>0$, $U\in C^1([0,T]\times[0,\infty);\R^m)$ i
$F\in C([0,T]\times[0,\infty);\R^m)$ spełniają
$U_t+A U_x=F$ na $(0,T)\times(0,\infty)$.
Załóżmy, że $U$, $F$ i ich pochodne mają zanik pozwalający
na całkowanie przez części przy $x\to\infty$ oraz że
$U(t,\cdot),F(t,\cdot)\in L^2(0,\infty)$.
Niech $P_+$ będzie rzutem ortogonalnym na sumę przestrzeni
własnych $A$ odpowiadających dodatnim wartościom własnym,
a $\lambda_+=\max(0,\lambda_{\max}(A))$. Jeżeli
$P_+U(t,0)=g(t)\in\operatorname{Ran}P_+$, to
\[
 \|U(t)\|_{L^2(0,\infty)}\le
 \|U(0)\|_{L^2(0,\infty)}+
 \sqrt{\lambda_+}\,\|g\|_{L^2(0,t)}+
 \int_0^t\|F(r)\|_{L^2(0,\infty)}\,\mathrm dr.
\]
W szczególności dla $g=F=0$ energia nie rośnie.
Nie narzuca się niezależnych danych na ujemnej podprzestrzeni
spektralnej $A$; warunek dla zerowej wartości własnej
wynika z równania zwyczajnego w każdym punkcie przestrzeni.
\end{twierdzenie}
\begin{proof}
Tożsamość lokalna dla macierzy stałej ma postać
$\partial_t|U|^2+\partial_x(U^{\mathsf T}AU)=2U^{\mathsf T}F$.
Po scałkowaniu po półprostej otrzymujemy
\[
 \frac{\mathrm d}{\mathrm dt}\|U(t)\|_2^2
 =U(t,0)^{\mathsf T}AU(t,0)+
 2\langle U(t),F(t)\rangle.
\]
Rozkład spektralny $A=A_+-A_-$, gdzie $A_\pm\ge0$
i ich nośniki spektralne są ortogonalne, daje
$U^{\mathsf T}AU\le\lambda_+|P_+U|^2
=\lambda_+|g|^2$ na $x=0$. Oznaczmy
$E(t)=\|U(t)\|_2^2$, $b(t)=\|F(t)\|_2$ oraz
\[
 a(t)=E(0)+\lambda_+\int_0^t|g(r)|^2\,\mathrm dr.
\]
Z nierówności Cauchy'ego--Schwarza dostajemy
$E'(t)\le a'(t)+2\sqrt{E(t)}\,b(t)$.
Dla $\delta>0$ wprowadzamy
\[
 z_\delta(t)=\left(a(t)+\delta^2+
 2\int_0^t\sqrt{E(r)}b(r)\,\mathrm dr\right)^{1/2}.
\]
Wtedy $\sqrt{E(t)}\le z_\delta(t)$ i
$z_\delta(t)\ge\sqrt{a(t)+\delta^2}$. Dlatego
\[
 z_\delta'(t)=\frac{a'(t)+2\sqrt{E(t)}b(t)}
 {2z_\delta(t)}
 \le\frac{a'(t)}{2\sqrt{a(t)+\delta^2}}+b(t).
\]
Całkując od zera do $t$, wykorzystujemy tożsamość
$\frac{\mathrm d}{\mathrm dt}\sqrt{a(t)+\delta^2}
=a'(t)/(2\sqrt{a(t)+\delta^2})$.
Otrzymujemy
$\sqrt{E(t)}\le\sqrt{a(t)+\delta^2}+
\int_0^t b(r)\,\mathrm dr$. Granica $\delta\downarrow0$
oraz nierówność $\sqrt{r+s}\le\sqrt r+\sqrt s$
kończą dowód. Wartości własne dodatnie odpowiadają
charakterystykom wchodzącym do $x>0$ przez $x=0$,
a ujemne --- wychodzącym, na których nie można niezależnie
zadawać śladu. Dla zerowej wartości własnej projekcja
równania ma postać $(P_0U)_t=P_0F$.
\end{proof}


\section{Podsumowanie}
Symetria części głównej daje tożsamość energii, z której wynikają
jednoznaczność, stabilność i skończona prędkość propagacji.
Po zróżniczkowaniu równania komutatory nie zawierają pochodnych
wyższego rzędu niż $s$; dzięki temu otrzymujemy poprawność
w $H^s$ bez utraty pochodnych. Nie każda macierz o rzeczywistych
wartościach własnych ma te własności; pokazuje to blok Jordana.
W następnym rozdziale porównamy tę metodę z innym
mechanizmem ewolucji: unitarnością dla równania Schrödingera.
Nie jest ono ogólnym układem symetryczno-hiperbolicznym
pierwszego rzędu w przestrzeni, toteż stożkowa zasada
propagacji nie przechodzi na nie bezpośrednio.

\chapter[Równanie Schrödingera]{Równanie Schrödingera: hamiltoniany samosprzężone, ewolucja unitarna i dyspersja}
\label{chap:pde22}

\section{Od równania falowego do ewolucji unitarnej}
Niech $d\ge1$. W tym rozdziale niewiadoma jest zespolona:
$u\colon\R\times\R^d\to\C$. Przyjmujemy jednostki, w których
$\hbar=1$ oraz masa cząstki jest równa $1/2$. Równanie
Schrödingera przyjmuje wówczas postać
\begin{equation}\label{eq:pde22-schrod}
 i\partial_tu=Hu,\qquad H=-\Delta+V(x),\qquad u(0)=u_0.
\end{equation}
Potencjał $V$ ma być rzeczywisty, chyba że wyraźnie podamy inne
założenia. Symbolem $H$ oznaczamy \emph{operator wraz z jego
dziedziną}, a nie tylko formalne wyrażenie różniczkowe.
Na $L^2(\R^d;\C)$ przyjmujemy iloczyn skalarny
$\langle f,g\rangle=\int_{\R^d}f(x)\overline{g(x)}\,\mathrm d\lambda_d(x)$,
liniowy w pierwszym argumencie.

Istotna różnica względem równania ciepła jest już widoczna w
transformacie Fouriera. Dla swobodnego równania $i u_t=-\Delta u$
mnożnik wynosi $e^{-it|\xi|^2}$: jego moduł jest stale równy $1$.
Nie ma więc tłumienia wysokich częstotliwości ani analogicznego
do ciepła wygładzania w skali Sobolewa. Za to oscylacja fazy
prowadzi do \emph{dyspersji} przestrzennej.

\begin{uwaga}[Porównanie trzech ewolucji]
\label{rem:pde-v43-evolution-contrast}
Równanie ciepła (rozdziały~\ref{chap:pde18} i~\ref{chap:pde20})
prowadzi dla dodatniego operatora samosprzężonego do
mnożnika $e^{-t\lambda}$, który tłumi wysokie wartości
własne. Układ hiperboliczny z rozdziału~\ref{chap:pde21}
ma skończoną prędkość propagacji przy odpowiednich
założeniach. Swobodne równanie Schrödingera zachowuje
normę $L^2$, ale nie ma ogólnej skończonej prędkości
propagacji; samo zachowanie normy nie implikuje
zwartości ani wygładzania. Wszystkie trzy stwierdzenia
odnoszą się do \emph{różnych} operatorów i klas rozwiązań.
\end{uwaga}

\section{Samosprzężoność swobodnego hamiltonianu}
Używamy unitarnej normalizacji transformaty Fouriera z
Rozdziału~8:
\[\widehat f(\xi)=(2\pi)^{-d/2}\int_{\R^d}e^{-ix\cdot\xi}f(x)\,\mathrm d\lambda_d(x).\]
Wzór ten przyjmujemy początkowo dla $f\in\mathcal S(\R^d)$.

\begin{twierdzenie}[Swobodny hamiltonian]\label{thm:pde22-free-sa}
Niech $H_0\colon D(H_0)\subset L^2(\R^d;\C)\to L^2(\R^d;\C)$
będzie operatorem $H_0u=-\Delta u$ z dziedziną
$D(H_0)=H^2(\R^d;\C)$. Wtedy $H_0$ jest samosprzężony,
nieujemny i domknięty, a $\mathcal F H_0\mathcal F^{-1}$
jest operatorem mnożenia przez $|\xi|^2$ na dziedzinie
$\{g\in L^2:|\xi|^2g\in L^2\}$.
\end{twierdzenie}
\begin{proof}
Równość $\widehat{-\Delta u}=|\xi|^2\widehat u$ najpierw
wynika z całkowania przez części dla $u\in\mathcal S$,
a następnie z gęstości i definicji pochodnych słabych.
Ponieważ $(1+|\xi|^2)^2$ jest porównywalne z
$1+|\xi|^4$, otrzymujemy dokładnie wskazaną dziedzinę.
Wystarczy zbadać operator mnożenia $Mg=mg$, $m(\xi)=|\xi|^2$.
Jest on symetryczny, ponieważ $m$ jest rzeczywista. Jeżeli
$h\in D(M^*)$, to istnieje $k\in L^2$ takie, że
$\langle Mg,h\rangle=\langle g,k\rangle$ dla wszystkich
$g\in D(M)$. Testując funkcjami $g$ o zwartym nośniku
w $\{|\xi|\le R\}$ dostajemy $mh=k$ prawie wszędzie na tym
zbiorze. Ponieważ $R$ jest dowolne, $mh\in L^2$; zatem
$h\in D(M)$ i $M^*h=Mh$. To dowodzi samosprzężoności.
Nieujemność wynika z $\langle Mg,g\rangle=\int_{\R^d}m(\xi)|g(\xi)|^2\,\mathrm d\lambda_d(\xi)\ge0$,
a domkniętość z ogólnego faktu, że operator samosprzężony
jest domknięty; tutaj można ją też zobaczyć bezpośrednio:
ze zbieżności $g_n\to g$ i $mg_n\to h$ w $L^2$ wynika,
po wybraniu wspólnego podciągu zbieżnego prawie wszędzie,
$h=mg$. Unitarność transformaty Fouriera kończy dowód.
\end{proof}

\begin{uwaga}[Dziedzina jest częścią twierdzenia]
Formalny laplasjan na $C_c^\infty(\Omega)$ dla obszaru
z brzegiem jest zwykle jedynie symetryczny, a nie samosprzężony.
Realizacje Dirichleta i Neumanna mogą dawać różne ewolucje.
Na $\R^d$ wybór $H^2$ jest kanoniczną samosprzężoną
realizacją swobodnego hamiltonianu.
\end{uwaga}

\section{Ewolucja unitarna i poprawność zagadnienia Cauchy'ego}
\begin{twierdzenie}[Grupa swobodnego Schrödingera]\label{thm:pde22-group}
Dla $t\in\R$ i $f\in L^2(\R^d;\C)$ niech
$U_0(t)f=\mathcal F^{-1}(e^{-it|\xi|^2}\widehat f)$.
Wtedy $(U_0(t))_{t\in\R}$ jest silnie ciągłą grupą
unitarną. Dla każdego $s\in\R$ zachodzą
$\|U_0(t)f\|_{H^s}=\|f\|_{H^s}$ i
$U_0(t)H^s=H^s$. Jeśli $f\in H^2$, to
$u(t)=U_0(t)f$ należy do
$C(\R;H^2)\cap C^1(\R;L^2)$ i spełnia
$i u_t=-\Delta u$. Jest to jedyne rozwiązanie w tej klasie
z daną $u(0)=f$.
\end{twierdzenie}
\begin{proof}
Ponieważ $|e^{-it|\xi|^2}|=1$, twierdzenie Plancherela
daje unitarność i zachowanie każdej normy $H^s$.
Tożsamość mnożników
$e^{-i(t+r)|\xi|^2}=e^{-it|\xi|^2}e^{-ir|\xi|^2}$
daje prawo grupowe. Ponadto
$\|(U_0(t)-I)f\|_2^2=
\int_{\R^d}|e^{-it|\xi|^2}-1|^2|\widehat f(\xi)|^2\,\mathrm d\lambda_d(\xi)$.
Całkowana funkcja dąży punktowo do zera i jest ograniczona
przez $4|\widehat f|^2$; twierdzenie o zbieżności zdominowanej
zapewnia silną ciągłość w zerze, a prawo grupowe w każdym $t$.
Tę samą argumentację można wykonać z wagą
$(1+|\xi|^2)^s$.
Jeśli $f\in H^2$, iloraz różnicowy w dziedzinie Fouriera
ma postać
$\frac{e^{-i(t+h)|\xi|^2}-e^{-it|\xi|^2}}h\widehat f$.
Dąży do $-i|\xi|^2e^{-it|\xi|^2}\widehat f$ i ma moduł
nie większy niż $|\xi|^2|\widehat f|$. Ponownie stosujemy
zbieżność zdominowaną w $L^2$. Otrzymujemy pochodną czasową
$u_t=-iH_0u$, a więc równanie. Dla dwóch rozwiązań
różnica $w$ spełnia to samo równanie z zerowymi danymi.
Różniczkując normę, uzyskujemy
$\frac{\mathrm d}{\mathrm dt}\|w\|_2^2=
2\operatorname{Re}\langle-iH_0w,w\rangle=0$,
ponieważ $\langle H_0w,w\rangle$ jest rzeczywiste.
Zatem $w=0$.
\end{proof}

\begin{stwierdzenie}[Granica regularności]
Dla każdej liczby $s\in\R$ i $t\in\R$ grupa $U_0(t)$
jest izometrią na $H^s(\R^d)$, ale dla $t\ne0$ nie
działa w ogólności jako operator wygładzający
$H^s\to H^{s+\eps}$ dla żadnego $\eps>0$.
\end{stwierdzenie}
\begin{proof}
Gdyby $U_0(t)f\in H^{s+\eps}$ dla każdego $f\in H^s$,
to ponieważ mnożnik odwrotny $U_0(-t)$ zachowuje
$H^{s+\eps}$, mielibyśmy $f=U_0(-t)U_0(t)f\in H^{s+\eps}$
dla wszystkich $f\in H^s$, co jest fałszem.
Dla ścisłości funkcje o gęstości spektralnej
$|\widehat f(\xi)|^2\asymp
(1+|\xi|^2)^{-s-d/2}(\log(2+|\xi|))^{-2}$
dla dużych $|\xi|$ należą do $H^s$, lecz nie do
$H^{s+\eps}$. Tę samą obserwację można też wyprowadzić
wprost z zachowania norm Sobolewa.
\end{proof}

\begin{twierdzenie}[Zasada Duhamela]\label{thm:pde22-duhamel}
Niech $T>0$, $f\in L^1((0,T);L^2(\R^d))$ i
$u_0\in L^2(\R^d)$. Funkcja
\[
 u(t)=U_0(t)u_0-i\int_0^tU_0(t-r)f(r)\,\mathrm dr
\]
należy do $C([0,T];L^2)$ i jest jedynym rozwiązaniem
łagodnym $i u_t=H_0u+f$ z warunkiem $u(0)=u_0$.
Spełnia oszacowanie
$\|u(t)\|_2\le\|u_0\|_2+\int_0^t\|f(r)\|_2\,\mathrm dr$.
Jeżeli ponadto $u_0\in H^2$ i
$f\in C([0,T];H^2)$, to rozwiązanie należy do
$C([0,T];H^2)\cap C^1([0,T];L^2)$ i spełnia równanie w $L^2$.
\end{twierdzenie}
\begin{proof}
Całka jest całką Bochnera w $L^2$. Z unitarności
$\|U_0(t-r)f(r)\|_2=\|f(r)\|_2$, skąd nierówność.
Ciągłość całki względem $t$ otrzymujemy przez odjęcie
całek dla chwil $t+h,t$, wykorzystując silną ciągłość $U_0$
i zbieżność zdominowaną na części wspólnej przedziałów;
pozostała część ma normę ograniczoną przez
$\int_t^{t+h}\|f(r)\|_2\,\mathrm dr$.
Dla mocniejszych danych pochodną pod znakiem całki
obliczamy w $L^2$ przy użyciu $H_0 f\in C_tL^2$:
$\partial_t\int_0^tU_0(t-r)f(r)\,\mathrm dr
=f(t)-iH_0\int_0^tU_0(t-r)f(r)\,\mathrm dr$.
Wtedy $i u_t=H_0u+f$. Jednoznaczność łagodna wynika
z samej formuły, a w klasie silnej także z zachowania normy
dla różnicy dwóch rozwiązań.
\end{proof}

\section{Potencjały rzeczywiste i zachowanie energii}
\begin{twierdzenie}[Ograniczony potencjał rzeczywisty]\label{thm:pde22-potential}
Niech $V\in L^\infty(\R^d;\R)$ i niech $V$ działa jako
operator mnożenia na $L^2(\R^d;\C)$. Wtedy
$H=H_0+V$, $D(H)=H^2(\R^d)$, jest samosprzężony.
Grupa $U_V(t)=e^{-itH}$ jest unitarna i silnie ciągła.
Dla $u_0\in H^2$ rozwiązanie $u(t)=U_V(t)u_0$
zachowuje zarówno normę $L^2$, jak i energię
$E(u)=\langle Hu,u\rangle=
\int_{\R^d}(|\nabla u|^2+V|u|^2)\,\mathrm d\lambda_d$.
\end{twierdzenie}
\begin{proof}
Operator $V$ jest ograniczony i samosprzężony.
Suma $H_0+V$ jest symetryczna na $H^2$.
Pokażemy jej samosprzężoność bez przyjmowania
ogólnego twierdzenia o perturbacjach.
Dla $a>\|V\|_\infty$ mamy
$\|(H_0\pm ia)^{-1}\|_{L^2\to L^2}\le a^{-1}$,
co wprost wynika z mnożnika $(|\xi|^2\pm ia)^{-1}$.
Faktoryzacja
$H\pm ia=(I+V(H_0\pm ia)^{-1})(H_0\pm ia)$
pokazuje, że $H\pm ia$ są surjektywne, gdyż
$\|V(H_0\pm ia)^{-1}\|<1$ i odwrotność pierwszego
czynnika otrzymujemy z szeregu Neumanna.
Symetryczny operator o obu surjektywnych operatorach
$H\pm ia$ jest samosprzężony: istotnie dla
$z\in D(H^*)$ i $h=(H^*+ia)z$ dobieramy
$y\in D(H)$ takie, że $(H+ia)y=h$.
Wtedy $(H^*+ia)(z-y)=0$, a
$\ker(H^*+ia)=\operatorname{Ran}(H-ia)^\perp=\{0\}$;
zatem $z=y\in D(H)$. To dowodzi $H=H^*$.
Z twierdzenia spektralnego dla operatorów samosprzężonych
wynika istnienie silnie ciągłej grupy $e^{-itH}$ oraz
$U_V(t)D(H)=D(H)$, $HU_V(t)=U_V(t)H$.
Unitarność daje zachowanie normy, a ostatnia komutacja
prowadzi do
$\langle HU_V(t)u_0,U_V(t)u_0\rangle=
\langle Hu_0,u_0\rangle$. Całkowanie przez części
jest uzasadnione dla $u_0\in H^2$ i $V$ ograniczone.
\end{proof}

\begin{stwierdzenie}[Lokalny bilans prawdopodobieństwa]
Niech $V\in L^\infty(\R^d;\R)$ oraz
$u_0\in H^2(\R^d)$. Dla rozwiązania z
Twierdzenia~\ref{thm:pde22-potential} zachodzi w sensie
dystrybucji na $\R\times\R^d$
\[
 \partial_t|u|^2+\operatorname{div}j=0,
 \qquad j=2\operatorname{Im}(\overline u\,\nabla u).
\]
\end{stwierdzenie}
\begin{proof}
Ponieważ $u_t=-iHu=i\Delta u-iVu$,
$\partial_t|u|^2=2\operatorname{Re}(\overline u\,u_t)
=-2\operatorname{Im}(\overline u\,\Delta u)$.
Wyraz $-iV|u|^2$ nie wnosi wkładu do części rzeczywistej,
bo $V$ jest rzeczywisty. Z kolei
$\operatorname{Im}(\nabla\overline u\cdot\nabla u)=0$,
więc $\operatorname{Im}(\overline u\,\Delta u)
=\operatorname{div}\operatorname{Im}(\overline u\,\nabla u)$.
Tożsamość jest poprawna dla gładkich funkcji, a następnie
w sensie dystrybucji przez aproksymację w $H^2$.
\end{proof}


\begin{twierdzenie}[Zależny od czasu, ograniczony potencjał rzeczywisty]
\label{thm:pde22-time-potential}
Niech $T>0$, $V\in L^\infty((0,T)\times\R^d;\R)$,
$F\in L^1(0,T;L^2(\R^d;\C))$ oraz
$u_0\in L^2(\R^d;\C)$. Istnieje dokładnie jedna funkcja
$u\in C([0,T];L^2(\R^d;\C))$ spełniająca formułę Duhamela
\[
 u(t)=U_0(t)u_0-i\int_0^tU_0(t-s)
       \bigl(V(s,\cdot)u(s)+F(s)\bigr)\,\mathrm ds.
\]
Zachodzi oszacowanie
$\|u(t)\|_2\le\|u_0\|_2+\int_0^t\|F(s)\|_2\,\mathrm ds$.
Jeżeli $F=0$, to $\|u(t)\|_2=\|u_0\|_2$ dla każdego
$t\in[0,T]$, mimo że operator $H(t)=-\Delta+V(t,\cdot)$
może zależeć od czasu. Nie twierdzimy, że funkcja ta jest
rozwiązaniem silnym o wartościach w $H^2$.
\end{twierdzenie}
\begin{proof}
Stosujemy obraz oddziaływania $w(t)=U_0(-t)u(t)$.
Szukana równość jest równoważna równaniu całkowemu
\[
 w(t)=u_0-i\int_0^t\bigl(K(s)w(s)+G(s)\bigr)\,\mathrm ds,
 \quad K(s)=U_0(-s)M_{V(s)}U_0(s),\quad
 G(s)=U_0(-s)F(s).
\]
Operator $M_{V(s)}$ jest mnożeniem przez rzeczywistą
funkcję $V(s,\cdot)$. Dla prawie każdego $s$ operator $K(s)$
jest ograniczony, samosprzężony i
$\|K(s)\|\le\|V\|_\infty$. Funkcja
$s\mapsto K(s)z$ jest silnie mierzalna dla każdego $z\in L^2$;
wynika to z mierzalności $V$, separowalności $L^2$ i silnej
ciągłości $U_0(s)$.

Iteracje Picarda równania całkowego zbiegają jednostajnie
w $C([0,T];L^2)$: różnica kolejnych iteracji jest
ograniczona przez stałą pomnożoną przez
$(\|V\|_\infty T)^n/n!$. Dla dwóch rozwiązań nierówność
Gronwalla daje jednoznaczność. Ponadto
$w\in W^{1,1}(0,T;L^2)$ i prawie wszędzie
$w'=-iKw-iG$. Dla naszego iloczynu skalarnego, liniowego
w pierwszym argumencie, samosprzężoność $K$ daje
\[
 \frac{\mathrm d}{\mathrm dt}\|w(t)\|_2^2
 =2\operatorname{Re}\langle-iK(t)w(t)-iG(t),w(t)\rangle
 =2\operatorname{Re}\langle-iG(t),w(t)\rangle.
\]
Przy $G=0$ norma jest stała. Ogólnie, regularizując
pierwiastek $\sqrt{\|w\|_2^2+\delta^2}$ i przechodząc
$\delta\downarrow0$, otrzymujemy
$\|w(t)\|_2\le\|u_0\|_2+\int_0^t\|G(s)\|_2\,\mathrm ds$.
Unitarność $U_0(t)$ oraz $\|G(s)\|_2=\|F(s)\|_2$
kończy dowód.
\end{proof}


\begin{uwaga}[Potencjały nieograniczone]
Potencjały Coulomba i oscylatora harmonicznego wymagają
innych twierdzeń o samosprzężoności lub konstrukcji
przez formy kwadratowe. Nie wynikają one z powyższego
twierdzenia, ponieważ nie należą do $L^\infty$.
\end{uwaga}

\section{Jądro swobodnego propagatora i oszacowanie dyspersyjne}
\begin{twierdzenie}[Wzór na propagator i zanik dyspersyjny]
\label{thm:pde22-disp}
Dla $t\in\R\setminus\{0\}$ i $f\in\mathcal S(\R^d)$ zachodzi
\begin{equation}\label{eq:pde22-kernel}
 U_0(t)f(x)=(4\pi it)^{-d/2}
 \int_{\R^d}\exp\!\left(\frac{i|x-y|^2}{4t}\right)
 f(y)\,\mathrm d\lambda_d(y),
\end{equation}
gdzie potęgę zespoloną wybiera się zgodnie z ciągłą
granicą regularizacji gaussowskiej. W szczególności
\begin{equation}\label{eq:pde22-L1Linf}
 \|U_0(t)f\|_\infty\le(4\pi|t|)^{-d/2}\|f\|_1.
\end{equation}
Dla $2\le r\le\infty$ i $r'$ sprzężonego do $r$
zachodzi ponadto
\[
 \|U_0(t)f\|_{L^r}\le C_{d,r}
 |t|^{-d(1/2-1/r)}\|f\|_{L^{r'}}
\]
dla $f\in L^{r'}\cap L^2$, z przedłużeniem do wskazanych
klas przez gęstość.
\end{twierdzenie}
\begin{proof}
Wstawiamy czynnik $e^{-\eps|\xi|^2}$ do mnożnika
Fouriera i otrzymujemy, dla $\eps>0$, splot z funkcją
\[
 K_{t,\eps}(x)=[4\pi(\eps+it)]^{-d/2}
 \exp\!\left(-\frac{|x|^2}{4(\eps+it)}\right).
\]
Formuła wynika przez obliczenie transformaty Gaussa,
początkowo dla rzeczywistego parametru dodatniego,
a następnie przez analityczną kontynuację do półpłaszczyzny
$\operatorname{Re}(\eps+it)>0$. W dziedzinie Fouriera
mamy $e^{-(\eps+it)|\xi|^2}\widehat f\to
 e^{-it|\xi|^2}\widehat f$ w $L^2$ na mocy dominacji.
Dla $f\in\mathcal S$ i ustalonego $t\ne0$ dominacja
w całce przestrzennej daje także zbieżność punktową
$K_{t,\eps}*f\to K_t*f$, gdzie
$K_t(x)=(4\pi it)^{-d/2}\exp(i|x|^2/(4t))$.
Zbieżności muszą opisywać tę samą dystrybucję,
co dowodzi wzoru~\eqref{eq:pde22-kernel}.
Ponieważ $|K_t(x)|=(4\pi|t|)^{-d/2}$,
nierówność~\eqref{eq:pde22-L1Linf} jest natychmiastowa.
Drugie oszacowanie wynika przez interpolację
Riesza--Thorina między $L^1\to L^\infty$
a unitarnością $L^2\to L^2$. Przy parametrze
$\theta\in[0,1]$ mamy $1/r=\theta/2$
i $1/r'=1-\theta/2$; potęga czasu to
$-d(1-\theta)/2=-d(1/2-1/r)$.
\end{proof}

\begin{przyklad}[Pakiet gaussowski i brak stożka propagacji]
Dla $u_0(x)=e^{-|x|^2/2}$ otrzymujemy
\[
 U_0(t)u_0(x)=(1+2it)^{-d/2}
 \exp\!\left(-\frac{|x|^2}{2(1+2it)}\right).
\]
Dla każdego $t\ne0$ funkcja ta jest niezerowa we wszystkich
punktach $\R^d$. Zatem równanie Schrödingera, w przeciwieństwie
do równania falowego z rozdziału 21, nie ma skończonej prędkości
propagacji w sensie zachowywania zwartego nośnika.
Ściśle rzecz biorąc, pakiet Gaussa ma już w chwili początkowej
niezwarty nośnik. Poprawny kontrprzykład ze zwartymi danymi
otrzymujemy jednak ze wzoru~\eqref{eq:pde22-kernel}:
jeżeli $f\in C_c^\infty(\R^d)$ jest niezerowa, to dla $t\ne0$
$U_0(t)f$ jest, po wyłączeniu czynnika
$e^{i|x|^2/(4t)}$, transformatą Fouriera funkcji
o zwartym nośniku. Jest więc rzeczywiście analityczna
jako funkcja zmiennej rzeczywistej i nie może znikać na
niepustym zbiorze otwartym bez identycznego zerowania.
Unitarność wyklucza identyczne zerowanie. W szczególności
$U_0(t)f$ nie może mieć zwartego nośnika.
\end{przyklad}

\section[Oszacowania Strichartza w dwóch wymiarach]{Oszacowania Strichartza w modelowym przypadku $d=2$}
Oszacowanie dyspersyjne dotyczy ustalonego czasu. Pokażemy,
jak z niego uzyskać normę czasowo-przestrzenną.
Wybrany przypadek $(q,r)=(4,4)$ jest wystarczająco bogaty,
aby zobaczyć całą ideę metody $TT^*$, a zarazem pozwala
uniknąć najtrudniejszych zagadnień końcowych.

\begin{lemat}[Oszacowanie całkowe w czasie]\label{lem:pde22-hls}
Dla każdego $g\in L^{4/3}(\R)$ zachodzi
\begin{equation}\label{eq:pde22-hls}
 \left\|\int_{\R}|t-s|^{-1/2}|g(s)|\,\mathrm ds
 \right\|_{L^4_t}\le C\|g\|_{L^{4/3}_t}.
\end{equation}
\end{lemat}
\begin{proof}
Niech $Mg(t)=\sup_{R>0}(2R)^{-1}\int_{t-R}^{t+R}|g(s)|\,\mathrm ds$
będzie operatorem maksymalnym Hardy'ego--Littlewooda.
Najpierw przypomnimy wraz z uzasadnieniem jego ograniczoność
na $L^p(\R)$ dla $p>1$. Dla skończenie wielu przedziałów
wybranych z rodziny $\{I:|I|^{-1}\int_I|g|>\lambda\}$
można wybrać rozłączne $I_j$ tak, aby ich trzykrotne
powiększenia pokrywały całą rozważaną skończoną rodzinę.
Wynika stąd oszacowanie słabego typu
$|\{Mg>\lambda\}|\le C\|g\|_1/\lambda$
przez aproksymację otwartego zbioru poziomicowego.
Rozdzielmy teraz $g=g_1+g_2$,
$g_1=g\,\chi_{\{|g|>\lambda/2\}}$ oraz
$g_2=g\,\chi_{\{|g|\le\lambda/2\}}$.
Wtedy $Mg_2\le\lambda/2$, więc
$|\{Mg>\lambda\}|\le C\lambda^{-1}
\int_{\{|g|>\lambda/2\}}|g(s)|\,\mathrm ds$.
Po pomnożeniu przez $p\lambda^{p-1}$, całkowaniu
w $\lambda$ i zastosowaniu twierdzenia Tonellego
dostajemy $\|Mg\|_p^p\le C_p\|g\|_p^p$.

Dla ustalonego $R>0$ rozbijamy całkę w~\eqref{eq:pde22-hls}
na $|t-s|\le R$ i $|t-s|>R$. W pierwszej części
rozkład na pierścienie dyadyczne daje oszacowanie
$C R^{1/2}Mg(t)$. W drugiej nierówność Höldera z
wykładnikami $4$ i $4/3$ daje
$C R^{-1/4}\|g\|_{4/3}$, gdyż
$\int_{|r|>R}|r|^{-2}\,\mathrm dr=2/R$.
Wybieramy $R>0$ tak, aby
$R^{3/4}=\|g\|_{4/3}/Mg(t)$, jeśli oba człony
są dodatnie (przypadki zerowe otrzymujemy granicznie).
W rezultacie lewa strona całki punktowej jest ograniczona
przez $C(Mg(t))^{1/3}\|g\|_{4/3}^{2/3}$.
Podnosząc do czwartej potęgi i całkując, otrzymujemy
\[
 \|I_{1/2}|g|\|_4^4\le C\|Mg\|_{4/3}^{4/3}
 \|g\|_{4/3}^{8/3}\le C\|g\|_{4/3}^4.
\]
To dowodzi tezy.
\end{proof}

\begin{twierdzenie}[Oszacowanie Strichartza $L^4_{t,x}$ w $\R^2$]
\label{thm:pde22-strichartz}
Niech $d=2$, $u_0\in L^2(\R^2)$ i
$F\in L^{4/3}(\R\times\R^2)$. Wtedy
\begin{align}
 \|U_0(t)u_0\|_{L^4(\R\times\R^2)}&\le C\|u_0\|_2,
 \label{eq:pde22-str-hom}\\
 \left\|\int_{s<t}U_0(t-s)F(s)\,\mathrm ds
 \right\|_{L^4(\R\times\R^2)}&\le
 C\|F\|_{L^{4/3}(\R\times\R^2)}.
 \label{eq:pde22-str-inhom}
\end{align}
Drugą całkę definiujemy początkowo dla gładkich danych
o zwartym nośniku, a następnie przez domknięcie.
\end{twierdzenie}
\begin{proof}
Oznaczmy $T\varphi(t)=U_0(t)\varphi$ dla
$\varphi\in\mathcal S(\R^2)$. Formalnym operatorem
sprzężonym jest
$T^*G=\int_{\R}U_0(-s)G(s)\,\mathrm ds$
(dla $G\in C_c^\infty(\R\times\R^2)$).
Stąd
$TT^*G(t)=\int_{\R}U_0(t-s)G(s)\,\mathrm ds$.
Interpolacja oszacowania dyspersyjnego dla $d=2$
między $L^1\to L^\infty$ i $L^2\to L^2$ daje
$\|U_0(t-s)G(s)\|_{L^4_x}\le
 C|t-s|^{-1/2}\|G(s)\|_{L^{4/3}_x}$.
Z nierówności Minkowskiego oraz
Lematu~\ref{lem:pde22-hls} wynika zatem
\begin{equation}\label{eq:pde22-ttstar}
 \|TT^*G\|_{L^4_{t,x}}\le C\|G\|_{L^{4/3}_{t,x}}.
\end{equation}
Wykorzystujemy teraz dodatniość formy
$\langle TT^*G,G\rangle_{t,x}=\|T^*G\|_{L^2_x}^2$.
Nierówność Höldera wraz z~\eqref{eq:pde22-ttstar}
daje $\|T^*G\|_2^2\le C\|G\|_{4/3}^2$.
Dualność pomiędzy $L^{4/3}_{t,x}$ i $L^4_{t,x}$
implikuje zatem
$\|T\varphi\|_4\le C^{1/2}\|\varphi\|_2$.
Przez gęstość otrzymujemy~\eqref{eq:pde22-str-hom}.

Dla operatora opóźnionego zastępujemy całkę
$\int_{\R}$ przez $\int_{s<t}$. To samo oszacowanie
punktowe pozostaje prawdziwe, ponieważ
$\int_{s<t}|t-s|^{-1/2}g(s)\,\mathrm ds\le
\int_{\R}|t-s|^{-1/2}g(s)\,\mathrm ds$
dla nieujemnego $g$; Lemat~\ref{lem:pde22-hls}
prowadzi wprost do~\eqref{eq:pde22-str-inhom}.
\end{proof}

\begin{uwaga}[Admisybilność i zakres twierdzenia]
W wymiarze $d$ para wykładników $2\le q,r\le\infty$
jest formalnie zgodna ze skalowaniem, gdy
$2/q+d/r=d/2$. Powyżej \emph{nie} dowodziliśmy
wszystkich takich oszacowań ani przypadków końcowych.
W szczególności przypadek $(q,r)=(2,\infty)$ w wymiarze
$d=2$ nie spełnia standardowego twierdzenia Strichartza.
Oszacowania dla operatorów z potencjałem wymagają dodatkowych
założeń spektralnych; nie wynikają z samej samosprzężoności.
\end{uwaga}

\section{Stan związany i brak dyspersji dla pewnych hamiltonianów}
\begin{przyklad}[Oscylator harmoniczny]
W przestrzeni $L^2(\R^d)$ rozważmy nieujemny operator
$H_{\mathrm{osc}}=-\Delta+|x|^2$ w jego standardowej
samosprzężonej realizacji. Funkcja
$\psi_0(x)=e^{-|x|^2/2}$ spełnia
$H_{\mathrm{osc}}\psi_0=d\psi_0$.
W szczególności $e^{-itH_{\mathrm{osc}}}\psi_0
=e^{-itd}\psi_0$: jej gęstość przestrzenna
$|\psi_0(x)|^2$ nie zmienia się w czasie.
Nie może więc istnieć uniwersalne oszacowanie
$\|e^{-itH_{\mathrm{osc}}}f\|_\infty\le C|t|^{-d/2}\|f\|_1$
dla wszystkich $t\ne0$ i wszystkich $f\in L^1\cap L^2$,
bo zastosowanie go do $\psi_0$ i przejście $t\to\infty$
prowadzi do sprzeczności.
\end{przyklad}

\section{Zależność od warunków brzegowych i dalsze pytania}
Na ograniczonym obszarze gładkim laplasjan Dirichleta
ma dyskretne widmo. Jego ewolucja Schrödingera jest nadal
unitarna, lecz na funkcji własnej działa tylko przez zmianę
fazy, a zatem nie występuje globalny zanik dyspersyjny
analogiczny do~\eqref{eq:pde22-L1Linf}. Jest to kolejny
powód, dla którego należy wyraźnie oddzielać twierdzenia
na $\R^d$ od zagadnień z brzegiem.

Pozostają ważne kierunki dalszego rozwoju:
bezpośrednia konstrukcja hamiltonianów z osobliwymi potencjałami,
lokalne i globalne oszacowania Strichartza dla nieliniowych
równań Schrödingera, zjawiska rozpraszania oraz analiza
rezonansów spektralnych. Zagadnienia te wykraczają poza
podstawowy model liniowy niniejszego rozdziału.


\chapter[Równania Hamiltona--Jacobiego]{Równania Hamiltona--Jacobiego: charakterystyki, rozwiązania lepkościowe i wzór Hopfa--Laxa}
\label{chap:pde23}

\section{Zagadnienie Cauchy'ego i geometryczne znaczenie hamiltonianu}
Niech $T>0$, $d\ge1$, $g\colon\R^d\to\R$, a
$H\colon\R^d\times\R^d\to\R$ będzie hamiltonianem.
Rozważamy równanie
\begin{equation}\label{eq:pde23-hj}
 \partial_tu(t,x)+H(x,\nabla_xu(t,x))=0,
 \qquad u(0,x)=g(x).
\end{equation}
W pierwszej części zakładamy gładkość danych, potem będziemy
rozpatrywać funkcje ciągłe, niekoniecznie różniczkowalne.
Odróżniamy gradient względem $x$ od gradientu $D_pH$
względem zmiennej pędowej $p$.

Równanie~\eqref{eq:pde23-hj} występuje w mechanice
Hamiltonowskiej, optymalnym sterowaniu i prawach zachowania.
Nawet jeśli $g$ i $H$ są gładkie, rozwiązanie klasyczne
może utracić różniczkowalność po skończonym czasie.
Nie jest to problem złej definicji pochodnej słabej:
po przecięciu charakterystyk potrzebna jest dodatkowa
zasada wyboru rozwiązania. Lokalną konstrukcję klasycznych
charakterystyk poznaliśmy w rozdziale~\ref{chap:pde-characteristics};
poniżej wyjaśnimy, dlaczego samo jej przedłużanie nie
rozwiązuje kwestii jednoznaczności.

\begin{definicja}[Rozwiązanie klasyczne]
Niech $H$ będzie ciągły, a $g$ ciągła. Funkcja
$u\in C([0,T]\times\R^d)\cap C^1((0,T)\times\R^d)$
jest rozwiązaniem klasycznym~\eqref{eq:pde23-hj}, jeżeli
spełnia równanie punktowo dla $0<t<T$ oraz $u(0,x)=g(x)$.
\end{definicja}

\section{Charakterystyki układu Hamiltona i utrata gładkości}
\begin{twierdzenie}[Układ charakterystyczny]
\label{thm:pde23-char}
Niech $H\in C^2(\R^d_x\times\R^d_p)$ i niech
$u\in C^2((0,T)\times\R^d)$ spełnia
$u_t+H(x,\nabla u)=0$. Dla dowolnego punktu $(t_0,x_0)$
rozważmy lokalne rozwiązanie układu
\begin{equation}\label{eq:pde23-hamilton}
 \dot X=D_pH(X,P),\qquad \dot P=-D_xH(X,P),
 \qquad X(t_0)=x_0,\quad P(t_0)=\nabla u(t_0,x_0).
\end{equation}
Wtedy $P(t)=\nabla u(t,X(t))$ i
\begin{equation}\label{eq:pde23-action}
 \frac{\mathrm d}{\mathrm dt}u(t,X(t))
 =P(t)\cdot D_pH(X(t),P(t))-H(X(t),P(t)).
\end{equation}
\end{twierdzenie}
\begin{proof}
Różniczkujemy równanie Hamiltona--Jacobiego względem $x_j$:
\[
 \partial_tu_{x_j}+\sum_k H_{p_k}(x,\nabla u)u_{x_jx_k}
 +H_{x_j}(x,\nabla u)=0.
\]
Jeżeli $X$ rozwiązuje
$\dot X=D_pH(X,\nabla u(t,X))$, to reguła łańcuchowa daje
$\frac{\mathrm d}{\mathrm dt}\nabla u(t,X(t))=-D_xH(X(t),\nabla u(t,X(t)))$.
Wobec jednoznaczności układu ODE otrzymujemy dokładnie
$P(t)=\nabla u(t,X(t))$. Dalsze różniczkowanie prowadzi do
$\frac{\mathrm d}{\mathrm dt}u=u_t+\nabla u\cdot\dot X
=-H+P\cdot D_pH$.
\end{proof}

Dla $H=H(p)$ druga część układu daje $P(t)=P(0)$.
Jeśli $g\in C^2(\R^d)$, charakterystyka wychodząca
z $y$ ma zatem postać
\begin{equation}\label{eq:pde23-explicit-char}
 X(t,y)=y+tD_pH(\nabla g(y)),\qquad P(t,y)=\nabla g(y).
\end{equation}
Przy odwracalności $y\mapsto X(t,y)$ wzór
\eqref{eq:pde23-action} wyznacza lokalne rozwiązanie
klasyczne. Jego wartość na charakterystyce wynosi
$g(y)+t[\nabla g(y)\cdot D_pH(\nabla g(y))-H(\nabla g(y))]$.
Warunek lokalnej odwracalności jest kontrolowany przez
\begin{equation}\label{eq:pde23-jacobian}
 D_yX(t,y)=I+tD_{pp}^2H(\nabla g(y))D^2g(y).
\end{equation}
Jeśli wyznacznik tej macierzy znika, metoda klasyczna
nie gwarantuje już jednowartościowego gradientu.

\begin{przyklad}[Powstanie osobliwości z gładkich danych]
\label{ex:pde23-cos}
Niech $d=1$, $H(p)=p^2/2$ i $g(y)=\cos y$. Wtedy
$X(t,y)=y-t\sin y$, a $\partial_yX=1-t\cos y$.
Dla $0\le t<1$ pochodna jest dodatnia dla wszystkich $y$.
W chwili $t=1$ zeruje się przy $y\in2\pi\mathbb Z$;
po czasie $t=1$ charakterystyki przecinają się i
klasyczna metoda nie wyznacza już globalnej funkcji $C^1$.
Poniżej wyprowadzimy jednoznaczne rozwiązanie lepkościowe
istniejące również dla $t>1$.
\end{przyklad}

\section{Rozwiązania lepkościowe: definicja bez pochodnych rozwiązania}
\begin{definicja}[Rozwiązanie lepkościowe]
\label{def:pde23-visc}
Niech $Q=(0,T)\times\Omega$, gdzie $\Omega\subset\R^d$
jest obszarem, a $H\colon\Omega\times\R^d\to\R$ jest
ciągły. Funkcja $u\in C(Q)$ jest \emph{podrozwiązaniem
lepkościowym} równania $u_t+H(x,\nabla u)=0$, jeśli dla
każdego $(t_0,x_0)\in Q$ i każdej funkcji
$\varphi\in C^1(Q)$, dla której $u-\varphi$ ma
lokalne maksimum w $(t_0,x_0)$, zachodzi
\begin{equation}\label{eq:pde23-sub}
 \varphi_t(t_0,x_0)+H(x_0,\nabla\varphi(t_0,x_0))\le0.
\end{equation}
\emph{Nadrozwiązanie lepkościowe} to funkcja, dla której
w każdym punkcie lokalnego minimum $u-\varphi$ zachodzi
odwrotna nierówność $\ge0$. \emph{Rozwiązanie lepkościowe}
jest jednocześnie podrozwiązaniem i nadrozwiązaniem.
Warunek początkowy $u(0,x)=g(x)$ rozumiemy w sensie ciągłości.
\end{definicja}

\begin{uwaga}
Słowo „lepkościowe” pochodzi z przybliżenia równań pierwszego
rzędu równaniami $u_t+H(x,\nabla u)=\eps\Delta u$,
$\eps\downarrow0$. Definicja nie wymaga istnienia takiego
przybliżenia ani pochodnych $u$. Korzysta wyłącznie
z pochodnych gładkiej funkcji testowej.
Nie utożsamiamy jej automatycznie z rozwiązaniem słabym
w sensie dystrybucyjnym.
\end{uwaga}

\begin{stwierdzenie}[Zgodność z rozwiązaniem klasycznym]
\label{prop:pde23-smooth-visc}
Jeżeli $u\in C^1(Q)$, to $u$ jest rozwiązaniem
lepkościowym równania $u_t+H(x,\nabla u)=0$
wtedy i tylko wtedy, gdy spełnia je punktowo.
\end{stwierdzenie}
\begin{proof}
Jeżeli $u$ spełnia równanie punktowo i $u-\varphi$ ma
lokalne ekstremum, pochodne $u-\varphi$ znikają:
$\varphi_t=u_t$ i $\nabla\varphi=\nabla u$. Dostajemy
właściwą nierówność, a nawet równość. W drugą stronę
w ustalonym punkcie $(t_0,x_0)$ bierzemy $\varphi=u$.
Różnica $u-\varphi\equiv0$ ma zarazem maksimum i minimum.
Obie nierówności wymuszają równość.
\end{proof}

\section{Zasada porównania na torusie}
Przez $\T^d=\R^d/(2\pi\mathbb Z)^d$ rozumiemy płaski torus,
a przez $\lambda_d$ na $\T^d$ miarę indukowaną przez miarę
Lebesgue'a na $[0,2\pi)^d$, bez normalizacji do jedynki;
w szczególności $\lambda_d(\T^d)=(2\pi)^d$.
Najważniejszym argumentem zapewniającym jednoznaczność
jest \emph{podwajanie zmiennych}. Podajemy pełny dowód
modelowego twierdzenia, w którym brak brzegu pozwala
oddzielić istotę argumentu od technicznych barier.

\begin{twierdzenie}[Porównanie rozwiązań lepkościowych]
\label{thm:pde23-comparison}
Niech $T>0$, a $H\colon\T^d\times\R^d\to\R$ będzie
ciągły i niech istnieje $L>0$ takie, że dla wszystkich
$x,y\in\T^d$, $p\in\R^d$ i dla dostatecznie bliskich
$x,y$, przy użyciu tej samej lokalnej reprezentacji
współrzędnych,
\begin{equation}\label{eq:pde23-HxLip}
 |H(x,p)-H(y,p)|\le L|x-y|(1+|p|).
\end{equation}
Niech $u,v\in C([0,T]\times\T^d)$ będą odpowiednio
podrozwiązaniem i nadrozwiązaniem lepkościowym
$w_t+H(x,\nabla w)=0$ w $(0,T)\times\T^d$.
Jeśli $u(0,x)\le v(0,x)$ dla każdego $x$, to
$u(t,x)\le v(t,x)$ dla wszystkich $(t,x)\in[0,T]\times\T^d$.
\end{twierdzenie}
\begin{proof}
Przypuśćmy, że nierówność jest fałszywa. Wybieramy
$0<\tau<T$ oraz $(t_*,x_*)\in(0,\tau)\times\T^d$
takie, że $u(t_*,x_*)-v(t_*,x_*)>0$.
Wybieramy małe $\delta,\kappa>0$ spełniające
$u(t_*,x_*)-v(t_*,x_*)-\delta t_*-\kappa/(\tau-t_*)>0$.
Dla $\eps,\alpha>0$ maksymalizujemy dla
$t\in[0,\tau)$, $s\in[0,\tau]$ i $x,y\in\T^d$ funkcję
\begin{equation}\label{eq:pde23-doubling}
 \Phi_{\eps,\alpha}(t,x,s,y)=u(t,x)-v(s,y)
 -\frac{d_{\T}(x,y)^2}{2\eps}
 -\frac{(t-s)^2}{2\alpha}
 -\delta t-\frac{\kappa}{\tau-t}.
\end{equation}
Z powodu ostatniego składnika maksimum jest osiągane
przy $t<\tau$; przy ustalonym $t$ wybór $s\in[0,\tau]$
nie wymaga osobnej bariery. Maksimum jest dodatnie,
bo można wstawić $t=s=t_*$ i $x=y=x_*$.
Niech $C_0=2\max(\|u\|_\infty,\|v\|_\infty)+1$.
W punkcie dodatniego maksimum wszystkie składniki kary
są nieujemne, toteż $\kappa/(\tau-t_\eps)\le C_0$.
A zatem $t_\eps\le\tau-\kappa/C_0$ niezależnie od
$\eps,\alpha$. Po wykazaniu, że
$|t_\eps-s_\alpha|\to0$, dostajemy również $s_\alpha<\tau$
dla dostatecznie małych parametrów. Odrębnie poniżej
wykluczamy ich zerowanie się przy $t=0$.

Niech $(t_\eps,x_\eps,s_\alpha,y_\eps)$ będzie
punktem maksimum; dla czytelności pomijamy w indeksach
część parametrów. Porównanie z punktem
$(t_\eps,x_\eps,s_\alpha,x_\eps)$ oraz jednostajna
ciągłość $v$ dają
$\frac{d_{\T}(x_\eps,y_\eps)^2}{2\eps}
\le \omega_v(d_{\T}(x_\eps,y_\eps))$.
Z ograniczoności $u,v$ wynika
$d_{\T}(x_\eps,y_\eps)\to0$, a z powyższej
nierówności
\begin{equation}\label{eq:pde23-pen-small}
 \frac{d_{\T}(x_\eps,y_\eps)^2}{\eps}\longrightarrow0.
\end{equation}
Analogicznie porównanie z $s=t$ i jednostajna ciągłość
$v$ w czasie dają $|t_\eps-s_\alpha|\to0$ i
$(t_\eps-s_\alpha)^2/\alpha\to0$. Wybieramy
kolejność granic $\eps\downarrow0$, $\alpha\downarrow0$
lub wspólną kolejność zbieżną do zera.

Dla dostatecznie małych parametrów zarówno $t_\eps$,
jak i $s_\alpha$ należą do $(0,\tau)$.
Gdyby na przykład $t_\eps=0$ dla pewnego ciągu parametrów,
to $s_\alpha\to0$, $d_{\T}(x_\eps,y_\eps)\to0$;
jednostajna ciągłość i warunek początkowy dawałyby
graniczną wartość maksimum $\le0$, w sprzeczności
z ustaloną dodatnią wartością w punkcie $(t_*,x_*)$.
Tak samo wykluczamy $s_\alpha=0$.

Ponieważ punkty przestrzenne są bliskie, kwadrat
odległości ma w ich otoczeniu gładką postać euklidesową.
Oznaczmy $p=(x_\eps-y_\eps)/\eps$ w lokalnych
współrzędnych. Z definicji maksimum skonstruowane
przez zamrożenie $(s,y)$ funkcje testowe dotykają $u$
z góry, a przez zamrożenie $(t,x)$ dotykają $v$ z dołu.
Otrzymujemy zatem
\begin{align*}
 \frac{t_\eps-s_\alpha}{\alpha}+\delta
  +\frac{\kappa}{(\tau-t_\eps)^2}+H(x_\eps,p)&\le0,\\
 \frac{t_\eps-s_\alpha}{\alpha}+H(y_\eps,p)&\ge0.
\end{align*}
Po odjęciu dostajemy
$\delta\le H(y_\eps,p)-H(x_\eps,p)$.
Założenie~\eqref{eq:pde23-HxLip} pozwala ograniczyć
prawą stronę przez
$L|x_\eps-y_\eps|+L|x_\eps-y_\eps|^2/\eps$,
co dąży do zera na mocy~\eqref{eq:pde23-pen-small}.
Sprzeczność dowodzi twierdzenia dla $t<\tau$,
a dowolność $\tau$ i ciągłość dają tezę do czasu $T$.
\end{proof}

\begin{wniosek}[Jednoznaczność i kontrakcja]
\label{cor:pde23-contract}
Przy założeniach Twierdzenia~\ref{thm:pde23-comparison}
rozwiązanie lepkościowe z danymi $g\in C(\T^d)$ jest
co najwyżej jedno. Jeżeli $u,v$ są dwoma rozwiązaniami
dla danych $g,h\in C(\T^d)$, to
$\|u(t)-v(t)\|_{L^\infty(\T^d)}\le
\|g-h\|_{L^\infty(\T^d)}$ dla $0\le t\le T$.
\end{wniosek}
\begin{proof}
Równanie nie zależy od wartości $u$, zatem $v+c$
jest rozwiązaniem, ilekroć $v$ nim jest. Stosujemy
porównanie do pary $(u,v+c)$, gdzie
$c=\|g-h\|_\infty$, a następnie zamieniamy rolami $u,v$.
\end{proof}

\section{Transformata Legendre'a i wzór Hopfa--Laxa}
\begin{definicja}[Funkcja sprzężona Legendre'a]
Niech $H\colon\R^d\to\R$ będzie funkcją wypukłą.
Jej funkcję sprzężoną definiujemy wzorem
$L(v)=\sup_{p\in\R^d}(p\cdot v-H(p))$.
Funkcję $H$ nazywamy \emph{superliniową}, gdy
$H(p)/|p|\to\infty$ przy $|p|\to\infty$.
\end{definicja}

\begin{lemat}[Dualność i kontrola prędkości minimalizatorów]
\label{lem:pde23-legendre}
Niech $H\in C^1(\R^d)$ będzie wypukła i superliniowa.
Wtedy $L$ jest skończona, ciągła, wypukła i superliniowa,
a dla każdego $p\in\R^d$ zachodzi
\begin{equation}\label{eq:pde23-duality}
 H(p)=\sup_{v\in\R^d}(p\cdot v-L(v)).
\end{equation}
Jeżeli $g$ jest globalnie $K$-lipschitzowska, to
$y\mapsto g(y)+tL((x-y)/t)$ osiąga minimum, a każda
jego minimalizująca wartość $v=(x-y)/t$ spełnia
$L(v)-K|v|\le L(0)$. W szczególności prędkości
minimalizatorów są ograniczone liczbą zależną tylko od
$H$ i $K$.
\end{lemat}
\begin{proof}
Ponieważ $H$ jest superliniowa, funkcja
$p\mapsto p\cdot v-H(p)$ dąży do $-\infty$ dla
$|p|\to\infty$, więc supremum definiujące $L(v)$
jest skończone i osiągane. Jako supremum funkcji
afinicznych $L$ jest wypukła i dolnie półciągła,
a skończona funkcja wypukła na całej $\R^d$ jest ciągła.
Dla dowolnego $R>0$ i $v\ne0$ wybieramy
$p=R v/|v|$, zatem $L(v)\ge R|v|-\max_{|p|=R}H(p)$.
Po podzieleniu przez $|v|$ i przejściu $|v|\to\infty$
wnioskujemy o superliniowości $L$, ponieważ $R$ jest
dowolne.

Z definicji $L$ mamy $p\cdot v-L(v)\le H(p)$.
Jeśli jednak $v=D H(p)$, wypukłość daje dla każdego $q$
$H(q)\ge H(p)+v\cdot(q-p)$, stąd
$L(v)=p\cdot v-H(p)$; powstaje równość w~\eqref{eq:pde23-duality}.

Dla ustalonych $x,t>0$ i $y=x-tv$ warunek Lipschitza
daje $g(x-tv)\ge g(x)-Kt|v|$. Ponieważ $L$ jest
superliniowa, funkcja $g(x-tv)+tL(v)$ dąży do
$+\infty$ przy $|v|\to\infty$; ciągłość zapewnia
osiągnięcie minimum. Porównując je z próbą $v=0$,
dostajemy
$g(x-tv)+tL(v)\le g(x)+tL(0)$.
Wykorzystanie nierówności Lipschitza daje ostatnią tezę.
\end{proof}

\begin{twierdzenie}[Wzór Hopfa--Laxa]
\label{thm:pde23-hopflax}
Niech $H\in C^1(\R^d)$ będzie wypukła i superliniowa,
a $g\colon\R^d\to\R$ niech będzie funkcją ograniczoną
i globalnie $K$-lipschitzowską. Dla $t>0$ zdefiniujmy
\begin{equation}\label{eq:pde23-hopflax}
 u(t,x)=\min_{y\in\R^d}\left\{
 g(y)+tL\!\left(\frac{x-y}{t}\right)\right\},
 \qquad u(0,x)=g(x).
\end{equation}
Wtedy $u$ jest ograniczona i globalnie Lipschitzowska
na $[0,T]\times\R^d$ dla każdego $T>0$ oraz jest
rozwiązaniem lepkościowym równania
$u_t+H(\nabla u)=0$. Ponadto
\begin{equation}\label{eq:pde23-dpp}
 u(t,x)=\min_{z\in\R^d}
 \left\{u(s,z)+(t-s)L\!\left(\frac{x-z}{t-s}\right)\right\}
 \quad (0\le s<t).
\end{equation}
\end{twierdzenie}
\begin{proof}
\textbf{Krok 1: ograniczenie minimalizatorów i ciągłość.}
Lemat~\ref{lem:pde23-legendre} daje osiąganie minimum
oraz wspólne ograniczenie $|x-y|/t\le R$.
Wybór dowolnego $v$ w równoważnej reprezentacji
$u(t,x)=\min_v\{g(x-tv)+tL(v)\}$ daje
$u(t,x)\le g(x)+tL(0)$, a nierówność Lipschitza
$g(x-tv)\ge g(x)-Kt|v|$ daje
$u(t,x)\ge g(x)+t\min_v(L(v)-K|v|)$.
Wynika stąd $\|u(t)-g\|_\infty\le Ct$.
Przy przesunięciu $x\mapsto x+h$ każda próba $v$
zmienia składnik $g(x-tv)$ o co najwyżej $K|h|$;
po wzięciu minimów otrzymujemy
$|u(t,x+h)-u(t,x)|\le K|h|$.

\textbf{Krok 2: zasada programowania dynamicznego.}
Oznaczmy $a=t-s$, $b=s$ i użyjmy wypukłości $L$.
Dla dowolnych $y,z\in\R^d$ mamy
\[
 aL\!\left(\frac{x-z}{a}\right)
 +bL\!\left(\frac{z-y}{b}\right)
 \ge (a+b)L\!\left(\frac{x-y}{a+b}\right).
\]
Minimalizacja kolejno po $y,z$ daje kierunek „$\ge$”
w~\eqref{eq:pde23-dpp}. Kierunek przeciwny otrzymujemy
przez wybór $z=y+\frac{s}{t}(x-y)$ na prostym odcinku
łączącym $y$ z $x$, kiedy obie prędkości są równe
$(x-y)/t$ i nierówność wypukłości staje się równością.
Przypadek $s=0$ wynika bezpośrednio z definicji.

Stosując wzór dynamiczny na przedziale długości $h>0$
i argument kroku pierwszego do funkcji $u(s,\cdot)$,
dostajemy $|u(s+h,x)-u(s,x)|\le Ch$ ze stałą
zależną tylko od $K,H$. Razem z oszacowaniem
przestrzennym dowodzi to globalnej ciągłości Lipschitza.

\textbf{Krok 3: nierówność dla funkcji dotykającej z góry.}
Niech $u-\varphi$ ma lokalne maksimum równe zero
w $(t,x)$, $t>0$. Ustalmy dowolne $v\in\R^d$.
Z~\eqref{eq:pde23-dpp} dla małych $h>0$ dostajemy
$u(t,x)\le u(t-h,x-hv)+hL(v)$.
Ponieważ $u(t-h,x-hv)\le\varphi(t-h,x-hv)$,
mamy $\varphi(t,x)\le\varphi(t-h,x-hv)+hL(v)$.
Po podzieleniu przez $h$ i przejściu do granicy
$\varphi_t+v\cdot\nabla\varphi\le L(v)$.
Warunek zachodzi dla każdego $v$; dzięki
\eqref{eq:pde23-duality} otrzymujemy
$\varphi_t+H(\nabla\varphi)\le0$.

\textbf{Krok 4: nierówność dla funkcji dotykającej z dołu.}
Niech $u-\varphi$ ma lokalne minimum równe zero
w $(t,x)$. Wybieramy minimalizator $z_h=x-hv_h$
w~\eqref{eq:pde23-dpp}. Wspólne oszacowanie Lipschitza
funkcji $u(t-h,\cdot)$ i Lemat~\ref{lem:pde23-legendre}
dają ograniczenie $|v_h|\le R$.
Z minimalności i dotykania z dołu wynika
$\varphi(t,x)\ge\varphi(t-h,x-hv_h)+hL(v_h)$.
Taylor z resztą $o(h)$ jednostajną dla $|v_h|\le R$
daje
$\varphi_t+v_h\cdot\nabla\varphi\ge L(v_h)+o(1)$.
Ale $L(v_h)-v_h\cdot\nabla\varphi\ge-H(\nabla\varphi)$,
więc $\varphi_t+H(\nabla\varphi)\ge o(1)$.
Przechodzimy do $h\downarrow0$.
Obie nierówności dowodzą rozwiązania lepkościowego.
\end{proof}

\begin{uwaga}[O jednoznaczności na całej przestrzeni]
Dla autonomicznego $H=H(p)$ powyższe rozwiązanie jest
jedynym ograniczonym, globalnie Lipschitzowskim
rozwiązaniem lepkościowym o danych $g$.
Dowód porównania na $\R^d$ otrzymuje się z dowodu
Twierdzenia~\ref{thm:pde23-comparison}, dodając do
funkcji podwajanej zmiennej składnik
$-\beta\sqrt{1+|x|^2}$, który zapewnia osiąganie
maksimum. Dla funkcji Lipschitzowskich gradienty funkcji
testowych w punktach styczności są jednostajnie
ograniczone. Nierówność dla $H$ wynika więc z jego
lokalnej lipschitzowskości w zmiennej pędowej, a dodatkowy
błąd ma rozmiar $O(\beta)$. Najpierw przechodzimy
$\eps,\alpha\downarrow0$, później $\beta\downarrow0$.
Jest to ta sama zasada porównania, z dodatkową lokalizacją
maksimum; nie korzysta z wypukłości $H$.
\end{uwaga}


\begin{wniosek}[Skończony obszar zależności we wzorze Hopfa--Laxa]
\label{cor:pde23-hopflax-local}
Niech $H$ spełnia założenia Twierdzenia~\ref{thm:pde23-hopflax},
a $g,h\colon\R^d\to\R$ będą ograniczone i globalnie
$K$-lipschitzowskie. Niech $u_g,u_h$ oznaczają funkcje
zdefiniowane przez wzór Hopfa--Laxa odpowiednio dla $g$ i $h$.
Istnieje $R_{H,K}<\infty$, zależne tylko od $H,K$, takie, że
jeśli $g=h$ na domkniętej kuli $\overline B_R(x_0)$, to
\[
 u_g(t,x)=u_h(t,x)
 \quad\text{dla }t>0,\qquad |x-x_0|+tR_{H,K}\le R.
\]
Jeżeli $R_{H,K}=0$, nierówność interpretujemy bez
ograniczenia czasu. Analogiczna własność zachodzi w chwili $t=0$.
\end{wniosek}
\begin{proof}
Z Lematu~\ref{lem:pde23-legendre} wszyscy minimalizatorzy
dla obu danych są postaci $y=x-tv$, gdzie
$L(v)-K|v|\le L(0)$. Ponieważ $L$ jest superliniowa,
ten podpoziom jest zwarty; wybieramy $R_{H,K}$ nie mniejsze
niż maksimum $|v|$ na nim. Gdy
$|x-x_0|+tR_{H,K}\le R$, każdy minimalizator dla $g$
i każdy minimalizator dla $h$ leży w $\overline B_R(x_0)$.
Na tej kuli obie funkcje minimalizowane są identyczne.
Minimum globalne każdej jest osiągane wewnątrz kuli,
więc obie wartości minimalne są równe.
\end{proof}


\section{Rozwiązanie na przeciętych charakterystykach i półwklęsłość}
\begin{twierdzenie}[Półwklęsłość w przypadku kwadratowym]
\label{thm:pde23-semiconcave}
Niech $H(p)=|p|^2/2$, $g$ będzie ograniczona i
Lipschitzowska, a $u$ będzie dane przez
\eqref{eq:pde23-hopflax}. Dla każdego $t>0$
funkcja $x\mapsto u(t,x)-|x|^2/(2t)$ jest wklęsła.
W szczególności $u(t,\cdot)$ jest lokalnie Lipschitzowska
i dwukrotnie różniczkowalna prawie wszędzie w sensie
Alexandrowa (ostatni fakt jest klasycznym twierdzeniem
o funkcjach półwklęsłych, którego tutaj nie dowodzimy).
\end{twierdzenie}
\begin{proof}
Ponieważ $L(v)=|v|^2/2$, mamy
\[
 u(t,x)-\frac{|x|^2}{2t}
 =\inf_{y\in\R^d}\left\{
 g(y)+\frac{|y|^2}{2t}-\frac{x\cdot y}{t}\right\}.
\]
Dla ustalonego $y$ prawa strona pod nawiasem jest
funkcją afiniczną względem $x$. Infimum funkcji
afinicznych jest funkcją wklęsłą. Lokalna
lipschitzowskość wynika tu również bezpośrednio
z Twierdzenia~\ref{thm:pde23-hopflax}.
\end{proof}

\begin{przyklad}[Utrata różniczkowalności i wybór lepkościowy]
\label{ex:pde23-shock}
Dla $g(x)=\cos x$ i $H(p)=p^2/2$ otrzymujemy
$u(t,x)=\min_y\{\cos y+(x-y)^2/(2t)\}$.
Jeśli $0<t<1$, mapa charakterystyk
$y\mapsto y-t\sin y$ jest ściśle rosnąca.
Gdy $t>1$ i $x=0$, funkcja minimalizowana jest parzysta;
punkt $y=0$ jest lokalnym maksimum, ponieważ druga pochodna
w $0$ wynosi $-1+1/t<0$. Dla $t$ dostatecznie bliskiego $1$
ma ona dwa symetryczne globalne minima $\pm a(t)$,
$a(t)>0$, spełniające $a(t)=t\sin a(t)$.
Dla $x$ dodatnich bliskich zera minimalizatorem jest
gałąź $+a(t)$, a dla ujemnych gałąź $-a(t)$.
Z twierdzenia o różniczkowaniu wartości minimalnej
(lub bezpośrednio z różnic ilorazowych) pochodne jednostronne
$u_x(t,0+)$ i $u_x(t,0-)$ wynoszą odpowiednio
$-a(t)/t$ i $a(t)/t$. Funkcja jest więc ciągła, lecz
nie jest różniczkowalna w $x=0$. Rozwiązanie lepkościowe
istnieje mimo przecięcia charakterystyk i jest jednoznaczne
w klasie ograniczonych rozwiązań Lipschitzowskich.
\end{przyklad}

\section{Związek z równaniem Burgersa i warunkiem entropii}
Jeżeli $u$ jest gładka i spełnia
$u_t+(u_x)^2/2=0$, to $w=u_x$ spełnia
$w_t+(w^2/2)_x=0$. Po powstaniu osobliwości
różniczkowanie można interpretować dystrybucyjnie.
Pochodna rozwiązania lepkościowego odpowiada
rozwiązaniu entropijnemu równania Burgersa dla
odpowiednich danych Lipschitzowskich.
Warunek entropii jest więc odpowiednikiem wyboru
lepkościowego na poziomie prawa zachowania.
Nie wyprowadzamy tutaj pełnego twierdzenia Krużkowa;
zostanie ono udowodnione w rozdziale o prawach zachowania.

\begin{stwierdzenie}[Różniczkowanie lepkiego równania Hamiltona--Jacobiego]
\label{prop:pde-v43-hj-viscous}
Niech $f\in C^2(\R)$, $\varepsilon>0$ oraz
$U\in C^{1,3}((0,T)\times\R;\R)$ spełnia
$U_t+f(U_x)=\varepsilon U_{xx}$ punktowo.
Wówczas $w=U_x\in C^{1,2}$ spełnia
$w_t+\partial_x f(w)=\varepsilon w_{xx}$.
Ponadto dla każdej wypukłej $\eta\in C^2(\R)$ i
strumienia $q$ określonego przez $q'=\eta' f'$ zachodzi
\[
 \partial_t\eta(w)+\partial_x q(w)
 =\varepsilon\partial_{xx}\eta(w)
 -\varepsilon\eta''(w)|w_x|^2
 \le\varepsilon\partial_{xx}\eta(w).
\]
\end{stwierdzenie}
\begin{proof}
Różniczkujemy równanie dla $U$ względem $x$:
$w_t+f'(w)w_x=\varepsilon w_{xx}$.
Po pomnożeniu przez $\eta'(w)$ otrzymujemy
$\partial_t\eta(w)+\partial_xq(w)
=\varepsilon\eta'(w)w_{xx}$.
Reguła łańcuchowa daje
$\eta'(w)w_{xx}=\partial_{xx}\eta(w)-\eta''(w)|w_x|^2$.
Ostatni wyraz jest niedodatni przez wypukłość $\eta$.
\end{proof}

\begin{uwaga}[Różniczkowanie nie jest jeszcze przejściem do granicy]
Powyższa tożsamość obowiązuje dla gładkich rozwiązań
\emph{z tą samą lepkością} $\varepsilon$ i łączy
rachunek w rozdziałach~\ref{chap:pde23} i~\ref{chap:pde24}.
Aby stwierdzić, że $\partial_xU$ jest rozwiązaniem
entropijnym dla granicznego $U$, trzeba ponadto
uzasadnić zbieżność pochodnych przy
$\varepsilon\downarrow0$ i zgodność danych początkowych;
sama zbieżność jednostajna $U^\varepsilon\to U$
nie zapewnia silnej zbieżności $U_x^\varepsilon$.
\end{uwaga}


\begin{uwaga}[Zakres wyników]
Twierdzenie o porównaniu udowodniliśmy na torusie
w założeniach~\eqref{eq:pde23-HxLip}; wzór Hopfa--Laxa
udowodniliśmy dla autonomicznego hamiltonianu wypukłego
i superliniowego. Ogólna teoria dla hamiltonianów
niewypukłych, zależnych od $x$ i sterowania optymalnego
wymaga osobnych hipotez i metod. Samo stosowanie
wzoru~\eqref{eq:pde23-hopflax} poza tymi założeniami
nie jest uzasadnione.
\end{uwaga}

\section{Wskazówki do dalszej lektury rozdziału}
W następnym rozdziale przejdziemy od potencjału $u$
do skalarnego prawa zachowania dla $w=u_x$, gdzie
pojawią się rozwiązania słabe, fale uderzeniowe,
warunki Rankine'a--Hugoniota i selekcja entropijna.
Charakterystyki i geometria ich przecinania
pozostaną podstawowym źródłem przykładów.

\chapter[Skalarne prawa zachowania]{Skalarne prawa zachowania: fale uderzeniowe, entropia i twierdzenie Krużkowa}
\label{chap:pde24}

\section{Od charakterystyk do rozwiązań słabych}
Niech $T>0$, $f\in C^2(\R)$ i $u_0\colon\R\to\R$ będzie
funkcją mierzalną. Rozważamy zagadnienie Cauchy'ego
\begin{equation}\label{eq:pde24-cl}
 u_t+\partial_x f(u)=0\quad\text{w }(0,T)\times\R,
 \qquad u(0,x)=u_0(x).
\end{equation}
Funkcję $f$ nazywamy \emph{strumieniem}. Dla gładkiego rozwiązania
równanie ma postać $u_t+f'(u)u_x=0$: wartości $u$ są stałe
wzdłuż charakterystyk o prędkości $f'(u)$. Po przecięciu
charakterystyk nie można na ogół kontynuować rozwiązania
w klasie $C^1$; musimy dopuścić nieciągłości.

Jest to sformułowanie całkowe wprowadzone już
w definicji~\ref{def:pde2-weak}, lecz w poniższym
rozdziale przyjmujemy silniejsze założenie lokalnej
ograniczoności rozwiązania, odpowiednie dla teorii entropijnej.
Nie jest to nowa, niezgodna definicja pojęcia słabości.

\begin{uwaga}[Potencjał a prawo zachowania]
\label{rem:pde-v43-hj-entropy-bridge}
Dla $f(w)=w^2/2$ związek z równaniem
Hamiltona--Jacobiego opisano w
rozdziale~\ref{chap:pde23}. Pełna rachunkowa zgodność
\emph{lepkich} równań jest udowodniona w
Stwierdzeniu~\ref{prop:pde-v43-hj-viscous}.
Selekcja rozwiązania entropijnego po zaniku lepkości
wymaga dodatkowej zwartości i jest odrębnym etapem,
którego nie zastępuje formalne różniczkowanie
rozwinięcia charakterystyk.
\end{uwaga}

\begin{definicja}[Rozwiązanie słabe]
Niech $u_0\in L^1_{\mathrm{loc}}(\R)$ i
$u\in L^\infty_{\mathrm{loc}}([0,T)\times\R)$.
Funkcja $u$ jest słabym rozwiązaniem~\eqref{eq:pde24-cl}, jeżeli
\begin{equation}\label{eq:pde24-weak}
 \int_0^T\!\int_\R
  \bigl(u\varphi_t+f(u)\varphi_x\bigr)\,\mathrm d\lambda_1(x)\,\mathrm d\lambda_1(t)
 +\int_\R u_0(x)\varphi(0,x)\,\mathrm d\lambda_1(x)=0
\end{equation}
dla każdej $\varphi\in C_c^\infty([0,T)\times\R)$.
Dla porównywania dwóch rozwiązań będziemy dodatkowo wymagać
$u\in C([0,T];L^1(\R))\cap L^\infty((0,T)\times\R)$.
\end{definicja}

\begin{stwierdzenie}[Bilans na ruchomym przedziale]
\label{prop:pde24-balance}
Niech $u\in C^1$ spełnia~\eqref{eq:pde24-cl} i niech
$a,b\in C^1([0,T])$ oraz $a(t)<b(t)$. Wówczas
\[
 \frac{\mathrm d}{\mathrm dt}\int_{a(t)}^{b(t)}u(t,x)\,\mathrm d\lambda_1(x)
 =f(u(t,a(t)))-f(u(t,b(t)))
 +b'(t)u(t,b(t))-a'(t)u(t,a(t)).
\]
\end{stwierdzenie}
\begin{proof}
Reguła Leibniza daje dwa składniki pochodzące z przesuwających
się końców przedziału. W pozostałej całce podstawiamy
$u_t=-\partial_xf(u)$ i stosujemy twierdzenie Newtona--Leibniza.
\end{proof}

\begin{definicja}[Wariacja całkowita i przestrzeń $BV$]
Niech $w\in L^1(\R)$. Definiujemy wariację całkowitą
\[
 \operatorname{TV}(w)=\sup\left\{
 \int_\R w(x)\phi'(x)\,\mathrm d\lambda_1(x):
 \phi\in C_c^1(\R),\ \|\phi\|_\infty\le1\right\}.
\]
Piszemy $w\in BV(\R)$, jeżeli $\operatorname{TV}(w)<\infty$.
Równoważnie pochodna dystrybucyjna $Dw$ jest skończoną
miarą ze znakiem, której wariacja całkowita jest równa
$\operatorname{TV}(w)$. Dla takiej funkcji zachodzi
$\|w(\cdot+h)-w\|_1\le |h|\operatorname{TV}(w)$.
Normę miary o skończonej wariacji oznaczamy przez
$\|\mu\|_{\mathcal M}=|\mu|(\R)$.
\end{definicja}

\section{Skoki i warunek Rankine'a--Hugoniota}
\begin{twierdzenie}[Warunek skoku]
\label{thm:pde24-RH}
Niech $s\in C^1((0,T))$ i niech $u_-,u_+$ będą funkcjami
klasy $C^1$ odpowiednio po lewej i prawej stronie krzywej
$x=s(t)$, mającymi jednostronne ślady $u_L(t),u_R(t)$.
Załóżmy, że każda gałąź spełnia klasycznie
$u_t+\partial_xf(u)=0$. Funkcja sklejona z tych gałęzi
jest rozwiązaniem dystrybucyjnym w otoczeniu krzywej wtedy
i tylko wtedy, gdy dla każdego $t$ zachodzi
\begin{equation}\label{eq:pde24-RH}
 s'(t)\bigl(u_R(t)-u_L(t)\bigr)
 =f(u_R(t))-f(u_L(t)).
\end{equation}
\end{twierdzenie}
\begin{proof}
Piszemy $u=u_-+(u_+-u_-)\chi_{\{x>s(t)\}}$.
Różniczkowanie dystrybucyjne daje człon osobliwy
$-s'(t)[u]\delta_{x=s(t)}$ w $u_t$ oraz
$[f(u)]\delta_{x=s(t)}$ w $\partial_xf(u)$,
gdzie $[h]=h_R-h_L$. Współczynnik przy delcie wynosi
$[f(u)]-s'[u]$. Znika on dokładnie przy~\eqref{eq:pde24-RH}.
Pozostałe części równania znikają na mocy założenia o obu
gałęziach. Równoważnie można scałkować~\eqref{eq:pde24-weak}
osobno po obu stronach krzywej i zebrać składniki brzegowe.
\end{proof}

Dla strumienia Burgersa $f(u)=u^2/2$, przy $u_L\ne u_R$,
prędkość skoku jest równa $(u_L+u_R)/2$. Sam warunek
Rankine'a--Hugoniota nie rozstrzyga jednak, czy taka fala
jest zgodna z fizycznym kierunkiem przepływu informacji.

\begin{kontrprzyklad}[Słabe rozwiązanie nie jest jednoznaczne]
\label{ex:pde24-expansion}
Dla $f(u)=u^2/2$ i danych Riemanna $u_0(x)=0$ dla $x<0$
oraz $u_0(x)=1$ dla $x>0$ funkcja
$u(t,x)=0$ dla $x<t/2$ i $u(t,x)=1$ dla $x>t/2$
spełnia warunek Rankine'a--Hugoniota, więc jest rozwiązaniem
słabym. Innym rozwiązaniem słabym tych samych danych jest
ciągły wachlarz $u(t,x)=0$ dla $x\le0$,
$u(t,x)=x/t$ dla $0<x<t$ oraz $u(t,x)=1$ dla $x\ge t$.
Dodatkowa zasada selekcji jest konieczna.
\end{kontrprzyklad}

\section{Entropie i poprawny kierunek fal uderzeniowych}
\begin{definicja}[Para entropijna]
Niech $f\in C^1(\R)$. Parą entropijną nazywamy funkcje
$\eta\in C^2(\R)$ oraz $q\in C^1(\R)$ takie, że $\eta$ jest
wypukła i $q'(z)=\eta'(z)f'(z)$. Rozwiązanie spełnia
nierówność entropijną, jeśli
$\partial_t\eta(u)+\partial_x q(u)\le0$ w sensie dystrybucji.
\end{definicja}

\begin{lemat}[Produkcja entropii przez lepkość]
\label{lem:pde24-viscentropy}
Jeżeli $u^\eps$ jest gładkim rozwiązaniem
$u_t^\eps+\partial_xf(u^\eps)=\eps u_{xx}^\eps$,
$\eps>0$, to dla każdej pary entropijnej zachodzi tożsamość
\begin{equation}\label{eq:pde24-entropyvisc}
 \partial_t\eta(u^\eps)+\partial_xq(u^\eps)
 =\eps\partial_{xx}\eta(u^\eps)
 -\eps\eta''(u^\eps)|u_x^\eps|^2
 \le\eps\partial_{xx}\eta(u^\eps).
\end{equation}
\end{lemat}
\begin{proof}
Mnożymy równanie przez $\eta'(u^\eps)$ i używamy reguły
łańcuchowej. Z równości
$\eta'(u)u_{xx}=\partial_{xx}\eta(u)-\eta''(u)u_x^2$
wynika wzór; znak ostatniego członu wynika z wypukłości.
\end{proof}

\begin{definicja}[Rozwiązanie entropijne Krużkowa]
\label{def:pde24-kru}
Niech $u_0\in L^1(\R)\cap L^\infty(\R)$.
Funkcja $u\in C([0,T];L^1(\R))\cap L^\infty((0,T)\times\R)$
jest \emph{rozwiązaniem entropijnym}~\eqref{eq:pde24-cl}, jeżeli
$u(0)=u_0$ w $L^1$ i dla każdego $k\in\R$ zachodzi
\begin{equation}\label{eq:pde24-kruzineq}
 \partial_t|u-k|+\partial_x\left(
   \operatorname{sgn}(u-k)(f(u)-f(k))\right)\le0
 \quad\text{w }\mathcal D'((0,T)\times\R),
\end{equation}
gdzie $\operatorname{sgn}(0)=0$. Rozumiemy to jako
nierówność po przetestowaniu nieujemną funkcją
$\varphi\in C_c^\infty((0,T)\times\R)$.
\end{definicja}

\begin{uwaga}
Użycie funkcji $|z-k|$ jest dopuszczalne przez aproksymację
gładkimi funkcjami wypukłymi. Z warunków dla wszystkich $k$
wynika również słabe równanie: jeśli $k$ jest większe od
istotnego supremum $u$, w~\eqref{eq:pde24-kruzineq}
uzyskujemy jedną nierówność liniową, a jeśli jest mniejsze
od istotnego infimum, nierówność przeciwną. Po zsumowaniu
otrzymujemy równość dystrybucyjną. Warunek początkowy
jest zadany oddzielnie przez ciągłość w $L^1$.
\end{uwaga}

\begin{twierdzenie}[Warunek entropijny na skoku]
\label{thm:pde24-entropyshock}
Niech $u_L\ne u_R$, $f\in C^2(\R)$ i niech skok między
stałymi stanami $u_L,u_R$ porusza się z prędkością
$\sigma=(f(u_R)-f(u_L))/(u_R-u_L)$.
Jest on entropijny w sensie lokalnych nierówności
Krużkowa wtedy i tylko wtedy, gdy dla każdego $k$ pomiędzy
$u_L$ i $u_R$ zachodzi
\begin{equation}\label{eq:pde24-chord}
 \begin{cases}
 f(k)\le \ell(k),&u_L>u_R,\\
 f(k)\ge \ell(k),&u_L<u_R,
 \end{cases}
 \qquad
 \ell(k)=f(u_L)+\sigma(k-u_L).
\end{equation}
W szczególności gdy $f$ jest ściśle wypukła, skok z
$u_L>u_R$ jest entropijny, natomiast skok z $u_L<u_R$
nie jest entropijny.
\end{twierdzenie}
\begin{proof}
Przyjmujemy $[h]=h_R-h_L$. Z rachunku dystrybucyjnego
nierówność dla entropii $\eta$ jest równoważna
$[q]-\sigma[\eta]\le0$. Jeżeli $k$ leży poza odcinkiem
wyznaczonym przez stany, wyrażenie to znika, ponieważ
$[q]=\pm[f]$ i $[\eta]=\pm[u]$.
Gdy $u_L>u_R$ i $u_R<k<u_L$, bezpośrednie podstawienie
daje
\[
 [q_k]-\sigma[|u-k|]=2(f(k)-\ell(k)).
\]
Przy $u_L<u_R$ otrzymujemy przeciwny znak.
Te warunki są konieczne i wystarczające. Ścisła wypukłość
umieszcza wykres $f$ ściśle poniżej cięciwy łączącej
odpowiednie punkty, co kończy dowód.
\end{proof}

\section{Zagadnienie Riemanna dla ściśle wypukłego strumienia}
\begin{twierdzenie}[Fala uderzeniowa i rozrzedzenia]
\label{thm:pde24-riemann}
Niech $f\in C^2(\R)$ i $f''(z)>0$ dla $z$ pomiędzy
$u_L,u_R$. Rozważmy dane $u_0=u_L$ dla $x<0$ oraz
$u_0=u_R$ dla $x>0$ (nie muszą należeć do $L^1$).
Dla $t>0$ istnieje samopodobne lokalnie ograniczone
rozwiązanie entropijne $u(t,x)=U(x/t)$:
\begin{enumerate}
\item Jeśli $u_L>u_R$, jest to skok przemieszczający się
z prędkością $\sigma=(f(u_L)-f(u_R))/(u_L-u_R)$.
\item Jeśli $u_L<u_R$, to
\[
 U(\xi)=\begin{cases}
 u_L,&\xi\le f'(u_L),\\
 (f')^{-1}(\xi),&f'(u_L)<\xi<f'(u_R),\\
 u_R,&\xi\ge f'(u_R).
 \end{cases}
\]
\end{enumerate}
Jeżeli $u_L=u_R$, rozwiązaniem jest funkcja stała $u(t,x)=u_L$.
Rozwiązania te są jednoznaczne w klasie lokalnie ograniczonych
rozwiązań entropijnych z lokalnym śladem początkowym,
gdy korzystamy z lokalnej wersji twierdzenia o kontrakcji
udowodnionego poniżej.
\end{twierdzenie}
\begin{proof}
W pierwszym przypadku warunek Rankine'a--Hugoniota
sprawia, że równanie jest spełnione słabo, a
Twierdzenie~\ref{thm:pde24-entropyshock} zapewnia
nierówności entropijne. W drugim przypadku $U$ jest
ciągła, a na wachlarzu spełnia $f'(U(\xi))=\xi$.
Obliczamy $u_t=-\xi U'(\xi)/t$ oraz
$\partial_x f(u)=f'(U)U'(\xi)/t$; składniki się znoszą.
Poza wachlarzem rozwiązanie jest stałe. Ponieważ funkcja
jest ciągła, w przejściu przez brzegi wachlarza nie
powstają wyrazy typu delta. Dla każdej wypukłej entropii
na wnętrzu wachlarza zachodzi równość
$\partial_t\eta(u)+\partial_xq(u)=0$; przejście do
entropii Krużkowa odbywa się przez aproksymację.
Warunek początkowy wynika z lokalnej zbieżności przy
$t\downarrow0$. Jednoznaczność wynika z lokalnej zasady
kontrakcji podanej na końcu rozdziału; dla równych stanów
stosujemy rozwiązanie stałe.
\end{proof}

\section{Nierówność Kato dla dwóch rozwiązań entropijnych}
\begin{twierdzenie}[Nierówność Kato--Krużkowa]
\label{thm:pde24-kato}
Niech $f\in C^1(\R)$ oraz niech
$u,v\in C([0,T];L^1(\R))\cap L^\infty((0,T)\times\R)$
będą rozwiązaniami entropijnymi z danymi $u_0,v_0$.
Wówczas w sensie dystrybucji na $(0,T)\times\R$
\begin{equation}\label{eq:pde24-kato}
 \partial_t|u-v|+\partial_x\bigl[
 \operatorname{sgn}(u-v)(f(u)-f(v))\bigr]\le0.
\end{equation}
\end{twierdzenie}
\begin{proof}
\textbf{Krok 1: podwojenie zmiennych.}
Dla ustalonej pary $(s,y)$ w nierówności entropijnej
dla $u(t,x)$ wybieramy $k=v(s,y)$, a dla ustalonej pary
$(t,x)$ w nierówności dla $v(s,y)$ wybieramy $k=u(t,x)$.
Są to dopuszczalne stałe, ponieważ w każdej nierówności
całkujemy tylko po zmiennych jej rozwiązania. Po
scalaniu nierówności twierdzeniem Fubiniego otrzymujemy
sumę dwóch nierówności na $(t,x,s,y)$.

\textbf{Krok 2: funkcja testowa.}
Niech $\rho\in C_c^\infty(\R)$ jest nieujemna, parzysta
i ma całkę równą $1$. Dla $\delta>0$ kładziemy
$\rho_\delta(z)=\delta^{-1}\rho(z/\delta)$.
Bierzemy $\psi\in C_c^\infty((0,T)\times\R)$, $\psi\ge0$,
i testujemy obie nierówności funkcją
\[
 \Phi_\delta(t,x,s,y)=
 \rho_\delta(t-s)\rho_\delta(x-y)\,
 \psi\!\left(\frac{t+s}{2},\frac{x+y}{2}\right).
\]
Dla małego $\delta$ nośnik leży we wnętrzu
$(0,T)^2\times\R^2$, zatem testowanie jest dozwolone.
Suma pochodnych spełnia dokładnie
\[
 (\partial_t+\partial_s)\Phi_\delta
 =\rho_\delta(t-s)\rho_\delta(x-y)\psi_t,
 \qquad
 (\partial_x+\partial_y)\Phi_\delta
 =\rho_\delta(t-s)\rho_\delta(x-y)\psi_x,
\]
gdzie pochodne $\psi$ obliczamy w zmiennych średnich.
W obu nierównościach występują te same wartości
$|u(t,x)-v(s,y)|$ i ten sam strumień
$\operatorname{sgn}(u(t,x)-v(s,y))
(f(u(t,x))-f(v(s,y)))$; w nierówności dla $v$
oba znaki zmieniają się jednocześnie.

\textbf{Krok 3: granica diagonalna.}
Dla $\delta\downarrow0$ funkcje
$u(t,x),v(s,y)$ można zastąpić ich przesunięciami
w czasie i przestrzeni. Twierdzenie o ciągłości translacji
w $L^1_{\mathrm{loc}}$ daje
\[
 |u(t,x)-v(s,y)|\to|u(t,x)-v(t,x)|
\]
w całkach z wagą $\Phi_\delta$. Strumień przechodzi
do granicy z tego samego powodu, gdyż na wspólnym
ograniczonym zakresie wartości $f$ jest Lipschitzowska,
a odwzorowanie
$(a,b)\mapsto\operatorname{sgn}(a-b)(f(a)-f(b))$
jest ciągłe i Lipschitzowskie. Ostatecznie
\[
 \int_0^T\!\int_\R
 \bigl(|u-v|\psi_t+
 \operatorname{sgn}(u-v)(f(u)-f(v))\psi_x\bigr)
 \,\mathrm d\lambda_1(x)\,\mathrm d\lambda_1(t)\ge0.
\]
Jest to~\eqref{eq:pde24-kato} w sensie dystrybucyjnym.
\end{proof}

\begin{twierdzenie}[Kontrakcja w $L^1$ i jednoznaczność]
\label{thm:pde24-L1}
Przy założeniach Twierdzenia~\ref{thm:pde24-kato} zachodzi
\begin{equation}\label{eq:pde24-L1}
 \|u(t)-v(t)\|_{L^1(\R)}
 \le\|u_0-v_0\|_{L^1(\R)},\qquad 0\le t\le T.
\end{equation}
W szczególności dane początkowe wyznaczają co najwyżej
jedno rozwiązanie entropijne w powyższej klasie.
\end{twierdzenie}
\begin{proof}
Oznaczmy $w=|u-v|\in C([0,T];L^1)$ i
$q=\operatorname{sgn}(u-v)(f(u)-f(v))$. Jeśli
$M=\max_{|z|\le B}|f'(z)|$, gdzie $B$ ogranicza
$|u|,|v|$, to $|q|\le Mw$. W nierówności Kato
testujemy nieujemną funkcją aproksymującą
$\chi_{[s,t]}(\tau)\zeta_R(x)$, gdzie $\zeta_R$ jest
równa $1$ na $[-R,R]$, równa $0$ poza $[-2R,2R]$
i spełnia $|\zeta_R'|\le C/R$. Przechodząc do granicy
w aproksymacji czasowej, otrzymujemy
\[
 \int_\R w(t,x)\zeta_R(x)\,\mathrm d\lambda_1(x)
 \le \int_\R w(s,x)\zeta_R(x)\,\mathrm d\lambda_1(x)
 +\frac{CM}{R}\int_s^t\!\int_\R w(\tau,x)
 \,\mathrm d\lambda_1(x)\,\mathrm d\tau.
\]
Gdy $R\to\infty$, ostatni człon znika, ponieważ
$w$ jest ciągła w $L^1$ na zwartym przedziale czasowym.
Otrzymujemy kontrakcję między dowolnymi $s,t$,
a następnie podstawiamy $s=0$.
\end{proof}

\begin{lemat}[Lokalna postać nierówności Kato]
\label{lem:pde24-local-kato}
Niech $f\in C^1(\R)$, a funkcje $u,v$ niech należą do
\[
 C([0,T];L^1_{\mathrm{loc}}(\R))\cap
 L^\infty_{\mathrm{loc}}([0,T]\times\R).
\]
Załóżmy, że spełniają lokalnie wszystkie nierówności entropijne
Krużkowa i mają lokalne ślady początkowe $u_0,v_0$.
Wówczas nierówność~\eqref{eq:pde24-kato} zachodzi
w $\mathcal D'((0,T)\times\R)$. Do wyprowadzenia
lokalnej kontrakcji na zwężających się przedziałach
wystarcza ograniczenie $|f'|$ na zakresie wartości
$u,v$ w rozważanym stożku.
\end{lemat}
\begin{proof}
Każda funkcja testowa $\psi\in C_c^\infty((0,T)\times\R)$
ma nośnik wewnątrz pewnego zwartego walca.
W dowodzie Twierdzenia~\ref{thm:pde24-kato}
całkujemy cztery zmienne tylko w niewielkim powiększeniu
tego walca, bo funkcje
$\rho_\delta(t-s)\rho_\delta(x-y)$ mają małe nośniki.
Zatem potrzebujemy jedynie ograniczeń $L^\infty$
i ciągłości translacji w $L^1$ na tym walcu,
a nie całkowalności $u,v$ na całej prostej.
Przejście diagonalne daje nierówność Kato lokalnie.
W testowaniu przez przybliżenie wskaźnika stożka
stosujemy analogiczne odcięcie przestrzenne; używane
granice czasu wynikają z założonej ciągłości w
$L^1_{\mathrm{loc}}$. Warunek $|f'|\le M$ na
wspólnym zakresie wartości zapewnia $|q|\le M|u-v|$.
\end{proof}

\begin{wniosek}[Skończona prędkość propagacji]
\label{cor:pde24-finite}
Niech $u,v$ będą lokalnie ograniczonymi rozwiązaniami
entropijnymi z lokalnymi śladami początkowymi w $L^1$.
Jeśli $M$ ogranicza $|f'|$ na zakresach ich wartości,
to dla każdych $a<b$ i $0\le t\le T$ zachodzi
\begin{equation}\label{eq:pde24-cone}
 \int_a^b|u(t,x)-v(t,x)|\,\mathrm d\lambda_1(x)
 \le\int_{a-Mt}^{b+Mt}|u_0(x)-v_0(x)|\,\mathrm d\lambda_1(x).
\end{equation}
\end{wniosek}
\begin{proof}
Na mocy Lematu~\ref{lem:pde24-local-kato}
możemy zastosować nierówność Kato lokalnie. Używamy
funkcji testowych przybliżających wskaźnik zwężającego się przedziału
$[a-M(t-\tau),b+M(t-\tau)]$ dla $0\le\tau\le t$.
Jego lewy koniec $L(\tau)=a-M(t-\tau)$ porusza się
z prędkością $M$, a prawy $R(\tau)=b+M(t-\tau)$
z prędkością $-M$. Oznaczając
$\zeta(\tau,x)=\chi_{[L(\tau),R(\tau)]}(x)$,
otrzymujemy w sensie dystrybucji względem $x$:
$\partial_\tau\zeta=-M(\delta_{L(\tau)}+\delta_{R(\tau)})$
oraz $\partial_x\zeta=\delta_{L(\tau)}-\delta_{R(\tau)}$.
Dlatego na obu brzegach współczynnik przy nierówności
Kato ma właściwy znak: odpowiednio $q-Mw\le0$
oraz $-q-Mw\le0$, ponieważ $|q|\le Mw$.
Aby nie podstawiać do dystrybucji funkcji niegładkiej,
wygładzamy oba końce przedziału przestrzennie i chwilę
końcową czasowo, stosujemy nierówność Kato do nieujemnych
funkcji testowych o zwartym nośniku, a potem przechodzimy
do granicy. Wszystkie całki są lokalne, a ciągłość
$u,v$ w $L^1$ na ograniczonych przedziałach zapewnia
zbieżność wyrazów czasowych. Otrzymujemy~\eqref{eq:pde24-cone}.
\end{proof}

\section{Istnienie przez zanikającą lepkość}
\begin{twierdzenie}[Istnienie dla danych $BV$]
\label{thm:pde24-existence}
Niech $f\in C^2(\R)$ oraz
$u_0\in BV(\R)\cap L^1(\R)\cap L^\infty(\R)$.
Wtedy dla każdego $T>0$ istnieje dokładnie jedno
rozwiązanie entropijne $u\in C([0,T];L^1(\R))\cap L^\infty$
z danymi $u_0$. Spełnia ono
\begin{equation}\label{eq:pde24-BVbound}
 \|u(t)\|_\infty\le\|u_0\|_\infty,
 \qquad \operatorname{TV}(u(t))\le\operatorname{TV}(u_0),
 \qquad \|u(t)\|_1\le\|u_0\|_1.
\end{equation}
\end{twierdzenie}
\begin{proof}
\textbf{Krok 1: aproksymacja lepkościowa.}
Zastępujemy $f$ przez $f-f(0)$, co nie zmienia równania,
i oznaczamy nadal tę funkcję przez $f$.
Niech $u_0^\eps=G_{\eps^2}*u_0$, gdzie $G_{\eps^2}$ jest
nieujemnym gładkim jądrem o masie $1$ i szerokości
$O(\eps)$. Rozwiązujemy
\begin{equation}\label{eq:pde24-visc}
 u_t^\eps+\partial_x f(u^\eps)=\eps u_{xx}^\eps,
 \qquad u^\eps(0)=u_0^\eps.
\end{equation}
Dla ustalonego $\eps>0$ rozwiązanie otrzymujemy
lokalnie z wzoru Duhamela
\[
 u^\eps(t)=G_{\eps t}*u_0^\eps
 -\int_0^t\partial_xG_{\eps(t-s)}*f(u^\eps(s))\,\mathrm ds,
\]
gdzie $G_t(x)=(4\pi t)^{-1/2}e^{-x^2/(4t)}$.
Oszacowanie $\|\partial_xG_{\eps t}\|_1=C(\eps t)^{-1/2}$
zapewnia kontrakcję odwzorowania Picarda na krótkim
przedziale w normie supremum, po odpowiednim lokalnym
ograniczeniu $f'$; zasada maksimum daje
$\|u^\eps(t)\|_\infty\le\|u_0\|_\infty$, co pozwala
kontynuować rozwiązanie do dowolnego czasu.

\textbf{Krok 2: kontrakcja lepkościowa.}
Dla dwóch gładkich rozwiązań $a,b$ równania lepkościowego
ich różnica $w=a-b$ spełnia
$w_t+\partial_x(f(a)-f(b))=\eps w_{xx}$.
Przyjmijmy $\eta_\delta(w)=\sqrt{w^2+\delta^2}-\delta$.
Po pomnożeniu przez $\eta_\delta'(w)$ i całkowaniu przez
części, korzystając z $\eta_\delta''\ge0$, otrzymujemy
w granicy $\delta\downarrow0$ nierówność Kato
\[
 \partial_t|a-b|+\partial_x\bigl[
 \operatorname{sgn}(a-b)(f(a)-f(b))\bigr]
 \le\eps\partial_{xx}|a-b|.
\]
Aby uzasadnić przejście do granicy, piszemy
$f(a)-f(b)=A(a,b)(a-b)$, gdzie
$A(a,b)=\int_0^1f'(b+\theta(a-b))\,\mathrm d\theta$.
Pozostały po różniczkowaniu strumienia błąd zawiera
$[\eta_\delta'(w)w-\eta_\delta(w)]\partial_xA$;
wartość nawiasu jest nie większa niż $\delta$.
Dla gładkich $a,b$ i na zbiorach zwartych błąd dąży
do zera. Testowanie przez odcięcia przestrzenne
i granica odcięć dowodzą
$\|a(t)-b(t)\|_1\le\|a(0)-b(0)\|_1$.
Stosując to do $b\equiv0$, otrzymujemy ograniczenie
normy $L^1$. Stosując do przesunięć $u^\eps(t,x+h)$
i $u^\eps(t,x)$, dostajemy
\[
 \|u^\eps(t,\cdot+h)-u^\eps(t,\cdot)\|_1
 \le\|u_0^\eps(\cdot+h)-u_0^\eps\|_1
 \le |h|\operatorname{TV}(u_0).
\]
Charakterystyka funkcji $BV$ przez przesunięcia
implikuje $\operatorname{TV}(u^\eps(t))\le\operatorname{TV}(u_0)$.

\textbf{Krok 3: jednostajny moduł ciągłości w czasie.}
Z tożsamości Duhamela stosowanej od czasu $t$ do $t+h$
i z $\|\partial_x f(u^\eps(t))\|_{\mathcal M}\le
M\operatorname{TV}(u_0)$, gdzie
$M=\max_{|z|\le\|u_0\|_\infty}|f'(z)|$, wynika
\[
 \|u^\eps(t+h)-u^\eps(t)\|_1
 \le C\operatorname{TV}(u_0)\sqrt{\eps h}
 +Mh\operatorname{TV}(u_0),\qquad 0<\eps\le1.
\]
Pierwszy składnik to oszacowanie
$\|G_{\eps h}*w-w\|_1\le
\operatorname{TV}(w)\int_\R|z|G_{\eps h}(z)\,\mathrm dz$;
w drugim korzystamy z tego, że półgrupa Gaussa jest
kontrakcją na miarach o skończonej wariacji.

\textbf{Krok 4: zwartość i granica.}
Ograniczenie w $BV$ i $L^1$ zapewnia względną zwartość
zbioru $\{u^\eps(t)\}$ w $L^1$ na każdym ograniczonym
przedziale przestrzennym. Wspólny moduł ciągłości w
czasie oraz diagonalny argument Arzeli--Ascoliego dają
zbieżność podciągu w $C([0,T];L^1_{\mathrm{loc}}(\R))$.
Aby przejść do zbieżności globalnej, wybieramy
$\zeta_R\in C^\infty(\R)$ równą $0$ na $[-R,R]$ i
równą $1$ poza $[-2R,2R]$, przy
$|\zeta_R'|\le C/R$, $|\zeta_R''|\le C/R^2$.
Nierówność Kato dla $a=u^\eps$, $b=0$ daje
\[
 \int_\R|u^\eps(t,x)|\zeta_R(x)\,\mathrm d\lambda_1(x)
 \le\int_\R|u_0^\eps(x)|\zeta_R(x)\,\mathrm d\lambda_1(x)
 +CT\|u_0\|_1\left(\frac M R+\frac1{R^2}\right).
\]
Prawa strona dąży do zera dla $R\to\infty$ jednostajnie
w $\eps\le1$ (najpierw dla danych początkowych, potem
przez jednostajną zbieżność $u_0^\eps\to u_0$ w $L^1$).
Zatem zbieżność zachodzi w $C([0,T];L^1(\R))$.

\textbf{Krok 5: równanie i nierówności entropijne.}
Wzór słaby dla $u^\eps$ przechodzi do granicy, gdyż
$f$ jest Lipschitzowska na wspólnym zakresie wartości,
a składnik lepkościowy ma postać
$\eps\int_0^T\int_{\R}u^\eps(t,x)\varphi_{xx}(t,x)\,\mathrm d\lambda_1(x)\,\mathrm d\lambda_1(t)$ i dąży do zera.
Dla gładkiej wypukłej entropii korzystamy
z~\eqref{eq:pde24-entropyvisc}. Po przetestowaniu
nieujemną $\varphi$ ostatnia prawa strona
$\eps\int_0^T\int_{\R}\eta(u^\eps(t,x))\varphi_{xx}(t,x)\,\mathrm d\lambda_1(x)\,\mathrm d\lambda_1(t)$ dąży do zera,
a silna zbieżność w $L^1_{\mathrm{loc}}$ pozwala
przenieść lewą stronę do granicy. Aproksymacja
$\eta(z)=|z-k|$ gładkimi funkcjami wypukłymi daje
warunki Krużkowa. Zbieżność w $C([0,T];L^1)$
gwarantuje $u(0)=u_0$. Ograniczenia~\eqref{eq:pde24-BVbound}
przechodzą do granicy przez dolną półciągłość wariacji
i zbieżność $L^1$, a jednoznaczność wynika z
Twierdzenia~\ref{thm:pde24-L1}. Ponieważ limit jest
jednoznaczny, zbiega cała rodzina $u^\eps$.
\end{proof}


\begin{twierdzenie}[Istnienie dla danych $L^1\cap L^\infty$ bez założenia $BV$]
\label{thm:pde24-existence-L1Linf}
Niech $f\in C^2(\R)$, $u_0\in L^1(\R)\cap L^\infty(\R)$
i $T>0$. Wtedy istnieje dokładnie jedno rozwiązanie entropijne
$u\in C([0,T];L^1(\R))\cap L^\infty((0,T)\times\R)$
z warunkiem $u(0)=u_0$. Spełnia ono
\[
 \|u(t)\|_\infty\le\|u_0\|_\infty,\qquad
 \|u(t)\|_1\le\|u_0\|_1
 \quad(0\le t\le T).
\]
Jeżeli dodatkowo $u_0\in BV(\R)$, to
$\operatorname{TV}(u(t))\le\operatorname{TV}(u_0)$.
\end{twierdzenie}
\begin{proof}
Wybieramy ciąg $u_0^n\in C_c^\infty(\R)$ taki, że
$u_0^n\to u_0$ w $L^1(\R)$ oraz
$\|u_0^n\|_\infty\le\|u_0\|_\infty$: wystarczy
obciąć $u_0$ przez funkcje z $[0,1]$, a następnie
zastosować splot z nieujemnym jądrem aproksymującym jedność.
Każde $u_0^n$ należy do $BV\cap L^1\cap L^\infty$,
a z Twierdzenia~\ref{thm:pde24-existence} istnieje
odpowiadające mu rozwiązanie entropijne $u^n$.
Oszacowanie~\eqref{eq:pde24-L1} daje
\[
 \sup_{0\le t\le T}\|u^n(t)-u^m(t)\|_1
 \le\|u_0^n-u_0^m\|_1.
\]
Zupełność $C([0,T];L^1(\R))$ zapewnia więc istnienie
$u$ z $u^n\to u$ w tej przestrzeni. Z jednostajnej
granicy w $L^1$ i ograniczeń $L^\infty$ rozwiązań $u^n$
wynika, po przejściu do podciągu prawie wszędzie,
$\|u\|_{L^\infty((0,T)\times\R)}\le\|u_0\|_\infty$.
Ponieważ $f$ i strumienie Krużkowa
$q_k(z)=\operatorname{sgn}(z-k)(f(z)-f(k))$
są Lipschitzowskie na wspólnym ograniczonym zakresie
wartości (dla każdego ustalonego $k$), możemy przejść
do granicy we wszystkich nierównościach entropijnych
na testach o zwartym nośniku. Ciągłość w $L^1$
zapewnia ślad $u(0)=u_0$. Ograniczenie normy $L^1$
przechodzi do granicy; jednoznaczność jest konsekwencją
Twierdzenia~\ref{thm:pde24-L1}. Jeżeli $u_0\in BV$,
możemy też skorzystać wprost z Twierdzenia
\ref{thm:pde24-existence}, które daje ograniczenie wariacji.
\end{proof}


\begin{uwaga}[O danych wyłącznie ograniczonych]
Twierdzenie~\ref{thm:pde24-existence-L1Linf} usuwa założenie
$BV$ dla danych całkowalnych. W ogólnej teorii Krużkowa
istnienie i jednoznaczność zachodzą także dla danych jedynie
$L^\infty$, przy odpowiednim rozumieniu lokalnych śladów.
Tego ostatniego przypadku nie wyprowadziliśmy z kontrakcji
globalnej $L^1$; wymaga on lokalizacji i osobnego argumentu.
\end{uwaga}

\section{Związek z równaniem Hamiltona--Jacobiego i dalsze problemy}
Dla $f(w)=w^2/2$ i funkcji $g\in W^{1,\infty}(\R)$
rozwiązanie lepkościowe
$U_t+\tfrac12(U_x)^2=0$ ma pochodną
$w=U_x$, która jest rozwiązaniem entropijnym
równania Burgersa, w odpowiednim sensie dystrybucyjnym.
Przyjmujemy tutaj ten fakt jako konsekwencję
przybliżenia zanikającą lepkością i jednoznaczności
wyboru entropijnego; dla ogólnych danych wymagany jest
odrębny argument zgodności dwóch granic.

\begin{przyklad}[Dyssypacja energii w fali Burgersa]
Dla $f(w)=w^2/2$ wybieramy $\eta(w)=w^2/2$, wtedy
$q(w)=w^3/3$. Przy skoku z $u_L=1$ do $u_R=0$
prędkość wynosi $\sigma=1/2$, a współczynnik
entropijnej delty jest równy
$[q]-\sigma[\eta]=-1/3+1/4=-1/12<0$.
Dla odwróconego skoku współczynnik jest dodatni,
więc warunek entropii zostaje naruszony.
\end{przyklad}

\begin{uwaga}[Zakres twierdzeń]
Twierdzenia o kontrakcji $L^1$ i istnieniu dotyczą
rozwiązań na prostej o danych z $L^1\cap L^\infty$;
dla ograniczenia wariacji początkowej dodatkowo zakładamy $BV$. Rozwiązania
Riemanna, których stany nie zanikają w nieskończoności,
rozpatrujemy w klasie lokalnej i stosujemy do nich
oszacowanie w stożku. Wyniki dla układów praw zachowania,
strumieni zależnych od $x$, wyższych wymiarów lub
nie-wypukłych zadań Riemanna wymagają dodatkowych metod.
\end{uwaga}



% ROZDZIAL 25: nowe rownania pololiniowe
\chapter[Równania półliniowe i eksplozja]{Równania półliniowe: lokalne istnienie, przedłużanie, eksplozja i energia}
\label{chap:pde25}

\section{Model, pojęcia i konieczność doboru przestrzeni}
Równanie jest \emph{półliniowe}, gdy najwyższe pochodne niewiadomej
występują liniowo, a ich współczynniki nie zależą od niewiadomej.
Przykładami są $u_t-\Delta u=f(u)$ oraz
$u_{tt}-\Delta u=f(u)$; natomiast $u_t-\diver(a(u)\nabla u)=0$
jest na ogół równaniem quasi-liniowym.
Dla zwięzłości pełny dowód lokalnego istnienia przeprowadzimy
najpierw na płaskim torusie $\T^d=\R^d/(2\pi\mathbb Z)^d$.
Przez $C(\T^d)$ rozumiemy rzeczywistą przestrzeń Banacha funkcji
ciągłych z normą supremum. W przeciwieństwie do znormalizowanej
normy $L^2(\T)$ z rozdziału~\ref{chap:pde-fourier},
tu używamy nieznormalizowanej miary objętości. Miara $\lambda_d$ na torusie jest
miarą Lebesgue'a przeniesioną z kostki $[0,2\pi)^d$;
$|\T^d|=(2\pi)^d$. Na torusie nie ma warunków brzegowych,
więc nieliniowość i mechanizmy rozwoju osobliwości nie są zasłaniane
kwestią regularności brzegu.

\begin{definicja}[Rozwiązanie łagodne]
\label{def:pde25-mild}
Niech $u_0\in C(\T^d)$, $f\in C^1(\R)$ i $T>0$.
Funkcję $u\in C([0,T];C(\T^d))$ nazywamy
\emph{rozwiązaniem łagodnym} równania $u_t-\Delta u=f(u)$
z warunkiem $u(0)=u_0$, jeżeli
\begin{equation}\label{eq:pde25-duhamel}
 u(t)=S(t)u_0+\int_0^tS(t-s)f(u(s))\,\mathrm ds,
 \qquad 0\le t\le T,
\end{equation}
gdzie $S(t)=e^{t\Delta}$ jest półgrupą cieplną na $\T^d$,
a $f(u(s))$ oznacza złożenie funkcji $f$ z $u(s,\cdot)$.
Całka w~\eqref{eq:pde25-duhamel} jest całką Bochnera
w przestrzeni Banacha $C(\T^d)$.
\end{definicja}

\begin{lemat}[Półgrupa ciepła na torusie]
\label{lem:pde25-heat}
Dla każdego $d\ge1$ operatory $S(t)$, $t\ge0$, tworzą silnie
ciągłą półgrupę na $C(\T^d)$. Są dodatnie, zachowują stałe
i spełniają $\|S(t)v\|_\infty\le\|v\|_\infty$.
Ponadto $\int_{\T^d}S(t)v\,\mathrm d\lambda_d
=\int_{\T^d}v\,\mathrm d\lambda_d$.
\end{lemat}
\begin{proof}
Dla $t>0$ periodyzujemy jądro Gaussa:
$K_t(z)=\sum_{k\in\mathbb Z^d}(4\pi t)^{-d/2}
\exp(-|z+2\pi k|^2/(4t))$.
Szereg wraz z dowolną liczbą pochodnych zbiega jednostajnie
dla ustalonego $t>0$. Mamy $K_t\ge0$ i
$\int_{\T^d}K_t(z)\,\mathrm d\lambda_d(z)=1$;
stąd dodatniość, kontrakcja i zachowanie stałych dla
$S(t)v=K_t*v$. Splot jąder Gaussa na $\R^d$ oraz
periodyzacja dają $S(t)S(s)=S(t+s)$. Dla $v$ ciągłej
i jednostajnie ciągłej na torusie rozdzielamy całkę splotową
na $|z|<\delta$ i $|z|\ge\delta$ (odległość liczona na torusie).
Pierwszą część oszacowujemy przez moduł ciągłości $v$,
a druga dąży do zera, ponieważ masy jąder poza każdym
ustalonym otoczeniem zera zanikają przy $t\downarrow0$.
Zatem $\|S(t)v-v\|_\infty\to0$. Tożsamość całek wynika
z twierdzenia Fubiniego i jednostkowej masy jądra.
\end{proof}

\begin{uwaga}[Dlaczego nie zaczynamy od $L^2$?]
Dla wielomianowej nieliniowości $f(u)=u^3$ operator
Nemyckiego $u\mapsto f(u)$ nie jest lokalnie Lipschitzowski
z $L^2(\T^d)$ do $L^2(\T^d)$: funkcje o coraz wyższych
wąskich pikach mogą mieć małą normę $L^2$, a ich sześciany
ogromną normę $L^2$. Twierdzenie Picarda w tej przestrzeni
nie byłoby więc uzasadnione. Przestrzeń $C(\T^d)$ pozwala
przeprowadzić argument bez założeń o wykładnikach Sobolewa.
\end{uwaga}

\section{Twierdzenie Picarda i kryterium przedłużania}
\begin{twierdzenie}[Lokalne istnienie i jednoznaczność]
\label{thm:pde25-local}
Niech $d\ge1$, $u_0\in C(\T^d)$ i $f\in C^1(\R)$.
Istnieje $T_0>0$ i dokładnie jedno rozwiązanie łagodne
$u\in C([0,T_0];C(\T^d))$ równania~\eqref{eq:pde25-duhamel}.
Rozwiązanie posiada jednoznaczne maksymalne przedłużenie
na przedział $[0,T_{\max})$, gdzie $T_{\max}\in(0,\infty]$.
Jeżeli $T_{\max}<\infty$, to
\begin{equation}\label{eq:pde25-blowalt}
 \limsup_{t\uparrow T_{\max}}\|u(t)\|_\infty=\infty.
\end{equation}
W rzeczywistości w takim przypadku norma dąży do
nieskończoności w sensie granicy.
\end{twierdzenie}
\begin{proof}
\textbf{Krok 1: wybór kuli.} Ustalmy $R>\|u_0\|_\infty+1$.
Na $[-R,R]$ przyjmijmy
$L_R=\max_{|z|\le R}|f'(z)|$ i
$M_R=\max_{|z|\le R}|f(z)|$.
W zupełnej przestrzeni $Z_T=C([0,T];C(\T^d))$
rozważmy domkniętą kulę
$B_R=\{v:\|v\|_{Z_T}\le R\}$.
Definiujemy $\Phi(v)$ jako prawą stronę~\eqref{eq:pde25-duhamel}.
Lemat~\ref{lem:pde25-heat} daje
$\|\Phi(v)(t)\|_\infty\le\|u_0\|_\infty+TM_R$.
Wybierając $T$ tak, by $TM_R\le R-\|u_0\|_\infty$,
otrzymujemy $\Phi(B_R)\subseteq B_R$.

\textbf{Krok 2: kontrakcja.} Twierdzenie o wartości średniej
i kontrakcja półgrupy implikują
$\|\Phi(v)-\Phi(w)\|_{Z_T}\le TL_R\|v-w\|_{Z_T}$.
Jeśli $TL_R<1$, zasada Banacha daje jedyny punkt stały
w $B_R$. Jednoznaczność nie zależy od wybranej kuli:
dla dowolnych dwóch rozwiązań na wspólnym przedziale
ich zakresy są ograniczone przez pewne $R$, a nierówność
$\|u(t)-v(t)\|_\infty\le
L_R\int_0^t\|u(s)-v(s)\|_\infty\,\mathrm ds$
i lemat Gronwalla wymuszają równość.

\textbf{Krok 3: sklejanie.} Półgrupowość pozwala przepisać
wzór Duhamela od dowolnego czasu $t_0$:
$u(t_0+t)=S(t)u(t_0)+\int_0^tS(t-s)f(u(t_0+s))\,\mathrm ds$.
To daje przedłużanie i maksymalny czas istnienia.

\textbf{Krok 4: alternatywa eksplozji.} Jeśli $T_{\max}<\infty$
i $\|u(t)\|_\infty\le R$ dla wszystkich $t$ bliskich
$T_{\max}$, to w konstrukcji z kroku pierwszego można
zastosować stale tę samą większą kulę oraz ten sam czas
lokalnego istnienia $\tau_R>0$, niezależny od chwili restartu.
Wybierając $t_0\in(T_{\max}-\tau_R/2,T_{\max})$
przedłużamy rozwiązanie poza $T_{\max}$, sprzeczność.
Ten sam argument pokazuje, że nie może istnieć ciąg
$t_n\uparrow T_{\max}$ z $\|u(t_n)\|_\infty\le R$,
ponieważ restart w $t_n$ daje przedłużenie o stały czas.
Zatem norma rzeczywiście dąży do nieskończoności.
\end{proof}

\begin{twierdzenie}[Stabilność i ciągła zależność od danych]
\label{thm:pde25-stability}
Niech $u,v$ będą rozwiązaniami łagodnymi na $[0,T]$,
z początkami $u_0,v_0\in C(\T^d)$, dla tej samej
$f\in C^1(\R)$. Jeśli
$\max\{\|u\|_{Z_T},\|v\|_{Z_T}\}\le R$, to
$\|u(t)-v(t)\|_\infty\le e^{L_Rt}\|u_0-v_0\|_\infty$
dla każdego $0\le t\le T$.
\end{twierdzenie}
\begin{proof}
Odejmujemy wzory Duhamela, używamy oszacowania
$|f(a)-f(b)|\le L_R|a-b|$ dla $|a|,|b|\le R$,
a następnie lematu Gronwalla w postaci całkowej.
\end{proof}


\begin{uwaga}[Przenoszenie konstrukcji między rozdziałami]
\label{rem:pde-v44-picard-scaling}
Dowód kontrakcji z twierdzenia~\ref{thm:pde25-local} opiera się na
silnej ciągłości i kontraktywności półgrupy w $C(\T^d)$ oraz na
lokalnej lipschitzowskości operatora Nemyckiego w tej normie.
W rozdziale~\ref{chap:pde27} użyjemy dokładnie tego samego
mechanizmu na torusie o okresie $L$ i z dyfuzją $D>0$;
zmiana okresu i czasu zostanie tam uzasadniona wzorem
\eqref{eq:pde-v44-torus-scale}. Nie jest natomiast poprawne
przeniesienie dowodu bezpośrednio na $L^2$ dla dowolnej
nieliniowości wielomianowej: odwzorowanie $v\mapsto v^3$
nie musi przyjmować wartości w $L^2$ ani być lokalnie
lipschitzowskie w tej przestrzeni.
\end{uwaga}


\section{Zasada porównania, dodatniość i istnienie globalne}
\begin{twierdzenie}[Porównanie rozwiązań półliniowych]
\label{thm:pde25-comparison}
Niech $f\in C^1(\R)$, $T>0$, a $u,v\in Z_T$
będą rozwiązaniami łagodnymi z danymi $u_0,v_0$.
Jeżeli $u_0(x)\le v_0(x)$ dla każdego $x\in\T^d$,
to $u(t,x)\le v(t,x)$ dla wszystkich
$(t,x)\in[0,T]\times\T^d$.
Nie zakładamy monotoniczności $f$.
\end{twierdzenie}
\begin{proof}
Oznaczmy $w=v-u$ i
$a(t,x)=\int_0^1f'(u+\theta(v-u))\,\mathrm d\theta$.
Wtedy $f(v)-f(u)=aw$ i $|a|\le L$ dla pewnej
stałej $L$ zależnej od zakresów rozwiązań.
Dla $z(t,x)=e^{Lt}w(t,x)$ wzór Duhamela ma postać
\[
 z(t)=S(t)w(0)+\int_0^tS(t-s)((a(s)+L)z(s))\,\mathrm ds.
\]
Tu $a+L\ge0$. Kolejne iteracje tego liniowego równania
Volterry, zaczynając od $z_0(t)=S(t)w(0)\ge0$, pozostają
nieujemne, gdyż półgrupa jest dodatnia. Iteracje zbiegają
jednostajnie na dostatecznie krótkim przedziale przez
kontrakcję ($T\|a+L\|_\infty<1$). Na pozostałych
częściach $[0,T]$ powtarzamy ten sam argument.
Otrzymujemy $w\ge0$.
\end{proof}

\begin{wniosek}[Globalność przy wzroście liniowym]
\label{cor:pde25-global}
Jeśli $f\in C^1(\R)$ i dla pewnych $A,B\ge0$ zachodzi
$|f(s)|\le A+B|s|$ dla każdego $s\in\R$, to każde
rozwiązanie łagodne na $\T^d$ istnieje dla wszystkich
$t\ge0$. Ponadto
\[
 \|u(t)\|_\infty\le
 e^{Bt}\|u_0\|_\infty+\frac AB(e^{Bt}-1)
\]
dla $B>0$; przy $B=0$ ostatni składnik należy zastąpić
przez $At$.
\end{wniosek}
\begin{proof}
Z kontrakcji półgrupy i wzoru Duhamela wynika
$\|u(t)\|_\infty\le\|u_0\|_\infty+
\int_0^t(A+B\|u(s)\|_\infty)\,\mathrm ds$.
Gronwall daje żądane ograniczenie na każdym skończonym
przedziale. Alternatywa~\eqref{eq:pde25-blowalt}
wyklucza skończony czas maksymalnego istnienia.
\end{proof}

\begin{przyklad}[Równanie logistyczne]
\label{ex:pde25-logistic}
Dla $u_t=\Delta u+u(1-u)$ i
$0\le u_0\in C(\T^d)$ otrzymujemy
$0\le u(t,x)\le M=\max\{1,\|u_0\|_\infty\}$.
Funkcja stała $0$ jest rozwiązaniem, a $M$ jest
supersolucją, bo $M(1-M)\le0$. Porównanie z funkcją
stałą $M$ można uzasadnić bez dodatkowej teorii
supersolucji: niech $m(t)$ rozwiązuje
$m'=m(1-m)$, $m(0)=M$; wtedy $m(t)\le M$.
Twierdzenie~\ref{thm:pde25-comparison} zastosowane
do rozwiązań $u$ i $m$ daje ograniczenie i globalność.
Choć $f(s)=s(1-s)$ rośnie kwadratowo w module,
nie występuje eksplozja dla takich danych nieujemnych.
\end{przyklad}


\begin{twierdzenie}[Oszacowanie przez równanie skalarne; warunek Osgooda]
\label{thm:pde25-osgood}
Niech $d\ge1$, $u_0\in C(\T^d)$ i $f\in C^1(\R)$.
Załóżmy, że funkcja $h:[0,\infty)\to(0,\infty)$ jest
niemalejąca, lokalnie lipschitzowska i spełnia
$|f(s)|\le h(|s|)$ dla wszystkich $s\in\R$.
Niech $M=\|u_0\|_\infty$, a $y:[0,\tau_M)\to[0,\infty)$
będzie maksymalnym rozwiązaniem $y'=h(y)$, $y(0)=M$.
Wówczas maksymalne rozwiązanie łagodne równania
$u_t=\Delta u+f(u)$ spełnia
$T_{\max}\ge\tau_M$ oraz
$\|u(t)\|_\infty\le y(t)$ dla $0\le t<\tau_M$.
Ponadto
\begin{equation}\label{eq:pde25-osgood}
 \tau_M=\int_M^\infty\frac{\mathrm ds}{h(s)},
\end{equation}
gdzie dopuszczamy wartość $+\infty$.
W szczególności rozbieżność całki w~\eqref{eq:pde25-osgood}
gwarantuje globalne istnienie rozwiązania $u$.
\end{twierdzenie}
\begin{proof}
Ponieważ $h>0$ i jest lokalnie lipschitzowska,
rozwiązanie $y$ jest ściśle rosnące. Rozdzielenie zmiennych
daje $t=\int_M^{y(t)}h(s)^{-1}\,\mathrm ds$; standardowe
kryterium przedłużania ODE dowodzi~\eqref{eq:pde25-osgood}.

Ustalmy $T<\min\{T_{\max},\tau_M\}$. Funkcje stałe
w przestrzeni $\overline y(t,x)=y(t)$ oraz
$\underline y(t,x)=-y(t)$ spełniają odpowiednio
$\overline y_t-\Delta\overline y-f(\overline y)\ge0$
i $\underline y_t-\Delta\underline y-f(\underline y)\le0$.
Dla uzasadnienia porównania w klasie łagodnej zauważmy,
że $w=\overline y-u$ spełnia wzór Duhamela z nieujemnym
danym początkowym $M-u_0$ oraz dodatkowym nieujemnym
źródłem $h(y)-f(y)$. Pozostała część źródła jest równa
$a(t,x)w$, gdzie
$a=\int_0^1f'(u+\theta(\overline y-u))\,\mathrm d\theta$
jest ograniczona. Mnożnik $e^{Lt}$, $L\ge\|a\|_\infty$,
przekształca to równanie w równanie Volterry z dodatnim
jądrem i nieujemnym współczynnikiem. Iteracja jak w dowodzie
twierdzenia~\ref{thm:pde25-comparison} daje $w\ge0$.
Analogicznie $u-\underline y\ge0$ (dodatkowe źródło wynosi
$f(-y)+h(y)\ge0$). Zatem $|u(t,x)|\le y(t)$.

Gdyby $T_{\max}<\tau_M$, funkcja $y$ byłaby ograniczona
na $[0,T_{\max}]$. Uzyskane oszacowanie przeczyłoby
alternatywie eksplozji z twierdzenia~\ref{thm:pde25-local}.
Wobec tego $T_{\max}\ge\tau_M$.
\end{proof}

\begin{przyklad}[Wzrost szybszy niż liniowy bez eksplozji z tego kryterium]
Dla $h(s)=(1+s)\log(e+s)$ całka
$\int_M^\infty h(s)^{-1}\,\mathrm ds$ jest rozbieżna.
Każda nieliniowość $f\in C^1(\R)$ spełniająca
$|f(s)|\le(1+|s|)\log(e+|s|)$ daje więc globalne
rozwiązania na torusie dla dowolnych danych ciągłych.
Pokazuje to, że warunek wzrostu liniowego z
wniosku~\ref{cor:pde25-global} nie jest konieczny.
\end{przyklad}


\section{Eksplozja w skończonym czasie: zasada średniej}
\begin{twierdzenie}[Wybuch dla dodatniego źródła na torusie]
\label{thm:pde25-blowup}
Niech $d\ge1$, $p>1$ oraz
$u_0\in C(\T^d)$, $u_0\ge0$, $u_0\not\equiv0$.
Niech $u$ będzie maksymalnym rozwiązaniem łagodnym
równania $u_t=\Delta u+(u_+)^p$, gdzie $u_+=\max\{u,0\}$.
Oznaczmy $m_0=|\T^d|^{-1}\int_{\T^d}u_0(x)
\,\mathrm d\lambda_d(x)>0$. Wówczas $u\ge0$ i
\begin{equation}\label{eq:pde25-mass-blowup}
 T_{\max}\le\frac{m_0^{1-p}}{p-1}<\infty.
\end{equation}
\end{twierdzenie}
\begin{proof}
Funkcja $s\mapsto(s_+)^p$ należy do $C^1(\R)$,
więc twierdzenie lokalne obowiązuje. Porównanie z
rozwiązaniem zerowym daje $u\ge0$. Z zachowania masy
przez półgrupę i~\eqref{eq:pde25-duhamel} wynika
\[
 m(t)=m_0+\int_0^t
 \frac1{|\T^d|}\int_{\T^d}u(s,x)^p\,
 \mathrm d\lambda_d(x)\,\mathrm ds.
\]
Funkcja $m$ jest klasy $C^1$ i, na mocy nierówności
Jensena dla wypukłej funkcji $s\mapsto s^p$ na
$[0,\infty)$, spełnia $m'(t)\ge m(t)^p$.
Ponieważ $m(t)>0$, możemy podzielić przez $m(t)^p$:
$\frac{\mathrm d}{\mathrm dt}m(t)^{1-p}
=(1-p)m(t)^{-p}m'(t)\le-(p-1)$.
Stąd $m(t)^{1-p}\le m_0^{1-p}-(p-1)t$.
Lewa strona jest dodatnia na całym przedziale istnienia,
co wymusza~\eqref{eq:pde25-mass-blowup}.
\end{proof}

\begin{uwaga}[Źródło eksplozji]
Wzrost normy wynika tu z nierówności różniczkowej dla
średniej, a nie z nieciągłości współczynników lub
braku regularności danych. Co więcej, nawet niewielkie,
niezerowe dane nieujemne prowadzą do eksplozji na torusie.
To twierdzenie nie przenosi się bez zmian na $\R^d$,
gdzie rozpraszanie masy może zapobiegać wybuchowi.
Wymagane są tam inne argumenty i krytyczne wykładniki.
\end{uwaga}

\begin{twierdzenie}[Wybuch dla dużych danych Dirichleta]
\label{thm:pde25-kaplan}
Niech $\Omega\subset\R^d$ będzie ograniczonym obszarem
klasy $C^2$, $p>1$ i $u$ nieujemnym rozwiązaniem klasycznym
na $[0,T_{\max})\times\overline\Omega$ zagadnienia
$u_t-\Delta u=u^p$, $u|_{\partial\Omega}=0$.
Niech $\lambda_1>0$ będzie pierwszą wartością własną
$-\Delta$ z warunkiem Dirichleta, a $\phi_1>0$
odpowiadającą jej funkcją własną,
$\int_\Omega\phi_1\,\mathrm d\lambda_d=1$.
Jeśli
$y_0=\int_\Omega u_0(x)\phi_1(x)\,\mathrm d\lambda_d(x)
>\lambda_1^{1/(p-1)}$,
to $T_{\max}<\infty$ i
$\|u(t)\|_\infty\to\infty$ dla $t\uparrow T_{\max}$,
o ile rozwiązanie należy do klasy, dla której obowiązuje
lokalne przedłużanie po oszacowaniu supremum.
\end{twierdzenie}
\begin{proof}
Dla $y(t)=\int_\Omega u(t,x)\phi_1(x)
\,\mathrm d\lambda_d(x)$ całkowanie przez części,
warunki Dirichleta i równanie na $\phi_1$ dają
$y'=-\lambda_1y+\int_\Omega u^p\phi_1\,
\mathrm d\lambda_d$.
Miara $\phi_1\,\mathrm d\lambda_d$ ma masę $1$,
więc Jensen zapewnia $y'\ge-\lambda_1y+y^p$.
Dla $y\ge y_0$ prawa strona jest dodatnia; stąd $y$
rośnie. Ustalmy $K>\max\{y_0,(2\lambda_1)^{1/(p-1)}\}$.
Dopóki $y\in[y_0,K]$, pochodna jest ograniczona
od dołu przez dodatnią stałą
$\min_{s\in[y_0,K]}(s^p-\lambda_1s)$,
a zatem $y$ osiąga $K$ w skończonym czasie.
Później $y'\ge\frac12 y^p$ i argument z poprzedniego
twierdzenia daje wybuch w skończonym czasie.
Ponieważ $0\le y(t)\le\|u(t)\|_\infty$,
nie może istnieć globalne rozwiązanie ograniczone.
\end{proof}

\section{Funkcjonał energii i równanie Allena--Cahna}
\begin{twierdzenie}[Dyssypacja energii]
\label{thm:pde25-energy}
Niech $V\in C^2(\R)$ i niech
$u\in C^1([0,T];C^1(\T^d))\cap C([0,T];C^2(\T^d))$ będzie rzeczywistym
rozwiązaniem $u_t=\Delta u-V'(u)$.
Definiujemy
\[
 E(u)=\int_{\T^d}\left(
 \frac12|\nabla u(x)|^2+V(u(x))\right)
 \,\mathrm d\lambda_d(x).
\]
Wtedy $t\mapsto E(u(t))$ jest różniczkowalna i
\begin{equation}\label{eq:pde25-energy}
 \frac{\mathrm d}{\mathrm dt}E(u(t))
 =-\int_{\T^d}|u_t(t,x)|^2\,\mathrm d\lambda_d(x)\le0.
\end{equation}
\end{twierdzenie}
\begin{proof}
Przy przyjętej regularności wolno różniczkować gradient w czasie i pod znakiem całki. Periodyczność i
całkowanie przez części dają
$\frac{\mathrm d}{\mathrm dt}E(u(t))
=\int_{\T^d}(-\Delta u+V'(u))u_t
\,\mathrm d\lambda_d$.
Z równania mamy $-\Delta u+V'(u)=-u_t$,
więc otrzymujemy~\eqref{eq:pde25-energy}.
\end{proof}

\begin{przyklad}[Równanie Allena--Cahna]
Dla $V(s)=\frac14(s^2-1)^2$ równanie
$u_t=\Delta u+u-u^3$ ma stabilne, stałe stany $u\equiv\pm1$
i niestabilny stan $u\equiv0$ w sensie liniaryzacji
jednorodnych perturbacji. Porównanie z rozwiązaniami
stałymi $\pm M$, gdzie $M\ge\max\{1,\|u_0\|_\infty\}$,
pokazuje globalne istnienie przy danych ciągłych.
Energia maleje zgodnie z~\eqref{eq:pde25-energy};
sama monotoniczność energii nie oznacza jednak,
że rozwiązanie musi zbiegać do konkretnego stanu
stałego bez dalszych argumentów zwartości.
\end{przyklad}

\section{Półliniowe równanie falowe w jednym wymiarze}
Dotychczas wykorzystywaliśmy wygładzającą półgrupę ciepła.
Dla równania falowego nieliniowość nie jest wygładzana:
wymagana jest przestrzeń, w której operator Nemyckiego
jest lokalnie Lipschitzowski. Pracujemy na $\T^1$.
Niech $A=1-\partial_x^2$ w przestrzeni Hilberta
$L^2(\T^1)$, a
$Y=H^1(\T^1)\times L^2(\T^1)$ z normą
$\|(u,v)\|_Y^2=\|A^{1/2}u\|_2^2+\|v\|_2^2$.

\begin{twierdzenie}[Lokalne istnienie dla równania falowego]
\label{thm:pde25-wave}
Niech $g\in C^1(\R)$, $(u_0,u_1)\in Y$
i $T>0$ będzie wystarczająco mały, zależnie od normy
danych i $g$. Wtedy zagadnienie
\[
 u_{tt}-u_{xx}+u=g(u),\quad
 u(0)=u_0,\quad u_t(0)=u_1
\]
ma dokładnie jedno rozwiązanie łagodne
$(u,u_t)\in C([0,T];Y)$. Rozwiązanie zależy
lokalnie Lipschitzowsko od danych w normie $Y$.
\end{twierdzenie}
\begin{proof}
\textbf{Krok 1: grupa liniowa.} W bazie Fouriera operator
$A$ ma wartości własne $1+k^2$, $k\in\mathbb Z$.
Definiujemy $C(t)=\cos(tA^{1/2})$ i
$D(t)=A^{-1/2}\sin(tA^{1/2})$ przez mnożniki Fouriera.
Równanie liniowe ma postać
$u(t)=C(t)u_0+D(t)u_1$,
$u_t(t)=-A^{1/2}\sin(tA^{1/2})u_0+C(t)u_1$.
Z tożsamości $\cos^2+\sin^2=1$ wynika, że operator
przejścia $U(t)$ jest izometrią na $Y$.

\textbf{Krok 2: lokalna lipschitzowskość.} Wymiar jeden
zapewnia zanurzenie $H^1(\T^1)\hookrightarrow C(\T^1)$.
Jeśli $\|v\|_{H^1},\|w\|_{H^1}\le R$, to
$\|g(v)-g(w)\|_2\le L_{CR}\|v-w\|_2
\le L_{CR}\|v-w\|_{H^1}$, gdzie
$L_{CR}=\max_{|s|\le CR}|g'(s)|$.

\textbf{Krok 3: punkt stały.} Dla $z=(u,v)$ definiujemy
\[
 \Psi(z)(t)=U(t)(u_0,u_1)
 +\int_0^tU(t-s)(0,g(u(s)))\,\mathrm ds.
\]
Na kuli $\|z\|_{C([0,T];Y)}\le R$
operator $\Psi$ jest odwzorowaniem w siebie, gdy
$T$ jest dość małe, ponieważ
$\|U(t-s)(0,h)\|_Y=\|h\|_2$.
Ponadto $\|\Psi(z)-\Psi(w)\|_{C Y}
\le TL_{CR}\|z-w\|_{C Y}$.
Po wybraniu $TL_{CR}<1$ zasada Banacha daje punkt stały.
Wzór Duhamela identyfikuje drugą składową z pochodną
czasową pierwszej w $L^2$. Różnicę dwóch rozwiązań
szacujemy jak w twierdzeniu~\ref{thm:pde25-stability}.
\end{proof}

\begin{stwierdzenie}[Energia półliniowego równania falowego]
\label{prop:pde25-waveenergy}
Niech $g\in C^1(\R)$, $G(s)=\int_0^s g(r)\,\mathrm dr$,
a $u$ będzie wystarczająco gładkim rzeczywistym
rozwiązaniem na $\T^1$ równania z poprzedniego twierdzenia.
Wtedy
\[
 \mathcal E(t)=\int_{\T^1}\left(
 \tfrac12u_t^2+\tfrac12u_x^2+\tfrac12u^2-G(u)\right)
 \,\mathrm d\lambda_1(x)
\]
jest stała w czasie.
\end{stwierdzenie}
\begin{proof}
Po pomnożeniu równania przez $u_t$ i scałkowaniu
po torusie składnik $-u_{xx}u_t$ po całkowaniu
przez części jest pochodną $\tfrac12\|u_x\|_2^2$;
pozostałe składniki są odpowiednio pochodnymi
$\tfrac12\|u_t\|_2^2$, $\tfrac12\|u\|_2^2$
i $-\int_{\T^1}G(u(t,x))\,\mathrm d\lambda_1(x)$. Suma pochodnych wynosi zero.
\end{proof}

\begin{uwaga}[Energia nie zawsze kontroluje normę]
Jeżeli $G$ rośnie szybciej niż kwadratowo i z dodatnim
znakiem, składnik $-G(u)$ może sprawić, że energia
nie jest ograniczona od dołu. Jej zachowanie samo w sobie
nie daje globalnego istnienia. Dla nieliniowości
$g(u)=-u^3$ energia ma natomiast dodatni człon
$\int_{\T^1}u(t,x)^4/4\,\mathrm d\lambda_1(x)$ i kontroluje normę energetyczną.
\end{uwaga}

\section{Eliptyczne równanie półliniowe i tożsamość Pohozaeva}
Rozważmy $-\Delta u=u^p$ w ograniczonym obszarze
$\Omega\subset\R^d$, $d\ge3$, z zerowym warunkiem
Dirichleta i $u>0$ wewnątrz. W przeciwieństwie do teorii
liniowej rozwiązanie może nie istnieć z przyczyn
wynikających z geometrii obszaru i wykładnika $p$.

\begin{lemat}[Tożsamość Pohozaeva]
\label{lem:pde25-pohozaev}
Niech $d\ge2$, $\Omega\subset\R^d$ będzie ograniczonym
obszarem klasy $C^2$, $F\in C^1(\R)$, $F(0)=0$,
a $u\in C^2(\overline\Omega)$ spełnia
$-\Delta u=F'(u)$ w $\Omega$ i $u=0$ na $\partial\Omega$.
Jeśli $\nu$ jest jednostkową normalną zewnętrzną, to
\begin{equation}\label{eq:pde25-pohozaev}
 \frac{d-2}{2}\int_\Omega|\nabla u|^2\,\mathrm d\lambda_d
 -d\int_\Omega F(u)\,\mathrm d\lambda_d
 =-\frac12\int_{\partial\Omega}(x\cdot\nu)
 (\partial_\nu u)^2\,\mathrm d\sigma(x),
\end{equation}
gdzie $\sigma$ jest miarą powierzchniową na brzegu.
\end{lemat}
\begin{proof}
Mnożymy równanie przez $x\cdot\nabla u$ i całkujemy.
Z jednej strony, ponieważ $\nabla F(u)=F'(u)\nabla u$,
$\int_\Omega F'(u)(x\cdot\nabla u)\,\mathrm d\lambda_d
=-d\int_\Omega F(u)\,\mathrm d\lambda_d$;
wyraz brzegowy $\int_{\partial\Omega}(x\cdot\nu)F(u)
\,\mathrm d\sigma$ znika z uwagi na $F(0)=0$.
Z drugiej strony całkowanie przez części daje
\[
 \begin{aligned}
 \int_\Omega(-\Delta u)(x\cdot\nabla u)\,\mathrm d\lambda_d
 &= -\int_{\partial\Omega}\partial_\nu u
 (x\cdot\nabla u)\,\mathrm d\sigma
 +\int_\Omega\nabla u\cdot\nabla(x\cdot\nabla u)
 \,\mathrm d\lambda_d\\
 &= -\int_{\partial\Omega}(x\cdot\nu)(\partial_\nu u)^2
 \,\mathrm d\sigma
 +\int_\Omega\bigl(|\nabla u|^2
 +\tfrac12 x\cdot\nabla|\nabla u|^2\bigr)
 \,\mathrm d\lambda_d.
 \end{aligned}
\]
Ponieważ $u$ ma stały ślad na brzegu, jego pochodne
styczne znikają tam, zatem $\nabla u=(\partial_\nu u)\nu$.
Twierdzenie o dywergencji zastosowane do
$x|\nabla u|^2$ zamienia ostatni składnik objętościowy
na $-\frac d2\int_\Omega|\nabla u|^2\,
\mathrm d\lambda_d+\frac12\int_{\partial\Omega}
(x\cdot\nu)(\partial_\nu u)^2\,\mathrm d\sigma$.
Po porównaniu z pierwszą stroną otrzymujemy
\eqref{eq:pde25-pohozaev}.
\end{proof}

\begin{twierdzenie}[Nieistnienie dodatnich rozwiązań przy wykładniku krytycznym]
\label{thm:pde25-pohozaev-nonexist}
Niech $d\ge3$, $p\ge(d+2)/(d-2)$ oraz
$\Omega\subset\R^d$ będzie ograniczonym obszarem klasy
$C^2$, \emph{ściśle gwiaździstym względem zera}:
$x\cdot\nu(x)>0$ dla każdego $x\in\partial\Omega$.
Nie istnieje funkcja $u\in C^2(\overline\Omega)$
taka, że $u>0$ w $\Omega$, $u=0$ na $\partial\Omega$
i $-\Delta u=u^p$ w $\Omega$.
\end{twierdzenie}
\begin{proof}
Mnożenie równania przez $u$ i całkowanie przez części
daje $\int_\Omega|\nabla u|^2\,\mathrm d\lambda_d
=\int_\Omega u^{p+1}\,\mathrm d\lambda_d$.
Stosujemy Lemat~\ref{lem:pde25-pohozaev}
do $F(s)=s^{p+1}/(p+1)$ dla $s\ge0$,
z dowolnym przedłużeniem $C^1$ na $\R$ i $F(0)=0$.
Otrzymujemy
\[
 \left(\frac{d-2}{2}-\frac{d}{p+1}\right)
 \int_\Omega|\nabla u|^2\,\mathrm d\lambda_d
 =-\frac12\int_{\partial\Omega}(x\cdot\nu)
 (\partial_\nu u)^2\,\mathrm d\sigma.
\]
Dla $p>(d+2)/(d-2)$ lewy współczynnik jest dodatni,
a prawa strona niedodatnia, więc $u\equiv0$,
sprzeczność. W przypadku krytycznym współczynnik
znika; ścisła gwiaździstość wymusza
$\partial_\nu u\equiv0$ na brzegu. Tymczasem $-\Delta u=u^p>0$
i $u>0$ wewnątrz, więc lemat Hopfa dla obszaru
klasy $C^2$ daje $\partial_\nu u<0$ na brzegu.
To sprzeczność kończąca dowód.
\end{proof}

\begin{uwaga}[Wykładnik krytyczny i ograniczenia twierdzenia]
Liczba $(d+2)/(d-2)$ odpowiada krytycznemu wykładnikowi
zanurzenia Sobolewa $H_0^1\hookrightarrow L^{2d/(d-2)}$.
Nie twierdzimy, że nie istnieją rozwiązania zmieniające znak
w przypadku krytycznym, ani że zachodzi ten sam wynik
na dowolnych obszarach: użyliśmy dodatniości oraz
ściśle gwiaździstej geometrii w twierdzeniu Pohozaeva.
\end{uwaga}

\section{Porównanie trzech mechanizmów i dalszy plan}
Dla równania ciepła z nieliniowością odpowiedni dobór
przestrzeni i kontrakcja dają lokalną jednoznaczność,
ale wzrost reakcji może wymusić eksplozję. Dla równania
falowego nieliniowość działa w układzie bez wygładzania,
co ogranicza bezpośredni argument do przestrzeni
zamkniętej na składanie z $g$. Dla zagadnienia
eliptycznego nawet bardzo gładkie dane i brzeg nie
zapewniają istnienia z powodu tożsamości całkowej.
W następnym rozdziale wykorzystamy operatory monotoniczne
i $p$-laplasjan, które prowadzą do innego mechanizmu istnienia
rozwiązań nieliniowych, bez kontrakcji Picarda.


% ROZDZIAL 26: monotonicznosc i plaplasjan
\chapter[Operatory monotoniczne i p-laplasjan]{Operatory monotoniczne, twierdzenie Minty'ego--Browdera i nieliniowe równania eliptyczne}
\label{chap:pde26}

\section{Dlaczego metoda Laxa--Milgrama przestaje wystarczać?}
W rozdziałach~\ref{chap:pde25} oraz wcześniejszych zagadnieniach
eliptycznych występowały albo operatory liniowe, albo nieliniowości
niższego rzędu. Teraz najwyższa pochodna występuje nieliniowo:
\begin{equation}\label{eq:pde26-main}
 -\diver\bigl(|\nabla u|^{p-2}\nabla u\bigr)=f,\qquad p>1.
\end{equation}
Dla $p\ne2$ nie istnieje forma dwuliniowa pozwalająca zastosować
wprost twierdzenie Laxa--Milgrama. Zastępujemy dwuliniowość
\emph{monotonicznością}, która w teorii nieliniowej odgrywa rolę
dodatniości formy kwadratowej.

W całym rozdziale $\Omega\subset\R^d$ oznacza ograniczony obszar
Lipschitzowski, o ile przy konkretnym wyniku nie zaznaczono inaczej.
Ustalamy $1<p<\infty$, $p'=p/(p-1)$, oraz
$V=W_0^{1,p}(\Omega)$, wyposażoną w normę
$\|u\|_V=\|\nabla u\|_{L^p(\Omega)}$.
Nierówność Poincarégo zapewnia jej równoważność ze zwykłą normą
Sobolewa. Przestrzeń $V$ jest separowalna i refleksywna, a jej
przestrzeń dualną oznaczamy przez $V^*=W^{-1,p'}(\Omega)$.
Przez $\langle f,v\rangle$ oznaczamy wartość funkcjonału
$f\in V^*$ na $v\in V$; w przypadku $f\in L^{p'}(\Omega)$
jest to całka $\int_\Omega f(x)v(x)\,\mathrm d\lambda_d(x)$.
Dla funkcji $v\in V$ oznaczamy $v_+=\max\{v,0\}$ i
$v_-=\max\{-v,0\}$; obie należą do $V$, a
$\nabla v_+=\chi_{\{v>0\}}\nabla v$ prawie wszędzie.

\begin{definicja}[Monotoniczność i hemiciągłość]\label{def:pde26-monotone}
Niech $X$ będzie rzeczywistą przestrzenią Banacha, a
$A:X\to X^*$ odwzorowaniem określonym na całym $X$.
Mówimy, że $A$ jest
\emph{monotoniczne}, jeżeli
$\langle Au-Av,u-v\rangle\ge0$ dla wszystkich $u,v\in X$,
zaś \emph{ściśle monotoniczne}, jeżeli nierówność jest ostra
przy $u\ne v$. Operator jest \emph{hemiciągły}, jeśli dla
wszystkich $u,v,w\in X$ odwzorowanie
$t\mapsto\langle A(u+tv),w\rangle$ jest ciągłe na $\R$.
Operator $A$ jest \emph{ograniczony na zbiorach ograniczonych},
jeśli dla każdego $R>0$ zachodzi
$\sup_{\|u\|_X\le R}\|Au\|_{X^*}<\infty$.
Nazywamy go \emph{koercywnym}, jeśli
$\langle Au,u\rangle/\|u\|_X\to+\infty$
przy $\|u\|_X\to\infty$.
\end{definicja}

\begin{uwaga}[Warianty słowa koercywność]
Spotyka się również warunek silniejszy
$\langle Au,u\rangle\ge c\|u\|_X^p-C$ dla pewnego $p>1$.
Jest on wystarczający do naszej definicji, ale nie jest z nią
równoważny. Koercywność nie oznacza monotoniczności;
obie własności są niezależnymi hipotezami twierdzenia
Minty'ego--Browdera.
\end{uwaga}

\section{Nierówności monotoniczności dla wektorów}
\begin{lemat}[Monotoniczność strumienia potęgowego]\label{lem:pde26-vector}
Dla $p>1$ odwzorowanie $a_p:\R^d\to\R^d$,
$a_p(\xi)=|\xi|^{p-2}\xi$ dla $\xi\ne0$ i $a_p(0)=0$,
jest ciągłe i ściśle monotoniczne. Co więcej, istnieją
stałe $c_p,C_p>0$ zależne tylko od $p$ takie, że dla
wszystkich $\xi,\eta\in\R^d$ zachodzą:
\begin{align}
 (a_p(\xi)-a_p(\eta))\cdot(\xi-\eta)
 &\ge \frac{2^{2-p}}{p-1}|\xi-\eta|^p &&(p\ge2),\label{eq:pde26-strong}\\
\intertext{a dla $1<p<2$ i $\xi\ne\eta$}
 (a_p(\xi)-a_p(\eta))\cdot(\xi-\eta)
 &\ge(p-1)(|\xi|+|\eta|)^{p-2}|\xi-\eta|^2.\label{eq:pde26-subquad}
\end{align}
W obu przypadkach prawa strona jest dodatnia przy $\xi\ne\eta$.
\end{lemat}
\begin{proof}
Funkcja $\Phi(\xi)=|\xi|^p/p$ jest ściśle wypukła,
a jej gradientem jest $a_p$. Wypukłość daje monotoniczność;
wyjaśniamy dodatkowo jej ilościowy charakter.
Dla $\zeta\ne0$ macierz pochodnej $a_p$ wynosi
$Da_p(\zeta)=|\zeta|^{p-2}I+(p-2)|\zeta|^{p-4}\zeta\otimes\zeta$.
Jej najmniejsza wartość własna jest równa
$\min\{1,p-1\}|\zeta|^{p-2}$. Po scałkowaniu po
odcinku $\eta+t(\xi-\eta)$ otrzymujemy (także dla odcinka
przechodzącego przez zero, przez przybliżenie)
\[
 (a_p(\xi)-a_p(\eta))\cdot(\xi-\eta)
 \ge \min\{1,p-1\}|\xi-\eta|^2
 \int_0^1|\eta+t(\xi-\eta)|^{p-2}\,\mathrm dt.
\]
Jeśli $1<p<2$, na całym odcinku
$|\eta+t(\xi-\eta)|\le|\xi|+|\eta|$, a wykładnik
$p-2$ jest ujemny. Otrzymujemy~\eqref{eq:pde26-subquad}.
Dla $p\ge2$ korzystamy z elementarnej nierówności
\[
 \int_0^1|\eta+t(\xi-\eta)|^{p-2}\,\mathrm dt
 \ge \frac{2^{2-p}}{p-1}|\xi-\eta|^{p-2}.
\]
Dla przejrzystości można ją uzasadnić po rozłożeniu
$\eta$ na składową równoległą i prostopadłą do
$h=\xi-\eta$: usunięcie składowej prostopadłej
nie zwiększa lewej strony, a funkcja
$b\mapsto\int_0^1|b+t|^{p-2}\,\mathrm dt$
przyjmuje minimum dla $b=-1/2$ z uwagi na symetrię
i wypukłość funkcji pierwotnej na półprostych;
minimum wynosi $2^{2-p}/(p-1)$.
To prowadzi do stałej $2^{2-p}/(p-1)$ w
\eqref{eq:pde26-strong}. Nie jest to najlepsza możliwa
stała, ale wystarcza do dalszych zastosowań.
Ścisłość wynika z dodatniości uzyskanych prawych stron.
\end{proof}

\begin{przyklad}[Dlaczego $p=2$ jest szczególny]
Dla $p=2$ mamy $a_2(\xi)=\xi$ i
$(a_2(\xi)-a_2(\eta))\cdot(\xi-\eta)=|\xi-\eta|^2$.
Dla $p>2$ różniczka $Da_p(0)$ jest zerowa, zatem
jednostajna eliptyczność linearyzacji zanika przy
$\nabla u=0$. Dla $1<p<2$ pochodna $a_p$ ma
osobliwość przy zerze. Żadne z tych zjawisk nie
niszczy jednak monotoniczności nieliniowego operatora.
\end{przyklad}

\section{Twierdzenie Minty'ego--Browdera z dowodem}
\begin{lemat}[Skończeniewymiarowa zasada istnienia]
\label{lem:pde26-brouwer}
Jeśli $F:\R^n\to\R^n$ jest ciągłe, a dla pewnego
$R>0$ zachodzi $F(x)\cdot x>0$ na sferze $|x|=R$,
to istnieje $x\in\R^n$ takie, że $|x|<R$ i $F(x)=0$.
\end{lemat}
\begin{proof}
Załóżmy przeciwnie, że $F$ nie zeruje się na
$\overline B(0,R)$. Definiujemy na tej kuli
odwzorowanie ciągłe
\[
 T(x)=-\frac{RF(x)}{|F(x)|}.
\]
Jego wartości należą do sfery o promieniu $R$,
a więc do rozważanej kuli. Twierdzenie Brouwera
o punkcie stałym daje $x=-RF(x)/|F(x)|$.
Wtedy $|x|=R$ oraz $F(x)\cdot x=-R|F(x)|<0$,
co jest sprzeczne z hipotezą. Wniosek jest prawdziwy.
\end{proof}

\begin{twierdzenie}[Twierdzenie Minty'ego--Browdera; wersja z pełną dziedziną]
\label{thm:pde26-minty}
Niech $X$ będzie rzeczywistą, separowalną, refleksywną
przestrzenią Banacha. Załóżmy, że $A:X\to X^*$ jest
monotoniczny, hemiciągły, ograniczony na zbiorach ograniczonych
i koercywny w sensie definicji~\ref{def:pde26-monotone}.
Wówczas dla każdego $f\in X^*$ istnieje $u\in X$ takie, że
$Au=f$. Jeśli ponadto $A$ jest ściśle monotoniczny,
to rozwiązanie jest jedyne.
\end{twierdzenie}
\begin{proof}
\textbf{Krok 1: przybliżenia Galerkina.}
Wybieramy liniowo niezależny ciąg $(e_j)$,
którego powłoka liniowa jest gęsta w $X$, i kładziemy
$X_n=\operatorname{span}\{e_1,\ldots,e_n\}$.
Szukamy $u_n\in X_n$ spełniającego
\begin{equation}\label{eq:pde26-galerkin}
 \langle Au_n,v\rangle=\langle f,v\rangle
 \quad\text{dla każdego }v\in X_n.
\end{equation}
W bazie $(e_1,\ldots,e_n)$ tworzymy układ
$F_j(c)=\langle A(\sum_kc_ke_k)-f,e_j\rangle=0$.
Hemiciągłość wraz z monotonicznością i lokalną
ograniczonością implikuje ciągłość $A$ od normy $X$
do słabej topologii $X^*$: istotnie, jeśli $z_m\to z$,
to każda słaba granica $Az_m$ (istniejąca na podciągu)
spełnia, po przejściu do granicy w nierówności
$\langle Az_m-Av,z_m-v\rangle\ge0$,
relację $\langle g-Av,z-v\rangle\ge0$ dla wszystkich $v$;
podstawiając $v=z-tw$ i przechodząc z $t\downarrow0$
przez hemiciągłość, otrzymujemy $g=Az$.
W przestrzeni o skończonym wymiarze słaba ciągłość
współrzędnych wystarcza do ciągłości $F$.

Koercywność i nierówność
$\langle Au-f,u\rangle\ge\langle Au,u\rangle
-\|f\|_{X^*}\|u\|_X$ pokazują, że
$\langle Au-f,u\rangle>0$ na sferze
$\|u\|_X=R$ dla dostatecznie dużego $R$,
niezależnego od $n$.
Dla każdego $n$ równoważność norm na $X_n$ pozwala
wybrać promień euklidesowy $r_n$ na tyle duży, że
$\|\sum_{j=1}^n c_je_j\|_X\ge R$ dla $|c|=r_n$.
Wówczas $F(c)\cdot c=\langle Au-f,u\rangle>0$.
Lemat~\ref{lem:pde26-brouwer} daje rozwiązanie
\eqref{eq:pde26-galerkin}.

\textbf{Krok 2: oszacowania i granice słabe.}
Testując~\eqref{eq:pde26-galerkin} przez $u_n$, mamy
$\langle Au_n,u_n\rangle=\langle f,u_n\rangle$.
Koercywność wymusza jednostajne ograniczenie
$\|u_n\|_X\le R$. Lokalna ograniczoność $A$ daje
$\|Au_n\|_{X^*}\le M_R$. Refleksywność $X$ oraz $X^*$
pozwala wybrać podciąg, dla którego
$u_n\rightharpoonup u$ w $X$ oraz
$Au_n\rightharpoonup g$ w $X^*$.
Dla dowolnego $j$ i $n\ge j$ mamy
$\langle Au_n,e_j\rangle=\langle f,e_j\rangle$.
Po przejściu do granicy i wykorzystaniu gęstości
otrzymujemy $g=f$.

\textbf{Krok 3: kluczowa nierówność graniczna.}
Wybieramy $v_n\in X_n$ takie, że $v_n\to u$ silnie
w $X$. Istnienie takiego ciągu wynika z gęstości
rosnących podprzestrzeni $X_n$: dla każdego $k$ wybieramy
$z_k\in X_{n_k}$ z $\|z_k-u\|_X<1/k$ i
$n_k\uparrow\infty$; kładziemy $v_n=z_k$ dla
$n_k\le n<n_{k+1}$. Dzięki równaniu Galerkina
$\langle Au_n,u_n-v_n\rangle=\langle f,u_n-v_n\rangle$.
Ponieważ $u_n\rightharpoonup u$, $v_n\to u$ i $Au_n$
pozostaje ograniczone, otrzymujemy
\begin{equation}\label{eq:pde26-limsup}
 \lim_{n\to\infty}\langle Au_n,u_n-u\rangle=0.
\end{equation}
Dla ustalonego $v\in X$ monotoniczność daje
$\langle Au_n-Av,u_n-v\rangle\ge0$.
Rozkładamy pierwszy człon jako
$\langle Au_n,u_n-u\rangle+\langle Au_n,u-v\rangle$
i korzystamy z~\eqref{eq:pde26-limsup}. W granicy
otrzymujemy $\langle f-Av,u-v\rangle\ge0$.

\textbf{Krok 4: identyfikacja operatora.}
Dla dowolnego $w\in X$ i $t>0$ podstawiamy $v=u-tw$.
Wynika stąd $\langle f-A(u-tw),w\rangle\ge0$.
Przy $t\downarrow0$ hemiciągłość daje
$\langle f-Au,w\rangle\ge0$. Stosując nierówność także
do $-w$, otrzymujemy $Au=f$.
Jeśli $Au=Av=f$, ścisła monotoniczność wymusza $u=v$.
\end{proof}

\begin{uwaga}[Dokładny zakres twierdzenia]
Nie wymagamy liniowości $A$ ani istnienia potencjału
wariacyjnego. Warunek ograniczoności na zbiorach
ograniczonych jest tu przyjęty jawnie, ponieważ używamy
go w dowodzie do wydzielenia słabo zbieżnego podciągu $Au_n$.
Twierdzenie nie obejmuje automatycznie operatorów
określonych jedynie na właściwej części $X$ ani równań,
w których zawodzi monotoniczność.
\end{uwaga}

\section{p-laplasjan jako operator na przestrzeni Sobolewa}
\begin{definicja}[Rozwiązanie słabe $p$-laplasjanu]\label{def:pde26-weak}
Niech $f\in V^*$. Funkcję $u\in V=W_0^{1,p}(\Omega)$
nazywamy słabym rozwiązaniem zagadnienia
\eqref{eq:pde26-main} z zerowym warunkiem Dirichleta,
jeśli dla każdego $v\in V$ zachodzi
\begin{equation}\label{eq:pde26-weak}
 \int_\Omega|\nabla u|^{p-2}\nabla u\cdot\nabla v
 \,\mathrm d\lambda_d(x)=\langle f,v\rangle.
\end{equation}
Operator $A_p:V\to V^*$ definiujemy przez lewą stronę
tego wzoru; zapis $-\Delta_pu=f$ oznacza równość
$A_pu=f$ w $V^*$.
\end{definicja}

\begin{twierdzenie}[Istnienie i jednoznaczność]\label{thm:pde26-plap}
Niech $1<p<\infty$ i niech $\Omega\subset\R^d$ będzie
ograniczonym obszarem Lipschitzowskim. Dla każdego
$f\in W^{-1,p'}(\Omega)$ istnieje dokładnie jedno
$u\in W_0^{1,p}(\Omega)$ spełniające
\eqref{eq:pde26-weak}. Ponadto
\begin{equation}\label{eq:pde26-apriori}
 \|\nabla u\|_{L^p(\Omega)}^{p-1}
 \le\|f\|_{W^{-1,p'}(\Omega)}.
\end{equation}
\end{twierdzenie}
\begin{proof}
Sprawdzamy każdą hipotezę twierdzenia
\ref{thm:pde26-minty} dla $A_p$.
Nierówność Höldera daje
$|\langle A_pu,v\rangle|\le
\|\nabla u\|_p^{p-1}\|\nabla v\|_p$, więc
$\|A_pu\|_{V^*}\le\|u\|_V^{p-1}$; operator jest
określony na całym $V$ i ograniczony na zbiorach
ograniczonych. Ciągłość strumienia $a_p$ oraz
oszacowanie
$|a_p(\nabla u+t\nabla w)|\le
C_p(|\nabla u|+|\nabla w|)^{p-1}$ dla $|t|\le1$
pozwalają zastosować zbieżność zdominowaną do
$\langle A_p(u+tw),v\rangle$, co daje hemiciągłość.
Monotoniczność i jej ścisła wersja wynikają
z lematu~\ref{lem:pde26-vector} po scałkowaniu.
Koercywność jest dokładna:
$\langle A_pu,u\rangle=\|\nabla u\|_p^p$,
a więc iloraz przez $\|u\|_V$ dąży do nieskończoności.
Twierdzenie Minty'ego--Browdera zapewnia jedyne rozwiązanie.
Testując jego równanie przez $u$, otrzymujemy
$\|u\|_V^p=\langle f,u\rangle\le
\|f\|_{V^*}\|u\|_V$. Dzielenie przez $\|u\|_V$
(dla $u=0$ teza jest oczywista) kończy dowód.
\end{proof}

\section{Metoda bezpośrednia rachunku wariacyjnego}
\begin{twierdzenie}[Charakteryzacja minimizacyjna]\label{thm:pde26-minimizer}
Niech $p>1$, $\Omega$ będzie ograniczonym obszarem
Lipschitzowskim, a $f\in V^*$. Funkcjonał
\[
 J_f(v)=\frac1p\int_\Omega|\nabla v|^p
 \,\mathrm d\lambda_d(x)-\langle f,v\rangle,\qquad v\in V,
\]
ma dokładnie jeden punkt minimum. Jest nim jedyne
rozwiązanie zagadnienia~\eqref{eq:pde26-weak}.
\end{twierdzenie}
\begin{proof}
Z nierówności Younga otrzymujemy
$J_f(v)\ge\|v\|_V^p/p-\|f\|_{V^*}\|v\|_V$,
więc $J_f$ jest koercywny i ograniczony od dołu.
Ciąg minimalizujący jest ograniczony w refleksywnej
przestrzeni $V$; wybrany podciąg słabo zbiega do $u$.
Norma w przestrzeni $L^p$ jest słabo półciągła z dołu,
a funkcjonał liniowy $v\mapsto\langle f,v\rangle$
jest słabo ciągły. Stąd $J_f(u)\le\liminf J_f(v_n)$,
czyli $u$ jest minimizatorem. Ścisła wypukłość
$z\mapsto|z|^p$ daje ścisłą wypukłość $J_f$ na $V$:
dwa elementy o równych gradientach różnią się o stałą,
a w $W_0^{1,p}$ stała ta musi być zerowa.
Różniczkując $J_f(u+tv)$ w $t=0$ (zdominowanie jak
w poprzednim dowodzie) dostajemy dokładnie
\eqref{eq:pde26-weak}. Odwrotnie, dla rozwiązania
$A_pu=f$ wypukłość $J_f$ i zerowanie jego pochodnej
zapewniają minimalność.
\end{proof}

\begin{uwaga}[Dlaczego przedstawiliśmy dwa dowody?]
Metoda minimizacji wykorzystuje szczególną strukturę
potencjałową $A_p=J_f' + f$. Twierdzenie Minty'ego--Browdera
działa również dla monotonicznych operatorów bez potencjału.
Przykładem jest strumień $a(x,\xi)$ monotoniczny w $\xi$,
który nie musi być gradientem funkcji wypukłej względem $\xi$.
\end{uwaga}


\begin{twierdzenie}[Strumienie typu Leraya--Lionsa]
\label{thm:pde26-leraylions}
Niech $\Omega\subset\R^d$ będzie ograniczonym obszarem
Lipschitzowskim, $1<p<\infty$ i $V=W_0^{1,p}(\Omega)$.
Niech $a:\Omega\times\R^d\to\R^d$ będzie funkcją
Carathéodory'ego: $x\mapsto a(x,\xi)$ jest mierzalna
przy ustalonym $\xi$, a $\xi\mapsto a(x,\xi)$ jest
ciągła dla prawie każdego $x$. Załóżmy, że istnieją
$C,\alpha>0$, $h\in L^{p'}(\Omega)$, $h\ge0$,
i $\beta\in L^1(\Omega)$, $\beta\ge0$, takie że
prawie wszędzie w $x$ i dla wszystkich $\xi,\eta\in\R^d$
zachodzą
\begin{align}
 |a(x,\xi)|&\le h(x)+C|\xi|^{p-1},\label{eq:pde26-ll-growth}\\
 a(x,\xi)\cdot\xi&\ge\alpha|\xi|^p-\beta(x),\label{eq:pde26-ll-coerc}\\
 (a(x,\xi)-a(x,\eta))\cdot(\xi-\eta)&\ge0.\label{eq:pde26-ll-mono}
\end{align}
Wówczas dla każdego $f\in V^*$ istnieje $u\in V$
spełniające
\begin{equation}\label{eq:pde26-ll-weak}
 \int_\Omega a(x,\nabla u(x))\cdot\nabla v(x)
 \,\mathrm d\lambda_d(x)=\langle f,v\rangle
 \qquad(v\in V).
\end{equation}
Jeżeli nierówność w~\eqref{eq:pde26-ll-mono} jest ostra
przy $\xi\ne\eta$ dla prawie każdego $x$, rozwiązanie
jest jednoznaczne.
\end{twierdzenie}
\begin{proof}
Definiujemy $A:V\to V^*$ lewą stroną
\eqref{eq:pde26-ll-weak}. Warunek wzrostu i nierówność
Höldera zapewniają
$\|Au\|_{V^*}\le\|h\|_{p'}+C\|u\|_V^{p-1}$;
odwzorowanie jest więc dobrze określone i ograniczone
na zbiorach ograniczonych. Monotoniczność wynika
z całkowania~\eqref{eq:pde26-ll-mono}.
Dla $u,v,w\in V$ oraz $|t|\le1$ zachodzi prawie wszędzie
\[
 |a(x,\nabla u+t\nabla v)\cdot\nabla w|
 \le\bigl(h(x)+C_pC(|\nabla u|+|\nabla v|)^{p-1}\bigr)
 |\nabla w|.
\]
Prawa strona jest całkowalna przez nierówność Höldera.
Ciągłość w zmiennej $\xi$ i twierdzenie o zbieżności
zdominowanej dowodzą hemiciągłości.
Z~\eqref{eq:pde26-ll-coerc} otrzymujemy
$\langle Au,u\rangle\ge\alpha\|u\|_V^p-\|\beta\|_1$,
a zatem $A$ jest koercywny.
Wszystkie hipotezy twierdzenia~\ref{thm:pde26-minty}
są spełnione, więc istnieje rozwiązanie.
Dla dwóch rozwiązań $u_1,u_2$ testujemy różnicę
równań przez $u_1-u_2$. Jeśli monotoniczność jest
ostra, nieujemna całkowana funkcja może mieć całkę
zero tylko gdy $\nabla u_1=\nabla u_2$ prawie wszędzie.
Nierówność Poincarégo daje $u_1=u_2$.
\end{proof}

\begin{przyklad}[Monotoniczny strumień bez potencjału]
\label{ex:pde26-skewflux}
Niech $d=2$, $p=2$ i
$J=\begin{psmallmatrix}0&-1\\1&0\end{psmallmatrix}$.
Dla niekonstantnej funkcji $\kappa\in W^{1,\infty}(\Omega)$
kładziemy $a(x,\xi)=(I+\kappa(x)J)\xi$.
Wówczas $(a(x,\xi)-a(x,\eta))\cdot(\xi-\eta)
=|\xi-\eta|^2$ oraz $a(x,\xi)\cdot\xi=|\xi|^2$.
Twierdzenie~\ref{thm:pde26-leraylions} zapewnia jednoznaczne
rozwiązanie dla każdego $f\in H^{-1}(\Omega)$.
Macierz $I+\kappa J$ nie jest symetryczna, zatem
strumień $a(x,\cdot)$ nie jest gradientem funkcji
skalarnej względem $\xi$. Niesymetrycznego członu
nie odzyskalibyśmy, minimalizując samą energię
$\frac12\int_\Omega|\nabla u|^2\,\mathrm d\lambda_2$.
\end{przyklad}



\begin{uwaga}[Monotoniczność, energia i ograniczenia jednostronne]
\label{rem:pde-v44-minty-obstacle}
Twierdzenie~\ref{thm:pde26-leraylions} nie wymaga istnienia
potencjału: operator monotoniczny może być niesymetryczny,
jak w przykładzie~\ref{ex:pde26-skewflux}. Nie wolno więc
wnioskować, że każde takie równanie jest warunkiem stacjonarności
funkcjonału energii. W rozdziale~\ref{chap:pde29} korzystamy
z innej, dodatkowej struktury: \emph{symetryczna} forma
koercywna określa iloczyn skalarny na $H_0^1(\Omega)$,
a ograniczenie $u\ge\psi$ prowadzi do nierówności wariacyjnej.
Reprezentację jej rozwiązania jako rzutu w normie energetycznej
podamy w stwierdzeniu~\ref{prop:pde-v44-obstacle-projection}.
\end{uwaga}


\section{Porównanie, dodatniość i stabilność względem prawej strony}
\begin{twierdzenie}[Zasada porównania]\label{thm:pde26-comparison}
Niech $f,g\in V^*$, a $u,v\in V$ będą odpowiednio
jedynymi rozwiązaniami $A_pu=f$ i $A_pv=g$.
Jeśli $\langle f-g,\varphi\rangle\le0$ dla wszystkich
nieujemnych $\varphi\in V$, to $u\le v$ prawie wszędzie
w $\Omega$. W szczególności z $f\ge0$ w powyższym
porządku dualnym wynika $u\ge0$ prawie wszędzie.
\end{twierdzenie}
\begin{proof}
Kładziemy $w=(u-v)_+\in V$ i odejmujemy równania słabe:
\[
 \int_{\{u>v\}}(a_p(\nabla u)-a_p(\nabla v))\cdot
 (\nabla u-\nabla v)\,\mathrm d\lambda_d(x)
 =\langle f-g,w\rangle\le0.
\]
Całkowana funkcja jest nieujemna, a przy nierównych
gradientach dodatni, na mocy lematu~\ref{lem:pde26-vector}.
Stąd $\nabla w=0$ prawie wszędzie, a nierówność Poincarégo
daje $w=0$. Dla drugiego twierdzenia porównujemy z
zerowym rozwiązaniem równania $A_p0=0$.
\end{proof}

\begin{twierdzenie}[Ilościowa stabilność]\label{thm:pde26-stability}
Przy założeniach twierdzenia~\ref{thm:pde26-comparison}
niech $u,v$ odpowiadają dowolnym $f,g\in V^*$.
Jeśli $p\ge2$, to
\begin{equation}\label{eq:pde26-holder-data}
 \|u-v\|_V\le
 \bigl((p-1)2^{p-2}\|f-g\|_{V^*}\bigr)^{1/(p-1)}.
\end{equation}
Jeżeli $1<p<2$ i $\max\{\|u\|_V,\|v\|_V\}\le R$,
to istnieje $C_p>0$ takie, że
\begin{equation}\label{eq:pde26-lipschitz-data}
 \|u-v\|_V\le C_p(2R)^{2-p}\|f-g\|_{V^*}.
\end{equation}
W każdym przypadku odwzorowanie $f\mapsto u$
jest ciągłe z $V^*$ do $V$.
\end{twierdzenie}
\begin{proof}
Odejmując równania i testując przez $u-v$ otrzymujemy
\[
 \langle A_pu-A_pv,u-v\rangle
 =\langle f-g,u-v\rangle
 \le\|f-g\|_{V^*}\|u-v\|_V.
\]
Dla $p\ge2$ całkujemy~\eqref{eq:pde26-strong},
co daje $2^{2-p}\|u-v\|_V^p/(p-1)$ po lewej stronie,
a zatem~\eqref{eq:pde26-holder-data}.
Dla $1<p<2$ wprowadzamy
$W(x)=|\nabla u(x)|+|\nabla v(x)|$ i
$h=|\nabla u-\nabla v|$. Na zbiorze $W>0$ stosujemy
nierówność Höldera z wykładnikami $2/p$ i $2/(2-p)$:
\[
 \int_\Omega h^p\,\mathrm d\lambda_d
 \le \left(\int_\Omega W^{p-2}h^2\,\mathrm d\lambda_d\right)^{p/2}
 \left(\int_\Omega W^p\,\mathrm d\lambda_d\right)^{(2-p)/2}.
\]
Na $W=0$ mamy $h=0$, więc niczego nie tracimy.
Po podniesieniu do potęgi $2/p$ i wykorzystaniu
\eqref{eq:pde26-subquad} dostajemy
$\langle A_pu-A_pv,u-v\rangle
\ge(p-1)(2R)^{p-2}\|u-v\|_V^2$.
Dzielimy przez $\|u-v\|_V$ i otrzymujemy
\eqref{eq:pde26-lipschitz-data} z $C_p=1/(p-1)$.
Wreszcie oszacowanie~\eqref{eq:pde26-apriori} zapewnia
wspólne ograniczenie norm rozwiązań, gdy $f,g$
pozostają w ograniczonym zbiorze $V^*$.
\end{proof}

\section{Jawne rozwiązania: kule i pierścienie}
\begin{przyklad}[Rozwiązanie skrętnościowe w kuli]\label{ex:pde26-torsion}
Niech $B_R=\{x\in\R^d:|x|<R\}$, $R>0$, oraz
$\alpha=p/(p-1)$. Definiujemy
\begin{equation}\label{eq:pde26-torsion}
 u(x)=\frac{p-1}{p}\,d^{-1/(p-1)}
 \bigl(R^\alpha-|x|^\alpha\bigr),\qquad x\in\overline B_R.
\end{equation}
Wtedy $u\in W_0^{1,p}(B_R)$ i $-\Delta_pu=1$
w sensie słabym. Rzeczywiście
$\nabla u=-d^{-1/(p-1)}|x|^{\alpha-2}x$, a zatem
$|\nabla u|^{p-2}\nabla u=-x/d$.
Dywergencja ostatniego pola jest równa $-1$,
również w punkcie $x=0$ w sensie dystrybucyjnym.
Dla $p>2$ mamy $1<\alpha<2$:
rozwiązanie jest różniczkowalne z ciągłym gradientem,
ale nie należy do $C^2$ w środku kuli.
To przykład pokazujący, dlaczego nie można automatycznie
przenosić regularności $C^2$ z przypadku liniowego $p=2$.
\end{przyklad}

\begin{przyklad}[Rozwiązanie radialne w pierścieniu]\label{ex:pde26-annulus}
Niech $0<r_0<r_1<\infty$, $d\ge2$ i $u(x)=U(|x|)$
rozwiązuje $\Delta_pu=0$ dla $r_0<|x|<r_1$.
Wtedy, dla $r\in(r_0,r_1)$,
$\frac{\mathrm d}{\mathrm dr}(r^{d-1}|U'|^{p-2}U')=0$.
Stąd $U'(r)=c\,r^{-(d-1)/(p-1)}$ z odpowiednio
zdefiniowaną stałą $c\in\R$. Po całkowaniu
\[
 U(r)=\begin{cases}
 a+b\,r^{(p-d)/(p-1)},&p\ne d,\\
 a+b\log r,&p=d.
 \end{cases}
\]
Stałe $a,b$ dobiera się z dwóch wartości brzegowych.
Przejście do $p=d$ ma sens jako granica po odjęciu
stałej, a nie przez bezpośrednie podstawienie do pierwszego wzoru.
Pojedynczy punkt $r=0$ został wykluczony celowo:
radialne rozwiązanie w pierścieniu może mieć tam
osobliwość w sensie dystrybucyjnym.
\end{przyklad}

\begin{przyklad}[Jednowymiarowe równanie z prawą stroną]\label{ex:pde26-1d}
Niech $\Omega=(0,1)$ i $f\equiv1$. Jednoznaczne rozwiązanie
$-(|u'|^{p-2}u')'=1$ z $u(0)=u(1)=0$ wynosi
\[
 u(x)=\frac{p-1}{p}\left(2^{-p/(p-1)}
 -|x-\tfrac12|^{p/(p-1)}\right).
\]
W istocie strumień $q=|u'|^{p-2}u'$ jest dokładnie
$q(x)=\tfrac12-x$, więc $-q'=1$; warunki brzegowe
wynikają z symetrii profilu. Dla $p=2$ odzyskujemy
$u(x)=x(1-x)/2$.
\end{przyklad}

\section{Pierwsza wartość własna nieliniowego operatora}
\begin{definicja}[Wartość własna $p$-laplasjanu]
\label{def:pde26-eigen}
Liczbę $\lambda\in\R$ nazywamy wartością własną
Dirichleta $p$-laplasjanu na $\Omega$, jeśli istnieje
niezerowa funkcja $u\in W_0^{1,p}(\Omega)$ spełniająca
$-\Delta_pu=\lambda|u|^{p-2}u$ w $V^*$.
Definiujemy też liczbę
\[
 \lambda_{1,p}(\Omega)=
 \inf_{v\in V\setminus\{0\}}
 \frac{\int_\Omega|\nabla v|^p\,\mathrm d\lambda_d(x)}
 {\int_\Omega|v|^p\,\mathrm d\lambda_d(x)}.
\]
\end{definicja}
\begin{twierdzenie}[Osiągalność minimum i równanie własne]
\label{thm:pde26-eigen}
Niech $\Omega\subset\R^d$ będzie ograniczonym obszarem
Lipschitzowskim i $1<p<\infty$.
Wówczas $0<\lambda_{1,p}(\Omega)<\infty$ oraz
istnieje $u\in W_0^{1,p}(\Omega)$ z
$\|u\|_{L^p(\Omega)}=1$, $u\ge0$ prawie wszędzie,
takie, że
\[
 \int_\Omega|\nabla u|^{p-2}\nabla u\cdot\nabla v
 \,\mathrm d\lambda_d(x)=
 \lambda_{1,p}(\Omega)
 \int_\Omega |u|^{p-2}uv\,\mathrm d\lambda_d(x)
 \quad(v\in V).
\]
Dla $r>0$ zachodzi ponadto prawo skalowania
$\lambda_{1,p}(r\Omega)=r^{-p}\lambda_{1,p}(\Omega)$.
\end{twierdzenie}
\begin{proof}
Z nierówności Poincarégo mamy
$\|v\|_p\le C\|\nabla v\|_p$,
więc iloraz Rayleigha jest ograniczony od dołu przez
$C^{-p}>0$. Skończoność otrzymujemy po podstawieniu
dowolnej niezerowej funkcji testowej z $C_c^\infty(\Omega)$.
Bierzemy ciąg minimalizujący z $\|v_n\|_p=1$.
Jego gradienty są ograniczone w $L^p$. Dzięki
refleksywności, po przejściu do podciągu,
$v_n\rightharpoonup u$ w $V$. Zwarte zanurzenie
$W_0^{1,p}(\Omega)\hookrightarrow L^p(\Omega)$
daje $v_n\to u$ silnie w $L^p$, a więc $\|u\|_p=1$.
Słaba półciągłość normy gradientu dowodzi osiągalności.
Ponieważ $|u|\in V$, $\||u|\|_p=1$ oraz
$\|\nabla|u|\|_p=\|\nabla u\|_p$, możemy wybrać
minimizer nieujemny. Dla ustalonego $v\in V$
różniczkujemy przy $t=0$ iloraz Rayleigha funkcji
$u+tv$; warunek zerowania pochodnej, znormalizowanie
$\|u\|_p=1$ i $\int_\Omega|\nabla u|^p\,\mathrm d\lambda_d(x)=\lambda_{1,p}$
dają wypisaną tożsamość słabą.
Wreszcie podstawienie $v(x)=w(x/r)$
w definicji wartości własnej na $r\Omega$ i
zmiana zmiennych $x=ry$ mnoży licznik przez
$r^{d-p}$, a mianownik przez $r^d$. Stąd
prawo skalowania. Nie twierdzimy tutaj, że
nieujemny minimizer jest punktowo ściśle dodatni;
taki wniosek wymaga dodatkowej teorii regularności.
\end{proof}

\section{Monotoniczny schemat różnicowy w jednym wymiarze}
Dzięki jednowymiarowej postaci strumienia możemy
zbudować schemat numeryczny, którego jednoznaczność
wynika z tej samej monotoniczności co twierdzenie
\ref{thm:pde26-plap}. Co istotne, dla stałej prawej
strony nie musimy korzystać z iteracyjnego solvera
nieliniowego.

\begin{twierdzenie}[Dokładna konstrukcja dyskretnego strumienia]
\label{thm:pde26-discrete}
Niech $p>1$ i $N\ge2$ będzie parzystą liczbą całkowitą.
Ustalmy $h=1/N$ oraz $x_i=ih$ dla $0\le i\le N$.
Szukamy liczb $U_0,\ldots,U_N$ spełniających
$U_0=U_N=0$ oraz dla $1\le i\le N-1$ układ
\begin{equation}\label{eq:pde26-discrete}
 -\frac{Q_{i+1/2}-Q_{i-1/2}}h=1,\qquad
 Q_{i+1/2}=\left|\frac{U_{i+1}-U_i}h\right|^{p-2}
 \frac{U_{i+1}-U_i}h.
\end{equation}
Układ ma dokładnie jedno rozwiązanie, określone przez
\begin{equation}\label{eq:pde26-flux}
 Q_{i+1/2}=\frac12-\left(i+\frac12\right)h,\qquad
 U_i=h\sum_{j=0}^{i-1}\operatorname{sgn}(Q_{j+1/2})
 |Q_{j+1/2}|^{1/(p-1)}.
\end{equation}
Jeśli $u$ jest rozwiązaniem jawnym z
przykładu~\ref{ex:pde26-1d}, to przy
$\beta=1/(p-1)$ zachodzi
\[
 \max_{0\le i\le N}|U_i-u(x_i)|
 \le C_p h^{\min\{2,1+\beta\}},
\]
gdzie $C_p$ nie zależy od $N$ (dla $p=2$ błąd
w węzłach jest dokładnie zerowy).
\end{twierdzenie}
\begin{proof}
\textbf{Krok 1: bilans strumienia.} Pierwsze równanie
\eqref{eq:pde26-discrete} daje
$Q_{i+1/2}=C-(i+1/2)h$. Odwrotnością funkcji
$a_p(s)=|s|^{p-2}s$ jest
$b_p(q)=\operatorname{sgn}(q)|q|^{1/(p-1)}$.
Zatem $U_{i+1}-U_i=h\,b_p(Q_{i+1/2})$.
Warunek $U_N=U_0$ wymaga
$\sum_{i=0}^{N-1}b_p(C-(i+1/2)h)=0$.
Lewa strona jest ściśle rosnącą i ciągłą funkcją
$C$, a dla $C=1/2$ wyrazy sumy znoszą się parami
przez symetrię względem środka odcinka. Stąd jedyne
$C$ jest równe $1/2$. Formuła~\eqref{eq:pde26-flux}
spełnia więc oba warunki brzegowe i wszystkie równania.

\textbf{Krok 2: niezależny dowód jednoznaczności.}
Dla dwóch rozwiązań $U,V$ pomnóżmy różnicę
równań przez $h(U_i-V_i)$ i zsumujmy po węzłach
wewnętrznych. Dyskretne całkowanie przez części,
wykorzystujące zerowe wartości na końcach, daje
\[
 \sum_{i=0}^{N-1}h\bigl(a_p(D_iU)-a_p(D_iV)\bigr)
 (D_iU-D_iV)=0,
\]
gdzie $D_iU=(U_{i+1}-U_i)/h$. Ścisła monotoniczność
$a_p$ implikuje $D_iU=D_iV$ dla wszystkich $i$,
a zatem $U=V$.

\textbf{Krok 3: zbieżność.} Wzór~\eqref{eq:pde26-flux}
wyraża $U_i$ jako sumę środków przedziałów dla
całki $\int_0^{x_i}b_p(1/2-x)\,\mathrm dx$,
która jest równa $u(x_i)$.
Na przedziałach odległych od $1/2$ druga pochodna
funkcji $g(x)=b_p(1/2-x)$ spełnia
$|g''(x)|\le C_p|x-1/2|^{\beta-2}$.
Błąd pojedynczego przedziału metody punktu środkowego
jest nie większy niż $C h^3\sup|g''|$.
Po zsumowaniu przedziałów odległych o co najmniej $jh$
od środka otrzymujemy oszacowanie
$C_ph^{1+\beta}\sum_{j=1}^{N/2}j^{\beta-2}$.
Przedziały najbliższe środka oszacowujemy bezpośrednio
przez całkę z $|g|$, co daje $C_ph^{1+\beta}$.
Dla $\beta<1$ wspomniany szereg jest jednostajnie
ograniczony; dla $\beta>1$ jego suma jest
$O(N^{\beta-1})$, co daje $O(h^2)$.
Gdy $\beta=1$, funkcja $g$ jest afiniczna i metoda
punktu środkowego jest dokładna. Dowód jest zakończony.
\end{proof}

\section{Granice twierdzeń i przypadki niemonotoniczne}
\begin{kontrprzyklad}[Zanik monotoniczności po odjęciu potencjału]
\label{ex:pde26-nonmonotone}
Dla $\lambda>0$ rozważmy
$Bu=A_pu-\lambda|u|^{p-2}u$ w $V^*$,
przy czym drugi składnik rozumiemy przez
$\langle |u|^{p-2}u,v\rangle=
\int_\Omega|u|^{p-2}uv\,\mathrm d\lambda_d(x)$.
Jeśli $\lambda>\lambda_{1,p}(\Omega)$, bierzemy
minimizer $\varphi$ z twierdzenia~\ref{thm:pde26-eigen}.
Wtedy
$\langle B\varphi-B0,\varphi-0\rangle
=\lambda_{1,p}-\lambda<0$, więc $B$ nie jest monotoniczny.
Dla $\lambda=\lambda_{1,p}$ równanie $Bu=0$ ma
co najmniej dwa różne rozwiązania: $u=0$ i $u=\varphi$.
Nie wolno więc automatycznie stosować twierdzenia
Minty'ego--Browdera po dodaniu wyrazu o niekorzystnym znaku.
\end{kontrprzyklad}

\begin{uwaga}[Co zostaje na później?]
Nie udowodniliśmy regularności $C^{1,\alpha}$
ani nierówności Harnacka dla $p$-harmonicznych funkcji;
nie wynikają one z samej monotoniczności. Również
nierówności wariacyjne z przeszkodą wymagają
modyfikacji relacji $Au=f$ na nierówność w zbiorze wypukłym;
wrócimy do nich w rozdziale 29.
Zaproponowana tu metoda Galerkina obejmuje również
strumienie typu Leraya--Lionsa przy właściwych
hipotezach wzrostu, monotoniczności i koercywności,
ale pełnego twierdzenia z zależnością od $x$ i $u$
nie ogłaszamy bez osobnej weryfikacji warunków.
\end{uwaga}


% ROZDZIAŁ 27: reakcja i dyfuzja
\chapter[Reakcja--dyfuzja i fale biegnące]{Równania reakcji--dyfuzji, fale biegnące, dynamika długoczasowa i niestabilność Turinga}
\label{chap:pde27}

\section{Mechanizmy reakcji i dyfuzji: zakres rozdziału}
Rozważamy skalarne równanie
\begin{equation}\label{eq:pde27-main}
 u_t=D\Delta u+f(u),\qquad D>0,
\end{equation}
gdzie dyfuzja wygładza rozkład przestrzenny, a reakcja opisuje
lokalną dynamikę. Badamy zarówno torus
$\T_L^d=(\R/L\mathbb Z)^d$, jak i prostą $\R$.
Na torusie stosujemy periodyczne warunki brzegowe;
na prostej badamy fronty $u(t,x)=U(x-ct)$.
Jeżeli $f$ jest tylko lokalnie lipschitzowska, rozwiązanie
nie musi istnieć dla wszystkich czasów (por. rozdział~\ref{chap:pde25}).
W szczególności nie wolno wyprowadzać zachowania długoczasowego
z samego twierdzenia o lokalnym istnieniu.

W obliczeniach dla modelu Fishera--KPP użyjemy
$f(s)=rs(1-s)$, $r>0$. Rozwiązanie z zakresu $0\le u\le1$
pozostaje w tym przedziale, mimo że $f$ nie jest monotoniczna
na całej osi. Jest to zastosowanie zasady porównania,
nie monotoniczności funkcji $f$.

\section{Istnienie, ciągłość względem danych i zasada porównania}
\begin{twierdzenie}[Istnienie lokalne i alternatywa eksplozji]\label{thm:pde27-local}
Niech $L,D>0$, $d\ge1$, $f\in C^1(\R)$ oraz
$u_0\in C(\T_L^d)$. Istnieje dokładnie jedno maksymalne
rozwiązanie łagodne $u\in C([0,T_{\max});C(\T_L^d))$
równania~\eqref{eq:pde27-main} z $u(0)=u_0$.
Dla każdego $T<T_{\max}$ zachodzi wzór Duhamela
\[
 u(t)=S_D(t)u_0+\int_0^tS_D(t-s)f(u(s))\,\mathrm ds,
 \qquad S_D(t)=e^{tD\Delta}.
\]
Jeżeli $T_{\max}<\infty$, to
$\|u(t)\|_{\infty}\to\infty$ przy $t\uparrow T_{\max}$.
Jeśli $u,v$ są rozwiązaniami o wartościach w $[-M,M]$
na $[0,T]$, to dla $K=\max_{|s|\le M}|f'(s)|$ mamy
\[
 \|u(t)-v(t)\|_{\infty}\le e^{Kt}\|u_0-v_0\|_{\infty},
 \qquad 0\le t\le T.
\]
\end{twierdzenie}
\begin{proof}
Na przestrzeni $X_T=C([0,T];C(\T_L^d))$
rozważamy operator prawej strony wzoru Duhamela.
Półgrupa cieplna jest kontrakcją w normie maksimum,
a na kuli $\|w\|_{X_T}\le M$ pochodna $f'$ jest ograniczona.
Jeśli $M>\|u_0\|_{\infty}+1$, wybieramy $T$ tak,
aby $\|f\|_{L^{\infty}([-M,M])}T\le
M-\|u_0\|_\infty$ i $KT<1$. Operator odwzorowuje
kulę w siebie i jest kontrakcją. Twierdzenie Banacha
daje istnienie i jednoznaczność.
Czas istnienia zależy od pułapu normy, nie od momentu
rozpoczęcia ewolucji. Gdyby przy skończonym $T_{\max}$
norma rozwiązania była ograniczona wzdłuż pewnego ciągu
tendującego do $T_{\max}$, można by uruchamiać kontrakcję
od tych chwil na wspólnym przedziale o dodatniej długości,
sprzecznie z maksymalnością. To dowodzi alternatywy.
Dla różnicy $w=u-v$ wzór Duhamela i twierdzenie o wartości
średniej dają
$\|w(t)\|_\infty\le\|w(0)\|_\infty+
K\int_0^t\|w(s)\|_\infty\,\mathrm ds$.
Lemat Gronwalla kończy dowód stabilności.
\end{proof}


\begin{stwierdzenie}[Skalowanie półgrupy cieplnej na torusie]
\label{prop:pde-v44-torus-scaling}
Niech $L,D>0$ i $d\ge1$. Dla $v\in C(\T_L^d)$ połóżmy
$(Rv)(\theta)=v(L\theta/(2\pi))$, $\theta\in\T^d$.
Wówczas $R:C(\T_L^d)\to C(\T^d)$ jest izometrią i
\begin{equation}\label{eq:pde-v44-torus-scale}
 R\,e^{tD\Delta_x}\,R^{-1}
 =e^{\,tD(2\pi/L)^2\Delta_\theta},\qquad t\ge0.
\end{equation}
Po podstawieniu $\tau=D(2\pi/L)^2t$ oraz
$w(\tau,\theta)=u(t,L\theta/(2\pi))$ równanie
$u_t=D\Delta_xu+f(u)$ przyjmuje postać
$w_\tau=\Delta_\theta w+\bigl[D(2\pi/L)^2\bigr]^{-1}f(w)$.
W konsekwencji konstrukcja lokalna i alternatywa eksplozji
z rozdziału~\ref{chap:pde25} stosują się do twierdzenia
\ref{thm:pde27-local} po tej zmianie zmiennych.
\end{stwierdzenie}
\begin{proof}
Z reguły łańcuchowej dla funkcji gładkiej
$R(\Delta_xv)=(2\pi/L)^2\Delta_\theta(Rv)$.
Dla wielomianów trygonometrycznych tożsamość półgrup
wynika także z działania na każdej częstotliwości
$k\in\mathbb Z^d$: obie strony mnożą współczynnik
Fouriera przez $\exp[-tD(2\pi/L)^2|k|^2]$.
Gęstość wielomianów trygonometrycznych w $C(\T^d)$
i kontraktywność półgrup rozszerzają równość na każdą
funkcję ciągłą. Zmiana czasu daje ostatni wzór.
Ponieważ zmodyfikowana nieliniowość nadal należy do
$C^1(\R)$, stosujemy twierdzenie~\ref{thm:pde25-local}.
\end{proof}


\begin{twierdzenie}[Zasada porównania i przedziały niezmiennicze]
\label{thm:pde27-comparison}
Niech $f\in C^1(\R)$, a $u,v$ będą rozwiązaniami z
poprzedniego twierdzenia, istniejącymi na $[0,T]$.
Jeżeli $u_0\le v_0$, to $u(t,\cdot)\le v(t,\cdot)$
dla każdego $0\le t\le T$. Jeśli $a\le b$,
$f(a)\ge0$ i $f(b)\le0$, to dane $a\le u_0\le b$
implikują $a\le u(t,\cdot)\le b$ na całym
przedziale istnienia; wówczas $T_{\max}=\infty$.
\end{twierdzenie}
\begin{proof}
Najpierw dla gładkich danych rozwiązania są klasyczne
dla $t>0$. Różnica $w=u-v$ spełnia
$w_t-D\Delta w=a(t,x)w$, gdzie
\[
 a(t,x)=\int_0^1 f'\bigl(v(t,x)+\theta(u(t,x)-v(t,x))\bigr)
 \,\mathrm d\theta
\]
jest ograniczona. Po pomnożeniu równania przez
$w_+=\max\{w,0\}$, scałkowaniu na torusie i zastosowaniu
nierówności Gronwalla otrzymujemy $w_+=0$, gdy $w_+(0)=0$.
Dla ogólnych ciągłych danych przechodzimy do granicy,
korzystając ze stabilności z twierdzenia~\ref{thm:pde27-local}.
Stałe $a,b$ są odpowiednio pod- i nadrozwiązaniami,
gdyż $0\le f(a)$ i $0\ge f(b)$; ich porównanie z $u$
można przeprowadzić tą samą metodą dla nierówności.
Przedział $[a,b]$ jest niezmienniczy, a ograniczenie normy
wyklucza skończony czas maksymalny.
\end{proof}

\begin{kontrprzyklad}[Porównanie nie wymaga $f'\ge0$]
Dla $f(s)=s(1-s)$ pochodna $f'(s)=1-2s$
jest ujemna na $(1/2,\infty)$.
Mimo to zasada porównania działa, ponieważ różnica
$f(u)-f(v)$ ma postać $a(t,x)(u-v)$ z ograniczonym
współczynnikiem na każdym ograniczonym przedziale czasu.
Dla układów reakcji--dyfuzji takie porównanie
składowa po składowej na ogół nie zachodzi bez
warunku kooperatywności.
\end{kontrprzyklad}

\section{Równanie Fishera--KPP: zbieżność do stanu nasycenia na torusie}
\begin{twierdzenie}[Zbieżność do stanu nasycenia]\label{thm:pde27-kpp-torus}
Niech $u_0\in C(\T_L^d)$, $0\le u_0\le1$ i
$u_0\not\equiv0$. Dla $D,r>0$ jedyne rozwiązanie
zagadnienia
$u_t=D\Delta u+ru(1-u)$ z danymi $u_0$
istnieje globalnie i
\[
 \lim_{t\to\infty}\|u(t,\cdot)-1\|_{C(\T_L^d)}=0.
\]
Jeśli $u_0\equiv0$, to $u(t,\cdot)\equiv0$.
\end{twierdzenie}
\begin{proof}
Porównanie ze stałymi rozwiązaniami $0$ i $1$
daje $0\le u\le1$ i istnienie globalne.
Dodatniość jądra cieplnego na spójnym torusie oraz
wzór Duhamela z nieujemnym źródłem $ru(1-u)$
pokazują, że dla każdego $t_0>0$ funkcja
$u(t_0,\cdot)$ jest ściśle dodatnia. Zwarta przestrzeń
$\T_L^d$ zapewnia minimum $m_0>0$.
Niech $y'=ry(1-y)$ i $z'=rz(1-z)$,
$y(t_0)=m_0$, $z(t_0)=\max u(t_0,\cdot)\le1$.
Porównanie daje $y(t)\le u(t,x)\le z(t)$.
Z jawnego rozwiązania równania logistycznego
\[
 y(t)=\frac{1}{1+(m_0^{-1}-1)e^{-r(t-t_0)}}
\]
wynika $y(t)\uparrow1$, podczas gdy $z(t)\le1$.
To dowodzi zbieżności jednostajnej.
\end{proof}

\begin{uwaga}[Dlaczego torus, a nie cała prosta?]
Na $\R$ rozwiązanie wychodzące z danych o zwartym nośniku
nie może zbiegać jednostajnie do $1$ na całej prostej
w skończonym czasie: daleko od początkowego nośnika
pozostaje bardzo małe. Właściwą asymptotyką jest
\emph{propagacja frontu}, z którą zapoznamy się dalej.
\end{uwaga}


\begin{wniosek}[Ilościowa zbieżność rozwiązania Fishera--KPP na torusie]
\label{cor:pde27-kpp-rate}
Przy hipotezach twierdzenia~\ref{thm:pde27-kpp-torus}
niech $t_0>0$ i
$m_{t_0}=\min_{x\in\T_L^d}(S_D(t_0)u_0)(x)>0$.
Wówczas dla wszystkich $t\ge t_0$
\begin{equation}\label{eq:pde27-kpp-rate}
 \|1-u(t,\cdot)\|_\infty\le
 \frac{1-m_{t_0}}{m_{t_0}}e^{-r(t-t_0)}.
\end{equation}
Ponadto jeśli $K_{D,t_0}^L$ oznacza dodatnie
periodyczne jądro cieplne na $\T_L^d$, to
\[
 m_{t_0}\ge\left(\min_{z\in\T_L^d}K_{D,t_0}^L(z)\right)
 \int_{\T_L^d}u_0(x)\,\mathrm d\lambda_d(x)>0.
\]
\end{wniosek}
\begin{proof}
Nieujemne źródło $ru(1-u)$ i wzór Duhamela dają
$u(t_0)\ge S_D(t_0)u_0$, stąd
$\min u(t_0)\ge m_{t_0}>0$.
Rozwiązanie skalarnego równania
$y'=ry(1-y)$ z $y(t_0)=m_{t_0}$ wynosi
$y(t)=(1+(m_{t_0}^{-1}-1)e^{-r(t-t_0)})^{-1}$.
Porównanie ze stałym w przestrzeni rozwiązaniem $y$
i nadrozwiązaniem $1$ daje
$0\le1-u(t,x)\le1-y(t)$.
Ponieważ $1-(1+a)^{-1}=a/(1+a)\le a$
dla $a\ge0$, otrzymujemy~\eqref{eq:pde27-kpp-rate}.
Ostatnia nierówność wynika bezpośrednio z dodatniości
i ciągłości jądra na zwartym torusie.
\end{proof}


\section{Fale biegnące: od PDE do układu równań zwyczajnych}
Niech $d=1$ i rozważmy równanie
$u_t=Du_{xx}+ru(1-u)$ na $\R$.
Dla postaci fali biegnącej $u(t,x)=U(z)$, $z=x-ct$,
otrzymujemy równanie profilu
\begin{equation}\label{eq:pde27-wave}
 DU''+cU'+rU(1-U)=0.
\end{equation}
Szukamy funkcji $U\in C^2(\R)$, dla której
$U(-\infty)=1$, $U(+\infty)=0$ i $U'<0$.
Prędkość $c$ jest prędkością przemieszczania się
profilu w prawo. Sama zamiana zmiennych
nie dowodzi jeszcze istnienia takiego profilu.

\begin{lemat}[Obszar niezmienniczy dla układu profilowego]
\label{lem:pde27-wedge}
Niech $D,r>0$, $c\ge2\sqrt{Dr}$ oraz
\[
 k=\frac{c-\sqrt{c^2-4Dr}}{2D}>0.
\]
Wtedy obszar
$\mathcal W=\{(U,q):0<U<1,\ 0<q<kU\}$
jest dodatnio niezmienniczy dla układu
\begin{equation}\label{eq:pde27-phase}
 U'=-q,\qquad q'=-\frac cD q+\frac rD U(1-U).
\end{equation}
Oznacza to, że trajektorie rozpoczynające się w $\mathcal W$
pozostają w nim dla dodatnich czasów istnienia.
Ponadto rozwiązanie wychodzące z punktu
$U(0)=1-\delta$, $q(0)=0$, gdzie $0<\delta<1/2$,
zachowuje $0<U<1$, $0<q<kU$ dla $z>0$ i
spełnia $U(z)\to0$, $q(z)\to0$ przy $z\to\infty$.
\end{lemat}
\begin{proof}
Na brzegu $q=0$, gdy $0<U<1$, mamy
$q'=rU(1-U)/D>0$. Na brzegu $q=kU$ otrzymujemy
\[
 (q-kU)'=\bigl(k^2-\tfrac cDk+\tfrac rD(1-U)\bigr)U
 =-\tfrac rD U^2<0,
\]
ponieważ $k^2-(c/D)k+r/D=0$.
Na brzegu $U=1$ wektor wskazuje do mniejszych wartości
$U$; z kolei $q<kU$ nie pozwala osiągnąć $U=0$
w skończonym czasie, gdyż $U'=-q\ge-kU$.
Standardowa jednoznaczność ODE i powyższe nierówności
zapewniają niezmienniczość. Dla danych z $q(0)=0$
natychmiast wchodzimy do wnętrza $\mathcal W$.
Funkcja $U$ jest malejąca, więc ma granicę
$\ell\in[0,1-\delta)$. Ponieważ $q=-U'\ge0$,
$\int_0^\infty q(z)\,\mathrm dz<\infty$.
Pochodna $q'$ jest ograniczona w $\mathcal W$;
stąd jednostajna ciągłość $q$ wraz z całkowalnością
daje $q(z)\to0$. Gdyby $\ell>0$, drugie
równanie~\eqref{eq:pde27-phase} prowadziłoby do
$q'(z)\to r\ell(1-\ell)/D>0$, co jest sprzeczne
z $q(z)\to0$. W konsekwencji $\ell=0$.
\end{proof}

\begin{twierdzenie}[Prędkości fal biegnących równania Fishera--KPP]
\label{thm:pde27-wave}
Niech $D,r>0$. Dla każdej $c\ge2\sqrt{Dr}$
istnieje $U\in C^2(\R)$ spełniająca
\eqref{eq:pde27-wave}, $0<U<1$, $U'<0$ oraz
$U(-\infty)=1$, $U(+\infty)=0$.
Natomiast dla żadnej $c<2\sqrt{Dr}$
taki profil nie istnieje. W szczególności
minimalna prędkość monotonicznej fali wynosi
$c_*=2\sqrt{Dr}$.
\end{twierdzenie}
\begin{proof}
\textbf{Krok 1: skonstruowanie pełnej trajektorii.}
Ustalmy $c\ge2\sqrt{Dr}$ i $k$ z
lematu~\ref{lem:pde27-wedge}. Dla $\delta\downarrow0$
rozważmy rozwiązania $(U_\delta,q_\delta)$
rozpoczynające się w $(1-\delta,0)$.
Każda trajektoria osiąga poziom $U=1/2$ dokładnie raz;
niech $t_\delta>0$ oznacza odpowiedni czas.
Pokażemy, że $t_\delta\to\infty$.
Kładziemy $v=1-U$ i korzystamy z równań
$v'=q$, $q'=-(c/D)q+(r/D)(1-v)v\le(r/D)v$.
Dla $v,q\ge0$ dostajemy
$(v+q)'\le C(v+q)$ z $C=\max\{1,r/D\}$,
a zatem $v(t)+q(t)\le\delta e^{Ct}$.
Gdy $U(t_\delta)=1/2$, mamy $v(t_\delta)=1/2$,
więc $t_\delta\ge C^{-1}\log(1/(2\delta))\to\infty$.

Przesuwamy trajektorie: $\widetilde U_\delta(z)
=U_\delta(t_\delta+z)$, analogicznie dla $q$.
Są one określone na $(-t_\delta,\infty)$,
mają wartości w zwartym domknięciu $\mathcal W$
i jednostajnie ograniczone pochodne.
Równania ODE ograniczają również drugie pochodne na
przedziałach zwartych. Twierdzenie Arzeli--Ascolego
i diagonalizacja dają podciąg zbieżny w $C^1_{\mathrm{loc}}(\R)$
do pełnej trajektorii $(U,q)$ układu~\eqref{eq:pde27-phase},
z $U(0)=1/2$. Ponieważ $q\ge0$ i $U'=-q$,
$U$ jest nierosnąca i ma granice $\ell_\pm$
przy $z\to\pm\infty$. Analogicznie jak w lemacie,
całkowalność $q$ na obu półprostych i jego
jednostajna ciągłość implikują $q(z)\to0$
na obu krańcach. Granice układu muszą zatem
spełniać $\ell_\pm(1-\ell_\pm)=0$.
Z $U(0)=1/2$ i monotoniczności wynika
$\ell_-=1$, $\ell_+=0$. Ponieważ $U$ ma różne granice,
nie może być stałe; z równań wynika wówczas $0<U<1$.
Gdyby $q(z_0)=0$ przy $0<U(z_0)<1$,
to $q'(z_0)>0$, co jest niemożliwe przy $q\ge0$
na całej prostej. Stąd $q>0$ i $U'<0$.

\textbf{Krok 2: wykluczenie prędkości mniejszych od krytycznej.}
Jeżeli $c\le0$, po podstawieniu $q=-U'>0$
do~\eqref{eq:pde27-phase} otrzymujemy $q'>0$,
a to przeczy jednocześnie całkowalności
$\int_\R q(z)\,\mathrm dz=1$ i monotoniczności $q$.
Niech teraz $0<c<2\sqrt{Dr}$. Dla $z$ dostatecznie dużych
$U(z)$ jest tak małe, że
$\mu(z)=\frac rD(1-U(z))-\frac{c^2}{4D^2}\ge\kappa^2>0$.
Po podstawieniu $y(z)=e^{cz/(2D)}U(z)>0$
równanie~\eqref{eq:pde27-wave} przyjmuje postać
$y''+\mu(z)y=0$.
Na przedziale $[a,a+\pi/\kappa]$, gdzie $a$ jest
dostatecznie duże, mnożymy równanie przez
$\sin(\kappa(z-a))$ i całkujemy względem $z$.
Po dwukrotnym całkowaniu przez części dostajemy
\[
 \int_a^{a+\pi/\kappa}(y''+\kappa^2y)
 \sin(\kappa(z-a))\,\mathrm dz
 =\kappa\bigl(y(a)+y(a+\pi/\kappa)\bigr)>0.
\]
Tymczasem $y''+\kappa^2y=-(\mu-\kappa^2)y\le0$,
a sinus jest nieujemny na całym przedziale.
Sprzeczność kończy dowód.
\end{proof}

\begin{przyklad}[Jawna fala o prędkości $5/\sqrt6$]
\label{ex:pde27-explicit}
Dla $D=r=1$ i $c=5/\sqrt6$ profil
\[
 U(z)=\frac{1}{(1+e^{z/\sqrt6})^2}
\]
rozwiązuje~\eqref{eq:pde27-wave} i spełnia
$U(-\infty)=1$, $U(+\infty)=0$, $U'<0$.
Podstawiając $w=e^{z/\sqrt6}$ mamy
$U'=-2w/(\sqrt6(1+w)^3)$ oraz
$U''=-w/(3(1+w)^3)+w^2/(1+w)^4$;
po sprowadzeniu do wspólnego mianownika
wyrażenie $U''+(5/\sqrt6)U'+U(1-U)$ jest zerowe.
Skalowanie $z\mapsto\sqrt{r/D}\,z$ i
$c\mapsto c/\sqrt{Dr}$ daje analogiczną falę
dla ogólnych $D,r>0$ z $c=(5/\sqrt6)\sqrt{Dr}$.
\end{przyklad}

\section{Granice propagacji na prostej}
\begin{twierdzenie}[Górne oszacowanie prędkości ekspansji]
\label{thm:pde27-spreading-upper}
Niech $D,r>0$, a $u_0\in C_c(\R)$ spełnia
$0\le u_0\le1$. Niech $u$ będzie globalnym
rozwiązaniem równania $u_t=Du_{xx}+ru(1-u)$.
Wówczas dla każdej $c>2\sqrt{Dr}$ istnieją
stałe $C,\gamma>0$, zależne od $c,u_0,D,r$, takie, że
\[
 \sup_{|x|\ge ct}u(t,x)\le C e^{-\gamma t}
 \quad\text{dla wszystkich }t\ge0.
\]
\end{twierdzenie}
\begin{proof}
Ponieważ $0\le u\le1$ oraz $ru(1-u)\le ru$, funkcja
$w(t,x)=e^{-rt}u(t,x)$ spełnia nierówność
$w_t-Dw_{xx}\le0$ oraz $w(0,x)=u_0(x)$.
Z zasady porównania dla równania ciepła, najpierw
dla gładkich danych, a potem przez aproksymację,
uzyskujemy na całej prostej
\[
 0\le u(t,x)\le e^{rt}\int_{\R}
 \frac{e^{-(x-y)^2/(4Dt)}}{\sqrt{4\pi Dt}}u_0(y)\,\mathrm dy,
 \qquad t>0.
\]
Porównanie to można również wyprowadzić bezpośrednio
ze wzoru Duhamela: $w_t-Dw_{xx}\le0$ oznacza,
że $S_D(t)u_0-w(t)$ jest całką splotową
nieujemnej funkcji z dodatnim jądrem cieplnym.

Wybierzmy $R>0$ tak, aby $\supp u_0\subset[-R,R]$,
oraz liczbę $c_0$ z $2\sqrt{Dr}<c_0<c$.
Niech $t_0=\max\{1,R/(c-c_0)\}$. Dla $t\ge t_0$
i $|x|\ge ct$ zachodzi $|x-y|\ge c_0t$
dla każdego $y\in\supp u_0$. Stąd
\[
 u(t,x)\le\frac{\|u_0\|_{L^1(\R)}}{\sqrt{4\pi Dt}}
 \exp\!\left(-\left(\frac{c_0^2}{4D}-r\right)t\right).
\]
Liczba $\gamma=c_0^2/(4D)-r$ jest dodatnia,
a czynnik $t^{-1/2}\le1$ dla $t\ge t_0$.
Na odcinku $0\le t<t_0$ korzystamy z $u\le1$.
Po zwiększeniu stałej $C$ otrzymujemy
$\sup_{|x|\ge ct}u(t,x)\le Ce^{-\gamma t}$
dla wszystkich $t\ge0$.\end{proof}

\begin{uwaga}[Co udowodniono, a czego nie]
Twierdzenie~\ref{thm:pde27-spreading-upper} daje wyłącznie
górne oszacowanie szybkości rozprzestrzeniania się.
Dopasowana dolna granica prędkości i pełne twierdzenie
o asymptotyce rozwiązania na całej prostej wymagają
osobnej konstrukcji podrozwiązań. Nie wynikają
same z istnienia fal ani z powyższej bariery.
\end{uwaga}

\section{Funkcjonały Lapunowa i równanie Allena--Cahna}
\begin{twierdzenie}[Dyssypacja energii i stany graniczne]
\label{thm:pde27-allen}
Niech $D,L>0$, $d\ge1$, a $u_0\in C^2(\T_L^d)$.
Niech $u$ będzie globalnym rozwiązaniem
$u_t=D\Delta u+u-u^3$ z danymi $u_0$.
Dla funkcjonału
\[
 E(u)=\int_{\T_L^d}\left(\frac D2|\nabla u|^2
 +\frac14(u^2-1)^2\right)\,\mathrm d\lambda_d(x)
\]
zachodzi dla $t\ge0$ tożsamość
$E(u(t))+\int_0^t\|u_s(s)\|_2^2\,\mathrm ds=E(u_0)$.
Ponadto istnieje ciąg $t_n\to\infty$ i funkcja
$u_\infty\in H^1(\T_L^d)$ takie, że
$u(t_n)\to u_\infty$ silnie w $L^2$ i
$-D\Delta u_\infty=u_\infty-u_\infty^3$
w sensie słabym. Nie twierdzimy bez dodatkowych
założeń o zbieżności całej trajektorii do jednego
stanu stacjonarnego.
\end{twierdzenie}
\begin{proof}
Wybieramy $M\ge\max\{1,\|u_0\|_\infty\}$.
Stałe $-M$ i $M$ są odpowiednio pod- i
nadrozwiązaniem, gdyż $f(-M)=-M+M^3\ge0$
i $f(M)=M-M^3\le0$. Z zasady porównania
otrzymujemy $|u(t,x)|\le M$ i globalność.
Dla gładkich rozwiązań różniczkujemy energię
i całkujemy przez części na torusie:
\[
 \frac{\mathrm d}{\mathrm dt}E(u(t))
 =\int_{\T_L^d}\bigl(-D\Delta u+u^3-u\bigr)u_t
 \,\mathrm d\lambda_d(x)
 =-\|u_t(t)\|_2^2.
\]
Aproksymacja i wygładzanie półgrupowe pozwalają
zachować tożsamość dla rozwiązań o podanej regularności.
W szczególności
$\sup_{t\ge0}\|u(t)\|_{H^1}<\infty$ oraz
$\int_0^\infty\|u_t(t)\|_2^2\,\mathrm dt<\infty$.
Istnieją więc $t_n\to\infty$, dla których
$\|u_t(t_n)\|_2\to0$: w przeciwnym razie istniałyby
$\eta>0$ i $t_*>0$ takie, że $\|u_t(t)\|_2\ge\eta$
dla prawie każdego $t\ge t_*$, wbrew całkowalności
kwadratu tej normy. Chwile $t_n$ wybieramy z pełnomiarowego
zbioru, na którym obowiązuje słaba postać równania.
Kompaktowość zanurzenia $H^1(\T_L^d)\hookrightarrow L^2$
pozwala po przejściu do podciągu założyć
$u(t_n)\rightharpoonup u_\infty$ w $H^1$ i
$u(t_n)\to u_\infty$ w $L^2$.
Ponieważ $|u(t_n)|\le M$ i $s\mapsto s-s^3$
jest lipschitzowska na $[-M,M]$, mamy
$f(u(t_n))\to f(u_\infty)$ w $L^2$.
Przechodząc do granicy w równaniu testowanym przez
$\varphi\in H^1(\T_L^d)$ otrzymujemy tezę.
\end{proof}

\section{Układy i niestabilność Turinga}
Dla wektorowej funkcji $U=(u_1,u_2)^{\mathsf T}$
rozważamy $U_t=\mathcal D\Delta U+F(U)$,
$\mathcal D=\operatorname{diag}(D_1,D_2)$, $D_i>0$.
Stan jednorodny $U_*$ spełnia $F(U_*)=0$.
Linearyzacja wokół tego stanu ma postać
$V_t=\mathcal D\Delta V+JV$, $J=DF(U_*)$.
W odróżnieniu od przypadku skalarnego dyfuzja może
wywołać niestabilność stanu stabilnego dla samego ODE.


\begin{twierdzenie}[Kryterium liniowej niestabilności Turinga]
\label{thm:pde27-turing-general}
Niech $J=\begin{psmallmatrix}a&b\\c&d\end{psmallmatrix}$
będzie rzeczywistą macierzą $2\times2$ oraz
$\mathcal D=\operatorname{diag}(D_1,D_2)$, $D_1,D_2>0$.
Załóżmy $\operatorname{tr}J<0$ i $\det J>0$,
tak że jednorodny układ liniowy $V_t=JV$ jest
asymptotycznie stabilny. Dla $\mu>0$ macierz
$J-\mu\mathcal D$ ma wartość własną o dodatniej
części rzeczywistej wtedy i tylko wtedy, gdy
\begin{equation}\label{eq:pde27-turing-poly}
 D_1D_2\mu^2-(D_2a+D_1d)\mu+\det J<0.
\end{equation}
Takie $\mu>0$ istnieje wtedy i tylko wtedy, gdy
\[
 S:=D_2a+D_1d>0,\qquad S^2>4D_1D_2\det J.
\]
Na konkretnym ograniczonym obszarze z jednorodnym
warunkiem Neumanna niestabilność występuje dokładnie
wtedy, gdy co najmniej jedna dodatnia wartość własna
$\mu_k$ operatora $-\Delta$ spełnia
\eqref{eq:pde27-turing-poly}; warunki na $S$ i wyróżnik
nie wystarczają bez tego warunku spektralnego.
\end{twierdzenie}
\begin{proof}
Na modzie Laplasjanu o wartości własnej $-\mu$
ewolucja liniowa ma macierz $M_\mu=J-\mu\mathcal D$.
Z hipotez wynika
$\operatorname{tr}M_\mu=\operatorname{tr}J-
\mu(D_1+D_2)<0$. Dla rzeczywistej macierzy $2\times2$
o ujemnym śladzie dodatnią część rzeczywistą
wartości własnej otrzymujemy dokładnie przy
$\det M_\mu<0$: wartości własne są wtedy
rzeczywiste i mają przeciwne znaki; w pozostałych
przypadkach ich części rzeczywiste są niedodatnie.
Bezpośredni rachunek daje wielomian
\eqref{eq:pde27-turing-poly}.
Ponieważ jest to trójmian o dodatnim współczynniku
kwadratowym i dodatnim wyrazie wolnym, przyjmuje
wartość ujemną dla pewnego $\mu>0$ dokładnie gdy
$S>0$ i jego wyróżnik $S^2-4D_1D_2\det J$
jest dodatni. Na obszarze ograniczonym dopuszczalne
częstości $\mu$ są dyskretne, więc ostatni warunek
należy sprawdzić dla rzeczywistego widma Neumanna.
\end{proof}

\begin{uwaga}[Przedział niestabilnych długości fal]
Jeżeli warunki z poprzedniego twierdzenia są spełnione,
niestałe mody rosną dokładnie dla
$\mu\in(\mu_-,\mu_+)$, gdzie
$\mu_\pm=(S\pm\sqrt{S^2-4D_1D_2\det J})/(2D_1D_2)$.
W przykładzie poniżej $\mu_-=1/5$ i $\mu_+=1/2$;
wybór $\mu_1=3/10$ rzeczywiście trafia do tego przedziału.
\end{uwaga}


\begin{stwierdzenie}[Kryterium niestabilności konkretnego modu]
\label{prop:pde27-turing}
Niech $\Omega=(0,L)$, na końcach niech obowiązują
jednorodne warunki Neumanna, oraz
\[
 J=\begin{pmatrix}1&-2\\2&-3\end{pmatrix},\qquad
 \mathcal D=\begin{pmatrix}1&0\\0&10\end{pmatrix}.
\]
Macierz $J$ jest stabilna w sensie liniowym dla ODE
$\dot V=JV$: obie jej wartości własne wynoszą $-1$.
Jeżeli $L=\pi/\sqrt{3/10}$, to odpowiedni
jednomodowy układ reakcji--dyfuzji ma dodatnią
rzeczywistą wartość własną. W szczególności
zero jest liniowo niestabilne dla PDE, mimo
stabilności przestrzennie jednorodnej dynamiki.
\end{stwierdzenie}
\begin{proof}
Funkcje $\cos(n\pi x/L)$ są funkcjami własnymi
laplasjanu Neumanna z wartościami $-\mu_n$,
$\mu_n=(n\pi/L)^2$. Na takim modzie ewolucja
jest wyznaczona przez macierz $J-\mu_n\mathcal D$.
Dla $n=0$ mamy
$\operatorname{tr}J=-2$, $\det J=1$, więc
wielomian charakterystyczny jest $(\lambda+1)^2$.
Dla $n=1$, $\mu_1=3/10$, otrzymujemy
\[
 \operatorname{tr}(J-\mu_1\mathcal D)=-\frac{53}{10},
 \qquad \det(J-\mu_1\mathcal D)
 =1-7\frac3{10}+10\frac9{100}=-\frac15.
\]
Ujemny wyznacznik oznacza jedną dodatnią i jedną
ujemną rzeczywistą wartość własną. Odpowiednia
funkcja własna $V(x)=v\cos(\pi x/L)$
rośnie wykładniczo w czasie. To dowodzi twierdzenia.
\end{proof}

\begin{uwaga}[Stabilność liniowa a nieliniowa]
Powyższy rachunek dowodzi niestabilności
\emph{zlinearyzowanego} PDE. Aby wnioskować o
nieliniowej niestabilności układu
$U_t=\mathcal D\Delta U+F(U)$, trzeba dodatkowo
sprecyzować regularność $F$ oraz zastosować
twierdzenie o niestabilności ze spektralnej wartości
własnej o dodatniej części rzeczywistej.
Takiego twierdzenia nie ukrywamy w argumencie liniowym.
\end{uwaga}

\section{Podsumowanie pojęciowe}
Monotoniczność operatorów z rozdziału~\ref{chap:pde26}
prowadziła do istnienia nieliniowych rozwiązań eliptycznych.
W obecnym rozdziale kluczową rolę odgrywają
\emph{porównanie} i \emph{dynamika półgrupy}.
W skalarnej teorii KPP dyfuzja i reakcja prowadzą
do frontu o minimalnej prędkości $2\sqrt{Dr}$;
w układach z różnymi współczynnikami dyfuzji
mogą natomiast powstawać niestabilne mody przestrzenne.
Tych dwóch mechanizmów nie należy utożsamiać.

\chapter[Równania Naviera--Stokesa]{Równania Naviera--Stokesa: rozkład Helmholtza, rozwiązania Leraya--Hopfa i struktura wirów}
\label{chap:pde28}

\section{Równanie nieściśliwego płynu i wybór przestrzeni}
Niech $d\in\{2,3\}$, $\nu>0$, $T>0$ i $\T^d=(\R/2\pi\mathbb Z)^d$.
Szukamy prędkości $u:(0,T)\times\T^d\to\R^d$ i ciśnienia
$p:(0,T)\times\T^d\to\R$ spełniających
\begin{equation}\label{eq:pde28-ns}
 \partial_tu+(u\cdot\nabla)u-\nu\Delta u+\nabla p=0,
 \qquad \diver u=0,
 \qquad u(0)=u_0.
\end{equation}
Współczynnik $\nu$ jest lepkością kinematyczną, $u$ jest \emph{polem
prędkości}, natomiast ciśnienie jest mnożnikiem wymuszającym warunek
nieściśliwości. Przyjęcie torusa pozwala wyprowadzić zasadnicze
argumenty bez dodatkowych założeń o regularności brzegu. Dla
pól wektorowych symbol $\nabla u$ oznacza macierz
$(\partial_j u_i)_{i,j}$, a $\diver(u\otimes u)$ jest wektorem
o współrzędnych $\sum_j\partial_j(u_i u_j)$.

Zakładamy $\int_{\T^d}u_0\,\mathrm d\lambda_d(x)=0$; średnia
prędkości zachowuje się w czasie i można ją wyeliminować przez
przejście do układu poruszającego się ze stałą prędkością.
Dla pól o średniej zero używamy przestrzeni
\[
 H=\{v\in L^2(\T^d;\R^d):\diver v=0\text{ w }\mathcal D',\ \textstyle\int_{\T^d}v\,\mathrm d\lambda_d=0\},
 \quad V=H\cap H^1(\T^d;\R^d).
\]
Na $V$ normy $\|\nabla v\|_2$ i $\|v\|_{H^1}$ są równoważne
na mocy nierówności Poincarégo. Dalej $V'$ jest przestrzenią
dualną $V$, a $V\subset H\simeq H'\subset V'$ oznacza trójkę
przestrzeni Hilberta z gęstymi, ciągłymi zanurzeniami.

\section{Rozkład Helmholtza i rzut Leraya}
\begin{twierdzenie}[Rozkład Helmholtza na torusie]\label{thm:pde28-helmholtz}
Niech $d\ge2$ i $F\in L^2(\T^d;\R^d)$ ma średnią zero.
Istnieje dokładnie jedna para $(v,\phi)$, gdzie
$v\in H$, $\phi\in H^1(\T^d)$ ma średnią zero, taka że
$F=v+\nabla\phi$. Składniki są ortogonalne w $L^2$.
Przyporządkowanie $PF=v$ jest ortogonalnym rzutem
$L^2_0(\T^d;\R^d)\to H$, przy czym dla
$k\in\mathbb Z^d\setminus\{0\}$
\begin{equation}\label{eq:pde28-P}
 \widehat{PF}(k)=\left(I-\frac{k\otimes k}{|k|^2}\right)\widehat F(k),
 \qquad \widehat\phi(k)=-\frac{i k\cdot\widehat F(k)}{|k|^2}.
\end{equation}
\end{twierdzenie}
\begin{proof}
Przyjmujemy konwencję $\widehat F(k)=(2\pi)^{-d}
\int_{\T^d}F(x)e^{-ik\cdot x}\,\mathrm d\lambda_d(x)$.
Tak jak dla szeregów Fouriera w rozdziale~\ref{chap:pde-fourier}
współczynniki są uśredniane; normy $L^2$ użyte przy oszacowaniach
energii płynu odnoszą się jednak do pełnej miary $\lambda_d$,
bez dzielenia przez $(2\pi)^d$.
Dla każdego $k\ne0$ macierz w~\eqref{eq:pde28-P} jest rzutem
ortogonalnym na podprzestrzeń prostopadłą do $k$.
Ponieważ jej norma operatorowa nie przekracza $1$, wzór definiuje
$PF\in L^2$. Mamy $k\cdot\widehat{PF}(k)=0$, zatem
$\diver PF=0$, a średnia $PF$ jest zerowa.
Z drugiej strony $ik\widehat\phi(k)=\widehat F(k)-\widehat{PF}(k)$.
Nierówność $|k||\widehat\phi(k)|\le|\widehat F(k)|$
i twierdzenie Parsevala pokazują, że $\phi\in H^1$.
Dla $v\in H$ i $\phi\in H^1$ zachodzi
$\int_{\T^d}v\cdot\nabla\phi\,\mathrm d\lambda_d=0$
na mocy definicji dywergencji dystrybucyjnej; jest to
ortogonalność składników. Jeżeli $F=v_1+\nabla\phi_1
=v_2+\nabla\phi_2$, to różnica obu rozkładów leży jednocześnie
w $H$ i w jego ortogonalnym dopełnieniu, zatem znika.
Warunek zerowej średniej wyznacza jednoznacznie $\phi$.
\end{proof}

\begin{uwaga}[Operator Stokesa]
Na $H$ określamy $A=-P\Delta$ z dziedziną
$D(A)=H^2(\T^d;\R^d)\cap H$. Ponieważ $P$ komutuje
z $\Delta$ na torusie, mamy $Av=-\Delta v$, a
$(Av,v)_H=\|\nabla v\|_2^2$. Własności bazy Fouriera
dają dodatnią samosprzężoność $A$ oraz zwartość $A^{-1}$.
Rzut $P$ nie jest mnożeniem przez stałą macierz w przestrzeni
fizycznej, lecz nielokalnym operatorem Fouriera.
\end{uwaga}

\section{Tożsamość energii i słaba postać równania}
Dla $a,b,c\in V$ określamy formę trójliniową
\[
 b(a,b,c)=\int_{\T^d}(a\cdot\nabla)b\cdot c\,\mathrm d\lambda_d(x).
\]
W wymiarach $2$ i $3$ jest ona dobrze określona na $V^3$ przez
nierówności Sobolewa i Höldera. Przy $\diver a=0$ zachodzi
$b(a,b,c)=-b(a,c,b)$, a w szczególności $b(a,b,b)=0$.
Ta ostatnia równość jest podstawowym mechanizmem dyssypacji energii.

\begin{definicja}[Rozwiązanie Leraya--Hopfa na torusie]\label{def:pde28-leray}
Niech $u_0\in H$, $\nu>0$ i $T>0$. Pole $u$ nazywamy
rozwiązaniem Leraya--Hopfa zagadnienia~\eqref{eq:pde28-ns} na $[0,T]$,
jeżeli $u\in L^\infty(0,T;H)\cap L^2(0,T;V)$,
$u\in C([0,T];H_{\mathrm w})$, $u(0)=u_0$ oraz dla każdego
$\varphi\in V$ i $t\in[0,T]$ spełnia tożsamość
\begin{equation}\label{eq:pde28-weak}
 (u(t),\varphi)_H-(u_0,\varphi)_H
 +\nu\int_0^t(\nabla u(s),\nabla\varphi)_{L^2}\,\mathrm ds
 +\int_0^t b(u(s),u(s),\varphi)\,\mathrm ds=0.
\end{equation}
Zapis $C([0,T];H_{\mathrm w})$ oznacza ciągłość wszystkich
funkcji skalarnych $t\mapsto(u(t),h)_H$, $h\in H$.
Wymagamy ponadto nierówności energii
\begin{equation}\label{eq:pde28-energy}
 \frac12\|u(t)\|_2^2+\nu\int_0^t\|\nabla u(s)\|_2^2\,\mathrm ds
 \le\frac12\|u_0\|_2^2\qquad (0\le t\le T).
\end{equation}
\end{definicja}

\begin{lemat}[Oszacowanie nieliniowości w przestrzeni dualnej]
\label{lem:pde28-B-est}
Niech $B(u,v)\in V'$ będzie określony przez
$\langle B(u,v),\varphi\rangle=b(u,v,\varphi)$.
Dla $u\in V$ zachodzą oszacowania
\[
 \|B(u,u)\|_{V'}\le
 \begin{cases}
 C\|u\|_2\|\nabla u\|_2,&d=2,\\
 C\|u\|_2^{1/2}\|\nabla u\|_2^{3/2},&d=3.
 \end{cases}
\]
Stała $C$ zależy tylko od wymiaru i wielkości torusa.
\end{lemat}
\begin{proof}
Wykorzystujemy antysymetrię w dwóch ostatnich zmiennych:
$|b(u,u,\varphi)|=|b(u,\varphi,u)|
\le\|u\|_4^2\|\nabla\varphi\|_2$.
Nierówność Gagliarda--Nirenberga na torusie, wraz z zerową średnią,
daje $\|u\|_4^2\le C\|u\|_2\|\nabla u\|_2$ w wymiarze $2$
oraz $\|u\|_4^2\le C\|u\|_2^{1/2}
\|\nabla u\|_2^{3/2}$ w wymiarze $3$.
Po podzieleniu przez $\|\varphi\|_V=\|\nabla\varphi\|_2$
i wzięciu supremum otrzymujemy tezę.
\end{proof}

\section{Globalne istnienie rozwiązania Leraya--Hopfa}
\begin{twierdzenie}[Istnienie słabego rozwiązania w wymiarach $2$ i $3$]
\label{thm:pde28-existence}
Niech $d\in\{2,3\}$, $\nu>0$, $T>0$ i $u_0\in H$.
Istnieje rozwiązanie Leraya--Hopfa zagadnienia~\eqref{eq:pde28-ns}
na $[0,T]$. Można je wybrać tak, że
$\partial_tu\in L^{4/3}(0,T;V')$ dla $d=3$,
a $\partial_tu\in L^2(0,T;V')$ dla $d=2$.
Ponadto $u(t)\to u_0$ silnie w $H$ przy $t\downarrow0$.
\end{twierdzenie}
\begin{proof}
\textbf{Krok 1: układy Galerkina.}
Niech $(w_j)_{j\ge1}$ będzie ortonormalną bazą $H$
złożoną z gładkich, rzeczywistych, bezdywergencyjnych
funkcji własnych operatora Stokesa, z uwzględnieniem krotności.
Oznaczmy $H_m=\operatorname{span}\{w_1,\dots,w_m\}$
i niech $P_m$ będzie rzutem ortogonalnym na $H_m$.
Szukamy $u_m:[0,T]\to H_m$ o współrzędnych
spełniających dla $1\le j\le m$
\begin{equation}\label{eq:pde28-galerkin}
 (u_m',w_j)_H+\nu(\nabla u_m,\nabla w_j)_{L^2}
 +b(u_m,u_m,w_j)=0,\quad u_m(0)=P_m u_0.
\end{equation}
Jest to skończony układ ODE o gładkiej prawej stronie;
przynajmniej lokalnie ma jednoznaczne rozwiązanie.

\textbf{Krok 2: energia i istnienie globalne przybliżeń.}
Mnożymy $j$-te równanie przez współrzędną $u_m$ przy $w_j$
i sumujemy. Wykorzystując $b(u_m,u_m,u_m)=0$, otrzymujemy
\begin{equation}\label{eq:pde28-gal-energy}
 \frac12\|u_m(t)\|_2^2+\nu\int_0^t\|\nabla u_m(s)\|_2^2
 \,\mathrm ds=\frac12\|P_m u_0\|_2^2\le\frac12\|u_0\|_2^2.
\end{equation}
Norma $u_m$ nie może eksplodować w skończonym czasie,
a skończony układ ODE można przedłużyć na dowolny $[0,T]$.
Oszacowania są jednostajne względem $m$.

\textbf{Krok 3: pochodne czasowe.}
Rzut $P_m$ jest kontrakcją na $V$ i $V'$:
bazę wybrano z funkcji własnych operatora Stokesa,
które są ortogonalne również dla iloczynu skalarnego
$(\nabla v,\nabla w)_{L^2}$. Wobec tego
$\|P_m\varphi\|_V\le\|\varphi\|_V$;
następnie z definicji normy dualnej
$\|P_m F\|_{V'}=\sup_{\|\varphi\|_V\le1}
|\langle F,P_m\varphi\rangle|\le\|F\|_{V'}$.
Z równania Galerkina oraz lematu~\ref{lem:pde28-B-est} mamy
\[
 \|u_m'\|_{V'}\le\nu\|\nabla u_m\|_2+
 \begin{cases}
 C\|u_m\|_2\|\nabla u_m\|_2,&d=2,\\
 C\|u_m\|_2^{1/2}\|\nabla u_m\|_2^{3/2},&d=3.
 \end{cases}
\]
Prawa strona jest ograniczona odpowiednio w
$L^2(0,T)$ i $L^{4/3}(0,T)$; dla pierwszego składnika
w drugim przypadku korzystamy z $T<\infty$.

\textbf{Krok 4: zwartość w przestrzeni i czasie.}
Z oszacowań wynika, po wybraniu podciągu,
$u_m\stackrel{*}{\rightharpoonup}u$ w $L^\infty(0,T;H)$,
$u_m\rightharpoonup u$ w $L^2(0,T;V)$ i
$u_m'\rightharpoonup u_t$ we wskazanej przestrzeni dualnej.
Zwartość $V\hookrightarrow H$ i ciągłość $H\hookrightarrow V'$
pozwalają zastosować twierdzenie Aubina--Lionsa
(zob. rozdział~\ref{chap:pde20}):
\begin{equation}\label{eq:pde28-strong}
 u_m\longrightarrow u\quad\text{silnie w }L^2(0,T;H).
\end{equation}
W szczególności $u_m\otimes u_m\to u\otimes u$ w
$L^1((0,T)\times\T^d)$, gdyż
$\|a\otimes a-b\otimes b\|_1
\le(\|a\|_2+\|b\|_2)\|a-b\|_2$ na czasoprzestrzeni.

\textbf{Krok 5: przejście do granicy w równaniu.}
Niech $\varphi\in\bigcup_{j\ge1}H_j$ i
$\eta\in C_c^\infty(0,T)$. Testujemy~\eqref{eq:pde28-galerkin}
przez $\eta(t)\varphi$ i całkujemy w czasie.
Człon liniowy przechodzi do granicy słabo, a człon nieliniowy
po całkowaniu przez części przestrzennym ma postać
$-\int_0^T\int_{\T^d}(u_m\otimes u_m):\nabla\varphi\,\mathrm d\lambda_d(x)\,\mathrm dt$,
więc przechodzi do granicy na mocy~\eqref{eq:pde28-strong}.
Gęstość sumy $H_j$ w $V$ i jednostajne oszacowanie
$B(u_m,u_m)$ w przestrzeni $L^{4/3}(0,T;V')$
($L^2$ dla $d=2$) rozszerzają tożsamość do wszystkich
$\varphi\in V$ i dają~\eqref{eq:pde28-weak}.

\textbf{Krok 6: warunek początkowy i nierówność energii.}
Dla każdego $\varphi\in V$ funkcje
$t\mapsto(u_m(t),\varphi)_H$ są jednostajnie ograniczone
i jednostajnie Hölderowskie w czasie: z kroku 3,
$|(u_m(t)-u_m(s),\varphi)|\le
C\|\varphi\|_V|t-s|^{1-1/q}$, gdzie $q=4/3$ lub $2$.
Diagonale dla przeliczalnego gęstego zbioru testów,
wraz z jednostajnym ograniczeniem w $H$,
dają $u_m(t)\rightharpoonup u(t)$ w $H$ dla
\emph{każdego} $t\in[0,T]$, a także
$u\in C([0,T];H_{\mathrm w})$.
Ponieważ $u_m(0)=P_m u_0\to u_0$ w $H$,
warunek początkowy jest zachowany.
Dla ustalonego $t$ stosujemy słabą półciągłość dolną
normy $\|u_m(t)\|_2$ i normy
$\|\nabla u_m\|_{L^2(0,t;L^2)}$ do~\eqref{eq:pde28-gal-energy},
otrzymując~\eqref{eq:pde28-energy} dla każdego $t$.
Wreszcie $u(t)\rightharpoonup u_0$ w $H$ przy $t\downarrow0$,
a nierówność energii daje
$\limsup_{t\downarrow0}\|u(t)\|_2\le\|u_0\|_2$.
Słaba zbieżność wraz ze zbieżnością norm dowodzi
$u(t)\to u_0$ silnie w $H$.
\end{proof}


\begin{uwaga}[Zwartość paraboliczna a nieliniowość Naviera--Stokesa]
\label{rem:pde-v44-ns-compactness}
Zwartość z kroku czwartego dowodu twierdzenia
\ref{thm:pde28-existence} ma dokładnie tę samą konstrukcję
funkcjonalną co w rozdziale~\ref{chap:pde20}:
$V\hookrightarrow H$ zwarte i $H\hookrightarrow V'$ ciągłe,
$u_m$ ograniczone w $L^2(0,T;V)$ oraz $\partial_tu_m$
ograniczone w $L^q(0,T;V')$, gdzie $q=2$ dla $d=2$
i $q=4/3$ dla $d=3$. Sama słaba zbieżność w przestrzeni
energii nie wystarczałaby do przejścia do granicy
w $u_m\otimes u_m$. Silna zbieżność w $L^2(0,T;H)$
zapewnia bowiem bezpośrednio
\[
 \|u_m\otimes u_m-u\otimes u\|_{L^1((0,T)\times\T^d)}
 \le (\|u_m\|_{L^2}+\|u\|_{L^2})\|u_m-u\|_{L^2}\longrightarrow0.
\]
Nie wnioskujemy stąd o silnej zbieżności gradientów ani
o równości energii w trzech wymiarach.
\end{uwaga}


\begin{uwaga}[Ciśnienie rozwiązania słabego]
Dla gładkiego rozwiązania warunek $\diver u=0$
i równanie~\eqref{eq:pde28-ns} dają
$-\Delta p=\sum_{i,j}\partial_i\partial_j(u_i u_j)$.
Na torusie wybieramy normalizację
$\int_{\T^d}p(t,x)\,\mathrm d\lambda_d(x)=0$.
Dla rozwiązania Leraya--Hopfa wzór ten określa $p$
w sensie dystrybucji przestrzennych i czasowych;
nie należy bez dodatkowych założeń traktować go
jako tożsamości klasycznych funkcji. Dla ustalonego czasu w sensie dystrybucji
przestrzennych rozwiązujemy równanie Poissona dla
$p$ o średniej zero; w bazie Fouriera, dla $k\ne0$,
$\widehat p(k)=-|k|^{-2}\sum_{i,j}k_i k_j
\widehat{(u_i u_j)}(k)$.
Współczynniki iloczynów $u_i u_j$ są dobrze określone,
gdy $u(t)\in L^2$. Stąd $p$ jest dobrze określoną
dystrybucją z dokładnością do funkcji samego czasu.
Z projekcji równania przez $P$ wynika
$\nabla p=-(I-P)\diver(u\otimes u)$.
\end{uwaga}


\begin{twierdzenie}[Równość energii i ciągłość silna w wymiarze dwa]
\label{thm:pde28-energy-equality-2d}
Niech $\nu>0$, $T>0$, $d=2$ i $u_0\in H$.
Każde rozwiązanie Leraya--Hopfa $u$ na $[0,T]$
należy do $C([0,T];H)$ i spełnia dla wszystkich
$t\in[0,T]$ dokładną równość
\begin{equation}\label{eq:pde28-energy-equality-2d}
 \frac12\|u(t)\|_2^2+
 \nu\int_0^t\|\nabla u(s)\|_2^2\,\mathrm ds
 =\frac12\|u_0\|_2^2.
\end{equation}
Ponadto, ponieważ na torusie $2\pi$ funkcje w $H$
mają zerową średnią,
$\|u(t)\|_2^2\le e^{-2\nu t}\|u_0\|_2^2$.
\end{twierdzenie}
\begin{proof}
Z lematu~\ref{lem:pde28-B-est} i warunków
$u\in L^\infty(0,T;H)\cap L^2(0,T;V)$ wynika
$\|B(u,u)\|_{V'}\le C\|u\|_2\|\nabla u\|_2\in
L^2(0,T)$. Niech $\mathcal A:V\to V'$ oznacza rozszerzenie
operatora Stokesa określone przez
$\langle\mathcal Au,\varphi\rangle=
(\nabla u,\nabla\varphi)_{L^2}$.
Wówczas $\|\mathcal Au\|_{V'}\le\|u\|_V$,
a równość słaba daje
$u_t=-\nu\mathcal Au-B(u,u)\in L^2(0,T;V')$.
W trójce Hilberta $V\subset H\subset V'$ mamy
zatem $u\in L^2(0,T;V)$ oraz $u_t\in L^2(0,T;V')$.
Lemat o regule łańcuchowej w trójce Hilberta,
omówiony w rozdziale~\ref{chap:pde20}, zapewnia
$u\in C([0,T];H)$ i
$\frac12\frac{\mathrm d}{\mathrm dt}\|u(t)\|_2^2
=\langle u_t(t),u(t)\rangle$ prawie wszędzie.
Testujemy równanie słabe przez $u(t)\in V$.
Dla takich chwil $b(u,u,u)=0$ na mocy
antysymetrii i gęstości gładkich pól bezdywergencyjnych.
Otrzymujemy pochodną
$\frac12(\|u\|_2^2)'=-\nu\|\nabla u\|_2^2$;
całkowanie w czasie daje~\eqref{eq:pde28-energy-equality-2d}.

Na torusie $\T^2$ najmniejsza dodatnia wartość
własna $-\Delta$ wynosi $1$, więc nierówność
Poincarégo ma postać $\|u\|_2\le\|\nabla u\|_2$.
Dla $E(t)=\|u(t)\|_2^2$ otrzymujemy
$E'(t)\le-2\nu E(t)$ prawie wszędzie.
Całkowanie tej nierówności daje wskazany zanik.
\end{proof}

\begin{wniosek}[Wykładniczy zanik energii rozwiązania Galerkina w 3D]
\label{cor:pde28-decay-3d}
Niech $\nu>0$, $T>0$, $u_0\in H$ i $d=3$.
Rozwiązanie Leraya--Hopfa otrzymane w dowodzie
twierdzenia~\ref{thm:pde28-existence} można wybrać tak, że
\[
 \|u(t)\|_2^2\le e^{-2\nu t}\|u_0\|_2^2,
 \qquad 0\le t\le T.
\]
Nie jest tu zakładana gładkość ani jednoznaczność
rozwiązań trójwymiarowych.
\end{wniosek}
\begin{proof}
Dla przybliżeń Galerkina z~\eqref{eq:pde28-gal-energy}
różniczkujemy energię i używamy
$\|\nabla u_m\|_2^2\ge\|u_m\|_2^2$:
$\frac{\mathrm d}{\mathrm dt}\|u_m(t)\|_2^2
\le-2\nu\|u_m(t)\|_2^2$.
Stąd $\|u_m(t)\|_2^2\le
 e^{-2\nu t}\|P_mu_0\|_2^2\le e^{-2\nu t}\|u_0\|_2^2$.
W kroku 6 dowodu~\ref{thm:pde28-existence} wybrano
jeden podciąg zbieżny słabo w $H$ dla każdej chwili
$t\in[0,T]$. Słaba półciągłość dolna normy przy
ustalonym $t$ daje pożądane oszacowanie dla tej granicy.
W szczególności w tym wniosku nie korzystamy z
nieudowodnionej równości energii dla wszystkich
rozwiązań Leraya--Hopfa w wymiarze trzy.
\end{proof}


\section{Wymiar dwa: jednoznaczność i wirowość}
\begin{twierdzenie}[Jednoznaczność słabych rozwiązań w wymiarze dwa]
\label{thm:pde28-unique2d}
Niech $\nu>0$, $T>0$ i $u_0\in H$ dla $d=2$.
Rozwiązanie Leraya--Hopfa na $[0,T]$ jest jednoznaczne.
Jeśli $u,v$ odpowiadają danym $u_0,v_0\in H$, to
\[
 \|u(t)-v(t)\|_2^2\le\|u_0-v_0\|_2^2
 \exp\left(\frac C\nu\int_0^t\|\nabla v(s)\|_2^2\,\mathrm ds\right),
 \qquad 0\le t\le T.
\]
\end{twierdzenie}
\begin{proof}
Z kroku 3 poprzedniego dowodu wynika
$u_t,v_t\in L^2(0,T;V')$, a więc tożsamość energii
w trójce $V\subset H\subset V'$ stosuje się do
$w=u-v$ (por. rozdział~\ref{chap:pde20}).
Odejmując równania i wykorzystując
$b(u,w,w)=0$, otrzymujemy
\[
 \frac12\frac{\mathrm d}{\mathrm dt}\|w\|_2^2
 +\nu\|\nabla w\|_2^2=-b(w,v,w).
\]
Wymiar dwa pozwala użyć nierówności Ladyżenskiej:
$|b(w,v,w)|\le\|w\|_4^2\|\nabla v\|_2
\le C\|w\|_2\|\nabla w\|_2\|\nabla v\|_2$.
Z nierówności Younga dostajemy
$\frac{\mathrm d}{\mathrm dt}\|w\|_2^2
\le (C/\nu)\|\nabla v\|_2^2\|w\|_2^2$.
Ponieważ $\nabla v\in L^2(0,T;L^2)$,
lemat Gronwalla daje podane oszacowanie, a dla równych
danych otrzymujemy $u=v$.
\end{proof}

\begin{stwierdzenie}[Równanie wirowości i zanik enstrofii]
\label{prop:pde28-vorticity}
Niech $d=2$, a $(u,p)$ będzie gładkim rozwiązaniem
\eqref{eq:pde28-ns} na $[0,T]\times\T^2$.
Przyjmijmy $\omega=\partial_1u_2-\partial_2u_1$
oraz $\mathcal E_\omega(t)=\frac12\|\omega(t)\|_2^2$.
Wówczas
\[
 \partial_t\omega+u\cdot\nabla\omega=\nu\Delta\omega,
 \qquad \mathcal E_\omega(t)+\nu\int_0^t\|\nabla\omega(s)\|_2^2
 \,\mathrm ds=\mathcal E_\omega(0).
\]
Dodatkowo dla bezdywergencyjnych pól o zerowej średniej
$\|\nabla u(t)\|_2=\|\omega(t)\|_2$.
\end{stwierdzenie}
\begin{proof}
Stosując operator $\partial_1$ do drugiego równania
wektorowego i odejmując $\partial_2$ pierwszego,
zwijamy gradient ciśnienia. Rozwijając pochodne
składników konwekcyjnych, korzystamy z
$\partial_1u_1+\partial_2u_2=0$;
pozostaje $u\cdot\nabla\omega$.
Mnożymy równanie przez $\omega$ i całkujemy po torusie.
Składnik konwekcyjny znika:
$\int_{\T^2}u\cdot\nabla(\omega^2/2)\,\mathrm d\lambda_2(x)=0$.
Równość norm wynika z tożsamości
$\int_{\T^2}|\nabla u|^2\,\mathrm d\lambda_2(x)=\int_{\T^2}(|\diver u|^2+|\operatorname{curl}u|^2)\,\mathrm d\lambda_2(x)$,
którą otrzymujemy przez całkowanie przez części.
\end{proof}

\begin{twierdzenie}[Globalne oszacowanie $H^1$ w wymiarze dwa]
\label{thm:pde28-H1-2d}
Niech $d=2$, $\nu>0$, $T>0$ oraz $u_0\in V$.
Jedyne rozwiązanie Leraya--Hopfa $u$ na $[0,T]$ spełnia
$u\in L^\infty(0,T;V)\cap L^2(0,T;H^2(\T^2;\R^2))$
oraz dla prawie każdego $t\in(0,T)$
\[
 \|\nabla u(t)\|_2^2+2\nu\int_0^t\|\Delta u(s)\|_2^2
 \,\mathrm ds\le\|\nabla u_0\|_2^2.
\]
\end{twierdzenie}
\begin{proof}
Wybieramy w dowodzie twierdzenia~\ref{thm:pde28-existence}
przestrzenie Galerkina złożone ze wszystkich bezdywergencyjnych
modów Fouriera $0<|k|\le N$. Dla $\omega_N=\partial_1u_{N,2}
-\partial_2u_{N,1}$ równanie rzutowane ma postać
$\partial_t\omega_N+Q_N(u_N\cdot\nabla\omega_N)
=\nu\Delta\omega_N$, gdzie $Q_N$ jest ortogonalnym
rzutem na skalarne mody z $0<|k|\le N$.
Ponieważ $Q_N\omega_N=\omega_N$, mnożenie przez
$\omega_N$ daje
\[
 \frac12\frac{\mathrm d}{\mathrm dt}\|\omega_N\|_2^2
 +\nu\|\nabla\omega_N\|_2^2=0.
\]
Przy tym człon transportowy znika przez $\diver u_N=0$.
Dla bezdywergencyjnego pola o zerowej średniej zachodzą
$\|\omega_N\|_2=\|\nabla u_N\|_2$ i
$\|\nabla\omega_N\|_2=\|\Delta u_N\|_2$,
co można sprawdzić oddzielnie na każdym modzie Fouriera.
Otrzymujemy zatem jednostajne ograniczenia w
$L^\infty(0,T;V)$ i $L^2(0,T;H^2)$.
Słaba zwartość i słaba półciągłość dolna norm
dają tezę dla granicy Galerkina. Na mocy
jednoznaczności z twierdzenia~\ref{thm:pde28-unique2d}
jest to wybrane rozwiązanie $u$.
\end{proof}

\begin{uwaga}[Dlaczego wymiar trzy jest inny]
W wymiarze $3$ wirowość $\boldsymbol\omega=\nabla\times u$
spełnia dla gładkiego rozwiązania
$\partial_t\boldsymbol\omega+(u\cdot\nabla)\boldsymbol\omega
=(\boldsymbol\omega\cdot\nabla)u+\nu\Delta\boldsymbol\omega$.
Pojawia się człon \emph{rozciągania wirów}, który nie ma
odpowiednika w wymiarze $2$. Z tożsamości energii $L^2$
nie otrzymuje się wprost ograniczenia normy $\|\nabla u\|_2$
we wszystkich czasach w wymiarze $3$.
\end{uwaga}

\section{Wymiar trzy: kryterium warunkowej jednoznaczności}
\begin{twierdzenie}[Jednoznaczność przy dodatkowej całkowalności]
\label{thm:pde28-serrin}
Niech $d=3$, $\nu>0$, $T>0$ i $u_0\in H$.
Niech $u,v$ będą rozwiązaniami Leraya--Hopfa na $[0,T]$
o tych samych danych początkowych. Załóżmy dodatkowo, że
$v\in L^\infty(0,T;V)\cap L^2(0,T;D(A))$,
$\partial_tv\in L^2(0,T;H)$ oraz
$v\in L^q(0,T;L^p(\T^3))$, gdzie $p>3$ i
$2/q+3/p\le1$. Wówczas $u=v$ prawie wszędzie na
$(0,T)\times\T^3$.
\end{twierdzenie}
\begin{proof}
Podane silniejsze założenia regularności $v$ pozwalają
uzasadnić tożsamość mieszaną dla iloczynu $(u(t),v(t))_H$.
Istotnie, zastępujemy $v$ skończonym rzutem spektralnym
$P_nv$ i wygładzamy go w czasie. W równaniu dla $u$
możemy używać takiego pola jako testu; przy przejściu do
granicy $P_nv\to v$ silnie w $L^4(0,T;V)$,
ponieważ $v\in L^\infty(0,T;V)$, a
$\partial_tu\in L^{4/3}(0,T;V')$.
Z kolei $\partial_tv\in L^2(0,T;H)$
pozwala przejść do granicy w pochodnej iloczynu.
Pozostałe składniki przechodzą przez nierówności
Höldera i oszacowanie z lematu~\ref{lem:pde28-B-est}.
Łącząc otrzymaną tożsamość mieszaną z nierównością
energii dla $u$ i równością energii dla $v$,
uzyskujemy dla $w=u-v$ nierówność
\[
 \frac12\|w(t)\|_2^2+\nu\int_0^t\|\nabla w\|_2^2\,\mathrm ds
 \le-\int_0^t b(w(s),v(s),w(s))\,\mathrm ds.
\]
To jest nierówność względnej energii potrzebna dalej.
Z nierówności Höldera i interpolacji
$L^2\!\!\;\text{--}\!L^6$ mamy
\[
 |b(w,w,v)|\le\|v\|_p\|w\|_{2p/(p-2)}\|\nabla w\|_2
 \le C\|v\|_p\|w\|_2^{1-3/p}
 \|\nabla w\|_2^{1+3/p}.
\]
Nierówność Younga daje, przy $r=2p/(p-3)$,
$|b(w,w,v)|\le(\nu/2)\|\nabla w\|_2^2
+C_{\nu,p}\|v\|_p^r\|w\|_2^2$.
Warunek $2/q+3/p\le1$ zapewnia $q\ge r$,
więc $\|v\|_p^r$ jest całkowalna w czasie.
Z lematu Gronwalla i $w(0)=0$ otrzymujemy $w=0$.
\end{proof}

\begin{uwaga}[Co pozostaje otwarte]
Dla $d=3$ twierdzenie~\ref{thm:pde28-existence} gwarantuje
istnienie globalnych rozwiązań słabych, ale nie dowodzi
ich jednoznaczności ani gładkości dla dowolnych danych w $H$.
Dowód twierdzenia~\ref{thm:pde28-serrin} jest dowodem
jednoznaczności słabego rozwiązania w obecności rozwiązania
o silniejszej regularności. Usunięcie tych dodatkowych założeń
i pozostawienie jedynie warunku $v\in L^q_tL^p_x$
daje znacznie silniejsze twierdzenie Prodiego--Serrina,
którego pełnego dowodu tutaj nie podajemy.
Argument wymiaru dwa wykorzystuje całkowalność
$\|\nabla v\|_2^2$ i nierówność Ladyżenskiej, której
odpowiednik w wymiarze trzy daje zbyt wysoką potęgę
$\|\nabla v\|_2$. Przejście od warunkowej jednoznaczności
do pełnej wymagałoby wykazania dodatkowej regularności;
nie należy tego zakładać bez dowodu.
\end{uwaga}

\section{Lokalne oszacowanie silnych rozwiązań i skala problemu}
\begin{stwierdzenie}[Oszacowanie normy $H^1$ w wymiarze trzy]
\label{prop:pde28-H1}
Niech $d=3$, $\nu>0$ i $(u,p)$ będzie gładkim rozwiązaniem
\eqref{eq:pde28-ns} na $[0,T]\times\T^3$, o średniej $u$ równej zero.
Istnieje $C>0$, zależna tylko od torusa, taka że
\[
 \frac{\mathrm d}{\mathrm dt}\|\nabla u(t)\|_2^2
 +\nu\|\Delta u(t)\|_2^2
 \le C\nu^{-3}\|\nabla u(t)\|_2^6.
\]
\end{stwierdzenie}
\begin{proof}
Iloczyn skalarny równania z $-\Delta u$ daje
$\frac12\frac{\mathrm d}{\mathrm dt}\|\nabla u\|_2^2
+\nu\|\Delta u\|_2^2
\le\|u\|_6\|\nabla u\|_3\|\Delta u\|_2$.
Z nierówności Sobolewa i Gagliarda--Nirenberga mamy
$\|u\|_6\le C\|\nabla u\|_2$ i
$\|\nabla u\|_3\le C\|\nabla u\|_2^{1/2}
\|\Delta u\|_2^{1/2}$.
Stąd prawa strona jest ograniczona przez
$C\|\nabla u\|_2^{3/2}\|\Delta u\|_2^{3/2}$.
Nierówność Younga z wykładnikami $4/3$ i $4$
pozwala pochłonąć połowę składnika
$\nu\|\Delta u\|_2^2$ i pozostawia
$C\nu^{-3}\|\nabla u\|_2^6$; pomnożenie przez $2$
i zmianę stałej kończy dowód.
\end{proof}

\begin{kontrprzyklad}[Dlaczego sama energia $L^2$ nie wystarcza]
Funkcje $u_n(x)=\sin(nx_1)e_2$ na $\T^3$, gdzie
$e_2=(0,1,0)$, mają zerową dywergencję i średnią.
Ich normy $\|u_n\|_2$ są jednakowo ograniczone,
podczas gdy $\|\nabla u_n\|_2\to\infty$.
Wobec tego kontrola samej energii $L^2$ nie daje
jednostajnego oszacowania $H^1$ w pojedynczych chwilach.
Nie jest to kontrprzykład dla regularności rozwiązań
równań Naviera--Stokesa, lecz przykład braku
odpowiedniego wniosku funkcjonalno-analitycznego.
\end{kontrprzyklad}

\section{Dokładne rozwiązanie: wir Taylora--Greena}
\begin{przyklad}[Wir Taylora--Greena na torusie dwuwymiarowym]
\label{ex:pde28-TG}
Niech $\nu>0$ i $(x,y)\in\T^2$. Funkcje
\begin{align*}
 u(t,x,y)&=e^{-2\nu t}(\sin x\cos y,-\cos x\sin y),\\
 p(t,x,y)&=\tfrac14e^{-4\nu t}(\cos 2x+\cos 2y)
\end{align*}
rozwiązują~\eqref{eq:pde28-ns}. Rzeczywiście,
$\diver u=0$, $\Delta u=-2u$, a
$(u\cdot\nabla)u=\tfrac12e^{-4\nu t}(\sin2x,\sin2y)
=-\nabla p$. Ponadto
\[
 \frac12\|u(t)\|_2^2=\pi^2e^{-4\nu t},\qquad
 \|\nabla u(t)\|_2^2=4\pi^2e^{-4\nu t},\qquad
 \omega(t,x,y)=2e^{-2\nu t}\sin x\sin y.
\]
Stąd tożsamość energii jest dokładna i pozwala
sprawdzić schematy numeryczne niezależnie od siatki.
\end{przyklad}


\begin{przyklad}[Trójwymiarowy wir Taylora--Greena: dane początkowe]
\label{ex:pde28-TG3D}
Na torusie $\T^3=([0,2\pi]/\sim)^3$ rozważmy
bezdywergencyjne pole początkowe
\[
 u_0(x,y,z)=(\sin x\cos y\cos z,\,-\cos x\sin y\cos z,\,0).
\]
Istotnie, $\partial_x u_{0,1}+\partial_y u_{0,2}=0$.
Jego wirowość jest \emph{wektorem}
\[
 \nabla\times u_0=(-\cos x\sin y\sin z,\,
 -\sin x\cos y\sin z,\,2\sin x\sin y\cos z).
\]
W szczególności trzecia składowa wirowości wynosi
$\omega_{0,z}=2\sin x\sin y\cos z$. Energia początkowa
względem miary Lebesgue'a na sześcianie okresowym wynosi
\[
 E(0)=\frac12\int_{[0,2\pi]^3}|u_0(x,y,z)|^2
 \,\mathrm d\lambda_3(x,y,z)=\pi^3.
\]
Ponieważ $u_{0,3}=0$, \emph{linie prądu pola początkowego}
leżą w płaszczyznach $z=\mathrm{const}$. Dla ustalonego
$z$ są liniami poziomicowymi funkcji prądu
$\psi_z(x,y)=\sin x\sin y\cos z$. Nie twierdzimy,
że trzecia składowa pozostanie zerowa podczas ewolucji
trójwymiarowych równań Naviera--Stokesa: jest to warunek
początkowy, nie jawny wzór rozwiązania w czasie.

\begin{center}
\includegraphics[width=.91\textwidth]{figures/taylor_green_3d_slices.png}

{\small\emph{Rysunek:} Obliczone z podanego pola przekroje
$z=\pi/4,3\pi/4,5\pi/4$: poziome strzałki wskazują kierunek
przepływu, barwy oznaczają podpisaną składową $\omega_{0,z}$
wirowości, a krzywe są poziomicami $\psi_z$. Obliczenie można
powtórzyć za pomocą pliku \texttt{scripts/taylor\_green\_3d\_plot.py};
interaktywny widok Plotly jest dołączony jako HTML.}
\end{center}

\begin{center}
\begin{minipage}{.90\textwidth}
\centering
\includegraphics[width=.90\linewidth]{figures/taylor_green_3d_schematic.png}

{\small\emph{Ilustracja schematyczna:} Przestrzenna wizualizacja
wirów Taylora--Greena. Jest to obraz koncepcyjny,
\emph{nie} wynik całkowania trajektorii powyższego pola
$u_0$ ani poprawna mapa podpisanej wielkości całej wirowości.
Fizycznie wierną wizualizacją danych początkowych jest
poprzedni rysunek.}
\end{minipage}
\end{center}
\end{przyklad}

\section{Porównanie zagadnień brzegowych i dalsze kierunki}
Na ograniczonym obszarze Lipschitzowskim $\Omega\subset\R^d$
z warunkiem przylegania $u|_{\partial\Omega}=0$
przestrzenie $H,V$ tworzy się przez domknięcie
gładkich bezdywergencyjnych pól o zwartym nośniku.
Odzyskiwanie ciśnienia wymaga wówczas analizy warunków
brzegowych, a operator Stokesa nie sprowadza się do
prostego mnożnika Fouriera. Dlatego powyższe dowody
są sformułowane na torusie; ich bezpośrednie przeniesienie
na dowolny obszar bez dodatkowych twierdzeń byłoby nieuzasadnione.

Podstawowa różnica między wymiarami jest następująca:
energia $L^2$ i globalne istnienie rozwiązań Leraya--Hopfa
działają w $d=2,3$; w wymiarze dwa mamy dodatkowo
jednoznaczność i kontrolę wirowości, podczas gdy w wymiarze trzy
pojawia się rozciąganie wirów i warunkowe kryteria regularności.
W kolejnych rozdziałach będziemy badać inne mechanizmy
nieliniowe: swobodne brzegi i homogenizację.

\chapter[Nierówności wariacyjne i przeszkoda]{Nierówności wariacyjne, zagadnienie przeszkody i swobodny brzeg}
\label{chap:pde29}

\section{Motywacja: energia z ograniczeniem jednostronnym}
W poprzednich rozdziałach minimalizowaliśmy funkcjonały energii na całych
przestrzeniach Sobolewa. W wielu zagadnieniach mechanicznych funkcja
$u$ nie może jednak przekroczyć przeszkody w ustalonym kierunku.
Na przykład odkształcenie membrany nie może być mniejsze niż wysokość
sztywnej podpory $\psi$. W takim przypadku warunek Eulera--Lagrange'a
nie musi być równaniem w całym obszarze. Na zbiorze kontaktu pojawia
się dodatkowa reakcja, a sam zbiór kontaktu nie jest znany z góry.

Niech $\Omega\subset\R^d$ będzie ograniczonym obszarem Lipschitzowskim.
Oznaczamy $V=H_0^1(\Omega)$, wyposażony w normę
$\|v\|_V=\|\nabla v\|_{L^2(\Omega)}$, i $V'=H^{-1}(\Omega)$.
Wszystkie relacje $v\geq w$ rozumiemy prawie wszędzie w $\Omega$.
Przyjmujemy $\psi\in H^1(\Omega)$ oraz $\psi^+\in V$;
ostatni warunek gwarantuje istnienie funkcji spełniającej
nierówność przeszkody i zerowy warunek brzegowy.
Definiujemy
\[
 K_\psi=\{v\in V:v\geq\psi\ \text{prawie wszędzie w }\Omega\}.
\]
Zbiór ten jest niepusty, bo $\psi^+\in K_\psi$, oraz wypukły.
Przez $f\in V'$ oznaczamy obciążenie, a przez
$\langle f,v\rangle$ wartość funkcjonału $f$ na $v$.

\begin{lemat}[Domkniętość zbioru funkcji dopuszczalnych]
\label{lem:pde29-K}
Dla ograniczonego obszaru Lipschitzowskiego $\Omega\subset\R^d$
i dowolnej funkcji $\psi\in H^1(\Omega)$ spełniającej
$\psi^+\in H_0^1(\Omega)$ zbiór $K_\psi$ jest niepusty,
wypukły i słabo domknięty w $H_0^1(\Omega)$.
\end{lemat}
\begin{proof}
Niepustość i wypukłość wynikają bezpośrednio z definicji.
Jeżeli $v_n\to v$ silnie w $V$, to z ciągłości zanurzenia
$V\hookrightarrow L^2(\Omega)$ mamy
$\|(\psi-v)^+\|_2\le\|v-v_n\|_2$, ponieważ
$v_n\ge\psi$. Zatem $v\in K_\psi$, więc zbiór jest normowo
domknięty. Każdy domknięty i wypukły podzbiór przestrzeni
Hilberta jest słabo domknięty: w przeciwnym razie twierdzenie
o oddzielaniu dałoby ciągły funkcjonał liniowy ściśle oddzielający
punkt słabego domknięcia od zbioru. Można też zastosować
bezpośrednio twierdzenie Mazura do słabo zbieżnego ciągu w $K_\psi$.
\end{proof}

\section{Rozwiązanie wariacyjne i jego jednoznaczność}
Niech $A(x)=(a_{ij}(x))_{i,j=1}^d$ będzie mierzalną macierzą
symetryczną, $A\in L^\infty(\Omega;\R^{d\times d})$, oraz
$\lambda,\Lambda>0$ takie, że dla prawie każdego $x\in\Omega$
i każdego $\xi\in\R^d$ zachodzi
$\lambda|\xi|^2\le\xi^{\mathsf T}A(x)\xi\le\Lambda|\xi|^2$.
Wprowadzamy formę $a(u,v)=\int_\Omega A\nabla u\cdot\nabla v\,
\mathrm d\lambda_d(x)$ i funkcjonał
$J_f(v)=\frac12a(v,v)-\langle f,v\rangle$.

\begin{twierdzenie}[Istnienie i jednoznaczność w zagadnieniu przeszkody]
\label{thm:pde29-existence}
Niech $\Omega$, $A$, $\psi$ spełniają powyższe założenia,
a $f\in H^{-1}(\Omega)$. Istnieje dokładnie jedna funkcja
$u\in K_\psi$, dla której
\begin{equation}\label{eq:pde29-VI}
 a(u,v-u)\ge\langle f,v-u\rangle
 \qquad\text{dla każdego }v\in K_\psi.
\end{equation}
Jest ona jedynym minimum $J_f$ na $K_\psi$.
\end{twierdzenie}
\begin{proof}
Z eliptyczności i nierówności Poincarégo mamy
$J_f(v)\ge\frac\lambda2\|v\|_V^2-\|f\|_{V'}\|v\|_V$;
$J_f$ jest więc ograniczony od dołu na $K_\psi$, a każdy ciąg
minimalizujący jest ograniczony w $V$. Wybieramy podciąg
zbieżny słabo w $V$ do pewnego $u$. Z lematu~\ref{lem:pde29-K}
wynika $u\in K_\psi$. Norma wyznaczona przez $a$ jest
słabo półciągła z dołu, a składnik liniowy jest słabo ciągły.
Wobec tego $J_f(u)\le\liminf J_f(u_n)$, co dowodzi istnienia.

Dla $v\in K_\psi$ odcinek $u+t(v-u)$ należy do $K_\psi$
dla $t\in[0,1]$. Minimalność daje
$0\le[J_f(u+t(v-u))-J_f(u)]/t
 =a(u,v-u)-\langle f,v-u\rangle+
 \frac t2a(v-u,v-u)$. Przechodząc z $t$ do zera,
otrzymujemy~\eqref{eq:pde29-VI}. Odwrotnie, dla funkcji
spełniającej tę nierówność
$J_f(v)-J_f(u)=a(u,v-u)-\langle f,v-u\rangle
+\frac12a(v-u,v-u)\ge0$.
Jeżeli $u_1,u_2$ spełniają nierówność, testujemy pierwszą przez
$v=u_2$, drugą przez $v=u_1$ i dodajemy. Otrzymujemy
$a(u_1-u_2,u_1-u_2)\le0$, skąd $u_1=u_2$.
\end{proof}


\begin{stwierdzenie}[Zagadnienie przeszkody jako rzut energetyczny]
\label{prop:pde-v44-obstacle-projection}
Niech $V=H_0^1(\Omega)$, $a$ będzie symetryczną,
ciągłą i koercywną formą z twierdzenia~\ref{thm:pde29-existence},
a $K_\psi\subset V$ niepustym, domkniętym i wypukłym
zbiorem dopuszczalnym. Niech $z\in V$ będzie jedynym
rozwiązaniem równania $a(z,v)=\langle f,v\rangle$ dla
każdego $v\in V$. Wtedy rozwiązanie przeszkody $u$
jest rzutem metrycznym $z$ na $K_\psi$ w normie
$\|v\|_a=a(v,v)^{1/2}$. Co więcej, dla ustalonego
$K_\psi$ i dwóch danych $f_i\in V'$ zachodzi
\begin{equation}\label{eq:pde-v44-projection-nonexpansive}
 \|u_1-u_2\|_a\le \|z_1-z_2\|_a
 \le\lambda^{-1/2}\|f_1-f_2\|_{V'},
\end{equation}
gdzie $z_i$ rozwiązuje równanie liniowe dla $f_i$,
a $u_i$ odpowiednie zagadnienie przeszkody.
\end{stwierdzenie}
\begin{proof}
Istnienie $z$ wynika z twierdzenia Laxa--Milgrama;
forma $a$ jest iloczynem skalarnym równoważnym
standardowemu na $V$. Z warunku
$a(u,v-u)\ge\langle f,v-u\rangle=a(z,v-u)$
dla $v\in K_\psi$ otrzymujemy
$a(u-z,v-u)\ge0$. Jest to dokładnie
charakteryzacja rzutu na domknięty zbiór wypukły
w przestrzeni Hilberta $(V,a)$.
Dla dwóch rzutów wykorzystujemy w nierównościach
wariacyjnych testy $v=u_2$ i $v=u_1$, a następnie
dodajemy. Dla $w=u_1-u_2$ dostajemy
$a(w,w)\le a(z_1-z_2,w)$; nierówność Cauchy'ego--Schwarza
daje pierwsze oszacowanie. Wreszcie, z
$a(z_1-z_2,v)=\langle f_1-f_2,v\rangle$ oraz
$a(v,v)\ge\lambda\|v\|_V^2$ mamy
$\|z_1-z_2\|_a\le\lambda^{-1/2}\|f_1-f_2\|_{V'}$.
\end{proof}


\begin{stwierdzenie}[Stabilność względem obciążenia]
\label{prop:pde29-stable}
Przy ustalonych $\Omega$, $A$, $\psi$ z
 twierdzenia~\ref{thm:pde29-existence}, jeżeli $u_i$ odpowiadają
$f_i\in V'$ ($i=1,2$), to
$\|u_1-u_2\|_V\le\lambda^{-1}\|f_1-f_2\|_{V'}$.
\end{stwierdzenie}
\begin{proof}
W obu nierównościach wariacyjnych wybieramy jako test funkcję
będącą drugim rozwiązaniem. Po dodaniu i zastosowaniu
eliptyczności otrzymujemy
$\lambda\|u_1-u_2\|_V^2\le
\langle f_1-f_2,u_1-u_2\rangle$.
Dalej stosujemy nierówność dualną; przypadek $u_1=u_2$
nie wymaga dzielenia przez normę.
\end{proof}

\begin{twierdzenie}[Zasada porównania dla przeszkód]
\label{thm:pde29-order}
Niech $A$, $\Omega$ spełniają powyższe założenia,
a $\psi_i\in H^1(\Omega)$, $\psi_i^+\in V$ oraz
$f_i\in L^2(\Omega)$, $i=1,2$. Jeśli
$\psi_1\le\psi_2$ i $f_1\le f_2$ prawie wszędzie,
to odpowiadające rozwiązania spełniają $u_1\le u_2$
prawie wszędzie w $\Omega$.
\end{twierdzenie}
\begin{proof}
Niech $w=(u_1-u_2)^+\in V$. Ponieważ
$\min(u_1,u_2)\ge\psi_1$ i
$\max(u_1,u_2)\ge\psi_2$,
w nierówności dla $u_1$ wolno wziąć $v=u_1-w$,
a w nierówności dla $u_2$ funkcję $v=u_2+w$.
Daje to $a(u_1,w)\le\int_\Omega f_1w\,\mathrm d\lambda_d(x)$
oraz $a(u_2,w)\ge\int_\Omega f_2w\,\mathrm d\lambda_d(x)$.
Po odjęciu i wykorzystaniu $\nabla w=\chi_{\{u_1>u_2\}}
\nabla(u_1-u_2)$ otrzymujemy
$\lambda\|w\|_V^2\le\int_\Omega(f_1-f_2)w\,\mathrm d\lambda_d(x)
\le0$. Stąd $w=0$.
\end{proof}


\begin{twierdzenie}[Ciągła zależność od przeszkody w normie energetycznej]
\label{thm:pde29-obstacle-stability}
Niech $\Omega\subset\R^d$ będzie ograniczonym obszarem
Lipschitzowskim, $V=H_0^1(\Omega)$, a $A$ niech będzie
symetryczna, mierzalna i niech spełnia warunki
$\lambda|\xi|^2\le\xi^{\mathsf T}A(x)\xi\le\Lambda|\xi|^2$
dla $\lambda,\Lambda>0$. Niech $f\in V'$ będzie ustalone.
Dla $j=1,2$ wybierzmy przeszkodę $\psi_j\in V$
i oznaczmy przez $u_j\in K_{\psi_j}$ rozwiązanie
nierówności wariacyjnej~\eqref{eq:pde29-VI}.
Wtedy
\begin{equation}\label{eq:pde29-obstacle-stability}
 \|u_1-u_2\|_V\le\frac{\Lambda}{\lambda}
 \|\psi_1-\psi_2\|_V.
\end{equation}
\end{twierdzenie}
\begin{proof}
Niech $h=\psi_1-\psi_2\in V$.
Ponieważ $u_2\ge\psi_2$, mamy $u_2+h\ge\psi_1$;
podobnie $u_1-h\ge\psi_2$. Są to zatem poprawne
funkcje testowe w obu nierównościach wariacyjnych.
Pierwszą testujemy przez $u_2+h$, a drugą przez $u_1-h$.
Po dodaniu, przy $w=u_1-u_2$, otrzymujemy
\[
 a(w,w-h)\le\langle f,w-h\rangle
                 -\langle f,w-h\rangle=0.
\]
Stąd $\lambda\|w\|_V^2\le a(w,w)\le a(w,h)
\le\Lambda\|w\|_V\|h\|_V$.
Dzieląc przez $\|w\|_V$ dla $w\ne0$,
otrzymujemy~\eqref{eq:pde29-obstacle-stability};
przypadek $w=0$ jest oczywisty.
\end{proof}

\begin{uwaga}[Ograniczenie założenia o przeszkodach]
Twierdzenie dotyczy przeszkód $\psi_j\in H_0^1(\Omega)$,
a nie tylko warunku $\psi_j^+\in H_0^1(\Omega)$.
Jest to istotne: przesunięte funkcje testowe
$u_2+\psi_1-\psi_2$ i $u_1-\psi_1+\psi_2$
muszą należeć do $H_0^1(\Omega)$.
\end{uwaga}

\section{Reakcja przeszkody i warunki komplementarności}
W tej sekcji rozważamy $A=I$. Oznaczamy przez $-\Delta u$
dystrybucję określoną wzorem
$\langle-\Delta u,\varphi\rangle
=\int_\Omega\nabla u\cdot\nabla\varphi\,\mathrm d\lambda_d(x)$.

\begin{stwierdzenie}[Dodatnia reakcja jako miara]
\label{prop:pde29-lambda-measure}
Niech $u\in K_\psi$ rozwiązuje~\eqref{eq:pde29-VI}
dla $A=I$ i $f\in V'$. Dystrybucja
$\mu=-\Delta u-f$ jest nieujemna, czyli
$\langle\mu,\varphi\rangle\ge0$ dla każdego
$\varphi\in C_c^\infty(\Omega)$, $\varphi\ge0$.
Wobec tego $\mu$ jest nieujemną miarą Radona w $\Omega$.
\end{stwierdzenie}
\begin{proof}
Jeśli $\varphi\ge0$, to $u+\varphi\in K_\psi$,
więc~\eqref{eq:pde29-VI} daje
$\langle-\Delta u-f,\varphi\rangle\ge0$.
Dla dowolnego zwartego $K\Subset\Omega$ wybierzmy
$\eta\in C_c^\infty(\Omega)$, $\eta\ge0$, $\eta=1$
w otoczeniu $K$. Jeżeli $\varphi\in C_c^\infty(K)$,
to $-\|\varphi\|_\infty\eta\le\varphi\le
\|\varphi\|_\infty\eta$. Dodatniość daje
$|\langle\mu,\varphi\rangle|\le
\|\varphi\|_\infty\langle\mu,\eta\rangle$.
Zatem $\mu$ jest lokalnie funkcjonałem ograniczonym
w normie jednostajnej; lokalne twierdzenie Riesza o reprezentacji
daje miarę Radona. Ta konstrukcja nie uprawnia jeszcze do mnożenia
$\mu$ przez dowolną funkcję Sobolewa: do tego potrzebna jest
kontrola regularności, którą uzyskamy poniżej.
\end{proof}


\begin{lemat}[Pochodne funkcji Sobolewa na zbiorze jej zer]
\label{lem:pde29-zero-level}
Niech $\Omega\subset\R^d$ będzie zbiorem otwartym.
Jeżeli $z\in H^1_{\mathrm{loc}}(\Omega)$ i
$E=\{x\in\Omega:z(x)=0\}$, to $\nabla z=0$
prawie wszędzie na $E$. Jeśli ponadto $z\in H^2_{\mathrm{loc}}(\Omega)$,
to $D^2z=0$ prawie wszędzie na $E$.
\end{lemat}
\begin{proof}
W jednym wymiarze funkcja bezwzględnie ciągła ma pochodną
równą zeru prawie wszędzie na zbiorze, na którym sama znika.
Istotnie, w każdym punkcie gęstości tego zbioru, w którym
istnieje pochodna, możemy wybierać punkty zbioru zer
zbiegające do niego z obu stron; ilorazy różnicowe są wtedy
zerowe, zatem pochodna także znika. Przez twierdzenie
Fubiniego i własność absolutnej ciągłości reprezentantów
funkcji Sobolewa na prawie wszystkich prostych równoległych
do osi otrzymujemy $\partial_jz=0$ prawie wszędzie na $E$
dla każdego $j$. Jeżeli $z\in H^2_{\mathrm{loc}}$, to
$\partial_jz\in H^1_{\mathrm{loc}}$ i znika prawie wszędzie
na $E$. Stosując pierwszy wniosek do każdej funkcji
$\partial_jz$, dostajemy $\partial_k\partial_jz=0$
prawie wszędzie na $E$, dla wszystkich $j,k$.
\end{proof}

\begin{twierdzenie}[Regularność $H^2$ i nierówność Lewy'ego--Stampacchii]
\label{thm:pde29-LS}
Niech $\Omega\subset\R^d$ będzie ograniczonym obszarem
klasy $C^{1,1}$, $f\in L^2(\Omega)$ i
$\psi\in H^2(\Omega)$, przy czym ślad $\psi$ na
$\partial\Omega$ jest niedodatni. Niech $u\in H_0^1(\Omega)$
będzie rozwiązaniem zagadnienia przeszkody dla $A=I$.
Wówczas $u\in H^2(\Omega)$, a $\mu=-\Delta u-f$ jest funkcją
z $L^2(\Omega)$ i prawie wszędzie zachodzi
\begin{equation}\label{eq:pde29-LS}
 0\le\mu\le(-\Delta\psi-f)^+,
 \qquad f\le-\Delta u\le\max\{f,-\Delta\psi\}.
\end{equation}
Ponadto $(u-\psi)\mu=0$ prawie wszędzie w $\Omega$.
\end{twierdzenie}
\begin{proof}
\textbf{1. Zadanie z karą.}
Dla $\eps>0$ minimalizujemy na $V=H_0^1(\Omega)$
funkcjonał
\[
 J_\eps(v)=\frac12\|\nabla v\|_2^2
 -\int_\Omega fv\,\mathrm d\lambda_d(x)
 +\frac1{2\eps}\int_\Omega((\psi-v)^+)^2\,\mathrm d\lambda_d(x).
\]
Jest on ściśle wypukły, słabo półciągły z dołu i koercywny,
więc ma jedyne minimum $u_\eps$. Funkcja
$t\mapsto((\psi-t)^+)^2$ jest różniczkowalna z pochodną
$-2(\psi-t)^+$; nierówność Lipschitza dla funkcji
$z\mapsto z^+$ i zbieżność zdominowana uzasadniają
różniczkowanie całki. Otrzymujemy
$-\Delta u_\eps=f+\mu_\eps$, gdzie
$\mu_\eps=\eps^{-1}(\psi-u_\eps)^+\ge0$.

\textbf{2. Jednostajna kontrola reakcji.}
Niech $w_\eps=(\psi-u_\eps)^+$. Ponieważ ślad
$\psi$ jest niedodatni, $w_\eps\in H_0^1(\Omega)$.
Wprowadzamy $h=-\Delta\psi-f\in L^2(\Omega)$.
Odejmujemy równanie $-\Delta\psi=f+h$ od
równania dla $u_\eps$ i testujemy różnicę przez
$w_\eps$. Na zbiorze $\{w_\eps>0\}$ mamy
$\nabla(u_\eps-\psi)=-\nabla w_\eps$. Dostajemy
\begin{equation}\label{eq:pde29-penalty}
 \|\nabla w_\eps\|_2^2+\eps\|\mu_\eps\|_2^2
 =\eps\int_\Omega h\mu_\eps\,\mathrm d\lambda_d(x)
 \le\eps\|h^+\|_2\|\mu_\eps\|_2.
\end{equation}
W szczególności $\|\mu_\eps\|_2\le\|h^+\|_2$,
$\|w_\eps\|_2\le\eps\|h^+\|_2$ oraz
$\|\nabla w_\eps\|_2\to0$.
Testujemy równanie $-\Delta u_\eps=f+\mu_\eps$ przez $u_\eps$.
Z nierówności Poincarégo i oszacowania reakcji otrzymujemy
\[
 \|\nabla u_\eps\|_2^2
 =\int_\Omega(f+\mu_\eps)u_\eps\,\mathrm d\lambda_d(x)
 \le C_P(\|f\|_2+\|h^+\|_2)\|\nabla u_\eps\|_2.
\]
Zatem $\|u_\eps\|_V\le C_P(\|f\|_2+\|h^+\|_2)$
niezależnie od $\eps$, co uzasadnia wybranie słabo zbieżnego
podciągu w następnym kroku.

\textbf{3. Granica i warunek komplementarności.}
Po wybraniu podciągu $u_\eps\rightharpoonup u$ w $V$,
$u_\eps\to u$ w $L^2(\Omega)$ i
$\mu_\eps\rightharpoonup\mu$ w $L^2(\Omega)$.
Ponieważ $w_\eps\to0$ w $L^2$, mamy $u\ge\psi$.
Przechodzimy do granicy w równaniu, dostając
$-\Delta u=f+\mu$ w $V'$, z $\mu\ge0$.
Dodatkowo
\[
 \int_\Omega(u_\eps-\psi)\mu_\eps\,\mathrm d\lambda_d(x)
 =-\eps\|\mu_\eps\|_2^2\longrightarrow0.
\]
Zbieżność silna pierwszego czynnika w $L^2$ i słaba drugiego
pozwalają przejść do granicy. Stąd
$\int_\Omega(u-\psi)\mu\,\mathrm d\lambda_d(x)=0$;
nieujemność czynników daje $(u-\psi)\mu=0$ prawie wszędzie.
Dla $v\in K_\psi$ zachodzi więc
$\int_\Omega\mu(v-u)\,\mathrm d\lambda_d(x)
 =\int_\Omega\mu(v-\psi)\,\mathrm d\lambda_d(x)\ge0$.
Wynika stąd~\eqref{eq:pde29-VI}, a jednoznaczność z
twierdzenia~\ref{thm:pde29-existence} identyfikuje $u$.

\textbf{4. Regularność i nierówność punktowa.}
Z równania $-\Delta u=f+\mu\in L^2(\Omega)$ oraz
globalnego oszacowania regularności Dirichleta na obszarze
klasy $C^{1,1}$ (rozdział~\ref{chap:pde15}) otrzymujemy
$u\in H^2(\Omega)$. Niech $z=u-\psi\in H^2(\Omega)$.
Z lematu~\ref{lem:pde29-zero-level} wynika
$D^2z=0$ prawie wszędzie na zbiorze $\{z=0\}$.
Na zbiorze kontaktu $\{u=\psi\}$ mamy więc
$-\Delta u=-\Delta\psi$ prawie wszędzie;
na jego dopełnieniu $\mu=0$ na mocy komplementarności.
To daje $0\le\mu\le h^+$ oraz drugi zapis~\eqref{eq:pde29-LS}.
\end{proof}

\begin{uwaga}[Granica regularności]
Oszacowanie $H^2$ zależy od regularności brzegu, a
nierówność~\eqref{eq:pde29-LS} od drugich pochodnych przeszkody.
Nie wolno ich przypisywać rozwiązaniu dla dowolnej
przeszkody $H^1$ ani dowolnego obszaru Lipschitzowskiego.
Na ogólnych obszarach $\mu$ jest na początku jedynie miarą
Radona, a swobodny brzeg może być nieregularny.
\end{uwaga}

\section{Zbiór kontaktu i swobodny brzeg}
Załóżmy dodatkowo $d\le3$, $\Omega$ klasy $C^{1,1}$ i
$\psi\in H^2(\Omega)$, $f\in L^2(\Omega)$.
Z zanurzenia $H^2(\Omega)\hookrightarrow C^0(\overline\Omega)$
dla $d\le3$ funkcje $u$ i $\psi$ mają ciągłe reprezentacje.
Definiujemy wtedy
$\mathcal C=\{x\in\Omega:u(x)=\psi(x)\}$ (zbiór kontaktu)
oraz $\mathcal N=\{x\in\Omega:u(x)>\psi(x)\}$ (zbiór braku kontaktu).
Zbiory $\mathcal C$ i $\mathcal N$ są odpowiednio względnie domknięty
i otwarty w $\Omega$. Zbiorem swobodnego brzegu nazywamy
$\Gamma=\partial_\Omega\mathcal N\cap\Omega$.
W $\mathcal N$ zachodzi $-\Delta u=f$ w sensie dystrybucji,
gdyż reakcja $\mu$ tam znika. Zbiór $\Gamma$ nie jest dany
z góry i jego regularność wymaga dodatkowej analizy.

\begin{przyklad}[Swobodny brzeg i skok drugiej pochodnej w wymiarze jeden]
\label{ex:pde29-explicit}
Niech $\Omega=(-1,1)$, $f=0$ i
$\psi(x)=(1-x^2)^2$. Jedynym rozwiązaniem zadania
przeszkody jest funkcja
\begin{equation}\label{eq:pde29-explicit}
 u(x)=\begin{cases}
 (1-x^2)^2,& |x|\le\frac13,\\
 \frac{32}{27}(1-|x|),&\frac13<|x|\le1.
 \end{cases}
\end{equation}
Istotnie, $u(\pm1)=0$, $u$ i $u'$ są ciągłe w punktach
$\pm1/3$. Dla $x\in[1/3,1]$ zachodzi
\[
 u(x)-\psi(x)=\frac{(1-x)(3x-1)^2(3x+5)}{27}\ge0;
\]
na lewej połowie korzystamy z parzystości. Ponadto
$-u''=(4-12x^2)\chi_{\{|x|<1/3\}}\ge0$ prawie wszędzie.
Reakcja jest skupiona na zbiorze kontaktu, więc dla każdej
$v\in K_\psi$ zachodzi
$\int_{-1}^1u'(v'-u')\,\mathrm d\lambda_1(x)
=\int_{-1}^1(-u'')(v-u)\,\mathrm d\lambda_1(x)\ge0$.
Z jednoznaczności wynika wzór~\eqref{eq:pde29-explicit}.
Swobodny brzeg składa się z dwóch punktów
$\Gamma=\{-1/3,1/3\}$. Funkcja $u$ należy do
$C^{1,1}([-1,1])$, ale nie do $C^2([-1,1])$,
ponieważ $u''$ ma skok w tych punktach mimo gładkości $\psi$.
\end{przyklad}

\begin{center}
\begin{tikzpicture}[x=4.4cm,y=2.5cm]
 \draw[-{Latex}] (-1.06,0)--(1.08,0) node[right] {$x$};
 \draw[-{Latex}] (-1,-.12)--(-1,1.15) node[above] {$u,\psi$};
 \draw[dashed,gray] (-.3333,0)--(-.3333,1.07);
 \draw[dashed,gray] (.3333,0)--(.3333,1.07);
 \draw[blue!70!black,thick,samples=180,domain=-1:1,smooth]
   plot (\x,{(1-\x*\x)^2});
 \draw[red!75!black,very thick,samples=60,domain=-.3333:.3333]
   plot (\x,{(1-\x*\x)^2});
 \draw[red!75!black,very thick] (-1,0)--(-.3333,.79012)
   (.3333,.79012)--(1,0);
 \node[below] at (-1,0) {$-1$};
 \node[below] at (-.3333,0) {$-\tfrac13$};
 \node[below] at (.3333,0) {$\tfrac13$};
 \node[below] at (1,0) {$1$};
 \node[blue!65!black,right] at (.58,.26) {$\psi$};
 \node[red!70!black,right] at (.62,.54) {$u$};
 \node[above] at (0,1.02) {zbiór kontaktu};
\end{tikzpicture}

\small Rysunek \ref{ex:pde29-explicit}: przeszkoda (niebieska)
i rozwiązanie (czerwone). W środku zachodzi $u=\psi$;
poza tym przedziałem $-u''=0$.
\end{center}

\section{Projekcja na zbiór wypukły i mnożniki ograniczeń}
Dla funkcjonału $f\in V'$ niech $w\in V$ będzie jedynym
rozwiązaniem $a(w,v)=\langle f,v\rangle$ dla $v\in V$
(twierdzenie Laxa--Milgrama). Wówczas
$J_f(v)=\frac12a(v-w,v-w)-\frac12a(w,w)$.
Rozwiązanie nierówności wariacyjnej jest zatem rzutem $w$
na $K_\psi$ w metryce określonej przez iloczyn skalarny $a$.
Nie jest to na ogół rzut punktowy $\max\{w,\psi\}$:
energia kontroluje pochodne, więc zmiana w jednym miejscu
wpływa na rozwiązanie w całym obszarze.

Dla domkniętego wypukłego zbioru $K\subset V$ definiujemy
stożek normalny
$N_K(u)=\{\eta\in V':\langle\eta,v-u\rangle\le0
\text{ dla każdego }v\in K\}$, jeżeli $u\in K$,
i $N_K(u)=\varnothing$ w przeciwnym przypadku.
Jeżeli $\mathcal A:V\to V'$ jest operatorem odpowiadającym $a$,
to nierówność~\eqref{eq:pde29-VI} jest równoważna
$0\in\mathcal Au-f+N_{K_\psi}(u)$.
Ta postać pozwala przenieść argumenty na inne problemy
jednostronne, także na przeszkody ruchome oraz przepływy
gradientowe z ograniczeniami.

\section{Schemat numeryczny i warunki punktowe}
W wymiarze jeden, dla $\Omega=(-1,1)$ i równomiernej siatki
$x_i=-1+ih$, $h=2/N$, $i=0,\ldots,N$, oznaczamy
$U_i\approx u(x_i)$ i $\Psi_i=\psi(x_i)$.
Dla $f=0$ dyskretna energia jest równa
$E_h(U)=\frac1{2h}\sum_{i=0}^{N-1}(U_{i+1}-U_i)^2$
przy $U_0=U_N=0$. Minimalizujemy ją na zbiorze
$U_i\ge\Psi_i$ dla $i=1,\dots,N-1$.
W węzłach wewnętrznych warunki optymalności mają postać
\begin{equation}\label{eq:pde29-discrete}
 U_i\ge\Psi_i,\qquad
 (2U_i-U_{i-1}-U_{i+1})/h^2\ge0,\qquad
 (U_i-\Psi_i)(2U_i-U_{i-1}-U_{i+1})=0.
\end{equation}
Dla przeszkody z przykładu~\ref{ex:pde29-explicit}
można zastosować metodę Gaussa--Seidela z projekcją,
zastępując każdy nowy węzeł przez większą z wartości
$\Psi_i$ i średniej sąsiednich wartości. Sprawdzanie
nierówności~\eqref{eq:pde29-discrete} jest ważniejsze niż
sam obraz krzywej: wykres może wyglądać poprawnie również
wtedy, gdy warunek komplementarności jest naruszony.

\section{Dalsze pytania i granice wykazanej teorii}
Udowodniliśmy istnienie i jednoznaczność nierówności
wariacyjnej w $H_0^1$, a przy mocniejszych założeniach
regularność $H^2$ i nierówność Lewy'ego--Stampacchii.
Nie udowodniliśmy ogólnej regularności $C^{1,1}$
przy dowolnej gładkiej przeszkodzie ani regularności
swobodnego brzegu. Analiza jego punktów regularnych
i osobliwych wymaga dodatkowych oszacowań wzrostu,
monotoniczności i argumentów geometrycznych;
nie wynika z samego sformułowania wariacyjnego.


\chapter[Homogenizacja okresowa równań eliptycznych]{Homogenizacja równań eliptycznych: zadanie komórkowe, macierz efektywna i granica oscylacyjna}
\label{chap:pde30}

\section{Dlaczego współczynników nie można po prostu uśrednić?}
Rozważmy ograniczony obszar Lipschitzowski $\Omega\subset\R^d$, $d\ge1$,
liczbę $\eps>0$ i macierz $A(y)$, która zmienia się w skali
pojedynczej komórki $Y=(0,1)^d$. W zagadnieniu Dirichleta
\begin{equation}\label{eq:pde30-micro}
 -\diver\bigl(A(x/\eps)\nabla u_\eps\bigr)=f\quad\text{w }\Omega,
 \qquad u_\eps\in H_0^1(\Omega),
\end{equation}
rozwiązanie odczuwa coraz drobniejsze niejednorodności. Pytamy,
czy istnieje granica $u_0$ i czy spełnia równanie tej samej postaci,
z pewną stałą macierzą $A^{\mathrm{hom}}$. Samo zastąpienie $A$
jego średnią arytmetyczną jest na ogół błędne: w kierunku
prostopadłym do warstw o przepływie decyduje średnia harmoniczna.

Przyjmujemy okres komórki równy $1$ w każdym kierunku.
Jest to inna normalizacja geometryczna niż dla torusa
$\T^d=\R^d/(2\pi\mathbb Z)^d$ z rozdziału~\ref{chap:pde25};
przy porównywaniu wzorów trzeba przeskalować zmienne.
Przez $H^1_{\mathrm{per}}(Y)$ rozumiemy przestrzeń ograniczeń do
$Y$ funkcji z $H^1_{\mathrm{loc}}(\R^d)$, których przedłużenie jest
$\mathbb Z^d$-okresowe. Jej podprzestrzenią o średniej zero jest
$H^1_{\mathrm{per},0}(Y)=\{v\in H^1_{\mathrm{per}}(Y):
\int_Yv\,\mathrm d\lambda_d(y)=0\}$.
Niech $A\in L^\infty_{\mathrm{per}}(Y;\R^{d\times d})$ będzie
\emph{symetryczna} oraz niech $0<\lambda\le\Lambda$ spełniają
\begin{equation}\label{eq:pde30-ellipticity}
 \lambda|\xi|^2\le\xi^{\mathsf T}A(y)\xi\le\Lambda|\xi|^2
 \quad (\xi\in\R^d),\quad\text{dla prawie każdego }y\in Y.
\end{equation}
Macierz przedłużamy okresowo na $\R^d$. Z warunku~\eqref{eq:pde30-ellipticity}
i twierdzenia Laxa--Milgrama zagadnienie~\eqref{eq:pde30-micro}
ma dokładnie jedno rozwiązanie dla każdego $f\in H^{-1}(\Omega)$;
ponadto $\|\nabla u_\eps\|_{L^2}\le\lambda^{-1}\|f\|_{H^{-1}}$.
Zapis $\diver$ oznacza dywergencję, a $\mathrm d\lambda_d$ miarę Lebesgue'a.

\section{Przypadek jednowymiarowy: dokładny wzór i tempo zbieżności}
Niech $\Omega=(0,L)$, $L>0$, $a\in L^\infty_{\mathrm{per}}(0,1)$
oraz $0<\lambda\le a\le\Lambda$ prawie wszędzie. Definiujemy
$b=1/a$ i $\bar b=\int_0^1b(y)\,\mathrm d\lambda_1(y)$, a także
$a_{\mathrm{hom}}=\bar b^{-1}$. Dla $f\in L^2(0,L)$ połóżmy
$F(x)=\int_0^xf(t)\,\mathrm d\lambda_1(t)$.

\begin{twierdzenie}[Jednowymiarowe równanie uśrednione]\label{thm:pde30-1d}
Dla każdego $\eps>0$ istnieje dokładnie jedno $u_\eps\in H_0^1(0,L)$
spełniające $-(a(x/\eps)u_\eps')'=f$ w sensie dystrybucji.
Niech $u_0$ będzie jedynym rozwiązaniem
$-a_{\mathrm{hom}}u_0''=f$ w $H_0^1(0,L)$.
Wtedy, gdy $\eps\downarrow0$,
$u_\eps\to u_0$ w $L^\infty(0,L)$ oraz
$u_\eps\rightharpoonup u_0$ w $H_0^1(0,L)$, a strumienie
$a(x/\eps)u_\eps'$ zbiegają silnie w $L^2(0,L)$ do
$a_{\mathrm{hom}}u_0'$. Istnieje stała $C$, zależna tylko od
$L,\lambda,\Lambda$, taka że dla $0<\eps\le1$
\begin{equation}\label{eq:pde30-1d-rate}
 \|u_\eps-u_0\|_{L^\infty(0,L)}
 +\|a(x/\eps)u_\eps'-a_{\mathrm{hom}}u_0'\|_{L^2(0,L)}
 \le C\eps\|f\|_{L^2(0,L)}.
\end{equation}
\end{twierdzenie}
\begin{proof}
Ponieważ $-(a(x/\eps)u_\eps')'=f$, funkcja
$q_\eps=a(x/\eps)u_\eps'$ jest bezwzględnie ciągła i
$q_\eps(x)=C_\eps-F(x)$. Całkując $u_\eps'$ oraz
wykorzystując $u_\eps(L)=u_\eps(0)=0$, otrzymujemy
\begin{equation}\label{eq:pde30-explicit-1d}
 C_\eps=\frac{\int_0^L b(t/\eps)F(t)\,\mathrm d\lambda_1(t)}
 {\int_0^L b(t/\eps)\,\mathrm d\lambda_1(t)},\qquad
 u_\eps(x)=\int_0^xb(t/\eps)(C_\eps-F(t))\,\mathrm d\lambda_1(t).
\end{equation}
W rozwiązaniu granicznym $u_0'=\bar b(C_0-F)$,
gdzie $C_0=L^{-1}\int_0^LF(t)\,\mathrm d\lambda_1(t)$.

Wyjaśnijmy ilościowo uśrednianie. Jeśli $g\in L^\infty_{\mathrm{per}}(0,1)$
ma średnią zero, to jej funkcja pierwotna
$G(s)=\int_0^sg(y)\,\mathrm d\lambda_1(y)$ jest ograniczona
i okresowa. Dla $h\in W^{1,1}(0,L)$ całkowanie przez części
daje
\begin{equation}\label{eq:pde30-average-rate}
 \left|\int_0^Lg(x/\eps)h(x)\,\mathrm d\lambda_1(x)\right|
 \le\eps\|G\|_\infty
 \bigl(|h(0)|+|h(L)|+\|h'\|_{L^1(0,L)}\bigr).
\end{equation}
Stosując ten wzór do $g=b-\bar b$ oraz $h=1,F$, otrzymujemy
$\int_0^Lb(t/\eps)\,\mathrm d\lambda_1(t)=L\bar b+O(\eps)$ i
$\int_0^Lb(t/\eps)F(t)\,\mathrm d\lambda_1(t)
=\bar b\int_0^LF(t)\,\mathrm d\lambda_1(t)+O(\eps\|f\|_2)$.
Mianownik we wzorze na $C_\eps$ jest nie mniejszy od
$L/\Lambda$, zatem $|C_\eps-C_0|\le C\eps\|f\|_2$.
Stosujemy~\eqref{eq:pde30-average-rate} ponownie do
$h=C_0-F$, także na każdym przedziale $(0,x)$,
ze stałą niezależną od $x$. Rozkładając
$u_\eps-u_0$ na składnik
$b(x/\eps)(C_\eps-C_0)$ i oscylującą różnicę
$(b(x/\eps)-\bar b)(C_0-F)$, otrzymujemy pierwszą normę
w~\eqref{eq:pde30-1d-rate}. Druga wynika natychmiast z
$q_\eps-q_0=C_\eps-C_0$.
Jednostajne ograniczenie gradientów i mocna zbieżność funkcji
identyfikują słabą granicę w $H_0^1$ jako $u_0$.
\end{proof}

\begin{stwierdzenie}[Poprawka do gradientu w wymiarze jeden]\label{prop:pde30-corrector1d}
Przy założeniach twierdzenia~\ref{thm:pde30-1d} istnieje jedyna
funkcja $\chi\in W^{1,\infty}_{\mathrm{per}}(0,1)$ o średniej zero,
której pochodna spełnia $\chi'(y)=a_{\mathrm{hom}}/a(y)-1$.
Ponadto
\[
\left\|u_\eps'-\bigl(1+\chi'(x/\eps)\bigr)u_0'\right\|_{L^2(0,L)}
\le C\eps\|f\|_2.
\]
Jeśli $a$ nie jest prawie wszędzie stała i $f\not\equiv0$,
to w ogólności nie zachodzi silna zbieżność $u_\eps'\to u_0'$.
\end{stwierdzenie}
\begin{proof}
Funkcja po prawej stronie równania na $\chi'$ ma średnią zero;
można ją całkować, a następnie odjąć średnią funkcji pierwotnej.
Z tożsamości $u_0'=\bar b(C_0-F)$ wynika
$(1+\chi'(x/\eps))u_0'=b(x/\eps)(C_0-F)$.
Różnica z $u_\eps'$ to $b(x/\eps)(C_\eps-C_0)$,
więc twierdzenie~\ref{thm:pde30-1d} daje oszacowanie.
Jeśli $a$ nie jest stała, to średnia z
$(a_{\mathrm{hom}}/a-1)^2$ jest dodatnia.
Przy $f\not\equiv0$ mamy $u_0'\not\equiv0$, stąd okresowe
uśrednianie w całce z wagą $|u_0'|^2$ daje
\[
\lim_{\eps\downarrow0}\|u_\eps'-u_0'\|_2^2=
\left(\int_0^1|a_{\mathrm{hom}}/a(y)-1|^2
\,\mathrm d\lambda_1(y)\right)\|u_0'\|_2^2>0.
\]
Istotne są więc poprawki oscylacyjne, a nie tylko granica $u_0$.
\end{proof}

\begin{przyklad}[Dlaczego średnia arytmetyczna jest błędna]
Niech na połowie komórki $a=1$, a na drugiej połowie $a=4$.
Wówczas średnia arytmetyczna wynosi $5/2$, a efektywny
współczynnik to $a_{\mathrm{hom}}=(\tfrac12+\tfrac18)^{-1}=8/5$.
Dla $f=1$ na $(0,1)$ rozwiązanie graniczne ma postać
$u_0(x)=x(1-x)/(2a_{\mathrm{hom}})$.
Zastąpienie $a_{\mathrm{hom}}$ przez $5/2$ daje inną funkcję;
różnica nie znika mimo dowolnie drobnych warstw.
\end{przyklad}

\section{Uśrednianie funkcji okresowych}
Do przejścia w wielu wymiarach potrzebujemy prostego lematu.
\begin{lemat}[Słaba granica funkcji okresowych]\label{lem:pde30-average}
Niech $D\subset\R^d$ będzie ograniczonym obszarem Lipschitzowskim,
a $g\in L^2_{\mathrm{per}}(Y)$. Wtedy
$g(x/\eps)\rightharpoonup\bar g:=\int_Yg(y)\,\mathrm d\lambda_d(y)$
słabo w $L^2(D)$, gdy $\eps\downarrow0$.
Ponadto $\eps g(x/\eps)\to0$ silnie w $L^2(D)$.
\end{lemat}
\begin{proof}
Pokrywamy $D$ skończonym prostopadłościanem. Dla $\eps\le1$
całkę $\int_D|g(x/\eps)|^2\,\mathrm d\lambda_d(x)$
oszacowujemy przez sumę po komórkach $\eps(k+Y)$;
liczba tych komórek jest rzędu $\eps^{-d}$, a całka
na każdej jest równa $\eps^d\|g\|_{L^2(Y)}^2$.
Zatem normy są jednostajnie ograniczone, co daje także drugą tezę.

Najpierw załóżmy, że $g$ jest wielomianem trygonometrycznym
$\sum_{k\in F}c_ke^{2\pi i k\cdot y}$. Wyraz stały jest $\bar g$,
a dla $k\ne0$ i $\varphi\in C_c^\infty(D)$ całka
$\int_D\varphi(x)e^{2\pi i k\cdot x/\eps}\,\mathrm d\lambda_d(x)$
zmierza do zera przez całkowanie przez części w kierunku $k$.
Wielomiany trygonometryczne są gęste w $L^2_{\mathrm{per}}(Y)$,
a wspomniane ograniczenie norm pozwala przybliżyć ogólne $g$.
Gęstość $C_c^\infty(D)$ w $L^2(D)$ domyka dowód.
\end{proof}

\section{Zadanie komórkowe i macierz efektywna}
\begin{twierdzenie}[Istnienie i jednoznaczność poprawek komórkowych]\label{thm:pde30-cell}
Niech $A$ spełnia~\eqref{eq:pde30-ellipticity}.
Dla każdego $j\in\{1,\ldots,d\}$ istnieje dokładnie jedna
funkcja $\chi_j\in H^1_{\mathrm{per},0}(Y)$ taka, że
\begin{equation}\label{eq:pde30-cell}
\int_Y A(y)(e_j+\nabla_y\chi_j(y))\cdot\nabla_y\varphi(y)
\,\mathrm d\lambda_d(y)=0
\quad\text{dla każdego }\varphi\in H^1_{\mathrm{per},0}(Y).
\end{equation}
Definiujemy wektory strumienia $Q_j(y)=A(y)(e_j+\nabla\chi_j(y))$
i macierz stałą
\begin{equation}\label{eq:pde30-ahom}
 A^{\mathrm{hom}}e_j=\int_YQ_j(y)\,\mathrm d\lambda_d(y).
\end{equation}
Wtedy $\diver_yQ_j=0$ w sensie dystrybucji na torusie,
$A^{\mathrm{hom}}$ jest symetryczna i
$\lambda|\xi|^2\le\xi^{\mathsf T}A^{\mathrm{hom}}\xi
\le\Lambda|\xi|^2$ dla każdego $\xi\in\R^d$.
\end{twierdzenie}
\begin{proof}
Na przestrzeni $H^1_{\mathrm{per},0}(Y)$ obowiązuje
okresowa nierówność Poincarégo, więc
$\|\nabla\varphi\|_2$ jest normą równoważną normie $H^1$.
Forma $\int_YA\nabla v\cdot\nabla\varphi$ jest ciągła
i koercywna, zaś funkcjonał
$\varphi\mapsto-\int_YAe_j\cdot\nabla\varphi$ jest ciągły.
Twierdzenie Laxa--Milgrama daje jedyne $\chi_j$.
Ponieważ testy o średniej zero wystarczają również dla
wszystkich funkcji okresowych (stała ma gradient zero),
$\diver_yQ_j=0$.

Dla $\xi\in\R^d$ niech $\chi_\xi=\sum_j\xi_j\chi_j$.
Równanie komórkowe i test $\varphi=\chi_\xi$ dają
\begin{equation}\label{eq:pde30-energycell}
 \xi^{\mathsf T}A^{\mathrm{hom}}\xi
 =\int_Y A(y)(\xi+\nabla\chi_\xi)\cdot
             (\xi+\nabla\chi_\xi)\,\mathrm d\lambda_d(y)
 =\min_{v\in H^1_{\mathrm{per},0}(Y)}
 \int_Y (\xi+\nabla v)^{\mathsf T}A(y)(\xi+\nabla v)\,\mathrm d\lambda_d(y).
\end{equation}
Minimalność wynika z rozwinięcia kwadratu i zaniku
wyrazu mieszanego na mocy~\eqref{eq:pde30-cell}.
Ponieważ $\int_Y\nabla v\,\mathrm d\lambda_d(y)=0$,
mamy $\int_Y|\xi+\nabla v|^2\,\mathrm d\lambda_d(y)
=|\xi|^2+\int_Y|\nabla v|^2\,\mathrm d\lambda_d(y)\ge|\xi|^2$.
To daje dolną granicę eliptyczności. Dla górnej wybieramy
$v=0$. Symetrię otrzymujemy z tożsamości
$e_i\cdot A^{\mathrm{hom}}e_j
=\int_Y(e_i+\nabla\chi_i)\cdot A(e_j+\nabla\chi_j)
\,\mathrm d\lambda_d(y)$, wynikającej z testu $\chi_i$
w zadaniu komórkowym. Prawa strona jest symetryczna
względem $i,j$, gdyż $A=A^{\mathsf T}$.
\end{proof}

\begin{uwaga}[Co jest mikroskopowe, a co makroskopowe?]
Funkcja $\chi_j$ jest okresową poprawką komórkową i zależy
od zmiennej szybkiej $y$. Macierz $A^{\mathrm{hom}}$ jest stała
w zmiennej wolnej $x$. Współczynniki $A(x/\eps)$ nie muszą
zbiegać punktowo ani silnie do $A^{\mathrm{hom}}$; zbiega
rozwiązanie i odpowiednio rozumiany strumień. Nawet słaba
granica $A(x/\eps)$ jest na ogół inna niż $A^{\mathrm{hom}}$.
\end{uwaga}

\section{Twierdzenie o homogenizacji: pełny dowód przez oscylujące testy}
\begin{twierdzenie}[Homogenizacja okresowa przy mierzalnych współczynnikach]
\label{thm:pde30-homogenization}
Niech $d\ge1$, $\Omega\subset\R^d$ będzie ograniczonym
obszarem Lipschitzowskim, a $A$ spełnia~\eqref{eq:pde30-ellipticity}.
Dla $f\in L^2(\Omega)$ niech $u_\eps\in H_0^1(\Omega)$ będzie
jedynym słabym rozwiązaniem~\eqref{eq:pde30-micro}.
Niech $u_0\in H_0^1(\Omega)$ spełnia
\begin{equation}\label{eq:pde30-macro}
 -\diver(A^{\mathrm{hom}}\nabla u_0)=f
 \quad\text{w }\Omega,\qquad u_0|_{\partial\Omega}=0.
\end{equation}
Wówczas dla $\eps\downarrow0$ zachodzą zbieżności
\[
 u_\eps\rightharpoonup u_0\text{ w }H_0^1(\Omega),\quad
 u_\eps\to u_0\text{ w }L^2(\Omega),\quad
 A(x/\eps)\nabla u_\eps\rightharpoonup
 A^{\mathrm{hom}}\nabla u_0\text{ w }L^2(\Omega;\R^d).
\]
Przy tych samych założeniach rezultat zachodzi również
dla $f\in H^{-1}(\Omega)$.
\end{twierdzenie}
\begin{proof}
\textbf{Krok 1: zwartość i strumień.}
Z oszacowania energetycznego
$\lambda\|\nabla u_\eps\|_2^2\le
\int_\Omega f u_\eps\,\mathrm d\lambda_d(x)$
i nierówności Poincarégo ciąg $(u_\eps)$ jest ograniczony
w $H_0^1(\Omega)$. Dla dowolnego ciągu $\eps_k\downarrow0$
wybieramy podciąg, dla którego $u_{\eps_k}\rightharpoonup u$
w $H_0^1$ oraz $u_{\eps_k}\to u$ silnie w $L^2$.
Strumienie $q_\eps=A(x/\eps)\nabla u_\eps$ są ograniczone
w $L^2(\Omega;\R^d)$; po dalszym podciągu $q_\eps\rightharpoonup q$.
Ponieważ $\diver q_\eps=-f$, otrzymujemy $\diver q=-f$.
W dalszym ciągu piszemy $\eps$ zamiast indeksów podciągu.

\textbf{Krok 2: oscylujące funkcje testowe.}
Dla $i\in\{1,\ldots,d\}$ zdefiniujmy
$w_i^\eps(x)=x_i+\eps\chi_i(x/\eps)$,
$Q_i^\eps(x)=Q_i(x/\eps)=A(x/\eps)\nabla w_i^\eps(x)$.
Z lematu~\ref{lem:pde30-average} wynika
$w_i^\eps\to x_i$ silnie w $L^2(\Omega)$ oraz
$Q_i^\eps\rightharpoonup A^{\mathrm{hom}}e_i$
w $L^2(\Omega;\R^d)$. Ponadto
$\diver Q_i^\eps=0$ w sensie dystrybucji. Aby to
uzasadnić przy $Q_i\in L^2_{\mathrm{per}}(Y)$, wygładzamy
$Q_i$ splotem z okresowym jądrem wygładzającym.
Otrzymane gładkie strumienie $Q_{i,n}$ nadal mają zerową
dywergencję, gdyż różniczkowanie komutuje ze splotem.
Dla każdego ustalonego $\eps>0$ mamy
$Q_{i,n}(x/\eps)\to Q_i(x/\eps)$ w $L^2_{\mathrm{loc}}(\Omega)$.
Równość $\operatorname{div}_x Q_{i,n}(x/\eps)=0$
przechodzi więc do granicy w dystrybucjach.
Nie wykorzystujemy przy tym klasycznych pochodnych $A$. Nie potrzebujemy ograniczoności funkcji $\chi_i$:
wystarczają jej okresowość i przynależność do $H^1$.

\textbf{Krok 3: identyfikacja granicy iloczynu.}
Dla $\varphi\in C_c^\infty(\Omega)$ z równania
$\diver q_\eps=-f$ otrzymujemy po całkowaniu przez części
\begin{align*}
\int_\Omega\varphi\,q_\eps\cdot\nabla w_i^\eps
\,\mathrm d\lambda_d(x)
&=\int_\Omega\varphi f w_i^\eps\,\mathrm d\lambda_d(x)
 -\int_\Omega w_i^\eps q_\eps\cdot\nabla\varphi
 \,\mathrm d\lambda_d(x)\\
&\longrightarrow \int_\Omega\varphi f x_i\,\mathrm d\lambda_d(x)
 -\int_\Omega x_i q\cdot\nabla\varphi\,\mathrm d\lambda_d(x)
 =\int_\Omega\varphi q_i\,\mathrm d\lambda_d(x).
\end{align*}
Użyliśmy silnej zbieżności $w_i^\eps$ w $L^2$, słabej
zbieżności $q_\eps$ oraz faktu $f\in L^2$.
Z drugiej strony, ponieważ $\diver Q_i^\eps=0$,
\begin{align*}
\int_\Omega\varphi\,Q_i^\eps\cdot\nabla u_\eps
\,\mathrm d\lambda_d(x)
&=-\int_\Omega u_\eps Q_i^\eps\cdot\nabla\varphi
\,\mathrm d\lambda_d(x)\\
&\longrightarrow
 -\int_\Omega u(A^{\mathrm{hom}}e_i)\cdot\nabla\varphi
 \,\mathrm d\lambda_d(x)
 =\int_\Omega\varphi(A^{\mathrm{hom}}e_i)\cdot\nabla u
 \,\mathrm d\lambda_d(x).
\end{align*}
Kluczowa jest dokładna tożsamość
$q_\eps\cdot\nabla w_i^\eps
=Q_i^\eps\cdot\nabla u_\eps$, wynikająca z symetrii $A$.
Porównując granice, dostajemy
$q_i=(A^{\mathrm{hom}}e_i)\cdot\nabla u$ dla każdego $i$.
Ponieważ $A^{\mathrm{hom}}$ jest symetryczna, oznacza to
$q=A^{\mathrm{hom}}\nabla u$.

\textbf{Krok 4: identyfikacja rozwiązania i cały ciąg.}
Równość $\diver q=-f$ i poprzednia identyfikacja dają
$-\diver(A^{\mathrm{hom}}\nabla u)=f$, $u\in H_0^1(\Omega)$.
Z eliptyczności $A^{\mathrm{hom}}$ rozwiązanie jest jedyne,
zatem $u=u_0$. Argument działa dla dowolnego podciągu,
więc zbieżności zachodzą bez wybierania podciągu.

\textbf{Krok 5: dane z $H^{-1}$.}
Dla $f\in H^{-1}(\Omega)$ wybieramy $f_n\in L^2(\Omega)$
tak, by $\|f_n-f\|_{H^{-1}}\to0$ (gęstość $L^2$ w $H^{-1}$).
Niech $u_\eps^{(n)}$ i $u_0^{(n)}$ będą rozwiązaniami
dla $f_n$. Jednostajna koercywność daje
$\|u_\eps^{(n)}-u_\eps\|_{H_0^1}
+\|u_0^{(n)}-u_0\|_{H_0^1}\le C\|f_n-f\|_{H^{-1}}$
ze stałą niezależną od $\eps,n$. Analogicznie
$\|A(x/\eps)\nabla(u_\eps^{(n)}-u_\eps)\|_2
\le C\|f_n-f\|_{H^{-1}}$. Dla ustalonego $n$ stosujemy
poprzednie kroki, a następnie przechodzimy z $n$ do
nieskończoności; uzyskujemy wszystkie trzy zbieżności.
\end{proof}


\begin{twierdzenie}[Zbieżność energii i stabilność względem danych]
\label{thm:pde30-energy-convergence}
Niech $\Omega\subset\R^d$ będzie ograniczonym obszarem
Lipschitzowskim, a $A$ będzie symetryczna, okresowa,
mierzalna i jednostajnie eliptyczna jak w
\eqref{eq:pde30-ellipticity}. Dla $f\in H^{-1}(\Omega)$
niech $u_\eps$ i $u_0$ będą rozwiązaniami zagadnień
\eqref{eq:pde30-micro} i~\eqref{eq:pde30-macro}.
Wówczas
\begin{equation}\label{eq:pde30-energy-convergence}
 \lim_{\eps\downarrow0}\int_\Omega
 A(x/\eps)\nabla u_\eps\cdot\nabla u_\eps
 \,\mathrm d\lambda_d(x)
 =\int_\Omega A^{\mathrm{hom}}\nabla u_0\cdot\nabla u_0
 \,\mathrm d\lambda_d(x).
\end{equation}
Co więcej, jeśli $f_\eps\to f$ silnie w $H^{-1}(\Omega)$,
a $v_\eps\in H_0^1(\Omega)$ rozwiązuje
$-\diver(A(x/\eps)\nabla v_\eps)=f_\eps$,
to $v_\eps\rightharpoonup u_0$ w $H_0^1$,
$v_\eps\to u_0$ w $L^2$, a odpowiednie strumienie
zbiegają słabo w $L^2$. Zachodzi także
\eqref{eq:pde30-energy-convergence} po zastąpieniu
$u_\eps$ przez $v_\eps$.
\end{twierdzenie}
\begin{proof}
Testując równanie dla $u_\eps$ funkcją $u_\eps$,
otrzymujemy, przy oznaczeniu parowania dualnego,
\[
 \int_\Omega A(x/\eps)\nabla u_\eps\cdot\nabla u_\eps
 \,\mathrm d\lambda_d(x)=\langle f,u_\eps\rangle.
\]
Słaba zbieżność w $H_0^1$ z
\ref{thm:pde30-homogenization} oznacza zbieżność
prawej strony do $\langle f,u_0\rangle$.
Test równania granicznego funkcją $u_0$ identyfikuje
tę liczbę z prawą stroną~\eqref{eq:pde30-energy-convergence}.

Dla zmiennego obciążenia niech $u_\eps$ nadal oznacza
rozwiązanie przy ustalonym $f$. Różnica
$z_\eps=v_\eps-u_\eps$ spełnia
\[
 \lambda\|z_\eps\|_{H_0^1}\le
 \|f_\eps-f\|_{H^{-1}},
\]
gdzie w normie $H_0^1$ bierzemy normę gradientu.
W szczególności $z_\eps\to0$ silnie w $H_0^1$,
a $\|A(x/\eps)\nabla z_\eps\|_{L^2}\to0$.
Wszystkie wymienione zbieżności dziedziczymy więc
z twierdzenia~\ref{thm:pde30-homogenization}.
Tożsamość energii dla $v_\eps$ ma postać
$\int_\Omega A(x/\eps)\nabla v_\eps\cdot\nabla v_\eps\,\mathrm d\lambda_d(x)
=\langle f_\eps,v_\eps\rangle$.
Z silnej zbieżności $f_\eps$ w $H^{-1}$,
ograniczoności $v_\eps$ w $H_0^1$ i jego słabej
zbieżności do $u_0$ wynika
$\langle f_\eps,v_\eps\rangle\to\langle f,u_0\rangle$,
co kończy dowód.
\end{proof}

\begin{uwaga}[Zbieżność energii nie oznacza silnej zbieżności gradientów]
W~\eqref{eq:pde30-energy-convergence} występują
\emph{zmienne formy kwadratowe}; nie jest to
zbieżność $\|\nabla u_\eps-\nabla u_0\|_2\to0$.
Przykład jednowymiarowy z
\ref{prop:pde30-corrector1d} pokazuje, że ta ostatnia
zbieżność może być fałszywa.
\end{uwaga}


\begin{uwaga}[Dlaczego w homogenizacji występuje energia?]
\label{rem:pde-v44-homogenization-monotone}
Dla każdego $\eps>0$ macierz $A(x/\eps)$
określa liniowy operator monotoniczny
$V\to V'$ z rozdziału~\ref{chap:pde26}. Istnienie
rozwiązania wynika z koercywności; nie wymaga ono
własnej konstrukcji energii homogenizacyjnej.
Jednak w twierdzeniu~\ref{thm:pde30-energy-convergence}
\emph{symetria} $A$ pozwala interpretować
$\frac12\int_\Omega A(x/\eps)\nabla v\cdot\nabla v\,
\mathrm d\lambda_d(x)-\langle f,v\rangle$ jako funkcjonał,
którego punkt minimalny rozwiązuje równanie. Dla
macierzy niesymetrycznych forma pozostaje koercywna,
lecz jej część antysymetryczna nie jest widoczna
w funkcjonale kwadratowym. Oszacowanie energii z testu
$v=u_\eps$ nadal ma sens; nie daje samo przez się
charakteryzacji minimizacyjnej całego operatora.
\end{uwaga}


\section{Warstwy w wymiarze dwa: anizotropia z izotropowych faz}
\begin{przyklad}[Efektywna macierz materiału warstwowego]\label{ex:pde30-layer}
Niech $d=2$, $A(y)=a(y_1)I_2$, gdzie $a$ jest okresowa,
$0<\lambda\le a\le\Lambda$. Dla $j=1$ poprawka zależy
tylko od $y_1$ i spełnia
$1+\chi_1'(y_1)=a_{\mathrm{hom}}/a(y_1)$.
Dla $j=2$ funkcja $\chi_2=0$ rozwiązuje zadanie komórkowe,
ponieważ $\diver_y(Ae_2)=\partial_{y_2}a(y_1)=0$.
Stąd
\begin{equation}\label{eq:pde30-layers}
 A^{\mathrm{hom}}=
 \begin{pmatrix}
  \bigl(\int_0^1 a(y)^{-1}\,\mathrm d\lambda_1(y)\bigr)^{-1}&0\\
  0&\int_0^1a(y)\,\mathrm d\lambda_1(y)
 \end{pmatrix}.
\end{equation}
Jeśli $a$ przyjmuje wartości $1$ i $4$ na dwóch równych
połówkach komórki, to $A^{\mathrm{hom}}=\operatorname{diag}(8/5,5/2)$.
W każdej fazie materiał jest lokalnie izotropowy,
ale jego zachowanie efektywne jest anizotropowe.
\end{przyklad}

\begin{stwierdzenie}[Granice średniej harmonicznej i arytmetycznej]
\label{prop:pde30-bounds}
Dla $d=1$ i współczynnika $a$ jak wyżej zachodzi
$\lambda\le a_{\mathrm{hom}}\le\int_0^1 a(y)
\,\mathrm d\lambda_1(y)\le\Lambda$,
a równość między średnią harmoniczną i arytmetyczną
zachodzi wtedy i tylko wtedy, gdy $a$ jest stała prawie wszędzie.
\end{stwierdzenie}
\begin{proof}
Pierwsza i ostatnia nierówność wynikają z oszacowań $a$.
Nierówność Cauchy'ego--Schwarza zastosowana do
$\sqrt a$ i $1/\sqrt a$ daje
$1\le(\int_0^1a)(\int_0^1a^{-1})$.
Równość zachodzi dokładnie, gdy $\sqrt a$ jest
proporcjonalna do $1/\sqrt a$ prawie wszędzie.
\end{proof}

\section{Pułapki zbieżności i poprawki pierwszego rzędu}
\begin{kontrprzyklad}[Słaba zbieżność współczynników nie określa równania granicznego]
Na $(0,1)$ niech $a(y)$ będzie niejednorodna i okresowa.
Wtedy $a(x/\eps)\rightharpoonup\int_0^1 a$ w $L^2(0,1)$,
lecz rozwiązania równań $-(a(x/\eps)u_\eps')'=f$
zbiegną do rozwiązania równania z $a_{\mathrm{hom}}
=(\int_0^1a^{-1})^{-1}$. Obie liczby różnią się,
gdy $a$ nie jest stała. Niedozwolone jest więc przejście
do granicy osobno w iloczynie dwóch słabo zbieżnych funkcji
$a(x/\eps)$ i $u_\eps'$.
\end{kontrprzyklad}

\begin{uwaga}[Granica zakresu twierdzenia]
Twierdzenie~\ref{thm:pde30-homogenization} dowodzi zbieżności
bez oszacowania jej szybkości w wymiarach $d\ge2$. Nie
implikuje samo z siebie globalnego szacowania
$\|\nabla u_\eps-(I+\nabla\chi(x/\eps))\nabla u_0\|_2\le C\eps$:
takie oszacowanie wymaga dodatkowej regularności danych,
współczynników i brzegu oraz analizy poprawek brzegowych.
Nie wystarcza formalne podstawienie ansatzu
$u_\eps\approx u_0+\eps\sum_j\chi_j(x/\eps)\partial_j u_0$.
W wymiarze jeden otrzymaliśmy natomiast
oszacowanie gradientowe z poprawką w
stwierdzeniu~\ref{prop:pde30-corrector1d} bez dodatkowej
gładkości współczynnika $a$.
\end{uwaga}

\section{Przykład numeryczny i interpretacja strumienia}
Dla $\Omega=(0,1)$ wybieramy $a(y)=2+\sin(2\pi y)$ i $f=1$.
Wtedy $\bar a=2$, podczas gdy
$a_{\mathrm{hom}}=\bigl(\int_0^1(2+\sin(2\pi y))^{-1}
\,\mathrm d\lambda_1(y)\bigr)^{-1}=\sqrt3$.
Rozwiązanie graniczne jest równe
$u_0(x)=x(1-x)/(2\sqrt3)$, a dokładne $u_\eps$ daje
wzór~\eqref{eq:pde30-explicit-1d}. W każdym mikrośrodowisku
strumień $q_\eps=a(x/\eps)u_\eps'$ jest funkcją afiniczną
$x$, mimo że sam gradient $u_\eps'$ szybko oscyluje.

W dodatku numerycznym~\ref{sec:pde30-code} przedstawiamy
kod Plotly, który wyznacza rodzinę rozwiązań $u_\eps$,
porównuje średnie arytmetyczną i harmoniczną, ilustruje
oscylacje gradientów oraz bada zbieżność numeryczną.



\chapter[Analiza mikrolokalna]{Analiza mikrolokalna: zbiór falowy, symbole i operatory pseudoróżniczkowe}
\label{chap:pde31}

\section{Motywacja i konwencja transformaty Fouriera}
Zbiór osobliwości dystrybucji mówi, \emph{gdzie} dystrybucja nie jest
funkcją gładką. Nie mówi jednak, w \emph{jakich kierunkach} jej
oscylacje wysokoczęstotliwościowe nie zanikają. Tę brakującą informację
koduje \emph{zbiór falowy} (ang. \emph{wavefront set}).
W analizie mikrolokalnej każdemu punktowi $x$ przyporządkowujemy
kowektor $\xi\ne0$; para $(x,\xi)$ należy do zbioru falowego, jeżeli
lokalna transformata Fouriera nie maleje szybko w stożku wokół $\xi$.

W tym rozdziale $d\ge1$, a $U\subset\R^d$ jest otwarty. Dla
$h\in\mathcal S(\R^d)$ ustalamy normalizację
\begin{equation}\label{eq:pde31-fourier}
 \widehat h(\xi)=\int_{\R^d}e^{-ix\cdot\xi}h(x)
 \,\mathrm d\lambda_d(x),\qquad
 h(x)=(2\pi)^{-d}\int_{\R^d}e^{ix\cdot\xi}\widehat h(\xi)
 \,\mathrm d\lambda_d(\xi).
\end{equation}
Przedłużamy ją na dystrybucje temperowane przez dualność,
zgodnie z definicją w rozdziale~\ref{chap:pde-fourier-Rd}.
Konwencja jest dokładnie tą z równania~\eqref{eq:pde8-convention};
nie zmieniamy znaku fazy ani czynników $(2\pi)^{-d}$.
Przez $\langle\xi\rangle=(1+|\xi|^2)^{1/2}$ oznaczamy standardową
wagę Sobolewa. \emph{Stożek otwarty} $\Gamma\subset\R^d\setminus\{0\}$
jest niezmienniczy względem mnożenia przez dodatnie liczby.

\begin{definicja}[Regularność mikrolokalna]\label{def:pde31-wf}
Niech $u\in\mathcal D'(U)$, $x_0\in U$ i
$\xi_0\in\R^d\setminus\{0\}$. Mówimy, że
$(x_0,\xi_0)\notin\operatorname{WF}(u)$, jeżeli istnieją
$\chi\in C_c^\infty(U)$ z $\chi=1$ w otoczeniu $x_0$
i otwarty stożek $\Gamma$ zawierający $\xi_0$, takie że dla każdego
$N\in\N$ istnieje $C_N>0$ spełniająca
\begin{equation}\label{eq:pde31-wf-decay}
 |\widehat{\chi u}(\xi)|\le C_N\langle\xi\rangle^{-N}
 \qquad(\xi\in\Gamma).
\end{equation}
Dopełnienie zbioru tak określonych par nazywamy zbiorem falowym $u$.
Jest to domknięty podzbiór $U\times(\R^d\setminus\{0\})$,
niezmienniczy względem dodatniego skalowania drugiej współrzędnej.
\end{definicja}

\begin{lemat}[Niezależność od dalszego obcięcia]\label{lem:pde31-cutoff}
Załóżmy, że $u\in\mathcal D'(U)$, $\chi\in C_c^\infty(U)$,
$\chi=1$ w otoczeniu $x_0$, i że
$\widehat{\chi u}$ szybko maleje w stożku $\Gamma$.
Jeżeli $\eta\in C_c^\infty(U)$ ma nośnik w tym otoczeniu,
a $\Gamma'$ jest stożkiem, którego domknięcie przecięte ze sferą
jednostkową leży wewnątrz $\Gamma$, to
$\widehat{\eta u}$ szybko maleje na $\Gamma'$.
\end{lemat}
\begin{proof}
Ponieważ $\eta u=\eta(\chi u)$, wzór na transformatę iloczynu daje
$\widehat{\eta u}=(2\pi)^{-d}\widehat\eta*\widehat{\chi u}$.
Transformata dystrybucji o zwartym nośniku ma wzrost wielomianowy:
$|\widehat{\chi u}(\zeta)|\le C\langle\zeta\rangle^M$.
W całce po $\zeta$ dzielimy przestrzeń na
$|\xi-\zeta|\le\delta|\xi|$ oraz jej dopełnienie.
Dla małego $\delta>0$ pierwszy obszar, przy $\xi\in\Gamma'$,
zawiera tylko $\zeta\in\Gamma$; stosujemy tam
\eqref{eq:pde31-wf-decay} oraz całkowalność $\widehat\eta$.
Na drugim obszarze $\widehat\eta(\xi-\zeta)$ maleje szybciej
niż dowolna potęga $\langle\xi-\zeta\rangle$.
Wybierając jej wykładnik większy niż $M+d+N+1$,
otrzymujemy $O(\langle\xi\rangle^{-N})$.
Ponieważ $N$ było dowolne, dowód jest zakończony.
\end{proof}

\begin{stwierdzenie}[Rzut zbioru falowego]\label{prop:pde31-singular-support}
Dla $u\in\mathcal D'(U)$ zachodzi
$\operatorname{sing\,supp}(u)=\pi_x(\operatorname{WF}(u))$,
gdzie $\pi_x(x,\xi)=x$, a $\operatorname{sing\,supp}(u)$
oznacza zbiór punktów, w których $u$ nie jest funkcją klasy $C^\infty$
lokalnie.
\end{stwierdzenie}
\begin{proof}
Jeżeli $u$ jest gładka w otoczeniu $x_0$, to dla
$\chi\in C_c^\infty$ o nośniku w tym otoczeniu
$\chi u\in C_c^\infty$; jej transformata maleje szybko we wszystkich
kierunkach, zatem $x_0$ nie leży w rzucie zbioru falowego.

Przeciwnie, załóżmy, że wszystkie pary $(x_0,\xi)$, $|\xi|=1$,
nie należą do zbioru falowego. Ze zwartości sfery wybieramy
skończoną liczbę stożków regularności i funkcji obcinających.
Zmniejszając wspólne otoczenie $x_0$ i stosując
lemat~\ref{lem:pde31-cutoff}, otrzymujemy jedną funkcję
$\eta\in C_c^\infty(U)$, $\eta=1$ blisko $x_0$,
której transformata maleje szybko we wszystkich kierunkach.
Wtedy $\eta u$ ma pochodne wszystkich rzędów zadane
całkami bezwzględnie zbieżnymi ze wzoru odwrotnego Fouriera,
a więc $u$ jest gładka blisko $x_0$.
\end{proof}

\section{Przykłady kierunkowych osobliwości}
\begin{przyklad}[Delta Diraca i skok]\label{ex:pde31-delta}
Dla $\delta_0\in\mathcal D'(\R^d)$ mamy
$\widehat{\chi\delta_0}=\chi(0)$, zatem
$\operatorname{WF}(\delta_0)=\{(0,\xi):\xi\ne0\}$.
Dla funkcji Heaviside'a $H(x)=\chi_{(0,\infty)}(x)$
w $\R$ zbiór falowy jest taki sam: po lokalizacji
w otoczeniu zera całkowanie przez części daje
$\widehat{\chi H}(\xi)=\chi(0)/(i\xi)+O(|\xi|^{-N})$
po wybraniu $\chi=1$ blisko zera, dla każdego $N$.
Zatem brak szybkiego zaniku zachodzi w obu kierunkach.
\end{przyklad}

\begin{przyklad}[Osobliwość hiperpłaszczyzny]\label{ex:pde31-conormal}
Niech $d\ge2$, $x=(x_1,x')$ i
$u=\delta(x_1)\otimes1(x')$, tj.
$\langle u,\varphi\rangle=\int_{\R^{d-1}}\varphi(0,x')
\,\mathrm d\lambda_{d-1}(x')$ dla $\varphi\in C_c^\infty(\R^d)$.
Wtedy
\[
 \operatorname{WF}(u)=
 \{((0,x'),(\xi_1,0)):\ x'\in\R^{d-1},\ \xi_1\ne0\}.
\]
Istotnie, przy lokalizacji $\chi(x)=\chi_1(x_1)\chi_2(x')$
transformata jest iloczynem $\chi_1(0)$ i
$\widehat\chi_2(\xi')$. Maleje szybko, gdy $|\xi'|$ jest
porównywalne z $|\xi|$, lecz nie maleje w kierunku normalnym
$\xi'=0$. Lemat~\ref{lem:pde31-cutoff} pozwala przejść
do dowolnej lokalizacji.
\end{przyklad}

\begin{center}
\begin{tikzpicture}[scale=0.83,>=Latex]
  \draw[->,gray!70] (-2.8,0)--(2.8,0) node[right] {$x_1$};
  \draw[->,gray!70] (0,-2.5)--(0,2.6) node[above] {$x_2$};
  \draw[blue!65!black,very thick] (0,-2.25)--(0,2.25);
  \foreach \y in {-1.5,0,1.5}{
    \draw[-{Latex[length=3mm]},red!65!black,thick] (0,\y)--(1.25,\y);
    \draw[-{Latex[length=3mm]},red!65!black,thick] (0,\y)--(-1.25,\y);
  }
  \node[blue!65!black,below right] at (0,-2.25) {nośnik $u$};
  \node[red!65!black,above right] at (1.15,1.5) {kowektory normalne};
\end{tikzpicture}

\small Rysunek: nośnik $\delta(x_1)$ i kierunki zbioru falowego
w wymiarze dwa (schemat).
\end{center}

\begin{uwaga}[Zbiór falowy a nośnik]
Dystrybucja $\delta(x_1)$ ma nośnik będący hiperpłaszczyzną,
lecz zbiór falowy składa się z \emph{kowektorów normalnych} do niej.
Dla funkcji gładkiej o niezwartym nośniku zbiór falowy może być pusty.
Nie należy utożsamiać tych dwóch pojęć.
\end{uwaga}


\begin{twierdzenie}[Dokładna propagacja zbioru falowego przy translacji]
\label{thm:pde31-transport-wf}
Niech $u_0\in\mathcal D'(\R^d)$, $c\in\R^d$,
a dla $t\in\R$ niech $u(t)=\tau_{ct}u_0$,
gdzie $\tau_a u_0$ oznacza translację dystrybucji,
która dla funkcji ma postać $(\tau_a u_0)(x)=u_0(x-a)$.
Wtedy $t\mapsto u(t)$ jest gładkim odwzorowaniem
$\R\to\mathcal D'(\R^d)$ i rozwiązuje
$\partial_tu+c\cdot\nabla_xu=0$ przy $u(0)=u_0$.
Dla każdego $t\in\R$ zachodzi dokładna tożsamość
\begin{equation}\label{eq:pde31-transport-wf}
 \operatorname{WF}(u(t))=
 \{(x+ct,\xi):(x,\xi)\in\operatorname{WF}(u_0)\}.
\end{equation}
\end{twierdzenie}
\begin{proof}
Definiujemy translację przez
$\langle\tau_a u_0,\varphi\rangle
=\langle u_0,\varphi(\cdot+a)\rangle$.
Różniczkowanie gładkich funkcji testowych po parametrze
$t$ pokazuje, że $t\mapsto\tau_{ct}u_0$ jest gładkie
w topologii dystrybucji i że
$\partial_tu=-c\cdot\nabla_xu$.

Dla dowolnej funkcji odcinającej
$\chi\in C_c^\infty(\R^d)$ definiujemy
$\eta(y)=\chi(y+a)$. Wtedy
$\chi\tau_a u_0=\tau_a(\eta u_0)$, więc
\[
 \widehat{\chi\tau_a u_0}(\xi)
 =e^{-ia\cdot\xi}\widehat{\eta u_0}(\xi).
\]
Współczynnik wykładniczy ma moduł jeden, zatem
szybki zanik w dowolnym stożku jest równoważny
po obu stronach. Funkcja $\chi$ jest równa jeden
w otoczeniu $x_0+a$ dokładnie wtedy, gdy $\eta$
jest równa jeden blisko $x_0$.
Stąd $(x_0+a,\xi_0)\notin\operatorname{WF}(\tau_a u_0)$
wtedy i tylko wtedy, gdy
$(x_0,\xi_0)\notin\operatorname{WF}(u_0)$.
Podstawiamy $a=ct$.
\end{proof}

\begin{przyklad}[Transport punktowej osobliwości]
\label{ex:pde31-moving-delta}
W jednym wymiarze dla $u_0=\delta_0$ otrzymujemy
$u(t,x)=\delta(x-ct)$. Z
\eqref{eq:pde31-transport-wf} wynika
$\operatorname{WF}(u(t))=\{(ct,\xi):\xi\ne0\}$.
Rozpatrując $u$ jako dystrybucję na czasoprzestrzeni
$\R^2$, można obliczyć także
\[
 \operatorname{WF}(u)=
 \{((t,ct),(-c\xi,\xi)):t\in\R,\ \xi\ne0\}.
\]
Aby to sprawdzić, dla $\rho\in C_c^\infty(\R^2)$
obliczamy
$\widehat{\rho u}(\tau,\xi)=
\int_\R e^{-it(\tau+c\xi)}\rho(t,ct)\,
\mathrm d\lambda_1(t)$.
Transformata maleje szybko poza kierunkami
$\tau+c\xi=0$. W kierunkach normalnych nie może
maleć szybko dla wszystkich lokalizacji wokół danego
punktu: wybór nieujemnej $\rho$ równej jeden blisko
punktu daje na prostej $\tau=-c\xi$ dodatnią stałą,
a lemat~\ref{lem:pde31-cutoff} wyklucza szybki zanik
przy dowolnej innej lokalizacji.
Nie dowodzi to ogólnego twierdzenia o propagacji
bicharakterystyk dla operatorów o zmiennych współczynnikach.
\end{przyklad}


\begin{stwierdzenie}[Wygładzanie cieplne usuwa zbiór falowy]
\label{prop:pde-v44-heat-wavefront}
Niech $T\in\mathcal E'(\R^d)$ będzie dystrybucją
o zwartym nośniku i $t>0$. Dla jądra ciepła
$G_t(x)=(4\pi t)^{-d/2}e^{-|x|^2/(4t)}$
funkcja $G_t*T$ należy do $\mathcal S(\R^d)$,
a więc $\operatorname{WF}(G_t*T)=\varnothing$.
W szczególności dla $T=\delta_0$ zbiór falowy
w chwili początkowej jest niepusty we wszystkich
kierunkach, ale znika dla każdego $t>0$.
\end{stwierdzenie}
\begin{proof}
Z teorii dystrybucji o zwartym nośniku
$\widehat T$ jest funkcją gładką na $\R^d$;
ona sama i wszystkie jej pochodne mają co najwyżej
wzrost wielomianowy. Ponieważ
$\widehat{G_t*T}(\xi)=e^{-t|\xi|^2}\widehat T(\xi)$,
reguła Leibniza pokazuje, że iloczyn należy
do przestrzeni Schwartza. Odwrotna transformata
Fouriera również należy do $\mathcal S$.
Z definicji~\ref{def:pde31-wf} każda funkcja gładka
ma pusty zbiór falowy.
\end{proof}
\begin{uwaga}[Nieodwracalność informacji wysokoczęstotliwościowej]
Powyższy fakt opisuje kierunki osobliwości, ale nie
jest oszacowaniem stabilności problemu odwrotnego.
Na obszarze ograniczonym mnożnik cieplny
$e^{-t\lambda_k}$ na funkcjach własnych laplasjanu
powoduje silną niestabilność odtwarzania warunku
początkowego; por. stwierdzenie
\ref{prop:pde-v44-heat-inverse} w rozdziale~\ref{chap:pde32}.
\end{uwaga}


\section{Symbole i kwantyzacja Kohna--Nirenberga}
\begin{definicja}[Klasa symboli $S^m_{1,0}$]\label{def:pde31-symbol}
Dla $m\in\R$ klasą $S^m_{1,0}(\R^d\times\R^d)$
nazywamy zbiór funkcji $a\in C^\infty(\R^d\times\R^d)$,
dla których dla każdych wielowskaźników $\alpha,\beta$
istnieje $C_{\alpha\beta}$ takie, że
\begin{equation}\label{eq:pde31-symbol}
 |\partial_x^\alpha\partial_\xi^\beta a(x,\xi)|
 \le C_{\alpha\beta}\langle\xi\rangle^{m-|\beta|}
 \quad(x,\xi\in\R^d).
\end{equation}
Dla $u\in\mathcal S(\R^d)$ operator pseudoróżniczkowy
w kwantyzacji Kohna--Nirenberga określamy wzorem
\begin{equation}\label{eq:pde31-op}
 \operatorname{Op}(a)u(x)=(2\pi)^{-d}
 \int_{\R^d}e^{ix\cdot\xi}a(x,\xi)\widehat u(\xi)
 \,\mathrm d\lambda_d(\xi).
\end{equation}
Rząd $m$ opisuje skalowanie operatora przy dużych częstotliwościach.
\end{definicja}

\begin{przyklad}[Mnożenie, pochodne, potencjały Bessela]
Jeżeli $b\in C_b^\infty(\R^d)$ (wszystkie pochodne są ograniczone),
to $a(x,\xi)=b(x)$ należy do $S^0_{1,0}$ i
$\operatorname{Op}(a)u=bu$.
Dla $a(\xi)=i\xi_j$ dostajemy $\partial_{x_j}$.
Symbol $\langle\xi\rangle^m$ definiuje operator
$(I-\Delta)^{m/2}$, który ma rząd $m$.
Wszystkie trzy wzory wynikają bezpośrednio z
\eqref{eq:pde31-fourier}.
\end{przyklad}

\begin{twierdzenie}[Oszacowanie Sobolewa dla symboli o zwartym nośniku w $x$]
\label{thm:pde31-compact-symbol}
Niech $m,s\in\R$ i $a\in S^m_{1,0}$ ma nośnik w zmiennej $x$
zawarty w ustalonym zwartym zbiorze $K\subset\R^d$.
Wtedy $\operatorname{Op}(a)$ przedłuża się jednoznacznie do
ograniczonego operatora $H^s(\R^d)\to H^{s-m}(\R^d)$.
Jego norma zależy od $m,s,d,K$ i skończonej liczby
stałych z~\eqref{eq:pde31-symbol}.
\end{twierdzenie}
\begin{proof}
Niech $\widehat a_x(\eta,\xi)$ oznacza transformatę
Fouriera $a$ po zmiennej $x$. Całkowanie przez części $N$ razy
po $x$ daje
$|\widehat a_x(\eta,\xi)|\le C_N\langle\eta\rangle^{-N}
\langle\xi\rangle^m$, ponieważ nośnik w $x$ jest zwarty.
Wzór na transformatę wyniku ma postać
\[
 \widehat{\operatorname{Op}(a)u}(\zeta)
 =(2\pi)^{-d}\int_{\R^d}\widehat a_x(\zeta-\xi,\xi)
 \widehat u(\xi)\,\mathrm d\lambda_d(\xi).
\]
Wprowadzamy $v(\xi)=\langle\xi\rangle^s\widehat u(\xi)$.
Jądro operatora przekształcającego $v$ w
$\langle\zeta\rangle^{s-m}\widehat{\operatorname{Op}(a)u}(\zeta)$
jest oszacowane przez
$C_N\langle\zeta-\xi\rangle^{-N+|s-m|}$, gdyż
$\langle\zeta\rangle^{s-m}\langle\xi\rangle^{m-s}
\le C\langle\zeta-\xi\rangle^{|s-m|}$.
Wybieramy $N>d+|s-m|$. Całki jądra względem obu zmiennych
są wtedy jednostajnie ograniczone. Nierówność Schura
oraz twierdzenie Plancherela dają
$\|\operatorname{Op}(a)u\|_{H^{s-m}}\le C\|u\|_{H^s}$.
Gęstość $\mathcal S$ kończy dowód.
\end{proof}

\begin{uwaga}[Zasięg twierdzenia]
Zwartego nośnika symbolu w zmiennej $x$ nie wymaga pełne
twierdzenie o ograniczoności operatorów pseudoróżniczkowych,
lecz nasze założenie umożliwia samodzielny dowód przez jądro Fouriera.
Dla symboli ogólnych potrzebna jest dodatkowa lokalizacja
lub pełna teoria ograniczoności operatorów pseudoróżniczkowych;
nie utożsamiamy tego skrótu z jej dowodem.
\end{uwaga}

\section{Jądro poza przekątną i lokalność osobliwości}
\begin{stwierdzenie}[Gładkość jądra poza przekątną]
\label{prop:pde31-kernel}
Niech $a\in S^m_{1,0}$. Jądro dystrybucyjne
$K_a(x,y)$ operatora~\eqref{eq:pde31-op} jest funkcją gładką
na zbiorze $\{(x,y):x\ne y\}$.
\end{stwierdzenie}
\begin{proof}
Formalnie
$K_a(x,y)=(2\pi)^{-d}\int_{\R^d}
 e^{i(x-y)\cdot\xi}a(x,\xi)\,\mathrm d\lambda_d(\xi)$.
Najpierw wstawiamy obcięcie częstotliwości
$\rho(\xi/R)\in C_c^\infty$ równe jeden dla $|\xi|\le R$.
Dla $z=x-y\ne0$ zachodzi
$e^{iz\cdot\xi}=(i|z|^2)^{-1}z\cdot\nabla_\xi
 e^{iz\cdot\xi}$. Po $N$ całkowaniach przez części
pochodne symbolu maleją jak $\langle\xi\rangle^{m-N}$.
Na każdym zwartym podzbiorze $\{|x-y|>0\}$ wybieramy
$N>m+d+k$ dla pochodnych do rzędu $k$, co zapewnia
jednostajną zbieżność całek i ich pochodnych przy $R\to\infty$.
Ponieważ $k$ jest dowolne, jądro jest gładkie poza przekątną.
\end{proof}


\begin{lemat}[Przedłużenie działania na dystrybucje o zwartym nośniku]
\label{lem:pde31-distribution-op}
Niech $a\in S^m_{1,0}(\R^d\times\R^d)$, $m\in\R$.
Operator $\operatorname{Op}(a)$ określony wzorem
\eqref{eq:pde31-op} dla funkcji Schwartza ma
kanoniczne przedłużenie
$\mathcal E'(\R^d)\to\mathcal D'(\R^d)$.
Dla $u\in\mathcal E'$ i $\varphi\in C_c^\infty$ jest ono
zdefiniowane przez
\[
\langle\operatorname{Op}(a)u,\varphi\rangle
=\langle u,\operatorname{Op}(a)^t\varphi\rangle,
\]
gdzie $t$ oznacza transpozycję względem biliniowego
parowania dystrybucji z funkcją testową.
\end{lemat}
\begin{proof}
Dla $\varphi\in C_c^\infty(\R^d)$ formalna transpozycja ma
postać
\[
 (\operatorname{Op}(a)^t\varphi)(y)
 =(2\pi)^{-d}\int_{\R^d}e^{-iy\cdot\xi}
 \left(\int_{\R^d}e^{ix\cdot\xi}a(x,\xi)\varphi(x)
 \,\mathrm d\lambda_d(x)\right)\mathrm d\lambda_d(\xi).
\]
Całka wewnętrzna jest szybko malejąca jako funkcja $\xi$.
W istocie, na zwartym nośniku $\varphi$ pochodne w $x$
symbolu $a$ są oszacowane przez $C\langle\xi\rangle^m$.
Wielokrotne całkowanie przez części względem $x$ daje dla
każdego $N$ oszacowanie $C_N\langle\xi\rangle^{m-N}$.
Ten sam argument działa po różniczkowaniu w $\xi$,
a zatem całka wewnętrzna należy do $\mathcal S(\R^d)$.
Jej odwrotna transformata jest funkcją gładką w $y$,
łącznie ze wszystkimi pochodnymi. Ponieważ $u$ ma zwarty
nośnik, działa w sposób ciągły na funkcjach gładkich
w otoczeniu swojego nośnika. Powyższy wzór definiuje
zatem dystrybucję $\operatorname{Op}(a)u$; ciągłość względem
$\varphi$ wynika z oszacowań na skończoną liczbę
półnorm funkcji testowych. Dla $u\in C_c^\infty$ dostajemy
pierwotną całkę przez twierdzenie Fubiniego.
\end{proof}

\begin{wniosek}[Własność pseudolokalna]
\label{cor:pde31-pseudolocal}
Jeśli $a\in S^m_{1,0}$, a $u\in\mathcal E'(\R^d)$
jest dystrybucją o zwartym nośniku, przy czym
$\operatorname{Op}(a)u$ rozumiemy w sensie
lematu~\ref{lem:pde31-distribution-op}, to
$\operatorname{sing\,supp}(\operatorname{Op}(a)u)
\subset\operatorname{sing\,supp}(u)$.
\end{wniosek}
\begin{proof}
Niech $x_0$ leży poza zbiorem osobliwości $u$.
Rozkładamy $u=\rho u+(1-\rho)u$, gdzie $\rho=1$ blisko $x_0$
i jest wsparta w otoczeniu gładkości $u$.
Funkcja $\rho u$ jest gładka i ma zwarty nośnik,
więc jej transformata jest szybko malejąca, a
\eqref{eq:pde31-op} można różniczkować dowolnie wiele razy.
Drugie wyrażenie daje funkcję gładką przy $x_0$ dzięki
stwierdzeniu~\ref{prop:pde31-kernel}, gdyż nośnik
$(1-\rho)u$ nie przecina otoczenia $x_0$.
\end{proof}

\section{Parametryks i regularność eliptyczna}
\begin{twierdzenie}[Parametryks operatora o stałych współczynnikach]
\label{thm:pde31-parametrix}
Niech $P(D)=\sum_{|\alpha|\le m}c_\alpha\partial^\alpha$
będzie operatorem o stałych współczynnikach o rzędzie $m\ge1$.
Załóżmy, że jego jednorodny symbol główny
$P_m(i\xi)=\sum_{|\alpha|=m}c_\alpha(i\xi)^\alpha$
nie znika dla $\xi\ne0$. Wtedy istnieje
$q\in S^{-m}_{1,0}$ i funkcja $r\in C_c^\infty(\R^d)$
takie, że
\begin{equation}\label{eq:pde31-param}
 q(\xi)P(i\xi)=1-r(\xi).
\end{equation}
Dla każdej dystrybucji temperowanej $u$ zachodzi
$u=q(D)P(D)u+r(D)u$. Operator $r(D)$ jest wygładzający,
tj. przekształca $\mathcal S'$ w $C^\infty$.
Jeśli $P(D)u\in H^s(\R^d)$ oraz $u\in H^t(\R^d)$ dla pewnego $t\in\R$,
to $u\in H^{s+m}(\R^d)$.
\end{twierdzenie}
\begin{proof}
Z eliptyczności symbolu głównego i zwartości sfery wynika
$|P_m(i\xi)|\ge c|\xi|^m$. Ponieważ symbole niższych
rzędów są $O(|\xi|^{m-1})$, mamy
$|P(i\xi)|\ge c'\langle\xi\rangle^m$ przy $|\xi|\ge R$.
Wybieramy $\rho\in C^\infty$ równą zero dla $|\xi|\le R$
i jeden dla $|\xi|\ge2R$. Kładziemy
$q(\xi)=\rho(\xi)/P(i\xi)$ i $r=1-\rho$.
Wielokrotne różniczkowanie ilorazu pokazuje
$|\partial_\xi^\beta q(\xi)|\le C_\beta
\langle\xi\rangle^{-m-|\beta|}$, a więc $q\in S^{-m}_{1,0}$.
Otrzymujemy~\eqref{eq:pde31-param}.
Ponieważ $r$ ma zwarty nośnik w częstotliwościach,
$\partial^\alpha r(D)u$ jest funkcją gładką dla każdego
wielowskaźnika $\alpha$, co wynika z reprezentacji Fouriera
dystrybucji o zwartym nośniku w przestrzeni częstotliwości.
Wreszcie $q(D):H^s\to H^{s+m}$ jest ograniczony
z definicji norm Sobolewa, a $r(D):H^t\to H^{s+m}$
dla dowolnych $s,t$, gdyż mnożnik ma zwarty nośnik i jest gładki.
\end{proof}

\begin{uwaga}[Eliptyczność mikrolokalna]
Dla gładkich operatorów o zmiennych współczynnikach analogiczny
parametryks konstruuje się iteracyjnie w rachunku symboli.
Standardowa konsekwencja to
$\operatorname{WF}(Pu)\subset\operatorname{WF}(u)$ oraz
$\operatorname{WF}(u)\subset\operatorname{WF}(Pu)\cup\operatorname{Char}(P)$,
gdzie $\operatorname{Char}(P)$ jest zbiorem zer symbolu głównego
poza przekrojem zerowym. Pełny dowód ogólnego rachunku składania
symboli wykracza poza przedstawiony tu samodzielny dowód
parametryksu dla stałych współczynników; nie stosujemy tego
wyniku do operatorów z jedynie mierzalnymi współczynnikami.
\end{uwaga}

\section{Zbiór falowy Sobolewa i powierzchnie charakterystyczne}
\begin{definicja}[Zbiór falowy Sobolewa rzędu $s$]
\label{def:pde31-wfs}
Dla $s\in\R$, $u\in\mathcal D'(U)$, para $(x_0,\xi_0)$
nie należy do $\operatorname{WF}^s(u)$, jeśli istnieje
$\chi\in C_c^\infty(U)$ równa jeden blisko $x_0$ i stożek
$\Gamma\ni\xi_0$ taki, że
$\int_\Gamma\langle\xi\rangle^{2s}|\widehat{\chi u}(\xi)|^2
\,\mathrm d\lambda_d(\xi)<\infty$.
\end{definicja}

\begin{przyklad}[Próg Sobolewa delty]
Dla $\delta_0$ lokalna transformata jest stała, więc
$\delta_0\in H^s_{\mathrm{loc}}(\R^d)$ dokładnie wtedy, gdy
$s<-d/2$. W istocie w współrzędnych biegunowych
zbieżność całki jest równoważna warunkowi
$\int_1^\infty r^{2s+d-1}\,\mathrm d\lambda_1(r)<\infty$.
Dla $s\ge-d/2$ cały stożek kierunków w punkcie zero
należy do $\operatorname{WF}^s(\delta_0)$.
\end{przyklad}

\begin{przyklad}[Charakterystyki równania falowego]
Symbol główny $\partial_t^2-c^2\Delta_x$ wynosi
$-\tau^2+c^2|\xi|^2$. Jego zbiór charakterystyczny jest stożkiem
$\tau^2=c^2|\xi|^2$, $(\tau,\xi)\ne0$.
W odróżnieniu od laplasjanu, który nie ma rzeczywistych
niezerowych kowektorów charakterystycznych, operator falowy
może przenosić osobliwości wzdłuż charakterystyk.
Twierdzenie o propagacji zbioru falowego wzdłuż bicharakterystyk
jest głębsze niż uzyskana tutaj własność pseudolokalna:
nie przedstawiamy jego pełnego dowodu bez rachunku operatorów
całkowych Fouriera.
\end{przyklad}

\section{Przybliżenie optyki geometrycznej}
\begin{stwierdzenie}[Równanie eikonalne jako warunek wysokich częstotliwości]
\label{prop:pde31-eikonal}
Niech $c>0$, a $a,\varphi$ będą rzeczywistymi funkcjami klasy
$C^2$ na otwartym podzbiorze $\R^{1+d}$.
Dla $u_\eps=a\exp(i\varphi/\eps)$ zachodzi
\[
 (\partial_t^2-c^2\Delta_x)u_\eps
 =e^{i\varphi/\eps}
 \Bigl[\eps^{-2}a(c^2|\nabla_x\varphi|^2-\varphi_t^2)
 +i\eps^{-1}\mathcal T(a,\varphi)
 +(\partial_t^2-c^2\Delta_x)a\Bigr],
\]
gdzie
$\mathcal T(a,\varphi)=2\varphi_ta_t+a\varphi_{tt}
 -2c^2\nabla_x\varphi\cdot\nabla_xa-c^2a\Delta_x\varphi$.
Zanik wyrazu rzędu $\eps^{-2}$ w punktach, gdzie $a\ne0$,
wymaga równania eikonalnego $\varphi_t^2=c^2|\nabla_x\varphi|^2$.
\end{stwierdzenie}
\begin{proof}
Różniczkujemy dwukrotnie iloczyn $ae^{i\varphi/\eps}$:
każda pochodna wykładnika wnosi $i\partial\varphi/\eps$.
Po zebraniu współczynników przy $\eps^{-2}$ otrzymujemy
$a(c^2|\nabla_x\varphi|^2-\varphi_t^2)$,
a przy $i\eps^{-1}$ dokładnie $\mathcal T(a,\varphi)$.
Pozostałe wyrazy odpowiadają operatorowi falowemu
zastosowanemu do amplitudy $a$.
\end{proof}

\section{Jak stosować analizę mikrolokalną w PDE?}
Praktyczny schemat polega na identyfikacji symbolu głównego,
wyznaczeniu zbioru charakterystycznego i zbadaniu,
czy rozważane pary $(x,\xi)$ są eliptyczne.
Tam, gdzie symbol nie znika, parametryks usuwa możliwość
występowania nowych osobliwości. Na zbiorze charakterystycznym
istotna staje się geometria przepływu hamiltonowskiego symbolu.
Przed użyciem tego schematu trzeba określić \emph{gładkość
współczynników}; metody z tego rozdziału nie zastępują teorii
De Giorgiego--Nasha--Mosera dla macierzy mierzalnych.



\chapter[Problemy odwrotne i sterowanie PDE]{Problemy odwrotne i sterowanie optymalne dla równań różniczkowych cząstkowych}
\label{chap:pde32}

\section{Problem prosty, odwrotny i sens stabilności}
Zagadnienie proste polega na wyznaczeniu rozwiązania $y$ PDE
z danych, np. źródła $f$ i współczynników operatora.
W zagadnieniu odwrotnym obserwujemy częściowo $y$ i próbujemy
odzyskać nieznane $f$, współczynniki albo warunki początkowe.
Problem sterowania polega natomiast na doborze danych wejściowych
$u$, aby wymusić pożądaną własność rozwiązania.

\begin{definicja}[Stabilność zagadnienia odwrotnego]
Niech $X,Y$ będą przestrzeniami unormowanymi, $K:X\to Y$
operatorem obserwacji, a $y=Kf$ danymi bez szumu.
Zagadnienie odtworzenia $f$ jest stabilne w tych normach,
jeśli istnieje stała $C>0$ spełniająca
$\|f_1-f_2\|_X\le C\|Kf_1-Kf_2\|_Y$ dla wszystkich
$f_1,f_2$ z rozważanej klasy. Sama jednoznaczność, tj.
$\ker K=\{0\}$, nie gwarantuje stabilności.
\end{definicja}

\begin{przyklad}[Odtwarzanie źródła równania eliptycznego]
\label{ex:pde32-elliptic-inverse}
Niech $\Omega\subset\R^d$ będzie ograniczonym obszarem
Lipschitzowskim, a $A=-\Delta+I$ z jednorodnym warunkiem
Dirichleta. Dla $f\in L^2(\Omega)$ niech $y=Kf=A^{-1}f$.
Operator $K:L^2(\Omega)\to L^2(\Omega)$ jest zwarty:
rozwiązania są jednostajnie ograniczone w $H_0^1(\Omega)$,
a zanurzenie $H_0^1\hookrightarrow L^2$ jest zwarte.
Jeśli $e_k$ jest znormalizowaną funkcją własną laplasjanu
Dirichleta z wartością własną $\lambda_k\to\infty$, to
$Ke_k=(1+\lambda_k)^{-1}e_k$. Mamy
$\|e_k\|_2=1$ i $\|Ke_k\|_2\to0$: nie istnieje
jednostajnie ciągła odwrotność $K^{-1}$ na jego obrazie,
mimo że $K$ jest różnowartościowy.
\end{przyklad}

\begin{uwaga}[Odtwarzanie warunku początkowego równania ciepła]
Dla operatora $-\Delta_D$ obserwacja rozwiązania po czasie
$T>0$ jest dana przez $K_T=\exp(T\Delta_D)$.
Na funkcjach własnych $K_Te_k=e^{-T\lambda_k}e_k$.
Współczynniki wysokoczęstotliwościowe są tłumione
wykładniczo; ich odzyskanie wzmacnia nawet niewielki szum.
Jest to przykład \emph{silnej niestabilności} zagadnienia odwrotnego.
\end{uwaga}


\begin{stwierdzenie}[Zwartość obserwacji cieplnej i brak stabilnego odwrócenia]
\label{prop:pde-v44-heat-inverse}
Niech $\Omega\subset\R^d$ będzie niepustym ograniczonym
obszarem Lipschitzowskim, $H=L^2(\Omega)$ i $T>0$.
Dla laplasjanu Dirichleta niech
$K_T=e^{T\Delta_D}:H\to H$. Wtedy $K_T$ jest
zwartym, samosprzężonym, dodatnim i różnowartościowym
operatorem. Jego odwrotność $K_T^{-1}:\operatorname{Ran}(K_T)\to H$
nie jest ograniczona, przy czym $\operatorname{Ran}(K_T)$
jest gęsty w $H$.
\end{stwierdzenie}
\begin{proof}
Z teorii spektralnej w rozdziale~\ref{chap:pde18}
istnieje baza ortonormalna $(e_k)$ przestrzeni $H$
i wartości własne $0<\lambda_k\to\infty$ operatora
$-\Delta_D$. Mamy $K_Te_k=e^{-T\lambda_k}e_k$.
Obcięcia skończonego rzędu
$K_T^{(N)}f=\sum_{k=1}^Ne^{-T\lambda_k}(f,e_k)e_k$
spełniają
$\|K_T-K_T^{(N)}\|=\sup_{k>N}e^{-T\lambda_k}\to0$,
zatem $K_T$ jest zwarty. Dodatnie rzeczywiste
mnożniki dają samosprzężoność i dodatniość;
ponieważ żaden nie znika, jądro jest zerowe.
Każde $e_k$ należy do obrazu, gdyż
$e_k=K_T(e^{T\lambda_k}e_k)$, więc obraz jest gęsty.
Gdyby odwrotność była ograniczona, mielibyśmy
$1=\|e_k\|_H\le C\|K_Te_k\|_H
=C e^{-T\lambda_k}$ dla wszystkich $k$, sprzeczność.
\end{proof}
\begin{uwaga}[Wygładzanie a identyfikowalność]
Zwartość operatora obserwacji wiąże się z wygładzaniem
parabolicznym z rozdziału~\ref{chap:pde20}; stwierdzenie
\ref{prop:pde-v44-heat-wavefront} pokazuje analogiczny
mechanizm z perspektywy zbioru falowego.
Jednoznaczność danych początkowych nie zapewnia ich
stabilnej rekonstrukcji. Twierdzenie~\ref{thm:pde32-tikhonov}
dodaje regularyzację, nie odtwarza utraconej
stabilności nieuregulowanej odwrotności.
\end{uwaga}


\section{Regularyzacja Tichonowa: istnienie i równanie normalne}
\begin{twierdzenie}[Regularyzacja kwadratowa]\label{thm:pde32-tikhonov}
Niech $X,Y$ będą rzeczywistymi przestrzeniami Hilberta,
$K\in\mathcal L(X,Y)$, $y^\delta\in Y$ oraz $\alpha>0$.
Funkcjonał
\begin{equation}\label{eq:pde32-tikhonov}
 J_\alpha(f)=\frac12\|Kf-y^\delta\|_Y^2
 +\frac\alpha2\|f\|_X^2
\end{equation}
ma dokładnie jeden minimizator $f_\alpha^\delta\in X$.
Jest on jedynym rozwiązaniem równania normalnego
\begin{equation}\label{eq:pde32-normal}
 (K^*K+\alpha I)f_\alpha^\delta=K^*y^\delta.
\end{equation}
Jeżeli $\|y^\delta-y\|_Y\le\delta$ i $y=Kf^\dagger$,
gdzie $f^\dagger$ jest rozwiązaniem równania $Kf=y$
o najmniejszej normie, to przy $\alpha\downarrow0$
i $\delta/\sqrt\alpha\to0$ zachodzi
$\|f_\alpha^\delta-f^\dagger\|_X\to0$.
\end{twierdzenie}
\begin{proof}
Funkcjonał jest ściśle wypukły, ciągły i koercywny.
Każdy ciąg minimalizujący jest ograniczony; z refleksywności
przestrzeni Hilberta wybieramy słabo zbieżny podciąg.
Normy po zastosowaniu ograniczonego operatora $K$ są słabo
dolnie półciągłe, więc granica minimalizuje $J_\alpha$.
Ścisła wypukłość składnika $\alpha\|f\|^2/2$ daje
jednoznaczność. Zróżniczkowanie w kierunku $h\in X$ daje
$\langle Kf-y^\delta,Kh\rangle_Y+
\alpha\langle f,h\rangle_X=0$; to właśnie
\eqref{eq:pde32-normal}. Alternatywnie jego rozwiązywalność
wynika z twierdzenia Laxa--Milgrama.

Niech $A_0=K^*K\ge0$ i $f_\alpha^0=(A_0+\alpha I)^{-1}A_0f^\dagger$.
Z równania normalnego
\[
 f_\alpha^\delta-f^\dagger
 =-\alpha(A_0+\alpha I)^{-1}f^\dagger
 +(A_0+\alpha I)^{-1}K^*(y^\delta-y).
\]
Rozkład spektralny operatora $A_0$ i fakt
$f^\dagger\perp\ker A_0$ pokazują, że pierwszy składnik
zbiega silnie do zera: mnożniki $\alpha/(t+\alpha)$
są ograniczone przez jeden i zbiegają do zera dla $t>0$.
Dla drugiego składnika korzystamy z tożsamości
$\|(A_0+\alpha I)^{-1}K^*\|
\le\sup_{t\ge0}\sqrt t/(t+\alpha)=1/(2\sqrt\alpha)$.
Stąd błąd od szumu nie przekracza $\delta/(2\sqrt\alpha)$,
co kończy dowód.
\end{proof}

\begin{stwierdzenie}[Oszacowanie ze źródłowym warunkiem regularności]
\label{prop:pde32-source}
Przy założeniach twierdzenia~\ref{thm:pde32-tikhonov}, jeśli
$f^\dagger=K^*w$ dla pewnego $w\in Y$, to
\begin{equation}\label{eq:pde32-rate}
 \|f_\alpha^\delta-f^\dagger\|_X
 \le \frac{\sqrt\alpha}{2}\|w\|_Y
 +\frac\delta{2\sqrt\alpha}.
\end{equation}
Wybór $\alpha\asymp\delta$ daje błąd $O(\sqrt\delta)$,
przy ustalonym $\|w\|_Y$.
\end{stwierdzenie}
\begin{proof}
Z komutacji funkcji spektralnych $KK^*$ z $K^*$ mamy
$\alpha(A_0+\alpha I)^{-1}K^*w
=K^*\alpha(KK^*+\alpha I)^{-1}w$.
Jego norma nie przekracza
$\sup_{t\ge0}\alpha\sqrt t/(t+\alpha)\|w\|
=(\sqrt\alpha/2)\|w\|$.
Część odpowiadającą szumowi oszacowano w poprzednim dowodzie.
\end{proof}


\begin{twierdzenie}[Ułamkowy warunek źródłowy i tempo Tichonowa]
\label{thm:pde32-fractional-source}
Niech $X,Y$ będą rzeczywistymi przestrzeniami Hilberta,
$K\in\mathcal L(X,Y)$, $A_0=K^*K$,
$0<\nu\le1$, $w\in X$ oraz
$f^\dagger=A_0^\nu w$. Przyjmijmy
$y=Kf^\dagger$ i $\|y^\delta-y\|_Y\le\delta$.
Wtedy $f^\dagger$ jest rozwiązaniem $Kf=y$
o najmniejszej normie, a minimizator
\eqref{eq:pde32-tikhonov} spełnia
\begin{equation}\label{eq:pde32-fractional-rate}
 \|f_\alpha^\delta-f^\dagger\|_X
 \le C_\nu\alpha^\nu\|w\|_X+
 \frac{\delta}{2\sqrt\alpha},\qquad
 C_\nu=\sup_{r\ge0}\frac{r^\nu}{1+r}<\infty.
\end{equation}
Jeżeli $\delta\downarrow0$ i $\|w\|_X$ pozostaje
ograniczona, to wybór
$\alpha\asymp\delta^{2/(2\nu+1)}$
daje błąd $O(\delta^{2\nu/(2\nu+1)})$.
\end{twierdzenie}
\begin{proof}
Z twierdzenia spektralnego dla dodatniego ograniczonego
operatora samosprzężonego $A_0$ definiujemy
$A_0^\nu$. Dla $\nu>0$ obraz $A_0^\nu$ jest
prostopadły do $\ker A_0=\ker K$; dlatego
$f^\dagger$ jest rozwiązaniem o najmniejszej normie.
Rozkład błędu w dowodzie
\ref{thm:pde32-tikhonov} ma postać
\[
 f_\alpha^\delta-f^\dagger
 =-\alpha(A_0+\alpha I)^{-1}A_0^\nu w
 +(A_0+\alpha I)^{-1}K^*(y^\delta-y).
\]
Twierdzenie spektralne daje
\[
 \|\alpha(A_0+\alpha I)^{-1}A_0^\nu\|
 \le\sup_{t\ge0}\frac{\alpha t^\nu}{t+\alpha}
 =\alpha^\nu\sup_{r\ge0}\frac{r^\nu}{1+r}
 =C_\nu\alpha^\nu.
\]
Stała $C_\nu$ jest skończona dla $0<\nu\le1$.
Część pochodząca od szumu ma normę nie większą
niż $\delta/(2\sqrt\alpha)$, jak wykazano w
\ref{thm:pde32-tikhonov}. To dowodzi oszacowania.
Porównanie wykładników $\alpha^\nu$ i
$\delta\alpha^{-1/2}$ prowadzi do
$\alpha^{\nu+1/2}\asymp\delta$, skąd podane tempo.
\end{proof}

\begin{uwaga}[Zakres tempa]
Jest to oszacowanie warunkowe: bez założenia
$f^\dagger\in\operatorname{Ran}(A_0^\nu)$
nie można przypisać wszystkim danym tempa
$\delta^{2\nu/(2\nu+1)}$. Ponadto w standardowej
regularyzacji Tichonowa wykładnik $\nu>1$ nie
prowadzi tą metodą do analogicznego polepszenia
członu obciążenia (zjawisko nasycenia).
\end{uwaga}

\begin{przyklad}[Zbyt mała regularyzacja]
Dla $K_T=\exp(T\Delta_D)$ mamy
$\sigma_k=e^{-T\lambda_k}$. Na obserwacji pojedynczego
modu $y^\delta=\sigma_ke_k+\delta e_k$ nieuregulowana
odwrotność mnoży szum przez $1/\sigma_k=e^{T\lambda_k}$.
Regularyzacja Tichonowa zastępuje ten mnożnik przez
$\sigma_k/(\sigma_k^2+\alpha)$, którego moduł jest
ograniczony przez $1/(2\sqrt\alpha)$. W celu uzyskania zbieżności do rozwiązania dokładnego przy
zanikającym szumie trzeba dobierać $\alpha=\alpha(\delta)$;
parametr stały ogranicza wpływ szumu, ale pozostawia
niezanikający błąd regularyzacji.
\end{przyklad}

\section{Sterowanie rozłożone w obszarze: operator stanu}
Niech $\Omega\subset\R^d$ będzie ograniczonym obszarem
Lipschitzowskim, $V=H_0^1(\Omega)$ i $H=L^2(\Omega)$.
Przyjmujemy mierzalną symetryczną macierz
$A\in L^\infty(\Omega;\R^{d\times d})$
i $\xi^{\mathsf T}A(x)\xi\ge\lambda|\xi|^2$ dla
każdego $\xi$ i prawie każdego $x$, z $\lambda>0$.
Niech $c\in L^\infty(\Omega)$, $c\ge0$.
Forma
$\mathfrak a(y,v)=\int_\Omega
(A\nabla y\cdot\nabla v+cyv)\,\mathrm d\lambda_d(x)$
jest koercywna i symetryczna na $V$.
Niech $f\in V'$ oraz $u\in H$. Przez $y(u)\in V$
oznaczamy jedyne rozwiązanie
\begin{equation}\label{eq:pde32-state}
 \mathfrak a(y(u),v)=\langle f,v\rangle_{V',V}
 +\int_\Omega uv\,\mathrm d\lambda_d(x)
 \quad\text{dla każdego }v\in V.
\end{equation}
Odwzorowanie $u\mapsto y(u)$ jest afiniczne i ciągłe
$H\to V$. Kładziemy $S(u)=y(u)-y(0)$; jest to liniowy
operator rozwiązania dla źródła $u$.

\begin{twierdzenie}[Istnienie sterowania optymalnego]\label{thm:pde32-control}
Niech $y_d\in H$, $\alpha>0$ i $U_{\mathrm{ad}}\subset H$
będzie niepustym, domkniętym i wypukłym podzbiorem.
Wtedy problem
\begin{equation}\label{eq:pde32-opt}
 \min_{u\in U_{\mathrm{ad}}}
 \mathcal J(u),\qquad
 \mathcal J(u)=\frac12\|y(u)-y_d\|_H^2
 +\frac\alpha2\|u\|_H^2
\end{equation}
ma dokładnie jedno rozwiązanie optymalne $\bar u$.
\end{twierdzenie}
\begin{proof}
Dla dowolnego ciągu minimalizującego $(u_n)$ nierówność
$\alpha\|u_n\|_H^2/2\le\mathcal J(u_n)$ daje ograniczoność.
Istnieje słabo zbieżny podciąg $u_n\rightharpoonup\bar u$ w $H$.
Każdy domknięty wypukły podzbiór przestrzeni Hilberta jest
słabo domknięty, więc $\bar u\in U_{\mathrm{ad}}$.
Ciągłość liniowa $S$ daje słabą zbieżność
$y(u_n)\rightharpoonup y(\bar u)$ w $H$.
Słaba dolna półciągłość norm daje
$\mathcal J(\bar u)\le\liminf\mathcal J(u_n)$.
Składnik z $\alpha>0$ zapewnia ścisłą wypukłość
$\mathcal J$, a wraz z wypukłością $U_{\mathrm{ad}}$
także jednoznaczność minimizatora.
\end{proof}

\section{Równanie sprzężone i warunki optymalności}
\begin{twierdzenie}[Warunek pierwszego rzędu i rzut na sterowania dopuszczalne]
\label{thm:pde32-adjoint}
Niech $\bar u$ będzie rozwiązaniem~\eqref{eq:pde32-opt},
$\bar y=y(\bar u)$, a $\bar p\in V$ niech będzie
jedynym rozwiązaniem równania sprzężonego
\begin{equation}\label{eq:pde32-adj}
 \mathfrak a(v,\bar p)
 =\int_\Omega(\bar y-y_d)v\,\mathrm d\lambda_d(x)
 \qquad(v\in V).
\end{equation}
Wówczas $\bar u$ charakteryzuje nierówność wariacyjna
\begin{equation}\label{eq:pde32-variational-control}
 \int_\Omega(\bar p+\alpha\bar u)(v-\bar u)
 \,\mathrm d\lambda_d(x)\ge0
 \quad\text{dla każdego }v\in U_{\mathrm{ad}},
\end{equation}
a równoważnie
$\bar u=P_{U_{\mathrm{ad}}}(-\bar p/\alpha)$,
gdzie $P_{U_{\mathrm{ad}}}$ jest rzutem metrycznym na
niepusty domknięty wypukły podzbiór przestrzeni Hilberta.
Jeżeli $U_{\mathrm{ad}}=\{v\in L^2(\Omega):a\le v\le b
\text{ prawie wszędzie}\}$ dla liczb $a<b$,
to $\bar u(x)=\min\{b,\max\{a,-\bar p(x)/\alpha\}\}$
dla prawie każdego $x\in\Omega$.
\end{twierdzenie}
\begin{proof}
Dla $h\in H$ pochodna stanu $z=S h$ spełnia
$\mathfrak a(z,v)=\int_\Omega hv\,\mathrm d\lambda_d(x)$.
Różniczkując funkcjonał otrzymujemy
$\mathcal J'(\bar u)h=(\bar y-y_d,z)_H+
\alpha(\bar u,h)_H$. Z równań dla $z$ i $\bar p$
wynika
$(\bar y-y_d,z)_H=\mathfrak a(z,\bar p)=
(h,\bar p)_H$. Stąd $\mathcal J'(\bar u)h=
(\bar p+\alpha\bar u,h)_H$.
Dla $v\in U_{\mathrm{ad}}$ funkcja
$t\mapsto\mathcal J(\bar u+t(v-\bar u))$ ma minimum
w $t=0$ na $[0,1]$, co daje~\eqref{eq:pde32-variational-control}.
Odwrotnie, nierówność i wypukłość $\mathcal J$ zapewniają minimum.
Charakterystyka rzutu Hilbertowskiego
$z=P_Cw\iff (z-w,v-z)_H\ge0$ dla każdego $v\in C$
daje formułę z $w=-\bar p/\alpha$.
Dla ograniczeń punktowych minimalizacja
$\|v-w\|_{L^2}^2$ rozdziela się punktowo,
co prowadzi do obcięcia wartości do przedziału $[a,b]$.
\end{proof}

\begin{uwaga}[Dlaczego równanie sprzężone jest użyteczne?]
Bez zmiennej sprzężonej obliczenie gradientu kosztu
wymagałoby rozwiązania nowego PDE dla każdego kierunku $h$.
Równanie~\eqref{eq:pde32-adj} wymaga jednego dodatkowego
rozwiązania PDE i daje gradient jako funkcję $\alpha u+p$.
To podstawowy mechanizm metod numerycznych sterowania.
\end{uwaga}

\section{Sterowanie równaniem ciepła i czas wsteczny}
\begin{stwierdzenie}[Równanie sprzężone dla parabolicznego sterowania rozłożonego]
\label{prop:pde32-heat-adjoint}
Niech $T>0$, $\Omega$ będzie ograniczonym obszarem
Lipschitzowskim, $u_0\in L^2(\Omega)$, $y_d\in
L^2((0,T)\times\Omega)$ oraz $\alpha>0$.
Dla $v\in L^2((0,T)\times\Omega)$ niech $y(v)$ będzie
słabym rozwiązaniem
$y_t-\Delta y=v$ w $(0,T)\times\Omega$,
$y=0$ na $(0,T)\times\partial\Omega$ i $y(0)=u_0$.
Koszt
$\mathcal J_T(v)=\frac12\|y(v)-y_d\|_{L^2((0,T)\times\Omega)}^2
+\frac\alpha2\|v\|_{L^2((0,T)\times\Omega)}^2$
ma pochodną
$\mathcal J_T'(v)h=
\int_0^T\!\int_\Omega(\alpha v+p)h
\,\mathrm d\lambda_d(x)\,\mathrm d\lambda_1(t)$,
gdzie $p$ jest rozwiązaniem równania sprzężonego
\[
 -p_t-\Delta p=y(v)-y_d,
 \qquad p(T)=0,\qquad p|_{\partial\Omega}=0.
\]
\end{stwierdzenie}
\begin{proof}
Pochodna stanu $z$ w kierunku $h$ spełnia
$z_t-\Delta z=h$, $z(0)=0$, z jednorodnym warunkiem brzegowym.
Słabe rozwiązania należą do
$L^2(0,T;H_0^1)$, a ich pochodne czasowe do
$L^2(0,T;H^{-1})$, więc wolno użyć formuły całkowania
przez części w trójce Gelfanda, udowodnionej w rozdziale 20.
Mnożąc równanie $z$ przez $p$ i równanie $p$ przez $z$,
a następnie odejmując, uzyskujemy
\[
 \int_0^T\!\int_\Omega(y(v)-y_d)z
 \,\mathrm d\lambda_d(x)\,\mathrm d\lambda_1(t)
 =\int_0^T\!\int_\Omega ph
 \,\mathrm d\lambda_d(x)\,\mathrm d\lambda_1(t).
\]
Składniki brzegowe w czasie znikają z $z(0)=0$ i $p(T)=0$.
Różniczkowanie kosztu oraz powyższa tożsamość
dają tezę.
\end{proof}

\section{Odtwarzanie współczynnika i granice identyfikowalności}
\begin{stwierdzenie}[Pochodna względem współczynnika dyfuzji]
\label{prop:pde32-coeff-gradient}
Niech $\Omega$ będzie ograniczonym obszarem Lipschitzowskim,
$0<\lambda\le a\le\Lambda$ prawie wszędzie,
$a\in L^\infty(\Omega)$, $f\in H^{-1}(\Omega)$.
Niech $y(a)\in H_0^1(\Omega)$ będzie jedynym rozwiązaniem
$-\diver(a\nabla y(a))=f$.
Dla $h\in L^\infty(\Omega)$ i $|\eps|$ dostatecznie małych
niech $a_\eps=a+\eps h$. Wtedy istnieje pochodna kierunkowa
$z=\lim_{\eps\to0}(y(a_\eps)-y(a))/\eps$ w $H_0^1(\Omega)$
i spełnia
\begin{equation}\label{eq:pde32-coeff-linear}
 \int_\Omega a\nabla z\cdot\nabla v
 \,\mathrm d\lambda_d(x)
 =-\int_\Omega h\nabla y(a)\cdot\nabla v
 \,\mathrm d\lambda_d(x)\qquad(v\in H_0^1(\Omega)).
\end{equation}
Jeżeli $y_d\in L^2(\Omega)$,
$\mathcal F(a)=\frac12\|y(a)-y_d\|_2^2$,
a $p$ spełnia
$-\diver(a\nabla p)=y(a)-y_d$ w $H^{-1}$, to
$\mathcal F'(a)h=-\int_\Omega
h\nabla y(a)\cdot\nabla p\,\mathrm d\lambda_d(x)$.
\end{stwierdzenie}
\begin{proof}
Z jednolitej eliptyczności $a_\eps$ wynika, że
$y(a_\eps)$ jest jednostajnie ograniczona w $H_0^1$.
Odejmując równania i testując różnicą rozwiązań
otrzymujemy
$\|y(a_\eps)-y(a)\|_{H_0^1}\le
C|\eps|\|h\|_\infty\|f\|_{H^{-1}}$.
Iloraz różnicowy $z_\eps$ spełnia
$\int_\Omega a\nabla z_\eps\cdot\nabla v\,\mathrm d\lambda_d(x)
=-\int_\Omega h\nabla y(a_\eps)\cdot\nabla v\,\mathrm d\lambda_d(x)$.
Prawa strona zbiega silnie w $H^{-1}$ do funkcjonału
$v\mapsto-\int_\Omega h\nabla y(a)\cdot\nabla v
\,\mathrm d\lambda_d(x)$;
koercywność daje silną zbieżność $z_\eps\to z$ w $H_0^1$.
Z definicji $p$ i~\eqref{eq:pde32-coeff-linear}
$\mathcal F'(a)h=(y(a)-y_d,z)_2
=\int_\Omega a\nabla p\cdot\nabla z\,\mathrm d\lambda_d(x)
=-\int_\Omega h\nabla y(a)\cdot\nabla p\,\mathrm d\lambda_d(x)$.
\end{proof}


\begin{wniosek}[Różniczkowalność w normie $L^\infty$]
\label{cor:pde32-frechet}
Przy założeniach stwierdzenia~\ref{prop:pde32-coeff-gradient}
operator rozwiązania $a\mapsto y(a)$ jest różniczkowalny
w sensie Frécheta jako odwzorowanie z otwartego zbioru
$\{a\in L^\infty(\Omega):\operatorname*{ess\,inf}a>0\}$
do $H_0^1(\Omega)$. Jego pochodna w kierunku $h\in L^\infty$
jest jedynym rozwiązaniem~\eqref{eq:pde32-coeff-linear}.
Dla $\|h\|_\infty\le\lambda/2$ istnieje stała $C$
zależna wyłącznie od $\lambda$, dla której
\[
 \|y(a+h)-y(a)-z_h\|_{H_0^1}
 \le C\|h\|_\infty^2\|f\|_{H^{-1}}.
\]
\end{wniosek}
\begin{proof}
Oznaczmy $y_h=y(a+h)$, $y=y(a)$ i $\delta y=y_h-y$.
Odejmowanie równań stanów daje
$\int_\Omega a\nabla\delta y\cdot\nabla v
=-\int_\Omega h\nabla y_h\cdot\nabla v$.
Ponieważ $a+h\ge\lambda/2$, z oszacowań energetycznych
wynikają $\|y_h\|_{H_0^1}\le2\lambda^{-1}\|f\|_{H^{-1}}$
i $\|\delta y\|_{H_0^1}\le
2\lambda^{-2}\|h\|_\infty\|f\|_{H^{-1}}$.
Pochodna kandydacka $z_h$ spełnia
$\int_\Omega a\nabla z_h\cdot\nabla v
=-\int_\Omega h\nabla y\cdot\nabla v$.
Dla reszty $r_h=\delta y-z_h$ otrzymujemy więc
\[
 \int_\Omega a\nabla r_h\cdot\nabla v
 =-\int_\Omega h\nabla\delta y\cdot\nabla v.
\]
Testując przez $v=r_h$ i wykorzystując koercywność,
otrzymujemy $\|r_h\|_{H_0^1}\le
2\lambda^{-3}\|h\|_\infty^2\|f\|_{H^{-1}}$.
Zależność $h\mapsto z_h$ jest liniowa i ograniczona,
a reszta jest kwadratowa względem normy $L^\infty$.
To dowodzi różniczkowalności Frécheta.
\end{proof}

\begin{kontrprzyklad}[Brak identyfikowalności z jednej obserwacji]
Dla $f=0$ i jednorodnych warunków Dirichleta rozwiązaniem
$-\diver(a\nabla y)=0$ jest $y=0$ dla dowolnego dodatniego,
ograniczonego współczynnika $a$. Zatem obserwacja całej
funkcji $y$ nie wyznacza $a$ bez odpowiedniego pobudzenia.
To nie jest wada regularyzacji, lecz strukturalny brak informacji.
\end{kontrprzyklad}

\section{Sterowanie skończenie wieloma modami równania ciepła}
\begin{twierdzenie}[Dokładny koszt sterowania pojedynczego modu]
\label{thm:pde32-gramian}
Niech $\lambda>0$, $T>0$ i $a\in\R$. Spośród wszystkich
$h\in L^2(0,T)$, dla których rozwiązanie
$z'(t)+\lambda z(t)=h(t)$, $z(0)=0$, spełnia $z(T)=a$,
minimalną normę ma
\[
 h_*(s)=\frac{a\,e^{-\lambda(T-s)}}{W_\lambda(T)},
 \qquad W_\lambda(T)=\frac{1-e^{-2\lambda T}}{2\lambda}.
\]
Ponadto $\int_0^T|h_*(s)|^2\,\mathrm d\lambda_1(s)
=a^2/W_\lambda(T)$.
\end{twierdzenie}
\begin{proof}
Ze wzoru wariacji stałej
$z(T)=\int_0^T e^{-\lambda(T-s)}h(s)
\,\mathrm d\lambda_1(s)$.
Nierówność Cauchy'ego--Schwarza daje
$|a|^2\le W_\lambda(T)\|h\|_{L^2(0,T)}^2$.
Równość zachodzi dokładnie dla $h$ proporcjonalnej
do $s\mapsto e^{-\lambda(T-s)}$.
Dobrany współczynnik proporcjonalności wynika z warunku
$z(T)=a$ i daje powyższy wzór.
\end{proof}

\begin{uwaga}[Zasięg twierdzenia o sterowalności]
Dla skończonej liczby modów laplasjanu wybieramy sterowania
wzdłuż odpowiadających im funkcji własnych i sumujemy koszty.
Jednak $W_{\lambda_k}(T)^{-1}\sim 2\lambda_k$,
więc ogólny cel $\sum_k a_ke_k\in L^2$ może wymagać
nieskończonego kosztu. Uzyskana dokładna sterowalność
każdego skończonego zestawu modów \emph{nie oznacza}
dokładnej sterowalności w całej przestrzeni $L^2$.
Na tym polega fundamentalna różnica pomiędzy problemem
skończenie wymiarowym a równaniem parabolicznym.
\end{uwaga}

\section{Podsumowanie metod i ograniczeń}
Regularyzacja stabilizuje problem odwrotny, lecz nie przywraca
informacji utraconej przez operator obserwacji.
W sterowaniu optymalnym istnienie i jednoznaczność wynikają
z wypukłości, a równanie sprzężone pozwala obliczyć pochodną
kosztu bez rozwiązywania osobnego PDE dla każdego kierunku.
Ogólna teoria identyfikacji współczynników z danych brzegowych
(np. problem Calderóna), nierówności obserwowalności,
sterowalność dokładna i przybliżona oraz problemy optymalizacji
nieliniowych układów PDE są bardziej zaawansowane;
nie przypisujemy im dowodów uzyskanych tylko dla
kwadratowego zadania eliptycznego lub skończonego zestawu modów.


\appendix
\chapter{Eksperymenty numeryczne do rozdziałów 2--32}
\label{chap:pde-code}

Każdy poniższy listing jest samodzielnym programem w języku Python,
korzystającym z bibliotek \texttt{numpy} i \texttt{plotly}.
Wizualizacje są interaktywne i zapisują się do plików \texttt{.html};
pliki zawierają również lokalną kopię biblioteki JavaScript Plotly,
więc otwierają się bez połączenia z internetem.
Programy uruchamia się przez \texttt{python nazwa\_pliku.py}.
W załączonym projekcie wszystkie 29 listingów zapisano także jako
osobne pliki w katalogu \texttt{scripts/listings/}. Aby uruchomić
przykład, należy wykonać odpowiadający mu plik \texttt{.py};
kompletną tabelę nazw umieszczono w pliku
\texttt{INSTRUKCJA\_URUCHOMIENIA.md}. Eksperymenty korzystają
z NumPy w wersji co najmniej 2.0 oraz Plotly. Część przykładów,
w szczególności wizualizacja 3D wiru Taylora--Greena, wymaga ponadto
SciPy i Matplotlib. Wykonanie programu testuje jego działanie
numeryczne, ale samo nie stanowi dowodu zbieżności schematu.

\section{Charakterystyki, równanie Poissona i ciepło}
W kodzie wykorzystujemy rozwiązanie charakterystyk
równania Burgersa dla danych $-\tanh x$, tłumienie
harmonicznych modów przez $r^k$ oraz jawny profil Gaussa.

\begin{lstlisting}[style=pythonpde,caption={Burgers, funkcje harmoniczne i dyfuzja: Plotly},label={code:pde-old}]
"""Interaktywne wykresy PDE: charakterystyki, profile, dyfuzja i Poisson.
Wymagania: numpy, plotly. Zapisuje samodzielne pliki HTML (offline).
"""
from pathlib import Path
import numpy as np
import plotly.graph_objects as go

OUT = Path(__file__).resolve().parent


def save(fig: go.Figure, name: str) -> None:
    fig.update_layout(template="plotly_white")
    fig.write_html(OUT / name, include_plotlyjs=True, full_html=True)


def data(x):
    return -np.tanh(x)


def position(t, xi):
    return xi + t * data(xi)


def burgers(t, x):
    xi = np.linspace(-8.0, 8.0, 100001)
    X = position(t, xi)
    if not np.all(np.diff(X) > 0):
        raise ValueError("Brak rozwiazania klasycznego dla tego czasu")
    return data(np.interp(x, X, xi))


def characteristic_plot():
    fig = go.Figure()
    ts = np.linspace(0, 0.99, 250)
    for xi in np.linspace(-2.4, 2.4, 13):
        fig.add_trace(go.Scatter(x=position(ts, xi), y=ts,
                                 mode="lines", name=f"xi={xi:.1f}"))
    fig.update_layout(title="Charakterystyki Burgersa", xaxis_title="x",
                      yaxis_title="t")
    save(fig, "burgers_characteristics.html")


def profile_plot():
    x = np.linspace(-3.0, 3.0, 1001)
    fig = go.Figure()
    for t in (0.0, 0.5, 0.9, 0.99):
        fig.add_trace(go.Scatter(x=x, y=burgers(t, x),
                                 mode="lines", name=f"t={t:.2f}"))
    fig.update_layout(title="Burgers: profile przed utrata regularnosci",
                      xaxis_title="x", yaxis_title="u(t,x)")
    save(fig, "burgers_profiles.html")


def poisson_modes():
    k = np.arange(0, 30)
    fig = go.Figure()
    for r in (0.2, 0.5, 0.8):
        fig.add_trace(go.Scatter(x=k, y=r**k, mode="lines+markers",
                                 name=f"r={r}"))
    fig.update_layout(title="Wygaszanie modow harmonicznych",
                      xaxis_title="czestosc k", yaxis_title="amplituda r^k",
                      yaxis_type="log")
    save(fig, "poisson_modes.html")


def heat_gaussian_plot():
    x = np.linspace(-4, 4, 1200)
    fig = go.Figure()
    for t in (0.0, 0.10, 0.40, 1.20):
        u = np.exp(-x**2 / (1 + 4*t)) / np.sqrt(1 + 4*t)
        fig.add_trace(go.Scatter(x=x, y=u, mode="lines", name=f"t={t:.2f}"))
    fig.update_layout(title="Dokladne rozwiazanie Gaussowskie (kappa=1)",
                      xaxis_title="x", yaxis_title="u(t,x)")
    save(fig, "heat_gaussian_profiles.html")


if __name__ == "__main__":
    characteristic_plot()
    profile_plot()
    poisson_modes()
    heat_gaussian_plot()
    for t in (0.0, 0.5, 0.9, 0.99):
        print(f"t={t:.2f}, pochodna w x=0: {-1/(1-t):.6f}")
\end{lstlisting}

\section{Schemat falowy i kontrola dokładności}
\begin{lstlisting}[style=pythonpde,caption={Schemat różnicowy równania falowego: Plotly i test błędu},label={code:pde-wave}]
"""Jawny schemat falowy: interaktywny wykres Plotly i blad max.
Wymagania: numpy, plotly. Zapisuje plik HTML do katalogu skryptu.
"""
from pathlib import Path
import numpy as np
import plotly.graph_objects as go

c, L, N = 1.0, 1.0, 300
dx = L / N
lam = 0.8
dt = lam * dx / c
x = np.linspace(0.0, L, N + 1)
T = 0.8
steps = int(round(T / dt))


def exact(t):
    return (np.cos(2*np.pi*c*t/L)*np.sin(2*np.pi*x/L)
            + 0.4*np.cos(7*np.pi*c*t/L)*np.sin(7*np.pi*x/L))


u_prev = exact(0.0)
u = u_prev.copy()
u[1:-1] += 0.5*lam**2*(u_prev[2:] - 2*u_prev[1:-1] + u_prev[:-2])
u[0] = u[-1] = 0.0
snapshots = {0: u_prev.copy()}
for m in range(1, steps):
    u_next = np.zeros_like(u)
    u_next[1:-1] = (2*u[1:-1] - u_prev[1:-1]
                    + lam**2*(u[2:] - 2*u[1:-1] + u[:-2]))
    u_prev, u = u, u_next
    if m + 1 in (steps//3, 2*steps//3, steps):
        snapshots[m + 1] = u.copy()

fig = go.Figure()
for m, values in snapshots.items():
    fig.add_trace(go.Scatter(x=x, y=values, mode="lines", name=f"t={m*dt:.3f}"))
fig.update_layout(template="plotly_white", title="Struna: lambda=0.8",
                  xaxis_title="x", yaxis_title="u(t,x)")
fig.write_html(Path(__file__).with_name("wave_leapfrog.html"),
               include_plotlyjs=True, full_html=True)
print("CFL:", lam, "blad max:", np.max(np.abs(u - exact(steps*dt))))
\end{lstlisting}

\section{Szeregi Fouriera i ewolucje modalne}
\begin{lstlisting}[style=pythonpde,caption={Średnie Fejéra, Gibbs oraz powierzchnie ewolucji: Plotly},label={code:pde-fourier}]
"""Szeregi Fouriera i PDE na odcinku, wykresy interaktywne Plotly.
Wymagania: numpy, plotly. Zapisuje samodzielne pliki HTML.
"""
from pathlib import Path
import numpy as np
import plotly.graph_objects as go

OUT = Path(__file__).resolve().parent


def save(fig, name):
    fig.update_layout(template="plotly_white")
    fig.write_html(OUT / name, include_plotlyjs=True, full_html=True)


def fourier_vs_fejer():
    # f(x) = sign(x), x in (-pi, pi). W wartosciach skoku: f=0.
    x = np.linspace(-np.pi, np.pi, 1400)
    f = np.sign(x)
    fig = go.Figure([go.Scatter(x=x, y=f, name="dane skokowe")])
    for N in (5, 15, 45):
        ks = np.arange(1, N+1, 2)
        terms = (4/np.pi) * np.sin(np.outer(ks, x))/ks[:, None]
        fourier = terms.sum(axis=0)
        fejer = ((1-ks/(N+1))[:, None] * terms).sum(axis=0)
        fig.add_trace(go.Scatter(x=x, y=fourier, name=f"S_{N}"))
        fig.add_trace(go.Scatter(x=x, y=fejer, name=f"sigma_{N}"))
    fig.update_layout(title="Zjawisko Gibbsa i srednie Fejera",
                      xaxis_title="x", yaxis_title="wartosc")
    save(fig, "fourier_vs_fejer.html")


def heat_modes():
    L, kappa = 1.0, 0.06
    x = np.linspace(0.0, L, 500)
    t = np.linspace(0.0, 0.6, 100)
    X, T = np.meshgrid(x, t)
    U = (np.exp(-kappa*np.pi**2*T)*np.sin(np.pi*X)
         + 0.45*np.exp(-kappa*(5*np.pi)**2*T)*np.sin(5*np.pi*X))
    fig = go.Figure(data=go.Surface(x=X, y=T, z=U, colorscale="Viridis"))
    fig.update_layout(title="Rownanie ciepla: tlumienie wysokich modow",
                      scene=dict(xaxis_title="x", yaxis_title="t",
                                 zaxis_title="u(t,x)"))
    save(fig, "heat_modes.html")


def wave_modes():
    L, c = 1.0, 1.0
    x = np.linspace(0.0, L, 500)
    t = np.linspace(0.0, 2.0, 140)
    X, T = np.meshgrid(x, t)
    U = (np.cos(np.pi*c*T/L)*np.sin(np.pi*X/L)
         + 0.45*np.cos(5*np.pi*c*T/L)*np.sin(5*np.pi*X/L))
    fig = go.Figure(data=go.Surface(x=X, y=T, z=U, colorscale="Viridis"))
    fig.update_layout(title="Rownanie falowe: utrzymywanie sie modow",
                      scene=dict(xaxis_title="x", yaxis_title="t",
                                 zaxis_title="u(t,x)"))
    save(fig, "wave_modes.html")


if __name__ == "__main__":
    fourier_vs_fejer()
    heat_modes()
    wave_modes()
    print("Wykresy Plotly zapisane w", OUT)
\end{lstlisting}

\section{Przybliżone jedynki, skok i jądro cieplne}
\begin{lstlisting}[style=pythonpde,caption={Dystrybucje: przybliżone jedynki, funkcja Heaviside'a i jądro cieplne (Plotly)},label={code:pde-distributions}]
"""Przyblizone jedynki, pochodne slabego skoku i jadro ciepla.
Wymagania: numpy, plotly. Eksportuje interaktywne HTML.
"""
from pathlib import Path
import numpy as np
import plotly.graph_objects as go

OUT = Path(__file__).resolve().parent
x = np.linspace(-2.0, 2.0, 10001)
dx = x[1] - x[0]


def bump(z):
    """Gladki profil o zwartym nosniku [-1,1], bez normalizacji."""
    z = np.asarray(z)
    y = np.zeros_like(z, dtype=float)
    mask = np.abs(z) < 1.0
    y[mask] = np.exp(-1.0 / (1.0 - z[mask] ** 2))
    return y


zq = np.linspace(-1.0, 1.0, 5001)
MASS = np.trapezoid(bump(zq), zq)


def mollifier(epsilon):
    return bump(x / epsilon) / (epsilon * MASS)


def save(fig, name):
    fig.update_layout(template="plotly_white", hovermode="x unified")
    fig.write_html(OUT / name, include_plotlyjs=True, full_html=True)


def delta_and_heaviside():
    fig_delta, fig_h = go.Figure(), go.Figure()
    eps_values = (0.5, 0.25, 0.10)
    for eps in eps_values:
        rho = mollifier(eps)
        fig_delta.add_trace(go.Scatter(x=x, y=rho, mode="lines",
                                       name=f"epsilon={eps}"))
        # H * rho = integral od -nieskonczonosci do x z rho(s) ds
        h = np.r_[0.0, np.cumsum((rho[:-1] + rho[1:]) * dx / 2)]
        fig_h.add_trace(go.Scatter(x=x, y=h, mode="lines",
                                   name=f"epsilon={eps}"))
        # Test dystrybucyjny: int rho_epsilon(x) * phi(x) dx -> phi(0)=1
        phi = np.exp(-x*x)
        pairing = np.trapezoid(rho * phi, x)
        print(f"epsilon={eps:5.2f}, masa={np.trapezoid(rho,x):.7f}, "
              f"<rho_epsilon,phi>={pairing:.7f}, phi(0)=1")
    fig_h.add_trace(go.Scatter(x=x, y=(x > 0).astype(float),
                               mode="lines", name="H (skok)",
                               line=dict(dash="dot", color="gray")))
    fig_delta.update_layout(title="Gladkie aproksymacje delty Diraca",
                            xaxis_title="x", yaxis_title="rho_epsilon(x)")
    fig_h.update_layout(title="Wygladzona funkcja Heaviside'a",
                        xaxis_title="x", yaxis_title="(H * rho_epsilon)(x)")
    save(fig_delta, "pde7_delta_plotly.html")
    save(fig_h, "pde7_heaviside_plotly.html")


def heat_kernel():
    kappa = 0.15
    fig = go.Figure()
    for t in (0.01, 0.04, 0.2):
        gamma = np.exp(-x*x/(4*kappa*t)) / np.sqrt(4*np.pi*kappa*t)
        fig.add_trace(go.Scatter(x=x, y=gamma, name=f"t={t}"))
        print(f"t={t:0.3f}, masa na [-2,2]={np.trapezoid(gamma,x):.7f}")
    fig.update_layout(title="Jadro cieplne w 1D jako aproksymacja delty",
                      xaxis_title="x", yaxis_title="Gamma(t,x)")
    save(fig, "pde7_heat_kernel_plotly.html")


if __name__ == "__main__":
    delta_and_heaviside()
    heat_kernel()
\end{lstlisting}


\section[Transformata, Parseval i mnożniki]{Transformata na $\R^d$, tożsamość Parsevala i mnożniki}
\begin{lstlisting}[style=pythonpde,
caption={Transformata Gaussa, dyskretna tożsamość Parsevala, ewolucje spektralne oraz normy $H^s$ (Plotly)},
label={code:pde8-plotly}]
"""PDE, rozdzial 8: FFT, Plancherel, ewolucja ciepla i fal, Plotly.
Wymagania: numpy>=2.0, plotly>=5. Zapisuje samodzielne pliki HTML.
Transformata ciagla: F f(xi)=int exp(-i x xi) f(x) dx.
Siatka FFT jest okresowa: L=N*dx, xi=2*pi*fftfreq(N,d=dx).
"""
from pathlib import Path
import numpy as np
import plotly.graph_objects as go
OUT = Path(__file__).resolve().parent
OUT.mkdir(parents=True, exist_ok=True)

def save_plot(fig, filename):
    fig.update_layout(template="plotly_white", hovermode="x unified",
                      margin=dict(l=60, r=25, t=65, b=60))
    fig.write_html(OUT / filename, include_plotlyjs=True, full_html=True)

def gauss_fft():
    n, length, sigma = 2048, 24.0, 1.1
    dx = length / n
    x = (np.arange(n) - n//2) * dx
    f = np.exp(-x*x/(2*sigma*sigma))
    xi = np.fft.fftshift(np.fft.fftfreq(n, d=dx))*2*np.pi
    fhat = dx * np.fft.fftshift(np.fft.fft(np.fft.ifftshift(f)))
    exact = np.sqrt(2*np.pi)*sigma*np.exp(-sigma*sigma*xi*xi/2)
    mask = np.abs(xi) <= 8.0
    err = np.max(np.abs(fhat[mask]-exact[mask]))
    parseval_x = dx * np.sum(np.abs(f)**2)
    dxi = 2*np.pi/length
    parseval_k = (dxi/(2*np.pi))*np.sum(np.abs(fhat)**2)
    print("FFT vs Gauss: blad max na |xi|<=8:", f"{err:.3e}")
    print("Parseval: przestrzen", f"{parseval_x:.10f}",
          "czestotliwosc", f"{parseval_k:.10f}")
    fig = go.Figure()
    fig.add_trace(go.Scatter(x=xi[mask], y=exact[mask], mode="lines",
                             name="transformata dokladna"))
    fig.add_trace(go.Scatter(x=xi[mask], y=np.real(fhat[mask]), mode="markers",
                             marker=dict(size=4), name="FFT * dx"))
    fig.update_layout(title="Transformata Gaussa: wzor dokladny i FFT",
                      xaxis_title="czestotliwosc xi",
                      yaxis_title="F[f](xi)")
    save_plot(fig, "pde8_gauss_fft.html")

def modes_heat_wave():
    n, length = 1024, 2*np.pi
    dx = length/n
    x = np.arange(n)*dx
    frequencies = 2*np.pi*np.fft.fftfreq(n, d=dx)
    # Dane na torusie: te same dla obu rownan, predkosc poczatkowa fali 0.
    f = np.sin(2*x) + 0.55*np.cos(9*x) + 0.25*np.sin(23*x)
    fhat = np.fft.fft(f)
    kappa, c = 0.09, 1.0
    times = (0.0, 0.08, 0.3, 0.85)
    fig_heat, fig_wave = go.Figure(), go.Figure()
    for t in times:
        heat = np.fft.ifft(np.exp(-kappa*t*frequencies**2)*fhat).real
        wave = np.fft.ifft(np.cos(c*t*np.abs(frequencies))*fhat).real
        fig_heat.add_trace(go.Scatter(x=x, y=heat, name=f"t={t:g}"))
        fig_wave.add_trace(go.Scatter(x=x, y=wave, name=f"t={t:g}"))
    fig_heat.update_layout(title="Rownanie ciepla: tlumienie wysokich modow",
                           xaxis_title="x", yaxis_title="u(t,x)")
    fig_wave.update_layout(title="Rownanie falowe: mody oscyluja",
                           xaxis_title="x", yaxis_title="u(t,x)")
    save_plot(fig_heat, "pde8_heat_fft.html")
    save_plot(fig_wave, "pde8_wave_fft.html")

def h_s_delta():
    # Obciecie czestotliwosci: dla s<-1/2 calka stabilizuje sie.
    s_values = (-0.8, -0.5, -0.2)
    cutoffs = np.geomspace(2, 250, 60)
    # Dla d=1, ||delta||_{H^s}^2=(2*pi)^-1 * int(1+xi^2)^s dxi.
    xi = np.linspace(0, 250, 20001)
    fig = go.Figure()
    for s in s_values:
        vals = (1+xi*xi)**s
        increments = 0.5*(vals[1:]+vals[:-1])*np.diff(xi)
        cumulative = np.r_[0.0, np.cumsum(increments)] / np.pi
        approx = np.interp(cutoffs, xi, cumulative)
        fig.add_trace(go.Scatter(x=cutoffs, y=approx, name=f"s={s}"))
    fig.update_layout(title="Delta w H^s(R): normy z obcieciem czestotliwosci",
                      xaxis_title="odciecie |xi| <= R", xaxis_type="log",
                      yaxis_title="(2*pi)^-1 calka do R")
    save_plot(fig, "pde8_delta_Hs.html")
if __name__ == "__main__":
    gauss_fft()
    modes_heat_wave()
    h_s_delta()
    print("Wyniki zapisane w", OUT)
\end{lstlisting}

\section{Wygładzanie Sobolewa i warunki brzegowe}
Poniższy eksperyment porównuje zbieżność splotu w normie $H^1$
ze zbieżnością obcięć brzegowych jedynie w $L^2$.
\begin{lstlisting}[style=pythonpde,caption={Wygładzanie w $H^1$, pochodna słaba i obcięcia brzegowe -- Plotly},label={code:pde-sobolev9}]
"""Iteracja 9: Sobolev W^{1,2}, periodyczny splot i warunki brzegowe.

Wymagania: numpy, plotly. Zapisuje trzy samodzielne pliki HTML.
"""
from pathlib import Path
import numpy as np
import plotly.graph_objects as go

HERE = Path(__file__).resolve().parent


def save(fig, filename):
    fig.update_layout(template="plotly_white", hovermode="x unified",
                      font=dict(size=13), legend=dict(orientation="h", y=1.11))
    fig.write_html(HERE / filename, include_plotlyjs=True, full_html=True)


def periodic_mollification():
    n = 4096
    x = -np.pi + 2 * np.pi * np.arange(n) / n
    h = 2 * np.pi / n
    u = np.abs(np.sin(x))
    du = np.sign(np.sin(x)) * np.cos(x)  # pochodna slaba a.e.
    k = np.fft.fftfreq(n, d=h) * (2 * np.pi)
    uhat = np.fft.fft(u)
    epsilons = [0.40, 0.20, 0.10, 0.05, 0.025, 0.0125]
    fig = go.Figure()
    fig.add_trace(go.Scatter(x=x, y=u, name="u=|sin(x)|", line=dict(width=3)))
    errors = []
    for eps in epsilons:
        mult = np.exp(-0.5 * eps**2 * k**2)
        smooth = np.fft.ifft(mult * uhat).real
        d_smooth = np.fft.ifft(1j * k * mult * uhat).real
        err_l2 = np.sqrt(h * np.sum((smooth - u)**2))
        err_grad = np.sqrt(h * np.sum((d_smooth - du)**2))
        errors.append((eps, np.hypot(err_l2, err_grad)))
        if eps in [0.40, 0.10, 0.025]:
            fig.add_trace(go.Scatter(x=x, y=smooth, name=f"eps={eps:.3f}"))
    fig.update_layout(title="Wygladzanie w H^1(T): u(x)=|sin x|",
                      xaxis_title="x", yaxis_title="u_epsilon(x)")
    save(fig, "pde9_wygladzanie_plotly.html")

    fig2 = go.Figure()
    fig2.add_trace(go.Scatter(x=x, y=du, name="pochodna slaba", line=dict(width=2)))
    for eps in [0.40, 0.10, 0.025]:
        derivative = np.fft.ifft(1j * k * np.exp(-0.5*eps**2*k**2) * uhat).real
        fig2.add_trace(go.Scatter(x=x, y=derivative, name=f"eps={eps:.3f}"))
    fig2.update_layout(title="Wygladzanie pochodnej slabej",
                       xaxis_title="x", yaxis_title="d u_epsilon / dx")
    save(fig2, "pde9_pochodne_plotly.html")
    print("Spadek bledu H1:", errors)
    assert errors[-1][1] < errors[0][1], "Blad powinien malec przy eps->0"


def smoothstep(z):
    v = np.clip(z, 0.0, 1.0)
    return 3*v*v - 2*v*v*v


def smoothstep_deriv(z):
    v = np.clip(z, 0.0, 1.0)
    return 6*v*(1-v)


def boundary_cutoffs():
    x = (np.arange(20001) + 0.5) / 20001
    vals = []
    fig = go.Figure()
    for n in [4, 8, 16, 32, 64]:
        a, b = smoothstep(n*x), smoothstep(n*(1-x))
        d_a = n*smoothstep_deriv(n*x)
        d_b = -n*smoothstep_deriv(n*(1-x))
        chi = a*b
        d_chi = d_a*b + a*d_b
        err = np.sqrt(np.mean((chi-1)**2))
        grad = np.sqrt(np.mean(d_chi**2))
        vals.append((n, err, grad))
        fig.add_trace(go.Scatter(x=x, y=chi, name=f"n={n}"))
    fig.update_layout(title="Obciecia daza do 1 w L^2, ale nie w H^1",
                      xaxis_title="x w (0,1)", yaxis_title="chi_n(x)")
    save(fig, "pde9_slady_obciecia_plotly.html")
    print("n, ||chi-1||L2, ||chi'||L2:", vals)
    assert vals[-1][1] < vals[0][1]
    assert vals[-1][2] > vals[0][2]


if __name__ == "__main__":
    periodic_mollification()
    boundary_cutoffs()
\end{lstlisting}


\section{Koncentracja krytyczna i ciąg Mosera — Plotly}
Poniższy program wizualizuje koncentrację przy wykładniku $p^*=6$ dla $d=3$, $p=2$, porównuje skalowanie norm oraz ciąg Mosera w wymiarze $2$. Wyniki są zapisywane w trzech interaktywnych plikach HTML.
\begin{lstlisting}[style=pythonpde,caption={Koncentracja i ostrosc zanurzen Sobolewa - Plotly},label={code:pde-embed10}]
"""PDE, rozdzial 10. Interaktywne wizualizacje zanurzen Sobolewa.
Uruchom: python demo_zanurzenia_plotly.py
Pliki HTML zapisuja lokalna kopie Plotly.js i dzialaja offline.
"""
from pathlib import Path
import numpy as np
import plotly.graph_objects as go

OUT = Path(__file__).resolve().parent
EPS = [0.8, 0.4, 0.2, 0.1, 0.05]


def bump(r):
    r = np.asarray(r, dtype=float)
    v = np.zeros_like(r)
    idx = np.abs(r) < 1
    v[idx] = np.exp(-1 / (1 - r[idx] ** 2))
    return v


def bump_prime(r):
    r = np.asarray(r, dtype=float)
    v = np.zeros_like(r)
    idx = np.abs(r) < 1
    v[idx] = bump(r[idx]) * (-2 * r[idx] / (1 - r[idx] ** 2) ** 2)
    return v


# d=3, p=2. Wykladnik krytyczny p*=6, skala eps^{-1/2}.
x = np.linspace(-1, 1, 1601)
fig = go.Figure()
for eps in EPS:
    fig.add_trace(go.Scatter(
        x=x, y=eps ** (-0.5) * bump(x / eps),
        name=f"eps={eps:g}", mode="lines"))
fig.update_layout(title="Koncentracja: W1,2 w wymiarze 3 (przekroj radialny)",
                  xaxis_title="x", yaxis_title="u_eps(x)", template="plotly_white")
fig.write_html(OUT / "pde10_koncentracja.html", include_plotlyjs=True)

# Calkowanie sferyczne: 4*pi int_0^1 |v(r)|^q r2 dr.
r = np.linspace(0, 1, 50001)
q_values = [2, 4, 6, 8]
baseline = {q: (4 * np.pi * np.trapezoid(bump(r) ** q * r ** 2, r)) ** (1 / q)
            for q in q_values}
grad = np.sqrt(4 * np.pi * np.trapezoid(bump_prime(r) ** 2 * r ** 2, r))
fig = go.Figure()
for q in q_values:
    vals = [e ** (-0.5 + 3 / q) * baseline[q] for e in EPS]
    fig.add_trace(go.Scatter(x=EPS, y=vals, mode="lines+markers", name=f"L^{q}"))
fig.add_trace(go.Scatter(x=EPS, y=[grad] * len(EPS),
                         mode="lines+markers", name="norma gradientu L2"))
fig.update_layout(title="Skalowanie norm: krytyczne L6 nie zalezy od eps",
                  xaxis_title="eps", yaxis_title="norma",
                  xaxis_type="log", yaxis_type="log", template="plotly_white")
fig.write_html(OUT / "pde10_skalowanie_norm.html", include_plotlyjs=True)

# Ciag Mosera w wymiarze 2: sup rosnie, norma gradientu stala.
r = np.geomspace(1e-5, 1, 900)
fig = go.Figure()
for eps in [0.1, 0.01, 0.001, 0.0001]:
    L = np.log(1 / eps)
    val = np.minimum(np.log(1 / r), L) / np.sqrt(L)
    fig.add_trace(go.Scatter(x=r, y=val, mode="lines", name=f"eps={eps:g}"))
fig.update_layout(title="Ciag Mosera: wzrost maksimum przy stalej normie gradu w L2",
                  xaxis_title="promien r", yaxis_title="u_eps(r)",
                  xaxis_type="log", template="plotly_white")
fig.write_html(OUT / "pde10_moser.html", include_plotlyjs=True)

for q in q_values:
    assert np.isfinite(baseline[q]) and baseline[q] > 0
assert np.isfinite(grad) and grad > 0
print(f"Kontrola norm: L6(eps=0.8)={baseline[6]:.8f}; "
      f"L6(eps=0.05)={baseline[6]:.8f}; grad L2={grad:.8f}")
print("Wygenerowano: pde10_koncentracja.html, pde10_skalowanie_norm.html, pde10_moser.html")
\end{lstlisting}


\section{Galerkin, zbieżność i energia — Plotly}
\label{sec:pde11-code}
Poniższy program tworzy macierz sztywności metody elementów liniowych,
porównuje aproksymację z rozwiązaniem dokładnym, weryfikuje wzory
na błędy w $L^2$ i seminormie $H^1$ oraz wizualizuje funkcjonał energii.
Wynikami są trzy interaktywne pliki HTML, otwierane bez połączenia
z Internetem.
\begin{lstlisting}[style=pythonpde,caption={Galerkin, błędy i funkcjonał energii (Plotly)},label={code:pde11-galerkin}]
"""PDE - rozdzial 11: metoda Galerkina dla -u''=1, u(0)=u(1)=0.

Instalacja: pip install numpy plotly
Uruchom: python demo_wariacyjne_plotly.py
Wyniki: trzy samodzielne pliki HTML z Plotly (offline).
"""
from pathlib import Path

import numpy as np
import plotly.graph_objects as go

OUT = Path(__file__).resolve().parent


def exact(x):
    return x * (1.0 - x) / 2.0


def fem_solution(N):
    """Zwraca wezly i wartosci liniowych elementow skonczonych."""
    h = 1.0 / N
    diagonal = np.full(N - 1, 2.0 / h)
    K = np.diag(diagonal)
    if N > 2:
        offdiag = np.full(N - 2, -1.0 / h)
        K += np.diag(offdiag, 1) + np.diag(offdiag, -1)
    b = np.full(N - 1, h)
    U = np.linalg.solve(K, b)
    nodes = np.linspace(0.0, 1.0, N + 1)
    values = np.r_[0.0, U, 0.0]
    assert np.allclose(values, exact(nodes), atol=1e-12)
    assert np.linalg.eigvalsh(K)[0] > 0
    return nodes, values


def quadrature_errors(nodes, values):
    """Bledy L2 i H1 liczone calkowaniem Gaussa na kazdej komorce."""
    z, weights = np.polynomial.legendre.leggauss(4)
    sq_l2 = sq_h1 = 0.0
    for i in range(len(nodes) - 1):
        x0, x1 = nodes[i : i + 2]
        h = x1 - x0
        x = (x0 + x1) / 2.0 + h * z / 2.0
        slope = (values[i + 1] - values[i]) / h
        approx = values[i] + slope * (x - x0)
        sq_l2 += h / 2.0 * np.sum(weights * (exact(x) - approx) ** 2)
        sq_h1 += h / 2.0 * np.sum(weights * ((0.5 - x) - slope) ** 2)
    return np.sqrt(sq_l2), np.sqrt(sq_h1)


def save(fig, name):
    fig.update_layout(template="plotly_white", font=dict(size=13))
    fig.write_html(OUT / name, include_plotlyjs=True, full_html=True)


def main():
    Ns = [4, 8, 16, 32, 64]
    h_values, errors_l2, errors_h1 = [], [], []
    fine = np.linspace(0.0, 1.0, 1200)

    profile = go.Figure()
    profile.add_trace(go.Scatter(x=fine, y=exact(fine), name="Rozwiazanie dokladne"))

    for N in Ns:
        nodes, values = fem_solution(N)
        err_l2, err_h1 = quadrature_errors(nodes, values)
        h = 1.0 / N
        expected_l2 = h**2 / np.sqrt(120.0)
        expected_h1 = h / np.sqrt(12.0)
        assert np.isclose(err_l2, expected_l2, rtol=1e-9, atol=1e-12)
        assert np.isclose(err_h1, expected_h1, rtol=1e-9, atol=1e-12)
        h_values.append(h)
        errors_l2.append(err_l2)
        errors_h1.append(err_h1)
        if N <= 16:
            profile.add_trace(go.Scatter(x=nodes, y=values, name=f"FEM N={N}",
                                         mode="lines+markers"))
        print(f"N={N:2d}, h={h:.5f}, blad L2={err_l2:.8e}, blad H1={err_h1:.8e}")

    profile.update_layout(title="Galerkin: rozwiazanie dokladne i aproksymacje liniowe",
                          xaxis_title="x", yaxis_title="u(x)")
    save(profile, "pde11_fem_rozwiazania.html")

    convergence = go.Figure()
    convergence.add_trace(go.Scatter(x=h_values, y=errors_l2,
                                     name="||u-u_h||_L2 (rzad 2)", mode="lines+markers"))
    convergence.add_trace(go.Scatter(x=h_values, y=errors_h1,
                                     name="||u'-u_h'||_L2 (rzad 1)", mode="lines+markers"))
    convergence.update_layout(title="Zbieznosc FEM: dokladne bledy dla -u''=1",
                              xaxis_title="h", yaxis_title="norma bledu",
                              xaxis_type="log", yaxis_type="log")
    save(convergence, "pde11_fem_zbieznosc.html")

    # Dla u(x)=x(1-x)/2 mamy int_0^1 |u'|^2 dx = int_0^1 u dx = 1/12.
    t = np.linspace(-0.25, 2.25, 300)
    energy = (0.5 * t * t - t) / 12.0
    energy_plot = go.Figure()
    energy_plot.add_trace(go.Scatter(x=t, y=energy, name="J(tu)"))
    energy_plot.add_trace(go.Scatter(x=[1], y=[-1.0 / 24.0], mode="markers",
                                     marker=dict(size=11), name="Minimum t=1"))
    energy_plot.update_layout(title="Zasada Dirichleta: energia na prostej tu",
                              xaxis_title="t", yaxis_title="J(tu)")
    save(energy_plot, "pde11_energia.html")

    orders = np.polyfit(np.log(h_values), np.log(errors_l2), 1)[0], \
             np.polyfit(np.log(h_values), np.log(errors_h1), 1)[0]
    assert np.allclose(orders, (2.0, 1.0), atol=1e-9)
    print(f"Zweryfikowane rzedy: L2={orders[0]:.12f}, H1={orders[1]:.12f}")


if __name__ == "__main__":
    main()
\end{lstlisting}

\section{Spektrum, półgrupy i rezonans — Plotly}
\label{sec:pde12-code}
Poniższy program porównuje dokładne wartości własne operatora
$-\frac{\mathrm d^2}{\mathrm dx^2}$ na $(0,1)$ z widmem macierzy
różnicowej, wizualizuje ewolucję cieplną i falową tych samych
danych początkowych oraz wzrost normy rezolwenty w pobliżu
pierwszej wartości własnej. Sprawdza również rząd zbieżności
wartości własnych, równy $2$ dla zastosowanego schematu.
\begin{lstlisting}[style=pythonpde,caption={Widmo operatora Dirichleta, ewolucje i rezonans (Plotly)},label={code:pde12-operators}]
"""PDE, rozdzial 12: widmo laplasjanu Dirichleta i ewolucje.

Instalacja: pip install numpy plotly
Uruchom: python demo_operatory_plotly.py
Wyniki: trzy samodzielne pliki HTML, dzialajace offline.
"""
from pathlib import Path
import numpy as np
import plotly.graph_objects as go

OUT = Path(__file__).resolve().parent
PI = np.pi


def discrete_spectrum(n_intervals):
    h = 1.0 / n_intervals
    matrix = np.diag(np.full(n_intervals - 1, 2.0 / h**2))
    matrix += np.diag(np.full(n_intervals - 2, -1.0 / h**2), 1)
    matrix += np.diag(np.full(n_intervals - 2, -1.0 / h**2), -1)
    values = np.linalg.eigvalsh(matrix)
    j = np.arange(1, n_intervals, dtype=float)
    formula = 4.0 / h**2 * np.sin(j * PI * h / 2.0) ** 2
    assert np.allclose(values, formula, rtol=2e-12, atol=2e-11)
    return values


def save(fig, name):
    fig.update_layout(template="plotly_white", font=dict(size=13))
    fig.write_html(OUT / name, include_plotlyjs=True, full_html=True)


def main():
    sizes = [12, 24, 48, 96, 192]
    nplot = 6
    exact = (PI * np.arange(1, nplot + 1)) ** 2
    fig = go.Figure()
    fig.add_trace(go.Scatter(x=np.arange(1, nplot + 1), y=exact,
                             name="Wartosci dokladne", mode="lines+markers"))
    errors = []
    for n in sizes:
        eigen = discrete_spectrum(n)
        fig.add_trace(go.Scatter(x=np.arange(1, nplot + 1), y=eigen[:nplot],
                                 name=f"FD, N={n}", mode="lines+markers"))
        errors.append(abs(eigen[0] - PI**2))
    fig.update_layout(title="Laplasjan Dirichleta: widmo ciagle vs. macierz FD",
                      xaxis_title="Numer modu j", yaxis_title="lambda_j")
    save(fig, "pde12_widmo_plotly.html")
    h = 1.0 / np.array(sizes, dtype=float)
    order = np.polyfit(np.log(h), np.log(errors), 1)[0]
    print(f"Rzad bledu pierwszej wartosci wlasnej = {order:.4f}")
    assert 1.95 < order < 2.05

    x = np.linspace(0.0, 1.0, 1000)
    modes = np.array([1, 3, 6], dtype=float)
    coeffs = np.array([1.0, 0.35, -0.2])
    basis = np.sin(np.pi * modes[:, None] * x[None, :])
    lambdas = (np.pi * modes) ** 2
    fig2 = go.Figure()
    for t in (0.0, 0.01, 0.05, 0.12):
        y = np.sum((coeffs * np.exp(-t * lambdas))[:, None] * basis, axis=0)
        fig2.add_trace(go.Scatter(x=x, y=y, name=f"Cieplo t={t:g}"))
    for t in (0.05, 0.12):
        y = np.sum((coeffs * np.cos(t * np.sqrt(lambdas)))[:, None] * basis,
                   axis=0)
        fig2.add_trace(go.Scatter(x=x, y=y, name=f"Fala t={t:g}",
                                  line=dict(dash="dash")))
    fig2.update_layout(title="Te same dane: wygladzanie ciepla i oscylacje fali",
                       xaxis_title="x", yaxis_title="u(t,x)")
    save(fig2, "pde12_cieplo_fala_plotly.html")

    lambda1 = PI**2
    delta = np.geomspace(0.001, 4.0, 160)
    fig3 = go.Figure()
    fig3.add_trace(go.Scatter(x=delta, y=1.0 / delta,
                              name="|| (A-(lambda_1+delta)I)^-1 phi_1 ||"))
    fig3.update_layout(title="Alternatywa Fredholma: rezonans blisko lambda_1",
                       xaxis_title="delta = |zeta-lambda_1|",
                       yaxis_title="Norma rozwiazania", xaxis_type="log",
                       yaxis_type="log")
    save(fig3, "pde12_rezolwenta_plotly.html")
    print("Zweryfikowano formula FD, blad O(h^2) i trzy wizualizacje Plotly.")


if __name__ == "__main__":
    main()
\end{lstlisting}


\section{Eliptyczne zagadnienie transmisyjne — Plotly}
\label{sec:pde13-code}
Program rozwiązuje metodą elementów skończonych zadanie
$-(a u')'=1$ na $(-1,1)$, z warunkami Dirichleta i skokiem
współczynnika $a$ w zerze. Porównuje rozwiązanie numeryczne ze
wzorem~\eqref{eq:pde13-exact}, przedstawia strumień
$q=-a u'$ oraz bada błędy w normie $L^2$ i seminormie $H^1$.
Sprawdza także dokładne wzory na błędy interpolacji:
$\|u-u_h\|_{L^2}=h^2\sqrt{17/1920}$ i
$\|(u-u_h)'\|_{L^2}=h\sqrt{17/192}$.
\begin{lstlisting}[style=pythonpde,caption={Transmisja i zbieżność metody Galerkina (Plotly)},label={code:pde13-transmission}]
"""PDE, rozdzial 13: dywergencyjne zagadnienie transmisyjne.

Instalacja: pip install numpy plotly
Uruchom: python demo_eliptyczne_plotly.py
Wyniki: samodzielne pliki HTML (dzialaja offline).
"""
from pathlib import Path
import numpy as np
import plotly.graph_objects as go

OUT = Path(__file__).resolve().parent


def exact(x):
    x = np.asarray(x, dtype=float)
    return np.where(x <= 0, -x*x/2 - 0.3*x + 0.2,
                    -x*x/8 - 0.075*x + 0.2)


def exact_deriv(x):
    x = np.asarray(x, dtype=float)
    return np.where(x < 0, -x - 0.3, -x/4 - 0.075)


def solve_tridiagonal(subdiag, diag, rhs):
    """Algorytm Thomasa: macierz symetryczna trojdiagonalna."""
    b, d, c = diag.copy(), rhs.copy(), subdiag.copy()
    for i in range(1, len(b)):
        fac = c[i - 1] / b[i - 1]
        b[i] -= fac * c[i - 1]
        d[i] -= fac * d[i - 1]
    x = np.empty_like(d)
    x[-1] = d[-1] / b[-1]
    for i in range(len(x)-2, -1, -1):
        x[i] = (d[i] - c[i] * x[i+1]) / b[i]
    return x


def fem_solution(n_per_half):
    """Liniowe FEM dla -(a u')'=1, u(-1)=u(1)=0.

    Punkt 0 jest wezelm siatki. a=1 po lewej i a=4 po prawej.
    """
    nodes = np.linspace(-1., 1., 2 * n_per_half + 1)
    h = nodes[1] - nodes[0]
    a_elem = np.where((nodes[:-1] + nodes[1:]) / 2 < 0, 1., 4.)
    diag = (a_elem[:-1] + a_elem[1:]) / h
    subdiag = -a_elem[1:-1] / h
    rhs = np.full(len(nodes)-2, h)
    sol = np.r_[0., solve_tridiagonal(subdiag, diag, rhs), 0.]
    assert np.max(np.abs(sol - exact(nodes))) < 5e-10
    return nodes, sol, a_elem


def errors(nodes, sol):
    z, w = np.polynomial.legendre.leggauss(5)
    l2sq = h1sq = 0.
    for j in range(len(nodes) - 1):
        x0, x1 = nodes[j:j+2]
        h = x1 - x0
        xx = (x0+x1)/2 + h*z/2
        approx = sol[j] + (sol[j+1]-sol[j])*(xx-x0)/h
        slope = (sol[j+1]-sol[j])/h
        l2sq += (h/2) * np.dot(w, (approx-exact(xx))**2)
        h1sq += (h/2) * np.dot(w, (slope-exact_deriv(xx))**2)
    return np.sqrt(l2sq), np.sqrt(h1sq)


def save(fig, name):
    fig.update_layout(template="plotly_white", font=dict(size=13),
                      hovermode="x unified")
    fig.write_html(OUT / name, include_plotlyjs=True, full_html=True)


def main():
    hs, l2errors, h1errors = [], [], []
    fig_sol = go.Figure()
    x = np.linspace(-1, 1, 1201)
    fig_sol.add_trace(go.Scatter(x=x, y=exact(x), mode="lines",
                                 name="Rozwiazanie dokladne"))
    fig_flux = go.Figure()
    fig_flux.add_trace(go.Scatter(x=x, y=x + 0.3, mode="lines",
                                  name="Dokladny strumien q=-a u'"))
    for N in [4, 8, 16, 32, 64, 128]:
        nodes, vals, a_elem = fem_solution(N)
        h = nodes[1]-nodes[0]
        l2, h1 = errors(nodes, vals)
        assert np.isclose(l2, h*h * np.sqrt(17/1920), rtol=5e-8, atol=5e-11)
        assert np.isclose(h1, h * np.sqrt(17/192), rtol=5e-8, atol=5e-11)
        hs.append(h)
        l2errors.append(l2)
        h1errors.append(h1)
        if N in (4, 16):
            fig_sol.add_trace(go.Scatter(x=nodes, y=vals, mode="lines+markers",
                                         name=f"FEM, N={N} na polowe"))
        if N in (4, 16):
            mid = (nodes[:-1]+nodes[1:])/2
            qfem = -a_elem*np.diff(vals)/h
            fig_flux.add_trace(go.Scatter(x=mid, y=qfem, mode="markers",
                                          name=f"q_h w srodkach, N={N}"))
        print(f"N={N:3d} h={h:.5f} L2={l2:.9e} H1-semi={h1:.9e}")
    fig_sol.add_vline(x=0, line_dash="dash", annotation_text="Interfejs")
    fig_sol.update_layout(title="Transmisja przez dwa osrodki: rozwiazanie slabe",
                          xaxis_title="x", yaxis_title="u(x)")
    save(fig_sol, "pde13_transmisja_rozwiazanie.html")
    fig_flux.add_vline(x=0, line_dash="dash", annotation_text="Interfejs")
    fig_flux.update_layout(title="Strumien ciagly w rozwiazaniu dokladnym",
                           xaxis_title="x", yaxis_title="q=-a u'(x)")
    save(fig_flux, "pde13_transmisja_strumien.html")
    fig_error = go.Figure()
    fig_error.add_trace(go.Scatter(x=hs, y=l2errors, mode="lines+markers",
                                   name="Blad L2, rzad 2"))
    fig_error.add_trace(go.Scatter(x=hs, y=h1errors, mode="lines+markers",
                                   name="Blad seminormy H1, rzad 1"))
    fig_error.update_layout(title="Zbieznosc FEM dla rownania eliptycznego",
                            xaxis_title="h", yaxis_title="Blad",
                            xaxis_type="log", yaxis_type="log")
    save(fig_error, "pde13_transmisja_zbieznosc.html")
    p2 = np.polyfit(np.log(hs), np.log(l2errors), 1)[0]
    p1 = np.polyfit(np.log(hs), np.log(h1errors), 1)[0]
    assert np.isclose(p2, 2., atol=1e-6) and np.isclose(p1, 1., atol=1e-6)
    print(f"Rzedy: L2={p2:.6f}, H1-semi={p1:.6f}; testy zaliczone.")


if __name__ == "__main__":
    main()
\end{lstlisting}

\section{Zasady maksimum, nierówność Harnacka i lemat Hopfa — Plotly}
Kod oblicza harmoniczne przedłużenie danych brzegowych w dysku,
weryfikuje nierówność Harnacka oraz pochodną normalną w przykładzie
ilustrującym lemat Hopfa. Porównanie jądra Poissona z rozwiązaniem
analitycznym jest jednocześnie testem poprawności normalizacji jądra.
\begin{lstlisting}[style=pythonpde,caption={Jądro Poissona, nierówność Harnacka i lemat Hopfa (Plotly)},label={code:pde14-plotly}]
"""PDE, rozdzial 14: jadro Poissona, Harnack i Hopf.

Instalacja: pip install numpy plotly
Uruchomienie: python demo_maksimum_harnack_plotly.py
Wykresy HTML zawieraja Plotly offline.
"""
from pathlib import Path
import numpy as np
import plotly.graph_objects as go

OUT = Path(__file__).resolve().parent
TWO_PI = 2.0 * np.pi


def save(fig, filename):
    fig.update_layout(template="plotly_white", font=dict(size=13))
    fig.write_html(OUT / filename, include_plotlyjs=True, full_html=True)


def boundary(phi):
    return 0.5 + np.cos(2.0 * phi) + 0.3 * np.sin(3.0 * phi)


def poisson_disk(r, theta, phi):
    # Jadrowa gestosc wzgledem miary dlugosci na okregu jednostkowym.
    return (1.0 - r * r) / (TWO_PI * (1.0 + r * r - 2.0 * r * np.cos(theta - phi)))


def harmonic_exact(r, theta):
    return 0.5 + r**2 * np.cos(2.0 * theta) + 0.3 * r**3 * np.sin(3.0 * theta)


def main():
    phi = np.linspace(0.0, TWO_PI, 2048, endpoint=False)
    theta = np.linspace(0.0, TWO_PI, 160)
    g = boundary(phi)
    fig = go.Figure()
    max_error = 0.0
    for r in [0.0, 0.3, 0.6, 0.85]:
        p = poisson_disk(r, theta[:, None], phi[None, :])
        vals = TWO_PI * np.mean(p * g[None, :], axis=1)
        exact = harmonic_exact(r, theta)
        error = np.max(np.abs(vals - exact))
        max_error = max(max_error, error)
        fig.add_trace(go.Scatter(x=theta, y=vals, mode="lines", name=f"r={r:g}"))
        assert error < 1e-9, (r, error)
    fig.update_layout(title="Przedluzenie harmoniczne przez jadro Poissona",
                      xaxis_title="kat theta", yaxis_title="u(r, theta)")
    save(fig, "pde14_poisson_plotly.html")
    print(f"Poisson: maksymalny blad dyskretyzacji = {max_error:.3e}")

    # Jawna funkcja harmoniczna: u(x,y)=1+0.6*x w B_1(0).
    r = 0.5
    t = np.linspace(0.0, TWO_PI, 500)
    vals = 1.0 + 0.6 * r * np.cos(t)
    ratio = np.max(vals) / np.min(vals)
    lower = (1.0 - r**2) / (1.0 + r)**2
    upper = 1.0 / (1.0 - r)**2
    assert np.min(vals) >= lower and np.max(vals) <= upper
    fig = go.Figure()
    fig.add_trace(go.Scatter(x=t, y=vals, name="u(0.5 cos(theta), 0.5 sin(theta))"))
    fig.add_hline(y=lower, line_dash="dash", annotation_text="dolne Harnack")
    fig.add_hline(y=upper, line_dash="dash", annotation_text="gorne Harnack")
    fig.update_layout(title=f"Harnack: rzeczywisty sup/inf = {ratio:.4f}",
                      xaxis_title="kat theta", yaxis_title="u")
    save(fig, "pde14_harnack_plotly.html")
    print(f"Harnack: sup/inf = {ratio:.6f}, zakres teoretyczny = [{lower:.3f}, {upper:.3f}]")

    # Dodatnia funkcja harmoniczna u(x,y)=1-x; w x0=(1,0) minimum 0.
    x = np.linspace(0.75, 1.0, 150)
    u = 1.0 - x
    fig = go.Figure()
    fig.add_trace(go.Scatter(x=x, y=u, name="u(x,0)=1-x"))
    fig.add_trace(go.Scatter(x=x, y=(1.0-x), mode="markers", name="styczna w x0", marker=dict(size=3)))
    fig.update_layout(title="Lemat Hopfa: u(1,0)=0, pochodna normalna = -1",
                      xaxis_title="x (promien do brzegu)", yaxis_title="u(x,0)")
    save(fig, "pde14_hopf_plotly.html")
    assert np.allclose(np.gradient(u, x), -1.0, atol=1e-10)
    print("Hopf: pochodna zewnetrzna = -1 (blad < 1e-10)")


if __name__ == "__main__":
    main()
\end{lstlisting}

\section{Ilorazy różnicowe i skok pochodnej — Plotly}
Przykład pokazuje, jak skok pochodnej prowadzi do nieograniczonej normy dyskretnych drugich pochodnych.
\begin{lstlisting}[style=pythonpde,caption={Brak regularnosci H2 dla transmisji (Plotly)},label={code:pde15-plotly}]
"""PDE, rozdzial 15: H2-regularnosc i transmisja (Plotly).
Instalacja: pip install numpy plotly
"""
from pathlib import Path
import numpy as np
import plotly.graph_objects as go

OUT = Path(__file__).resolve().parent

def save(fig, name):
    fig.update_layout(template="plotly_white", font=dict(size=13))
    fig.write_html(OUT / name, include_plotlyjs=True, full_html=True)

def transmission(x):
    return np.where(x <= 0.5, x, 0.5 + 0.5*(x-0.5))

def main():
    x = np.linspace(0., 1., 1001)
    u = transmission(x)
    fig = go.Figure()
    fig.add_trace(go.Scatter(x=x, y=u, name="u: ciagla, z zalamaniem"))
    fig.add_trace(go.Scatter(x=x, y=x*(1-x), name="v: gladka"))
    fig.add_vline(x=0.5, line_dash="dot", annotation_text="granica osrodkow")
    fig.update_layout(title="Regularnosc rozwiazan: gladkie i transmisyjne",
                      xaxis_title="x", yaxis_title="wartosc")
    save(fig, "pde15_regularnosc_plotly.html")

    Ns = np.array([20, 40, 80, 160, 320, 640])
    hs = 1.0/Ns
    norms = []
    for N in Ns:
        xgrid = np.linspace(0, 1, N+1)
        uh = transmission(xgrid)
        second = (uh[2:] - 2*uh[1:-1] + uh[:-2]) / (1/N)**2
        norms.append(np.sqrt(np.sum(second**2) / N))
    norms = np.array(norms)
    reference = 0.5 / np.sqrt(hs)
    assert np.allclose(norms, reference, rtol=1e-10, atol=1e-10)
    slope = np.polyfit(np.log(hs), np.log(norms), 1)[0]
    assert abs(slope + .5) < 1e-10
    fig = go.Figure()
    fig.add_trace(go.Scatter(x=hs, y=norms, mode="lines+markers", name="dyskretna norma H2"))
    fig.add_trace(go.Scatter(x=hs, y=reference, mode="lines", name="0.5 / sqrt(h)"))
    fig.update_layout(title=f"Brak H2: norma drugiej roznicy ~ h^({slope:.2f})",
                      xaxis_type="log", yaxis_type="log",
                      xaxis_title="h", yaxis_title="norma L2 drugiej roznicy")
    save(fig, "pde15_skok_h2_plotly.html")
    print(f"Nachylenie log-log: {slope:.8f} (oczekiwane -0.5)")

if __name__ == "__main__":
    main()
\end{lstlisting}

\section[Drugie pochodne i ostrość oszacowań Schaudera — Plotly]{Drugie pochodne, brak oszacowania w $L^\infty$ i ostrość oszacowań Schaudera — Plotly}
Poniższy eksperyment na torusie dwuwymiarowym ilustruje kontrast
między normą średniokwadratową a normą supremum Hessianu.
Jest to eksperyment periodyczny, nie dowód twierdzenia dla $\R^d$.
Na końcu kod przedstawia dokładny przykład ostrości wykładnika Höldera.
\begin{lstlisting}[style=pythonpde,caption={Rachunek spektralny dla Hessianu i wykresy ilustrujące oszacowania Schaudera (Plotly)},label={code:pde16-plotly}]
"""Rozdzial 16: regularnosc W2p i oszacowania Schaudera; interaktywne wykresy Plotly.
Instalacja: pip install numpy plotly
Uruchomienie: python demo_calderon_schauder_plotly.py
"""
from pathlib import Path
import numpy as np
import plotly.graph_objects as go

OUT = Path(__file__).resolve().parent


def save(fig, name):
    fig.update_layout(template='plotly_white', font=dict(size=13))
    fig.write_html(OUT / name, include_plotlyjs=True, full_html=True)


def mixed_hessian(N):
    # Kwadrat w torusie [0,2pi)^2; z jego funkcji wskaznikowej
    # usuwamy srednia, by rozwiazanie -Delta u=f istnialo okresowo.
    x = 2 * np.pi * np.arange(N) / N
    square = ((x[:, None] < np.pi) & (x[None, :] < np.pi)).astype(float)
    f = square - square.mean()
    k = np.fft.fftfreq(N, d=1.0/N)  # calkowite czestosci
    K1, K2 = np.meshgrid(k, k, indexing='ij')
    k2 = K1*K1 + K2*K2
    m = np.zeros_like(k2, dtype=float)
    np.divide(-K1*K2, k2, out=m, where=k2 != 0)
    d12u = np.fft.ifft2(np.fft.fft2(f)*m).real
    return x, f, d12u


def demo_fft():
    N = 256
    x, f, d12u = mixed_hessian(N)
    fig = go.Figure(data=go.Heatmap(x=x, y=x, z=d12u.T,
                                     colorbar=dict(title='u_xy')))
    fig.update_layout(title='Drugie pochodne rozwiazania -Delta u=f, dane naroznikowe',
                      xaxis_title='x', yaxis_title='y', yaxis_scaleanchor='x')
    save(fig, 'pde16_fft_hessjan.html')

    Ns = [32, 64, 128, 256, 512]
    peaks = []
    rms = []
    for n in Ns:
        _, fn, hn = mixed_hessian(n)
        peaks.append(float(np.max(np.abs(hn))))
        rms.append(float(np.mean(hn**2)**0.5))
        # Mnoznik Hessianu ma modul <= 1 w normie L2:
        assert rms[-1] <= float(np.mean(fn**2)**0.5) + 1e-11
    fig = go.Figure()
    fig.add_trace(go.Scatter(x=Ns, y=peaks, mode='lines+markers', name='maksimum |u_xy|'))
    fig.add_trace(go.Scatter(x=Ns, y=rms, mode='lines+markers', name='norma sredniokwadratowa'))
    fig.update_layout(title='Wzrost maksimum przy zachowanym oszacowaniu L2',
                      xaxis_title='liczba wezlow N', xaxis_type='log',
                      yaxis_title='wartosc normy')
    save(fig, 'pde16_fft_normy.html')
    print('N / maksimum / RMS')
    for n, peak, r in zip(Ns, peaks, rms):
        print(f'{n:4d} / {peak:.5f} / {r:.5f}')
    assert peaks[-1] > peaks[0], 'Oczekiwano wzrostu maksimum przy narozniku'


def demo_holder():
    x = np.linspace(-1, 1, 1501)
    alpha = 0.6
    u = -np.abs(x)**(alpha+2)/((alpha+1)*(alpha+2))
    f = np.abs(x)**alpha
    fig = go.Figure()
    fig.add_trace(go.Scatter(x=x, y=u, name='u: rozwiazanie -u\'\'=f'))
    fig.add_trace(go.Scatter(x=x, y=f, name='f=|x|^alpha = -u\'\''))
    fig.update_layout(title=f'Ostrosc oszacowan Schaudera, alpha={alpha}',
                      xaxis_title='x', yaxis_title='wartosc')
    save(fig, 'pde16_holder_ostrosc.html')
    h = np.logspace(-4, -1, 30)
    beta = 0.8
    qa = h**alpha / h**alpha
    qb = h**alpha / h**beta
    fig = go.Figure()
    fig.add_trace(go.Scatter(x=h, y=qa, name='iloraz Holdera dla beta=alpha'))
    fig.add_trace(go.Scatter(x=h, y=qb, name='iloraz Holdera dla beta>alpha'))
    fig.update_layout(title='Wyzszy wykladnik Holdera jest nieosiagalny',
                      xaxis_type='log', yaxis_type='log',
                      xaxis_title='odleglosc h', yaxis_title='iloraz')
    save(fig, 'pde16_holder_ilorazy.html')
    assert np.allclose(qa, 1)
    assert qb[0] > qb[-1]


if __name__ == '__main__':
    demo_fft()
    demo_holder()
\end{lstlisting}

\section{Regularność De Giorgiego--Nasha--Mosera: przykład radialny — Plotly}
\label{sec:pde17-code}
Program wizualizuje słabe rozwiązanie o optymalnej regularności
$C^{0,\gamma}$, macierz współczynników oraz ilorazy Höldera.
Dodatkowo sprawdza wartości maksimum i minimum przykładowej
funkcji harmonicznej. Jest to eksperyment ilustracyjny,
nie dowód twierdzenia o regularności przy mierzalnych współczynnikach.
\begin{lstlisting}[style=pythonpde,caption={Mierzalne współczynniki, regularność Höldera i Harnack (Plotly)},label={code:pde17-plotly}]
"""PDE rozdzial 17: ostrosc Holdera przy mierzalnych wspolczynnikach.

Zaleznosci: pip install numpy plotly
Uruchomienie: python demo_degiorgi_moser_plotly.py
Wyniki: pde17_profil_holder.html, pde17_macierz.html,
        pde17_ilorazy_holdera.html, pde17_harnack_affine.html
"""
from pathlib import Path
import numpy as np
import plotly.graph_objects as go

OUT = Path(__file__).resolve().parent


def store(fig: go.Figure, filename: str) -> None:
    fig.update_layout(template="plotly_white", font=dict(size=13))
    fig.write_html(OUT / filename, include_plotlyjs=True, full_html=True)


def radial_example(gamma: float, n: int = 301):
    xx = np.linspace(-1.0, 1.0, n)
    x, y = np.meshgrid(xx, xx, indexing="xy")
    radius = np.hypot(x, y)
    theta_cos = np.divide(x, radius, out=np.zeros_like(x), where=radius > 0)
    u = radius**gamma * theta_cos
    u[radius > 1] = np.nan
    # A = gamma^2 I + (1-gamma^2) e_r \otimes e_r
    e1 = np.divide(x, radius, out=np.zeros_like(x), where=radius > 0)
    e2 = np.divide(y, radius, out=np.zeros_like(y), where=radius > 0)
    a11 = gamma**2 + (1.0-gamma**2) * e1**2
    a22 = gamma**2 + (1.0-gamma**2) * e2**2
    a12 = (1.0-gamma**2) * e1 * e2
    return xx, x, y, radius, u, a11, a12, a22


def main():
    gamma = 0.55
    xx, x, y, radius, u, a11, a12, a22 = radial_example(gamma)
    fig = go.Figure(go.Heatmap(x=xx, y=xx, z=u, colorscale="RdBu",
                               zmid=0, colorbar=dict(title="u")))
    fig.update_layout(title=f"Rozwiazanie r^gamma cos(theta), gamma={gamma}",
                      xaxis_title="x", yaxis_title="y", yaxis_scaleanchor="x")
    store(fig, "pde17_profil_holder.html")

    fig = go.Figure()
    fig.add_trace(go.Heatmap(x=xx, y=xx, z=np.where(radius <= 1, a12, np.nan),
                             colorscale="RdBu", zmid=0, colorbar=dict(title="a12")))
    fig.update_layout(title="Mierzalna macierz A: wspolczynnik a12",
                      xaxis_title="x", yaxis_title="y", yaxis_scaleanchor="x")
    store(fig, "pde17_macierz.html")

    r = np.logspace(-4, -1, 90)
    fig = go.Figure()
    for beta in [0.35, gamma, 0.75]:
        quotient = r**(gamma-beta)
        fig.add_trace(go.Scatter(x=r, y=quotient, mode="lines",
                                 name=f"beta={beta}"))
    fig.update_layout(title="Ilorazy Holdera w zerze: r^(gamma-beta)",
                      xaxis_title="r", yaxis_title="iloraz", xaxis_type="log",
                      yaxis_type="log")
    store(fig, "pde17_ilorazy_holdera.html")
    assert np.all(np.diff(r**(gamma-0.75)) < 0)

    rr = np.linspace(0.0, 0.95, 100)
    # harmonic positive v(x,y)=2+0.4x on B_1
    sup = 2.0 + 0.4 * rr
    inf = 2.0 - 0.4 * rr
    ratio = sup/inf
    assert np.max(ratio) < 1.5
    fig = go.Figure()
    fig.add_trace(go.Scatter(x=rr, y=sup, name="sup na B_r"))
    fig.add_trace(go.Scatter(x=rr, y=inf, name="inf na B_r"))
    fig.add_trace(go.Scatter(x=rr, y=ratio, name="sup/inf"))
    fig.update_layout(title="Przykladowa dodatnia funkcja harmoniczna i Harnack",
                      xaxis_title="promien r", yaxis_title="wartosc")
    store(fig, "pde17_harnack_affine.html")
    print("Zweryfikowano cztery wykresy i prosty test Harnacka.")


if __name__ == "__main__":
    main()
\end{lstlisting}


\section{Widmo eliptyczne: wartości własne, Weyl i ślad cieplny — Plotly}
\label{sec:pde18-code}
Poniższy program oblicza dyskretne wartości własne jednowymiarowego
laplasjanu Dirichleta, porównuje funkcję zliczającą kwadratowego
laplasjanu różnicowego z wyrazem głównym prawa Weyla oraz
wizualizuje ślad półgrupy cieplnej. Ostatnie porównanie wykorzystuje
jawne wartości własne ciągłego laplasjanu na kwadracie.
Przy ustalonej siatce nie należy interpretować końca widma
macierzy jako asymptotyki wysokich częstotliwości operatora ciągłego.
\begin{lstlisting}[style=pythonpde,caption={Wartości własne, prawo Weyla i ślad cieplny (Plotly)},label={code:pde18-plotly}]
"""PDE, rozdzial 18: min-max, Weyl, slad cieplny.

Zaleznosci: pip install numpy scipy plotly
Uruchom: python demo_widmo_weyl_plotly.py
Wyniki: pde18_wartosci_wlasne.html, pde18_funkcja_zliczajaca.html,
        pde18_slad_cieplny.html
"""
from pathlib import Path
import numpy as np
import plotly.graph_objects as go
from scipy.linalg import eigh_tridiagonal

OUT = Path(__file__).resolve().parent


def save(fig, filename):
    fig.update_layout(template='plotly_white', font=dict(size=13),
                      margin=dict(l=60, r=35, t=65, b=65))
    fig.write_html(OUT / filename, include_plotlyjs=True, full_html=True)


def dirichlet_interval():
    ns = np.array([12, 20, 32, 50, 80, 128, 200])
    errors = []
    for n in ns:
        h = 1.0 / (n + 1)
        vals = eigh_tridiagonal(np.full(n, 2 / h**2),
                                np.full(n - 1, -1 / h**2),
                                select='i', select_range=(0, 0),
                                eigvals_only=True)
        errors.append(abs(vals[0] - np.pi**2))
    err = np.asarray(errors)
    h = 1 / (ns + 1)
    fig = go.Figure()
    fig.add_trace(go.Scatter(x=h, y=err, mode='lines+markers',
                             name='Blad pierwszej wartosci wlasnej'))
    fig.add_trace(go.Scatter(x=h, y=err[-1] * (h/h[-1])**2,
                             mode='lines', name='Linia odniesienia: h^2'))
    fig.update_layout(title='Zbieznosc wartosci wlasnej laplasjanu (Dirichlet)',
                      xaxis_type='log', yaxis_type='log',
                      xaxis_title='Krok h', yaxis_title='Blad bezwzgledny')
    save(fig, 'pde18_wartosci_wlasne.html')
    slope = np.polyfit(np.log(h), np.log(err), 1)[0]
    assert 1.95 < slope < 2.05, f'Oczekiwany rzad 2, otrzymano {slope}'
    print('Rzad zbieznosci wartosci wlasnej:', round(slope, 6))


def weyl_square():
    m = 100
    h = 1 / (m + 1)
    n = np.arange(1, m+1)
    one_d = 4 * np.sin(np.pi * n/(2*(m+1)))**2 / h**2
    two_d = np.sort((one_d[:, None] + one_d[None, :]).ravel())
    cutoffs = np.linspace(0, 2400, 450)
    counts = np.searchsorted(two_d, cutoffs, side='right')
    leading = cutoffs / (4*np.pi)  # |Omega|=1, d=2
    fig = go.Figure()
    fig.add_trace(go.Scatter(x=cutoffs, y=counts, mode='lines',
                             line_shape='hv', name='Dyskretne zliczanie modow'))
    fig.add_trace(go.Scatter(x=cutoffs, y=leading, mode='lines',
                             name='Wyraz glowny Weyla: Lambda/(4 pi)'))
    fig.update_layout(title='Funkcja zliczajaca na kwadracie jednostkowym',
                      xaxis_title='Prog Lambda', yaxis_title='Liczba modow')
    save(fig, 'pde18_funkcja_zliczajaca.html')
    # m=100 makes low modes accurate; at high frequencies the grid
    # is fixed and hence Weyl is only illustrative.
    print('N_dyskretne(1000)=', np.searchsorted(two_d, 1000, side='right'))


def heat_trace():
    j = np.arange(1, 180)
    eigs = np.pi**2 * (j[:, None]**2 + j[None, :]**2)
    t = np.logspace(-3.3, -1.2, 170)
    exact = np.array([np.exp(-tau*eigs).sum() for tau in t])
    leading = 1 / (4*np.pi*t)  # |Omega|=1, d=2
    fig = go.Figure()
    fig.add_trace(go.Scatter(x=t, y=exact, mode='lines',
                             name='Slad cieplny (suma modow)'))
    fig.add_trace(go.Scatter(x=t, y=leading, mode='lines',
                             name='Asymptotyka 1/(4 pi t)'))
    fig.update_layout(title='Slad cieplny laplasjanu Dirichleta na kwadracie',
                      xaxis_type='log', yaxis_type='log',
                      xaxis_title='Czas t', yaxis_title='Tr exp(t Delta_D)')
    save(fig, 'pde18_slad_cieplny.html')
    ratio = exact[0] / leading[0]
    assert 0.8 < ratio < 1.0
    print('Stosunek sladu do wyrazu glownego dla malego t:', round(ratio, 6))


if __name__ == '__main__':
    dirichlet_interval()
    weyl_square()
    heat_trace()
\end{lstlisting}



\section{Pojemność kondensatora kulistego — Plotly}
\label{sec:pde19-code}
Poniższy program porównuje pojemności względne
$\operatorname{cap}_{2,B_1}(\overline B_r)$ w wymiarach $2$ i $3$.
Każda krzywa jest obliczana z dokładnego wzoru
\eqref{eq:pde19-ball-exact}; program nie zastępuje dowodu
minimalności energii ani nie numeryzuje równania Laplace'a.
Wykresy zapisywane są jako interaktywne pliki HTML.
\begin{lstlisting}[style=pythonpde,caption={Pojemność Sobolewa i tempo zaniku pojemności punktu (Plotly)},label={code:pde19-plotly}]
"""PDE, rozdzial 19: pojemnosc kondensatora kulistego.
Wymaga: numpy, plotly. Wyniki: interaktywne pliki HTML.
"""
from pathlib import Path
import numpy as np
import plotly.graph_objects as go

OUT = Path(__file__).resolve().parent
R = 1.0
r = np.geomspace(1e-6, 0.9, 240)
cap2 = 2 * np.pi / np.log(R / r)
cap3 = 4 * np.pi / (1 / r - 1 / R)


def save(fig, name):
    fig.update_layout(template="plotly_white", font=dict(size=13))
    fig.write_html(OUT / name, include_plotlyjs=True, full_html=True)


def condenser_plot(cap, d, name):
    fig = go.Figure()
    fig.add_trace(go.Scatter(x=r, y=cap, mode="lines",
                             name=f"pojemnosc d={d}"))
    fig.update_layout(title=f"Pojemnosc wzgledna kuli w B_1, d={d}",
                      xaxis_title="promien r", yaxis_title="energia minimalna")
    save(fig, name)


def comparison():
    fig = go.Figure()
    fig.add_trace(go.Scatter(x=r, y=cap2, name="d=2 (logarytm)"))
    fig.add_trace(go.Scatter(x=r, y=cap3, name="d=3 (potega)"))
    fig.update_layout(title="Zanik pojemnosci punktu, rozne tempa",
                      xaxis_type="log", yaxis_type="log",
                      xaxis_title="promien r", yaxis_title="pojemnosc")
    save(fig, "pde19_pojemnosc_porownanie.html")


def test_exact_values():
    assert np.all(np.diff(cap2) > 0)
    assert np.all(np.diff(cap3) > 0)
    assert np.isclose(2 * np.pi / np.log(2), 9.064720283654388)
    assert np.isclose(4 * np.pi / (2 - 1), 4 * np.pi)
    print("d=2, r=1/2:", 2 * np.pi / np.log(2))
    print("d=3, r=1/2:", 4 * np.pi)
    print("d=2, r=1e-6:", cap2[0])
    print("d=3, r=1e-6:", cap3[0])


if __name__ == "__main__":
    test_exact_values()
    condenser_plot(cap2, 2, "pde19_pojemnosc_2d.html")
    condenser_plot(cap3, 3, "pde19_pojemnosc_3d.html")
    comparison()
\end{lstlisting}


\section{Paraboliczne równania ewolucyjne: wsteczny Euler i stabilność — Plotly}
\label{sec:pde20-code}
Eksperyment porównuje niejawny schemat różnicowy z rozwiązaniem dokładnym
na odcinku, testuje rząd zbieżności przy $\tau\asymp h^2$ oraz
ilustruje tłumienie poszczególnych modów. W algorytmie stosujemy
dyskretną transformatę sinusową, która diagonalizuje laplasjan
z jednorodnymi warunkami Dirichleta. Jest to implementacja
wstecznego Eulera, a nie schematu dokładnej ewolucji macierzy.
\begin{lstlisting}[style=pythonpde,caption={Niejawny schemat cieplny, zbieżność i tłumienie modów (Plotly)},label={code:pde20-plotly}]
"""Iteracja 20, rownanie ciepa na (0,1) z Dirichletem.

Wsteczny Euler i centralny laplasjan roznicowy,
diagonalizacja macierzy przez dyskretna transformate sinusowa.
Instalacja: python -m pip install numpy scipy plotly
Uruchomienie: python demo_paraboliczne_plotly.py
"""
from pathlib import Path
import numpy as np
import plotly.graph_objects as go
from scipy.fft import dst, idst

OUT = Path(__file__).resolve().parent
T = 0.035


def initial(x):
    return np.sin(np.pi * x) + 0.5 * np.sin(3 * np.pi * x)


def exact(x, t):
    return np.exp(-np.pi**2 * t) * np.sin(np.pi*x) + 0.5 * np.exp(-9*np.pi**2*t) * np.sin(3*np.pi*x)


def backward_euler(n, at=T, cfl=0.32):
    """Rozwiazanie po nstep krokach: kazdy krok jest niejawny."""
    h = 1 / (n + 1)
    nstep = max(1, int(np.ceil(at / (cfl*h*h))))
    tau = at / nstep
    x = np.arange(1, n+1) * h
    j = np.arange(1, n+1)
    eig = 4 / h**2 * np.sin(j*np.pi/(2*(n+1)))**2
    coeffs = dst(initial(x), type=1, norm='ortho')
    ampl = (1 + tau*eig)**(-nstep)
    uh = idst(coeffs*ampl, type=1, norm='ortho')
    return x, uh, tau, eig


def save(fig, file):
    fig.update_layout(template='plotly_white', font=dict(size=13),
                      margin=dict(l=70, r=35, b=60, t=75))
    fig.write_html(OUT/file, include_plotlyjs=True, full_html=True)


def profiles():
    n=80
    h = 1/(n+1)
    x = np.arange(1,n+1)*h
    fig = go.Figure()
    for t in [0.0, 0.004, 0.012, T]:
        if t == 0:
            u=initial(x)
        else:
            _, u, _, _ = backward_euler(n, at=t)
        fig.add_trace(go.Scatter(x=x, y=u, mode='lines', name=f't={t:g}'))
    fig.update_layout(title='Rownanie ciepa: wygadzanie i zanikanie modow',
                      xaxis_title='x',yaxis_title='u(t,x)')
    save(fig, 'pde20_profile_plotly.html')
    # Zanik nie zwieksza normy w l2 dla danych bez wymuszenia.
    assert np.linalg.norm(backward_euler(n)[1]) <= np.linalg.norm(initial(x))


def convergence():
    sizes=np.array([12,18,27,40,60,90,135])
    errs=[]
    for n in sizes:
        x, u, _, _ = backward_euler(int(n))
        h = 1/(n+1)
        errs.append(np.sqrt(h*np.sum((u-exact(x,T))**2)))
    hs=1/(sizes+1)
    errs=np.array(errs)
    slope=np.polyfit(np.log(hs),np.log(errs),1)[0]
    print('Wsteczny Euler / laplasjan roznicowy: rzad bedu',round(slope,5))
    assert 1.85 < slope < 2.12, f'Nieoczekiwany rzad: {slope}'
    fig=go.Figure()
    fig.add_trace(go.Scatter(x=hs,y=errs,mode='lines+markers',name='Bad w L2_h'))
    fig.add_trace(go.Scatter(x=hs,y=errs[-1]*(hs/hs[-1])**2,name='Odniesienie: h2'))
    fig.update_layout(title=f'Zbieznosc dla 0.32h2, rzad {slope:.3f}',
                      xaxis_type='log',yaxis_type='log',
                      xaxis_title='h',yaxis_title='||U-u||_(L2_h)')
    save(fig, 'pde20_zbieznosc_plotly.html')


def damping():
    z=np.logspace(-3,2,500)
    fig=go.Figure()
    fig.add_trace(go.Scatter(x=z,y=np.exp(-z),name='Dokadny: exp(-z)'))
    fig.add_trace(go.Scatter(x=z,y=1/(1+z),name='Wsteczny Euler: 1/(1+z)'))
    fig.add_trace(go.Scatter(x=z,y=np.abs(1-z),name='Jawny Euler: |1-z|'))
    fig.update_layout(title='Tumienie jednego modu; z=',
                      xaxis_title='z',yaxis_title='wspoczynnik amplitudy',
                      xaxis_type='log',yaxis_type='log')
    save(fig,'pde20_tlumienie_plotly.html')
    assert np.all(1/(1+z)<=1)
    assert np.any(np.abs(1-z)>1)


if __name__ == '__main__':
    profiles()
    convergence()
    damping()
\end{lstlisting}


\section{Układ symetryczno-hiperboliczny: schemat pod prąd — Plotly}
\label{sec:pde21-code}
Program oblicza ewolucję dwóch modów o różnych prędkościach,
porównuje rozwiązanie numeryczne z dokładnym,
sprawdza nierosnącą energię dyskretną i wyznacza rząd
zbieżności jawnego schematu pod prąd. Rozdzielenie macierzy
$A=A_++A_-$ odbywa się przez diagonalizację ortogonalną,
a warunek CFL ma postać $\tau\max_k|a_k|/h\le1$.
\begin{lstlisting}[style=pythonpde,caption={Hiperboliczny układ, energia i zbieżność: Plotly},label={code:pde21-plotly}]
"""Rozdzial 21: dwa przemieszczajace sie mody ukadu symetrycznego.
Zaleznosci: numpy, plotly; uruchomienie: python demo_hiperboliczne_plotly.py
Schemat pod prad z rozdzieleniem macierzy A=A_++A_-.
"""
from pathlib import Path
import numpy as np
import plotly.graph_objects as go

OUT = Path(__file__).resolve().parent
THETA = 0.48
Q = np.array([[np.cos(THETA), -np.sin(THETA)],
              [np.sin(THETA), np.cos(THETA)]])
SPEEDS = np.array([-1.0, 0.6])
A = Q @ np.diag(SPEEDS) @ Q.T
A_PLUS = Q @ np.diag(np.maximum(SPEEDS, 0.0)) @ Q.T
A_MINUS = Q @ np.diag(np.minimum(SPEEDS, 0.0)) @ Q.T


def wrapped_gauss(x, center, width):
    y = (x - center + 0.5) % 1.0 - 0.5
    return np.exp(-(y / width) ** 2)


def initial(x):
    w = np.vstack([wrapped_gauss(x, 0.25, 0.065),
                   -0.8 * wrapped_gauss(x, 0.71, 0.075)])
    return Q @ w


def exact(x, t):
    w = np.vstack([wrapped_gauss(x - SPEEDS[0] * t, 0.25, 0.065),
                   -0.8 * wrapped_gauss(x - SPEEDS[1] * t, 0.71, 0.075)])
    return Q @ w


def upwind(n, T=0.35, cfl=0.8, record=False):
    h = 1.0 / n
    x = np.arange(n) * h
    steps = int(np.ceil(T * np.max(abs(SPEEDS)) / (cfl * h)))
    dt = T / steps
    U = initial(x)
    en = [h * np.sum(U * U)]
    for _ in range(steps):
        left = U - np.roll(U, 1, axis=1)
        right = np.roll(U, -1, axis=1) - U
        U = U - dt / h * (A_PLUS @ left + A_MINUS @ right)
        if record:
            en.append(h * np.sum(U * U))
    ref = exact(x, T)
    err = np.sqrt(h * np.sum((U - ref) ** 2))
    return x, U, ref, err, np.array(en), dt, h


def export(fig, name):
    fig.update_layout(template='plotly_white', font=dict(size=13),
                      margin=dict(l=75, r=35, b=55, t=75))
    fig.write_html(OUT / name, include_plotlyjs=True, full_html=True)


def main():
    n = 240
    x, U, target, err, energy, dt, h = upwind(n, record=True)
    fig = go.Figure()
    for k in range(2):
        fig.add_trace(go.Scatter(x=x, y=U[k], mode='lines',
                     name=f'numeryczne u{k+1}'))
        fig.add_trace(go.Scatter(x=x, y=target[k], mode='lines',
                     line=dict(dash='dash'), name=f'dokladne u{k+1}'))
    fig.update_layout(title=f'Ukad symetryczny: mody o predkosciach -1 i 0.6; blad={err:.3g}',
                      xaxis_title='x', yaxis_title='u(t,x)')
    export(fig, 'pde21_profile_plotly.html')

    fig = go.Figure()
    fig.add_trace(go.Scatter(x=np.arange(energy.size) * dt, y=energy,
                             mode='lines', name='Dyskretna energia l2'))
    fig.update_layout(title='Kontrakcja energii: schemat pod prad, CFL=0.8',
                      xaxis_title='czas', yaxis_title='h sum |U_j|^2')
    export(fig, 'pde21_energia_plotly.html')
    assert np.max(np.diff(energy)) <= 1e-12, 'Energia powinna byc nierosnaca'
    assert np.allclose(A, A.T)

    sizes = np.array([40, 60, 90, 135, 200, 300, 450])
    errors = np.array([upwind(int(k))[3] for k in sizes])
    hs = 1.0 / sizes
    slope = np.polyfit(np.log(hs[-4:]), np.log(errors[-4:]), 1)[0]
    print('Rzad zbieznosci, cztery ostatnie siatki:', round(slope, 4))
    print('Blad l2 dla N=240:', round(err, 7))
    assert 0.65 < slope < 1.15, f'Nieoczekiwany rzad {slope}'
    fig = go.Figure()
    fig.add_trace(go.Scatter(x=hs, y=errors, mode='lines+markers', name='Blad l2'))
    fig.add_trace(go.Scatter(x=hs, y=errors[-1] * hs / hs[-1], mode='lines',
                             line=dict(dash='dash'), name='Skala O(h)'))
    fig.update_layout(title=f'Zbieznosc schematu pod prad: rzad {slope:.3f}',
                      xaxis_title='h', yaxis_title='blad l2',
                      xaxis_type='log', yaxis_type='log')
    export(fig, 'pde21_zbieznosc_plotly.html')


if __name__ == '__main__':
    main()
\end{lstlisting}

\section{Swobodna ewolucja Schrödingera na torusie — Plotly}
\label{sec:pde22-code}
Kod stosuje reprezentację spektralną na torusie, a nie swobodne
jądro propagatora na $\R^d$. Sprawdza zachowanie normy
$L^2$, ruch pakietu i wzór Parsevala na błąd obcięcia.

\begin{lstlisting}[style=pythonpde,caption={Ewolucja Schrodingera: pakiet, faza i błąd widmowy (Plotly)},label={code:pde22-plotly}]
"""Rozdzial 22: swobodny propagator Schrodingera na torusie.

Zaleznosci: numpy, plotly
Uruchomienie: python demo_schrodinger_plotly.py
Wyniki: pde22_pakiet_plotly.html, pde22_faza_plotly.html,
        pde22_obciecie_plotly.html
"""
from pathlib import Path
import numpy as np
import plotly.graph_objects as go

OUT = Path(__file__).resolve().parent
N = 512
x = -np.pi + (2 * np.pi / N) * np.arange(N)
k = np.arange(-100, 101)
k0, sigma, x0 = 12.0, 5.0, -1.0
coef = np.exp(-0.5 * ((k - k0) / sigma)**2) * np.exp(-1j * k * x0)

def psi(t, cutoff=None):
    mask = np.ones(k.shape, dtype=bool) if cutoff is None else np.abs(k) <= cutoff
    z = (coef[mask] * np.exp(-1j * t * k[mask]**2))[:, None]
    return np.sum(z * np.exp(1j * k[mask, None] * x[None, :]), axis=0) / np.sqrt(2 * np.pi)

def store(fig, name):
    fig.update_layout(template="plotly_white", font=dict(size=13),
                      margin=dict(l=70, r=30, t=70, b=60))
    fig.write_html(OUT / name, include_plotlyjs=True, full_html=True)

def main():
    times = [0.0, 0.015, 0.03, 0.06, 0.09]
    fig = go.Figure()
    for t in times:
        p = psi(t)
        fig.add_trace(go.Scatter(x=x, y=np.abs(p)**2, mode='lines',
                                 name=f't={t:.3f}'))
    fig.update_layout(title='Pakiet swobodny na torusie: propagacja i dyspersja',
                      xaxis_title='x', yaxis_title='|u(t,x)|^2')
    store(fig, 'pde22_pakiet_plotly.html')

    # Przestrzenna siatka nie zawiera prawego konca okresu:
    # okresowa kwadratura trapezowa jest dokladna dla pasmowych iloczynow.
    norms_periodic = [(2*np.pi/N)*np.sum(np.abs(psi(t))**2) for t in times]
    parseval = np.sum(np.abs(coef)**2)
    print('Maksymalna zmiana normy (kwadrat):',
          max(abs(v - parseval) for v in norms_periodic))
    assert max(abs(v-parseval) for v in norms_periodic) < 1e-10

    t = 0.04
    fig = go.Figure()
    fig.add_trace(go.Scatter(x=k, y=np.abs(coef), name='|c_k|', mode='lines'))
    fig.add_trace(go.Scatter(x=k, y=np.angle(np.exp(-1j*t*k*k)),
                             name='faza propagatora', yaxis='y2', mode='markers',
                             marker=dict(size=3)))
    fig.update_layout(title='Zachowany modul i zmieniajaca sie faza',
                      xaxis_title='k', yaxis_title='|c_k|',
                      yaxis2=dict(title='faza [rad]', overlaying='y', side='right'))
    store(fig, 'pde22_faza_plotly.html')

    full = psi(0.05)
    cutoffs = [6, 10, 14, 18, 24, 32, 48, 64]
    errors = [(2*np.pi/N*np.sum(np.abs(full-psi(0.05,K))**2))**0.5
              for K in cutoffs]
    exact_tails = [np.sqrt(np.sum(np.abs(coef[np.abs(k)>K])**2))
                   for K in cutoffs]
    assert np.allclose(errors, exact_tails, rtol=1e-7, atol=1e-12)
    fig = go.Figure()
    fig.add_trace(go.Scatter(x=cutoffs, y=errors, name='blad numeryczny',
                             mode='lines+markers'))
    fig.add_trace(go.Scatter(x=cutoffs, y=exact_tails, name='ogon Parsevala',
                             mode='lines', line=dict(dash='dash')))
    fig.update_layout(title='Obciecie szeregu i tozsamosc Parsevala',
                      xaxis_title='maksymalne |k|', yaxis_title='blad L2',
                      yaxis_type='log')
    store(fig, 'pde22_obciecie_plotly.html')

if __name__ == '__main__':
    main()
\end{lstlisting}



\section{Równanie Hamiltona--Jacobiego: charakterystyki i rozwiązania lepkościowe (Plotly)}
Kod oblicza wartości wzoru Hopfa--Laxa przez minimalizację na siatce,
rysuje charakterystyki i pochodną rozwiązania po powstaniu osobliwości.
Wykresy ilustrują teorię; nie zastępują dowodów z Rozdziału~\ref{chap:pde23}.
\begin{lstlisting}[style=pythonpde,caption={Wzór Hopfa--Laxa i powstawanie narożnika (Plotly)},label={code:pde23-plotly}]
from pathlib import Path
import numpy as np
import plotly.graph_objects as go
OUT = Path(__file__).resolve().parent
def hopf_lax(t: float, x: np.ndarray, y: np.ndarray):
    """Minimalizacja dla g(y)=cos(y) i H(p)=p^2/2."""
    cost = np.cos(y)[None, :] + (x[:, None] - y[None, :])**2 / (2 * t)
    argmin = np.argmin(cost, axis=1)
    return cost[np.arange(len(x)), argmin], y[argmin]
def save(fig, filename):
    fig.update_layout(template="plotly_white", font=dict(size=13))
    fig.write_html(OUT / filename, include_plotlyjs=True, full_html=True)
def main():
    x = np.linspace(-np.pi, np.pi, 601)
    y = np.linspace(-np.pi - 2.2, np.pi + 2.2, 1801)
    fig = go.Figure()
    for t in (0.2, 0.8, 1.02, 1.35):
        u, minimizer = hopf_lax(t, x, y)
        fig.add_scatter(x=x, y=u, mode="lines", name=f"t={t:.2f}")
    fig.add_scatter(x=x, y=np.cos(x), mode="lines", name="t=0", line=dict(dash="dash"))
    fig.update_layout(title="Rozwiazanie lepkosciowe: g(x)=cos x, H(p)=p^2/2",
                      xaxis_title="x", yaxis_title="u(t,x)")
    save(fig, "pde23_profile_plotly.html")
    yy = np.linspace(-np.pi, np.pi, 35)
    tt = np.linspace(0, 1.6, 125)
    fig = go.Figure()
    for a in yy:
        fig.add_scatter(x=a - tt * np.sin(a), y=tt, mode="lines",
                        showlegend=False, line=dict(width=1.2))
    fig.add_hline(y=1, line_dash="dash", annotation_text="pierwsza degeneracja: t=1")
    fig.update_layout(title="Charakterystyki: x(t,y)=y-t sin y",
                      xaxis_title="poozenie x", yaxis_title="czas t")
    save(fig, "pde23_charakterystyki_plotly.html")
    t = 1.35
    dx = x[1] - x[0]
    u, minimizer = hopf_lax(t, x, y)
    grad = np.gradient(u, dx)
    fig = go.Figure()
    fig.add_scatter(x=x, y=grad, mode="lines", name="gradient numeryczny")
    fig.add_scatter(x=x, y=(x - minimizer) / t, mode="lines",
                    name="(x-y_min)/t", line=dict(dash="dash"))
    fig.update_layout(title="Skok gradientu rozwiazania po t=1",
                      xaxis_title="x", yaxis_title="pochodna przestrzenna")
    save(fig, "pde23_skok_gradientu_plotly.html")
    k = np.argmin(np.abs(x))
    xx = x[k]
    costs = np.cos(y) + (xx - y)**2 / (2*t)
    near = np.where(costs <= costs.min() + 1e-4)[0]
    assert np.any(y[near] < -0.4) and np.any(y[near] > 0.4)
    assert np.all(1 - 0.8*np.cos(yy) > 0)
    print("Trzy wizualizacje zapisane. Test symetrycznych minimow: OK.")
if __name__ == "__main__":
    main()
\end{lstlisting}


\section{Skalarne prawa zachowania: rozwiązania Riemanna — Plotly}
\label{sec:pde24-code}
Program implementuje monotoniczny schemat Godunowa dla strumienia
$f(u)=u^2/2$, porównuje go ze ścisłymi rozwiązaniami
entropijnymi zagadnienia Riemanna i wyznacza błąd w normie
$L^1$ na skończonym przedziale obliczeniowym.
Dane Riemanna nie są całkowalne na całej prostej; eksperyment
wykonujemy w skończonym oknie przestrzennym.
\begin{lstlisting}[style=pythonpde,caption={Fala uderzeniowa, fala rozrzedzenia i metoda Godunowa — Plotly},label={code:pde24-plotly}]
"""Skalarne prawa zachowania: fala uderzeniowa, rozrzedzenie, zbieznosc.

Uruchomienie: python demo_prawa_zachowania_plotly.py
Wymagania: numpy, plotly
"""
from pathlib import Path
import numpy as np
import plotly.graph_objects as go

OUT = Path(__file__).resolve().parent


def flux(u):
    return 0.5 * np.asarray(u) ** 2


def godunov(left, right):
    """Strumien Godunowa dla wypukego f(u)=u^2/2."""
    left = np.asarray(left)
    right = np.asarray(right)
    rarefaction = np.where(left >= 0, flux(left),
                          np.where(right <= 0, flux(right), 0.0))
    shock = np.maximum(flux(left), flux(right))
    return np.where(left <= right, rarefaction, shock)


def riemann_exact(x, t, left, right):
    if left > right:
        speed = (left + right) / 2
        return np.where(x < speed*t, left, right)
    if t == 0:
        return np.where(x < 0, left, right)
    return np.where(x <= left*t, left,
                    np.where(x >= right*t, right, x/t))


def numerical(n, tfinal, left, right):
    dx = 6.0/n
    x = -3.0 + (np.arange(n) + 0.5) * dx
    u = np.where(x < 0, left, right).astype(float)
    t = 0.0
    while t < tfinal - 1e-13:
        max_speed = max(1.0, np.max(np.abs(u)))
        dt = min(0.43*dx/max_speed, tfinal-t)
        padded = np.concatenate(([left], u, [right]))
        numerical_flux = godunov(padded[:-1], padded[1:])
        u -= dt/dx * (numerical_flux[1:] - numerical_flux[:-1])
        t += dt
    return x, u, dx


def graph(filename, title, traces, xaxis='x', yaxis='u'):
    fig = go.Figure()
    for x, y, name, style in traces:
        fig.add_trace(go.Scatter(x=x, y=y, name=name, mode='lines',
                                 line=dict(dash=style)))
    fig.update_layout(template='plotly_white', title=title,
                      xaxis_title=xaxis, yaxis_title=yaxis)
    if filename.endswith('zbieznosc.html'):
        fig.update_xaxes(type='log')
        fig.update_yaxes(type='log')
    fig.write_html(OUT / filename, include_plotlyjs=True, full_html=True)


def main():
    t = 0.8
    for left, right, name in [(1.0, 0.0, 'uderzenie'),
                               (0.0, 1.0, 'rozrzedzenie')]:
        x, u, dx = numerical(600, t, left, right)
        exact = riemann_exact(x, t, left, right)
        graph(f'pde24_{name}.html', f'Burgers: {name} (t={t})',
              [(x, exact, 'Rozwiazanie entropijne', 'solid'),
               (x, u, 'Schemat Godunowa', 'dash')])
        print(name, 'L1 error:', np.sum(abs(u-exact))*dx)

    ns = np.array([100, 200, 400, 800, 1600])
    trace_data = []
    for left, right, name, dash in [(1.0, 0.0, 'Uderzenie', 'solid'),
                                   (0.0, 1.0, 'Rozrzedzenie', 'dash')]:
        errors = []
        for n in ns:
            x, u, dx = numerical(int(n), t, left, right)
            errors.append(np.sum(abs(u-riemann_exact(x,t,left,right)))*dx)
        errors = np.asarray(errors)
        rate = np.polyfit(np.log(6.0/ns), np.log(errors), 1)[0]
        trace_data.append((6.0/ns, errors, f'{name}, nachylenie {rate:.2f}', dash))
        print(f'{name}: measured rate {rate:.4f}')
    graph('pde24_zbieznosc.html', 'Schemat Godunowa: bad L1',
          trace_data, xaxis='Krok siatki h', yaxis='Bad L1')

    # Kontrola zachowania ograniczenia maksimum przez monotoniczny schemat.
    x, u, _ = numerical(800, t, 1.0, 0.0)
    assert np.all(u >= -1e-12) and np.all(u <= 1+1e-12)


if __name__ == '__main__':
    main()
\end{lstlisting}


\section{Półliniowe równania ewolucyjne: energia i wybuch — Plotly}
\label{sec:pde25-code}
Program ilustruje spadek energii w równaniu Allena--Cahna
oraz dolną barierę na średnią rozwiązania równania
$u_t=\Delta u+u^2$ na torusie. Schemat Stranga jest
eksperymentem numerycznym, a nie dowodem twierdzeń.
\begin{lstlisting}[style=pythonpde,caption={Równania półliniowe: energia i eksplozja (Plotly)},label={code:pde25-plotly}]
"""PDE, rozdzial 25: reakcja, energia i eksplozja na torusie.

Uruchomienie: python demo_pololiniowe_plotly.py
Zaleznosci: numpy, plotly
Zapisuje trzy samodzielne wykresy interaktywne do katalogu skryptu.
"""
from pathlib import Path
import numpy as np
import plotly.graph_objects as go

OUT = Path(__file__).resolve().parent
N = 512
x = 2 * np.pi * np.arange(N) / N
k = np.fft.fftfreq(N, d=1 / N)
dx = 2 * np.pi / N

def diffuse(u, dt):
    return np.fft.ifft(np.fft.fft(u) * np.exp(-(k * k) * dt)).real

def save(fig, name):
    fig.update_layout(template="plotly_white", font=dict(size=13))
    fig.write_html(OUT / name, include_plotlyjs=True, full_html=True)

def allen_cahn():
    # Splitting Stranga: dyfuzja, dokladny lokalny przeplyw u'=u-u^3,
    # dyfuzja. Dokladny przeplyw reakcji zachowuje znak.
    u = 0.17 + 0.72 * np.cos(x) + 0.14 * np.cos(3 * x)
    dt = 0.001
    snapshots = {0.0: u.copy()}
    ts, energies = [], []
    def energy(v):
        vx = np.fft.ifft(1j * k * np.fft.fft(v)).real
        return dx * np.sum(0.5 * vx**2 + 0.25 * (v*v - 1)**2)
    for j in range(1601):
        if j % 20 == 0:
            ts.append(j * dt)
            energies.append(energy(u))
        if j in (400, 800, 1600):
            snapshots[round(j * dt, 2)] = u.copy()
        if j == 1600:
            break
        u = diffuse(u, dt / 2)
        e2 = np.exp(2 * dt)
        u = e2**0.5 * u / np.sqrt(1 + u*u * (e2 - 1))
        u = diffuse(u, dt / 2)
    # Wybieramy dostatecznie dokladny krok, by energia byla nierosnaca.
    assert np.max(np.diff(energies)) < 1e-6
    fig = go.Figure()
    for t, v in snapshots.items():
        fig.add_trace(go.Scatter(x=x, y=v, name=f"t={t:g}"))
    fig.update_layout(title="Rownanie Allena-Cahna: wygadzanie i reakcja", xaxis_title="x", yaxis_title="u(t,x)")
    save(fig, "pde25_allen_cahn_profile.html")
    fig2 = go.Figure(go.Scatter(x=ts, y=energies, name="E(u(t))"))
    fig2.update_layout(title="Rownanie Allena-Cahna: dyssypacja energii", xaxis_title="czas t", yaxis_title="energia")
    save(fig2, "pde25_allen_cahn_energia.html")
    print("Rownanie Allena-Cahna: energia poczatkowa i koncowa:", energies[0], energies[-1])

def blowup_mean():
    # u_t=Delta u+u^2, u0=0.35+0.15 cos(x), dodatnia srednia 0.35.
    u = 0.35 + 0.15 * np.cos(x)
    m0 = np.mean(u)
    bound = 1.0 / m0
    dt = 0.0005
    tmax = 2.2
    ts, means, mins, maxs = [], [], [], []
    for j in range(int(round(tmax/dt)) + 1):
        if j % 40 == 0:
            ts.append(j * dt)
            means.append(float(np.mean(u)))
            mins.append(float(np.min(u)))
            maxs.append(float(np.max(u)))
        if j == int(round(tmax/dt)):
            break
        u = diffuse(u, dt / 2)
        assert dt * np.max(u) < 1
        u = u / (1 - dt * u)  # scisly przeplyw reakcji u'=u^2
        u = diffuse(u, dt / 2)
    ts = np.asarray(ts)
    barrier = m0 / (1 - m0 * ts)
    # Blad splittigu i kwadratury dopuszczamy ponizej 1e-3.
    assert np.min(np.asarray(means) - barrier) > -1e-3
    fig = go.Figure()
    fig.add_trace(go.Scatter(x=ts, y=means, name="srednia numeryczna"))
    fig.add_trace(go.Scatter(x=ts, y=barrier, name="bariera Jensena / ODE", line=dict(dash="dash")))
    fig.add_trace(go.Scatter(x=ts, y=maxs, name="maksimum"))
    fig.update_layout(title=f"u_t=Delta u + u^2: srednia i dolna bariera, T*  {bound:.3f}", xaxis_title="czas t", yaxis_title="wartosc")
    save(fig, "pde25_eksplozja_srednia.html")
    print("Bariera T* <=", bound, "roznica sredniej i bariery w koncu", means[-1] - barrier[-1])

if __name__ == "__main__":
    allen_cahn()
    blowup_mean()

\end{lstlisting}



\section{Nieliniowy strumień p-laplasjanu: przykłady — Plotly}
\label{sec:pde26-code}
Program rozwiązujący jednowymiarowe zagadnienie $-\Delta_pu=1$
z zerowymi warunkami Dirichleta konstruuje dyskretny strumień
z równania bilansu i odtwarza węzłowe wartości rozwiązania.
Osobno mierzymy błąd i wzrost przybliżonej drugiej pochodnej
w maksimum. Przy $p>2$ rozwiązanie jest klasy $C^1$,
lecz nie musi należeć do $C^2$. Obliczenia stanowią
ilustrację, a nie dowód twierdzenia Minty'ego--Browdera.
\begin{lstlisting}[style=pythonpde,caption={Nieliniowy strumień, błąd aproksymacji i krzywizna (Plotly)},label={code:pde26-plotly}]
"""PDE 26: p-Laplacian Dirichlet, finite-volume flux and Plotly.
Requires: numpy, plotly. Run: python demo_plaplasjan_plotly.py
"""
from pathlib import Path
import numpy as np
import plotly.graph_objects as go

OUT = Path(__file__).resolve().parent
PS = (1.5, 2.0, 3.0, 5.0)


def exact(x, p):
    exponent = p / (p - 1)
    return ((p - 1) / p) * (0.5**exponent - np.abs(x - 0.5)**exponent)


def discrete(p, n):
    """Solve -D_h(|D_h u|^(p-2) D_h u)=1 with zero end values.
    In 1D the numerical flux is q_{i+1/2}=1/2-(i+1/2)h.
    """
    x = np.linspace(0, 1, n + 1)
    h = 1.0 / n
    xm = (x[1:] + x[:-1]) / 2
    flux = 0.5 - xm
    grad = np.sign(flux) * np.abs(flux)**(1 / (p - 1))
    values = np.r_[0.0, np.cumsum(h * grad)]
    assert abs(values[-1]) < 5e-13
    q_recovered = np.sign(np.diff(values) / h) * np.abs(np.diff(values) / h)**(p - 1)
    residual = -(np.diff(q_recovered) / h) - 1
    assert np.max(np.abs(residual)) < 1e-8
    return x, values


def save(fig, filename):
    fig.update_layout(template='plotly_white', font=dict(size=13))
    fig.write_html(OUT / filename, include_plotlyjs=True, full_html=True)


def make_profiles():
    x_dense = np.linspace(0, 1, 1001)
    fig = go.Figure()
    for p in PS:
        xd, ud = discrete(p, 16)
        fig.add_trace(go.Scatter(x=x_dense, y=exact(x_dense, p),
                                 mode='lines', name=f'p={p:g}, exact'))
        fig.add_trace(go.Scatter(x=xd, y=ud, mode='markers',
                                 name=f'p={p:g}, grid n=16'))
    fig.update_layout(title='One-dimensional p-Laplacian: -Delta_p u = 1',
                      xaxis_title='x', yaxis_title='u(x)')
    save(fig, 'pde26_profile.html')


def make_convergence():
    ns = np.array((20, 40, 80, 160, 320, 640))
    fig = go.Figure()
    for p in (1.5, 3.0, 5.0):
        errors = []
        for n in ns:
            x, u = discrete(p, int(n))
            errors.append(np.max(np.abs(u - exact(x, p))))
        errors = np.asarray(errors)
        rate = -np.polyfit(np.log(ns), np.log(errors), 1)[0]
        print('p=', p, 'convergence rate=', float(rate))
        assert errors[-1] < errors[0]
        fig.add_trace(go.Scatter(x=1/ns, y=errors, mode='lines+markers',
                                 name=f'p={p:g}, slope={rate:.2f}'))
    fig.update_layout(title='Error of nonlinear flux discretization',
                      xaxis_title='grid step h', yaxis_title='max nodal error',
                      xaxis_type='log', yaxis_type='log')
    save(fig, 'pde26_zbieznosc.html')


def make_curvature():
    ns = np.array((20, 40, 80, 160, 320, 640))
    fig = go.Figure()
    for p in (2.0, 3.0, 5.0):
        alpha = p / (p - 1)
        c = (p - 1) / p
        hs = 1 / ns
        # Exact centered second difference of the continuum solution.
        curvature = 2 * c * hs**(alpha - 2)
        assert curvature[-1] >= curvature[0] - 1e-12
        fig.add_trace(go.Scatter(x=hs, y=curvature, mode='lines+markers',
                                 name=f'p={p:g}, power={alpha-2:.3f}'))
    fig.update_layout(title='Growth of the centered second difference at x=1/2',
                      xaxis_title='h', yaxis_title='-D_h^2 u(1/2)',
                      xaxis_type='log', yaxis_type='log')
    save(fig, 'pde26_krzywizna.html')


if __name__ == '__main__':
    make_profiles()
    make_convergence()
    make_curvature()
    print('Discrete flux residuals and endpoint conditions verified.')
\end{lstlisting}


\section{Reakcja--dyfuzja: fale biegnące, energia i niestabilność Turinga (Plotly)}
\label{sec:pde27-code}
Program symuluje fale biegnące dla równania Fishera--KPP,
rozpad energii w równaniu Allena--Cahna i wzrost modów
w niestabilności Turinga. Wizualizacje ilustrują wyniki analityczne;
nie zastępują ich dowodów.
\begin{lstlisting}[style=pythonpde,caption={Fale reakcji--dyfuzji, energia i niestabilność Turinga (Plotly)},label={code:pde27-plotly}]
"""PDE 27: Fisher--KPP, fala biegnaca, Allen--Cahn i Turing.

Uruchomienie: pip install numpy plotly; python demo_reakcja_dyfuzja_plotly.py
Wyniki: pde27_fronty.html, pde27_fala_jawna.html,
        pde27_energia.html, pde27_turing.html
"""
from pathlib import Path
import numpy as np
import plotly.graph_objects as go

OUT = Path(__file__).resolve().parent


def save(fig, filename):
    fig.update_layout(template="plotly_white", font=dict(size=12),
                      legend=dict(orientation="h", y=-0.22))
    fig.write_html(OUT / filename, include_plotlyjs=True, full_html=True)


def reaction_kpp(u, dt, r=1.0):
    e = np.exp(r * dt)
    return u * e / (1 + u * (e - 1))


def split_step_kpp():
    n, L, dt, tmax = 1024, 120.0, 0.04, 16.0
    x = np.linspace(-L/2, L/2, n, endpoint=False)
    k = 2 * np.pi * np.fft.fftfreq(n, d=L/n)
    diffusion_half = np.exp(-k*k*dt/2)
    u = .8 * np.exp(-x*x/6)
    fig = go.Figure()
    fig.add_scatter(x=x, y=u, name="t=0")
    for j in range(1, round(tmax/dt)+1):
        u = np.fft.ifft(np.fft.fft(u) * diffusion_half).real
        u = reaction_kpp(np.clip(u, 0, 1), dt)
        u = np.fft.ifft(np.fft.fft(u) * diffusion_half).real
        u = np.clip(u, 0, 1)
        if j % 100 == 0:
            fig.add_scatter(x=x, y=u, name=f"t={j*dt:g}")
    assert np.min(u) >= 0 and np.max(u) <= 1
    fig.update_layout(title="Rownanie Fishera-KPP: rozprzestrzenianie danych skupionych",
                      xaxis_title="x", yaxis_title="u(t,x)")
    save(fig, 'pde27_fronty.html')


def explicit_front():
    z = np.linspace(-28, 28, 1601)
    w = np.exp(z / np.sqrt(6))
    u = (1+w)**-2
    up = -2*w/(np.sqrt(6)*(1+w)**3)
    upp = -w/(3*(1+w)**3)+w*w/(1+w)**4
    c = 5/np.sqrt(6)
    residual = upp+c*up+u*(1-u)
    assert np.max(np.abs(residual)) < 5e-15
    fig = go.Figure()
    for t in (0, 3, 6, 9):
        x = z+c*t
        fig.add_scatter(x=x, y=u, name=f"t={t}")
    fig.update_layout(title=f"Fala rownania Fishera-KPP: c=5/sqrt6~{c:.4f}, residuum < 5e-15",
                      xaxis_title="x", yaxis_title="U(x-ct)")
    save(fig, 'pde27_fala_jawna.html')
    return float(np.max(np.abs(residual)))


def reaction_allen(u, dt):
    e = np.exp(dt)
    return u*e/np.sqrt(1+u*u*(e*e-1))


def allen_cahn_energy():
    n, L, D, dt = 512, 20.0, .3, .02
    x = np.linspace(0, L, n, endpoint=False)
    k = 2*np.pi*np.fft.fftfreq(n, d=L/n)
    half = np.exp(-D*k*k*dt/2)
    u = .6*np.cos(2*np.pi*x/L)+.15*np.cos(6*np.pi*x/L)
    energy, time = [], []
    for j in range(601):
        if j % 5 == 0:
            ux = np.fft.ifft(1j*k*np.fft.fft(u)).real
            E = (L/n)*np.sum(D*ux**2/2+(u*u-1)**2/4)
            energy.append(E)
            time.append(j*dt)
        if j < 600:
            u = np.fft.ifft(np.fft.fft(u)*half).real
            u = reaction_allen(u, dt)
            u = np.fft.ifft(np.fft.fft(u)*half).real
    energy = np.asarray(energy)
    assert energy[-1] < energy[0] and np.all(np.isfinite(energy))
    fig = go.Figure(go.Scatter(x=time, y=energy, mode="lines+markers", name="E(u)"))
    fig.update_layout(title="Rownanie Allena-Cahna: dyssypacja energii (rozszczepienie Stranga)",
                      xaxis_title="t", yaxis_title="energia")
    save(fig, 'pde27_energia.html')
    return float(np.max(np.diff(energy)))


def turing_mode():
    mu = np.linspace(0, 1.1, 400)
    lam = []
    for m in mu:
        A = np.array([[1-m,-2],[2,-3-10*m]])
        lam.append(np.max(np.linalg.eigvals(A).real))
    lam = np.asarray(lam)
    at = np.linalg.eigvals(np.array([[.7,-2],[2,-6]])).real
    assert max(at) > 0
    assert np.all(lam[(mu > .2) & (mu < .5)] > 0)
    fig = go.Figure(go.Scatter(x=mu, y=lam, name="max Re lambda(J-muD)"))
    fig.add_hline(y=0, line_dash="dash")
    fig.add_vrect(x0=.2, x1=.5, fillcolor="lightgray", opacity=.18,
                  line_width=0, annotation_text="niestabilne mody")
    fig.update_layout(title="Niestabilnosc Turinga: stabilna reakcja, niestabilna dyfuzja",
                      xaxis_title="mu (wartosc wasna -Delta)", yaxis_title="tempo wzrostu")
    save(fig, 'pde27_turing.html')
    return float(max(at))


if __name__ == '__main__':
    split_step_kpp()
    print('maks. residuum jawnej fali:', explicit_front())
    print('maks. przyrost energii:', allen_cahn_energy())
    print('dodatnia wartosc wlasna Turinga:', turing_mode())
    print('Zapisano 4 wykresy Plotly w:', OUT)

\end{lstlisting}


\section{Równania Naviera--Stokesa: wir Taylora--Greena i rzut Leraya — Plotly}
\label{sec:pde28-code}
Poniższy program sprawdza dokładną energię wiru Taylora--Greena,
przedstawia wirowość i ciśnienie oraz oblicza rzut Leraya
w przestrzeni Fouriera. W modelowym polu dywergencja po rzutowaniu
jest bliska zeru z dokładnością maszynową. To eksperyment numeryczny,
nie zamiennik dowodu istnienia rozwiązania Leraya--Hopfa.
\begin{lstlisting}[style=pythonpde,caption={Wir Taylora--Greena i rzut Leraya (Plotly)},label={code:pde28-plotly}]
"""Rownania Naviera-Stokesa: wir Taylora-Greena i rzut Leraya.

pip install numpy plotly
python demo_navier_stokes_plotly.py
Wyniki: pde28_wiry.html, pde28_energia.html,
        pde28_cisnienie.html, pde28_rzut_helmholtza.html
"""
from pathlib import Path
import numpy as np
import plotly.graph_objects as go

OUT = Path(__file__).resolve().parent
N, NU = 128, 0.08
x = np.linspace(0, 2*np.pi, N, endpoint=False)
X, Y = np.meshgrid(x, x, indexing='ij')
k = np.fft.fftfreq(N, d=1/N)  # Fourier integers, not physical cycles
KX, KY = np.meshgrid(k, k, indexing='ij')


def store(fig, name):
    fig.update_layout(template='plotly_white', font=dict(size=13))
    fig.write_html(OUT / name, include_plotlyjs=True, full_html=True)


def taylor_green(t):
    decay = np.exp(-2*NU*t)
    ux = decay*np.sin(X)*np.cos(Y)
    uy = -decay*np.cos(X)*np.sin(Y)
    p = 0.25*decay**2*(np.cos(2*X)+np.cos(2*Y))
    omega = 2*decay*np.sin(X)*np.sin(Y)
    return ux, uy, p, omega


def divergence_fft(vx, vy):
    return np.fft.ifft2(1j*KX*np.fft.fft2(vx)
                          +1j*KY*np.fft.fft2(vy)).real


def project(vx, vy):
    hx, hy = np.fft.fft2(vx), np.fft.fft2(vy)
    den = KX*KX+KY*KY
    safe = np.where(den == 0, 1, den)
    dot = KX*hx+KY*hy
    hx = hx-KX*dot/safe
    hy = hy-KY*dot/safe
    hx[0,0]=hy[0,0]=0
    return np.fft.ifft2(hx).real, np.fft.ifft2(hy).real


def plots():
    u1,u2,p,omega = taylor_green(0.5)
    fig = go.Figure(go.Contour(x=x,y=x,z=omega.T, colorscale='RdBu',
                 contours=dict(showlabels=True)))
    fig.update_layout(title='Wir Taylora-Greena: wirowosc w chwili t=0.5',
                      xaxis_title='x',yaxis_title='y', yaxis_scaleanchor='x')
    store(fig, 'pde28_wiry.html')

    times=np.linspace(0,10,150)
    e_exact=np.pi**2*np.exp(-4*NU*times)
    e_numeric=np.array([0.5*(taylor_green(t)[0]**2+
                           taylor_green(t)[1]**2).mean()*(2*np.pi)**2
                        for t in times])
    assert np.max(np.abs(e_exact-e_numeric))<2e-12
    fig=go.Figure()
    fig.add_trace(go.Scatter(x=times,y=e_numeric,name='Energia z siatki',mode='markers',
                             marker=dict(size=4)))
    fig.add_trace(go.Scatter(x=times,y=e_exact,name='Energia dokladna',mode='lines'))
    fig.update_layout(title='Tozsamosc energii: E(t)=pi^2 exp(-4 nu t)',
                      xaxis_title='czas t',yaxis_title='E(t)')
    store(fig,'pde28_energia.html')

    fig=go.Figure(go.Heatmap(x=x,y=x,z=p.T,colorscale='RdBu'))
    fig.update_layout(title='Cisnienie wiru Taylora-Greena w chwili t=0.5',
                      xaxis_title='x',yaxis_title='y',yaxis_scaleanchor='x')
    store(fig,'pde28_cisnienie.html')

    fx=np.sin(Y)+np.cos(X)
    fy=np.sin(X)+np.cos(Y)
    px,py=project(fx,fy)
    assert np.max(np.abs(px-np.sin(Y)))<2e-12
    assert np.max(np.abs(py-np.sin(X)))<2e-12
    div_before=np.max(np.abs(divergence_fft(fx,fy)))
    div_after=np.max(np.abs(divergence_fft(px,py)))
    assert div_before>1.9 and div_after<1e-10
    assert np.max(np.abs(divergence_fft(u1,u2)))<1e-10
    fig=go.Figure()
    step=9
    for vx,vy,name,dx in [(fx,fy,'F przed rzutem',0), (px,py,'PF po rzucie',2*np.pi+1)]:
        first=True
        for i in range(0,N,step):
            for j in range(0,N,step):
                xi,yi=x[i]+dx,x[j]
                fig.add_trace(go.Scatter(x=[xi,xi+0.27*vx[i,j],None],
                      y=[yi,yi+0.27*vy[i,j],None],mode='lines',
                      line=dict(width=1.1),name=name,showlegend=first))
                first=False
    fig.update_layout(title=f'Rzut Leraya: max div przed={div_before:.2f}, po={div_after:.1e}',
                      xaxis_title='dwa pola obok siebie',yaxis_title='y',
                      yaxis_scaleanchor='x')
    store(fig,'pde28_rzut_helmholtza.html')
    print('Energia: maksymalny blad', np.max(np.abs(e_exact-e_numeric)))
    print('Dywergencja przed i po rzucie',div_before,div_after)


if __name__ == '__main__':
    plots()
\end{lstlisting}


\section{Zagadnienie przeszkody i reakcja kontaktowa — Plotly}
\label{sec:pde29-code}
Program implementuje metodę Gaussa--Seidela z projekcją
na zbiór funkcji spełniających przeszkodę. Porównuje wyniki
z dokładnym rozwiązaniem z przykładu~\ref{ex:pde29-explicit}
i kontroluje nieujemność reakcji oraz warunek komplementarności.
Listing jest obramowany; samodzielny plik \texttt{demo\_przeszkoda\_plotly.py}
znajduje się również w pakiecie projektu.
\begin{lstlisting}[style=pythonpde,caption={Przeszkoda, reakcja i zbieżność schematu (Plotly)},label={code:pde29-plotly}]
"""Rozdzial 29 PDE: metoda Gaussa-Seidela z projekcja dla przeszkody.

Wymagania: pip install numpy plotly
Wykonanie: python demo_przeszkoda_plotly.py
Wyniki: pde29_przeszkoda.html, pde29_reakcja.html,
        pde29_zbieznosc.html
"""
from pathlib import Path
import numpy as np
import plotly.graph_objects as go

OUT = Path(__file__).resolve().parent


def obstacle(x: np.ndarray) -> np.ndarray:
    return (1.0 - x * x) ** 2


def exact(x: np.ndarray) -> np.ndarray:
    a = np.abs(x)
    return np.where(a <= 1 / 3, (1.0 - x * x) ** 2,
                    32 / 27 * (1.0 - a))


def projected_gauss_seidel(n: int, omega=1.75, tol=2e-12,
                            max_sweeps=18000):
    h = 2.0 / n
    x = np.linspace(-1.0, 1.0, n + 1)
    psi = obstacle(x)
    u = psi.copy()
    u[0] = u[-1] = 0
    history = []
    for k in range(max_sweeps):
        change = 0.0
        for i in range(1, n):
            new = max(psi[i], u[i] + omega *
                      (0.5 * (u[i - 1] + u[i + 1]) - u[i]))
            change = max(change, abs(new - u[i]))
            u[i] = new
        if k % 25 == 0:
            energy = np.sum(np.diff(u) ** 2) / (2 * h)
            history.append((k, energy, change))
        if change < tol:
            break
    else:
        raise RuntimeError(f'Projected GS did not converge for n={n}')
    reaction = (2 * u[1:-1] - u[:-2] - u[2:]) / (h * h)
    slack = u[1:-1] - psi[1:-1]
    return x, u, psi, reaction, slack, np.array(history), k + 1


def save(fig, name):
    fig.update_layout(template='plotly_white', font=dict(size=13))
    fig.write_html(OUT / name, include_plotlyjs=True, full_html=True)


def main():
    ns = [40, 80, 160]
    results = [projected_gauss_seidel(n) for n in ns]
    errors = []
    for n, (x, u, psi, reaction, slack, hist, sweeps) in zip(ns, results):
        err = np.max(np.abs(u - exact(x)))
        errors.append(err)
        assert np.min(slack) > -1e-11, (n, np.min(slack))
        assert np.min(reaction) > -2e-7, (n, np.min(reaction))
        assert np.max(np.abs(reaction * slack)) < 5e-7
        print(f'N={n}, sweeps={sweeps}, max error={err:.4e}, '
              f'min reaction={reaction.min():.2e}, '
              f'max complementarity={np.max(np.abs(reaction * slack)):.2e}')

    x, u, psi, reaction, slack, hist, _ = results[-1]
    xf = np.linspace(-1, 1, 901)
    fig = go.Figure()
    fig.add_trace(go.Scatter(x=xf, y=obstacle(xf),
                             name='Przeszkoda psi', mode='lines'))
    fig.add_trace(go.Scatter(x=xf, y=exact(xf),
                             name='Rozwi\u0105zanie dok\u0142adne', mode='lines'))
    fig.add_trace(go.Scatter(x=x, y=u,
                             name='Wynik numeryczny', mode='markers',
                             marker=dict(size=5)))
    for point in (-1 / 3, 1 / 3):
        fig.add_vline(x=point, line_dash='dash', line_color='gray')
    fig.update_layout(title='Zagadnienie przeszkody: kontakt i swobodny brzeg',
                      xaxis_title='x', yaxis_title='Wysoko\u015b\u0107')
    save(fig, 'pde29_przeszkoda.html')

    xnodes = x[1:-1]
    reaction_exact = np.where(np.abs(xnodes) < 1 / 3,
                              4 - 12 * xnodes * xnodes, 0.0)
    fig = go.Figure()
    fig.add_trace(go.Scatter(x=xnodes, y=reaction,
                             mode='lines+markers',
                             name='Reakcja dyskretna',
                             marker=dict(size=4)))
    fig.add_trace(go.Scatter(x=xnodes, y=reaction_exact,
                             name='Reakcja dok\u0142adna', mode='lines'))
    fig.update_layout(title='Reakcja przeszkody: -u" >= 0',
                      xaxis_title='x', yaxis_title='Mno\u017cnik ograniczenia')
    save(fig, 'pde29_reakcja.html')

    fig = go.Figure()
    fig.add_trace(go.Scatter(x=np.array(ns), y=errors, mode='lines+markers',
                             name='Maksymalny b\u0142\u0105d w\u0119z\u0142owy'))
    fig.update_layout(title='Zbie\u017cno\u015b\u0107 metody z projekcj\u0105',
                      xaxis_type='log', yaxis_type='log',
                      xaxis_title='Liczba przedzia\u0142\xf3w N',
                      yaxis_title='B\u0142\u0105d maksimum')
    save(fig, 'pde29_zbieznosc.html')
    if all(err > 0 for err in errors):
        orders = np.log(np.array(errors[:-1]) / np.array(errors[1:])) / np.log(2)
        print('Obserwowane rz\u0119dy zbie\u017cno\u015bci:', orders)


if __name__ == '__main__':
    main()
\end{lstlisting}

\section{Homogenizacja periodyczna: rozwiązanie i strumień — Plotly}
\label{sec:pde30-code}
Program porównuje dokładne rozwiązania jednowymiarowe z równaniem
uśrednionym, pokazuje oscylacje gradientu oraz zbieżność błędu.
Ostatni wykres przedstawia anizotropię macierzy efektywnej
materiału dwuwarstwowego; jest ilustracją prawa
~\eqref{eq:pde30-layers}, a nie obliczeniem metodą elementów skończonych.
\begin{lstlisting}[style=pythonpde,caption={Homogenizacja okresowa i macierz efektywna (Plotly)},label={code:pde30-plotly}]
"""PDE, rozdzial 30: 1D homogenizacja okresowa i srednie efektywne.

Wymagania: pip install numpy plotly
Uruchomienie: python demo_homogenizacja_plotly.py
Pliki wynikowe: pde30_profile.html, pde30_gradient_strumien.html,
                pde30_zbieznosc.html, pde30_anizotropia.html
"""
from pathlib import Path
import numpy as np
import plotly.graph_objects as go

DIR = Path(__file__).resolve().parent
A_HOM = np.sqrt(3.0)


def cumulative_trapezoid(v, h):
    return np.concatenate(([0.0], np.cumsum((v[1:] + v[:-1]) * (h / 2.0))))


def micro_solution(periods, samples_per_period=240):
    eps = 1.0 / periods
    x = np.linspace(0.0, 1.0, periods * samples_per_period + 1)
    h = x[1] - x[0]
    a = 2.0 + np.sin(2.0 * np.pi * x / eps)
    b = 1.0 / a
    ib = cumulative_trapezoid(b, h)
    ibx = cumulative_trapezoid(b * x, h)
    c = ibx[-1] / ib[-1]
    u = c * ib - ibx
    gradient = b * (c - x)
    flux = c - x
    u0 = x * (1.0 - x) / (2.0 * A_HOM)
    grad0 = (0.5 - x) / A_HOM
    err = np.sqrt(np.trapezoid((u - u0) ** 2, x))
    return x, u, gradient, flux, u0, grad0, err


def export(fig, name):
    fig.update_layout(template="plotly_white", font=dict(size=13))
    fig.write_html(DIR / name, full_html=True, include_plotlyjs=True)


def main():
    first = go.Figure()
    for n in [4, 8, 16, 32]:
        x, u, du, q, u0, du0, err = micro_solution(n)
        skip = max(1, len(x) // 2000)
        first.add_trace(go.Scatter(x=x[::skip], y=u[::skip], name=f"eps=1/{n}", mode="lines"))
    first.add_trace(go.Scatter(x=x, y=u0, name="rozwiazanie graniczne", line=dict(color="black", dash="dash")))
    first.update_layout(title="Homogenizacja: rozwiazania mikroskopowe i graniczne",
                        xaxis_title="x", yaxis_title="u(x)")
    export(first, "pde30_profile.html")

    x, u, grad, flux, u0, grad0, err = micro_solution(24)
    fig = go.Figure()
    fig.add_trace(go.Scatter(x=x, y=grad, name="gradient u'eps"))
    fig.add_trace(go.Scatter(x=x, y=grad0, name="gradient u'0", line=dict(dash="dot")))
    fig.add_trace(go.Scatter(x=x, y=flux, name="strumien aepsu'eps", line=dict(width=3)))
    fig.add_trace(go.Scatter(x=x, y=0.5-x, name="strumien graniczny", line=dict(dash="dash")))
    fig.update_layout(title="Silna zbieznosc strumienia, oscylacje gradientu",
                      xaxis_title="x", yaxis_title="wartosc")
    export(fig, "pde30_gradient_strumien.html")

    ns = np.array([4, 8, 16, 32, 64], dtype=int)
    errors = np.array([micro_solution(int(n))[6] for n in ns])
    order = float(np.polyfit(np.log(1.0/ns), np.log(errors), 1)[0])
    fig = go.Figure()
    fig.add_trace(go.Scatter(x=1.0/ns, y=errors, mode="lines+markers", name="blad L^2"))
    fig.add_trace(go.Scatter(x=1.0/ns, y=errors[0]*(ns[0]/ns),
                             mode="lines", name="wzorzec O(eps)", line=dict(dash="dash")))
    fig.update_layout(title=f"Zbieznosc do rozwiazania granicznego; rzad ~{order:.3f}",
                      xaxis_type="log", yaxis_type="log",
                      xaxis_title="eps", yaxis_title="||ueps-u0|| L^2")
    export(fig, "pde30_zbieznosc.html")

    t = np.linspace(0, 2*np.pi, 360)
    fig = go.Figure()
    fig.add_trace(go.Scatter(x=np.sqrt(8.0/5.0)*np.cos(t),
                             y=np.sqrt(5.0/2.0)*np.sin(t),
                             name="warstwy: diag(8/5,5/2)"))
    fig.add_trace(go.Scatter(x=np.sqrt(2.5)*np.cos(t),
                             y=np.sqrt(2.5)*np.sin(t),
                             name="bledna srednia arytmetyczna", line=dict(dash="dash")))
    fig.update_yaxes(scaleanchor="x", scaleratio=1)
    fig.update_layout(title="Izotropowe fazy -> anizotropowa macierz efektywna",
                      xaxis_title="kierunek prostopadly do warstw (y1)",
                      yaxis_title="kierunek wzdluz warstw (y2)")
    export(fig, "pde30_anizotropia.html")
    assert abs(A_HOM - 1.7320508075688772) < 1e-12
    assert all(np.diff(errors) < 0), errors
    assert 0.8 < order < 1.2, order
    assert abs(flux[-1] - flux[0] + 1.0) < 1e-12
    assert abs(u[-1]) < 5e-12
    print("Efektywny wspolczynnik:", A_HOM)
    print("Rzad zbieznosci L2:", order)
    print("Bledy L2:", errors)


if __name__ == "__main__":
    main()
\end{lstlisting}

\section{Analiza mikrolokalna i symbole: eksperymenty Plotly}
\label{sec:pde31-code}
Rozkład częstości dla wąskiej warstwy, stożek charakterystyczny i parametryks.
\begin{lstlisting}[style=pythonpde,caption={Eksperymenty rozdziału 31 (Python i Plotly)},label={code:pde31-plotly}]
"""Rozdzial 31: kierunki osobliwosci i symbole Fourierowskie.
Wymaga numpy, plotly; zapisuje trzy interaktywne wykresy HTML.
Przyblizenie delta(x1) jest gladkie dla skonczonej szerokosci;
wykres ilustruje, a nie dowodzi twierdzenia o zbiorze falowym.
"""
from pathlib import Path
import numpy as np
import plotly.graph_objects as go

OUT = Path(__file__).resolve().parent


def save(fig, filename):
    fig.update_layout(template="plotly_white", font=dict(size=13))
    fig.write_html(OUT / filename, include_plotlyjs=True, full_html=True)


def directional_singularity():
    n = 256
    x = np.linspace(-np.pi, np.pi, n, endpoint=False)
    dx = x[1] - x[0]
    xx, yy = np.meshgrid(x, x, indexing="ij")
    angle = np.pi / 5
    normal = xx * np.cos(angle) + yy * np.sin(angle)
    # Okresowy profil waskiego pasa; zwezone przyblizenie delty.
    profile = np.exp(-0.5 * (np.sin(normal / 2) / 0.075) ** 2)
    window = np.exp(-0.5 * ((xx / 2.2) ** 8 + (yy / 2.2) ** 8))
    img = profile * window
    freq = np.fft.fftshift(np.fft.fftfreq(n, d=dx)) * 2 * np.pi
    amp = np.log1p(np.abs(np.fft.fftshift(np.fft.fft2(img))))
    fig = go.Figure(go.Heatmap(x=freq, y=freq, z=amp.T, colorscale="Viridis"))
    fig.update_layout(title="Kierunki wysokich czestotliwosci: waski pas",
                      xaxis_title="xi_1", yaxis_title="xi_2")
    fig.update_xaxes(range=[-45, 45]); fig.update_yaxes(range=[-45, 45])
    save(fig, "pde31_kierunki_fouriera.html")
    assert np.isfinite(amp).all()


def characteristic_cone():
    xi = np.linspace(-5, 5, 801)
    fig = go.Figure()
    fig.add_trace(go.Scatter(x=xi, y=2 * np.abs(xi), name="tau = +2|xi|"))
    fig.add_trace(go.Scatter(x=xi, y=-2 * np.abs(xi), name="tau = -2|xi|"))
    fig.update_layout(title="Zbior charakterystyczny rownania falowego, c=2",
                      xaxis_title="xi", yaxis_title="tau")
    save(fig, "pde31_charakterystyki.html")
    assert np.max(np.abs((2 * np.abs(xi)) ** 2 - 4 * xi ** 2)) < 1e-10


def parametrix_filter():
    xi = np.linspace(-20, 20, 1201)
    p = 1 + xi ** 2
    inverse = 1 / p
    fig = go.Figure()
    fig.add_trace(go.Scatter(x=xi, y=p, name="P(i xi)=1+xi^2"))
    fig.add_trace(go.Scatter(x=xi, y=inverse, name="q(xi)=1/(1+xi^2)"))
    fig.update_layout(title="Symbol eliptyczny i odwrotny mnoznik",
                      xaxis_title="xi", yaxis_title="wartosc")
    save(fig, "pde31_parametryks.html")
    assert np.allclose(p * inverse, 1)


if __name__ == "__main__":
    directional_singularity()
    characteristic_cone()
    parametrix_filter()
    print("Mikrolokalna: zapisano 3 wykresy Plotly.")
\end{lstlisting}


\section{Problemy odwrotne i sterowanie optymalne: eksperymenty Plotly}
\label{sec:pde32-code}
Zanik wartości singularnych, regularyzacja Tichonowa i sterowanie.
\begin{lstlisting}[style=pythonpde,caption={Eksperymenty rozdziału 32 (Python i Plotly)},label={code:pde32-plotly}]
"""Rozdzial 32: niestabilnosc, regularyzacja Tichonowa i sterowanie PDE.
Wymaga numpy, plotly. Zawiera testy formul dla modelu spektralnego.
"""
from pathlib import Path
import numpy as np
import plotly.graph_objects as go

OUT = Path(__file__).resolve().parent


def save(fig, filename):
    fig.update_layout(template="plotly_white", font=dict(size=13))
    fig.write_html(OUT / filename, include_plotlyjs=True, full_html=True)


def singular_values():
    n = np.arange(1, 41)
    sig_ell = 1 / (1 + n ** 2)
    sig_heat = np.exp(-0.01 * n ** 2)
    fig = go.Figure()
    fig.add_trace(go.Scatter(x=n, y=sig_ell, mode="lines+markers", name="(-Delta+I)^(-1)"))
    fig.add_trace(go.Scatter(x=n, y=sig_heat, mode="lines+markers", name="exp(T Delta), T=0.01"))
    fig.update_layout(title="Male wartosci singularne dwoch operatorow PDE",
                      xaxis_title="numer modu", yaxis_title="wartosc singularna",
                      yaxis_type="log")
    save(fig, "pde32_wartosci_singularne.html")


def inverse_heat():
    rng = np.random.default_rng(713)
    n = np.arange(1, 51)
    T = 0.007
    sigma = np.exp(-T * n ** 2)
    truth = np.exp(-n / 7) * np.sin(n / 3)
    delta = 0.002
    noise = rng.normal(0, 1, n.size)
    noise = delta * noise / np.linalg.norm(noise)
    data = sigma * truth + noise
    alpha = np.logspace(-10, -2, 130)
    recon = sigma[:, None] * data[:, None] / (sigma[:, None] ** 2 + alpha[None, :])
    error = np.linalg.norm(recon - truth[:, None], axis=0)
    idx = int(np.argmin(error)); alpha_best = alpha[idx]
    fig = go.Figure()
    fig.add_trace(go.Scatter(x=alpha, y=error, name="blad rekonstrukcji"))
    fig.add_vline(x=alpha_best, line_dash="dash",
                  annotation_text="minimum dla danych testowych")
    fig.update_layout(title="Regularyzacja Tichonowa: kompromis szum-gladzenie",
                      xaxis_title="alpha", xaxis_type="log",
                      yaxis_title="blad wspolczynnikow")
    save(fig, "pde32_regularyzacja.html")
    fig2 = go.Figure()
    fig2.add_trace(go.Scatter(x=n, y=truth, name="sygnal prawdziwy"))
    fig2.add_trace(go.Scatter(x=n, y=recon[:, idx], name="rekonstrukcja"))
    fig2.update_layout(title=f"Odtwarzanie modow, alpha={alpha_best:.2e}",
                       xaxis_title="numer modu", yaxis_title="wspolczynnik")
    save(fig2, "pde32_rekonstrukcja.html")
    assert error[idx] <= error[0] and error[idx] <= error[-1]
    return float(alpha_best), float(error[idx])


def elliptic_control():
    x = np.linspace(0, np.pi, 601)
    alpha = 0.15
    n = np.array([1.0, 3.0])
    desired_coeff = np.array([1.0, 0.35])
    sigma = 1 / (n ** 2)
    control_coeff = sigma * desired_coeff / (sigma ** 2 + alpha)
    observed_coeff = sigma * control_coeff
    sin_modes = np.sin(n[:, None] * x)
    y_desired = desired_coeff @ sin_modes
    y_state = observed_coeff @ sin_modes
    control = control_coeff @ sin_modes
    fig = go.Figure()
    for values, name in ((y_desired, "stan pozadany"),
                         (y_state, "stan optymalny"), (control, "sterowanie")):
        fig.add_trace(go.Scatter(x=x, y=values, name=name))
    fig.update_layout(title="Sterowanie rozlozone dla -y''=u",
                      xaxis_title="x", yaxis_title="wartosc")
    save(fig, "pde32_sterowanie.html")
    assert np.allclose((sigma ** 2 + alpha) * control_coeff,
                       sigma * desired_coeff)


if __name__ == "__main__":
    singular_values()
    a, err = inverse_heat()
    elliptic_control()
    print(f"Optymalne testowe alpha={a:.3e}, blad={err:.4f}; zapisano 4 wykresy.")
\end{lstlisting}


\end{document}
